\documentclass[a4paper, reqno
]{amsart}

\usepackage[utf8]{inputenc}
\usepackage[english]{babel}
\usepackage{amsmath}
\usepackage{amssymb}
\usepackage{graphicx}
\usepackage{xcolor}
\usepackage{caption}
\usepackage{tikz}
\usetikzlibrary{arrows.meta,calc}
\usepackage[all]{xy}
\usepackage{hyperref}
\usepackage[notcite, notref]{showkeys}

% Macros

\makeatletter
\newcommand*\bigcdot{\mathpalette\bigcdot@{.5}}
\newcommand*\bigcdot@[2]{\mathbin{\vcenter{\hbox{\scalebox{#2}{$\m@th#1\bullet$}}}}}
\makeatother

\newcommand{\ol}{\overline}
\newcommand{\pr}{\mathrm{pr}}
\DeclareMathOperator{\Hom}{Hom}
\DeclareMathOperator{\Fnc}{Fnc}
\DeclareMathOperator{\Rel}{Rel}
\DeclareMathOperator{\G}{G}
\DeclareMathOperator{\bconcat}{\curlywedge}
\newcommand{\concat}{\ensuremath{\text{\Large $\curlywedge$}}}
\newcommand{\bb}[1]{\ensuremath{\lvert #1 \rvert}}
\newcommand{\vs}[1]{\mathbin{\downarrow}#1}

% Signatures, algebras and equations
\newcommand{\sig}[1]{\mathbf{#1}}             % signature: \sig{\Sigma}
\newcommand{\calg}[1]{\boldsymbol{\mathsf{#1}}} % categorial algebra: \calg{A}
\newcommand{\eeq}{\overset{\mathrm e}{=}}      % existence equation
\newcommand{\QE}{\mathrm {QE}}
\newcommand{\Dom}{\mathrm {Dom}}
\newcommand{\id}{\mathrm {id}}
\newcommand{\Id}{\mathrm {Id}}
\newcommand{\gd}{\mathfrak {gd}}              % generalized hep-derivors
\newcommand{\gl}{\mathrm {gl}}                % globular

% Source and target of cells
\newcommand{\src}{\mathrm {sc}}
\newcommand{\tgt}{\mathrm {tg}}

% Graduation, unification, discretization, truncation: functors and derivors
\newcommand{\Gd}{\mathrm {Gd}}
\newcommand{\Uf}{\mathrm {Uf}}
\newcommand{\Ds}{\mathrm {Ds}}
\newcommand{\Tr}{\mathrm {Tr}}
\newcommand{\bGd}{\mathbf {Gd}}
\newcommand{\bUf}{\mathbf {Uf}}
\newcommand{\bDs}{\mathbf {Ds}}
\newcommand{\bTr}{\mathbf {Tr}}

% Arrows (xy-pic)

\NoCompileMatrices
\def\labelstyle{\textstyle}
\xymatrixrowsep = {8ex}
\xymatrixcolsep = {10ex}

\newdir{+>}{@{|}*@{-}!/-6.5pt/@{>}*!/-13pt/{}}
\newdir{ +}{{}*!/-5pt/@{+}}

\newsavebox{\xymor}
\newsavebox{\xydmor}

\newcommand{\xynode}{\makebox[0ex]{}}
\savebox{\xymor}{\ensuremath{%
		\xymatrix@1@C=19pt{\xynode \ar@{>}[r] & \xynode }}}
\savebox{\xydmor}{\ensuremath{%
		\xymatrix@1@C=19pt{\xynode \ar@{{}{-}_{>}}[r] & \xynode }}}

\newcommand{\mor}{\usebox{\xymor}}    % mapping         --->
\newcommand{\dmor}{\usebox{\xydmor}}  % partial mapping ---,

\newcommand{\function}[4]{
	\begin{array}{@{\:}c@{\:}c@{\:}l}
		#1 &\mor& #2 \\
		#3 &\longmapsto& #4
\end{array} }
\newcommand{\nfunction}[4]
{\left\{
	\function{#1}{#2}{#3}{#4}
	\right. }

% Theorem-like environments

\theoremstyle{plain}
\newtheorem{theorem}{Theorem}[section]
\newtheorem{proposition}[theorem]{Proposition}
\newtheorem{corollary}[theorem]{Corollary}
\newtheorem{lemma}[theorem]{Lemma}
\newtheorem{claim}[theorem]{Claim}
\newtheorem{remark}[theorem]{Remark}

\theoremstyle{definition}
\newtheorem{definition}[theorem]{Definition}
\newtheorem*{convention}{Convention}

\begin{document}

\title{Higher-order categorial algebras}

\author[Climent]{J. Climent Vidal}
\address{Universitat de Val\`{e}ncia\\
         Departament de L\`{o}gica i Filosofia de la Ci\`{e}ncia
         }
\email{Juan.B.Climent@uv.es}

\author[Cosme]{E. Cosme Ll\'{o}pez}
\address{Universitat de Val\`{e}ncia\\
         Departament de Matem\`{a}tiques\\ and 
         Nantong University\\
         School of Mathematics and Statistics
         }
\email{Enric.Cosme@uv.es}

\author[Ruiz]{R. Ruiz Mora}
\address{Universitat de Val\`{e}ncia\\
         Departament de Matem\`{a}tiques
         }
\email{Raul.Ruiz-Mora@uv.es}

\author[Vago]{M. Vago Marzal}

\begin{abstract}
We introduce generalized hep-derivors, morphisms between hierarchical equationally partial signatures that send every sort to a set of existence equations in one variable and every operation symbol to a term, and we prove that every generalized hep-derivor induces a functor between the corresponding categories of many-sorted partial algebras. The hierarchical equationally partial signatures and the generalized hep-derivors form a category over which the categories of weak ep-algebras constitute a normal lax indexed category, and its Grothendieck construction yields the category of the ep-algebras of all such signatures. We apply these notions to higher-order categorial algebras, i.e., many-sorted algebras whose underlying many-sorted sets are strict higher-order categories and whose operations are higher-order functors. We present the single-set and the globular $n$-categorial algebras, for $n\in\mathbb{N}\cup\{\omega\}$, as ep-varieties, and we prove that, in the single-set case, the condition that every cell has a finite dimension cannot be expressed by quasi-existence equations. By means of generalized hep-derivors we define the graduation, unification, discretization and truncation functors, we prove that graduation and unification are mutually quasi-inverse equivalences and that discretization is left adjoint to truncation, we show that the categories of $\omega$-categorial algebras are projective limits of the categories of $n$-categorial algebras, and we combine all these functors in algebraic squares.
\end{abstract}

\maketitle

\section{Introduction}\label{SIntroduction}

Higher-order categorial algebras are many-sorted algebras whose underlying many-sorted sets carry a structure of strict higher-order category and whose operations are higher-order functors. For a set of sorts $S$, an $S$-sorted signature $\Sigma$ and $n\in\mathbb{N}$, the category $\mathsf{nCatAlg}(\sig{\Sigma})$ of the $n$-categorial $\sig{\Sigma}$-algebras is such that the forgetful functor from it to $\mathsf{Set}^{S}$ factorizes through the category of $S$-sorted strict $n$-categories; for $S=1$ and the empty signature, it is the category $\mathsf{nCat}$ of strict $n$-categories. These algebras arise naturally in the study of higher-order many-sorted rewriting systems, since the paths of a rewriting system of order $n$, up to suitable congruences, form an $n$-categorial algebra, and the higher-order Curry-Howard isomorphisms of~\cite{CVCLRM26} are isomorphisms between such algebras.

As for strict higher-order categories, there are two natural ways of presenting the higher-order categorial algebras. In the \emph{single-set} presentation, which goes back to~\cite{Str87}, all the cells of a given sort, whatever their dimension, form a single set, and the dimension of a cell is read off from its sources and targets; in the \emph{globular} presentation, see, e.g.,~\cite{abgmmm23}, the cells of every dimension form a separate sort, and there are sources, targets and identities between dimensions. In both cases the composition of cells is only partially defined, so that the higher-order categorial algebras are partial algebras, and both presentations can be described by $\QE$-equations in the sense of Burmeister~\cite{pb86, pb02}. Moreover, the domains of the compositions are given by existence equations, so that the corresponding signatures are equationally partial signatures in the sense of Reichel~\cite{Rei87}.

The main aim of this work is to relate, by means of functors induced by morphisms between signatures, the different presentations and the different dimensions of the higher-order categorial algebras. To this end, we introduce the notion of generalized hep-derivor, which combines two ideas. On the one hand, the derivors of~\cite{CVST10b}, which send every operation symbol of a signature to a term of another, induce functors between the categories of algebras in the opposite direction. On the other hand, the general derivors of~\cite{LG18} send, moreover, every sort to a formula that delimits, inside the interpretation of its image, the elements that interpret it. A generalized hep-derivor between hierarchical equationally partial signatures sends every sort to a set of existence equations in one variable and every operation symbol to a term; the domains of the transformed operations are then determined by the domain conditions of the operation symbols, and the hierarchical character of the signatures allows us to define the transformed algebras by well-founded recursion. As is customary in the work of our group, see~\cite{CVCL16, CVCL18, CVCL23, CVCLRM26}, we have tried to make the presentation as complete and detailed as possible.

\subsection{Our work}

In Section~\ref{SPreliminaries} we collect the notions on many-sorted sets, free monoids, many-sorted algebras, terms and variables, Hall algebras and derivors, many-sorted partial algebras, existence and quasi-existence equations, adjunctions and the Grothendieck construction that we use in the rest of this work. In addition to the Grothendieck construction for split indexed categories, we state the Grothendieck construction for normal lax indexed categories, in which composites are only preserved up to a coherent family of natural transformations, and we prove that it gives rise to a category.

In Section~\ref{SGenHepDerivors} we introduce the ep-signatures, the ep-algebras and their varieties, the hierarchical ep-signatures, or hep-signatures, and the generalized hep-derivors between them. We prove that every generalized hep-derivor induces a functor between the categories of partial algebras, and we characterize when the transformed algebra of a partial algebra is an ep-algebra. Then we prove that the hep-signatures and the generalized hep-derivors form a category $\mathsf{hepSig}_{\gd}$ and we compare the functor induced by a composite with the composite of the induced functors: in general, the latter is only contained in the former. Consequently, the categories of weak ep-algebras, together with the induced functors and these inclusions, constitute a normal lax indexed category over $\mathsf{hepSig}_{\gd}$, and, by means of the Grothendieck construction, we obtain the category $\mathsf{hepAlg}_{\gd}$ of the ep-algebras of all the hep-signatures, which generalizes the category $\mathsf{Alg}_{\mathfrak{d}}$ of~\cite{CVST10b}.

In Sections~\ref{SSingleSet} and~\ref{SGlobular} we recall the single-set and the globular higher-order categorial algebras, and we present them as the categories canonically associated to the ep-varieties of suitable hep-specifications. In dimension $\omega$, the globular $\omega$-categorial algebras also form such a category, whereas, in the single-set presentation, the condition that every cell has a finite dimension cannot be expressed by $\QE$-equations, since the class of the single-set $\omega$-categorial algebras is not closed under products.

In Section~\ref{SGraduation} we define, by means of generalized hep-derivors, the graduation functor, which distributes the cells of a single-set $n$-categorial algebra into sorts according to their dimension, and the unification functor, which keeps only the cells of top dimension of a globular $n$-categorial algebra. We prove that they are mutually quasi-inverse equivalences of categories and, more precisely, that graduation is left adjoint to unification in an adjoint equivalence whose unit is the identity. In Section~\ref{STruncation} we define, in both presentations, the discretization functors, which regard an $m$-categorial algebra as an $n$-categorial algebra whose cells of dimension greater than $m$ are identities, and the truncation functors, which forget the cells of dimension greater than $m$. We prove that discretization is left adjoint to truncation, that these functors compose as expected and that they are compatible with graduation and unification. We end this section by defining the truncation and discretization functors in dimension $\omega$.

In Sections~\ref{SLimitSingleSet} and~\ref{SLimitGlobular} we prove that the categories of single-set and of globular $\omega$-categorial algebras are the projective limits of the projective systems formed by the categories of $n$-categorial algebras and the truncation functors. In Section~\ref{SOmegaFunctors} we obtain the graduation functor in dimension $\omega$ by the universal property of the projective limit and the unification functor in dimension $\omega$ as an inductive limit, since, as we show, it cannot be induced by a generalized hep-derivor; and we prove that, also in dimension $\omega$, graduation and unification are mutually quasi-inverse equivalences and discretization is left adjoint to truncation. Finally, in Section~\ref{SAlgSquare} we show that graduation, unification, discretization and truncation fit together in algebraic squares, both between two finite dimensions and between a finite dimension and $\omega$.

\subsection{Prerequisites, notation and terminology}

The only prerequisites for reading this work are familiarity with category theory, universal algebra, in particular many-sorted partial algebras, the theory of ordered sets and set theory. This can be achieved by reading, for category theory, e.g., \cite{hs73,sM98}, for universal algebra, e.g., \cite{gb15,pb82,pb86,pb02,bs81,gtw85,gG08,agk67,sch66,sch70,w92}, for the theory of ordered sets, e.g., \cite{bl05,DP02,R08}, and, for set theory, e.g., \cite{nB70,end77}.

\textbf{Set theory.} Our underlying set theory is $\mathbf{ZFSk}$, Zermelo-Fraenkel-Skolem set theory (also known as $\mathbf{ZFC}$, i.e., Zermelo-Fraenkel set theory with the axiom of choice) plus the existence of a Grothendieck universe $\boldsymbol{\mathcal{U}}$, fixed once and for all, see~\cite{sM98}, pp.~21--24. The elements of $\boldsymbol{\mathcal{U}}$ are called $\boldsymbol{\mathcal{U}}$-small sets and the subsets of $\boldsymbol{\mathcal{U}}$ are called $\boldsymbol{\mathcal{U}}$-large sets or classes.

\textbf{Ordinals.} An \emph{ordinal} $\alpha$ is a transitive set that is well-ordered by $\in$, thus $\alpha=\{\,\beta\mid\beta\in\alpha\,\}$. The first transfinite ordinal $\omega_{0}$ is denoted by $\mathbb{N}$, or by $\omega$, and it is the set of all \emph{natural numbers}; thus, for every $n\in\mathbb{N}$, $n=\{0,\ldots,n-1\}$.

\textbf{Functions and mappings.} For sets $A$ and $B$, a \emph{function} from $A$ to $B$ is a subset $F$ of $A\times B$ such that, for every $x\in A$, there exists a unique $y\in B$ such that $(x,y)\in F$, and a \emph{partial function} from $A$ to $B$ is a subset $F$ of $A\times B$ such that, for every $x\in A$, there exists at most one $y\in B$ such that $(x,y)\in F$. We denote by $\mathrm{Fnc}(A,B)$ and by $\mathrm{Pfnc}(A,B)$ the sets of all functions and of all partial functions from $A$ to $B$, respectively. The \emph{domain (of definition)} of a partial function $F$ is $\Dom(F)=\{x\in A\mid\exists\,y\in B\,((x,y)\in F)\}$; a partial function is usually denoted by $(F_{x})_{x\in\Dom(F)}$ and a function by $(F_{x})_{x\in A}$. A \emph{mapping} from $A$ to $B$ is an ordered triple $f=(A,F,B)$, denoted by $f\colon A\mor B$, in which $F$ is a function from $A$ to $B$, and a \emph{partial mapping} from $A$ to $B$ is an ordered triple $f=(A,F,B)$, denoted by $f\colon A\dmor B$, in which $F$ is a partial function from $A$ to $B$. We denote by $\mathrm{Hom}(A,B)$, or by $B^{A}$, and by $\mathrm{Hom}_{\mathrm{p}}(A,B)$ the sets of all mappings and of all partial mappings from $A$ to $B$, respectively. Thus $\mathrm{Hom}(A,B)=\{A\}\times\mathrm{Fnc}(A,B)\times\{B\}$ and $\mathrm{Hom}_{\mathrm{p}}(A,B)=\{A\}\times\mathrm{Pfnc}(A,B)\times\{B\}$. For a partial mapping $f=(A,F,B)$, we let $\Dom(f)$ stand for $\Dom(F)$; let us note that $\Dom(f)\subseteq A$, and that $\Dom(f)=A$ if $f$ is a mapping. For partial mappings $f\colon A\dmor B$ and $g\colon B\dmor C$, the composite $g\circ f$ is the partial mapping from $A$ to $C$ with $\Dom(g\circ f)=f^{-1}[\Dom(g)]$ that sends $x$ to $g(f(x))$.

\textbf{Subsets and images.} For a set $A$, we denote by $\mathrm{Sub}(A)$ the set of all subsets of $A$ and, for $X\in\mathrm{Sub}(A)$, by $\complement_{A}X$, or $A-X$, the complement of $X$ in $A$. For a mapping $f\colon A\mor B$, the mappings $f[\cdot]\colon\mathrm{Sub}(A)\mor\mathrm{Sub}(B)$, of $f$-\emph{direct image formation}, and $f^{-1}[\cdot]\colon\mathrm{Sub}(B)\mor\mathrm{Sub}(A)$, of $f$-\emph{inverse image formation}, send $X\subseteq A$ to $f[X]=\{f(x)\mid x\in X\}$ and $Y\subseteq B$ to $f^{-1}[Y]=\{x\in A\mid f(x)\in Y\}$, respectively. We write $\mathrm{Im}(f)$ for the image $f[A]$ of $f$, $\mathrm{Ker}(f)$ for its kernel and, for $X\subseteq A$, $f{\upharpoonright}_{X}$ for the restriction of $f$ to $X$.

\textbf{Categories.} We let $\mathsf{Set}$ stand for the category whose set of objects is $\boldsymbol{\mathcal{U}}$ and whose morphisms are the mappings between $\boldsymbol{\mathcal{U}}$-small sets, and $\mathsf{Cat}$ for the category of all $\boldsymbol{\mathcal{U}}$-categories, i.e., the category whose objects are the categories $\mathsf{C}$ such that $\mathrm{Ob}(\mathsf{C})\subseteq\boldsymbol{\mathcal{U}}$ and, for every $x,y\in\mathrm{Ob}(\mathsf{C})$, $\mathrm{Hom}(x,y)\in\boldsymbol{\mathcal{U}}$, and whose morphisms are the functors between them. For a category $\mathsf{C}$, we denote by $\mathrm{Ob}(\mathsf{C})$ its class of objects and, for $x,y\in\mathrm{Ob}(\mathsf{C})$, by $\mathrm{Hom}_{\mathsf{C}}(x,y)$, or by $\mathrm{Hom}(x,y)$, the set of morphisms from $x$ to $y$.

More specific notions and conventions are introduced in Section~\ref{SPreliminaries} and in the sections in which they are used.

\subsection{A note from the authors}
This project started in the fall of 2018 under the guidance of Juan Climent and Enric Cosme.  Ra\'{u}l Ruiz joined our team in October 2023, and his contributions  constitute an important part of his PhD thesis.

On February 2024, we publicly unveiled our first findings on ArXiv.  This delay in the public presentation of our results was due, on the one hand, to the fact that we wanted to make the presentation as clean as possible and, on the other hand, to the fact that we wanted to present the results in such a way that the relationship between the results of the first and the second part was seen as inevitable and fundamental.  On January 2026, we published the second version on ArXiv, which includes the third part of this work, on morphisms of rewriting systems.

Due to the meticulous and conceptually intricate nature of our project, which has necessitated numerous revisions to achieve a state satisfactory to its authors, the document has expanded beyond our initial expectations. We hope that the latter will be offset by the former.


\subsection{Acknowledgements}
The second and third authors were supported by the  PID2024-159495NB-I00 grant from the
Ministerio de Ciencia e Innovaci\'{o}n, Spain, and the CIAICO/2023/007 grant  from the
Conselleria d’Educaci\'{o}, Universitats i Ocupaci\'{o}, Generalitat Valenciana. The second author has held a Specially Appointed Professor position at Nantong University during the completion of this work. The third author was supported by the  CIACIF/2022/489 grant from the Conselleria d'Educaci\'{o}, Universitats i Ocupaci\'{o}, Generalitat Valenciana co-funded by the European Social Fund. Special thanks to Asunci\'{o}n de Montesa for her categorical knowledge. 


\section{Preliminaries}\label{SPreliminaries}

In this section we collect the notions on many-sorted sets, free monoids, many-sorted algebras, terms, derivors, many-sorted partial algebras, existence equations, adjunctions and the Grothendieck construction that we use in the rest of this work. The definitions are taken from~\cite{CVCLRM26}, where the reader will also find the proofs of the results that we state without proof. Sections~\ref{SSingleSet} and~\ref{SGlobular} present higher-order categorial algebras as models of sets of $\QE$-equations between partial algebras, and Section~\ref{SGenHepDerivors} relates such models through derivors. We follow~\cite{pb86, pb02} for partial algebras and $\QE$-equations.

\subsection{Many-sorted sets}\label{SManySortedSets}

We begin by recalling the notion of many-sorted set.

\begin{definition}\label{DSortedSet}
Let $S$ be a set of sorts in $\boldsymbol{\mathcal{U}}$. An $S$-\emph{sorted set} is a mapping $A = (A_{s})_{s\in S}$ from $S$ to $\boldsymbol{\mathcal{U}}$. If $A$ and $X$ are $S$-sorted sets, then we will say that $X$ is an $S$-\emph{sorted subset} or, to abbreviate, a \emph{subset} of $A$, denoted by $X\subseteq A$, if, for every $s\in S$, $X_{s}\subseteq A_{s}$. We will denote by $\mathrm{Sub}(A)$ the set of all $S$-sorted sets $X$ such that $X\subseteq A$.
\end{definition}

In the following definition we recall the many-sorted relations, functions and mappings.

\begin{definition}\label{DSortedMapping}
Let $S$ be a set of sorts and $A$ and $B$ $S$-sorted sets. The \emph{cartesian product of} $A$ \emph{and} $B$, denoted by $A\times B$, is the $S$-sorted set $\left(A_{s}\times B_{s}\right)_{s\in S}.$

Let $\Phi$ be an $S$-sorted set. We will say that $\Phi$ is an $S$-\emph{sorted relation from} $A$ \emph{to} $B$ if $\Phi\subseteq A\times B$. Thus, for every $s\in S$, $\Phi_{s}\subseteq A_{s}\times B_{s}$. We denote by $\Rel(A,B)$ the set of all $S$-sorted relations from $A$ to $B$. If $A=B$, then we write $\Rel(A)$ instead of $\Rel(A,A)$ and call its elements $S$-\emph{sorted relations on} $A$.

The \emph{diagonal of} $A$, denoted by $\Delta_{A}$, is the $S$-sorted relation on $A$ defined, for every $s\in S$, as $\Delta_{A_{s}}$, and the \emph{codiagonal of} $A$, denoted by $\nabla_{A}$, is the $S$-sorted relation on $A$ defined, for every $s\in S$, as $\nabla_{A_{s}} = A_{s}\times A_{s}$.

Let $\Phi$ be an $S$-sorted relation from $A$ to $B$ and $\Psi$ an $S$-sorted relation from $B$ to $C$, then the \emph{composition of} $\Phi$ \emph{and} $\Psi$, denoted by $\Psi\circ\Phi$, is the $S$-sorted relation from $A$ to $C$ defined, for every $s\in S$, as
$$
\Psi_{s}\circ\Phi_{s}
=\left\lbrace
(x,z)\in A_{s}\times C_{s}
\mid \exists\,y\in B_{s}\,
\left(
(x,y)\in \Phi_{s}\And (y,z)\in \Psi_{s}
\right)
\right\rbrace.
$$
This composition is associative and the diagonal relation $\Delta_{A}$ is a neutral element for it.

Let $\Phi$ be an $S$-sorted relation from $A$ to $B$. Then the \emph{inverse} of $\Phi$, denoted by $\Phi^{-1}$, is the $S$-sorted relation from $B$ to $A$ defined, for every $s\in S$, as
$$
\Phi_{s}^{-1}
=\left\lbrace(y,x)\in B_{s}\times A_{s}\mid (x,y)\in \Phi_{s}
\right\rbrace.
$$

Let $\Phi$ be an $S$-sorted relation from $A$ to $B$, $X$ an $S$-sorted subset of $A$, and $Y$ an $S$-sorted subset of $B$, then the \emph{direct image} of $X$ under $\Phi$, denoted by $\Phi[X]$, is the $S$-sorted subset of $B$ defined, for every $s\in S$, as
$$
\Phi[X]_{s} =\left\lbrace\,b\in B_{s}\mid \exists\, x\in X_{s}((x,b)\in \Phi_{s})\,\right\rbrace,
$$
therefore $\Phi[X] = (\Phi_{s}[X_{s}])_{s\in S}$, and the \emph{inverse image} of $Y$ under $\Phi$, denoted by $\Phi^{-1}[Y]$, is the $S$-sorted subset of $A$ defined, for every $s\in S$, as
$$
\Phi^{-1}[Y]_{s} =\left\lbrace\,a\in A_{s}\mid \exists\, y\in Y_{s}((a,y)\in \Phi_{s})\,\right\rbrace,
$$
therefore $\Phi^{-1}[Y] = (\Phi^{-1}_{s}[Y_{s}])_{s\in S}$. For an $S$-sorted relation $\Phi$ on an $S$-sorted set $A$ we will call $\Phi[A]$ and $\Phi^{-1}[A]$ the \emph{image} and the \emph{domain} of $\Phi$ respectively and we will denote them by $\mathrm{Im}(\Phi)$ and $\Dom(\Phi)$ respectively. Moreover, we will call $\Dom(\Phi)\cup\mathrm{Im}(\Phi)$ the \emph{field} of $\Phi$ and we will denote it by $\mathrm{Fld}(\Phi)$.

An $S$-\emph{sorted function from} $A$ \emph{to} $B$ is a functional $S$-sorted relation $F$ from $A$ to $B$, i.e., an $S$-sorted relation $F$ from $A$ to $B$ such that, for every $s\in S$, $F_{s}$ is a function from $A_{s}$ to $B_{s}$. We denote by $\Fnc(A,B)$ the set of all $S$-sorted functions from $A$ to $B$. The composition of $S$-sorted functions, which is a particular case of the composition of $S$-sorted relations, is an $S$-sorted function.

An $S$-\emph{sorted mapping from} $A$ \emph{to} $B$ is a triple $f = (A,F,B)$ where $F$ is an $S$-sorted function from $A$ to $B$. We denote by $\Hom(A,B)$ or by $B_{A}$ the set of all $S$-sorted mappings from $A$ to $B$. We consider the expressions $f\in\Hom(A,B)$, $f\in B_{A}$, and $f\colon A\mor B$ as synonymous. Moreover, given $f\colon A\mor B$ and $g\colon B\mor C$, $g\circ f = (A,G\circ F,C)$, the \emph{composition of} $f$ \emph{and} $g$, is an $S$-sorted mapping from $A$ to $C$, and $\id^{A} = (A,\Delta_{A},A)$ is an $S$-sorted endomapping of $A$, the identity $S$-sorted mapping at $A$. We denote by $\mathsf{Set}^{S}$ the category of $S$-sorted sets and $S$-sorted mappings.

Let $f\colon A\mor B$ be an $S$-sorted mapping. Then the mapping
$$
f[\cdot]\colon\mathrm{Sub}(A)\mor\mathrm{Sub}(B),
$$
of $f$-\emph{direct image formation}, sends $X \in \mathrm{Sub}(A)$ to $f[X] = (f_{s}[X_{s}])_{s\in S} \in \mathrm{Sub}(B)$, and the mapping
$$
f^{-1}[\cdot]\colon\mathrm{Sub}(B)\mor\mathrm{Sub}(A),
$$
of $f$-\emph{inverse image formation}, sends $Y \in \mathrm{Sub}(B)$ to $f^{-1}[Y] = (f_{s}^{-1}[Y_{s}] )_{s\in S} \in \mathrm{Sub}(A)$.

The \emph{image} of $f$, denoted by $\mathrm{Im}(f)$, is $f[A]$, i.e., $(f_{s}[A_{s}])_{s\in S}$. Moreover, if $X\subseteq A$, then the \emph{restriction of} $f$ \emph{to} $X$, denoted by $f{\upharpoonright}_{X}$, $f{\mid}_{X}$ or $\mathrm{res}_{X}(f)$, is $f\circ \mathrm{in}^{X,A}$, where $\mathrm{in}^{X,A} = (\mathrm{in}^{X_{s},A_{s}})_{s\in S}$ is the canonical embedding of $X$ into $A$.
\end{definition}

\begin{remark}\label{RSortedMapping}
To give an $S$-sorted mapping from $A$ to $B$, as in Definition~\ref{DSortedMapping}, is equivalent to give an $S$-indexed family $f = (f_{s})_{s\in S}$, where, for every $s$ in $S$, $f_{s}$ is a mapping from $A_{s}$ to $B_{s}$. Thus, an $S$-sorted mapping from $A$ to $B$ is, essentially, an element of $\prod_{s\in S}\mathrm{Hom}(A_{s}, B_{s})$.
\end{remark}

Every mapping between sets of sorts induces a functor between the corresponding categories of many-sorted sets.

\begin{proposition}\label{PDeltaPhiFunc}
Let $S$ and $T$ be sets and $\varphi \colon S \mor T$ a mapping. We let $\Delta_{\varphi}$ stand for the assignment from $\mathsf{Set}^{T}$ to $\mathsf{Set}^{S}$ defined as follows:
\begin{enumerate}
\item for every $T$-sorted set $A = (A_{t})_{t \in T}$, $\Delta_{\varphi}(A)$ is the $S$-sorted set $A_{\varphi} = (A_{\varphi(s)})_{s \in S}$, and
\item for every $T$-sorted sets $A$ and $B$, and every $T$-sorted mapping $f \colon A \mor B$, $\Delta_{\varphi}(f)$ is the $S$-sorted mapping $f_{\varphi} = (f_{\varphi(s)})_{s \in S}$.
\end{enumerate}
Then, $\Delta_{\varphi}$ is a covariant functor from $\mathsf{Set}^{T}$ to $\mathsf{Set}^{S}$.
\end{proposition}

We next recall the products and the coproducts of many-sorted sets.

\begin{definition}\label{DSortedProd}
Let $S$ be a set of sorts, $I$ a set in $\boldsymbol{\mathcal{U}}$ and $(A^{i})_{i\in I}$ an $I$-indexed family of $S$-sorted sets. Then the \emph{product} of $(A^{i})_{i\in I}$, denoted by $\prod_{i\in I}A^{i}$, is the $S$-sorted set defined, for every $s\in S$, as $\left(\prod\nolimits_{i\in I}A^{i}\right)_{s} = \prod\nolimits_{i\in I}A^{i}_{s}$, where
$$
\textstyle
\prod_{i\in I}A^{i}_{s} = \left\{(a_{i})_{i\in I}\in\mathrm{Fnc}\left(I,\bigcup_{i\in I}A^{i}_{s}\right)\mid \forall\,i\in I\,\left(a_{i}\in A^{i}_{s}\right)\right\}.
$$
For every $i\in I$, the \emph{i-th canonical projection}, $\mathrm{pr}^{i} = (\mathrm{pr}^{i}_{s})_{s\in S}$, is the $S$-sorted mapping from $\prod_{i\in I}A^{i}$ to $A^{i}$ that, for every $s\in S$, sends $(a_{i})_{i\in I}$ in $\prod_{i\in I}A^{i}_{s}$ to $a_{i}$ in $A^{i}_{s}$. The ordered pair $(\prod_{i\in I}A^{i},(\mathrm{pr}^{i})_{i\in I})$ has the following universal property: For every $S$-sorted set $B$ and every $I$-indexed family of $S$-sorted mappings $(f^{i})_{i\in I}$, where, for every $i\in I$, $f^{i}$ is an $S$-sorted mapping from $B$ to $A^{i}$, there exists a unique $S$-sorted mapping $\left<f^{i}\right>_{i\in I}$ from $B$ to $\prod_{i\in I}A^{i}$ such that, for every $i\in I$, $\mathrm{pr}^{i}\circ \left<f^{i}\right>_{i\in I} = f^{i}$.

The \emph{coproduct} of $(A^{i})_{i\in I}$, denoted by $\coprod_{i\in I}A^{i}$, is the $S$-sorted set defined, for every $s\in S$, as $\left(\coprod\nolimits_{i\in I}A^{i}\right)_{s} = \coprod\nolimits_{i\in I}A^{i}_{s}$, where
$$
\textstyle
\coprod_{i\in I}A^{i}_{s} = \bigcup_{i\in I}(A^{i}_{s}\times\{i\}).
$$
For every $i\in I$, the \emph{i-th canonical injection}, $\mathrm{in}^{i}$, is the $S$-sorted mapping from $A^{i}$ to $\coprod_{i\in I}A^{i}$ that, for every $s\in S$, sends $a$ in $A^{i}_{s}$ to $(a,i)$ in $\coprod_{i\in I}A^{i}_{s}$. The ordered pair $(\coprod_{i\in I}A^{i},(\mathrm{in}^{i})_{i\in I})$ has the following universal property: For every $S$-sorted set $B$ and every $I$-indexed family of $S$-sorted mappings $(f^{i})_{i\in I}$, where, for every $i\in I$, $f^{i}$ is an $S$-sorted mapping from $A^{i}$ to $B$, there exists a unique $S$-sorted mapping $[f^{i}]_{i\in I}$ from $\coprod_{i\in I}A^{i}$ to $B$ such that, for every $i\in I$, $ [f^{i}]_{i\in I}\circ\mathrm{in}^{i} = f^{i}$.

The remaining set-theoretic operations on $S$-sorted sets: $\amalg$ (binary coproduct), $\bigcup$ (union), $\cup$ (binary union), $\bigcap$ (intersection), $\cap$ (binary intersection), $-$ (difference), and $\complement_{A}$ (complement of an $S$-sorted set in a fixed $S$-sorted $A$), are defined in a similar way, i.e., componentwise.
\end{definition}

In the following proposition we recall the notion of corestriction of a many-sorted mapping.

\begin{figure}
$$\xymatrix{
  & A \ar[dl]_{h} \ar[d]^{f} \\
Y \ar @{{ +}{-}{>}}[r]_{\mathrm{in}^{Y,B}} & B
 }
$$
\caption{The corestriction of a many-sorted mapping.}
\label{FPUS}
\end{figure}

\begin{proposition}\label{PUS}
Let $S$ be a set of sorts, $A$ and $B$ $S$-sorted sets, $Y$ a subset of $B$, and $f$ an $S$-sorted mapping from $A$ to $B$. Then the following statements are equivalent:
\begin{enumerate}
\item $\mathrm{Im}(f)\subseteq Y$.
\item There exists an $S$-sorted mapping $h$ from $A$ to $Y$ such that the diagram in Figure~\ref{FPUS} commutes.
\end{enumerate}
If one of the above equivalent statements holds, then we will call $h$, which is univocally determined, the \emph{corestriction of} $f$ \emph{to} $Y$ and we denote it by $f|^{Y}$ or $\mathrm{cores}_{Y}(f)$.
\end{proposition}

\begin{proposition}\label{PPSet}
Let $S$ be a set of sorts and $\mathsf{PSet}^{S}$ the category whose objects are the \emph{$S$-sorted set pairs}, i.e., the ordered pairs $(A,X)$ where $A$ is an $S$-sorted set and $X\subseteq A$, and in which the set of morphisms from $(A,X)$ to $(B,Y)$ is the set of all $S$-sorted mappings $f$ from $A$ to $B$ such that $f[X]\subseteq Y$. Let $G$ be the functor from $\mathsf{Set}^{S}$ to $\mathsf{PSet}^{S}$ whose object mapping sends $A$ to $(A,A)$ and whose morphism mapping sends $f\colon A\mor B$ to $f\colon (A,A)\mor (B,B)$. Then, for every $S$-sorted pair $(B,Y)$, there exists a universal mapping from $G$ to $(B,Y)$, which is precisely the ordered pair $(Y,\mathrm{in}^{Y,B})$ with $\mathrm{in}^{Y,B}\colon (Y,Y)\mor (B,Y)$ the morphism in $\mathsf{PSet}^{S}$ associated to $\mathrm{in}^{Y,B}\colon Y\mor B$.
\end{proposition}

\begin{definition}\label{DBirestriction}
Let $S$ be a set of sorts and $f$ a $\mathsf{PSet}^{S}$-morphism from $(A,X)$ to $(B,Y)$. Then we denote by $f|_{X}^{Y}$, $\mathrm{bires}_{X,Y}(f)$, or, if no confusion can arise, $\widehat{f}$ the $S$-sorted mapping $\mathrm{cores}_{Y}(\mathrm{res}_{X}(f))$ (which is identical to $\mathrm{res}_{X}(\mathrm{cores}_{Y}(f))$). We will call this $S$-sorted mapping the \emph{birestriction of} $f$ to $X$ and $Y$.
\end{definition}

\subsection{Free monoids}\label{SFreeMonoids}

We next recall the concept of free monoid on a set and several notions associated with it, which are used to construct the free algebra on an $S$-sorted set and to define the variables of the terms.

\begin{definition}\label{DFreeMonoid}
Let $A$ be a set. The \emph{free monoid on} $A$, denoted by $\mathbf{A}^{\star}$, is $(A^{\star},\curlywedge,\lambda)$, where $A^{\star}$, the set of all \emph{words on} $A$, is $\bigcup_{n\in\mathbb{N}}\mathrm{Hom}(n,A)$, the set of all mappings $\mathbf{a}\colon n\mor A$ from some $n\in \mathbb{N}$ to $A$. A word $\mathbf{a}\in A^{\star}$ is usually denoted as a sequence $(a_{i})_{i\in\bb{\mathbf{a}}}$, where, for $i\in\bb{\mathbf{a}}$, $a_{i}$ is the letter in $A$ satisfying $\mathbf{a}(i)=a_{i}$. Furthermore, $\curlywedge$, the \emph{concatenation} of words on $A$, is the binary operation on $A^{\star}$ which sends a pair of words $(\mathbf{a},\mathbf{b})$ on $A$ to the mapping $\mathbf{a}\curlywedge \mathbf{b}$ from $\bb{\mathbf{a}}+\bb{\mathbf{b}}$ to $A$, where $\bb{\mathbf{a}}$ and $\bb{\mathbf{b}}$ are the lengths ($\equiv$ domains) of the mappings $\mathbf{a}$ and $\mathbf{b}$, respectively, defined as follows:
$$
\mathbf{a}\bconcat \mathbf{b}
\nfunction
{\bb{\mathbf{a}}+\bb{\mathbf{b}}}{A}
{i}{
\begin{cases}
  a_{i}, & \text{if $0\leq i < \bb{\mathbf{a}}$;}\\
  b_{i-\bb{\mathbf{a}}}, & \text{if $\bb{\mathbf{a}}\leq i < \bb{\mathbf{a}}+\bb{\mathbf{b}}$,}
\end{cases}
   }
$$
and $\lambda$, the \emph{empty word on} $A$, is the unique mapping from $\varnothing$ to $A$. We will denote by $\eta^{A}$ the mapping from $A$ to $A^{\star}$ that sends $a\in A$ to $(a)\in A^{\star}$, i.e., to the mapping $(a)\colon 1\mor A$ that sends $0$ to $a$. The ordered pair $(\mathbf{A}^{\star},\eta^{A})$ is a universal morphism from $A$ to the forgetful functor from the category $\mathsf{Mon}$, of monoids, to $\mathsf{Set}$.
\end{definition}

\begin{remark}\label{RLength}
For a word $\mathbf{a}\in A^{\star}$, $\bb{\mathbf{a}}$, the length of $\mathbf{a}$, is the value at $\mathbf{a}$ of the unique homomorphism $\bb{\,\cdot\,}$ from $\mathbf{A}^{\star}$ to $(\mathbb{N},+,0)$, the additive monoid of the natural numbers, such that $\bb{\,\cdot\,}\circ \eta^{A} = \kappa_{1}$, where $\kappa_{1}$ is the mapping from $A$ to $\mathbb{N}$ constantly $1$. Note that, for every $n\in \mathbb{N}$, $\bb{\,\cdot\,}$ sends $\mathbf{a}\in \mathrm{Hom}(n,A)$ to $n$. Thus, for the family of mappings $(\kappa_{n})_{n\in \mathbb{N}}$, where, for every $n\in \mathbb{N}$, $\kappa_{n}$ is the mapping from $\mathrm{Hom}(n,A)$ to $\mathbb{N}$ constantly $n$, and by applying the universal property of the coproduct, we have $\bb{\,\cdot\,} = [\kappa_{n}]_{n\in \mathbb{N}}$.
\end{remark}

\begin{definition}\label{DSubw}
Let $\mathbf{a}$ and $\mathbf{a}'$ be words in $A^{\star}$ and $k$, $l\in\bb{\mathbf{a}}$ such that $k\leq l$. We will say that $\mathbf{a}'$ is a \emph{subword} of $\mathbf{a}$ \emph{beginning at position} $k$ and \emph{ending at position} $l$ if there are words $\mathbf{b}$ and $\mathbf{c}$ such that $\mathbf{a} = \mathbf{b}\bconcat \mathbf{a}'\bconcat \mathbf{c}$, $\bb{\mathbf{b}} = k$, and $\bb{\mathbf{c}} = (\bb{\mathbf{a}}-1)-l$. For $\mathbf{a}\in A^{\star}$ and $k$, $l\in\bb{\mathbf{a}}$ such that $k\leq l$, we will denote by $\mathbf{a}^{k,l}$ the word $\mathbf{a}'$ beginning at position $k$ and ending at position $l$. In particular, we will denote by $\mathbf{a}^{0,l}$ the word $\mathbf{a}'$ beginning at position $0$ and ending at position $l$, which is an \emph{initial segment} of $\mathbf{a}$, and by $\mathbf{a}^{k,\bb{\mathbf{a}}-1}$ the word $\mathbf{a}'$ beginning at position $k$ and ending at position $\bb{\mathbf{a}}-1$, which is a \emph{final segment} of $\mathbf{a}$. We will call the word $\mathbf{a}^{0,\bb{\mathbf{a}}-2}$ the \emph{maximal prefix} of $\mathbf{a}$. Moreover, we will call $\mathbf{a}^{0,0}$, which is, essentially, the first letter of $\mathbf{a}$, the \emph{head} of $\mathbf{a}$ and $\mathbf{a}^{\bb{\mathbf{a}}-1,\bb{\mathbf{a}}-1}$, which is, essentially, the last letter of $\mathbf{a}$, the \emph{tail} of $\mathbf{a}$. A word $\mathbf{a}$ may have several subwords equal to $\mathbf{a}'$. In that case, the equation $\mathbf{a} = \mathbf{b}\bconcat \mathbf{a}'\bconcat \mathbf{c}$ has several solutions $(\mathbf{b},\mathbf{c})$. If the pairs $(\mathbf{b}_{i},\mathbf{c}_{i})$ $(i\in n)$ are all solutions of $\mathbf{a} = \mathbf{b}\bconcat \mathbf{a}'\bconcat \mathbf{c}$ and if $\bb{\mathbf{b}_{0}}< \bb{\mathbf{b}_{1}}<\cdots<\bb{\mathbf{b}_{n-1}}$, then $(\mathbf{b}_{i},\mathbf{c}_{i})$ determines the $i$-th occurrence of $\mathbf{a}'$ in $\mathbf{a}$. The solution $(\mathbf{b},\mathbf{c})$ in which either $\mathbf{b}$ or $\mathbf{c}$ is $\lambda$ is not excluded.

Let $\mathbf{a}$ be a word in $A^{\star}$ and $a\in A$. We will say that \emph{$a$ occurs in} $\mathbf{a}$ if there are words $\mathbf{b}$, $\mathbf{c}$ in $A^{\star}$ such that $\mathbf{a} = \mathbf{b}\bconcat(a)\bconcat \mathbf{c}$. Note that $a$ occurs in $\mathbf{a}$ if and only if there exists an $i\in \bb{\mathbf{a}}$ such that $\mathbf{a}(i) = a$. We will denote by $\bb{\mathbf{a}}_{a}$ the natural number $\mathrm{card}(\{i\in \bb{\mathbf{a}}\mid \mathbf{a}(i) = a\}) = \mathrm{card}(\mathbf{a}^{-1}[\{a\}])$, i.e., the number of occurrences of $a$ in $\mathbf{a}$. Moreover, we let $(i_{j})_{j\in\bb{\mathbf{a}}_{a}}$ stand for the enumeration in ascending order of the occurrences of $a$ in $\mathbf{a}$. Thus $(i_{j})_{j\in\bb{\mathbf{a}}_{a}}$ is the order embedding of $(\bb{\mathbf{a}}_{a},<)$ into $(\bb{\mathbf{a}},<)$ defined recursively as follows:
$$i_{0} = \mathrm{min}\{i\in \bb{\mathbf{a}}\mid \mathbf{a}(i) = a\},$$
and, for $j\in [1,\bb{\mathbf{a}}_{a}-1]$,
$$
i_{j} = \mathrm{min}\{i\in \bb{\mathbf{a}}-\{i_{0},\ldots,i_{j-1}\}\mid \mathbf{a}(i) = a\}.
$$
If, for every $j\in \bb{\mathbf{a}}_{a}$, the pairs $(\mathbf{b}_{i_{j}},\mathbf{c}_{i_{j}})$ are all solutions of $\mathbf{a} = \mathbf{b}_{i_{j}}\bconcat (a)\bconcat \mathbf{c}_{i_{j}}$ and if $\bb{\mathbf{b}_{i_{0}}}< \bb{\mathbf{b}_{i_{1}}}<\cdots<\bb{\mathbf{b}_{i_{\bb{\mathbf{a}}_{a}-1}}}$, then $(\mathbf{b}_{i_{j}},\mathbf{c}_{i_{j}})$ determines the $j$-th occurrence of $(a)$ in $\mathbf{a}$ and we will say that \emph{$a$ occurs at the $i_{j}$-th place of $\mathbf{a}$}.
\end{definition}

\begin{remark}\label{ROccurrences}
Let $a$ be an element of $A$. Then, for $\mathbf{a}\in A^{\star}$, $\bb{\mathbf{a}}_{a}$, the number of occurrences of $a$ in $\mathbf{a}$, is the value at $\mathbf{a}$ of the unique homomorphism $\bb{\,\cdot\,}_{a}$ from $\mathbf{A}^{\star}$ to $(\mathbb{N},+,0)$ such that $\bb{\,\cdot\,}_{a}\circ\eta^{A} = \delta_{a}$, where $\delta_{a}$ is the mapping from $A$ to $\mathbb{N}$ that sends $a\in A$ to $1\in\mathbb{N}$ and $b\in A-\{a\}$ to $0\in \mathbb{N}$. Therefore, since, for every $\mathbf{a}\in A^{\star}$, the $A$-indexed family $(\bb{\mathbf{a}}_{a})_{a\in A}$ in $\mathbb{N}$ is such that $\bb{\mathbf{a}}_{a} = 0$ for all but a finite number of elements $a$ in $A$, i.e., is such that $\mathrm{card}(\{a\in A\mid \bb{\mathbf{a}}_{a}\neq 0\})<\aleph_{0}$, we have that $\sum_{a\in A}\bb{\mathbf{a}}_{a}\in \mathbb{N}$ and, obviously, for every $\mathbf{a}\in A^{\star}$, $\bb{\mathbf{a}} = \sum_{a\in A}\bb{\mathbf{a}}_{a}$, i.e., $\bb{\,\cdot\,} = \sum_{a\in A}\bb{\,\cdot\,}_{a}$.
\end{remark}

For a mapping $\varphi\colon S\mor T$, we let $\varphi^{\star}$ stand for the unique homomorphism of monoids from $\mathbf{S}^{\star}$ to $\mathbf{T}^{\star}$ such that $\varphi^{\star}\circ\eta^{S}=\eta^{T}\circ\varphi$. Thus, for every $\mathbf{s}\in S^{\star}$, $\varphi^{\star}(\mathbf{s})=(\varphi(s_{j}))_{j\in\bb{\mathbf{s}}}$ and, in particular, $\varphi^{\star}(\lambda)=\lambda$.

\subsection{Many-sorted algebras}\label{SManySortedAlg}

We next recall the notions of many-sorted signature, many-sorted algebra and homomorphism between many-sorted algebras.

\begin{convention}
In what follows, for a set of sorts $S$, an arbitrary word on $S$ will be denoted by $\mathbf{s}$, i.e., a lower case bold type $s$. The letter $s$ will be used to represent an arbitrary letter in $S$.
\end{convention}

\begin{definition}\label{DSignature}
Let $S$ be a set of sorts. An $S$-\emph{sorted signature} is a mapping $\Sigma$ from $S^{\star}\times S$ to $\boldsymbol{\mathcal{U}}$ which sends a pair $(\mathbf{s},s)\in S^{\star}\times S$ to the set $\Sigma_{\mathbf{s},s}$ of the \emph{formal operations} of \emph{arity} $\mathbf{s}$, \emph{sort} (or \emph{coarity}) $s$, and \emph{rank} (or \emph{biarity}) $(\mathbf{s},s)$. A \emph{many-sorted signature} is an ordered pair $\sig{\Sigma}=(S,\Sigma)$, where $S$ is a set of sorts and $\Sigma$ an $S$-sorted signature.
\end{definition}

\begin{definition}\label{DAlg}
Let $S$ be a set of sorts and $\Sigma$ an $S$-sorted signature. The $S^{\star}\times S$-sorted set of the \emph{finitary operations on} an $S$-sorted set $A$ is $(\mathrm{Hom}(A_{\mathbf{s}},A_{s}))_{(\mathbf{s},s)\in S^{\star}\times S}$, where, for every $\mathbf{s}\in S^{\star}$, $A_{\mathbf{s}} = \prod_{j\in \bb{\mathbf{s}}}A_{s_{j}}$, with $\bb{\mathbf{s}}$ denoting the length of the word $\mathbf{s}$ (if $\mathbf{s} = \lambda$, then $A_{\lambda}$ is a final set). Sometimes we let $\mathrm{O}_{\mathrm{H}_{S}}(A)$ stand for $(\mathrm{Hom}(A_{\mathbf{s}},A_{s}))_{(\mathbf{s},s)\in S^{\star}\times S}$. A \emph{structure of} $\Sigma$-\emph{algebra on} an $S$-sorted set $A$ is a family $(F_{\mathbf{s},s})_{(\mathbf{s},s)\in S^{\star}\times S}$, denoted by $F$, where, for $(\mathbf{s},s)\in S^{\star}\times S$, $F_{\mathbf{s},s}$ is a mapping from $\Sigma_{\mathbf{s},s}$ to $\mathrm{Hom}(A_{\mathbf{s}},A_{s})$ (if $(\mathbf{s},s) = (\lambda,s)$ and $\sigma\in \Sigma_{\lambda,s}$, then $F_{\lambda,s}(\sigma)$ picks out an element of $A_{s}$). A \emph{many-sorted} $\Sigma$-\emph{algebra} (or, to abbreviate, $\Sigma$-\emph{algebra}) is a pair $(A,F)$, denoted by $\mathbf{A}$, where $A$ is an $S$-sorted set and $F$ a structure of $\Sigma$-algebra on $A$. For a pair $(\mathbf{s},s)\in S^{\star}\times S$ and a formal operation $\sigma\in \Sigma_{\mathbf{s},s}$, in order to simplify the notation, the operation $F_{\mathbf{s},s}(\sigma)$ from $A_{\mathbf{s}}$ to $A_{s}$ will be written as $F_{\sigma}$. In some cases, to avoid mistakes, we will denote by $F^{\mathbf{A}}$ the structure of $\Sigma$-algebra on $A$, and, for $(\mathbf{s},s)\in S^{\star}\times S$ and $\sigma\in \Sigma_{\mathbf{s},s}$, by $F^{\mathbf{A}}_{\sigma}$, or simply by $\sigma^{\mathbf{A}}$, the corresponding operation. Moreover, for $s\in S$ and $\sigma\in\Sigma_{\lambda,s}$, we will, usually, denote by $\sigma^{\mathbf{A}}$ the value of the mapping $F^{\mathbf{A}}_{\sigma}\colon A_{\lambda}\mor A_{s}$ at the unique element in $A_{\lambda}$.

A $\Sigma$-\emph{homomorphism} (or, to abbreviate, \emph{homomorphism}) from $\mathbf{A}$ to $\mathbf{B}$, where $\mathbf{B} = (B,G)$, is a triple $(\mathbf{A},f,\mathbf{B})$, denoted by $f\colon \mathbf{A}\mor \mathbf{B}$, where $f$ is an $S$-sorted mapping from $A$ to $B$ such that, for every $(\mathbf{s},s)\in S^{\star}\times S$, every $\sigma\in \Sigma_{\mathbf{s},s}$, and every $(a_{j})_{j\in \bb{\mathbf{s}}}\in A_{\mathbf{s}}$, we have
$
f_{s}(\sigma^{\mathbf{A}}((a_{j})_{j\in \bb{\mathbf{s}}})) = \sigma^{\mathbf{B}}(f_{\mathbf{s}}((a_{j})_{j\in \bb{\mathbf{s}}})),
$
where $f_{\mathbf{s}}$ is the mapping $\prod_{j\in \bb{\mathbf{s}}}f_{s_{j}}$ from $A_{\mathbf{s}}$ to $B_{\mathbf{s}}$ that sends $(a_{j})_{j\in \bb{\mathbf{s}}}$ in $A_{\mathbf{s}}$ to $(f_{s_{j}}(a_{j}))_{j\in \bb{\mathbf{s}}}$ in $B_{\mathbf{s}}$. For a $\Sigma$-algebra $\mathbf{A}$ the triple $(\mathbf{A},\id^{A},\mathbf{A})$, where $\id^{A}$ is the identity at $A$, which is a homomorphism from $\mathbf{A}$ to $\mathbf{A}$, will be denoted by $\id^{\mathbf{A}}$ and will be called the \emph{identity} homomorphism at $\mathbf{A}$. We will denote by $\mathsf{Alg}(\Sigma)$ the category of $\Sigma$-algebras and homomorphisms and by $\mathrm{Alg}(\Sigma)$ the set of objects of $\mathsf{Alg}(\Sigma)$.
\end{definition}

We next recall the subalgebras of a many-sorted algebra and the subalgebra generating operator, which we use to construct the free many-sorted algebras.

\begin{definition}\label{DAlgSub}
Let $S$ be a set of sorts, $\Sigma$ an $S$-sorted signature, $\mathbf{A}$ a $\Sigma$-algebra and $X\subseteq A$. Given $(\mathbf{s},s)\in S^{\star}\times S$ and $\sigma\in\Sigma_{\mathbf{s},s}$, we will say that $X$ is \emph{closed under the operation} $\sigma^{\mathbf{A}}\colon A_{\mathbf{s}}\mor A_{s}$ if, for every $(a_{j})_{j\in\bb{\mathbf{s}}}\in X_{\mathbf{s}}$, $\sigma^{\mathbf{A}}((a_{j})_{j\in\bb{\mathbf{s}}})\in X_{s}$. We will say that $X$ is a \emph{closed subset} of $\mathbf{A}$ if $X$ is closed under the operations of $\mathbf{A}$. We will denote by $\mathrm{Cl}(\mathbf{A})$ the set of all closed subsets of $\mathbf{A}$ (which is an algebraic closure system on $A$) and by $\mathbf{Cl}(\mathbf{A})$ the algebraic lattice $(\mathrm{Cl}(\mathbf{A}),\subseteq)$. We will say that a $\Sigma$-algebra $\mathbf{B}$ is a \emph{subalgebra} of $\mathbf{A}$ if $B\subseteq A$ and the canonical embedding of $B$ into $A$ determines an embedding of $\mathbf{B}$ into $\mathbf{A}$. We will denote by $\mathrm{Sub}(\mathbf{A})$ the set of all subalgebras of $\mathbf{A}$. Since $\mathrm{Cl}(\mathbf{A})$ and $\mathrm{Sub}(\mathbf{A})$ are isomorphic, we shall feel free to deal either with a closed subset of $\mathbf{A}$ or with the correlated subalgebra of $\mathbf{A}$, whichever is most convenient for the work at hand.
\end{definition}

\begin{definition}\label{DSg}
Let $S$ be a set of sorts, $\Sigma$ an $S$-sorted signature and $\mathbf{A}$ a $\Sigma$-algebra. Then we denote by $\mathrm{Sg}_{\mathbf{A}}$ the many-sorted closure operator on $A$ defined as follows:
$$\textstyle
  \mathrm{Sg}_{\mathbf{A}}\nfunction
  {\mathrm{Sub}(A)}
  {\mathrm{Sub}(A)}
  {X}
  {\bigcap \{\,C\in\mathrm{Sub}(\mathbf{A})\mid X\subseteq C\,\}}
$$
We call $\mathrm{Sg}_{\mathbf{A}}$ the \emph{subalgebra generating many-sorted operator on} $A$ \emph{determined by} $\mathbf{A}$. For every $X\subseteq A$, we call $\mathrm{Sg}_{\mathbf{A}}(X)$ the \emph{subalgebra of} $\mathbf{A}$ \emph{generated by} $X$. Moreover, if $X\subseteq A$ is such that $\mathrm{Sg}_{\mathbf{A}}(X) = A$, then we say that $X$ is an $S$-sorted set of \emph{generators} of $\mathbf{A}$, or that $X$ \emph{generates} $\mathbf{A}$.
\end{definition}

\begin{proposition}[Principle of Proof by Algebraic Induction]\label{PPAI}
Let $S$ be a set of sorts, $\Sigma$ an $S$-sorted signature and $\mathbf{A}$ a $\Sigma$-algebra generated by $X$. Then to prove that a subset $Y$ of $A$ is equal to $A$ it suffices to show: (1) $X\subseteq Y$ (algebraic induction basis) and (2) $Y$ is a subalgebra of $\mathbf{A}$ (algebraic induction step).
\end{proposition}

\subsection{The free algebra on a many-sorted set}\label{SFreeAlg}

We next state that the forgetful functor $\mathrm{G}_{\Sigma}$ from $\mathsf{Alg}(\Sigma)$ to $\mathsf{Set}^{S}$ has a left adjoint $\mathbf{T}_{\Sigma}$ which assigns to an $S$-sorted set $X$ the free $\Sigma$-algebra $\mathbf{T}_{\Sigma}(X)$ on $X$. Let us note that in what follows, to construct the algebra of $\Sigma$-rows in $X$, and the free $\Sigma$-algebra on $X$, since neither the $S$-sorted signature $\Sigma$ nor the $S$-sorted set $X$ are subject to any constraint, coproducts must necessarily be used.

\begin{definition}\label{DRows}
Let $S$ be a set of sorts, $\Sigma$ an $S$-sorted signature and $X$ an $S$-sorted set. The \emph{algebra of} $\Sigma$-\emph{rows in} $X$, denoted by $\mathbf{W}_{\Sigma}(X)$, is defined as follows:
\begin{enumerate}
\item The underlying $S$-sorted set of $\mathbf{W}_{\Sigma}(X)$, written as $\mathrm{W}_{\Sigma}(X)$, is precisely the $S$-sorted set $((\coprod\Sigma \amalg \coprod X)^{\star})_{s\in S}$, i.e., the mapping from $S$ to $\boldsymbol{\mathcal{U}}$ constantly $(\coprod\Sigma \amalg \coprod X)^{\star}$, where $(\coprod\Sigma \amalg \coprod X)^{\star}$ is the set of all words on the set $\coprod\Sigma \amalg \coprod X$, i.e., on the set
      $$
      \textstyle
      [(\bigcup_{(\mathbf{s},s)\in S^{\star}\times S}(\Sigma_{\mathbf{s},s}\times\{(\mathbf{s},s)\}))\times\{0\}]\cup
      [(\bigcup_{s\in S}(X_{s}\times\{s\}))\times\{1\}].
      $$
\item For every $(\mathbf{s},s)\in S^{\star}\times S$, and every $\sigma\in\Sigma_{\mathbf{s},s}$, the structural operation $F_{\sigma}$ associated to $\sigma$ is the mapping from $\mathrm{W}_{\Sigma}(X)_{\mathbf{s}}$ to ${\mathrm{W}_{\Sigma}(X)}_{s}$ which sends $(P_{j})_{j\in\bb{\mathbf{s}}} \in \mathrm{W}_{\Sigma}(X)_{\mathbf{s}}$ to $(\sigma)\curlywedge\concat_{j\in\bb{\mathbf{s}}}P_{j} \in {\mathrm{W}_{\Sigma}(X)}_{s}$, where, for every $(\mathbf{s},s)\in S^{\star}\times S$, and every $\sigma\in\Sigma_{\mathbf{s},s}$, $(\sigma)$ stands for $(((\sigma,(\mathbf{s},s)),0))$, which is the value at $\sigma$ of the canonical mapping from $\Sigma_{\mathbf{s},s}$ to $(\coprod\Sigma \amalg \coprod X)^{\star}$.
\end{enumerate}
\end{definition}

\begin{definition}\label{DAlgFree}
Let $S$ be a set of sorts, $\Sigma$ an $S$-sorted signature and $X$ an $S$-sorted set. The \emph{free} $\Sigma$-\emph{algebra on} $X$, denoted by $\mathbf{T}_{\Sigma}(X)$, is the $\Sigma$-algebra determined by $\mathrm{Sg}_{\mathbf{W}_{\Sigma}(X)}((\{(x)\mid x\in X_{s}\})_{s\in S})$, the subalgebra of $\mathbf{W}_{\Sigma}(X)$ generated by $(\{(x)\mid x\in X_{s}\})_{s\in S}$, where, for every $s\in S$ and every $x\in X_{s}$, $(x)$ stands for $(((x,s),1))$, which is the value at $x$ of the canonical mapping from $X_{s}$ to $(\coprod\Sigma \amalg \coprod X)^{\star}$. We will denote by $\mathrm{T}_{\Sigma}(X)$ the underlying $S$-sorted set of $\mathbf{T}_{\Sigma}(X)$ and, for $s\in S$, we will call the elements of $\mathrm{T}_{\Sigma}(X)_{s}$ \emph{terms of type} $s$ \emph{with variables in} $X$ or $(X,s)$-\emph{terms}.
\end{definition}

Since $(\{(x)\mid x\in X_{s}\})_{s\in S}$ is a generating subset of $\mathbf{T}_{\Sigma}(X)$, to prove that a subset $\mathcal{T}$ of $\mathrm{T}_{\Sigma}(X)$ is equal to $\mathrm{T}_{\Sigma}(X)$ it suffices, by Proposition~\ref{PPAI}, to show: (1) $(\{(x)\mid x\in X_{s}\})_{s\in S}\subseteq \mathcal{T}$ (algebraic induction basis) and (2) $\mathcal{T}$ is a subalgebra of $\mathbf{T}_{\Sigma}(X)$ (algebraic induction step). In the many-sorted case we have, as in the single-sorted case, the following characterization of the elements of $\mathrm{T}_{\Sigma}(X)_{s}$, for $s\in S$.

\begin{proposition}\label{PTermChar}
Let $S$ be a set of sorts, $\Sigma$ an $S$-sorted signature and $X$ an $S$-sorted set. Then, for every $s\in S$ and every $P\in \mathrm{T}_{\Sigma}(X)_{s}$, we have that $P$ is a term of type $s$ with variables in $X$ if and only if $P = (x)$, for a unique $x\in X_{s}$, or $P = (\sigma)$, for a unique $\sigma\in\Sigma_{\lambda,s}$, or $P = (\sigma)\curlywedge\concat(P_{j})_{j\in\bb{\mathbf{s}}}$, for a unique $\mathbf{s}\in S^{\star}-\{\lambda\}$, a unique $\sigma\in\Sigma_{\mathbf{s},s}$, and a unique family $(P_{j})_{j\in\bb{\mathbf{s}}}\in\mathrm{T}_{\Sigma}(X)_{\mathbf{s}}$.

Moreover, the three possibilities are mutually exclusive. From now on, for simplicity of notation, we will write $x$, $\sigma^{\mathbf{T}_{\Sigma}(X)}$, and $\sigma^{\mathbf{T}_{\Sigma}(X)}(P_{0},\ldots,P_{\bb{\mathbf{s}}-1})$ or $\sigma^{\mathbf{T}_{\Sigma}(X)}((P_{j})_{j\in \bb{\mathbf{s}}})$ instead of $(x)$, $(\sigma)$, and $(\sigma)\curlywedge\concat(P_{j})_{j\in\bb{\mathbf{s}}}$, respectively. Thus, in particular, the structural operation $F_{\sigma}$ (or more accurately $F_{\sigma}^{\mathbf{T}_{\Sigma}(X)}$) associated with $\sigma$ is identified with $\sigma^{\mathbf{T}_{\Sigma}(X)}$.
\end{proposition}

When no confusion can arise, we write $\sigma((P_{j})_{j\in\bb{\mathbf{s}}})$ for $\sigma^{\mathbf{T}_{\Sigma}(X)}((P_{j})_{j\in\bb{\mathbf{s}}})$. From Proposition~\ref{PTermChar} follows, immediately, the universal property of the free $\Sigma$-algebra on an $S$-sorted set $X$, as stated in the subsequent proposition.

\begin{proposition}\label{PPropUniv}
Let $S$ be a set of sorts, $\Sigma$ an $S$-sorted signature and $X$ an $S$-sorted set. Then the pair $(\eta^{X},\mathbf{T}_{\Sigma}(X))$, where $\eta^{X}$, the \emph{insertion of (the $S$-sorted set of generators)} $X$ \emph{into} $\mathrm{T}_{\Sigma}(X)$, is the corestriction to $\mathrm{T}_{\Sigma}(X)$ of the canonical embedding of $X$ into $\mathrm{W}_{\Sigma}(X)$, has the following universal property: for every $\Sigma$-algebra $\mathbf{A}$ and every $S$-sorted mapping $f\colon X\mor A$, there exists a unique homomorphism $f^{\sharp}\colon\mathbf{T}_{\Sigma}(X)\mor\mathbf{A}$ such that $f^{\sharp}\circ \eta^{X} = f$.
\end{proposition}
\begin{proof}
For every $s\in S$ and every $(X,s)$-term $P$, the $s$-th coordinate $f^{\sharp}_{s}$ of $f^{\sharp}$ is defined recursively as follows: $f^{\sharp}_{s}(x) = f_{s}(x)$, if $P = x$; $f^{\sharp}_{s}(\sigma^{\mathbf{T}_{\Sigma}(X)}) = \sigma^{\mathbf{A}}$, if $P = \sigma^{\mathbf{T}_{\Sigma}(X)}$; and, finally, $f^{\sharp}_{s}(\sigma^{\mathbf{T}_{\Sigma}(X)}((P_{j})_{j\in\bb{\mathbf{s}}})) = \sigma^{\mathbf{A}}((f^{\sharp}_{s_{j}}(P_{j}))_{j\in\bb{\mathbf{s}}})$, if $P = \sigma^{\mathbf{T}_{\Sigma}(X)}((P_{j})_{j\in\bb{\mathbf{s}}})$.
\end{proof}

The just stated proposition allows us to carry out definitions by algebraic recursion on a free many-sorted algebra, as indeed we will be doing throughout this work.

\begin{corollary}\label{CTermAdj}
Let $S$ be a set of sorts and $\Sigma$ an $S$-sorted signature. Then the functor $\mathbf{T}_{\Sigma}$, which sends an $S$-sorted set $X$ to $\mathbf{T}_{\Sigma}(X)$ and an $S$-sorted mapping $f$ from $X$ to $Y$ to $f^{@} (= (\eta^{Y}\circ f)^{\sharp})$, the unique homomorphism from $\mathbf{T}_{\Sigma}(X)$ to $\mathbf{T}_{\Sigma}(Y)$ such that $f^{@}\circ\eta^{X} = \eta^{Y}\circ f$, is left adjoint for the forgetful functor $\mathrm{G}_{\Sigma}$ from $\mathsf{Alg}(\Sigma)$ to $\mathsf{Set}^{S}$.
\end{corollary}

We introduce two small results that will be of interest later on.

\begin{proposition}\label{PPropUnivId}
Let $S$ be a set of sorts, $\Sigma$ an $S$-sorted signature and $X$ an $S$-sorted set. Then $(\eta^{X})^{\sharp} = \id^{\mathbf{T}_{\Sigma}(X)}$.
\end{proposition}

\begin{proposition}\label{PPropUnivComp}
Let $S$ be a set of sorts, $\Sigma$ an $S$-sorted signature, $X$ an $S$-sorted set, $\mathbf{A}$ and $\mathbf{B}$ $\Sigma$-algebras, $f$ a mapping from $X$ to $A$ and $g$ a $\Sigma$-homomorphism from $\mathbf{A}$ to $\mathbf{B}$. Then $(g \circ f)^{\sharp} = g \circ f^{\sharp}$.
\end{proposition}
\begin{proof}
By associativity and the universal property of $f^{\sharp}$, we have that
$$
(g \circ f^{\sharp}) \circ \eta^{X}
=
g \circ (f^{\sharp} \circ \eta^{X})
=
g \circ f.
$$
Thus, the result follows by the uniqueness of $(g \circ f)^{\sharp}$.
\end{proof}

\subsection{Variables and substitutions}\label{SVariables}

We next recall the $S$-sorted sets of variables that we associate to the words on $S$, and the identifications of the families of elements indexed by a word with the valuations of these variables.

\begin{definition}\label{DVariables}
Let $S$ be a set of sorts. For every word $\mathbf{s}\in S^{\star}$, we let $\vs{\mathbf{s}}$ stand for the $S$-sorted set defined, for every $s\in S$, as $(\vs{\mathbf{s}})_{s} = \mathbf{s}^{-1}[\{s\}]$, i.e., the set of the positions at which $s$ occurs in $\mathbf{s}$.

Let $V^{S} = (V^{S}_{s})_{s\in S}$ be an arbitrary, but fixed, $S$-sorted set of variables with a countably infinity of variables in each component, say $V^{S}_{s} = \{\,v^{s}_{n}\mid n\in \mathbb{N}\,\}$, for every $s\in S$. Then, for every word $\mathbf{s}\in S^{\star}$, we let $V^{S}_{\mathbf{s}}$ stand for the finite subset of $V^{S}$ defined, for every $s\in S$, as
$$
V^{S}_{\mathbf{s},s} = \{\,v^{s}_{i}\in V^{S}_{s}\mid i\in \bb{\mathbf{s}}_{s}\,\}
$$
(observe that $V^{S}_{\mathbf{s},s}$ is empty for those sorts $s$ that do not occur in the word $\mathbf{s}$).
\end{definition}

\begin{proposition}\label{PVariables}
Let $S$ be a set of sorts and $\mathbf{s}$ an element of $S^{\star}$. Then the following statements hold:
\begin{enumerate}
\item there exists an isomorphism $\vartheta_{\mathbf{s}}$ from $\vs{\mathbf{s}}$ to $V^{S}_{\mathbf{s}}$, which, for every $s\in S$ and every $j\in\bb{\mathbf{s}}_{s}$, sends $i_{j}$ to $v^{s}_{j}$, where $(i_{j})_{j\in\bb{\mathbf{s}}_{s}}$ is the enumeration in ascending order of the occurrences of $s$ in $\mathbf{s}$, see Definition~\ref{DSubw}; and
\item for every $S$-sorted set $A$, there exists an isomorphism $\theta_{\mathbf{s},A}$ from $A_{\vs{\mathbf{s}}}=\mathrm{Hom}(\vs{\mathbf{s}},A)$ to $A_{\mathbf{s}}$, which sends $a$ to $(a_{s_{j}}(j))_{j\in\bb{\mathbf{s}}}$, and, for every $S$-sorted mapping $f\colon A\mor B$, $f_{\mathbf{s}}\circ\theta_{\mathbf{s},A}=\theta_{\mathbf{s},B}\circ f_{\vs{\mathbf{s}}}$, where $f_{\vs{\mathbf{s}}}$ sends $a\in A_{\vs{\mathbf{s}}}$ to $f\circ a\in B_{\vs{\mathbf{s}}}$.
\end{enumerate}
\end{proposition}

In virtue of the previous proposition, for every word $\mathbf{s}\in S^{\star}$ and every $S$-sorted set $A$, the sets $A_{\vs{\mathbf{s}}}$, $A_{\mathbf{s}}$ and $A_{V^{S}_{\mathbf{s}}}$ are naturally isomorphic. In what follows, if there is no ambiguity, we will not distinguish notationally between the elements of $A_{\vs{\mathbf{s}}}$ and those of $A_{\mathbf{s}}$ or those of $A_{V^{S}_{\mathbf{s}}}$. Also, for simplicity of notation, instead of using $V^{S}_{\mathbf{s}}$ we will use $\vs{\mathbf{s}}$. Accordingly, for $j\in\bb{\mathbf{s}}$, we let $v^{\mathbf{s}}_{j}$ stand for the variable $\vartheta_{\mathbf{s},s_{j}}(j)$ of $V^{S}_{\mathbf{s}}$ that corresponds to the $j$-th letter of $\mathbf{s}$, i.e., $v^{\mathbf{s}}_{j}=v^{s_{j}}_{i}$ if the $j$-th letter of $\mathbf{s}$ is the $i$-th occurrence, counting from $0$, of the sort $s_{j}$ in $\mathbf{s}$; and a family $\mathbf{a}=(a_{j})_{j\in\bb{\mathbf{s}}}\in A_{\mathbf{s}}$ is identified with the valuation of $\vs{\mathbf{s}}$ in $A$ that sends $v^{\mathbf{s}}_{j}$ to $a_{j}$, for every $j\in\bb{\mathbf{s}}$. For instance, $\vs{(s)}=\{v^{s}_{0}\}$, so that an element $a\in A_{s}$ is identified with the valuation of $\vs{(s)}$ in $A$ that sends $v^{s}_{0}$ to $a$; and the variables of $\vs{(s,s)}$ are $v^{(s,s)}_{0}=v^{s}_{0}$ and $v^{(s,s)}_{1}=v^{s}_{1}$.

We next fix the notation for substitutions. Let $S$ be a set of sorts, $\Sigma$ an $S$-sorted signature, $X$ and $Y$ $S$-sorted sets and $f\colon Y\mor\mathrm{T}_{\Sigma}(X)$ an $S$-sorted mapping. By Proposition~\ref{PPropUniv}, $f$ determines a unique homomorphism $f^{\sharp}\colon\mathbf{T}_{\Sigma}(Y)\mor\mathbf{T}_{\Sigma}(X)$ such that $f^{\sharp}\circ\eta^{Y}=f$, the \emph{substitution} determined by $f$. When $Y=\vs{\mathbf{s}}$, for a word $\mathbf{s}\in S^{\star}$, $f$ is identified with the family $\mathcal{Q}=(Q_{j})_{j\in\bb{\mathbf{s}}}=(f(v^{\mathbf{s}}_{j}))_{j\in\bb{\mathbf{s}}}\in\mathrm{T}_{\Sigma}(X)_{\mathbf{s}}$, and the substitution $f^{\sharp}$, which sends $v^{\mathbf{s}}_{j}$ to $Q_{j}$, for every $j\in\bb{\mathbf{s}}$, is also denoted by $\left(\left(\!\begin{smallmatrix}v^{\mathbf{s}}_{j}\\Q_{j}\end{smallmatrix}\!\right)_{j\in \bb{\mathbf{s}}}\right)^{\sharp}$, by $\mathcal{Q}^{\sharp}$ or by $(Q_{j})^{\sharp}_{j\in\bb{\mathbf{s}}}$. Finally, for every set of $\mathrm{E}$-equations $\mathcal{E}\subseteq\mathrm{Eeqt}(\Sigma)_{Y}$, see Definition~\ref{DEEq} below, we let $f^{\sharp}[\mathcal{E}]\subseteq\mathrm{Eeqt}(\Sigma)_{X}$ stand for the $S$-sorted set $(\{(f^{\sharp}_{s}(P),f^{\sharp}_{s}(Q))\mid(P,Q)\in\mathcal{E}_{s}\})_{s\in S}$.

\subsection{Derivors}\label{SPrelimDerivors}

In this subsection, after defining the variety of Hall algebras, we introduce the notion of derivor between many-sorted signatures and the category $\mathsf{Sig}_{\mathfrak{d}}$ of many-sorted signatures and derivors, see~\cite{CVST10b}. Before defining the notion of Hall algebra, we recall that
\begin{enumerate}
	\item a finitary specification is an ordered triple $(S,\Sigma,\mathcal{E})$, where $S$ is a set of sorts, $\Sigma$ an $S$-sorted signature, and $\mathcal{E}\subseteq\mathrm{Eq}_{V^{S}}(\Sigma)=
	(\mathrm{T}_{\Sigma}(X)_{s}^{2})_{(X,s)\in \mathrm{Sub}_{\mathrm{f}}(V^{S})\times S}$, i.e., a set of finitary $\Sigma$-equations, where $V^{S}$, the $S$-sorted set of the variables, is a fixed $S$-countably infinite $S$-sorted set and $\mathrm{Sub}_{\mathrm{f}}(V^{S})$ is the set of its finite subsets; and that
	\item a $\Sigma$-algebra $\mathbf{A}$ is a $(S,\Sigma,\mathcal{E})$-algebra if $\mathbf{A}\models^{\Sigma} \mathcal{E}$, i.e., if, for every $(X,s)\in \mathrm{Sub}_{\mathrm{f}}(V^{S})\times S$ and every $(P,Q)\in \mathcal{E}_{X,s}$, $\mathbf{A}$ is a model of $(P,Q)$, in symbols $\mathbf{A}\models^{\Sigma}_{X,s} (P,Q)$, which in turn means that, for every $a\in \mathrm{Hom}(X,A)$, $a^{\sharp}_{s}(P) = a^{\sharp}_{s}(Q)$.
\end{enumerate}

\begin{definition}\label{DHall}
	Let $S$ be a set of sorts and let $V^{\mathrm{H}_{S}} = (V^{\mathrm{H}_{S}}_{\mathbf{s},s})_{(\mathbf{s},s)\in S^{\star}\times S}$ be a fixed $S^{\star}\times S$-sorted set of variables with a countably infinity of variables in each component, say $V^{\mathrm{H}_{S}}_{\mathbf{s},s} = \{\,v^{\mathbf{s},s}_{n}\mid n\in \mathbb{N}\,\}$ (which is isomorphic to $\{(\mathbf{s},s)\}\times\mathbb{N}$), for every $(\mathbf{s},s)\in S^{\star}\times S$. A \emph{Hall algebra for} $S$ is a $\mathrm{H}_{S} = (S^{\star}\times S,\Sigma^{\mathrm{H}_{S}},\mathcal{E}^{\mathrm{H}_{S}})$-algebra, where $\Sigma^{\mathrm{H}_{S}}$ is the $S^{\star}\times S$-sorted signature, i.e., the $(S^{\star}\times S)^{\star}\times (S^{\star}\times S)$-sorted set, defined as follows:
	\begin{enumerate}
		\item[$\mathrm{HS}_{1}$.] For every $\mathbf{s}\in S^{\star}$ and $i\in\bb{\mathbf{s}}$,
		$$
		\pi^{\mathbf{s}}_{i}\colon\lambda\mor(\mathbf{s},s_{i}),
		$$
		where $\bb{\mathbf{s}}$ is the \emph{length} of the word $\mathbf{s}$ and $\lambda$ the \emph{empty word} in the underlying set of the free monoid on $S^{\star}\times S$.
		\item[$\mathrm{HS}_{2}$.] For every $\mathbf{r}$, $\mathbf{s}\in S^{\star}$ and $s\in S$,
		$$
		\xi_{\mathbf{r},\mathbf{s},s}\colon ((\mathbf{s},s),(\mathbf{r},s_{0}),\ldots,(\mathbf{r},s_{\bb{\mathbf{s}}-1}))\mor(\mathbf{r},s);
		$$
	\end{enumerate}
	while $\mathcal{E}^{\mathrm{H}_{S}}$ is the sub-$(S^{\star}\times S)^{\star}\times (S^{\star}\times S)$-sorted set of $\mathrm{Eq}(\Sigma^{\mathrm{H}_{S}})$, where
	$$
	\mathrm{Eq}\left(\Sigma^{\mathrm{H}_{S}}\right) =
	\left(\mathrm{T}_{\Sigma^{\mathrm{H}_{S}}}(\vs{\ol{\mathbf{s}}})_{(\mathbf{r},s)}^{2}\right)_{(\ol{\mathbf{s}},(\mathbf{r},s))\in
		(S^{\star}\times S)^{\star}\times (S^{\star}\times S)},
	$$
	defined as follows:
	\begin{enumerate}
		\item[$\mathrm{H}_{1}$.] \emph{Projection}. For every $\mathbf{r}$, $\mathbf{s}\in S^{\star}$ and $i\in\bb{\mathbf{s}}$, the equation
		$$
		\xi_{\mathbf{r},\mathbf{s},s_{i}}(\pi^{\mathbf{s}}_{i},v^{\mathbf{r},s_{0}}_{0},\ldots,
		v^{\mathbf{r},s_{\bb{\mathbf{s}}-1}}_{\bb{\mathbf{s}}-1})=v^{\mathbf{r},s_{i}}_{i}
		$$
		of type $(((\mathbf{r},s_{0}),\ldots,(\mathbf{r},s_{\bb{\mathbf{s}}-1})),(\mathbf{r},s_{i})).$
		\item[$\mathrm{H}_{2}$.] \emph{Identity}. For every $\mathbf{r}\in S^{\star}$ and $j\in \bb{\mathbf{r}}$, the equation
		$$
		\xi_{\mathbf{r},\mathbf{r},r_{j}}(v^{\mathbf{r},r_{j}}_{j},\pi^{\mathbf{r}}_{0},\ldots,\pi^{\mathbf{r}}_{\bb{\mathbf{r}}-1})=
		v^{\mathbf{r},r_{j}}_{j}
		$$
		of type $(((\mathbf{r},r_{j})),(\mathbf{r},r_{j})).$
		\item[$\mathrm{H}_{3}$.] \emph{Associativity}. For every $\mathbf{q}$, $\mathbf{r}$, $\mathbf{s}\in S^{\star}$ and $s\in S$, the equation
		\begin{multline*}
			\xi_{\mathbf{q},\mathbf{r},s}(
			\xi_{\mathbf{r},\mathbf{s},s}(v^{\mathbf{s},s}_{0},v^{\mathbf{r},s_{0}}_{1},\ldots,
			v^{\mathbf{r},s_{\bb{\mathbf{s}}-1}}_{\bb{\mathbf{s}}}),
			v^{\mathbf{q},r_{0}}_{\bb{\mathbf{s}}+1},
			\ldots,v^{\mathbf{q},r_{\bb{\mathbf{r}}-1}}_{\bb{\mathbf{s}}+\bb{\mathbf{r}}}) = \\
			\xi_{\mathbf{q},\mathbf{s},s}(v^{\mathbf{s},s}_{0},
			\xi_{\mathbf{q},\mathbf{r},s_{0}}(v^{\mathbf{r},s_{0}}_{1},v^{\mathbf{q},r_{0}}_{\bb{\mathbf{s}}+1},
			\ldots,v^{\mathbf{q},r_{\bb{\mathbf{r}}-1}}_{\bb{\mathbf{s}}+\bb{\mathbf{r}}}),
			\\
			\ldots,
			\xi_{\mathbf{q},\mathbf{r},s_{\bb{\mathbf{s}}-1}}(v^{\mathbf{r},s_{\bb{\mathbf{s}}-1}}_{\bb{\mathbf{s}}},
			v^{\mathbf{q},r_{0}}_{\bb{\mathbf{s}}+1},
			\ldots,v^{\mathbf{q},r_{\bb{\mathbf{r}}-1}}_{\bb{\mathbf{s}}+\bb{\mathbf{r}}}))
		\end{multline*}
		of type $(((\mathbf{s},s),(\mathbf{r},s_{0}),\ldots,(\mathbf{r},s_{\bb{\mathbf{s}}-1}),
		(\mathbf{q},r_{0}),\ldots,(\mathbf{q},r_{\bb{\mathbf{r}}-1})),(\mathbf{q},s)).$
	\end{enumerate}
	We call the formal constants $\pi^{\mathbf{s}}_{i}$ \emph{projections}, and the formal operations $\xi_{\mathbf{r},\mathbf{s},s}$ \emph{substitution operators}. Furthermore, we will denote by $\mathsf{Alg}(\mathrm{H}_{S})$ the category of Hall algebras for $S$ and homomorphisms between Hall algebras. Since $\mathsf{Alg}(\mathrm{H}_{S})$ is a variety, the forgetful functor $\mathrm{G}_{\mathrm{H}_{S}}$ from $\mathsf{Alg}(\mathrm{H}_{S})$ to $\mathsf{Set}^{S^{\star}\times S}$ has a left adjoint $\mathbf{T}_{\mathrm{H}_{S}}$, which assigns to an $S^{\star}\times S$-sorted set $\Sigma$ the corresponding free Hall algebra $\mathbf{T}_{\mathrm{H}_{S}}(\Sigma)$.
\end{definition}

For every $S$-sorted set $A$, $\mathrm{O}_{\mathrm{H}_{S}}(A) = (\mathrm{Hom}(A_{\mathbf{s}},A_{s}))_{(\mathbf{s},s)\in S^{\star}\times S}$, the $S^{\star}\times S$-sorted set of operations for $A$, is naturally equipped with a structure of Hall algebra, as stated in the following proposition, if we realize the projections as the true projections and the substitution operators as the generalized composition of mappings.

\begin{proposition}\label{PHallOp}
	Let $S$ be a set of sorts, $A$ an $S$-sorted set and $\mathbf{O}_{\mathrm{H}_{S}}(A)$ the $\Sigma^{\mathrm{H}_{S}}$-algebra with underlying many-sorted set $\mathrm{O}_{\mathrm{H}_{S}}(A)$ and algebraic structure defined as follows:
	\begin{enumerate}
		\item For every $\mathbf{s}\in S^{\star}$ and $i\in\bb{\mathbf{s}}$, $(\pi^{\mathbf{s}}_{i})^{\mathbf{O}_{\mathrm{H}_{S}}(A)} = \mathrm{pr}^{A}_{\mathbf{s},i}\colon A_{\mathbf{s}}\mor A_{s_{i}}$.
		\item For every $\mathbf{r},\mathbf{s}\in S^{\star}$ and $s\in S$, $\xi_{\mathbf{r},\mathbf{s},s}^{\mathbf{O}_{\mathrm{H}_{S}}(A)}$ is defined, for every $f\in A_{s}^{ A_{\mathbf{s}} }$ and $g\in A^{A_{\mathbf{r}}}_{\mathbf{s}}$, as $\xi_{\mathbf{r},\mathbf{s},s}^{\mathbf{O}_{\mathrm{H}_{S}}(A)}(f,g_{0},\ldots,g_{\bb{\mathbf{s}}-1}) = f\circ\langle g_{i}\rangle_{i\in\bb{\mathbf{s}} }$, where $\langle g_{i}\rangle_{i\in\bb{\mathbf{s}}}$ is the unique mapping from $A_{\mathbf{r}}$ to $A_{\mathbf{s}}$ such that, for every $i\in \bb{\mathbf{s}}$, we have that
		$$
		\mathrm{pr}^{A}_{\mathbf{s},i}\circ\langle g_{i}\rangle_{i\in\bb{\mathbf{s}}} = g_{i}.
		$$
	\end{enumerate}
	Then $\mathbf{O}_{\mathrm{H}_{S}}(A)$ is a Hall algebra, the \emph{Hall algebra for} $(S,A)$.
\end{proposition}

For every $S$-sorted signature $\Sigma$, $\mathrm{Tm}_{\mathrm{H}_{S}}(\Sigma) = (\mathrm{T}_{\Sigma}(\vs{\mathbf{s}})_{s})_{(\mathbf{s},s)\in S^{\star}\times S}$ is also equipped with a structure of Hall algebra that formalizes the concept of substitution as stated in the following proposition.

\begin{proposition}\label{PHallTm}
	Let $S$ be a set of sorts, $\Sigma$ an $S$-sorted signature and $\mathbf{Tm}_{\mathrm{H}_{S}}(\Sigma)$ the $\Sigma^{\mathrm{H}_{S}}$-algebra with underlying many-sorted set $\mathrm{Tm}_{\mathrm{H}_{S}}(\Sigma)$ and algebraic structure defined as follows:
	\begin{enumerate}
		\item For every $\mathbf{s}\in S^{\star}$ and $i\in\bb{\mathbf{s}}$, $(\pi^{\mathbf{s}}_{i})^{\mathbf{Tm}_{\mathrm{H}_{S}}(\Sigma)}$ is the image under $\eta^{\vs{\mathbf{s}}}_{s_{i}}$ of the variable $v^{\mathbf{s}}_{i}$, where $\eta^{\vs{\mathbf{s}}} = (\eta^{\vs{\mathbf{s}}}_{s})_{s\in S}$ is the canonical embedding of $\vs{\mathbf{s}}$ into $\mathrm{T}_{\Sigma}(\vs{\mathbf{s}})$. Sometimes, to abbreviate, we will write $\pi^{\mathbf{s}}_{i}$ instead of $(\pi^{\mathbf{s}}_{i})^{\mathbf{Tm}_{\mathrm{H}_{S}}(\Sigma)}$.
		\item For every $\mathbf{r},\mathbf{s}\in S^{\star}$ and $s\in S$, $\xi_{\mathbf{r},\mathbf{s},s}^{\mathbf{Tm}_{\mathrm{H}_{S}}(\Sigma)}$ is the mapping
		$$\xi_{\mathbf{r},\mathbf{s},s}^{\mathbf{Tm}_{\mathrm{H}_{S}}(\Sigma)}\nfunction
		{\mathrm{T}_{\Sigma}(\vs{\mathbf{s}})_{s} \times
			\mathrm{T}_{\Sigma}(\vs{\mathbf{r}})_{s_{0}} \times \cdots \times
			\mathrm{T}_{\Sigma}(\vs{\mathbf{r}})_{s_{\bb{\mathbf{s}}-1}}}
		{\mathrm{T}_{\Sigma}(\vs{\mathbf{r}})_{s}}
		{(P,(Q_{i})_{i\in\bb{\mathbf{s}}})}
		{\left(\left(\!\begin{smallmatrix}v^{\mathbf{s}}_{i}\\Q_{i}\end{smallmatrix}\!\right)_{i\in \bb{\mathbf{s}}}\right)^{\sharp}_{s}(P)}
		$$
		where, for $\left(\!\begin{smallmatrix}v^{\mathbf{s}}_{i}\\Q_{i}\end{smallmatrix}\!\right)_{i\in \bb{\mathbf{s}}}$, the $S$-sorted mapping from $\vs{\mathbf{s}}$ to $\mathrm{T}_{\Sigma}(\vs{\mathbf{r}})$ canonically associated to the family $(Q_{i})_{i\in\bb{\mathbf{s}}}$, $\left(\left(\!\begin{smallmatrix}v^{\mathbf{s}}_{i}\\Q_{i}\end{smallmatrix}\!\right)_{i\in \bb{\mathbf{s}}}\right)^{\sharp}$, also denoted by $\mathcal{Q}^{\sharp}$, is the unique homomorphism from $\mathbf{T}_{\Sigma}(\vs{\mathbf{s}})$ into $\mathbf{T}_{\Sigma}(\vs{\mathbf{r}})$ such that
		$$
		\left(\left(\!\begin{smallmatrix}v^{\mathbf{s}}_{i}\\Q_{i}\end{smallmatrix}\!\right)_{i\in \bb{\mathbf{s}}}\right)^{\sharp}\circ \eta^{\vs{\mathbf{s}}} = \left(\!\begin{smallmatrix}v^{\mathbf{s}}_{i}\\Q_{i}\end{smallmatrix}\!\right)_{i\in \bb{\mathbf{s}}}.
		$$
		Sometimes, to abbreviate, we will write $\xi_{\mathbf{r},\mathbf{s},s}$ instead of $\xi_{\mathbf{r},\mathbf{s},s}^{\mathbf{Tm}_{\mathrm{H}_{S}}(\Sigma)}$.
	\end{enumerate}
	Then $\mathbf{Tm}_{\mathrm{H}_{S}}(\Sigma)$ is a Hall algebra, the \emph{Hall algebra for} $(S,\Sigma)$.
\end{proposition}

The Hall algebra for $(S,\Sigma)$ is, up to isomorphism, the free Hall algebra on $\Sigma$, see~\cite{CVCLRM26}.

\begin{proposition}\label{PHallFree}
	Let $S$ be a set of sorts and $\Sigma$ an $S$-sorted signature, i.e., an $S^{\star}\times S$-sorted set. Then the Hall algebra $\mathbf{Tm}_{\mathrm{H}_{S}}(\Sigma)$ is isomorphic to $\mathbf{T}_{\mathrm{H}_{S}}(\Sigma)$, the free Hall algebra on $\Sigma$. More precisely, for the $S^{\star}\times S$-sorted mapping $\eta^{\Sigma}$ from $\Sigma$ to $\mathrm{Tm}_{\mathrm{H}_{S}}(\Sigma)$ defined, for every $(\mathbf{s},s)\in S^{\star}\times S$, as
	$$
	\eta^{\Sigma}_{\mathbf{s},s}\nfunction
	{\Sigma_{\mathbf{s},s}}
	{\mathrm{T}_{\Sigma}(\vs{\mathbf{s}})_{s}}
	{\sigma}
	{\sigma\left((v^{\mathbf{s}}_{j})_{j\in\bb{\mathbf{s}}}\right),}
	$$
	for every Hall algebra $\mathbf{A}$ and every $S^{\star}\times S$-sorted mapping $f$ from $\Sigma$ to $A$, there is a unique homomorphism $f^{\sharp}$ from $\mathbf{Tm}_{\mathrm{H}_{S}}(\Sigma)$ to $\mathbf{A}$ such that $f^{\sharp}\circ \eta^{\Sigma} = f$.
\end{proposition}

For a set of sorts $S$ and an $S$-sorted signature $\Sigma$, the derived operation symbols of $\Sigma$, i.e., the elements of the components $\mathrm{T}_{\Sigma}(\vs{\mathbf{s}})_{s}$ of $\mathrm{Tm}_{\mathrm{H}_{S}}(\Sigma)$, can be considered as the operation symbols of a new signature. The mappings $\varphi$ from a set of sorts $S$ to another $T$, together with the $S^{\star}\times S$-sorted mappings $d$ that assign to operation symbols of an $S$-sorted signature $\Sigma$ derived operation symbols of a $T$-sorted signature $\Lambda$, i.e., terms of the components $\mathrm{T}_{\Lambda}(\vs{\varphi^{\star}(\mathbf{s})})_{\varphi(s)}$ of $\mathrm{Tm}_{\mathrm{H}_{T}}(\Lambda)_{\varphi^{\star}\times\varphi}$, form a new class of morphisms between many-sorted signatures which we will call \emph{derivors}.

\begin{definition}\label{DDerivor}
	Let $\sig{\Sigma} = (S,\Sigma)$ and $\sig{\Lambda} = (T,\Lambda)$ be many-sorted signatures. A \emph{derivor} from $\sig{\Sigma}$ to $\sig{\Lambda}$ is an ordered pair $\mathbf{d} = (\varphi,d)$, where $\varphi\colon S\mor T$ and $d\colon\Sigma\mor\mathrm{Tm}_{\mathrm{H}_{T}}(\Lambda)_{\varphi^{\star}\times\varphi}$. Thus, for every $(\mathbf{s},s)\in S^{\star}\times S$,
	$$
	d_{\mathbf{s},s}\colon \Sigma_{\mathbf{s},s}\mor\mathrm{Tm}_{\mathrm{H}_{T}}(\Lambda)_{\varphi^{\star}(\mathbf{s}),\varphi(s)}=
	\mathrm{T}_{\Lambda}(\vs{\varphi^{\star}(\mathbf{s})})_{\varphi(s)}.
	$$
	Hence, for every $\sigma\in \Sigma_{\mathbf{s},s}$, $d_{\mathbf{s},s}(\sigma)\in \mathrm{T}_{\Lambda}(\vs{\varphi^{\star}(\mathbf{s})})_{\varphi(s)}$, i.e., $d_{\mathbf{s},s}(\sigma)$ is a term of sort $\varphi(s)$ with variables in $\vs{\varphi^{\star}(\mathbf{s})}$. We will write $d(\sigma)$ instead of $d_{\mathbf{s},s}(\sigma)$ when this is unlikely to cause confusion.
\end{definition}

Every mapping $\varphi\colon S\mor T$ determines a functor from the category $\mathsf{Alg}(\mathrm{H}_{T})$ to the category $\mathsf{Alg}(\mathrm{H}_{S})$, so that $\mathrm{Tm}_{\mathrm{H}_{T}}(\Lambda)_{\varphi^{\star}\times\varphi}$ is in its turn equipped with a structure of Hall algebra for $S$, which allows us to define the composition of derivors. In fact, the $S^{\star}\times S$-sorted mapping $h^{\varphi}$ from $\Sigma^{\mathrm{H}_{S}}$ to $\Sigma^{\mathrm{H}_{T}}_{\varphi^{\star}\times \varphi}$ that sends $\pi^{\mathbf{s}}_{i}$ to $\pi^{\varphi^{\star}(\mathbf{s})}_{i}$ and $\xi_{\mathbf{r},\mathbf{s},s}$ to $\xi_{\varphi^{\star}(\mathbf{r}),\varphi^{\star}(\mathbf{s}),\varphi(s)}$ determines a morphism of algebraic presentations $(\varphi^{\star}\times \varphi, h^{\varphi})$ from $(S^{\star}\times S,\Sigma^{\mathrm{H}_{S}},\mathcal{E}^{\mathrm{H}_{S}})$ to $(T^{\star}\times T,\Sigma^{\mathrm{H}_{T}},\mathcal{E}^{\mathrm{H}_{T}})$, and the morphisms of algebraic presentations determine functors in the opposite direction between the associated categories of algebras, see~\cite{CVCLRM26}.

If $\mathbf{d}\colon\sig{\Sigma}\mor\sig{\Lambda}$ is a derivor, then $d\colon \Sigma\mor \mathrm{Tm}_{\mathrm{H}_{T}}(\Lambda)_{\varphi^{\star}\times\varphi}$ determines a homomorphism of Hall algebras $d^{\sharp}\colon\mathbf{Tm}_{\mathrm{H}_{S}}(\Sigma)\mor
\mathbf{Tm}_{\mathrm{H}_{T}}(\Lambda)_{\varphi^{\star}\times\varphi}$. Thus, for every $(\mathbf{s},s)\in S^{\star}\times S$, $d^{\sharp}_{\mathbf{s},s}$ sends terms in $\mathrm{T}_{\Sigma}(\vs{\mathbf{s}})_{s}$ to terms in $\mathrm{T}_{\Lambda}(\vs{\varphi^{\star}(\mathbf{s})})_{\varphi(s)}$.

\begin{definition}\label{DDerCompId}
	Let $\mathbf{d}=(\varphi,d)\colon\sig{\Sigma}\mor\sig{\Lambda}$ and $\mathbf{e}=(\psi,e)\colon\sig{\Lambda}\mor\sig{\Omega}$ be derivors, where $\sig{\Omega}=(U,\Omega)$ is a many-sorted signature. Then $\mathbf{e}\circ \mathbf{d} = (\psi,e)\circ(\varphi,d)$, the \emph{composition} of $\mathbf{d}$ and $\mathbf{e}$, is the derivor $(\psi\circ\varphi,e^{\sharp}_{\varphi^{\star}\times\varphi}\circ d)$, where $e^{\sharp}$ is the canonical extension of $e$ to the free Hall algebra on $\Lambda$. For every many-sorted signature $\sig{\Sigma} = (S,\Sigma)$, the \emph{identity} at $(S,\Sigma)$ is $(\id^{S},\eta^{\Sigma})$. When this is unlikely to cause confusion, we write $e^{\sharp}\circ d$ instead of $e^{\sharp}_{\varphi^{\star}\times\varphi}\circ d$.
\end{definition}

The just stated definition allows us to form a category of many-sorted signatures whose morphisms are the derivors.

\begin{proposition}\label{PSigdCat}
	The many-sorted signatures together with the derivors constitute a category, denoted by $\mathsf{Sig}_{\mathfrak{d}}$.
\end{proposition}

In the following remark we describe the homomorphisms $d^{\sharp}$, the composition and the identities of derivors in terms of substitutions, which is the form in which we use them in this work.

\begin{remark}\label{RDerivorTerms}
Let $\mathbf{d}=(\varphi,d)$ be a derivor from $\sig{\Sigma}=(S,\Sigma)$ to $\sig{\Lambda}=(T,\Lambda)$ and $\mathbf{s}\in S^{\star}$. Then, by Propositions~\ref{PHallTm} and~\ref{PHallFree}, the mapping $d^{\sharp}$, which sends, for every $s'\in S$, the terms in $\mathrm{T}_{\Sigma}(\vs{\mathbf{s}})_{s'}$ to terms in $\mathrm{T}_{\Lambda}(\vs{\varphi^{\star}(\mathbf{s})})_{\varphi(s')}$, satisfies the following equalities:
\begin{enumerate}
\item[(Der1)] $d^{\sharp}(v^{\mathbf{s}}_{j})=v^{\varphi^{\star}(\mathbf{s})}_{j}$, for every $j\in\bb{\mathbf{s}}$;
\item[(Der2)] $d^{\sharp}(\sigma((R_{j})_{j\in\bb{\mathbf{r}}}))=(d^{\sharp}(R_{j}))^{\sharp}_{j\in\bb{\mathbf{r}}}(d(\sigma))$, for every $(\mathbf{r},s')\in S^{\star}\times S$, every $\sigma\in\Sigma_{\mathbf{r},s'}$ and every $(R_{j})_{j\in\bb{\mathbf{r}}}\in\mathrm{T}_{\Sigma}(\vs{\mathbf{s}})_{\mathbf{r}}$, i.e., $d(\sigma)$ with $d^{\sharp}(R_{j})$ substituted for $v^{\varphi^{\star}(\mathbf{r})}_{j}$.
\end{enumerate}
In particular, $d^{\sharp}(\sigma((v^{\mathbf{r}}_{j})_{j\in\bb{\mathbf{r}}}))=d(\sigma)$. Moreover:
\begin{enumerate}
\item[(Der3)] the identity derivor at $(S,\Sigma)$ is $(\id^{S},\eta^{\Sigma})$, where $\eta^{\Sigma}(\sigma)=\sigma((v^{\mathbf{s}}_{j})_{j\in\bb{\mathbf{s}}})$, for $\sigma\in\Sigma_{\mathbf{s},s}$, and $(\eta^{\Sigma})^{\sharp}$ is the identity; and the composition of $(\varphi,d)$ and a derivor $(\psi,e)$ from $(T,\Lambda)$ to a many-sorted signature $(U,\Omega)$ is the derivor $(\psi\circ\varphi,e^{\sharp}\circ d)$ from $(S,\Sigma)$ to $(U,\Omega)$; and
\item[(Der4)] the composition of derivors is associative and the identity derivors are neutral for it, i.e., for every derivor $(\chi,g)$ from $(U,\Omega)$ to a many-sorted signature $(W,\Xi)$, $(g^{\sharp}\circ e)^{\sharp}\circ d=g^{\sharp}\circ(e^{\sharp}\circ d)$, which follows from the equality $(g^{\sharp}\circ e)^{\sharp}=g^{\sharp}\circ e^{\sharp}$, and $(\eta^{\Lambda})^{\sharp}\circ d=d=d^{\sharp}\circ\eta^{\Sigma}$, see Proposition~\ref{PSigdCat}.
\end{enumerate}
Finally, for every $\mathbf{s}\in S^{\star}$ and every set of $\mathrm{E}$-equations $\mathcal{H}\subseteq\mathrm{Eeqt}(\Sigma)_{\vs{\mathbf{s}}}$, see Definition~\ref{DEEq} below, we let $d^{\sharp}[\mathcal{H}]\subseteq\mathrm{Eeqt}(\Lambda)_{\vs{\varphi^{\star}(\mathbf{s})}}$ stand for the set of $\mathrm{E}$-equations
$$
d^{\sharp}[\mathcal{H}]=\left\lbrace\left(d^{\sharp}(P),d^{\sharp}(Q)\right)\,\middle|\,(P,Q)\in\mathcal{H}\right\rbrace.
$$
\end{remark}

\subsection{Many-sorted partial algebras}\label{SPartialAlg}

We next recall the notions of many-sorted partial algebra and homomorphism between many-sorted partial algebras.

\begin{definition}\label{DPartialAlg}
Let $S$ be a set of sorts and $\Sigma$ an $S$-sorted signature. A \emph{structure of partial} $\Sigma$-\emph{algebra on} an $S$-sorted set $A$ is a family $(F_{\mathbf{s},s})_{(\mathbf{s},s)\in S^{\star}\times S}$, denoted by $F$, where, for $(\mathbf{s},s)\in S^{\star}\times S$, $F_{\mathbf{s},s}$ is a mapping from $\Sigma_{\mathbf{s},s}$ to $\mathrm{Hom}_{\mathrm{p}}(A_{\mathbf{s}},A_{s})$. If $\Dom(F_{\mathbf{s},s}(\sigma)) = A_{\mathbf{s}}$, then $F_{\mathbf{s},s}(\sigma)$ is a \emph{total operation} from $A_{\mathbf{s}}$ to $A_{s}$. If $\Dom(F_{\mathbf{s},s}(\sigma)) = \varnothing$, then $F_{\mathbf{s},s}(\sigma)$ is called a \emph{discrete operation} or an \emph{empty (partial) operation} from $A_{\mathbf{s}}$ to $A_{s}$. If $(\mathbf{s},s) = (\lambda,s)$, $\sigma\in \Sigma_{\lambda,s}$, and $\Dom(F_{\lambda,s}(\sigma))\neq\varnothing$, then $F_{\lambda,s}(\sigma)$ picks out an element of $A_{s}$. Therefore a partial operation $F_{\lambda,s}(\sigma)$ from $A_{\lambda}$ to $A_{s}$, also called a \emph{partial constant} $F_{\lambda,s}(\sigma)$ from $A_{\lambda}$ to $A_{s}$, is either discrete or distinguishes exactly an element of $A_{s}$. A \emph{many-sorted partial} $\Sigma$-\emph{algebra} (or, to abbreviate, \emph{partial $\Sigma$-algebra}) is a pair $(A,F)$, abbreviated to $\mathbf{A}$, where $A$ is an $S$-sorted set and $F$ a structure of partial $\Sigma$-algebra on $A$. For a pair $(\mathbf{s},s)\in S^{\star}\times S$ and a formal operation $\sigma\in \Sigma_{\mathbf{s},s}$, in order to simplify the notation, the partial operation $F_{\mathbf{s},s}(\sigma)$ from $A_{\mathbf{s}}$ to $A_{s}$ will be written as $F_{\sigma}$. In some cases, to avoid mistakes, we will denote by $F^{\mathbf{A}}$ the structure of partial $\Sigma$-algebra on $A$, and for $(\mathbf{s},s)\in S^{\star}\times S$ and $\sigma\in \Sigma_{\mathbf{s},s}$, by $F^{\mathbf{A}}_{\sigma}$, or simply by $\sigma^{\mathbf{A}}$, the corresponding partial operation.

A $\Sigma$-\emph{homomorphism} (or, to abbreviate, \emph{homomorphism}) from $\mathbf{A}$ to $\mathbf{B}$, where $\mathbf{B} = (B,G)$, is a triple $(\mathbf{A},f,\mathbf{B})$, abbreviated to $f\colon \mathbf{A}\mor \mathbf{B}$, where $f$ is an $S$-sorted mapping from $A$ to $B$ such that, for every $(\mathbf{s},s)\in S^{\star}\times S$, every $\sigma\in \Sigma_{\mathbf{s},s}$, and every $(a_{j})_{j\in \bb{\mathbf{s}}}\in A_{\mathbf{s}}$, if $(a_{j})_{j\in \bb{\mathbf{s}}}\in\Dom(\sigma^{\mathbf{A}})$, then $(f_{s_{j}}(a_{j}))_{j\in \bb{\mathbf{s}}}\in\Dom(\sigma^{\mathbf{B}})$ and
$
f_{s}(\sigma^{\mathbf{A}}((a_{j})_{j\in \bb{\mathbf{s}}})) = \sigma^{\mathbf{B}}((f_{s_{j}}(a_{j}))_{j\in \bb{\mathbf{s}}}).
$
For a partial $\Sigma$-algebra $\mathbf{A}$ the triple $(\mathbf{A},\id^{A},\mathbf{A})$, where $\id^{A}$ is the identity at $A$, which is a homomorphism from $\mathbf{A}$ to $\mathbf{A}$, will be denoted by $\id^{\mathbf{A}}$ and will be called the \emph{identity} homomorphism at $\mathbf{A}$. We will denote by $\mathsf{PAlg}(\Sigma)$ the category of partial $\Sigma$-algebras and homomorphisms and by $\mathrm{PAlg}(\Sigma)$ the set of objects of $\mathsf{PAlg}(\Sigma)$.
\end{definition}

Let us note that homomorphisms preserve, but need not reflect, definedness. The category $\mathsf{Alg}(\Sigma)$ is a full subcategory of the category $\mathsf{PAlg}(\Sigma)$. In the following proposition we recall a characterization of the isomorphisms of partial algebras.

\begin{proposition}\label{PPAlgIso}
Let $S$ be a set of sorts, $\Sigma$ an $S$-sorted signature and $f$ a homomorphism from a partial $\Sigma$-algebra $\mathbf{A}$ to a partial $\Sigma$-algebra $\mathbf{B}$. Then the following statements are equivalent:
\begin{enumerate}
\item $f$ is an isomorphism from $\mathbf{A}$ to $\mathbf{B}$, i.e., $f$ is a section and a retraction.
\item $f$ is bijective and $(\mathbf{B},f^{-1},\mathbf{A})$, abbreviated to $f^{-1}\colon \mathbf{B}\mor \mathbf{A}$, is also a homomorphism from $\mathbf{B}$ to $\mathbf{A}$.
\end{enumerate}
\end{proposition}

We next recall the product of a family of many-sorted partial algebras.

\begin{definition}\label{DPAlgProd}
Let $S$ be a set of sorts, $\Sigma$ an $S$-sorted signature and $(\mathbf{A}^{i})_{i\in I}$ a family of partial $\Sigma$-algebras. Then
\begin{enumerate}
\item The \emph{product} of $(\mathbf{A}^{i})_{i\in I}$, $\prod_{i\in I}\mathbf{A}^{i}$, is the partial $\Sigma$-algebra which has as $S$-sorted underlying set $\prod_{i\in I}A^{i}$, and where, for every $(\mathbf{s},s)\in S^{\star}\times S$ and every $\sigma\in \Sigma_{\mathbf{s},s}$, the structural operation $\sigma^{\prod_{i\in I}\mathbf{A}^{i}}$ is defined as follows:
      \begin{enumerate}
      \item $\Dom(\sigma^{\prod_{i\in I}\mathbf{A}^{i}}) = \{(a_{j})_{j\in \bb{\mathbf{s}}}\in
      (\prod_{i\in I}A^{i})_{\mathbf{s}}\mid \forall\,i\in I\,
      ((a_{j}(i))_{j\in \bb{\mathbf{s}}}\in \Dom(\sigma^{\mathbf{A}^{i}}))\}$.
      \item
      $$
      \sigma^{\prod_{i\in I}\mathbf{A}^{i}}
      \nfunction
      {\Dom(\sigma^{\prod_{i\in I}\mathbf{A}^{i}})}
      {\prod_{i\in I}A^{i}_{s}}
      {(a_{j})_{j\in \bb{\mathbf{s}}}}
      {(\sigma^{\mathbf{A}^{i}}((a_{j}(i))_{j\in \bb{\mathbf{s}}})_{i\in I}}
      $$
      \end{enumerate}
\item For every $i\in I$, the $i$-th \emph{canonical projection} is the homomorphism from $\prod_{i\in I}\mathbf{A}^{i}$ to $\mathbf{A}^{i}$ determined by the $S$-sorted mapping $\pr^{i}$ which, for every $s\in S$, is defined as follows
      $$
      \pr^{i}_{s}
      \nfunction
      {\prod_{i\in I}A^{i}_{s}}
      {A^{i}_{s}}
      {(a_{i})_{ i\in I}}
      {a_{i}}
      $$
\end{enumerate}
\end{definition}

\begin{proposition}\label{PPAlgProd}
Let $S$ be a set of sorts, $\Sigma$ an $S$-sorted signature and $(\mathbf{A}^{i})_{i\in I}$ a family of partial $\Sigma$-algebras. Then the pair $(\prod_{i\in I}\mathbf{A}^{i},(\pr^{i})_{i\in I})$ is a product in $\mathsf{PAlg}(\Sigma)$.
\end{proposition}

\subsection{Terms and equations}\label{STermsEq}

We next recall the partial term operations of a partial algebra. In~\cite{CVCLRM26} the value of a term at a valuation in a partial algebra is defined by means of the Schmidt construction, see~\cite{sch70}; for the purposes of this work, the following equivalent recursive description suffices.

\begin{definition}\label{DTermOp}
Let $S$ be a set of sorts, $\Sigma$ an $S$-sorted signature and $X$ an $S$-sorted set. For an $S$-sorted set $A$, we let $A_{X}=\mathrm{Hom}(X,A)$ stand for the set of the valuations of $X$ in $A$. For a partial $\Sigma$-algebra $\mathbf{A}$, $s\in S$ and $P\in\mathrm{T}_{\Sigma}(X)_{s}$, the \emph{partial term operation} $P^{\mathbf{A}}$ induced by $P$ is the partial mapping from $A_{X}=\mathrm{Hom}(X,A)$ to $A_{s}$ defined by recursion on $P$ as follows:
\begin{enumerate}
\item if $P=x$ is a variable, then $P^{\mathbf{A}}$ is total and sends $a$ to $a_{s}(x)$; and
\item if $P=\sigma((P_{j})_{j\in\bb{\mathbf{s}}})$, then $a\in\Dom(P^{\mathbf{A}})$ if, and only if, $a\in\Dom(P_{j}^{\mathbf{A}})$ for every $j\in\bb{\mathbf{s}}$ and $(P_{j}^{\mathbf{A}}(a))_{j\in\bb{\mathbf{s}}}\in\Dom(\sigma^{\mathbf{A}})$, in which case
$$
P^{\mathbf{A}}(a)=\sigma^{\mathbf{A}}\left((P_{j}^{\mathbf{A}}(a))_{j\in\bb{\mathbf{s}}}\right).
$$
\end{enumerate}
\end{definition}

For $(X,s)\in\boldsymbol{\mathcal{U}}^{S}\times S$, we next define the notion of existence equation of type $(X,s)$.

\begin{definition}\label{DEEq}
Let $S$ be a set of sorts and $\Sigma$ an $S$-sorted signature.
\begin{enumerate}
\item If $(X,s)\in\boldsymbol{\mathcal{U}}^{S}\times S$, then an \emph{existence equation} (briefly $\mathrm{E}$-\emph{equation} or $\mathrm{Eeqt}$) of type $(X,s)$ is a pair $(P,Q)$, also denoted by $P\eeq Q$, where $P$ and $Q$ are $\Sigma$-terms of type $(X,s)$, i.e., elements of $\mathrm{T}_{\Sigma}(X)_{s}$. If an $\mathrm{E}$-equation of type $(X,s)$ is such that $X$ is $S$-finite, i.e., for every $s\in S$, $X_{s}$ is finite, then we will call it $S$\emph{-finitary}, and if it is such that $X$ is finite, then we will call it \emph{finitary}.
\item We denote by $\mathrm{Eeqt}(\Sigma)_{X}$ the $S$-sorted set $(\mathrm{T}_{\Sigma}(X)_{s}^{2})_{s\in S}$ of all $\mathrm{E}$-equations with variables in an $S$-sorted set $X$. A \emph{set of $\mathrm{E}$-equations} with variables in $X$ is an $S$-sorted subset $\mathcal{E}\subseteq\mathrm{Eeqt}(\Sigma)_{X}$.
\end{enumerate}
\end{definition}

For $(X,s)\in\boldsymbol{\mathcal{U}}^{S}\times S$, we next define the ternary relation of satisfiability $\cdot \models^{\Sigma}_{X,s}\!\!\cdot\, [\cdot]$ between a partial $\Sigma$-algebra, an $\mathrm{E}$-equation of type $(X,s)$, and a valuation of $X$ in the underlying $S$-sorted set of the partial $\Sigma$-algebra, and the derived binary relation of validity.

\begin{definition}\label{DEEqSat}
Let $S$ be a set of sorts, $\Sigma$ an $S$-sorted signature, $\mathbf{A}$ a partial $\Sigma$-algebra, $(X,s)\in\boldsymbol{\mathcal{U}}^{S}\times S$, $(P,Q)\in \mathrm{Eeqt}(\Sigma)_{X,s}$, $\mathcal{E}\subseteq \mathrm{Eeqt}(\Sigma)_{X}$, and $a\in A_{X} = \mathrm{Hom}(X,A)$.
\begin{enumerate}
\item We will say that $\mathbf{A}$ \emph{satisfies the $\mathrm{E}$-equation} $(P,Q)$ \emph{with respect to the valuation} $a$, denoted by
      $$
      \mathbf{A}\models^{\Sigma}_{X,s}(P,Q) [a],
      $$
      if and only if $a\in\Dom(P^{\mathbf{A}})$, $a\in\Dom(Q^{\mathbf{A}})$, and $P^{\mathbf{A}}(a) = Q^{\mathbf{A}}(a)$, see Definition~\ref{DTermOp}.
\item We will say that $\mathbf{A}$ \emph{satisfies} $\mathcal{E}$ \emph{with respect to the valuation} $a$, denoted by $\mathbf{A}\models^{\Sigma}_{X}\mathcal{E}[a]$, if, for every $s\in S$ and every $(P,Q)\in\mathcal{E}_{s}$, $\mathbf{A}\models^{\Sigma}_{X,s}(P,Q)[a]$.
\item We will say that $(P,Q)$ is \emph{valid} in $\mathbf{A}$, denoted by $\mathbf{A}\models^{\Sigma}_{X,s}(P,Q)$, or by $\mathbf{A}\models(P,Q)$, when this is unlikely to cause confusion, if, for every $a\in A_{X}$, $\mathbf{A}\models^{\Sigma}_{X,s}(P,Q) [a]$.
\end{enumerate}
We omit the sub- and superscripts when this is unlikely to cause confusion.
\end{definition}

Let us note that $\mathbf{A}\models^{\Sigma}_{X,s}(P,P)[a]$ if, and only if, $a\in\Dom(P^{\mathbf{A}})$, i.e., the $\mathrm{E}$-equation $P\eeq P$ states that $P$ is defined.

For $(X,(\mathbf{s},s))$ in $\boldsymbol{\mathcal{U}}^{S}\times (S^{\star}\times S)$, we next define the notion of quasi-existence equation of type $(X,(\mathbf{s},s))$, its satisfaction and the varieties determined by sets of them.

\begin{definition}\label{DQEEq}
Let $S$ be a set of sorts and $\Sigma$ an $S$-sorted signature.
\begin{enumerate}
\item If $(X,(\mathbf{s},s))\in\boldsymbol{\mathcal{U}}^{S}\times (S^{\star}\times S)$, then a \emph{quasi-existence equation} (briefly $\QE$-\emph{equation} or $\mathrm{QEeqt}$) of type $(X,(\mathbf{s},s))$ is a pair
$$
\left(\left(
P_{j},Q_{j}
\right)_{j\in \bb{\mathbf{s}}},(P,Q)\right),
$$
also denoted by
$$
\textstyle
\bigwedge_{j\in\bb{\mathbf{s}}}(P_{j},Q_{j})\to (P,Q)\, \text{ or by }\,
\bigwedge_{j\in\bb{\mathbf{s}}}(P_{j}\eeq Q_{j})\to (P\eeq Q),
$$
where, for every $j\in \bb{\mathbf{s}}$, $(P_{j},Q_{j})\in\mathrm{T}_{\Sigma}(X)^{2}_{s_{j}}$ and $(P,Q)\in \mathrm{T}_{\Sigma}(X)^{2}_{s}$. The pairs $(P_{j},Q_{j})$, for $j\in \bb{\mathbf{s}}$, are called the \emph{premisses} and $(P,Q)$ the \emph{conclusion} of the $\QE$-equation. If a $\QE$-equation of type $(X,(\mathbf{s}, s))$ is such that $X$ is $S$-finite, then we will call it $S$-\emph{finitary}, and if it is such that $X$ is finite, then we will call it \emph{finitary}.
\item We denote by $\mathrm{QEeqt}(\Sigma)_{X}$ the $S^{\star}\times S$-sorted set of all $\QE$-equations with variables in $X$.
\item Let $\mathbf{A}$ be a partial $\Sigma$-algebra and $a\in A_{X}$. We will say that $\mathbf{A}$ \emph{satisfies} the $\QE$-equation $\bigwedge_{j\in\bb{\mathbf{s}}}(P_{j},Q_{j})\to (P,Q)$ \emph{with respect to the valuation} $a$, denoted by
      $$
      \textstyle
      \mathbf{A}\models^{\Sigma}_{X,\mathbf{s},s}\bigwedge_{j\in\bb{\mathbf{s}}}(P_{j},Q_{j})\to (P,Q) [a],
      $$
      if and only if, if, for every $j\in \bb{\mathbf{s}}$, $\mathbf{A}\models^{\Sigma}_{X,s_{j}} (P_{j},Q_{j})[a]$, then $\mathbf{A}\models^{\Sigma}_{X,s} (P,Q)[a]$. We will say that the $\QE$-equation \emph{holds} in $\mathbf{A}$, or that it is \emph{valid} in $\mathbf{A}$, if $\mathbf{A}$ satisfies it with respect to every valuation $a\in A_{X}$.
\item A \emph{$\QE$-specification} is a quadruple $(S,\Sigma,X,\mathcal{E})$ with $\mathcal{E}\subseteq\mathrm{QEeqt}(\Sigma)_{X}$. We will denote by $\mathrm{Mod}_{\Sigma}(\mathcal{E})$ the class of all partial $\Sigma$-algebras in which every $\QE$-equation of $\mathcal{E}$ holds, the \emph{models} of the specification. We will say that a class $\boldsymbol{\mathcal{K}}$ of partial $\Sigma$-algebras is a \emph{finitary} $\QE$-\emph{variety} if there exists a set $\mathcal{E}$ of finitary $\QE$-equations such that $\boldsymbol{\mathcal{K}} = \mathrm{Mod}_{\Sigma}(\mathcal{E})$. For a $\QE$-specification $(S,\Sigma,X,\mathcal{E})$, we also denote $\mathrm{Mod}_{\Sigma}(\mathcal{E})$ by $\mathcal{V}(\mathcal{E})$ and call it the \emph{$\QE$-variety} determined by the specification.
\end{enumerate}
\end{definition}

Every $\QE$-variety is closed under products, see~\cite{pb86}. In fact, if $(\mathbf{A}^{i})_{i\in I}$ is a family of partial $\Sigma$-algebras in which a $\QE$-equation holds and $a$ is a valuation in $\prod_{i\in I}\mathbf{A}^{i}$ that satisfies its premisses, then, by Definitions~\ref{DPAlgProd} and~\ref{DTermOp}, for every $i\in I$, $\pr^{i}\circ a$ satisfies the premisses in $\mathbf{A}^{i}$, hence also the conclusion, and, again by Definitions~\ref{DPAlgProd} and~\ref{DTermOp}, $a$ satisfies the conclusion in $\prod_{i\in I}\mathbf{A}^{i}$.

The following two facts on the partial term operations are proved by induction on terms.

\begin{remark}\label{RSubst}
Let $S$ be a set of sorts, $\Sigma$ an $S$-sorted signature, $\mathbf{A}$ a partial $\Sigma$-algebra, $X$ and $Y$ $S$-sorted sets, $f\colon Y\mor\mathrm{T}_{\Sigma}(X)$ an $S$-sorted mapping and $a\in A_{X}$ a valuation such that, for every $s\in S$ and every $y\in Y_{s}$, $a\in\Dom(f_{s}(y)^{\mathbf{A}})$. Let $f^{\mathbf{A}}(a)\in A_{Y}$ be the valuation that sends $y\in Y_{s}$ to $f_{s}(y)^{\mathbf{A}}(a)$. Then, for every $s\in S$ and every $R\in\mathrm{T}_{\Sigma}(Y)_{s}$, $a\in\Dom(f^{\sharp}_{s}(R)^{\mathbf{A}})$ if, and only if, $f^{\mathbf{A}}(a)\in\Dom(R^{\mathbf{A}})$, and then $f^{\sharp}_{s}(R)^{\mathbf{A}}(a)=R^{\mathbf{A}}(f^{\mathbf{A}}(a))$. Consequently, for every $\mathcal{E}\subseteq\mathrm{Eeqt}(\Sigma)_{Y}$, $\mathbf{A}\models f^{\sharp}[\mathcal{E}][a]$ if, and only if, $\mathbf{A}\models\mathcal{E}[f^{\mathbf{A}}(a)]$.
\end{remark}

\begin{remark}\label{RHomTerm}
Let $S$ be a set of sorts, $\Sigma$ an $S$-sorted signature, $h\colon\mathbf{A}\mor\mathbf{B}$ a homomorphism of partial $\Sigma$-algebras, $X$ an $S$-sorted set, $a\in A_{X}$ a valuation, $s\in S$ and $P\in\mathrm{T}_{\Sigma}(X)_{s}$. If $a\in\Dom(P^{\mathbf{A}})$, then $h\circ a\in\Dom(P^{\mathbf{B}})$ and $P^{\mathbf{B}}(h\circ a)=h_{s}(P^{\mathbf{A}}(a))$. Consequently, for every $\mathcal{E}\subseteq\mathrm{Eeqt}(\Sigma)_{X}$, if $\mathbf{A}\models\mathcal{E}[a]$, then $\mathbf{B}\models\mathcal{E}[h\circ a]$.
\end{remark}

\subsection{Adjunctions, corestrictions and the Grothendieck construction}\label{SAdjunctions}

We recall the notions of adjunction and of the Grothendieck construction, which we use in the rest of this work. We let ``$\ast$'' and ``$\circ$'' stand for the horizontal and vertical composition of natural transformations, respectively; we denote by $\Id_{\mathsf{X}}$ the identity functor at $\mathsf{X}$; and we denote by $\id_{G}$ the identity natural transformation at $G\colon \mathsf{A}\mor \mathsf{X}$. Moreover, for categories $\mathsf{A}$ and $\mathsf{X}$, functors $F,G\colon\mathsf{A}\mor\mathsf{X}$ and a natural transformation $\theta$ from $F$ to $G$, we write $\theta=(\theta_{a})_{a\in\mathrm{Ob}(\mathsf{A})}$, where, for every object $a$ of $\mathsf{A}$, $\theta_{a}\colon F(a)\mor G(a)$ is the component of $\theta$ at $a$.

\begin{definition}\label{DAdj}
An \emph{adjointness} is a $6$-tuple $(\mathsf{A},\mathsf{X},G,F,\eta,\varepsilon)$ where $\mathsf{A}$ and $\mathsf{X}$ are categories, $G$ a functor from $\mathsf{A}$ to $\mathsf{X}$, $F$ a functor from $\mathsf{X}$ to $\mathsf{A}$, $\eta$ a natural transformation from $\Id_{\mathsf{X}}$ to $G\circ F$ and $\varepsilon$ a natural transformation from $F\circ G$ to $\Id_{\mathsf{A}}$ subject to the so called triangular identities
$$
(G\ast \varepsilon)\circ (\eta\ast G) = \id_{G} \text{ and }
(\varepsilon\ast F)\circ (F\ast\eta) = \id_{F}.
$$
Thus, for every object $a$ of $\mathsf{A}$, the first equation asserts that the morphism $\eta_{G(a)}$ from $G(a)$ to $G(F(G(a)))$ when composed with the morphism $G(\varepsilon_{a})$ from $G(F(G(a)))$ to $G(a)$ is $\id_{G(a)}$, the identity morphism at $G(a)$, and, for every object $x$ of $\mathsf{X}$, the second equation asserts that the morphism $F(\eta_{x})$ from $F(x)$ to $F(G(F(x)))$ when composed with the morphism $\varepsilon_{F(x)}$ from $F(G(F(x)))$ to $F(x)$ is $\id_{F(x)}$. If $(\mathsf{A},\mathsf{X},G,F,\eta,\varepsilon)$ is an adjointness, then we will say that $F$ is a \emph{left adjoint} of $G$, that $G$ is a \emph{right adjoint} of $F$, that $\eta$ is the \emph{unit} of the adjunction and that $\varepsilon$ is the \emph{counit} of the adjunction. We will also speak of an \emph{adjunction from} $\mathsf{X}$ \emph{to} $\mathsf{A}$, meaning by it a $4$-tuple $(G,F,\eta,\varepsilon)$ such that $(\mathsf{A},\mathsf{X},G,F,\eta,\varepsilon)$ is an adjointness, and we will write $F\dashv G$ to indicate that $F$ is a left adjoint of $G$.
\end{definition}

We next recall the restrictions and the corestrictions of a functor to full subcategories.

\begin{definition}\label{DCorestriction}
Let $F$ be a functor from a category $\mathsf{A}$ to a category $\mathsf{X}$, $\mathsf{A}'$ a full subcategory of $\mathsf{A}$ and $\mathsf{X}'$ a full subcategory of $\mathsf{X}$. The \emph{restriction} of $F$ to $\mathsf{A}'$ is the functor $F\circ\mathrm{In}_{\mathsf{A}'}$ from $\mathsf{A}'$ to $\mathsf{X}$, where $\mathrm{In}_{\mathsf{A}'}$ is the inclusion functor of $\mathsf{A}'$ into $\mathsf{A}$. We will say that $F$ \emph{corestricts} to $\mathsf{X}'$ if, for every object $a$ of $\mathsf{A}$, $F(a)$ is an object of $\mathsf{X}'$. In this case, since $\mathsf{X}'$ is full, there exists a unique functor $F'$ from $\mathsf{A}$ to $\mathsf{X}'$ such that $\mathrm{In}_{\mathsf{X}'}\circ F'=F$, the \emph{corestriction} of $F$ to $\mathsf{X}'$. When this is unlikely to cause confusion, we will denote the corestriction of the restriction of $F$ to $\mathsf{A}'$, if it exists, by $F$ itself, i.e., we will write $F\colon\mathsf{A}'\mor\mathsf{X}'$.
\end{definition}

We next recall the split indexed categories and the Grothendieck construction, which assigns to a split indexed category a split fibration, see~\cite{Gro71} and, for split fibrations, \cite{BW12}, pp.~327--329. We follow the notation of~\cite{JY21}, Chapter~10, but add the domain of $F$ as a superscript to indicate the contravariance of $F$.

\begin{definition}\label{DSplitIndexed}
Let $\mathsf{I}$ be a category. A \emph{split} $\mathsf{I}$-\emph{indexed category} is a contravariant functor $F$ from $\mathsf{I}$ to $\mathsf{Cat}$. Given an object $i\in \mathrm{Ob}(\mathsf{I})$, we write $\mathsf{F}_{i}$ for the category $F(i)$, and given a morphism $\varphi\in \mathrm{Hom}_{\mathsf{I}}(i,j)$, we write $F_{\varphi}$ for the functor $F(\varphi)\colon \mathsf{F}_{j}\mor \mathsf{F}_{i}$. A \emph{split indexed category} is an ordered pair $(\mathsf{I},F)$ in which $\mathsf{I}$ is a category and $F$ a split $\mathsf{I}$-indexed category.
\end{definition}

\begin{definition}\label{DGConst}
Let $(\mathsf{I},F)$ be a split indexed category. Then the \emph{Grothen\-dieck construction} at $(\mathsf{I},F)$ is $(\int^{\mathsf{I}}F,\pi_{F})$, where $\int^{\mathsf{I}}F$ is the category defined as follows:
\begin{enumerate}
\item $\mathrm{Ob}(\int^{\mathsf{I}}F) = \bigcup_{i\in\mathrm{Ob}(\mathsf{I})}(\{i\}\times\mathrm{Ob}(\mathsf{F}_{i}))$.
\item For every $(i,x), (j,y)\in \mathrm{Ob}(\int^{\mathsf{I}}F)$, $\mathrm{Hom}_{\int^{\mathsf{I}}F}((i,x),(j,y))$ is the set of all ordered pairs $(\varphi,f)$, where $\varphi\in \mathrm{Hom}_{\mathsf{I}}(i,j)$ and $f\in \mathrm{Hom}_{\mathsf{F}_{i}}(x,F_{\varphi}(y))$.
\item For every $(i,x)\in \mathrm{Ob}(\int^{\mathsf{I}}F)$, the identity morphism at $(i,x)$ is $(\id_{i},\id_{x})$.
\item For every $(i,x), (j,y), (k,z)\in \mathrm{Ob}(\int^{\mathsf{I}}F)$, every $(\varphi,f)\colon (i,x)\mor (j,y)$, and every $(\psi,g)\colon (j,y)\mor (k,z)$, the composite morphism $(\psi,g)\circ (\varphi,f)$ from $(i,x)$ to $(k,z)$ is
$$
(\psi,g)\circ (\varphi,f) = (\psi\circ \varphi,F_{\varphi}(g)\circ f)\colon (i,x)\mor (k,z).
$$
Notice that $f\colon x\mor F_{\varphi}(y)$ and $g\colon y\mor F_{\psi}(z)$, hence, since $F_{\varphi}$ is a functor from $\mathsf{F}_{j}$ to $\mathsf{F}_{i}$, $F_{\varphi}(g)\colon F_{\varphi}(y)\mor F_{\varphi}(F_{\psi}(z))$. Therefore $F_{\varphi}(g)\circ f\colon x\mor F_{\varphi}(F_{\psi}(z)) = F_{\psi\circ\varphi}(z)$.
\end{enumerate}
And $\pi_{F}$ is the functor from $\int^{\mathsf{I}}F$ to $\mathsf{I}$ that sends an object $(i,x)$ of $\int^{\mathsf{I}}F$ to $i$ and a morphism $(\varphi,f)\colon (i,x)\mor (j,y)$ of $\int^{\mathsf{I}}F$ to $\varphi\colon i\mor j$. The functor $\pi_{F}$ is a split fibration, the \emph{canonical split fibration} from $\int^{\mathsf{I}}F$ to $\mathsf{I}$ (determined by $(\mathsf{I},F)$).
\end{definition}

In Section~\ref{SGenHepDerivors} we will need the Grothendieck construction for assignments that preserve identities, but only preserve composites up to a coherent family of natural transformations, which need not be isomorphisms. We next state the corresponding notion and the corresponding construction.

\begin{definition}\label{DLaxIndexed}
Let $\mathsf{I}$ be a category. A \emph{normal lax} $\mathsf{I}$-\emph{indexed category} is an ordered pair $(F,\kappa)$ in which
\begin{enumerate}
\item $F$ assigns to every object $i$ of $\mathsf{I}$ a category $\mathsf{F}_{i}$, and to every morphism $\varphi\colon i\mor j$ of $\mathsf{I}$ a functor $F_{\varphi}\colon\mathsf{F}_{j}\mor\mathsf{F}_{i}$, in such a way that, for every object $i$ of $\mathsf{I}$, $F_{\id_{i}}=\Id_{\mathsf{F}_{i}}$; and
\item $\kappa$ assigns to every pair of morphisms $\varphi\colon i\mor j$ and $\psi\colon j\mor k$ of $\mathsf{I}$ a natural transformation
$$
\kappa^{\varphi,\psi}=\left(\kappa^{\varphi,\psi}_{z}\right)_{z\in\mathrm{Ob}(\mathsf{F}_{k})}\colon F_{\varphi}\circ F_{\psi}\Longrightarrow F_{\psi\circ\varphi},
$$
\end{enumerate}
subject to the following conditions:
\begin{enumerate}
\item[(L1)] for every morphism $\varphi\colon i\mor j$ of $\mathsf{I}$, $\kappa^{\id_{i},\varphi}=\id_{F_{\varphi}}=\kappa^{\varphi,\id_{j}}$; and
\item[(L2)] for every triple of morphisms $\varphi\colon i\mor j$, $\psi\colon j\mor k$ and $\chi\colon k\mor l$ of $\mathsf{I}$, the following equality holds
$$
\kappa^{\varphi,\chi\circ\psi}\circ\left(F_{\varphi}\ast\kappa^{\psi,\chi}\right)
=
\kappa^{\psi\circ\varphi,\chi}\circ\left(\kappa^{\varphi,\psi}\ast F_{\chi}\right),
$$
i.e., for every object $w$ of $\mathsf{F}_{l}$, $\kappa^{\varphi,\chi\circ\psi}_{w}\circ F_{\varphi}(\kappa^{\psi,\chi}_{w})=\kappa^{\psi\circ\varphi,\chi}_{w}\circ\kappa^{\varphi,\psi}_{F_{\chi}(w)}$.
\end{enumerate}
Let us note that, if every $\kappa^{\varphi,\psi}$ is an identity, then $F$ is a split $\mathsf{I}$-indexed category, see Definition~\ref{DSplitIndexed}.
\end{definition}

\begin{definition}\label{DLaxGConst}
Let $\mathsf{I}$ be a category and $(F,\kappa)$ a normal lax $\mathsf{I}$-indexed category. Then $\int^{\mathsf{I}}(F,\kappa)$ stands for the following data:
\begin{enumerate}
\item the objects are the ordered pairs $(i,x)$ with $i\in\mathrm{Ob}(\mathsf{I})$ and $x\in\mathrm{Ob}(\mathsf{F}_{i})$;
\item the morphisms from $(i,x)$ to $(j,y)$ are the ordered pairs $(\varphi,f)$, where $\varphi\in \mathrm{Hom}_{\mathsf{I}}(i,j)$ and $f\in \mathrm{Hom}_{\mathsf{F}_{i}}(x,F_{\varphi}(y))$;
\item the identity morphism at $(i,x)$ is $(\id_{i},\id_{x})$; and
\item for every $(\varphi,f)\colon (i,x)\mor (j,y)$ and every $(\psi,g)\colon (j,y)\mor (k,z)$, the composite morphism $(\psi,g)\circ (\varphi,f)$ from $(i,x)$ to $(k,z)$ is
$$
(\psi,g)\circ (\varphi,f) = \left(\psi\circ \varphi,\kappa^{\varphi,\psi}_{z}\circ F_{\varphi}(g)\circ f\right).
$$
\end{enumerate}
Moreover, $\pi_{(F,\kappa)}$ stands for the assignment that sends an object $(i,x)$ to $i$ and a morphism $(\varphi,f)$ to $\varphi$.
\end{definition}

Let us note that, since $f\colon x\mor F_{\varphi}(y)$ and $g\colon y\mor F_{\psi}(z)$, the composite $\kappa^{\varphi,\psi}_{z}\circ F_{\varphi}(g)\circ f$ is a morphism of $\mathsf{F}_{i}$ from $x$ to $F_{\psi\circ\varphi}(z)$, so that the above definition is correct. If every $\kappa^{\varphi,\psi}$ is an identity, then $\int^{\mathsf{I}}(F,\kappa)$ is the Grothendieck construction of Definition~\ref{DGConst}.

\begin{proposition}\label{PLaxGConst}
Let $\mathsf{I}$ be a category and $(F,\kappa)$ a normal lax $\mathsf{I}$-indexed category. Then $\int^{\mathsf{I}}(F,\kappa)$ is a category and $\pi_{(F,\kappa)}$ is a functor from $\int^{\mathsf{I}}(F,\kappa)$ to $\mathsf{I}$.
\end{proposition}
\begin{proof}
\textsf{(Identities)} Let $(\varphi,f)\colon(i,x)\mor(j,y)$ be a morphism. Then the following chains of equalities hold
\allowdisplaybreaks
\begin{align*}
(\varphi,f)\circ(\id_{i},\id_{x})
&=
\left(\varphi\circ\id_{i},\kappa^{\id_{i},\varphi}_{y}\circ F_{\id_{i}}(f)\circ\id_{x}\right)
\tag{1}
\\&=
\left(\varphi,\id_{F_{\varphi}(y)}\circ f\circ\id_{x}\right)
\tag{2}
\\&=
(\varphi,f),
\tag{3}
\\
(\id_{j},\id_{y})\circ(\varphi,f)
&=
\left(\id_{j}\circ\varphi,\kappa^{\varphi,\id_{j}}_{y}\circ F_{\varphi}(\id_{y})\circ f\right)
\tag{4}
\\&=
\left(\varphi,\id_{F_{\varphi}(y)}\circ\id_{F_{\varphi}(y)}\circ f\right)
\tag{5}
\\&=
(\varphi,f).
\tag{6}
\end{align*}

The first and the fourth equalities unravel Definition~\ref{DLaxGConst}; the second and the fifth equalities follow from condition~(L1) of Definition~\ref{DLaxIndexed}, from the fact that $F_{\id_{i}}=\Id_{\mathsf{F}_{i}}$ and from the fact that $F_{\varphi}$ is a functor; finally, the third and the last equalities follow from the fact that $\mathsf{F}_{i}$ is a category.

\textsf{(Associativity)} Let $(\varphi,f)\colon(i,x)\mor(j,y)$, $(\psi,g)\colon(j,y)\mor(k,z)$ and $(\chi,h)\colon(k,z)\mor(l,w)$ be morphisms. Both $(\chi,h)\circ((\psi,g)\circ(\varphi,f))$ and $((\chi,h)\circ(\psi,g))\circ(\varphi,f)$ have $\chi\circ\psi\circ\varphi$ as first coordinate. For the second coordinates, the following chain of equalities holds
\allowdisplaybreaks
\begin{align*}
\kappa^{\psi\circ\varphi,\chi}_{w}\circ F_{\psi\circ\varphi}(h)\circ\kappa^{\varphi,\psi}_{z}\circ F_{\varphi}(g)\circ f
&=
\kappa^{\psi\circ\varphi,\chi}_{w}\circ\kappa^{\varphi,\psi}_{F_{\chi}(w)}\circ F_{\varphi}(F_{\psi}(h))\circ F_{\varphi}(g)\circ f
\tag{1}
\\&=
\kappa^{\varphi,\chi\circ\psi}_{w}\circ F_{\varphi}(\kappa^{\psi,\chi}_{w})\circ F_{\varphi}(F_{\psi}(h))\circ F_{\varphi}(g)\circ f
\tag{2}
\\&=
\kappa^{\varphi,\chi\circ\psi}_{w}\circ F_{\varphi}\left(\kappa^{\psi,\chi}_{w}\circ F_{\psi}(h)\circ g\right)\circ f.
\tag{3}
\end{align*}

The first equality follows from the naturality of $\kappa^{\varphi,\psi}$, applied to the morphism $h\colon z\mor F_{\chi}(w)$ of $\mathsf{F}_{k}$; the second equality follows from condition~(L2) of Definition~\ref{DLaxIndexed}; finally, the last equality follows from the fact that $F_{\varphi}$ is a functor. By Definition~\ref{DLaxGConst}, the left-hand side of the above chain is the second coordinate of $(\chi,h)\circ((\psi,g)\circ(\varphi,f))$ and its right-hand side is the second coordinate of $((\chi,h)\circ(\psi,g))\circ(\varphi,f)$.

\textsf{(The projection)} By Definition~\ref{DLaxGConst}, $\pi_{(F,\kappa)}$ preserves identities and composites.

This completes the proof.
\end{proof}

We will call $(\int^{\mathsf{I}}(F,\kappa),\pi_{(F,\kappa)})$ the \emph{Grothendieck construction} at $(\mathsf{I},(F,\kappa))$.

% \section{Derivors}\label{SDerivors}
%
% In this subsection, after defining the variety of Hall algebras, we introduce the notion of derivor between many-sorted signatures. This will allow us, using the homomorphisms between Hall algebras, to obtain $\mathsf{Sig}_{\mathfrak{d}}$, the category of many-sorted signatures and derivors. Next, after defining a suitable contravariant functor from $\mathsf{Sig}_{\mathfrak{d}}$ to $\mathsf{Cat}$, the category of $\boldsymbol{\mathcal{U}}$-locally small categories and functors, we obtain, by means of the Grothendieck construction, $\mathsf{Alg}_{\mathfrak{d}}$, the category with objects the ordered pairs $((S,\Sigma),(A,F))$, with $(S,\Sigma)$ a many-sorted signature and $(A,F)$ a $\Sigma$-algebra, and morphisms from $((S,\Sigma),(A,F))$ to  $((T,\Lambda),(B,G))$ the ordered pairs $((\varphi,d),f)$, with $(\varphi,d)$ a derivor from $(S,\Sigma)$ to $(T,\Lambda)$ and $f$ a homomorphism of $\Sigma$-algebras from $(A,F)$ to $(B_{\varphi},G^{(\varphi,d)})$, a canonically derived algebra of $(B,G)$. 
%
% Before defining the notion of Hall algebra, we recall that
% \begin{enumerate}
% 	\item a finitary specification is an ordered triple $(S,\Sigma,\mathcal{E})$, where $S$ is a set of sorts, $\Sigma$ an $S$-sorted signature, and $\mathcal{E}\subseteq\mathrm{Eq}_{V^{S}}(\Sigma)=
% 	(\mathrm{T}_{\Sigma}(X)_{s}^{2})_{(X,s)\in \mathrm{Sub}_{\mathrm{f}}(V^{S})\times S}$, i.e., a set of finitary  $\Sigma$-equations, where $V^{S}$, the $S$-sorted set of the variables, is a fixed $S$-countably infinite $S$-sorted set; and that
% 	\item a $\Sigma$-algebra $\mathbf{A}$ is a $(S,\Sigma,\mathcal{E})$-algebra if $\mathbf{A}\models^{\Sigma} \mathcal{E}$, i.e., if, for every $(X,s)\in \mathrm{Sub}_{\mathrm{f}}(V^{S})\times S$ and every $(P,Q)\in \mathcal{E}_{X,s}$, $\mathbf{A}$ is a model of $(P,Q)$, in symbols $\mathbf{A}\models^{\Sigma}_{X,s} (P,Q)$, which in turn means that, for every $a\in \mathrm{Hom}(X,A)$, $a^{\sharp}_{s}(P) = a^{\sharp}_{s}(Q)$.
% \end{enumerate}
%
% \begin{definition}
% 	Let $S$ be a set of sorts and let $V^{\mathrm{H}_{S}} = (V^{\mathrm{H}_{S}}_{\mathbf{s},s})_{(\mathbf{s},s)\in S^{\star}\times S}$ be a fixed $S^{\star}\times S$-sorted set of variables with a countably infinity of variables in each component, say $V^{\mathrm{H}_{S}}_{\mathbf{s},s} = \{\,v^{\mathbf{s},s}_{n}\mid n\in \mathbb{N}\,\}$ (which is isomorphic to 
% 	$\{(\mathbf{s},s)\}\times\mathbb{N}$), for every $(\mathbf{s},s)\in S^{\star}\times S$.  A \emph{Hall algebra for} $S$ is a $\mathrm{H}_{S} = (S^{\star}\times S,\Sigma^{\mathrm{H}_{S}},\mathcal{E}^{\mathrm{H}_{S}})$-algebra, where $\Sigma^{\mathrm{H}_{S}}$ is the $S^{\star}\times S$-sorted signature, i.e., the $(S^{\star}\times S)^{\star}\times (S^{\star}\times S)$-sorted set, defined as follows:
% 	\begin{enumerate}
% 		\item[$\mathrm{HS}_{1}$.] For every $\mathbf{s}\in S^{\star}$ and $i\in\bb{\mathbf{s}}$,
% 		$$
% 		\pi^{\mathbf{s}}_{i}\colon\lambda\mor(\mathbf{s},s_{i}),
% 		$$
% 		where $\bb{\mathbf{s}}$ is the \emph{length} of the word $\mathbf{s}$ and $\lambda$ the
% 		\emph{empty word} in the underlying set of the free monoid on
% 		$S^{\star}\times S$.
%
% 		\item[$\mathrm{HS}_{2}$.] For every $\mathbf{r}$, $\mathbf{s}\in S^{\star}$ and $s\in S$,
% 		$$
% 		\xi_{\mathbf{r},\mathbf{s},s}\colon ((\mathbf{s},s),(\mathbf{r},s_{0}),\ldots,(\mathbf{r},s_{\bb{\mathbf{s}}-1}))\mor(\mathbf{r},s);
% 		$$
% 	\end{enumerate}
% 	while $\mathcal{E}^{\mathrm{H}_{S}}$ is the sub-$(S^{\star}\times
% 	S)^{\star}\times (S^{\star}\times S)$-sorted set of
% 	$\mathrm{Eq}(\Sigma^{\mathrm{H}_{S}})$, where
% 	$$
% 	\mathrm{Eq}\left(\Sigma^{\mathrm{H}_{S}}\right) =
% 	\left(\mathrm{T}_{\Sigma^{\mathrm{H}_{S}}}(\vs{\ol{\mathbf{s}}})_{(\mathbf{r},s)}^{2}\right)_{(\ol{\mathbf{s}},(\mathbf{r},s))\in
% 		(S^{\star}\times S)^{\star}\times (S^{\star}\times S)},
% 	$$
% 	defined as follows:
% 	\begin{enumerate}
% 		\item[$\mathrm{H}_{1}$.] \emph{Projection}.
% 		For every $\mathbf{r}$, $\mathbf{s}\in S^{\star}$ and $i\in\bb{\mathbf{s}}$, the equation
% 		$$
% 		\xi_{\mathbf{r},\mathbf{s},s_{i}}(\pi^{\mathbf{s}}_{i},v^{\mathbf{r},s_{0}}_{0},\ldots,
% 		v^{\mathbf{r},s_{\bb{\mathbf{s}}-1}}_{\bb{\mathbf{s}}-1})=v^{\mathbf{r},s_{i}}_{i}
% 		$$
% 		of type
% 		$(((\mathbf{r},s_{0}),\ldots,(\mathbf{r},s_{\bb{\mathbf{s}}-1})),(\mathbf{r},s_{i})).$
% 		\item[$\mathrm{H}_{2}$.] \emph{Identity}.
% 		For every $\mathbf{r}\in S^{\star}$ and $j\in \bb{\mathbf{r}}$, the equation
% 		$$
% 		\xi_{\mathbf{r},\mathbf{r},r_{j}}(v^{\mathbf{r},r_{j}}_{j},\pi^{\mathbf{r}}_{0},\ldots,\pi^{\mathbf{r}}_{\bb{\mathbf{r}}-1})=
% 		v^{\mathbf{r},r_{j}}_{j}
% 		$$
% 		of type
% 		$(((\mathbf{r},r_{j})),(\mathbf{r},r_{j})).$
% 		\item[$\mathrm{H}_{3}$.] \emph{Associativity}.
% 		For every $\mathbf{q}$, $\mathbf{r}$, $\mathbf{s}\in S^{\star}$ and $s\in S$, the
% 		equation
% 		\begin{multline*}
% 			\xi_{\mathbf{q},\mathbf{r},s}(
% 			\xi_{\mathbf{r},\mathbf{s},s}(v^{\mathbf{s},s}_{0},v^{\mathbf{r},s_{0}}_{1},\ldots,
% 			v^{\mathbf{r},s_{\bb{\mathbf{s}}-1}}_{\bb{\mathbf{s}}}),
% 			v^{\mathbf{q},r_{0}}_{\bb{\mathbf{s}}+1},
% 			\ldots,v^{\mathbf{q},r_{\bb{\mathbf{r}}-1}}_{\bb{\mathbf{s}}+\bb{\mathbf{r}}}) = \\
% 			\xi_{\mathbf{q},\mathbf{s},s}(v^{\mathbf{s},s}_{0},
% 			\xi_{\mathbf{q},\mathbf{r},s_{0}}(v^{\mathbf{r},s_{0}}_{1},v^{\mathbf{q},r_{0}}_{\bb{\mathbf{s}}+1},
% 			\ldots,v^{\mathbf{q},r_{\bb{\mathbf{r}}-1}}_{\bb{\mathbf{s}}+\bb{\mathbf{r}}}),
% 			\\
% 			\ldots, 
% 			\xi_{\mathbf{q},\mathbf{r},s_{\bb{\mathbf{s}}-1}}(v^{\mathbf{r},s_{\bb{\mathbf{s}}-1}}_{\bb{\mathbf{s}}},
% 			v^{\mathbf{q},r_{0}}_{\bb{\mathbf{s}}+1},
% 			\ldots,v^{\mathbf{q},r_{\bb{\mathbf{r}}-1}}_{\bb{\mathbf{s}}+\bb{\mathbf{r}}}))
% 		\end{multline*}
% 		of type
% 		$(((\mathbf{s},s),(\mathbf{r},s_{0}),\ldots,(\mathbf{r},s_{\bb{\mathbf{s}}-1}),
% 		(\mathbf{q},r_{0}),\ldots,(\mathbf{q},r_{\bb{\mathbf{r}}-1})),(\mathbf{q},s)).$
% 	\end{enumerate}
%
% 	We call the formal constants $\pi^{\mathbf{s}}_{i}$ \emph{projections}, and the formal operations $\xi_{\mathbf{r},\mathbf{s},s}$ \emph{substitution operators}. Furthermore, we will denote by $\mathsf{Alg}(\mathrm{H}_{S})$ the category of Hall algebras for $S$ and homomorphisms between Hall algebras.  Since $\mathsf{Alg}(\mathrm{H}_{S})$ is a variety, the forgetful functor $\mathrm{G}_{\mathrm{H}_{S}}$ from $\mathsf{Alg}(\mathrm{H}_{S})$ to $\mathsf{Set}^{S^{\star}\times S}$ has a left adjoint $\mathbf{T}_{\mathrm{H}_{S}}$, situation denoted by $\mathbf{T}_{\mathrm{H}_{S}}\dashv \mathrm{G}_{\mathrm{H}_{S}}$, or diagrammatically by
% 	$$
% 	\xymatrix@=5pc{ \mathsf{Alg}(\mathrm{H}_{S})
% 		\ar@<1.5ex>[r]^{\mathrm{G}_{\mathrm{H}_{S}}} \ar@{}[r]|{\uadj} &
% 		\mathsf{Set}^{S^{\star}\times S}
% 		\ar@<1.5ex>[l]^{\mathbf{T}_{\mathrm{H}_{S}}} }
% 	$$
% 	which assigns to an $S^{\star}\times S$-sorted set $\Sigma$ the corresponding free Hall algebra
% 	$\mathbf{T}_{\mathrm{H}_{S}}(\Sigma)$.
% \end{definition}
%
% \begin{remark}
% 	The category $\mathsf{Alg}(\mathrm{H}_{S})$, of Hall algebras for $S$, is an algebraic version of many-sorted clone theory.
% \end{remark}
%
% \begin{remark}
% 	From $\mathrm{H}_{3}$, for $\mathbf{s}=\lambda$, the empty word on $S$, we get the invariance of constant functions axiom in~\cite{gm85}: For every $\mathbf{q}$, $\mathbf{r}\in S^{\star}$ and $s\in S$, we have the equation
% 	$$
% 	\xi_{\mathbf{q},\mathbf{r},s}(\xi_{\mathbf{r},\lambda,s}(v^{\lambda,s}_{0}),v^{\mathbf{q},r_{0}}_{1},
% 	\ldots,v^{\mathbf{q},r_{\bb{\mathbf{r}}-1}}_{\bb{\mathbf{r}}})=
% 	\xi_{\mathbf{q},\lambda,s}(v^{\lambda,s}_{0})
% 	$$
% 	of type $(((\lambda,s),(\mathbf{q},r_{0}),\ldots,(\mathbf{q},r_{\bb{\mathbf{r}}-1})),(\mathbf{q},s))$.
% \end{remark}
%
% For every $S$-sorted set $A$, $\mathrm{O}_{\mathrm{H}_{S}}(A) = (\mathrm{Hom}(A_{\mathbf{s}},A_{s}))_{(\mathbf{s},s)\in S^{\star}\times S}$, the $S^{\star}\times S$-sorted set of operation for $A$, is naturally equipped with a structure of Hall algebra, as stated in the following proposition, if we realize the projections as the true projections and the substitution operators as the generalized composition of mappings.
%
%
%
%
%
% \begin{proposition}
% 	Let $A$ be an $S$-sorted set and $\mathbf{O}_{\mathrm{H}_{S}}(A)$ the $\Sigma^{\mathrm{H}_{S}}$-algebra with underlying many-sorted set $\mathrm{O}_{\mathrm{H}_{S}}(A)$ and algebraic structure defined as follows:
% 	\begin{enumerate}
% 		\item For every $\mathbf{s}\in S^{\star}$ and $i\in\bb{\mathbf{s}}$,
% 		$(\pi^{\mathbf{s}}_{i})^{\mathbf{O}_{\mathrm{H}_{S}}(A)} =
% 		\mathrm{pr}^{A}_{\mathbf{s},i}\colon A_{\mathbf{s}}\mor A_{s_{i}}$.
%
% 		\item For every $\mathbf{r},\mathbf{s}\in S^{\star}$ and $s\in S$,
% 		$\xi_{\mathbf{r},\mathbf{s},s}^{\mathbf{O}_{\mathrm{H}_{S}}(A)}$ is defined, for every $f\in
% 		A_{s}^{ A_{\mathbf{s}} }$ and $g\in A^{A_{\mathbf{r}}}_{\mathbf{s}}$, as
% 		$\xi_{\mathbf{r},\mathbf{s},s}^{\mathbf{O}_{\mathrm{H}_{S}}(A)}(f,g_{0},\ldots,g_{\bb{\mathbf{s}}-1})
% 		= f\circ\langle g_{i}\rangle_{i\in\bb{\mathbf{s}} }$, where $\langle
% 		g_{i}\rangle_{i\in\bb{\mathbf{s}}}$ is the unique mapping from $A_{\mathbf{r}}$ to
% 		$A_{\mathbf{s}}$ such that, for every $i\in \bb{\mathbf{s}}$, we have that
% 		$$
% 		\mathrm{pr}^{A}_{\mathbf{s},i}\circ\langle g_{i}\rangle_{i\in\bb{\mathbf{s}}} = g_{i}.
% 		$$
% 	\end{enumerate}
% 	Then $\mathbf{O}_{\mathrm{H}_{S}}(A)$ is a Hall algebra, the \emph{Hall algebra for} $(S,A)$.
% \end{proposition}
%
% \begin{remark}
% 	The closed sets of the Hall algebra $\mathbf{O}_{\mathrm{H}_{S}}(A)$ for $(S,A)$ are precisely the clones of (many-sorted) operations on the $S$-sorted set $A$. On the other hand, every $\Sigma$-algebra $\mathbf{A}$ has associated a Hall algebra. In fact, it suffices to consider $\mathbf{O}_{\mathrm{H}_{S}}(A)$, denoted by $\mathbf{O}_{\mathrm{H}_{S}}(\mathbf{A})$. Moreover, the finitary derived operations on $\mathbf{A}$ is a subalgebra of the Hall algebra $\mathbf{O}_{\mathrm{H}_{S}}(\mathbf{A})$.
% \end{remark}
%
%
%
% For every $S$-sorted signature $\Sigma$, $\mathrm{Tm}_{\mathrm{H}_{S}}(\Sigma) = (\mathrm{T}_{\Sigma}(\vs{\mathbf{s}})_{s})_{(\mathbf{s},s)\in S^{\star}\times S}$ is also e\-quipped with a structure of Hall algebra that formalizes the concept of substitution as stated in the following proposition.
%
% \begin{proposition}
% 	Let $\Sigma$ be an $S$-sorted signature and $\mathbf{Tm}_{\mathrm{H}_{S}}(\Sigma)$ the $\Sigma^{\mathrm{H}_{S}}$-algebra
% 	with underlying many-sorted set $\mathrm{Tm}_{\mathrm{H}_{S}}(\Sigma)$ and algebraic structure defined as follows:
% 	\begin{enumerate}
% 		\item For every $\mathbf{s}\in S^{\star}$ and $i\in\bb{\mathbf{s}}$,
% 		$(\pi^{\mathbf{s}}_{i})^{\mathbf{Tm}_{\mathrm{H}_{S}}(\Sigma)}$
% 		is the image under $\eta^{\vs{\mathbf{s}}}_{s_{i}}$ of
% 		the variable $v_{i}^{s_{i}}$, where $\eta^{\vs{\mathbf{s}}} =
% 		(\eta^{\vs{\mathbf{s}}}_{s})_{s\in S}$ is the canonical embedding
% 		of $\vs{\mathbf{s}}$ into $\mathrm{T}_{\Sigma}(\vs{\mathbf{s}})$. Sometimes,
% 		to abbreviate, we will write $\pi^{\mathbf{s}}_{i}$ instead of
% 		$(\pi^{\mathbf{s}}_{i})^{\mathbf{Tm}_{\mathrm{H}_{S}}(\Sigma)}$.
%
% 		\item For every $\mathbf{r},\mathbf{s}\in S^{\star}$ and $s\in S$,
% 		$\xi_{\mathbf{r},\mathbf{s},s}^{\mathbf{Tm}_{\mathrm{H}_{S}}(\Sigma)}$
% 		is the mapping
% 		$$\xi_{\mathbf{r},\mathbf{s},s}^{\mathbf{Tm}_{\mathrm{H}_{S}}(\Sigma)}\nfunction
% 		{\mathrm{T}_{\Sigma}(\vs{\mathbf{s}})_{s} \times
% 			\mathrm{T}_{\Sigma}(\vs{\mathbf{r}})_{s_{0}} \times \cdots \times
% 			\mathrm{T}_{\Sigma}(\vs{\mathbf{r}})_{s_{\bb{\mathbf{s}}-1}}}
% 		{\mathrm{T}_{\Sigma}(\vs{\mathbf{r}})_{s}}
% 		{(P,(Q_{i})_{i\in\bb{\mathbf{s}}})}
% 		{\left(\left(\!\begin{smallmatrix}v^{s_{i}}_{i}\\Q_{i}\end{smallmatrix}\!\right)_{i\in \bb{\mathbf{s}}}\right)^{\sharp}_{s}(P)}
% 		$$
% 		where, for $\left(\!\begin{smallmatrix}v^{s_{i}}_{i}\\Q_{i}\end{smallmatrix}\!\right)_{i\in \bb{\mathbf{s}}}$, the $S$-sorted mapping from $\vs{\mathbf{s}}$
% 		to $\mathrm{T}_{\Sigma}(\vs{\mathbf{r}})$ canonically associated to the family
% 		$(Q_{i})_{i\in\bb{\mathbf{s}}}$, $\left(\left(\!\begin{smallmatrix}v^{s_{i}}_{i}\\Q_{i}\end{smallmatrix}\!\right)_{i\in \bb{\mathbf{s}}}\right)^{\sharp}$, also denoted by $\mathcal{Q}^{\sharp}$, is the
% 		unique homomorphism from $\mathbf{T}_{\Sigma}(\vs{\mathbf{s}})$ into
% 		$\mathbf{T}_{\Sigma}(\vs{\mathbf{r}})$ such that
% 		$$
% 		\left(\left(\!\begin{smallmatrix}v^{s_{i}}_{i}\\Q_{i}\end{smallmatrix}\!\right)_{i\in \bb{\mathbf{s}}}\right)^{\sharp}\circ \eta^{\vs{\mathbf{s}}} = \left(\!\begin{smallmatrix}v^{s_{i}}_{i}\\Q_{i}\end{smallmatrix}\!\right)_{i\in \bb{\mathbf{s}}}.
% 		$$
% 		Sometimes, to abbreviate, we will write $\xi_{\mathbf{r},\mathbf{s},s}$ instead of
% 		$\xi_{\mathbf{r},\mathbf{s},s}^{\mathbf{Tm}_{\mathrm{H}_{S}}(\Sigma)}$.
% 	\end{enumerate}
%
% 	Then $\mathbf{Tm}_{\mathrm{H}_{S}}(\Sigma)$ is a Hall algebra, the \emph{Hall algebra for} $(S,\Sigma)$.
% \end{proposition}
%
% \begin{remark}
% 	For every $\Sigma$-algebra $\mathbf{A}$ there exists a homomorphism of Hall algebras $\mathrm{d}^{\mathbf{A}} = (\mathrm{d}^{\mathbf{A}}_{\mathbf{\mathbf{s}},s})_{(\mathbf{\mathbf{s}},s)\in S^{\star}\times S}$ from the Hall algebra $\mathbf{Tm}_{\mathrm{H}_{S}}(\Sigma)$ to the Hall algebra $\mathbf{O}_{\mathrm{H}_{S}}(\mathbf{A})$ and its image is $\mathbf{D}_{\mathrm{H}_{S}}(\mathbf{A})$, the Hall subalgebra of $\mathbf{O}_{\mathrm{H}_{S}}(\mathbf{A})$ of the finitary derived operations on $\mathbf{A}$, which, for every $(\mathbf{\mathbf{s}},s)\in S^{\star}\times S$, is such that 
% 	$\mathrm{D}_{\mathrm{H}_{S}}(\mathbf{A})_{\mathbf{\mathbf{s}},s} = 
% 	\mathrm{D}_{\mathbf{\mathbf{s}}}(\mathbf{A})_{s}$, where $\mathrm{D}_{\mathbf{\mathbf{s}}}(\mathbf{A})_{s}$ is the $s$-th coordinate of the underlying $S$-sorted of 
% 	$\mathbf{D}_{\mathbf{\mathbf{s}}}(\mathbf{A})$, the $\Sigma$-algebra of the $\mathbf{\mathbf{s}}$-\emph{ary} derived operations on $\mathbf{A}$, introduced in Definition~\ref{wdop}.
% \end{remark}
%
% %\begin{proposition}
% %Let $\Sigma$ be an $S$-sorted signature. Then $\mathbf{Tm}_{\mathrm{H}_{S}}(\Sigma)^{\wp}$ is a Hall algebra.
% %\end{proposition}
%
% %recall that the underlying $S^{\star}\times S$-sorted of $\mathbf{Tm}_{\mathrm{H}_{S}}(\Sigma)^{\wp}$ is $(\mathrm{Sub}(\mathrm{T}_{\Sigma}(\vs{w})_{s}))_{(w,s)\in S^{\star}\times S}$
%
% Our next goal is to prove that, for every $S^{\star}\times S$-sorted set $\Sigma$, $\mathbf{T}_{\mathrm{H}_{S}}(\Sigma)$, the free Hall algebra on $\Sigma$, is isomorphic to $\mathbf{Tm}_{\mathrm{H}_{S}}(\Sigma)$. We remark that the existence of this isomorphism is interesting because it enables us to get a more tractable description of the terms in $\mathbf{T}_{\mathrm{H}_{S}}(\Sigma)$.
% %We remark that the existence of this isomorphism is interesting because it enables us, on the one hand, to get a more tractable description of the terms in $\mathbf{T}_{\mathrm{H}_{S}}(\Sigma)$, and, on the other hand, as we will show afterwards, to state, for every $\Sigma$-algebra $\mathbf{A}$, taking into account the adjunction $\mathbf{T}_{\mathrm{H}_{S}}\dashv \mathrm{G}_{\mathrm{H}_{S}}$, the existence of a homomorphism of Hall algebras
% %$\Tr^{\mathbf{A}}$ from $\mathbf{Tm}_{\mathrm{H}_{S}}(\Sigma)$ to $\mathbf{O}_{\mathrm{H}_{S}}(\mathbf{A}) = \mathbf{O}_{\mathrm{H}_{S}}(A)$ such that $\mathrm{Th}_{\Sigma}(\mathbf{A})$, the finitary $\Sigma$-equational theory determined by $\mathbf{A}$, is
% %precisely $\mathrm{Ker}(\Tr^{\mathbf{A}})$, the kernel of the homomorphism $\Tr^{\mathbf{A}}$.
%
% To attain the goal just stated we begin by defining, for a Hall algebra $\mathbf{A}$, an $S$-sorted signature $\Sigma$, an
% $S^{\star}\times S$-mapping $f\colon\Sigma\mor A$, and a word $\mathbf{r}\in S^{\star}$, the concept of derived $\Sigma$-algebra of
% $\mathbf{A}$ for $(f,\mathbf{r})$, since it will be used afterwards in the proof of the isomorphism between
% $\mathbf{T}_{\mathrm{H}_{S}}(\Sigma)$ and $\mathbf{Tm}_{\mathrm{H}_{S}}(\Sigma)$.
%
%
%
% \begin{definition}
% 	Let $\mathbf{A}$ be a Hall algebra and $\Sigma$ an $S$-sorted signature. Then, for every $f\colon\Sigma\mor A$ and $\mathbf{r}\in
% 	S^{\star}$, $\mathbf{A}^{f,\mathbf{r}}$, the \emph{derived} $\Sigma$-\emph{algebra} \emph{of} $\mathbf{A}$ \emph{for} $(f,\mathbf{r})$,
% 	is the $\Sigma$-algebra with underlying $S$-sorted set $A^{f,\mathbf{r}} = (A_{\mathbf{r},s})_{s\in S}$ and algebraic structure $F^{f,\mathbf{r}}$, defined, for every $(\mathbf{s},s)\in S^{\star}\times S$, as
% 	$$
% 	F^{f,\mathbf{r}}_{\mathbf{s},s} \nfunction
% 	{\Sigma_{\mathbf{s},s}}{\mathrm{O}_{\mathbf{s}}(A^{f,\mathbf{r}})_{s}} {\sigma}{\nfunction
% 		{\prod_{i\in \bb{\mathbf{s}}}A_{\mathbf{u},s_{i}}}
% 		{A_{\mathbf{u},s}}
% 		{(a_{0},\ldots,a_{\bb{\mathbf{s}}-1})}
% 		{\xi_{\mathbf{r},\mathbf{s},s}^{\mathbf{A}}(f_{(\mathbf{s},s)}(\sigma),a_{0},\ldots,a_{\bb{\mathbf{s}}-1})}
% 	}
% 	$$
% 	where $\mathrm{O}_{\mathbf{s}}(A^{f,\mathbf{r}})_{s} = A_{\mathbf{r},s}^{\prod_{i\in \bb{\mathbf{s}}}A_{\mathbf{r},s_{i}}}$.
%
% 	Furthermore, we will denote by $p^{\mathbf{r}}$ the $S$-sorted mapping from $\vs{\mathbf{r}}$ to $A^{f,\mathbf{r}}$ defined, for every $s\in S$ and $i\in\bb{\mathbf{r}}$, as $p^{\mathbf{r}}_{s}(v^{s}_{i}) = (\pi^{\mathbf{r}}_{i})^{\mathbf{A}}$, and by $(p^{\mathbf{r}})^{\sharp}$ the unique homomorphism from
% 	$\mathbf{T}_{\Sigma}(\vs{\mathbf{r}})$ to $\mathbf{A}^{f,\mathbf{r}}$ such that $(p^{\mathbf{r}})^{\sharp}\circ \eta^{\vs{\mathbf{r}}} = p^{\mathbf{r}}$.
% \end{definition}
%
% \begin{remark}
% 	For a $\Sigma$-algebra $\mathbf{B}=(B,G)$, we have that $G\colon\Sigma\mor\mathrm{O}_{\mathrm{H}_{S}}(B)$ and
% 	$\mathbf{B}\iso\mathbf{O}_{\mathrm{H}_{S}}(B)^{G,\lambda}$, where $\lambda$ is the empty word on $S$.  Besides, for every $\mathbf{r}\in S^{\star}$, we have that $\mathbf{B}^{B_{\mathbf{r}}}$, the direct $B_{\mathbf{r}}$-power of $\mathbf{B}$, is isomorphic to $\mathbf{O}_{\mathrm{H}_{S}}(B)^{G,\mathbf{r}}$.
% \end{remark}
%
%
%
% \begin{lemma}\label{L:aux}
% 	Let $\Sigma$ be an $S$-sorted signature, $\mathbf{A}$ a Hall algebra, $f\colon\Sigma\mor A$ and $\mathbf{q}\in S^{\star}$.  Then, for
% 	every $(\mathbf{r},s)\in S^{\star}\times S$, $P\in \mathrm{T}_{\Sigma}(\vs{\mathbf{r}})_{s}$ and $a\in \prod_{i\in \bb{\mathbf{r}}}A_{\mathbf{q},r_{i}}$, we have that
% 	$$
% 	P^{\mathbf{A}^{f,\mathbf{q}}}(a_{0},\ldots,a_{\bb{\mathbf{r}}-1}) =
% 	\xi_{\mathbf{q},\mathbf{r},s}^{\mathbf{A}}((p^{\mathbf{r}})^{\sharp}_{s}(P),a_{0},\ldots,a_{\bb{\mathbf{r}}-1}).
% 	$$
% \end{lemma}
%
% \begin{proof}
% 	By algebraic induction on the complexity of $P$.  If $P$ is a variable $v_{i}^{s}$, with $i\in\bb{\mathbf{r}}$, then
% 	\begin{align*}
% 		v_{i}^{s,\mathbf{A}^{f,\mathbf{q}}}(a_{0},\ldots,a_{\bb{\mathbf{r}}-1})\ &=
% 		a^{\sharp}_{r_{i}}(v^{s}_{i})\\
% 		&=
% 		a_{i} \\
% 		&=
% 		\xi_{\mathbf{q},\mathbf{r},s}^{\mathbf{A}}((\pi^{\mathbf{r}}_{i})^{\mathbf{A}},a_{0},\ldots,a_{\bb{\mathbf{r}}-1} )
% 		\tag{by $\mathrm{H}_{1}$} \\
% 		&=
% 		\xi_{\mathbf{q},\mathbf{r},s}^{\mathbf{A}}((p^{w})^{\sharp}_{s}(v^{s}_{i}),a_{0},\ldots,a_{\bb{\mathbf{r}}-1} ).
% 	\end{align*}
% 	Let us assume that $P = \sigma(Q_{0},\ldots,Q_{\bb{x}-1})$, with $\sigma\in \Sigma_{\mathbf{s},s}$ and that, for every $j\in\bb{x}$, $Q_{j}\in \mathrm{T}_{\Sigma}(\vs{\mathbf{r}})_{s_{j}}$ fulfils the induction hypothesis. Then we have that
% 	\begin{flushleft}
% 		$(\sigma(Q_{0},\ldots,Q_{\bb{\mathbf{s}}-1}))^{\mathbf{A}^{f,\mathbf{q}}}
% 		(a_{0},\ldots,a_{\bb{\mathbf{r}}-1})$
% 		\allowdisplaybreaks
% 		\begin{align*}
% 			&=
% 			\sigma^{\mathbf{A}^{f,\mathbf{q}}}
% 			(
% 			Q_{0}^{\mathbf{A}^{f,\mathbf{q}}}(a_{0},\ldots,a_{\bb{\mathbf{r}}-1}),
% 			\ldots,
% 			Q_{\bb{\mathbf{s}}-1}^{\mathbf{A}^{f,\mathbf{q}}}(a_{0},\ldots,a_{\bb{\mathbf{r}}-1})
% 			)\\
% 			&=
% 			\xi_{\mathbf{q},\mathbf{s},s}^{\mathbf{A}}
% 			(
% 			f(\sigma),
% 			Q_{0}^{\mathbf{A}^{f,\mathbf{q}}}(a_{0},\ldots,a_{\bb{\mathbf{r}}-1}),
% 			\ldots,
% 			Q_{\bb{\mathbf{s}}-1}^{\mathbf{A}^{f,\mathbf{q}}}(a_{0},\ldots,a_{\bb{\mathbf{r}}-1})
% 			)\\
% 			&=
% 			\xi_{\mathbf{q},\mathbf{s},s}^{\mathbf{A}}
% 			(
% 			f(\sigma),
% 			\xi_{\mathbf{q},\mathbf{r},s_{0}}^{\mathbf{A}}
% 			((p^{\mathbf{r}})^{\sharp}_{s_{0}}(Q_{0}),a_{0},\ldots,a_{\bb{\mathbf{r}}-1}),
% 			\ldots, 
% 			\\&\qquad\qquad\qquad\qquad\qquad\qquad
% 			\xi_{\mathbf{q},\mathbf{r},s_{\bb{\mathbf{s}}-1}}^{\mathbf{A}}
% 			((p^{\mathbf{r}})^{\sharp}_{s_{\bb{\mathbf{s}}-1}}(Q_{\bb{\mathbf{s}}-1}),a_{0},\ldots,a_{\bb{\mathbf{r}}-1})
% 			) \tag{by ind.}
% 			\\&=
% 			\xi_{\mathbf{q},\mathbf{r},s}^{\mathbf{A}}
% 			(
% 			\xi_{\mathbf{r},\mathbf{s},s}^{\mathbf{A}}
% 			(f(\sigma),
% 			(p^{\mathbf{r}})^{\sharp}_{s_{0}}(Q_{0}),
% 			\ldots,
% 			(p^{\mathbf{r}})^{\sharp}_{s_{\bb{\mathbf{s}}-1}}(Q_{\bb{\mathbf{s}}-1})
% 			),
% 			a_{0},
% 			\ldots,
% 			a_{\bb{\mathbf{r}}-1}
% 			) \tag{by $\mathrm{H}_{3}$}   \\
% 			&=
% 			\xi_{\mathbf{q},\mathbf{r},s}^{\mathbf{A}}
% 			(
% 			\sigma^{\mathbf{A}_{\mathbf{r}}}
% 			(
% 			(p^{\mathbf{r}})^{\sharp}_{s_{0}}(Q_{0}),
% 			\ldots,
% 			(p^{\mathbf{r}})^{\sharp}_{s_{\bb{\mathbf{s}}-1}}(Q_{\bb{\mathbf{s}}-1})
% 			),
% 			a_{0},
% 			\ldots,
% 			a_{\bb{\mathbf{r}}-1}
% 			) \\
% 			&=
% 			\xi_{\mathbf{q},\mathbf{r},s}^{\mathbf{A}}
% 			(
% 			(p^{\mathbf{r}})^{\sharp}_{s}(\sigma,Q_{0},\ldots,Q_{\bb{\mathbf{s}}-1}),
% 			a_{0},
% 			\ldots,
% 			a_{\bb{\mathbf{r}}-1}
% 			) \\
% 			&=
% 			\xi_{\mathbf{q},\mathbf{r},s}^{\mathbf{A}}
% 			(
% 			(p^{\mathbf{r}})^{\sharp}_{s}(P),
% 			a_{0},
% 			\ldots,
% 			a_{\bb{\mathbf{r}}-1}
% 			).
% 		\end{align*}
% 	\end{flushleft}
% 	This completes the proof.
% \end{proof}
%
% Next we prove that, for every $S^{\star}\times S$-sorted set $\Sigma$, the Hall algebra for $(S,\Sigma)$ is isomorphic to the free Hall algebra on $\Sigma$.
%
%
% \begin{proposition}\label{iso:FrH-TerH}
% 	Let $\Sigma$ be an $S$-sorted signature, i.e., an $S^{\star}\times S$-sorted set.  Then the Hall algebra $\mathbf{Tm}_{\mathrm{H}_{S}}(\Sigma)$ is isomorphic to $\mathbf{T}_{\mathrm{H}_{S}}(\Sigma)$, the free Hall algebra on $\Sigma$.
% \end{proposition}
%
% \begin{proof}
% 	It is enough to prove that $\mathbf{Tm}_{\mathrm{H}_{S}}(\Sigma)$ has the universal property of $\mathbf{T}_{\mathrm{H}_{S}}(\Sigma)$. Therefore we have to specify an $S^{\star}\times S$-sorted mapping $\eta^{\Sigma}$ from $\Sigma$ to $\mathrm{Tm}_{\mathrm{H}_{S}}(\Sigma)$ such that, for every Hall algebra $\mathbf{A}$ and $S^{\star}\times S$-sorted mapping $f$ from $\Sigma$ to $A$, there is a unique homomorphism $\widehat{f}$ from $\mathbf{Tm}_{\mathrm{H}_{S}}(\Sigma)$ to $\mathbf{A}$ such that
% 	$\widehat{f}\circ \eta^{\Sigma} = f$.  Let $\eta^{\Sigma}$ be the $S^{\star}\times S$-sorted mapping defined, for every
% 	$(\mathbf{s},s)\in S^{\star}\times S$, as%
% 	$$
% 	\eta^{\Sigma}_{\mathbf{s},s}\nfunction
% 	{\Sigma_{\mathbf{s},s}}
% 	{\mathrm{T}_{\Sigma}(\vs{\mathbf{s}})_{s}}
% 	{\sigma}
% 	{\sigma\left(v^{s_{0}}_{0},\ldots,v^{s_{\bb{\mathbf{s}}-1}}_{\bb{\mathbf{s}}-1}\right)}
% 	$$
% 	Let $\mathbf{A}$ be a Hall algebra, $f\colon\Sigma\mor A$ an $S^{\star}\times S$-sorted mapping and $\widehat{f}$ the $S^{\star}\times S$-sorted mapping from $\mathrm{Tm}_{\mathrm{H}_{S}}(\Sigma)$ to $A$ defined, for every $(\mathbf{s},s)\in S^{\star}\times S$, as $\widehat{f}_{\mathbf{s},s} = (p^{\mathbf{s}})^{\sharp}_{s}$, where, we recall, $(p^{\mathbf{s}})^{\sharp}$ is the unique homomorphism from $\mathbf{T}_{\Sigma}(\vs{\mathbf{s}})$ to $\mathbf{A}^{f,\mathbf{s}}$ such that $(p^{\mathbf{s}})^{\sharp}\circ \eta^{\vs{\mathbf{s}}} = p^{\mathbf{s}}$. Then $\widehat{f}$ is a homomorphism of Hall algebras, because, on the one hand, for $\mathbf{s}\in S^{\star}$ and $i\in\bb{\mathbf{s}}$ we have that%
% 	\begin{align*}
% 		\widehat{f}_{\mathbf{s},s_{i}}((\pi^{\mathbf{s}}_{i})^{\mathbf{Tm}_{\mathrm{H}_{S}}(\Sigma)})
% 		&=
% 		\widehat{f}_{\mathbf{s},s_{i}}(v^{s_{i}}_{i}) \\
% 		&=
% 		p^{\mathbf{s}}_{s_{i}}(v^{s_{i}}_{i}) \\
% 		&=
% 		(\pi^{\mathbf{s}}_{i})^{\mathbf{A}},
% 	\end{align*}
% 	and, on the other hand, for $P\in \mathrm{T}_{\Sigma}(\vs{\mathbf{s}})_{s}$ and
% 	$\mathcal{Q} = (Q_{i})_{i\in \bb{\mathbf{s}}}\in \mathrm{T}_{\Sigma}(\vs{\mathbf{r}})_{\mathbf{s}}$ we have that %
% 	\begin{flushleft}
% 		$\widehat{f}_{\mathbf{r},s}(
% 		\xi_{\mathbf{r},\mathbf{s},s}^{\mathbf{Tm}_{\mathrm{H}_{S}}(\Sigma)}
% 		(P,Q_{0},\ldots,Q_{\bb{\mathbf{s}}-1})
% 		)$
% 		\begin{align*}
% 			\quad    &=
% 			(p^{\mathbf{r}})^{\sharp}_{s}(
% 			\mathcal{Q}^{\sharp}_{s}(P)
% 			) \\
% 			&=
% 			((p^{\mathbf{r}})^{\sharp}\circ \mathcal{Q})^{\sharp}_{s} (P)
% 			\tag{$(p^{\mathbf{r}})^{\sharp}\circ\mathcal{Q}^{\sharp} =
% 				((p^{\mathbf{r}})^{\sharp}\circ \mathcal{Q})^{\sharp}$}
% 			\\
% 			&=
% 			P^{\mathbf{A}^{f,\mathbf{r}}}
% 			((p^{\mathbf{r}})^{\sharp}_{s_{0}}(Q_{0}),\ldots,
% 			(p^{\mathbf{r}})^{\sharp}_{s_{\bb{\mathbf{s}}-1}}(Q_{\bb{\mathbf{s}}-1})
% 			) \\
% 			&=
% 			\xi_{\mathbf{r},\mathbf{s},s}^{\mathbf{A}}
% 			((p^{\mathbf{r}})^{\sharp}_{s}(P),
% 			(p^{\mathbf{r}})^{\sharp}_{s_{0}}(Q_{0}),\ldots,
% 			(p^{\mathbf{r}})^{\sharp}_{s_{\bb{\mathbf{s}}-1}}(Q_{\bb{\mathbf{s}}-1})
% 			)
% 			\tag{by Lemma~\ref{L:aux}}
% 			\\
% 			&=
% 			\xi_{\mathbf{r},\mathbf{s},s}^{\mathbf{A}}
% 			( \widehat{f}_{\mathbf{s},s}(P),
% 			\widehat{f}_{\mathbf{r},s_{0}}(Q_{0}),\ldots,
% 			\widehat{f}_{\mathbf{r},s_{\bb{\mathbf{s}}-1}}(Q_{\bb{\mathbf{s}}-1})
% 			).
% 		\end{align*}
% 	\end{flushleft}
% 	Therefore the $S^{\star}\times S$-sorted mapping $\widehat{f}$ is a homomorphism. Furthermore, $\widehat{f}\circ \eta^{\Sigma} = f$, because, for every $\mathbf{s}\in S^{\star}$, $s\in S$, and $\sigma\in \Sigma_{\mathbf{s},s}$, we have that%
% 	\begin{align*}
% 		\widehat{f}_{\mathbf{s},s}(\eta^{\Sigma}_{\mathbf{s},s}(\sigma)) &=
% 		(p^{\mathbf{s}})^{\sharp}_{s}(\sigma(v^{s_{0}}_{0},\ldots,v^{s_{\bb{\mathbf{s}}-1}}_{\bb{\mathbf{s}}-1}))
% 		\\&=
% 		\sigma^{\mathbf{A}_{\mathbf{s}}}
% 		(p^{\mathbf{s}}_{{s}_{0}}(v^{s_{0}}_{0}),
% 		\ldots,
% 		p^{\mathbf{s}}_{{s}_{\bb{\mathbf{s}}-1}}(v^{s_{\bb{\mathbf{s}}-1}}_{\bb{\mathbf{s}}-1})
% 		) \\
% 		&=
% 		\xi_{\mathbf{s},\mathbf{s},s}^{\mathbf{A}}
% 		( f_{(\mathbf{s},s)}(\sigma),
% 		(\pi^{\mathbf{s}}_{0})^{\mathbf{A}},
% 		\ldots,
% 		(\pi^{\mathbf{s}}_{\bb{\mathbf{s}}-1})^{\mathbf{A}}
% 		) \\
% 		&=
% 		f_{\mathbf{s},s}(\sigma) \tag{by $\mathrm{H}_{2}$}.
% 	\end{align*}
% 	It is obvious that $\widehat{f}$ is the unique homomorphism such that $\widehat{f}\circ \eta^{\Sigma} = f$.  Henceforth $\mathbf{Tm}_{\mathrm{H}_{S}}(\Sigma)$ is isomorphic to $\mathbf{T}_{\mathrm{H}_{S}}(\Sigma)$.
% \end{proof}
%
%
% This isomorphism, together with the adjunction $\mathbf{T}_{\mathrm{H}_{S}}\dashv\mathrm{G}_{\mathrm{H}_{S}}$,  has as an immediate consequence that, for every $S$-sorted set $A$ and every $S$-sorted signature $\Sigma$, the sets $\mathrm{Hom}({\Sigma},\mathrm{G}_{\mathrm{H}_{S}}(\mathbf{O}_{\mathrm{H}_{S}}(A)))$, in 
% $\mathsf{Set}^{S^{\star}\times S}$, and $\mathrm{Hom}(\mathbf{Tm}_{\mathrm{H}_{S}}(\Sigma),\mathbf{O}_{\mathrm{H}_{S}}(A))$, in $\mathsf{Alg}(\mathrm{H}_{S})$, are naturally isomorphic.
%
% Actually, for an $S$-sorted set $A$ and an $S$-sorted signature $\Sigma$,
% \begin{enumerate}
% 	\item the mapping that assigns to a structure of $\Sigma$-algebra $F$ on $A$ (i.e., an $S^{\star}\times S$-sorted mapping $F$ from $\Sigma$ to $\mathrm{O}_{\mathrm{H}_{S}}(A)$) the unique  homomorphism of Hall algebras $\widehat{F} = (\widehat{F}_{w,s})_{(w,s)\in S^{\star}\times S}$ from $\mathbf{Tm}_{\mathrm{H}_{S}}(\Sigma)$ to $\mathbf{O}_{\mathrm{H}_{S}}(A)$ such that $\widehat{F}\circ \eta^{\Sigma} = F$, together with
% 	\item the mapping that assigns to a homomorphism $h$ from $\mathbf{Tm}_{\mathrm{H}_{S}}(\Sigma)$ to $\mathbf{O}_{\mathrm{H}_{S}}(A)$, essentially, the algebraic structure $\mathrm{G}_{\mathrm{H}_{S}}(h)\circ \eta^{\Sigma}$ on $A$, where $\eta^{\Sigma}$ is the canonical embedding of $\Sigma$ into $\mathrm{T}_{\mathrm{H}_{S}}(\Sigma)$, 
% \end{enumerate}
% which are natural in both variables and mutually inverse bijections, provide the mentioned natural isomorphism.
%
% \begin{remark}
% 	For an $S$-sorted set $A$, a structure of $\Sigma$-algebra $F$ on $A$ and a word $\mathbf{s}\in S^{\star}$, the subfamily $(\widehat{F}_{\mathbf{s},s})_{s\in S}$ of $\widehat{F} = (\widehat{F}_{\mathbf{s},s})_{(\mathbf{s},s)\in S^{\star}\times S}$,
% 	denoted by $\mathrm{d}^{(A,F)} = (\mathrm{d}^{(A,F)}_{\mathbf{s},s})_{s\in S}$, is the unique homomorphism from $\mathbf{T}_{\Sigma}(\vs{\mathbf{s}})$ to $(A,F)^{A_{\mathbf{s}}}$, the direct $A_{\mathbf{s}}$-power of $(A,F)$, such that $\mathrm{d}^{(A,F)}_{\mathbf{s}}\circ \eta^{\vs{\mathbf{s}}} = \mathrm{p}^{A}_{\mathbf{s}}$, where $\mathrm{p}^{A}_{\mathbf{s}}$ is the $S$-sorted mapping from $\vs{\mathbf{s}}$ to $A^{A_{\mathbf{s}}}$ defined, for every $s\in S$ and every $v^{s}_{i}\in(\vs{\mathbf{s}})_{s}$, as $\mathrm{p}^{A}_{\mathbf{s},s}(v^{s}_{i}) = \mathrm{pr}^{A}_{\mathbf{s},i}$.
% \end{remark}
%
% For a set of sorts $S$ and an $S$-sorted signature $\Sigma$, the derived operation symbols of $\Sigma$, i.e., the elements of the components $\mathrm{T}_{\Sigma}(\vs{\mathbf{s}})_{s}$ of $\mathrm{Tm}_{\mathrm{H}_{S}}(\Sigma)$ can be considered as the operation symbols of a new signature. The mappings $\varphi$ from a set of sorts $S$ to another $T$, together with the $S^{\star}\times S$-sorted mappings $d$ that assign to operation symbols of an $S$-sorted signature $\Sigma$ derived operations symbols of a $T$-sorted signature $\Lambda$, i.e., terms of the components $\mathrm{T}_{\Lambda}(\vs{\varphi^{\star}(\mathbf{s})})_{\varphi(s)}$ of $\mathrm{Tm}_{\mathrm{H}_{T}}(\Lambda)_{\varphi^{\star}\times\varphi}$, form a new class of morphisms between many-sorted signatures which we will call \emph{derivors}.
%
% \begin{definition}\label{DDerivor}
% 	Let $\sig{\Sigma} = (S,\Sigma)$ and $\sig{\Lambda} = (T,\Lambda)$ be many-sorted signatures. A \emph{derivor} from $\sig{\Sigma}$ to $\sig{\Lambda}$ is an ordered pair $\mathbf{d} = (\varphi,d)$, where $\varphi\colon S\mor T$ and  $d\colon\Sigma\mor\mathrm{Tm}_{\mathrm{H}_{T}}(\Lambda)_{\varphi^{\star}\times\varphi}$. Thus, for every $(\mathbf{s},s)\in S^{\star}\times S$, %
% 	$$
% 	d_{\mathbf{s},s}\colon \Sigma_{\mathbf{s},s}\mor\mathrm{Tm}_{\mathrm{H}_{T}}(\Lambda)_{\varphi^{\star}(\mathbf{s}),\varphi(s)}=
% 	\mathrm{T}_{\Lambda}(\vs{\varphi^{\star}(\mathbf{s})})_{\varphi(s)}.
% 	$$
% 	Hence, for every $\sigma\in \Sigma_{\mathbf{s},s}$, $d_{\mathbf{s},s}(\sigma)\in \mathrm{T}_{\Lambda}(\vs{\varphi^{\star}(\mathbf{s})})_{\varphi(s)}$, i.e., $d_{\mathbf{s},s}(\sigma)$ is a term of sort $\varphi(s)$ with variables in $\vs{\varphi^{\star}(\mathbf{s})}$.
% \end{definition}
%
%
%
% \begin{remark}\label{Rdop}
% 	Let $\varphi$ be a mapping from the set of sorts $S$ to the set of sorts $T$ and let $V^{T}$ be a fixed $T$-sorted set with a countably infinity of variables in each component, say $V^{T}_{t} = \{v^{t}_{n}\mid n\in \mathbb{N}\}$ (which is isomorphic to $\{t\}\times\mathbb{N}$). Then, for every $\mathbf{s}\in S^{\star}$, we let $\vs{\varphi^{\star}(\mathbf{s})} = ((\vs{\varphi^{\star}(\mathbf{s})})_{t})_{t\in T}$ stand for the $T$-sorted set  defined, for every $t\in T$, as:
% 	$$
% 	(\vs{\varphi^{\star}(\mathbf{s})})_{t} = \{v^{t}_{0},\ldots,v^{t}_{\bb{\varphi^{\star}(\mathbf{s})}_{t}-1}\} (\cong \varphi^{\star}(\mathbf{s})^{-1}[\{t\}] = \{i\in \bb{\varphi^{\star}(\mathbf{s})}\mid \varphi(w_{i}) = t\}).
% 	$$
% 	Thus, for every $t\in T$, $(\vs{\varphi^{\star}(\mathbf{s})})_{t}$ is an initial segment of $V^{T}_{t}$ and 
% 	$\vs{\varphi^{\star}(\mathbf{s})}$ is a finite subset of $V^{T}$. The latter is fulfilled because $\mathrm{supp}_{T}(\vs{\varphi^{\star}(\mathbf{s})}) = \mathrm{Im}(\varphi^{\star}(\mathbf{s}))$ and $\bb{\varphi^{\star}(\mathbf{s})} = \bb{\mathbf{s}}$ is finite. Moreover, $\mathrm{supp}_{T}(\vs{\varphi^{\star}(\mathbf{s})}) =\varphi[\mathrm{supp}_{S}(\vs{\mathbf{s}})]$. Recall  that $\varphi^{\star}(\mathbf{s}) = (\varphi(s_{i}))_{i\in \bb{\mathbf{s}}}$ and that, since $\varphi^{\star}(\lambda) = \lambda$, $\vs{\varphi^{\star}(\lambda)} = \varnothing^{T} = (\varnothing)_{t\in T}$.
%
% 	If $\mathbf{d} = (\varphi,d)$ is a derivor from $\sig{\Sigma}$ to $\sig{\Lambda}$, $\mathbf{B}$ a 
% 	$\sig{\Lambda}$-algebra, $(\mathbf{s},s)\in S^{\star}\times S$ and $\sigma\in \Sigma_{\mathbf{s},s}$, then the derived operation (or term operation or polynomial operation) of type $(\varphi^{\star}(\mathbf{s}),\varphi(s))$ determined by the term $d_{\mathbf{s},s}(\sigma)$ in $\mathbf{B}$, denoted by $d_{\mathbf{s},s}(\sigma)^{\mathbf{B}}$ or by $d(\sigma)^{\mathbf{B}}$ for short, is a mapping from $B_{\varphi^{\star}(\mathbf{s})}$ to $B_{\varphi(s)}$. In particular, for a $T$-sorted set $Y$, we have that the derived operation (or term operation or polynomial operation) of type $(\varphi^{\star}(\mathbf{s}),\varphi(s))$ determined by the term $d_{\mathbf{s},s}(\sigma)$ in $\mathbf{T}_{\Lambda}(Y)$, denoted by $d_{\mathbf{s},s}(\sigma)^{\mathbf{T}_{\Lambda}(Y)}$ or by  $d(\sigma)^{\mathbf{T}_{\Lambda}(Y)}$ for short, is a mapping from $\mathrm{T}_{\Lambda}(Y)_{\varphi^{\star}(\mathbf{s})}$ to $\mathrm{T}_{\Lambda}(Y)_{\varphi(s)}$. Notice that, for every $s\in S$ and every $\sigma\in \Sigma_{\lambda,s}$, $d_{\lambda,s}(\sigma)\in \mathrm{T}_{\Lambda}(\varnothing^{T})_{\varphi(s)}$, i.e., $d_{\lambda,s}(\sigma)$ is a closed term of sort $\varphi(s)$.  Therefore, for every $\sig{\Lambda}$-algebra $\mathbf{B}$, $d_{\lambda,s}(\sigma)^{\mathbf{B}}$ will be a mapping from $B_{\lambda}$ to $B_{\varphi(s)}$, i.e., essentially, a distinguished element of $B_{\varphi(s)}$. In particular, for a $T$-sorted set $Y$, $d_{\lambda,s}(\sigma)^{\mathbf{T}_{\Lambda}(Y)}$ will be a mapping from $\mathrm{T}_{\Lambda}(Y)_{\lambda}$ to $\mathrm{T}_{\Lambda}(Y)_{\varphi(s)}$, i.e., essentially, a distinguished element of $\mathrm{T}_{\Lambda}(Y)_{\varphi(s)}$. Let us note that $d(\sigma)^{\mathbf{B}}$ is the value of $\mathrm{d}^{\mathbf{B}}_{\varphi^{\star}(\mathbf{s}),s}$ at $d_{\mathbf{s},s}(\sigma)$ and that $d(\sigma)^{\mathbf{T}_{\Lambda}(Y)}$ is the value of $\mathrm{d}^{\mathbf{T}_{\Lambda}(Y)}_{\varphi^{\star}(\mathbf{s}),s}$ at $d_{\mathbf{s},s}(\sigma)$, where 
% 	$\mathrm{d}^{\mathbf{B}}_{\varphi^{\star}(\mathbf{s})}$ is the unique homomorphism from $\mathbf{T}_{\Lambda}(\vs{\varphi^{\star}(\mathbf{s})})$ to $\mathbf{B}^{B_{\varphi^{\star}(\mathbf{s})}}$ for which the left diagram below commutes, and $\mathrm{d}^{\mathbf{T}_{\Lambda}(Y)}_{\varphi^{\star}(\mathbf{s})}$ the unique homomorphism from $\mathbf{T}_{\Lambda}(\vs{\varphi^{\star}(\mathbf{s})})$ to $\mathbf{T}_{\Lambda}(Y)^{\mathrm{T}_{\Lambda}(Y)_{\varphi^{\star}(\mathbf{s})}}$ for which the right diagram below commutes.
% 	$$
% 	\begin{aligned}
% 		\xymatrix@C=40pt@R=33pt{
% 			\vs{\varphi^{\star}(\mathbf{s})}
% 			\ar[r]^-{\eta^{\vs{\varphi^{\star}(\mathbf{s})}}}
% 			\ar[rd]_-{\mathrm{p}^{B}_{\varphi^{\star}(\mathbf{s})}} &
% 			\mathrm{T}_{\Lambda}(\vs{\varphi^{\star}(\mathbf{s})}) \ar[d]^{\mathrm{d}^{\mathbf{B}}_{\varphi^{\star}(\mathbf{s})}}
% 			\\
% 			& B^{B_{\varphi^{\star}(\mathbf{s})}}
% 		}
% 		%
% 	\end{aligned}
% 	\quad \text{} \qquad\;
% 	\begin{aligned}
% 		%
% 		\xymatrix@C=40pt@R=33pt{
% 			\vs{\varphi^{\star}(\mathbf{s})}
% 			\ar[r]^-{\eta^{\vs{\varphi^{\star}(\mathbf{s})}}}
% 			\ar[rd]_-{\mathrm{p}^{\mathrm{T}_{\Lambda}(Y)}_{\varphi^{\star}(\mathbf{s})}} &
% 			\mathrm{T}_{\Lambda}(\vs{\varphi^{\star}(\mathbf{s})}) \ar[d]^{\mathrm{d}^{\mathbf{T}_{\Lambda}(Y)}_{\varphi^{\star}(\mathbf{s})}}
% 			\\
% 			& \mathrm{T}_{\Lambda}(Y)^{\mathrm{T}_{\Lambda}(Y)_{\varphi^{\star}(\mathbf{s})}}
% 		}
% 		%
% 	\end{aligned}
% 	$$
% 	Recall that, for every $t\in T$ and every $i\in (\vs{\varphi^{\star}(\mathbf{s})})_{t}$ (thus $(\varphi^{\star}(\mathbf{s}))(i) = \varphi(s_{i}) = t$), $\mathrm{p}^{B}_{\varphi^{\star}(\mathbf{s}),t}(i)$ is $\mathrm{pr}^{B}_{\varphi^{\star}(\mathbf{s}),i}$, i.e., the $\varphi^{\star}(\mathbf{s})$-ary, $i$-th canonical projection from $B_{\varphi^{\star}(\mathbf{s})}$ to $B_{t} = B_{(\varphi^{\star}(\mathbf{s}))(i)}$.
% \end{remark}
%
% \begin{remark}
% 	While derivors are not the most general type of morphism that might be considered between many-sorted signatures---for instance, one could consider polyderivors, see~\cite{CVST10b}---, they are an important class of such morphisms. One reason for its relevance is its formal properties (see below), another that there are many mathematical examples of them which are of interest (see~\cite{CVST10b}).
% \end{remark}
%
% Since by, Proposition~\ref{iso:FrH-TerH}, $\mathbf{Tm}_{\mathrm{H}_{T}}(\Lambda)$ is isomorphic to 
% $\mathbf{T}_{\mathrm{H}_{T}}(\Lambda)$, the derivors can be defined, alternative, but equivalently, as ordered pairs  $\mathbf{d} = (\varphi,d)$ with $d\colon\Sigma\mor \mathrm{T}_{\mathrm{H}_{T}}(\Lambda)$.
%
% On the other hand, every mapping $\varphi\colon S\mor T$ determines a functor from the category  $\mathsf{Alg}(\mathrm{H}_{T})$ to the category $\mathsf{Alg}(\mathrm{H}_{S})$, so $\mathrm{Tm}_{\mathrm{H}_{T}}(\Lambda)_{\varphi^{\star}\times\varphi}$ is in its turn e\-quipped with a structure of Hall algebra for $S$, which will allow us, in particular, to define the composition of derivors. We next show the existence of such a functor by defining a morphism of algebraic presentations from  $(\Sigma^{\mathrm{H}_{S}},\mathcal{E}^{\mathrm{H}_{S}})$ to $(\Sigma^{\mathrm{H}_{T}},\mathcal{E}^{\mathrm{H}_{T}})$.
%
%
%
% \begin{proposition}
% 	Let $\varphi\colon S\mor T$ be a mapping between sets of sorts and $h^{\varphi}$ the $S^{\star}\times S$-sorted mapping from $\Sigma^{\mathrm{H}_{S}}$ to $\Sigma^{\mathrm{H}_{T}}_{\varphi^{\star}\times \varphi}$ defined as follows: 
% 	\begin{enumerate}
% 		\item For every $\mathbf{s}\in S^{\star}$ and every $i\in\bb{\mathbf{s}}$, $h^{\varphi}(\pi^{\mathbf{s}}_{i}) = \pi^{\varphi^{\star}(\mathbf{s})}_{i}$.
%
% 		\item  For every $\mathbf{r},\,\mathbf{s}\in S^{\star}$ and every $s\in S$,
% 		$h^{\varphi}(\xi_{\mathbf{r},\mathbf{s},s}) = \xi_{\varphi^{\star}(\mathbf{r}),\varphi^{\star}(\mathbf{s}),\varphi(s)}$.
% 	\end{enumerate}
% 	Then $(\varphi^{\star}\times \varphi, h^{\varphi}) \colon(S^{\star}\times S,\Sigma^{\mathrm{H}_{S}},\mathcal{E}^{\mathrm{H}_{S}})\mor
% 	(T^{\star}\times T,\Sigma^{\mathrm{H}_{T}},\mathcal{E}^{\mathrm{H}_{T}})$ is a morphism of algebraic presentations.
% \end{proposition}
%
% The morphisms of algebraic presentations determine functors in the opposite direction between the associated categories of algebras. Therefore, each mapping $\varphi\colon S\mor T$ between sets of sorts, determines a functor $(\varphi^{\star}\times \varphi, h^{\varphi})^{\ast}$ from $\mathsf{Alg}(\mathrm{H}_{T})$ to $\mathsf{Alg}(\mathrm{H}_{S})$, which transforms Hall algebras for $T$ into Hall algebras for $S$. The action of the functor on the free Hall algebra on a $T$-sorted signature $\Lambda$ is a Hall algebra for $S$, whose underlying $S^{\star}\times S$-sorted set is
% $\mathrm{Tm}_{\mathrm{H}_{T}}(\Lambda)_{\varphi^{\star}\times \varphi}$.
%
% If $\mathbf{d}\colon\sig{\Sigma}\mor\sig{\Lambda}$ is a derivor, then
% $d\colon \Sigma\mor \mathrm{Tm}_{\mathrm{H}_{T}}(\Lambda)_{\varphi^{\star}\times\varphi}$ determines a homomorphism of Hall algebras $d^{\sharp}\colon\mathbf{Tm}_{\mathrm{H}_{S}}(\Sigma)\mor
% \mathbf{Tm}_{\mathrm{H}_{T}}(\Lambda)_{\varphi^{\star}\times\varphi}$. Thus, for every $(\mathbf{s},s)\in S^{\star}\times S$, $d^{\sharp}_{\mathbf{s},s}$ sends terms in $\mathrm{T}_{\Sigma}(\vs{\mathbf{s}})_{s}$ to terms in $\mathrm{T}_{\Lambda}(\vs{\varphi^{\star}(\mathbf{s})})_{\varphi(s)}$.
%
% \begin{definition}\label{DDerCompId}
% 	Let $\mathbf{d}\colon\sig{\Sigma}\mor\sig{\Lambda}$ and $\mathbf{e}\colon\sig{\Lambda}\mor\sig{\Omega}$ be  derivors. Then $\mathbf{e}\circ \mathbf{d} = (\psi,e)\circ(\varphi,d)$, the \emph{composition} of $\mathbf{d}$ and  $\mathbf{e}$, is the derivor $(\psi\circ\varphi,e^{\sharp}_{\varphi^{\star}\times
% 		\varphi}\circ d)$, where $e^{\sharp}_{\varphi^{\star}\times\varphi}\circ d$ is obtained from
% 	$$
% 	\begin{aligned}
% 		\xymatrix@C=40pt@R=30pt{
% 			\Lambda
% 			\ar[r]^-{\eta^{\Lambda}}
% 			\ar[rd]_-{e} &
% 			\mathrm{Tm}_{\mathrm{H}_{T}}(\Lambda) \ar[d]^{e^{\sharp}}
% 			\\
% 			& \mathrm{Tm}_{\mathrm{H}_{U}}(\Omega)_{\psi^{\star}\times \psi}
% 		}
% 		%
% 	\end{aligned}
% 	\qquad \text{as} \qquad\;\;
% 	\begin{aligned}
% 		%
% 		\xymatrix@C=40pt@R=30pt{
% 			\mathrm{Tm}_{\mathrm{H}_{T}}(\Lambda)_{\varphi^{\star}\times \varphi}
% 			\ar[d]^{e^{\sharp}_{\varphi^{\star}\times \varphi}} &
% 			\Sigma \ar[l]_-{d} \\
% 			{\mathrm{Tm}_{\mathrm{H}_{U}}(\Omega)_{\psi^{\star}\times\psi}}_{\varphi^{\star}\times\varphi}
% 		}
% 	\end{aligned}
% 	$$
% 	being $e^{\sharp}$ the canonical extension of $e$ to the free Hall algebra on $\Lambda$.
%
% 	For every many-sorted signature $\sig{\Sigma} = (S,\Sigma)$, the \emph{identity} at $(S,\Sigma)$
% 	is $(\id^{S},\eta^{\Sigma})$.
% \end{definition}
%
% The just stated definition allows us to form a category of many-sorted signatures whose morphisms are the derivors.
%
% \begin{proposition}\label{PSigdCat}
% 	The many-sorted signatures together with the derivors constitute a category, denoted by $\mathsf{Sig}_{\mathfrak{d}}$.
% \end{proposition}
%
%
%
% \begin{remark}
% 	Let $\mathsf{Sig}$ be the category whose objects are the many-sorted signatures and whose morphisms from $\sig{\Sigma}$ to $\sig{\Lambda}$ are the ordered pairs $(\varphi,d)$, where $\varphi$ is a mapping from $S$ to $T$ and $d$ an $S^{\star}\times S$-sorted mapping from $\Sigma$ to $\Lambda_{\varphi^{\star}\times \varphi}$ (thus a mapping in $\mathsf{Sig}(S)$). Note that $\mathsf{Sig}$ is the category obtained by means of the Grothendieck construction for an  appropriate contravariant functor from $\mathsf{Set}$ to $\mathsf{Cat}$.
% 	We next show that the category $\mathsf{Sig}_{\mathfrak{d}}$ can be obtained as the Kleisli category for a suitable monad. In fact, for every  set of sorts $S$, we have the adjoint situation $\mathbf{T}_{\mathrm{H}_{S}}\dashv\mathrm{G}_{\mathrm{H}_{S}}$
% 	%$$
% 	%\xymatrix@=5pc{ \mathsf{Alg}(\mathrm{H}_{S})
% 		%\ar@<1.5ex>[r]^{\mathrm{G}_{\mathrm{H}_{S}}} \ar@{}[r]|{\uadj} &
% 		%\mathbf{Set}^{S^{\star}\times S}
% 		%\ar@<1.5ex>[l]^{\mathbf{T}_{\mathrm{H}_{S}}} }
% 	%$$
% 	and, thus, a monad on $\mathsf{Sig}(S)$ which we will denote by $\mathbb{T}_{\mathrm{H}_{S}} = (\mathrm{T}_{\mathrm{H}_{S}},\eta^{\mathrm{H}_{S}},\mu^{\mathrm{H}_{S}})$. Then the ordered triple $\boldsymbol{\mathfrak{d}} = (\mathfrak{d},\eta,\mu)$ in which (1) $\mathfrak{d}$ is the endofunctor at $\mathsf{Sig}$ that sends $(S,\Sigma)$ to $(S,\mathrm{T}_{\mathrm{H}_{S}}(\Sigma))$ and a morphism $(\varphi,d)$ from $(S,\Sigma)$ to  $(T,\Lambda)$ to the morphism $(\varphi,d^{\sharp})$ from $(S,\mathrm{T}_{\mathrm{H}_{S}}(\Sigma))$ to  $(T,\mathrm{T}_{\mathrm{H}_{T}}(\Lambda))$, (2) $\eta$ the natural transformation from $\Id_{\mathsf{Sig}}$ to  $\mathfrak{d}$ that sends $(S,\Sigma)$ to $\eta_{(S,\Sigma)} = (\id,\eta^{\mathrm{H}_{S}}_{\Sigma})$, and (3) $\mu$ the natural transformation from $\mathfrak{d}\circ \mathfrak{d}$ to $\mathfrak{d}$ that sends $(S,\Sigma)$ to $\mu_{(S,\Sigma)} = (\id,\mu^{\mathrm{H}_{S}}_{\Sigma})$, is a monad on $\mathsf{Sig}$ and $\mathsf{Kl}(\boldsymbol{\mathfrak{d}})$, the Kleisli category for $\boldsymbol{\mathfrak{d}}$, is isomorphic to $\mathsf{Sig}_{\mathfrak{d}}$. This shows, in our opinion, the mathematical naturalness of the notion of derivor. Moreover, by defining a suitable notion of transformation between derivors one can equip the category $\mathsf{Sig}_{\mathfrak{d}}$ with a structure of $2$-category (this was, in fact, already done in~\cite{CVST10b} for polyderivors).
%
% \end{remark}
%
%
%
% \begin{remark}
% 	Since $\mathsf{Sig}$ has coproducts, $\mathsf{Kl}(\boldsymbol{\mathfrak{d}})\cong \mathsf{Sig}_{\mathfrak{d}}$ has coproducts. In fact, let $(\sig{\Sigma}^{i})_{i\in I}$ be an $I$-indexed family in $\mathrm{Ob}(\mathsf{Kl}(\boldsymbol{\mathfrak{d}}))$ and $(\coprod_{i\in I}\sig{\Sigma}^{i},(\mathrm{in}^{\sig{\Sigma}^{i}})_{i\in I})$ the coproduct of $(\sig{\Sigma}^{i})_{i\in I}$ in $\mathsf{Sig}$, where  the first coordinate of $\coprod_{i\in I}\sig{\Sigma}^{i}$ is $\coprod_{i\in I}S_{i}$ and the second coordinate of $\coprod_{i\in I}\sig{\Sigma}^{i}$, denoted by $\coprod_{i\in I}\mathrm{\Sigma}^{i}$, is defined, for every $(((s_{i_{j}},i_{j})_{j\in n}),(s,i))\in (\coprod_{i\in I}S_{i})^{\star}\times \coprod_{i\in I}S_{i}$, as follows: 
% 	\begin{equation*}
% 		\textstyle
% 		(\coprod_{i\in I}\mathrm{\Sigma}^{i})_{(((s_{i_{j}},i_{j})_{j\in n}),(s,i))} = 
% 		\begin{cases}
% 			\Sigma^{i}_{\mathbf{s},s}, \text{\,\,if $\exists\,\mathbf{s}\in S_{i}^{\star}$ such that $(s_{i_{j}},i_{j})_{j\in n} = \mathrm{in}_{i}^{\star}(\mathbf{s})$;}\\
% 			\varnothing, \text{\,\,otherwise}.
% 		\end{cases}
% 	\end{equation*}
% 	Let us recall that $\mathrm{in}_{i}^{\star}$ is the monoid homomorphism from the free monoid on $S_{i}$ to the free monoid on $\coprod_{i\in I}S_{i}$ that sends $\mathbf{s} = (s_{j})_{j\in n}\in S_{i}^{\star}$ to 
% 	$(s_{j},i)_{j\in n}\in (\coprod_{i\in I}S_{i})^{\star}$. Therefore, 
% 	$(s_{i_{j}},i_{j})_{j\in n} = \mathrm{in}_{i}^{\star}(\mathbf{s})$ if and only if, for every $j\in n$, $i_{j} = i$ and $s_{i_{j}} = w_{j}$.
% 	Then $(\boldsymbol{\mathfrak{d}}(\coprod_{i\in I}\sig{\Sigma}^{i}),(\eta_{\coprod_{i\in I}\sig{\Sigma}^{i}}\circ \mathrm{in}^{\sig{\Sigma}^{i}})_{i\in I})$ will be the required coproduct of $(\sig{\Sigma}^{i})_{i\in I}$ in $\mathsf{Kl}(\boldsymbol{\mathfrak{d}})$. This property of the category $\mathsf{Sig}_{\mathfrak{d}}$ will be essential in our work because to be able to compare Curry-Howard towers it will be necessary, as an intermediate step, to construct, by means of the formation of coproducts, from given many-sorted  signatures $(\sig{\Sigma}^{i})_{i\in I}$, a new many-sorted signature $\coprod_{i\in I}\sig{\Sigma}^{i}$ (finally, from $\coprod_{i\in I}\sig{\Sigma}^{i}$, we will construct the many-sorted signature $\biguplus_{i\in I}\sig{\Sigma}^{i}$ which will be the one of interest for us).
%
% 	%$\sig{\Sigma}$, $\sig{\Omega}_{\mathrm{c}}$, $\sig{\Omega}_{\mathrm{c},\mathcal{A}}$, $\sig{\Lambda}$, and $\sig{\Omega}_{\mathrm{c},\mathcal{B}}$ and $\mathsf{Sig}_{\mathfrak{d}}$-morphisms 
% 	%$\mathbf{d}$ from $\sig{\Sigma}$ to $\sig{\Lambda}$ and $\mathbf{e}$ from $\sig{\Omega}_{\mathrm{c},\mathcal{A}}$ to $\sig{\Omega}_{\mathrm{c},\mathcal{B}}$, as, e.g.,  
% 	%$$
% 	%\textstyle
% 	%\mathbf{d}\amalg\mathbf{id}_{\sig{\Omega}_{\mathrm{c}}}\colon \sig{\Sigma}\amalg\sig{\Omega}_{\mathrm{c}}\mor \sig{\Lambda}\amalg\sig{\Omega}_{\mathrm{c}} 
% 	%\quad\text{and}\quad
% 	%\mathbf{d}\amalg\mathbf{e}\colon \sig{\Sigma}\amalg\sig{\Omega}_{\mathrm{c},\mathcal{A}}\mor \sig{\Lambda}\amalg\sig{\Omega}_{\mathrm{c},\mathcal{B}}.
% 	%$$
% \end{remark}
%
% \section{Generalized derivors}\label{SGenDerivors}
%
% The general derivors of~\cite{LG18} are morphisms between order-sorted signatures with predicates which send every sort to a formula that delimits, inside the interpretation of the image sort, the elements that interpret it, and every operation symbol to a formula. Here we need neither the order on the sorts nor the predicates, and, instead of arbitrary formulas, we use existence equations. Thus our generalized derivors are derivors in the sense of Definition~\ref{DDerivor} together with a set of existence equations for every sort.
%
%
% We first fix some notation. For a set of sorts $T$, a word $\mathbf{t}\in T^{\star}$ and $j\in\bb{\mathbf{t}}$, we let $v^{\mathbf{t}}_{j}$ stand for the variable of $\vs{\mathbf{t}}$ that corresponds to the $j$-th letter of $\mathbf{t}$; thus $v^{\mathbf{t}}_{j}\in(\vs{\mathbf{t}})_{t_{j}}$ and, for a word $(t)$ of length one, $v^{(t)}_{0}=v^{t}_{0}$. For a $T$-sorted signature $\Lambda$ and a $T$-sorted set $X$, we let $\mathrm{Eq}_{\Lambda}(X)$ stand for $\coprod_{t\in T}\mathrm{T}_{\Lambda}(X)^{2}_{t}$, the set of existence equations of $\Lambda$ with variables in $X$, an element $(P,Q)$ of $\mathrm{T}_{\Lambda}(X)^{2}_{t}$ being the existence equation $P\eeq Q$ of Definition~\ref{DQEEq}. Note that $\mathrm{Sub}(\mathrm{Eq}_{\Lambda}(X))\cong\mathrm{Sub}((\mathrm{T}_{\Lambda}(X)^{2}_{t})_{t\in T})$, so that a set of existence equations can be regarded, equivalently, as a $T$-sorted set. If $\mathbf{B}$ is a partial $\Lambda$-algebra, $b\colon X\mor B$ a valuation and $(P,Q)\in\mathrm{Eq}_{\Lambda}(X)$, we write $\mathbf{B}\models(P,Q)[b]$ if $P\eeq Q$ holds in $\mathbf{B}$ under $b$, and, for $\mathcal{E}\subseteq\mathrm{Eq}_{\Lambda}(X)$, $\mathbf{B}\models\mathcal{E}[b]$ if $\mathbf{B}\models(P,Q)[b]$ for every $(P,Q)\in\mathcal{E}$. A family $\mathbf{b}=(b_{j})_{j\in\bb{\mathbf{t}}}\in B_{\mathbf{t}}$ is identified with the valuation of $\vs{\mathbf{t}}$ in $B$ that sends $v^{\mathbf{t}}_{j}$ to $b_{j}$. Finally, for a $T$-sorted mapping $f\colon X\mor\mathrm{T}_{\Lambda}(Y)$ and $\mathcal{E}\subseteq\mathrm{Eq}_{\Lambda}(X)$, we let $f^{\sharp}[\mathcal{E}]$ stand for $\{(f^{\sharp}(P),f^{\sharp}(Q))\mid(P,Q)\in\mathcal{E}\}\subseteq\mathrm{Eq}_{\Lambda}(Y)$, where $f^{\sharp}$ is the substitution determined by $f$; when $X=\vs{\mathbf{t}}$, $f$ is identified with the family $(f(v^{\mathbf{t}}_{j}))_{j\in\bb{\mathbf{t}}}$.
%
% \begin{definition}\label{DGeneralDerivor}
% Let $\sig{\Sigma}=(S,\Sigma)$ and $\sig{\Lambda}=(T,\Lambda)$ be many-sorted signatures. A \emph{generalized derivor} from $\sig{\Sigma}$ to $\sig{\Lambda}$ is an ordered triple $\mathbf{d}=(\varphi,D,d)$ such that
% \begin{enumerate}
% \item $(\varphi,d)$ is a derivor from $\sig{\Sigma}$ to $\sig{\Lambda}$, see Definition~\ref{DDerivor}, i.e., $\varphi\colon S\mor T$ and, for every $(\mathbf{s},s)\in S^{\star}\times S$, $d_{\mathbf{s},s}\colon\Sigma_{\mathbf{s},s}\mor\mathrm{T}_{\Lambda}(\vs{\varphi^{\star}(\mathbf{s})})_{\varphi(s)}$; and
% \item $D=(D_{s})_{s\in S}$ assigns to every sort $s\in S$ a set $D_{s}\subseteq\mathrm{Eq}_{\Lambda}(\vs{(\varphi(s))})$ of existence equations of $\Lambda$ in the single variable $v^{\varphi(s)}_{0}$, which we call the \emph{domain} of $s$.
% \end{enumerate}
% For a word $\mathbf{s}\in S^{\star}$, we let $D_{\mathbf{s}}\subseteq\mathrm{Eq}_{\Lambda}(\vs{\varphi^{\star}(\mathbf{s})})$ stand for $\bigcup_{j\in\bb{\mathbf{s}}}(\iota^{\mathbf{s}}_{j})^{\sharp}[D_{s_{j}}]$, where $\iota^{\mathbf{s}}_{j}$ sends $v^{\varphi(s_{j})}_{0}$ to $v^{\varphi^{\star}(\mathbf{s})}_{j}$, i.e., $D_{\mathbf{s}}$ is the set of the domain conditions of the letters of $\mathbf{s}$, each one stated for the corresponding variable of $\vs{\varphi^{\star}(\mathbf{s})}$.
% \end{definition}
%
% The set $D_{s}$ is the counterpart of the formula $\Delta_{s}(x)$ of~\cite{LG18}, and the term $d_{\mathbf{s},s}(\sigma)$ is the counterpart of the formula $\Phi^{\sigma}_{\mathbf{s},s}(x_{0},\ldots,x_{\bb{\mathbf{s}}-1},y)$, which for an ordinary derivor is, in~\cite{LG18}, the equation $y=d_{\mathbf{s},s}(\sigma)$. Every derivor $(\varphi,d)$ is the generalized derivor $(\varphi,\varnothing,d)$, where $\varnothing$ stands for the family $(\varnothing)_{s\in S}$.
%
%
% \begin{remark}\label{RGenDerivorSem}
% Let $\mathbf{d}=(\varphi,D,d)$ be a generalized derivor from $\sig{\Sigma}$ to $\sig{\Lambda}$ and $\mathbf{B}$ a $\Lambda$-algebra. For $s\in S$, let $B^{D}_{s}=\{b\in B_{\varphi(s)}\mid\mathbf{B}\models D_{s}[b]\}$ and, for $\mathbf{s}\in S^{\star}$, $B^{D}_{\mathbf{s}}=\prod_{j\in\bb{\mathbf{s}}}B^{D}_{s_{j}}$; since $\iota^{\mathbf{s}}_{j}$ only renames a variable, $B^{D}_{\mathbf{s}}=\{\mathbf{b}\in B_{\varphi^{\star}(\mathbf{s})}\mid\mathbf{B}\models D_{\mathbf{s}}[\mathbf{b}]\}$. The intended interpretation of $\mathbf{d}$ sends $\mathbf{B}$ to the $\Sigma$-algebra with underlying $S$-sorted set $(B^{D}_{s})_{s\in S}$ in which every $\sigma\in\Sigma_{\mathbf{s},s}$ is interpreted as the restriction of the derived operation $d_{\mathbf{s},s}(\sigma)^{\mathbf{B}}$ to $B^{D}_{\mathbf{s}}$. This is a $\Sigma$-algebra only if, for every $\sigma\in\Sigma_{\mathbf{s},s}$, $d_{\mathbf{s},s}(\sigma)^{\mathbf{B}}$ sends $B^{D}_{\mathbf{s}}$ into $B^{D}_{s}$. This closure condition is not part of Definition~\ref{DGeneralDerivor}; a syntactic condition that guarantees it for every $\mathbf{B}$ is condition~(GEP1) of Definition~\ref{DGenEpDerivor}, see Remark~\ref{RGenEpDerivor}. For $D=\varnothing$, $B^{D}_{s}=B_{\varphi(s)}$ and the condition always holds.
% \end{remark}
%
% \section{ep-derivors}\label{SEpDerivors}
%
% Equationally partial signatures, see~\cite{Rei87}, are signatures in which every operation symbol carries a set of existence equations that describes its domain of definition. In this section we state them in our notation, extend the domain conditions from the operation symbols to the terms, and define the morphisms between them that are induced by derivors. We use the notation fixed in Section~\ref{SGenDerivors}.
%
% \begin{definition}\label{DEpSignature}
% An \emph{equationally partial signature}, or \emph{ep-signature}, is an ordered triple $\sig{\Sigma}=(S,\Sigma,\mathrm{def}^{\sig{\Sigma}})$ in which $(S,\Sigma)$ is a many-sorted signature and $\mathrm{def}^{\sig{\Sigma}}$ assigns to every $(\mathbf{s},s)\in S^{\star}\times S$ and every $\sigma\in\Sigma_{\mathbf{s},s}$ a set $\mathrm{def}^{\sig{\Sigma}}(\sigma)\subseteq\mathrm{Eq}_{\Sigma}(\vs{\mathbf{s}})$ of existence equations with variables in $\vs{\mathbf{s}}$, the set of \emph{domain conditions} of $\sigma$.
%
% Let $\mathbf{A}=(A,F)$ be a partial $\Sigma$-algebra, see Definition~\ref{DPartialAlg}. For $(\mathbf{s},s)\in S^{\star}\times S$ and $\sigma\in\Sigma_{\mathbf{s},s}$, we let $A_{\mathrm{def}^{\sig{\Sigma}}(\sigma)}$ stand for the subset of $A_{\mathbf{s}}$ defined as
% $$
% A_{\mathrm{def}^{\sig{\Sigma}}(\sigma)}=\{\mathbf{a}\in A_{\mathbf{s}}\mid\mathbf{A}\models\mathrm{def}^{\sig{\Sigma}}(\sigma)[\mathbf{a}]\}.
% $$
% We say that $\mathbf{A}$ is an \emph{ep-algebra of type $\sig{\Sigma}$}, or a \emph{$\sig{\Sigma}$-algebra}, if, for every $(\mathbf{s},s)\in S^{\star}\times S$ and every $\sigma\in\Sigma_{\mathbf{s},s}$, $\Dom(F_{\sigma})=A_{\mathrm{def}^{\sig{\Sigma}}(\sigma)}$. We denote by $\mathsf{PAlg}(\sig{\Sigma})$ the category of $\sig{\Sigma}$-algebras and the homomorphisms of partial $\Sigma$-algebras between them.
% \end{definition}
%
%
% Note that the domain conditions of $\sigma$ are evaluated in $\mathbf{A}$ itself, so they may involve $\sigma$ and the other partial operations of $\mathbf{A}$. A many-sorted signature $(S,\Sigma)$ is identified with the ep-signature $(S,\Sigma,\varnothing)$ in which every operation symbol has an empty set of domain conditions; its ep-algebras are exactly the (total) $\Sigma$-algebras.
%
%
% \begin{remark}\label{RSubst}
% Let $\mathbf{A}$ be a partial $\Sigma$-algebra, $X$ and $Y$ $S$-sorted sets, $f\colon Y\mor\mathrm{T}_{\Sigma}(X)$ an $S$-sorted mapping and $a\colon X\mor A$ a valuation such that, for every $y\in Y$, $f(y)^{\mathbf{A}}[a]$ is defined. Let $f^{\mathbf{A}}[a]\colon Y\mor A$ be the valuation that sends $y$ to $f(y)^{\mathbf{A}}[a]$. Then, by induction on terms, for every term $R\in\mathrm{T}_{\Sigma}(Y)_{t}$, $f^{\sharp}(R)^{\mathbf{A}}[a]$ is defined if and only if $R^{\mathbf{A}}[f^{\mathbf{A}}[a]]$ is defined, and then they are equal. Consequently, for every $\mathcal{E}\subseteq\mathrm{Eq}_{\Sigma}(Y)$, $\mathbf{A}\models f^{\sharp}[\mathcal{E}][a]$ if and only if $\mathbf{A}\models\mathcal{E}[f^{\mathbf{A}}[a]]$. The hypothesis on $f$ cannot be dropped: if $y$ does not occur in $R$, then $f^{\sharp}(R)^{\mathbf{A}}[a]$ may be defined although $f(y)^{\mathbf{A}}[a]$ is not.
% \end{remark}
%
% To define the morphisms between ep-signatures we extend the domain conditions from the operation symbols to the terms by algebraic recursion, as in~\cite{Rei87}.
%
% \begin{definition}\label{DDeft}
% Let $\sig{\Sigma}=(S,\Sigma,\mathrm{def}^{\sig{\Sigma}})$ be an ep-signature and $X$ an $S$-sorted set. We let $\mathrm{deft}^{\sig{\Sigma}}_{X}$, or simply $\mathrm{deft}^{\sig{\Sigma}}$, stand for the $S$-sorted mapping that sends every term $P\in\mathrm{T}_{\Sigma}(X)_{t}$, $t\in S$, to the set $\mathrm{deft}^{\sig{\Sigma}}(P)\subseteq\mathrm{Eq}_{\Sigma}(X)$ defined by recursion on $P$ as follows:
% \begin{enumerate}
% \item if $P=x$, with $x\in X_{t}$, then $\mathrm{deft}^{\sig{\Sigma}}(x)=\varnothing$;
% \item if $P=\sigma((P_{j})_{j\in\bb{\mathbf{s}}})$, with $\sigma\in\Sigma_{\mathbf{s},t}$ and $(P_{j})_{j\in\bb{\mathbf{s}}}\in\mathrm{T}_{\Sigma}(X)_{\mathbf{s}}$, then
% $$
% \textstyle
% \mathrm{deft}^{\sig{\Sigma}}(\sigma((P_{j})_{j\in\bb{\mathbf{s}}}))=
% \bigcup_{j\in\bb{\mathbf{s}}}\mathrm{deft}^{\sig{\Sigma}}(P_{j})
% \cup
% (P_{j})_{j\in\bb{\mathbf{s}}}^{\sharp}[\mathrm{def}^{\sig{\Sigma}}(\sigma)],
% $$
% where $(P_{j})_{j\in\bb{\mathbf{s}}}^{\sharp}$ is the substitution that sends $v^{\mathbf{s}}_{j}$ to $P_{j}$, for every $j\in\bb{\mathbf{s}}$.
% \end{enumerate}
% \end{definition}
%
% \begin{lemma}\label{LDeft}
% Let $\sig{\Sigma}$ be an ep-signature, $\mathbf{A}$ a $\sig{\Sigma}$-algebra, $X$ an $S$-sorted set, $a\colon X\mor A$ a valuation and $P\in\mathrm{T}_{\Sigma}(X)_{t}$. Then $P^{\mathbf{A}}[a]$ is defined if and only if $\mathbf{A}\models\mathrm{deft}^{\sig{\Sigma}}(P)[a]$.
% \end{lemma}
% \begin{proof}
% By induction on $P$. For a variable both conditions hold. Let $P=\sigma((P_{j})_{j\in\bb{\mathbf{s}}})$. If $P^{\mathbf{A}}[a]$ is defined, then every $P_{j}^{\mathbf{A}}[a]$ is defined, hence, by the induction hypothesis, $\mathbf{A}\models\mathrm{deft}^{\sig{\Sigma}}(P_{j})[a]$, and $\mathbf{c}=(P_{j}^{\mathbf{A}}[a])_{j\in\bb{\mathbf{s}}}$ is in $\Dom(F_{\sigma})=A_{\mathrm{def}^{\sig{\Sigma}}(\sigma)}$, i.e., $\mathbf{A}\models\mathrm{def}^{\sig{\Sigma}}(\sigma)[\mathbf{c}]$, which, by Remark~\ref{RSubst}, means that $\mathbf{A}\models(P_{j})_{j\in\bb{\mathbf{s}}}^{\sharp}[\mathrm{def}^{\sig{\Sigma}}(\sigma)][a]$. Conversely, if $\mathbf{A}\models\mathrm{deft}^{\sig{\Sigma}}(P)[a]$, then, by the induction hypothesis, every $P_{j}^{\mathbf{A}}[a]$ is defined and, by Remark~\ref{RSubst}, $\mathbf{c}\in A_{\mathrm{def}^{\sig{\Sigma}}(\sigma)}=\Dom(F_{\sigma})$, so $P^{\mathbf{A}}[a]$ is defined.
% \end{proof}
%
%
% \begin{definition}\label{DCn}
% Let $\sig{\Sigma}=(S,\Sigma,\mathrm{def}^{\sig{\Sigma}})$ be an ep-signature and $X$ an $S$-sorted set. The \emph{semantic consequence operator} for $\sig{\Sigma}$ and $X$ is the mapping $\mathrm{Cn}^{\sig{\Sigma}}_{X}$ from $\mathrm{Sub}(\mathrm{Eq}_{\Sigma}(X))$ to $\mathrm{Sub}(\mathrm{Eq}_{\Sigma}(X))$ that sends $\mathcal{E}\subseteq\mathrm{Eq}_{\Sigma}(X)$ to
% \begin{multline*}
% \mathrm{Cn}^{\sig{\Sigma}}_{X}(\mathcal{E})=
% \{(P,Q)\in\mathrm{Eq}_{\Sigma}(X)\mid
% \text{for every $\sig{\Sigma}$-algebra $\mathbf{A}$ and every $a\colon X\mor A$,}\\
% \text{if $\mathbf{A}\models\mathcal{E}[a]$, then $\mathbf{A}\models(P,Q)[a]$}\}.
% \end{multline*}
% \end{definition}
%
% $\mathrm{Cn}^{\sig{\Sigma}}_{X}$ is a closure operator on $\mathrm{Eq}_{\Sigma}(X)$, i.e., it is extensive, monotone and idempotent. Moreover, for $\mathcal{E},\mathcal{E}'\subseteq\mathrm{Eq}_{\Sigma}(X)$, $\mathrm{Cn}^{\sig{\Sigma}}_{X}(\mathcal{E})=\mathrm{Cn}^{\sig{\Sigma}}_{X}(\mathcal{E}')$ if and only if, for every $\sig{\Sigma}$-algebra $\mathbf{A}$ and every valuation $a\colon X\mor A$, $\mathbf{A}\models\mathcal{E}[a]$ if and only if $\mathbf{A}\models\mathcal{E}'[a]$. In fact, if $\mathrm{Cn}^{\sig{\Sigma}}_{X}(\mathcal{E})=\mathrm{Cn}^{\sig{\Sigma}}_{X}(\mathcal{E}')$ and $\mathbf{A}\models\mathcal{E}'[a]$, then $\mathbf{A}\models(P,Q)[a]$ for every $(P,Q)\in\mathrm{Cn}^{\sig{\Sigma}}_{X}(\mathcal{E}')=\mathrm{Cn}^{\sig{\Sigma}}_{X}(\mathcal{E})\supseteq\mathcal{E}$; the converse is immediate.
%
% Let $\mathbf{d}=(\varphi,d)$ be a derivor from $(S,\Sigma)$ to $(T,\Lambda)$. Recall that $d^{\sharp}$, the homomorphism of Hall algebras determined by $d$, sends, for every $(\mathbf{s},t)\in S^{\star}\times S$, the terms in $\mathrm{T}_{\Sigma}(\vs{\mathbf{s}})_{t}$ to terms in $\mathrm{T}_{\Lambda}(\vs{\varphi^{\star}(\mathbf{s})})_{\varphi(t)}$; since it is a homomorphism of Hall algebras, $d^{\sharp}(v^{\mathbf{s}}_{j})=v^{\varphi^{\star}(\mathbf{s})}_{j}$ and $d^{\sharp}(\sigma((R_{j})_{j\in\bb{\mathbf{r}}}))=(d^{\sharp}(R_{j}))_{j\in\bb{\mathbf{r}}}^{\sharp}(d_{\mathbf{r},t}(\sigma))$, for $\sigma\in\Sigma_{\mathbf{r},t}$. For a partial $\Lambda$-algebra $\mathbf{B}$, we let $\mathbf{d}^{\ast}(\mathbf{B})$ stand for the partial $\Sigma$-algebra whose underlying $S$-sorted set is $(B_{\varphi(s)})_{s\in S}$ and in which every $\sigma\in\Sigma_{\mathbf{s},s}$ is interpreted as the partial term operation $d_{\mathbf{s},s}(\sigma)^{\mathbf{B}}\colon B_{\varphi^{\star}(\mathbf{s})}\dmor B_{\varphi(s)}$, which sends $\mathbf{b}$ to $d_{\mathbf{s},s}(\sigma)^{\mathbf{B}}[\mathbf{b}]$ whenever the latter is defined.
%
%
% As the last sentence of Remark~\ref{RSubst} shows, the existence equation $(d^{\sharp}(P),d^{\sharp}(Q))$ need not express, in $\mathbf{B}$, that $(P,Q)$ holds in $\mathbf{d}^{\ast}(\mathbf{B})$: if $d_{\mathbf{s},s}(\sigma)$ does not contain the variable $v^{\varphi^{\star}(\mathbf{s})}_{j}$, then $d^{\sharp}(\sigma((P_{i})_{i}))$ can be defined in $\mathbf{B}$ although $P_{j}$, and hence $\sigma((P_{i})_{i})$, is not defined in $\mathbf{d}^{\ast}(\mathbf{B})$. For this reason we also translate the definedness of every subterm. For a term $P$, we let $\mathrm{Sb}(P)$ stand for the set of its subterms, defined by $\mathrm{Sb}(x)=\{x\}$ and $\mathrm{Sb}(\sigma((P_{j})_{j}))=\{\sigma((P_{j})_{j})\}\cup\bigcup_{j}\mathrm{Sb}(P_{j})$.
%
% \begin{definition}\label{DTranslation}
% Let $\mathbf{d}=(\varphi,d)$ be a derivor from $(S,\Sigma)$ to $(T,\Lambda)$ and $\mathbf{s}\in S^{\star}$. The \emph{translation along $\mathbf{d}$} of an existence equation $(P,Q)\in\mathrm{Eq}_{\Sigma}(\vs{\mathbf{s}})$ is the set
% $$
% \mathrm{tr}_{\mathbf{d}}(P,Q)=\{(d^{\sharp}(P),d^{\sharp}(Q))\}\cup\{(d^{\sharp}(R),d^{\sharp}(R))\mid R\in\mathrm{Sb}(P)\cup\mathrm{Sb}(Q)\}\subseteq\mathrm{Eq}_{\Lambda}(\vs{\varphi^{\star}(\mathbf{s})}),
% $$
% and, for $\mathcal{E}\subseteq\mathrm{Eq}_{\Sigma}(\vs{\mathbf{s}})$, $\mathrm{tr}_{\mathbf{d}}[\mathcal{E}]=\bigcup_{(P,Q)\in\mathcal{E}}\mathrm{tr}_{\mathbf{d}}(P,Q)$.
% \end{definition}
%
% If, for every $\sigma\in\Sigma_{\mathbf{s},s}$, every variable of $\vs{\varphi^{\star}(\mathbf{s})}$ occurs in $d_{\mathbf{s},s}(\sigma)$, then every $d^{\sharp}(R)$, for $R\in\mathrm{Sb}(P)$, is a subterm of $d^{\sharp}(P)$, and the existence equation $(d^{\sharp}(P),d^{\sharp}(Q))$ alone has the same models as $\mathrm{tr}_{\mathbf{d}}(P,Q)$.
%
% \begin{lemma}\label{LTranslation}
% Let $\mathbf{d}=(\varphi,d)$ be a derivor from $(S,\Sigma)$ to $(T,\Lambda)$, $\mathbf{B}$ a partial $\Lambda$-algebra, $\mathbf{s}\in S^{\star}$ and $\mathbf{b}\in B_{\varphi^{\star}(\mathbf{s})}$. Then, for every $R\in\mathrm{T}_{\Sigma}(\vs{\mathbf{s}})_{t}$,
% \begin{enumerate}
% \item if $R^{\mathbf{d}^{\ast}(\mathbf{B})}[\mathbf{b}]$ is defined, then $d^{\sharp}(R)^{\mathbf{B}}[\mathbf{b}]$ is defined and equal to it;
% \item if $d^{\sharp}(R')^{\mathbf{B}}[\mathbf{b}]$ is defined for every $R'\in\mathrm{Sb}(R)$, then $R^{\mathbf{d}^{\ast}(\mathbf{B})}[\mathbf{b}]$ is defined.
% \end{enumerate}
% Consequently, for every $\mathcal{E}\subseteq\mathrm{Eq}_{\Sigma}(\vs{\mathbf{s}})$, $\mathbf{d}^{\ast}(\mathbf{B})\models\mathcal{E}[\mathbf{b}]$ if and only if $\mathbf{B}\models\mathrm{tr}_{\mathbf{d}}[\mathcal{E}][\mathbf{b}]$.
% \end{lemma}
% \begin{proof}
% Both statements are proved by induction on $R$. For $R=v^{\mathbf{s}}_{j}$ both values are $b_{j}$. Let $R=\sigma((R_{j})_{j\in\bb{\mathbf{r}}})$, with $\sigma\in\Sigma_{\mathbf{r},t}$, so that $d^{\sharp}(R)=(d^{\sharp}(R_{j}))_{j\in\bb{\mathbf{r}}}^{\sharp}(d_{\mathbf{r},t}(\sigma))$. (1) If $R^{\mathbf{d}^{\ast}(\mathbf{B})}[\mathbf{b}]$ is defined, then so is every $c_{j}=R_{j}^{\mathbf{d}^{\ast}(\mathbf{B})}[\mathbf{b}]$, which, by the induction hypothesis, is $d^{\sharp}(R_{j})^{\mathbf{B}}[\mathbf{b}]$, and $d_{\mathbf{r},t}(\sigma)^{\mathbf{B}}[\mathbf{c}]$ is defined and equal to $R^{\mathbf{d}^{\ast}(\mathbf{B})}[\mathbf{b}]$; by Remark~\ref{RSubst}, $d^{\sharp}(R)^{\mathbf{B}}[\mathbf{b}]$ is defined and equal to $d_{\mathbf{r},t}(\sigma)^{\mathbf{B}}[\mathbf{c}]$. (2) By the induction hypothesis every $c_{j}=R_{j}^{\mathbf{d}^{\ast}(\mathbf{B})}[\mathbf{b}]$ is defined and, by~(1), equal to $d^{\sharp}(R_{j})^{\mathbf{B}}[\mathbf{b}]$; since $d^{\sharp}(R)^{\mathbf{B}}[\mathbf{b}]$ is defined, Remark~\ref{RSubst} shows that $d_{\mathbf{r},t}(\sigma)^{\mathbf{B}}[\mathbf{c}]$ is defined, i.e., that $R^{\mathbf{d}^{\ast}(\mathbf{B})}[\mathbf{b}]$ is defined.
%
% For the last statement, if $(P,Q)$ holds in $\mathbf{d}^{\ast}(\mathbf{B})$ under $\mathbf{b}$, then every subterm of $P$ and of $Q$ is defined in $\mathbf{d}^{\ast}(\mathbf{B})$ under $\mathbf{b}$ and, by~(1), $\mathbf{B}\models\mathrm{tr}_{\mathbf{d}}(P,Q)[\mathbf{b}]$. Conversely, if $\mathbf{B}\models\mathrm{tr}_{\mathbf{d}}(P,Q)[\mathbf{b}]$, then, by~(2), $P$ and $Q$ are defined in $\mathbf{d}^{\ast}(\mathbf{B})$ under $\mathbf{b}$ and, by~(1), their values are $d^{\sharp}(P)^{\mathbf{B}}[\mathbf{b}]=d^{\sharp}(Q)^{\mathbf{B}}[\mathbf{b}]$.
% \end{proof}
%
% \begin{definition}\label{DDerivorEP}
% Let $\sig{\Sigma}=(S,\Sigma,\mathrm{def}^{\sig{\Sigma}})$ and $\sig{\Lambda}=(T,\Lambda,\mathrm{def}^{\sig{\Lambda}})$ be ep-signatures. An \emph{ep-derivor} from $\sig{\Sigma}$ to $\sig{\Lambda}$ is a derivor $\mathbf{d}=(\varphi,d)$ from $(S,\Sigma)$ to $(T,\Lambda)$ such that, for every $(\mathbf{s},s)\in S^{\star}\times S$ and every $\sigma\in\Sigma_{\mathbf{s},s}$,
% $$
% \mathrm{Cn}^{\sig{\Lambda}}_{\vs{\varphi^{\star}(\mathbf{s})}}\left(\mathrm{tr}_{\mathbf{d}}[\mathrm{def}^{\sig{\Sigma}}(\sigma)]\right)
% =
% \mathrm{Cn}^{\sig{\Lambda}}_{\vs{\varphi^{\star}(\mathbf{s})}}\left(\mathrm{deft}^{\sig{\Lambda}}(d_{\mathbf{s},s}(\sigma))\right).
% $$
% \end{definition}
%
% \begin{proposition}\label{PEpDerivorFunctor}
% Let $\sig{\Sigma}$ and $\sig{\Lambda}$ be ep-signatures and $\mathbf{d}=(\varphi,d)$ a derivor from $(S,\Sigma)$ to $(T,\Lambda)$. Then $\mathbf{d}$ is an ep-derivor from $\sig{\Sigma}$ to $\sig{\Lambda}$ if and only if, for every $\sig{\Lambda}$-algebra $\mathbf{B}$, $\mathbf{d}^{\ast}(\mathbf{B})$ is a $\sig{\Sigma}$-algebra. In this case, $\mathbf{d}^{\ast}$ is the object mapping of a functor $\mathbf{d}^{\ast}\colon\mathsf{PAlg}(\sig{\Lambda})\mor\mathsf{PAlg}(\sig{\Sigma})$ that sends a homomorphism $h\colon\mathbf{B}\mor\mathbf{B}'$ to $(h_{\varphi(s)})_{s\in S}$.
% \end{proposition}
% \begin{proof}
% Let $\mathbf{B}$ be a $\sig{\Lambda}$-algebra, $\sigma\in\Sigma_{\mathbf{s},s}$ and $\mathbf{b}\in B_{\varphi^{\star}(\mathbf{s})}$. By Lemma~\ref{LDeft}, $\mathbf{b}\in\Dom(d_{\mathbf{s},s}(\sigma)^{\mathbf{B}})$ if and only if $\mathbf{B}\models\mathrm{deft}^{\sig{\Lambda}}(d_{\mathbf{s},s}(\sigma))[\mathbf{b}]$, and, by Lemma~\ref{LTranslation}, $\mathbf{b}$ satisfies the domain conditions of $\sigma$ in $\mathbf{d}^{\ast}(\mathbf{B})$, i.e., $\mathbf{d}^{\ast}(\mathbf{B})\models\mathrm{def}^{\sig{\Sigma}}(\sigma)[\mathbf{b}]$, if and only if $\mathbf{B}\models\mathrm{tr}_{\mathbf{d}}[\mathrm{def}^{\sig{\Sigma}}(\sigma)][\mathbf{b}]$. Hence $\mathbf{d}^{\ast}(\mathbf{B})$ is a $\sig{\Sigma}$-algebra for every $\mathbf{B}$ if and only if, for every $\sigma$, the sets $\mathrm{tr}_{\mathbf{d}}[\mathrm{def}^{\sig{\Sigma}}(\sigma)]$ and $\mathrm{deft}^{\sig{\Lambda}}(d_{\mathbf{s},s}(\sigma))$ have the same models, which, as noted after Definition~\ref{DCn}, means that they have the same semantic consequences.
%
% Let $h\colon\mathbf{B}\mor\mathbf{B}'$ be a homomorphism of partial $\Lambda$-algebras. By induction on terms, if $R^{\mathbf{B}}[\mathbf{b}]$ is defined, then $R^{\mathbf{B}'}[h\circ\mathbf{b}]$ is defined and equal to $h(R^{\mathbf{B}}[\mathbf{b}])$. Applied to the terms $d_{\mathbf{s},s}(\sigma)$, this shows that $(h_{\varphi(s)})_{s\in S}$ is a homomorphism from $\mathbf{d}^{\ast}(\mathbf{B})$ to $\mathbf{d}^{\ast}(\mathbf{B}')$. Identities and composites are clearly preserved.
% \end{proof}
%
% \begin{remark}\label{REpDerivor}
% An earlier draft of Definition~\ref{DDerivorEP} only required that $(d^{\sharp}\times d^{\sharp})[\mathrm{def}^{\sig{\Sigma}}(\sigma)]\subseteq\mathrm{Cn}^{\sig{\Lambda}}(\mathrm{deft}^{\sig{\Lambda}}(d_{\mathbf{s},s}(\sigma)))$. Such an inclusion only guarantees that every family in $\Dom(d_{\mathbf{s},s}(\sigma)^{\mathbf{B}})$ satisfies the domain conditions of $\sigma$ in $\mathbf{d}^{\ast}(\mathbf{B})$, so the operations of $\mathbf{d}^{\ast}(\mathbf{B})$ could be undefined on families that satisfy their domain conditions, and $\mathbf{d}^{\ast}$ would not send $\sig{\Lambda}$-algebras to $\sig{\Sigma}$-algebras. By Proposition~\ref{PEpDerivorFunctor}, the equality of the semantic consequences, which is the condition required in~\cite{Rei87}, is exactly what is needed for $\mathbf{d}^{\ast}$ to be a functor from $\mathsf{PAlg}(\sig{\Lambda})$ to $\mathsf{PAlg}(\sig{\Sigma})$. Moreover, the translation $\mathrm{tr}_{\mathbf{d}}$ replaces $d^{\sharp}\times d^{\sharp}$ because of the phenomenon described before Definition~\ref{DTranslation}.
%
% If $\sig{\Sigma}$ and $\sig{\Lambda}$ are many-sorted signatures, regarded as ep-signatures with empty sets of domain conditions, then $\mathrm{tr}_{\mathbf{d}}[\varnothing]=\varnothing$ and $\mathrm{deft}^{\sig{\Lambda}}(P)=\varnothing$ for every term $P$, so that every derivor is an ep-derivor and $\mathbf{d}^{\ast}$ is the usual functor between categories of algebras.
% \end{remark}

\section{Generalized hep-derivors}\label{SGenHepDerivors}

The general derivors of~\cite{LG18} are morphisms between order-sorted signatures with predicates which send every sort to a formula that delimits, inside the interpretation of the image sort, the elements that interpret it, and every operation symbol to a formula. The equationally partial signatures of~\cite{Rei87} are signatures in which every operation symbol carries a set of existence equations that describes its domain of definition. In this section we combine both ideas in the notion of generalized hep-derivor. We need neither the order on the sorts nor the predicates of~\cite{LG18} and, instead of arbitrary formulas, we use existence equations: a generalized hep-derivor sends every sort to a set of existence equations in one variable, which delimits the elements that interpret the sort, and every operation symbol to a term, whose domain in the transformed algebras is determined by the domain conditions of the operation symbol.

We begin by recalling the notions of ep-signature and of ep-algebra, and by introducing the hierarchical ep-signatures, or hep-signatures, together with their specifications, their ep-varieties and the categories canonically associated to them. Next we define the generalized hep-derivors and we prove that every generalized hep-derivor induces a functor between the categories of partial algebras. Then we explain how we check that the restriction of such a functor to the category canonically associated to the ep-variety of a hep-specification corestricts to the category canonically associated to the ep-variety of another one. We end this section by proving that the hep-signatures and the generalized hep-derivors form a category, $\mathsf{hepSig}_{\gd}$, by comparing the functor induced by a composition with the composite of the induced functors, and by assembling, by means of the Grothendieck construction for a normal lax indexed category, the ep-algebras of all the hep-signatures into a category $\mathsf{hepAlg}_{\gd}$. We use the notions on terms, substitutions and derivors recalled in Subsections~\ref{SVariables}, \ref{SPrelimDerivors} and~\ref{STermsEq}.

\subsection{ep-signatures, hep-signatures and their specifications}\label{SEpSignatures}

We begin by defining the notion of ep-signature.

\begin{definition}\label{DEpSignature}
An \emph{equationally partial signature}, or \emph{ep-signature}, is an ordered triple $\sig{\Sigma}=(S,\Sigma,\mathrm{def}^{(S,\Sigma)})$ in which $(S,\Sigma)$ is a many-sorted signature and $\mathrm{def}^{(S,\Sigma)}$ assigns to every $(\mathbf{s},s)\in S^{\star}\times S$ and every $\sigma\in\Sigma_{\mathbf{s},s}$ a set
$$
\mathrm{def}^{(S,\Sigma)}(\sigma)\subseteq\mathrm{Eeqt}(\Sigma)_{\vs{\mathbf{s}}}
$$
of $\mathrm{E}$-equations with variables in $\vs{\mathbf{s}}$. We will call $\mathrm{def}^{(S,\Sigma)}(\sigma)$ the set of \emph{domain conditions} of $\sigma$.
\end{definition}

In the following definition we introduce the partial algebras whose operations are defined exactly on the families that satisfy the corresponding domain conditions.

\begin{definition}\label{DEpAlgebra}
Let $\sig{\Sigma}=(S,\Sigma,\mathrm{def}^{(S,\Sigma)})$ be an ep-signature and $\mathbf{A}$ a partial $\Sigma$-algebra, see Definition~\ref{DPartialAlg}. For every $(\mathbf{s},s)\in S^{\star}\times S$ and every $\sigma\in\Sigma_{\mathbf{s},s}$, we let $A_{\mathrm{def}^{(S,\Sigma)}(\sigma)}$ stand for the subset of $A_{\mathbf{s}}$ defined as
$$
A_{\mathrm{def}^{(S,\Sigma)}(\sigma)}=\left\lbrace\mathbf{a}\in A_{\mathbf{s}}\,\middle|\,\mathbf{A}\models\mathrm{def}^{(S,\Sigma)}(\sigma)[\mathbf{a}]\right\rbrace.
$$
We will say that $\mathbf{A}$ is an \emph{ep-algebra of type $\sig{\Sigma}$}, or a \emph{$\sig{\Sigma}$-algebra}, if, for every $(\mathbf{s},s)\in S^{\star}\times S$ and every $\sigma\in\Sigma_{\mathbf{s},s}$, the following equality holds
$$
\Dom\left(\sigma^{\mathbf{A}}\right)=A_{\mathrm{def}^{(S,\Sigma)}(\sigma)}.
$$
We will denote by $\mathsf{EPAlg}(\sig{\Sigma})$ the full subcategory of $\mathsf{PAlg}(\Sigma)$, the category of partial $\Sigma$-algebras and their homomorphisms, see Definition~\ref{DPartialAlg}, whose objects are the ep-algebras of type $\sig{\Sigma}$.
\end{definition}

\begin{remark}\label{REpAlgebra}
Let us note that the domain conditions of an operation symbol $\sigma$ are evaluated in $\mathbf{A}$ itself, so they may involve the other partial operations of $\mathbf{A}$, and even $\sigma$ itself. Let us also note that a many-sorted signature $(S,\Sigma)$ is identified with the ep-signature $(S,\Sigma,\varnothing)$, in which every operation symbol has an empty set of domain conditions. Its ep-algebras are exactly the (total) $\Sigma$-algebras, so that
$$
\mathsf{EPAlg}(S,\Sigma,\varnothing)=\mathsf{Alg}(\Sigma).
$$
\end{remark}

In order to assemble the ep-algebras of all the hep-signatures into a single category by means of the Grothendieck construction, see Subsection~\ref{SHepAlgCat}, we will need the following weakening of the notion of ep-algebra, in which the operations are only required to be defined on families that satisfy the corresponding domain conditions.

\begin{definition}\label{DWeakEpAlgebra}
Let $\sig{\Sigma}=(S,\Sigma,\mathrm{def}^{(S,\Sigma)})$ be an ep-signature and $\mathbf{A}$ a partial $\Sigma$-algebra. We will say that $\mathbf{A}$ is a \emph{weak ep-algebra of type $\sig{\Sigma}$}, or a \emph{weak $\sig{\Sigma}$-algebra}, if, for every $(\mathbf{s},s)\in S^{\star}\times S$ and every $\sigma\in\Sigma_{\mathbf{s},s}$, the following inclusion holds
$$
\Dom\left(\sigma^{\mathbf{A}}\right)\subseteq A_{\mathrm{def}^{(S,\Sigma)}(\sigma)}.
$$
We will denote by $\mathsf{wEPAlg}(\sig{\Sigma})$ the full subcategory of $\mathsf{PAlg}(\Sigma)$ whose objects are the weak ep-algebras of type $\sig{\Sigma}$.
\end{definition}

Let us note that every $\sig{\Sigma}$-algebra is a weak $\sig{\Sigma}$-algebra, so that $\mathsf{EPAlg}(\sig{\Sigma})$ is a full subcategory of $\mathsf{wEPAlg}(\sig{\Sigma})$.

In order to define the algebras transformed along a generalized hep-derivor by recursion, we will need the domain conditions to be organized hierarchically. To this end, we first assign to every term the set of the operation symbols that occur in it.

\begin{definition}\label{DOpSymbols}
Let $\Sigma$ be an $S$-sorted signature and $X$ an $S$-sorted set. We let $\mathrm{Op}$ stand for the mapping that assigns to every $s\in S$ and every term $P\in\mathrm{T}_{\Sigma}(X)_{s}$ the set $\mathrm{Op}(P)$ of the \emph{operation symbols that occur in $P$}, defined by recursion on $P$ as follows:
\begin{enumerate}
\item if $P=x$ is a variable, then $\mathrm{Op}(P)=\varnothing$; and
\item if $P=\sigma((P_{j})_{j\in\bb{\mathbf{s}}})$, with $(\mathbf{s},s)\in S^{\star}\times S$ and $\sigma\in\Sigma_{\mathbf{s},s}$, then
$$
\mathrm{Op}(P)=\{\sigma\}\cup\textstyle\bigcup_{j\in\bb{\mathbf{s}}}\mathrm{Op}(P_{j}).
$$
\end{enumerate}
Moreover, for every $s\in S$ and every $\mathrm{E}$-equation $(P,Q)\in\mathrm{Eeqt}(\Sigma)_{X,s}$, we let $\mathrm{Op}(P,Q)$ stand for $\mathrm{Op}(P)\cup\mathrm{Op}(Q)$ and, for every set of $\mathrm{E}$-equations $\mathcal{H}\subseteq\mathrm{Eeqt}(\Sigma)_{X}$, we let $\mathrm{Op}(\mathcal{H})$ stand for $\bigcup_{s\in S}\bigcup_{(P,Q)\in\mathcal{H}_{s}}\mathrm{Op}(P,Q)$.
\end{definition}

We are now in position to introduce the hierarchical ep-signatures.

\begin{definition}\label{DHepSignature}
Let $\sig{\Sigma}=(S,\Sigma,\mathrm{def}^{(S,\Sigma)})$ be an ep-signature. We let $\lhd$ stand for the binary relation on the set of the operation symbols of $\Sigma$ defined, for every pair of operation symbols $\tau$ and $\sigma$ of $\Sigma$, as follows:
\begin{align*}
\tau&\lhd\sigma
&\mbox{if, and only if,}&&
&\tau\in\mathrm{Op}\left(\mathrm{def}^{(S,\Sigma)}(\sigma)\right),
\end{align*}
see Definition~\ref{DOpSymbols}, i.e., if, and only if, $\tau$ occurs in some equation of $\mathrm{def}^{(S,\Sigma)}(\sigma)$. We will say that $\sig{\Sigma}$ is \emph{hierarchical}, or that it is a \emph{hierarchical equationally partial signature}, or a \emph{hep-signature}, if $\lhd$ is well-founded.
\end{definition}

Let us note that, in a hep-signature, no operation symbol occurs in its own domain conditions, since, otherwise, $\sigma\lhd\sigma$ would give rise to an infinite descending chain.

We next introduce the specifications on ep-signatures.

\begin{definition}\label{DEpSpec}
An \emph{ep-specification} is an ordered pair $(\sig{\Sigma},\mathcal{E})$, where $\sig{\Sigma}=(S,\Sigma,\mathrm{def}^{(S,\Sigma)})$ is an ep-signature and $\mathcal{E}\subseteq\mathrm{QEeqt}(\Sigma)_{V^{S}}$ is a set of $\QE$-equations, see Definition~\ref{DQEEq}. We will say that $(\sig{\Sigma},\mathcal{E})$ is a \emph{hep-specification} if $\sig{\Sigma}$ is a hep-signature, see Definition~\ref{DHepSignature}. An ep-signature $\sig{\Sigma}$ is identified with the ep-specification $(\sig{\Sigma},\varnothing)$, and a hep-signature with the hep-specification $(\sig{\Sigma},\varnothing)$.
\end{definition}

Following the case of the $\QE$-specifications, see Definition~\ref{DQEEq}, an ep-specification determines a class of ep-algebras, its ep-variety, and a category canonically associated to it.

\begin{definition}\label{DEpVariety}
Let $(\sig{\Sigma},\mathcal{E})$ be an ep-specification. We will say that a $\sig{\Sigma}$-algebra $\mathbf{A}$ is a \emph{model} of $(\sig{\Sigma},\mathcal{E})$ if every equation of $\mathcal{E}$ holds in $\mathbf{A}$. We let $\mathcal{V}(\sig{\Sigma},\mathcal{E})$ stand for the \emph{ep-variety} determined by the ep-specification $(\sig{\Sigma},\mathcal{E})$, i.e., the class of all the models of $(\sig{\Sigma},\mathcal{E})$, and $\mathsf{EPAlg}(\sig{\Sigma},\mathcal{E})$ for the category canonically associated to $\mathcal{V}(\sig{\Sigma},\mathcal{E})$, i.e., the full subcategory of $\mathsf{EPAlg}(\sig{\Sigma})$ whose objects are the $\sig{\Sigma}$-algebras in $\mathcal{V}(\sig{\Sigma},\mathcal{E})$.
\end{definition}

Let us note that $\mathcal{V}(\sig{\Sigma},\varnothing)$ is the class of all the $\sig{\Sigma}$-algebras, so that $\mathsf{EPAlg}(\sig{\Sigma},\varnothing)=\mathsf{EPAlg}(\sig{\Sigma})$.

\subsection{Generalized hep-derivors}\label{SGenHepDerivorsDef}

We next define the notion of generalized hep-derivor.

\begin{definition}\label{DGenHepDerivor}
Let $\sig{\Sigma}=(S,\Sigma,\mathrm{def}^{(S,\Sigma)})$ and $\sig{\Lambda}=(T,\Lambda,\mathrm{def}^{(T,\Lambda)})$ be hep-signatures, see Definition~\ref{DHepSignature}. A \emph{generalized hep-derivor} from $\sig{\Sigma}$ to $\sig{\Lambda}$ is an ordered triple $\mathbf{d}=(\varphi,d,\Dom^{(\varphi,d)})$ such that
\begin{enumerate}
\item $(\varphi,d)$ is a derivor from $(S,\Sigma)$ to $(T,\Lambda)$, see Subsection~\ref{SPrelimDerivors}, i.e., $\varphi$ is a mapping from $S$ to $T$ and $d$ assigns to every $(\mathbf{s},s)\in S^{\star}\times S$ and every $\sigma\in\Sigma_{\mathbf{s},s}$ a term $d(\sigma)\in\mathrm{T}_{\Lambda}(\vs{\varphi^{\star}(\mathbf{s})})_{\varphi(s)}$; and
\item $\Dom^{(\varphi,d)}=(\Dom^{(\varphi,d)}_{s})_{s\in S}$ assigns to every sort $s\in S$ a set
$$
\Dom^{(\varphi,d)}_{s}\subseteq\mathrm{Eeqt}(\Lambda)_{\vs{\varphi(s)}}
$$
of $\mathrm{E}$-equations in the single variable $v^{\varphi(s)}_{0}$, the \emph{domain} of $s$.
\end{enumerate}
\end{definition}

We next associate to every partial $\Lambda$-algebra the $S$-sorted set of the elements that satisfy the domains of a generalized hep-derivor. Since it is obtained from the $S$-sorted set $B_{\varphi}=(B_{\varphi(s)})_{s\in S}$ by restricting every $B_{\varphi(s)}$ to the elements that satisfy $\Dom^{(\varphi,d)}_{s}$, we denote it, as the sets $A_{\mathrm{def}^{(S,\Sigma)}(\sigma)}$ of Definition~\ref{DEpAlgebra}, by subscripting the domains. Let us recall that, by the identifications of Subsection~\ref{SVariables}, for every $s\in S$, an element $b\in B_{\varphi(s)}$ is identified with the valuation of $\vs{(\varphi(s))}$ in $B$ that sends $v^{\varphi(s)}_{0}$ to $b$.

\begin{definition}\label{DDomainSet}
Let $\sig{\Sigma}=(S,\Sigma,\mathrm{def}^{(S,\Sigma)})$ and $\sig{\Lambda}=(T,\Lambda,\mathrm{def}^{(T,\Lambda)})$ be hep-signatures, $\mathbf{d}=(\varphi,d,\Dom^{(\varphi,d)})$ a generalized hep-derivor from $\sig{\Sigma}$ to $\sig{\Lambda}$ and $\mathbf{B}$ a partial $\Lambda$-algebra. We let $B_{\Dom^{(\varphi,d)}}$ stand for the $S$-sorted subset of $B_{\varphi}$ defined, for every $s\in S$, as
$$
\left(B_{\Dom^{(\varphi,d)}}\right)_{s}=\left\lbrace b\in B_{\varphi(s)}\,\middle|\,\mathbf{B}\models\Dom^{(\varphi,d)}_{s}[b]\right\rbrace.
$$
The elements of $(B_{\Dom^{(\varphi,d)}})_{s}$ will be called the \emph{cells of sort $s$} of $\mathbf{B}$ relative to $\mathbf{d}$.
\end{definition}

Let $\sig{\Sigma}=(S,\Sigma,\mathrm{def}^{(S,\Sigma)})$ and $\sig{\Lambda}=(T,\Lambda,\mathrm{def}^{(T,\Lambda)})$ be hep-signatures, $\mathbf{d}=(\varphi,d,\Dom^{(\varphi,d)})$ a generalized hep-derivor from $\sig{\Sigma}$ to $\sig{\Lambda}$ and $f\colon\mathbf{B}\mor\mathbf{B}'$ a homomorphism of partial $\Lambda$-algebras. Then, for every $s\in S$, $f_{\varphi(s)}{\,\!\upharpoonright}_{(B_{\Dom^{(\varphi,d)}})_{s}}$ corestricts to $(B'_{\Dom^{(\varphi,d)}})_{s}$, see Proposition~\ref{PUS}. Indeed, let $b$ be a cell of sort $s$ of $\mathbf{B}$, i.e., $\mathbf{B}\models\Dom^{(\varphi,d)}_{s}[b]$. Then, by Remark~\ref{RHomTerm}, it happens that
$$
\mathbf{B}'\models\Dom^{(\varphi,d)}_{s}[f_{\varphi(s)}(b)],
$$
i.e., $f_{\varphi(s)}(b)\in(B'_{\Dom^{(\varphi,d)}})_{s}$. Taking this into account, we let $f_{\Dom^{(\varphi,d)}}$ stand for the $S$-sorted mapping
$$
f_{\Dom^{(\varphi,d)}}=f_{\varphi}\bigr|^{B'_{\Dom^{(\varphi,d)}}}_{B_{\Dom^{(\varphi,d)}}}=\left(f_{\varphi(s)}\bigr|^{(B'_{\Dom^{(\varphi,d)}})_{s}}_{(B_{\Dom^{(\varphi,d)}})_{s}}\right)_{s\in S},
$$
i.e., the family of the corresponding birestrictions. Thus $f_{\Dom^{(\varphi,d)}}\colon B_{\Dom^{(\varphi,d)}}\mor B'_{\Dom^{(\varphi,d)}}$.

In the following proposition we prove that the above construction is functorial.

\begin{proposition}\label{PDomainSetFunctor}
Let $\sig{\Sigma}=(S,\Sigma,\mathrm{def}^{(S,\Sigma)})$ and $\sig{\Lambda}=(T,\Lambda,\mathrm{def}^{(T,\Lambda)})$ be hep-signatures, $\mathbf{d}=(\varphi,d,\Dom^{(\varphi,d)})$ a generalized hep-derivor from $\sig{\Sigma}$ to $\sig{\Lambda}$. Then $\G_{\mathbf{d}}$, the assignment from $\mathsf{PAlg}(\Lambda)$ to $\mathsf{Set}^{S}$ defined as follows:
\begin{enumerate}
\item for every partial $\Lambda$-algebra $\mathbf{B}$, $\G_{\mathbf{d}}(\mathbf{B})$ is the $S$-sorted set $B_{\Dom^{(\varphi,d)}}$, introduced in Definition~\ref{DDomainSet}; and
\item for every homomorphism $f\colon\mathbf{B}\mor\mathbf{B}'$ of partial $\Lambda$-algebras, $\G_{\mathbf{d}}(f)$ is the $S$-sorted mapping $f_{\Dom^{(\varphi,d)}}$ from $B_{\Dom^{(\varphi,d)}}$ to $B'_{\Dom^{(\varphi,d)}}$;
\end{enumerate}
is a functor.
\end{proposition}
\begin{proof}
By the discussion after Definition~\ref{DDomainSet}, it suffices to prove that $\G_{\mathbf{d}}$ preserves identities and composites.

Let $\mathbf{B}$ be a partial $\Lambda$-algebra. For every $s\in S$, the following chain of equalities holds
\allowdisplaybreaks
\begin{align*}
\G_{\mathbf{d}}\left(\id_{\mathbf{B}}\right)_{s}
&=
\left(\id_{B}\right)_{\varphi(s)}\bigr|^{(B_{\Dom^{(\varphi,d)}})_{s}}_{(B_{\Dom^{(\varphi,d)}})_{s}}
\tag{1}
\\&=
\id_{(B_{\Dom^{(\varphi,d)}})_{s}}
\tag{2}
\\&=
\left(\id_{\G_{\mathbf{d}}(\mathbf{B})}\right)_{s}.
\tag{3}
\end{align*}

The first equality unravels the definition of $\G_{\mathbf{d}}$ on morphisms; the second equality follows from the fact that the birestriction of an identity mapping is an identity mapping; finally, the last equality unravels the definition of $\G_{\mathbf{d}}$ on objects.

Thus, $\G_{\mathbf{d}}$ preserves identities.

Let $f\colon\mathbf{B}\mor\mathbf{B}'$ and $g\colon\mathbf{B}'\mor\mathbf{B}''$ be homomorphisms of partial $\Lambda$-algebras. For every $s\in S$, the following chain of equalities holds
\allowdisplaybreaks
\begin{align*}
\G_{\mathbf{d}}\left(g\circ f\right)_{s}
&=
\left(g_{\varphi(s)}\circ f_{\varphi(s)}\right)\bigr|^{(B''_{\Dom^{(\varphi,d)}})_{s}}_{(B_{\Dom^{(\varphi,d)}})_{s}}
\tag{1}
\\&=
g_{\varphi(s)}\bigr|^{(B''_{\Dom^{(\varphi,d)}})_{s}}_{(B'_{\Dom^{(\varphi,d)}})_{s}}
\circ
f_{\varphi(s)}\bigr|^{(B'_{\Dom^{(\varphi,d)}})_{s}}_{(B_{\Dom^{(\varphi,d)}})_{s}}
\tag{2}
\\&=
\G_{\mathbf{d}}\left(g\right)_{s}\circ\G_{\mathbf{d}}\left(f\right)_{s}.
\tag{3}
\end{align*}

The first equality unravels the definition of $\G_{\mathbf{d}}$ on the composite $g\circ f$; the second equality follows from the fact that the birestriction of the composite of two mappings is the composite of their respective birestrictions, taking into account that, by the discussion after Definition~\ref{DDomainSet}, $f_{\varphi(s)}{\,\!\upharpoonright}_{(B_{\Dom^{(\varphi,d)}})_{s}}$ corestricts to $(B'_{\Dom^{(\varphi,d)}})_{s}$; finally, the last equality unravels the definition of $\G_{\mathbf{d}}$ on $g$ and on $f$.

Thus, $\G_{\mathbf{d}}$ preserves composites.

This completes the proof.
\end{proof}

\begin{remark}\label{RDomainSetEmpty}
Let $\sig{\Sigma}=(S,\Sigma,\mathrm{def}^{(S,\Sigma)})$ and $\sig{\Lambda}=(T,\Lambda,\mathrm{def}^{(T,\Lambda)})$ be hep-signatures, $\mathbf{d}=(\varphi,d,\Dom^{(\varphi,d)})$ a generalized hep-derivor from $\sig{\Sigma}$ to $\sig{\Lambda}$ such that $\Dom^{(\varphi,d)}_{s}=\varnothing$, for every $s\in S$. Then, for every partial $\Lambda$-algebra $\mathbf{B}$, $B_{\Dom^{(\varphi,d)}}=B_{\varphi}$, and, for every homomorphism $f$, $f_{\Dom^{(\varphi,d)}}=f_{\varphi}$. Therefore
$$
\G_{\mathbf{d}}=\Delta_{\varphi}\circ\G_{\Lambda},
$$
where $\G_{\Lambda}$ is the forgetful functor from $\mathsf{PAlg}(\Lambda)$ to $\mathsf{Set}^{T}$ and $\Delta_{\varphi}$ is the functor from $\mathsf{Set}^{T}$ to $\mathsf{Set}^{S}$ that sends a $T$-sorted set $B$ to $B_{\varphi}=(B_{\varphi(s)})_{s\in S}$.
\end{remark}

We next show that the $S$-sorted set of cells is equipped with a structure of partial $\Sigma$-algebra, which we call the transformed algebra.

\begin{proposition}\label{PTransformedAlg}
Let $\sig{\Sigma}=(S,\Sigma,\mathrm{def}^{(S,\Sigma)})$ and $\sig{\Lambda}=(T,\Lambda,\mathrm{def}^{(T,\Lambda)})$ be hep-signatures, $\mathbf{d}=(\varphi,d,\Dom^{(\varphi,d)})$ a generalized hep-derivor from $\sig{\Sigma}$ to $\sig{\Lambda}$ and $\mathbf{B}$ a partial $\Lambda$-algebra. Then there exists a partial $\Sigma$-algebra $\mathbf{d}^{\ast}_{\gd}(\mathbf{B})$, the \emph{transformed algebra} of $\mathbf{B}$ along $\mathbf{d}$, whose underlying $S$-sorted set is $B_{\Dom^{(\varphi,d)}}$.
\end{proposition}
\begin{proof}
For every $(\mathbf{s},s)\in S^{\star}\times S$ and every $\sigma\in\Sigma_{\mathbf{s},s}$, we interpret the operation symbol $\sigma$ as the partial mapping
$$
\sigma^{\mathbf{d}^{\ast}_{\gd}(\mathbf{B})}\colon\left(B_{\Dom^{(\varphi,d)}}\right)_{\mathbf{s}}\dmor\left(B_{\Dom^{(\varphi,d)}}\right)_{s}
$$
defined as follows:
\begin{enumerate}
\item its domain is the set of the families $\mathbf{b}\in(B_{\Dom^{(\varphi,d)}})_{\mathbf{s}}$ such that
\begin{enumerate}
\item[(D1)] $\mathbf{d}^{\ast}_{\gd}(\mathbf{B})\models\mathrm{def}^{(S,\Sigma)}(\sigma)[\mathbf{b}]$, i.e., $\mathbf{b}$ satisfies the domain conditions of $\sigma$ in $\mathbf{d}^{\ast}_{\gd}(\mathbf{B})$;
\item[(D2)] $\mathbf{b}\in\Dom(d(\sigma)^{\mathbf{B}})$, i.e., the partial term operation $d(\sigma)^{\mathbf{B}}$ is defined at $\mathbf{b}$; and
\item[(D3)] $d(\sigma)^{\mathbf{B}}(\mathbf{b})\in(B_{\Dom^{(\varphi,d)}})_{s}$, i.e., its value is a cell of sort $s$; and
\end{enumerate}
\item it sends every family $\mathbf{b}$ in its domain to $d(\sigma)^{\mathbf{B}}(\mathbf{b})$.
\end{enumerate}
Condition (D1) refers to the operations of $\mathbf{d}^{\ast}_{\gd}(\mathbf{B})$ itself. We prove, by well-founded recursion along $\lhd$, that the above clauses determine a unique partial mapping $\sigma^{\mathbf{d}^{\ast}_{\gd}(\mathbf{B})}$ for every operation symbol $\sigma$ of $\Sigma$.

Let $(\mathbf{s},s)\in S^{\star}\times S$ and $\sigma\in\Sigma_{\mathbf{s},s}$, and let us assume that $\tau^{\mathbf{d}^{\ast}_{\gd}(\mathbf{B})}$ has been uniquely determined for every $\tau\lhd\sigma$. By Definition~\ref{DHepSignature}, the operation symbols that occur in $\mathrm{def}^{(S,\Sigma)}(\sigma)$ are exactly the $\tau\lhd\sigma$. Therefore, by Definitions~\ref{DTermOp} and~\ref{DEEq}, whether a family $\mathbf{b}\in(B_{\Dom^{(\varphi,d)}})_{\mathbf{s}}$ satisfies (D1) only depends on the partial mappings already determined. Conditions (D2) and (D3) only depend on $\mathbf{B}$ and on $\mathbf{d}$. Hence the domain of $\sigma^{\mathbf{d}^{\ast}_{\gd}(\mathbf{B})}$ is uniquely determined. Moreover, by (D3), for every family $\mathbf{b}$ in it, $d(\sigma)^{\mathbf{B}}(\mathbf{b})$ is a cell of sort $s$, so that $\sigma^{\mathbf{d}^{\ast}_{\gd}(\mathbf{B})}$ is a well-defined partial mapping from $(B_{\Dom^{(\varphi,d)}})_{\mathbf{s}}$ to $(B_{\Dom^{(\varphi,d)}})_{s}$.

All in all, since $\lhd$ is well-founded, $\sigma^{\mathbf{d}^{\ast}_{\gd}(\mathbf{B})}$ is uniquely determined for every operation symbol $\sigma$ of $\Sigma$, and $\mathbf{d}^{\ast}_{\gd}(\mathbf{B})=(B_{\Dom^{(\varphi,d)}},(\sigma^{\mathbf{d}^{\ast}_{\gd}(\mathbf{B})})_{\sigma})$ is a partial $\Sigma$-algebra.

This completes the proof.
\end{proof}

In what follows, the conditions (D1), (D2) and (D3) of Proposition~\ref{PTransformedAlg} are those of its proof. Thus, for every $(\mathbf{s},s)\in S^{\star}\times S$ and every $\sigma\in\Sigma_{\mathbf{s},s}$, a family $\mathbf{b}\in(B_{\Dom^{(\varphi,d)}})_{\mathbf{s}}$ is in $\Dom(\sigma^{\mathbf{d}^{\ast}_{\gd}(\mathbf{B})})$ if, and only if, it satisfies (D1), (D2) and (D3), and, in this case, $\sigma^{\mathbf{d}^{\ast}_{\gd}(\mathbf{B})}(\mathbf{b})=d(\sigma)^{\mathbf{B}}(\mathbf{b})$. In particular, by (D1), for every $\sigma\in\Sigma_{\mathbf{s},s}$,
\begin{equation*}
\Dom\left(\sigma^{\mathbf{d}^{\ast}_{\gd}(\mathbf{B})}\right)
\subseteq
\left(B_{\Dom^{(\varphi,d)}}\right)_{\mathrm{def}^{(S,\Sigma)}(\sigma)},
\end{equation*}
where $(B_{\Dom^{(\varphi,d)}})_{\mathrm{def}^{(S,\Sigma)}(\sigma)}$ is the set of the families of cells that satisfy $\mathrm{def}^{(S,\Sigma)}(\sigma)$ in $\mathbf{d}^{\ast}_{\gd}(\mathbf{B})$, see Definition~\ref{DEpAlgebra}, i.e., $\mathbf{d}^{\ast}_{\gd}(\mathbf{B})$ is a weak $\sig{\Sigma}$-algebra, see Definition~\ref{DWeakEpAlgebra}.

Let us note that, in $\mathbf{d}^{\ast}_{\gd}(\mathbf{B})$, the domain of $\sigma$ is not part of the data, and it may be smaller than that of $d(\sigma)^{\mathbf{B}}$. For instance, for the discretization functors of Section~\ref{STruncation}, the term assigned to $\circ^{k}_{s}$, for $k$ greater than or equal to the dimension of the algebra, is the variable $v^{s}_{0}$, which is defined everywhere, whereas, in the transformed algebras, the domain condition of $\circ^{k}_{s}$ is satisfied only by the pairs of equal cells. In~\cite{Rei87}, the morphisms of ep-signatures are required, moreover, to translate the domain conditions of $\sigma$ into a set of equations equivalent to the definedness of $d(\sigma)$.

In the following proposition we characterize when the transformed algebra is a $\sig{\Sigma}$-algebra.

\begin{proposition}\label{PTransformedEpAlg}
Let $\sig{\Sigma}=(S,\Sigma,\mathrm{def}^{(S,\Sigma)})$ and $\sig{\Lambda}=(T,\Lambda,\mathrm{def}^{(T,\Lambda)})$ be hep-signatures, $\mathbf{d}=(\varphi,d,\Dom^{(\varphi,d)})$ a generalized hep-derivor from $\sig{\Sigma}$ to $\sig{\Lambda}$ and $\mathbf{B}$ a partial $\Lambda$-algebra. Then the following conditions are equivalent:
\begin{enumerate}
\item $\mathbf{d}^{\ast}_{\gd}(\mathbf{B})$ is a $\sig{\Sigma}$-algebra; and
\item for every $(\mathbf{s},s)\in S^{\star}\times S$, every $\sigma\in\Sigma_{\mathbf{s},s}$ and every $\mathbf{b}\in(B_{\Dom^{(\varphi,d)}})_{\mathbf{s}}$ such that $\mathbf{d}^{\ast}_{\gd}(\mathbf{B})\models\mathrm{def}^{(S,\Sigma)}(\sigma)[\mathbf{b}]$, we have that $\mathbf{b}\in\Dom(d(\sigma)^{\mathbf{B}})$ and $d(\sigma)^{\mathbf{B}}(\mathbf{b})\in(B_{\Dom^{(\varphi,d)}})_{s}$.
\end{enumerate}
\end{proposition}
\begin{proof}
Let $(\mathbf{s},s)\in S^{\star}\times S$ and $\sigma\in\Sigma_{\mathbf{s},s}$. We let $(B_{\Dom^{(\varphi,d)}})_{\mathrm{def}^{(S,\Sigma)}(\sigma)}$ stand for the set of the families of cells that satisfy $\mathrm{def}^{(S,\Sigma)}(\sigma)$ in $\mathbf{d}^{\ast}_{\gd}(\mathbf{B})$, see Definition~\ref{DEpAlgebra}. By condition (D1) of Proposition~\ref{PTransformedAlg}, the following inclusion holds
\begin{equation*}
\Dom\left(\sigma^{\mathbf{d}^{\ast}_{\gd}(\mathbf{B})}\right)
\subseteq
\left(B_{\Dom^{(\varphi,d)}}\right)_{\mathrm{def}^{(S,\Sigma)}(\sigma)}.
\tag{I}
\end{equation*}
By Definition~\ref{DEpAlgebra}, $\mathbf{d}^{\ast}_{\gd}(\mathbf{B})$ is a $\sig{\Sigma}$-algebra if, and only if, the inclusion~(I) is an equality for every operation symbol $\sigma$ of $\Sigma$.

\textsf{(1) implies (2)} Let us assume that the inclusion~(I) is an equality and let $\mathbf{b}$ be a family in $(B_{\Dom^{(\varphi,d)}})_{\mathrm{def}^{(S,\Sigma)}(\sigma)}$. Then $\mathbf{b}\in\Dom(\sigma^{\mathbf{d}^{\ast}_{\gd}(\mathbf{B})})$, so that, by conditions (D2) and (D3) of Proposition~\ref{PTransformedAlg}, $\mathbf{b}\in\Dom(d(\sigma)^{\mathbf{B}})$ and $d(\sigma)^{\mathbf{B}}(\mathbf{b})\in(B_{\Dom^{(\varphi,d)}})_{s}$.

\textsf{(2) implies (1)} Let us assume that condition (2) holds for $\sigma$ and let $\mathbf{b}$ be a family in $(B_{\Dom^{(\varphi,d)}})_{\mathrm{def}^{(S,\Sigma)}(\sigma)}$. Then $\mathbf{b}$ satisfies condition (D1) and, by condition (2), conditions (D2) and (D3) of Proposition~\ref{PTransformedAlg}. Hence $\mathbf{b}\in\Dom(\sigma^{\mathbf{d}^{\ast}_{\gd}(\mathbf{B})})$, and the inclusion~(I) is an equality.

This completes the proof.
\end{proof}

In the following proposition we prove that the birestrictions of homomorphisms introduced after Definition~\ref{DDomainSet} are homomorphisms between the transformed algebras.

\begin{proposition}\label{PGenHepDerivorHom}
Let $\sig{\Sigma}=(S,\Sigma,\mathrm{def}^{(S,\Sigma)})$ and $\sig{\Lambda}=(T,\Lambda,\mathrm{def}^{(T,\Lambda)})$ be hep-signatures, $\mathbf{d}=(\varphi,d,\Dom^{(\varphi,d)})$ a generalized hep-derivor from $\sig{\Sigma}$ to $\sig{\Lambda}$ and $f\colon\mathbf{B}\mor\mathbf{B}'$ a homomorphism of partial $\Lambda$-algebras. Then $f_{\Dom^{(\varphi,d)}}$ is a homomorphism of partial $\Sigma$-algebras from $\mathbf{d}^{\ast}_{\gd}(\mathbf{B})$ to $\mathbf{d}^{\ast}_{\gd}(\mathbf{B}')$.
\end{proposition}
\begin{proof}
Let $h$ stand for $f_{\Dom^{(\varphi,d)}}$, which, by the discussion after Definition~\ref{DDomainSet}, is an $S$-sorted mapping from $B_{\Dom^{(\varphi,d)}}$ to $B'_{\Dom^{(\varphi,d)}}$. We prove, by well-founded induction along $\lhd$, that, for every $(\mathbf{s},s)\in S^{\star}\times S$, every $\sigma\in\Sigma_{\mathbf{s},s}$ and every $\mathbf{b}\in\Dom(\sigma^{\mathbf{d}^{\ast}_{\gd}(\mathbf{B})})$, the following conditions hold:
\begin{enumerate}
\item[(H$_{\sigma}$1)] $h\circ\mathbf{b}\in\Dom(\sigma^{\mathbf{d}^{\ast}_{\gd}(\mathbf{B}')})$; and
\item[(H$_{\sigma}$2)] $h_{s}(\sigma^{\mathbf{d}^{\ast}_{\gd}(\mathbf{B})}(\mathbf{b}))=\sigma^{\mathbf{d}^{\ast}_{\gd}(\mathbf{B}')}(h\circ\mathbf{b})$.
\end{enumerate}

Let $\sigma\in\Sigma_{\mathbf{s},s}$ and assume that (H$_{\tau}$1) and (H$_{\tau}$2) hold for every $\tau\lhd\sigma$. Let $\mathbf{b}$ be a family in $\Dom(\sigma^{\mathbf{d}^{\ast}_{\gd}(\mathbf{B})})$.

\textsf{(H$_{\sigma}$1)} We check the conditions (D1), (D2) and (D3) that define the domain of $\sigma^{\mathbf{d}^{\ast}_{\gd}(\mathbf{B}')}$ in Proposition~\ref{PTransformedAlg}.

\textsf{(D1)} By the induction hypothesis, $h$ is a homomorphism between the reducts of $\mathbf{d}^{\ast}_{\gd}(\mathbf{B})$ and $\mathbf{d}^{\ast}_{\gd}(\mathbf{B}')$ to the operation symbols that occur in $\mathrm{def}^{(S,\Sigma)}(\sigma)$. Since $\mathbf{d}^{\ast}_{\gd}(\mathbf{B})\models\mathrm{def}^{(S,\Sigma)}(\sigma)[\mathbf{b}]$, by Remark~\ref{RHomTerm}, it follows that
$$
\mathbf{d}^{\ast}_{\gd}(\mathbf{B}')\models\mathrm{def}^{(S,\Sigma)}(\sigma)[h\circ\mathbf{b}].
$$

\textsf{(D2)} Since $\mathbf{b}\in\Dom(d(\sigma)^{\mathbf{B}})$ and $f$ is a homomorphism, by Remark~\ref{RHomTerm}, $f_{\varphi}\circ\mathbf{b}=h\circ\mathbf{b}$ is in $\Dom(d(\sigma)^{\mathbf{B}'})$ and
$$
d(\sigma)^{\mathbf{B}'}(h\circ\mathbf{b})=f_{\varphi(s)}\left(d(\sigma)^{\mathbf{B}}(\mathbf{b})\right).
$$

\textsf{(D3)} Since $d(\sigma)^{\mathbf{B}}(\mathbf{b})\in(B_{\Dom^{(\varphi,d)}})_{s}$, by the discussion after Definition~\ref{DDomainSet}, $f_{\varphi(s)}(d(\sigma)^{\mathbf{B}}(\mathbf{b}))\in(B'_{\Dom^{(\varphi,d)}})_{s}$, i.e., by (D2), $d(\sigma)^{\mathbf{B}'}(h\circ\mathbf{b})$ is a cell of sort $s$ of $\mathbf{B}'$.

Therefore, by Proposition~\ref{PTransformedAlg}, $h\circ\mathbf{b}\in\Dom(\sigma^{\mathbf{d}^{\ast}_{\gd}(\mathbf{B}')})$.

\textsf{(H$_{\sigma}$2)} The following chain of equalities holds
\allowdisplaybreaks
\begin{align*}
h_{s}\left(\sigma^{\mathbf{d}^{\ast}_{\gd}(\mathbf{B})}\left(\mathbf{b}\right)\right)
&=
f_{\varphi(s)}\left(d(\sigma)^{\mathbf{B}}\left(\mathbf{b}\right)\right)
\tag{1}
\\&=
d(\sigma)^{\mathbf{B}'}\left(h\circ\mathbf{b}\right)
\tag{2}
\\&=
\sigma^{\mathbf{d}^{\ast}_{\gd}(\mathbf{B}')}\left(h\circ\mathbf{b}\right).
\tag{3}
\end{align*}

The first equality unravels the definitions of $h$ and of $\sigma^{\mathbf{d}^{\ast}_{\gd}(\mathbf{B})}$, see Proposition~\ref{PTransformedAlg}; the second equality follows from Remark~\ref{RHomTerm}, as shown in the proof of (H$_{\sigma}$1); finally, the last equality follows from Proposition~\ref{PTransformedAlg}, taking into account that, by (H$_{\sigma}$1), $h\circ\mathbf{b}\in\Dom(\sigma^{\mathbf{d}^{\ast}_{\gd}(\mathbf{B}')})$.

All in all, since $\lhd$ is well-founded, (H$_{\sigma}$1) and (H$_{\sigma}$2) hold for every operation symbol $\sigma$ of $\Sigma$, i.e., $h$ is a homomorphism from $\mathbf{d}^{\ast}_{\gd}(\mathbf{B})$ to $\mathbf{d}^{\ast}_{\gd}(\mathbf{B}')$.

This completes the proof.
\end{proof}

We are now in position to prove that every generalized hep-derivor induces a functor between the categories of partial algebras. Let us note that, in general, the transformed algebra of a $\sig{\Lambda}$-algebra need not be a $\sig{\Sigma}$-algebra, see Proposition~\ref{PTransformedEpAlg} and Remark~\ref{RGenHepDerivorCases}. Hence a generalized hep-derivor does not induce, in general, a functor from $\mathsf{EPAlg}(\sig{\Lambda})$ to $\mathsf{EPAlg}(\sig{\Sigma})$, and we define the induced functor on all the partial $\Lambda$-algebras. In Subsection~\ref{SEpRestriction} we consider its restrictions to the categories canonically associated to ep-varieties.

\begin{proposition}\label{PGenHepDerivorFunctor}
Let $\sig{\Sigma}=(S,\Sigma,\mathrm{def}^{(S,\Sigma)})$ and $\sig{\Lambda}=(T,\Lambda,\mathrm{def}^{(T,\Lambda)})$ be hep-signatures, $\mathbf{d}=(\varphi,d,\Dom^{(\varphi,d)})$ a generalized hep-derivor from $\sig{\Sigma}$ to $\sig{\Lambda}$. Then $\mathbf{d}^{\ast}_{\gd}$, the assignment from $\mathsf{PAlg}(\Lambda)$ to $\mathsf{PAlg}(\Sigma)$ defined as follows:
\begin{enumerate}
\item for every partial $\Lambda$-algebra $\mathbf{B}$, $\mathbf{d}^{\ast}_{\gd}(\mathbf{B})$ is the partial $\Sigma$-algebra introduced in Proposition~\ref{PTransformedAlg}; and
\item for every homomorphism $f\colon\mathbf{B}\mor\mathbf{B}'$ of partial $\Lambda$-algebras, $\mathbf{d}^{\ast}_{\gd}(f)$ is the homomorphism $f_{\Dom^{(\varphi,d)}}$ from $\mathbf{d}^{\ast}_{\gd}(\mathbf{B})$ to $\mathbf{d}^{\ast}_{\gd}(\mathbf{B}')$, see Proposition~\ref{PGenHepDerivorHom};
\end{enumerate}
is a functor.
\end{proposition}
\begin{proof}
By Propositions~\ref{PTransformedAlg} and~\ref{PGenHepDerivorHom}, it suffices to prove that $\mathbf{d}^{\ast}_{\gd}$ preserves identities and composites. Since $\mathbf{d}^{\ast}_{\gd}$ and $\G_{\mathbf{d}}$ act in the same way on homomorphisms, this follows from Proposition~\ref{PDomainSetFunctor}.

This completes the proof.
\end{proof}

\begin{corollary}\label{CGenHepDerivorDiagram}
Let $\sig{\Sigma}=(S,\Sigma,\mathrm{def}^{(S,\Sigma)})$ and $\sig{\Lambda}=(T,\Lambda,\mathrm{def}^{(T,\Lambda)})$ be hep-signatures, $\mathbf{d}=(\varphi,d,\Dom^{(\varphi,d)})$ a generalized hep-derivor from $\sig{\Sigma}$ to $\sig{\Lambda}$. Then the following diagram commutes
$$
\xymatrix@C=50pt@R=35pt{
\mathsf{PAlg}(\Lambda)\ar[r]^-{\mathbf{d}^{\ast}_{\gd}}\ar[rd]_-{\G_{\mathbf{d}}} &
\mathsf{PAlg}(\Sigma)\ar[d]^-{\G_{\Sigma}}\\
& \mathsf{Set}^{S}
}
$$
where $\G_{\Sigma}$ is the forgetful functor.
\end{corollary}
\begin{proof}
It follows from Propositions~\ref{PDomainSetFunctor}, \ref{PTransformedAlg} and~\ref{PGenHepDerivorFunctor}, since the underlying $S$-sorted set of $\mathbf{d}^{\ast}_{\gd}(\mathbf{B})$ is $B_{\Dom^{(\varphi,d)}}$ and the underlying $S$-sorted mapping of $\mathbf{d}^{\ast}_{\gd}(f)$ is $f_{\Dom^{(\varphi,d)}}$.
\end{proof}

\subsection{Restriction to ep-varieties}\label{SEpRestriction}

Let $\sig{\Sigma}=(S,\Sigma,\mathrm{def}^{(S,\Sigma)})$ and $\sig{\Lambda}=(T,\Lambda,\mathrm{def}^{(T,\Lambda)})$ be hep-signatures, $\mathbf{d}=(\varphi,d,\Dom^{(\varphi,d)})$ a generalized hep-derivor from $\sig{\Sigma}$ to $\sig{\Lambda}$, and $(\sig{\Sigma},\mathcal{E})$ and $(\sig{\Lambda},\mathcal{F})$ hep-specifications. By Definitions~\ref{DCorestriction} and~\ref{DEpVariety}, the restriction of $\mathbf{d}^{\ast}_{\gd}$ to $\mathsf{EPAlg}(\sig{\Lambda},\mathcal{F})$ corestricts to $\mathsf{EPAlg}(\sig{\Sigma},\mathcal{E})$ if, and only if, for every object $\mathbf{B}$ in $\mathsf{EPAlg}(\sig{\Lambda},\mathcal{F})$, $\mathbf{d}^{\ast}_{\gd}(\mathbf{B})$ is a $\sig{\Sigma}$-algebra in which every $\QE$-equation of $\mathcal{E}$ holds. Accordingly, for each of the generalized hep-derivors of Sections~\ref{SGraduation} and~\ref{STruncation}, we first prove, by means of Proposition~\ref{PTransformedEpAlg}, that the transformed algebras are ep-algebras, and then that the axioms of the corresponding ep-variety hold in them. In this case, the corestriction of the restriction of $\mathbf{d}^{\ast}_{\gd}$ is denoted by
$$
\mathbf{d}^{\ast}_{\gd}\colon\mathsf{EPAlg}(\sig{\Lambda},\mathcal{F})\mor\mathsf{EPAlg}(\sig{\Sigma},\mathcal{E}).
$$

\begin{remark}\label{RGenHepDerivorCases}
\begin{enumerate}
\item If $\sig{\Sigma}=(S,\Sigma)$ and $\sig{\Lambda}=(T,\Lambda)$ are many-sorted signatures, regarded as hep-signatures with empty sets of domain conditions, then, for every derivor $(\varphi,d)$, the restriction of $(\varphi,d,\varnothing)^{\ast}_{\gd}$ to $\mathsf{EPAlg}(T,\Lambda,\varnothing)=\mathsf{Alg}(\Lambda)$ corestricts to $\mathsf{EPAlg}(S,\Sigma,\varnothing)=\mathsf{Alg}(\Sigma)$, since the partial term operations of a total algebra are total. Moreover, this corestriction is the usual functor that sends a $\Lambda$-algebra $\mathbf{B}$ to the $\Sigma$-algebra on $B_{\varphi}$ in which every $\sigma$ is interpreted as the term operation $d(\sigma)^{\mathbf{B}}$, see~\cite{CVST10b}.
\item Let $\sig{\Sigma}$ and $\sig{\Lambda}$ be hep-signatures and $\mathbf{d}$ a generalized hep-derivor from $\sig{\Sigma}$ to $\sig{\Lambda}$. The functor $\mathbf{d}^{\ast}_{\gd}$ is defined on all the partial $\Lambda$-algebras, but its restriction to $\mathsf{EPAlg}(\sig{\Lambda},\mathcal{F})$ may only corestrict to $\mathsf{EPAlg}(\sig{\Sigma},\mathcal{E})$ for some hep-specifications $(\sig{\Lambda},\mathcal{F})$ and $(\sig{\Sigma},\mathcal{E})$. For instance, the domains of the sorts of the graduation derivor of Section~\ref{SGraduation} are only preserved by the term operations in the objects of the category canonically associated to the ep-variety of the axioms of $n$-categories: in a partial algebra in which $\src^{j}_{s}$ is not idempotent, the $j$-source of a cell need not be a $j$-cell.
\end{enumerate}
\end{remark}

\subsection{The category \texorpdfstring{$\mathsf{hepSig}_{\gd}$}{hepSiggd}}\label{SHepSigCat}

In this subsection we prove that the hep-signatures and the generalized hep-derivors form a category. We use the notation $e^{\sharp}[\mathcal{H}]$ of Remark~\ref{RDerivorTerms} for the translation of a set of $\mathrm{E}$-equations $\mathcal{H}$ along a derivor $(\psi,e)$.

\begin{definition}\label{DHepComp}
Let $\sig{\Sigma}=(S,\Sigma,\mathrm{def}^{(S,\Sigma)})$, $\sig{\Lambda}=(T,\Lambda,\mathrm{def}^{(T,\Lambda)})$ and $\sig{\Omega}=(U,\Omega,\mathrm{def}^{(U,\Omega)})$ be hep-signatures, $\mathbf{d}=(\varphi,d,\Dom^{(\varphi,d)})$ a generalized hep-derivor from $\sig{\Sigma}$ to $\sig{\Lambda}$ and $\mathbf{e}=(\psi,e,\Dom^{(\psi,e)})$ a generalized hep-derivor from $\sig{\Lambda}$ to $\sig{\Omega}$. The \emph{composition} of $\mathbf{d}$ and $\mathbf{e}$ is the generalized hep-derivor
$$
\mathbf{e}\circ\mathbf{d}=\left(\psi\circ\varphi,e^{\sharp}\circ d,\Dom^{\mathbf{e}\circ\mathbf{d}}\right)
$$
from $\sig{\Sigma}$ to $\sig{\Omega}$, where $(\psi\circ\varphi,e^{\sharp}\circ d)$ is the composition of the derivors $(\varphi,d)$ and $(\psi,e)$, see~(Der3), and, for every $s\in S$,
$$
\Dom^{\mathbf{e}\circ\mathbf{d}}_{s}=\Dom^{(\psi,e)}_{\varphi(s)}\cup e^{\sharp}\left[\Dom^{(\varphi,d)}_{s}\right].
$$
The \emph{identity} at $\sig{\Sigma}$ is the generalized hep-derivor $\id_{\sig{\Sigma}}=(\id_{S},\eta^{\Sigma},\varnothing)$ from $\sig{\Sigma}$ to $\sig{\Sigma}$, where $(\id_{S},\eta^{\Sigma})$ is the identity derivor at $(S,\Sigma)$, see~(Der3), and every domain is empty.
\end{definition}

Let us note that $\Dom^{(\psi,e)}_{\varphi(s)}$ and $e^{\sharp}[\Dom^{(\varphi,d)}_{s}]$ are sets of $\mathrm{E}$-equations with variables in $\vs{\psi(\varphi(s))}$, since $\psi^{\star}((\varphi(s)))=(\psi(\varphi(s)))$, so that the above definition is correct.

\begin{proposition}\label{PHepSigCat}
The hep-signatures, together with the generalized hep-derivors, the composition and the identities of Definition~\ref{DHepComp}, form a category, which we denote by $\mathsf{hepSig}_{\gd}$.
\end{proposition}
\begin{proof}
Let $\sig{\Sigma}=(S,\Sigma,\mathrm{def}^{(S,\Sigma)})$, $\sig{\Lambda}=(T,\Lambda,\mathrm{def}^{(T,\Lambda)})$, $\sig{\Omega}=(U,\Omega,\mathrm{def}^{(U,\Omega)})$ and $\sig{\Xi}=(W,\Xi,\mathrm{def}^{(W,\Xi)})$ be hep-signatures, and $\mathbf{d}=(\varphi,d,\Dom^{(\varphi,d)})$, $\mathbf{e}=(\psi,e,\Dom^{(\psi,e)})$ and $\mathbf{g}=(\chi,g,\Dom^{(\chi,g)})$ generalized hep-derivors from $\sig{\Sigma}$ to $\sig{\Lambda}$, from $\sig{\Lambda}$ to $\sig{\Omega}$ and from $\sig{\Omega}$ to $\sig{\Xi}$, respectively. By~(Der4), the underlying derivors of $\mathbf{g}\circ(\mathbf{e}\circ\mathbf{d})$ and of $(\mathbf{g}\circ\mathbf{e})\circ\mathbf{d}$ are equal, and the underlying derivors of $\id_{\sig{\Lambda}}\circ\mathbf{d}$ and of $\mathbf{d}\circ\id_{\sig{\Sigma}}$ are $(\varphi,d)$. It remains to compare the domains. Let $s\in S$.

\textsf{(Associativity)} The following chain of equalities holds
\allowdisplaybreaks
\begin{align*}
\Dom^{\mathbf{g}\circ(\mathbf{e}\circ\mathbf{d})}_{s}
&=
\Dom^{(\chi,g)}_{\psi(\varphi(s))}\cup g^{\sharp}\left[\Dom^{(\psi,e)}_{\varphi(s)}\cup e^{\sharp}\left[\Dom^{(\varphi,d)}_{s}\right]\right]
\tag{1}
\\&=
\Dom^{(\chi,g)}_{\psi(\varphi(s))}\cup g^{\sharp}\left[\Dom^{(\psi,e)}_{\varphi(s)}\right]\cup g^{\sharp}\left[e^{\sharp}\left[\Dom^{(\varphi,d)}_{s}\right]\right]
\tag{2}
\\&=
\Dom^{\mathbf{g}\circ\mathbf{e}}_{\varphi(s)}\cup\left(g^{\sharp}\circ e\right)^{\sharp}\left[\Dom^{(\varphi,d)}_{s}\right]
\tag{3}
\\&=
\Dom^{(\mathbf{g}\circ\mathbf{e})\circ\mathbf{d}}_{s}.
\tag{4}
\end{align*}

The first equality unravels Definition~\ref{DHepComp}; the second equality follows from the fact that $g^{\sharp}[\,\cdot\,]$ preserves unions; the third equality follows from Definition~\ref{DHepComp} and the equality $(g^{\sharp}\circ e)^{\sharp}=g^{\sharp}\circ e^{\sharp}$, see~(Der4); finally, the last equality unravels Definition~\ref{DHepComp}.

\textsf{(Identities)} By Definition~\ref{DHepComp} and~(Der3), the following equalities hold
\begin{align*}
\Dom^{\id_{\sig{\Lambda}}\circ\mathbf{d}}_{s}
&=\varnothing\cup(\eta^{\Lambda})^{\sharp}\left[\Dom^{(\varphi,d)}_{s}\right]=\Dom^{(\varphi,d)}_{s},
\\
\Dom^{\mathbf{d}\circ\id_{\sig{\Sigma}}}_{s}
&=\Dom^{(\varphi,d)}_{s}\cup d^{\sharp}[\varnothing]=\Dom^{(\varphi,d)}_{s},
\end{align*}
since $(\eta^{\Lambda})^{\sharp}$ is the identity.

This completes the proof.
\end{proof}

\subsection{Composition and induced functors}\label{SHepCompFunctors}

In this subsection we compare the functor induced by the composition of two generalized hep-derivors with the composite of the induced functors. The composition of Definition~\ref{DHepComp} translates the domains of $\mathbf{d}$ by means of the terms of $\mathbf{e}$, but these terms may be defined in a partial $\Omega$-algebra $\mathbf{C}$ on more families than the corresponding operations of $\mathbf{e}^{\ast}_{\gd}(\mathbf{C})$. Therefore both functors need not coincide; we prove that the first one always contains the second one, and that they coincide under a condition that we will check in Sections~\ref{SGraduation} and~\ref{STruncation}. We begin by relating the partial term operations of $\mathbf{e}^{\ast}_{\gd}(\mathbf{C})$ with those of $\mathbf{C}$.

\begin{lemma}\label{LTermComp}
Let $\sig{\Lambda}=(T,\Lambda,\mathrm{def}^{(T,\Lambda)})$ and $\sig{\Omega}=(U,\Omega,\mathrm{def}^{(U,\Omega)})$ be hep-signatures, $\mathbf{e}=(\psi,e,\Dom^{(\psi,e)})$ a generalized hep-derivor from $\sig{\Lambda}$ to $\sig{\Omega}$ and $\mathbf{C}$ a partial $\Omega$-algebra. Let $\mathbf{t}\in T^{\star}$, $t\in T$, $P\in\mathrm{T}_{\Lambda}(\vs{\mathbf{t}})_{t}$ and $\mathbf{c}\in(C_{\Dom^{(\psi,e)}})_{\mathbf{t}}$. If $\mathbf{c}\in\Dom(P^{\mathbf{e}^{\ast}_{\gd}(\mathbf{C})})$, then $\mathbf{c}\in\Dom(e^{\sharp}(P)^{\mathbf{C}})$ and
$$
e^{\sharp}(P)^{\mathbf{C}}(\mathbf{c})=P^{\mathbf{e}^{\ast}_{\gd}(\mathbf{C})}(\mathbf{c}).
$$
Consequently, for every $\mathcal{H}\subseteq\mathrm{Eeqt}(\Lambda)_{\vs{\mathbf{t}}}$, if $\mathbf{e}^{\ast}_{\gd}(\mathbf{C})\models\mathcal{H}[\mathbf{c}]$, then $\mathbf{C}\models e^{\sharp}[\mathcal{H}][\mathbf{c}]$.
\end{lemma}
\begin{proof}
We proceed by induction on the term $P$.

\textsf{(Variables)} If $P=v^{\mathbf{t}}_{j}$, for some $j\in\bb{\mathbf{t}}$, then, by~(Der1), $e^{\sharp}(P)=v^{\psi^{\star}(\mathbf{t})}_{j}$, so that $e^{\sharp}(P)^{\mathbf{C}}(\mathbf{c})=c_{j}=P^{\mathbf{e}^{\ast}_{\gd}(\mathbf{C})}(\mathbf{c})$.

\textsf{(Operation symbols)} Let us assume that $P=\lambda((P_{j})_{j\in\bb{\mathbf{r}}})$, with $\lambda\in\Lambda_{\mathbf{r},t}$, and that the statement holds for every $P_{j}$. Let $\mathbf{c}\in\Dom(P^{\mathbf{e}^{\ast}_{\gd}(\mathbf{C})})$ and let $\mathbf{y}$ stand for the family $(P_{j}^{\mathbf{e}^{\ast}_{\gd}(\mathbf{C})}(\mathbf{c}))_{j\in\bb{\mathbf{r}}}$, which, by Definition~\ref{DTermOp}, is in $\Dom(\lambda^{\mathbf{e}^{\ast}_{\gd}(\mathbf{C})})$. By Proposition~\ref{PTransformedAlg}, $\mathbf{y}\in\Dom(e(\lambda)^{\mathbf{C}})$ and $\lambda^{\mathbf{e}^{\ast}_{\gd}(\mathbf{C})}(\mathbf{y})=e(\lambda)^{\mathbf{C}}(\mathbf{y})$. By~(Der2), $e^{\sharp}(P)=(e^{\sharp}(P_{j}))^{\sharp}_{j\in\bb{\mathbf{r}}}(e(\lambda))$. Therefore, by the induction hypothesis and Remark~\ref{RSubst}, $\mathbf{c}\in\Dom(e^{\sharp}(P)^{\mathbf{C}})$, and the following chain of equalities holds
\allowdisplaybreaks
\begin{align*}
e^{\sharp}(P)^{\mathbf{C}}\left(\mathbf{c}\right)
&=
\left(\left(e^{\sharp}(P_{j})\right)^{\sharp}_{j\in\bb{\mathbf{r}}}\left(e(\lambda)\right)\right)^{\mathbf{C}}\left(\mathbf{c}\right)
\tag{1}
\\&=
e(\lambda)^{\mathbf{C}}\left(\left(e^{\sharp}(P_{j})^{\mathbf{C}}(\mathbf{c})\right)_{j\in\bb{\mathbf{r}}}\right)
\tag{2}
\\&=
e(\lambda)^{\mathbf{C}}\left(\mathbf{y}\right)
\tag{3}
\\&=
\lambda^{\mathbf{e}^{\ast}_{\gd}(\mathbf{C})}\left(\mathbf{y}\right)
\tag{4}
\\&=
P^{\mathbf{e}^{\ast}_{\gd}(\mathbf{C})}\left(\mathbf{c}\right).
\tag{5}
\end{align*}

The first equality follows from~(Der2); the second equality follows from Remark~\ref{RSubst}; the third equality follows from the induction hypothesis; the fourth equality follows from Proposition~\ref{PTransformedAlg}; finally, the last equality unravels the definition of the partial term operation $P^{\mathbf{e}^{\ast}_{\gd}(\mathbf{C})}$, see Definition~\ref{DTermOp}.

\textsf{(Equations)} Let $(P,Q)\in\mathcal{H}$ be such that $\mathbf{e}^{\ast}_{\gd}(\mathbf{C})\models(P,Q)[\mathbf{c}]$. Then $P^{\mathbf{e}^{\ast}_{\gd}(\mathbf{C})}(\mathbf{c})$ and $Q^{\mathbf{e}^{\ast}_{\gd}(\mathbf{C})}(\mathbf{c})$ are defined and equal, so that, by the above, $e^{\sharp}(P)^{\mathbf{C}}(\mathbf{c})$ and $e^{\sharp}(Q)^{\mathbf{C}}(\mathbf{c})$ are defined and equal, i.e., $\mathbf{C}\models(e^{\sharp}(P),e^{\sharp}(Q))[\mathbf{c}]$.

This completes the proof.
\end{proof}

In the following lemma we prove that the transformed algebra only depends on the domains through the cells they determine.

\begin{lemma}\label{LSameCells}
Let $\sig{\Sigma}=(S,\Sigma,\mathrm{def}^{(S,\Sigma)})$ and $\sig{\Lambda}=(T,\Lambda,\mathrm{def}^{(T,\Lambda)})$ be hep-signatures, $\mathbf{d}=(\varphi,d,\Dom)$ and $\mathbf{d}'=(\varphi,d,\Dom')$ generalized hep-derivors from $\sig{\Sigma}$ to $\sig{\Lambda}$ with the same underlying derivor, and $\mathbf{B}$ a partial $\Lambda$-algebra such that $B_{\Dom}=B_{\Dom'}$. Then $\mathbf{d}^{\ast}_{\gd}(\mathbf{B})=\mathbf{d}'^{\ast}_{\gd}(\mathbf{B})$. Moreover, for every homomorphism $f\colon\mathbf{B}\mor\mathbf{B}'$ of partial $\Lambda$-algebras such that $B'_{\Dom}=B'_{\Dom'}$, $\mathbf{d}^{\ast}_{\gd}(f)=\mathbf{d}'^{\ast}_{\gd}(f)$.
\end{lemma}
\begin{proof}
Both partial $\Sigma$-algebras have the same underlying $S$-sorted set. We prove, by well-founded induction along $\lhd$, that every $\sigma\in\Sigma_{\mathbf{s},s}$ has the same interpretation in both. Let us assume that every $\tau\lhd\sigma$ does. Then, by Definition~\ref{DHepSignature}, a family of cells satisfies condition (D1) of Proposition~\ref{PTransformedAlg} for $\mathbf{d}$ if, and only if, it satisfies it for $\mathbf{d}'$. Conditions (D2) and (D3) only depend on $d(\sigma)$, on $\mathbf{B}$ and on the common cells. Hence $\sigma$ has the same domain in both algebras and, on it, both interpretations send $\mathbf{b}$ to $d(\sigma)^{\mathbf{B}}(\mathbf{b})$. Finally, by Proposition~\ref{PGenHepDerivorFunctor}, $\mathbf{d}^{\ast}_{\gd}(f)$ and $\mathbf{d}'^{\ast}_{\gd}(f)$ are the birestrictions of $f_{\varphi}$ to the same cells.

This completes the proof.
\end{proof}

In the following proposition we prove that the identity generalized hep-derivor induces, on the weak $\sig{\Sigma}$-algebras, and, in particular, on the $\sig{\Sigma}$-algebras, the identity functor.

\begin{proposition}\label{PHepId}
Let $\sig{\Sigma}=(S,\Sigma,\mathrm{def}^{(S,\Sigma)})$ be a hep-signature. Then, for every weak $\sig{\Sigma}$-algebra $\mathbf{A}$ and every homomorphism $f$ between weak $\sig{\Sigma}$-algebras, see Definition~\ref{DWeakEpAlgebra}, $(\id_{\sig{\Sigma}})^{\ast}_{\gd}(\mathbf{A})=\mathbf{A}$ and $(\id_{\sig{\Sigma}})^{\ast}_{\gd}(f)=f$. Thus the restriction of $(\id_{\sig{\Sigma}})^{\ast}_{\gd}$ to $\mathsf{wEPAlg}(\sig{\Sigma})$ is $\Id_{\mathsf{wEPAlg}(\sig{\Sigma})}$ and its restriction to $\mathsf{EPAlg}(\sig{\Sigma})$ is $\Id_{\mathsf{EPAlg}(\sig{\Sigma})}$.
\end{proposition}
\begin{proof}
Let $\mathbf{A}$ be a weak $\sig{\Sigma}$-algebra. Since the domains of $\id_{\sig{\Sigma}}$ are empty, by Remark~\ref{RDomainSetEmpty}, the underlying $S$-sorted set of $(\id_{\sig{\Sigma}})^{\ast}_{\gd}(\mathbf{A})$ is $A$, and $(\id_{\sig{\Sigma}})^{\ast}_{\gd}(f)=f$. We prove, by well-founded induction along $\lhd$, that every $\sigma\in\Sigma_{\mathbf{s},s}$ has the same interpretation in $\mathbf{A}$ and in $(\id_{\sig{\Sigma}})^{\ast}_{\gd}(\mathbf{A})$.

Let $\sigma\in\Sigma_{\mathbf{s},s}$ and assume that every $\tau\lhd\sigma$ has the same interpretation in both algebras. Then, by Definition~\ref{DHepSignature}, a family $\mathbf{a}\in A_{\mathbf{s}}$ satisfies $\mathrm{def}^{(S,\Sigma)}(\sigma)$ in $\mathbf{A}$ if, and only if, it satisfies it in $(\id_{\sig{\Sigma}})^{\ast}_{\gd}(\mathbf{A})$. Moreover, by~(Der3), $\eta^{\Sigma}(\sigma)^{\mathbf{A}}=\sigma^{\mathbf{A}}$. Therefore, the following chain of equalities holds
\allowdisplaybreaks
\begin{align*}
\Dom\left(\sigma^{(\id_{\sig{\Sigma}})^{\ast}_{\gd}(\mathbf{A})}\right)
&=
A_{\mathrm{def}^{(S,\Sigma)}(\sigma)}\cap\Dom\left(\eta^{\Sigma}(\sigma)^{\mathbf{A}}\right)
\tag{1}
\\&=
A_{\mathrm{def}^{(S,\Sigma)}(\sigma)}\cap\Dom\left(\sigma^{\mathbf{A}}\right)
\tag{2}
\\&=
\Dom\left(\sigma^{\mathbf{A}}\right),
\tag{3}
\end{align*}
and, on this domain, $\sigma^{(\id_{\sig{\Sigma}})^{\ast}_{\gd}(\mathbf{A})}$ acts as $\eta^{\Sigma}(\sigma)^{\mathbf{A}}=\sigma^{\mathbf{A}}$.

The first equality follows from Proposition~\ref{PTransformedAlg}, taking into account the induction hypothesis and the fact that all the elements of $A$ are cells, so that condition (D3) always holds; the second equality follows from the fact that $\eta^{\Sigma}(\sigma)^{\mathbf{A}}=\sigma^{\mathbf{A}}$; finally, the last equality follows from the fact that $\mathbf{A}$ is a weak $\sig{\Sigma}$-algebra, so that $\Dom(\sigma^{\mathbf{A}})\subseteq A_{\mathrm{def}^{(S,\Sigma)}(\sigma)}$, see Definition~\ref{DWeakEpAlgebra}.

Finally, since $\mathsf{wEPAlg}(\sig{\Sigma})$ and $\mathsf{EPAlg}(\sig{\Sigma})$ are full subcategories of $\mathsf{PAlg}(\Sigma)$ and every $\sig{\Sigma}$-algebra is a weak $\sig{\Sigma}$-algebra, the last statement follows from the above.

This completes the proof.
\end{proof}

We are now in position to compare the functor induced by a composition with the composite of the induced functors.

\begin{proposition}\label{PHepComp}
Let $\sig{\Sigma}=(S,\Sigma,\mathrm{def}^{(S,\Sigma)})$, $\sig{\Lambda}=(T,\Lambda,\mathrm{def}^{(T,\Lambda)})$ and $\sig{\Omega}=(U,\Omega,\mathrm{def}^{(U,\Omega)})$ be hep-signatures, $\mathbf{d}=(\varphi,d,\Dom^{(\varphi,d)})$ a generalized hep-derivor from $\sig{\Sigma}$ to $\sig{\Lambda}$, $\mathbf{e}=(\psi,e,\Dom^{(\psi,e)})$ a generalized hep-derivor from $\sig{\Lambda}$ to $\sig{\Omega}$ and $\mathbf{C}$ a partial $\Omega$-algebra. Then the following statements hold:
\begin{enumerate}
\item the underlying $S$-sorted set of $\mathbf{d}^{\ast}_{\gd}(\mathbf{e}^{\ast}_{\gd}(\mathbf{C}))$ is an $S$-sorted subset of that of $(\mathbf{e}\circ\mathbf{d})^{\ast}_{\gd}(\mathbf{C})$, and the inclusion $\kappa^{\mathbf{d},\mathbf{e}}_{\mathbf{C}}$ of the former into the latter is a homomorphism of partial $\Sigma$-algebras
$$
\kappa^{\mathbf{d},\mathbf{e}}_{\mathbf{C}}\colon\mathbf{d}^{\ast}_{\gd}\left(\mathbf{e}^{\ast}_{\gd}(\mathbf{C})\right)\mor(\mathbf{e}\circ\mathbf{d})^{\ast}_{\gd}(\mathbf{C});
$$
moreover, for every homomorphism $f\colon\mathbf{C}\mor\mathbf{C}'$ of partial $\Omega$-algebras,
$$
(\mathbf{e}\circ\mathbf{d})^{\ast}_{\gd}(f)\circ\kappa^{\mathbf{d},\mathbf{e}}_{\mathbf{C}}=\kappa^{\mathbf{d},\mathbf{e}}_{\mathbf{C}'}\circ\mathbf{d}^{\ast}_{\gd}\left(\mathbf{e}^{\ast}_{\gd}(f)\right);
$$
\item if $\mathbf{d}^{\ast}_{\gd}(\mathbf{e}^{\ast}_{\gd}(\mathbf{C}))$ and $(\mathbf{e}\circ\mathbf{d})^{\ast}_{\gd}(\mathbf{C})$ have the same underlying $S$-sorted set and $\mathbf{d}^{\ast}_{\gd}(\mathbf{e}^{\ast}_{\gd}(\mathbf{C}))$ is a $\sig{\Sigma}$-algebra, then
$$
(\mathbf{e}\circ\mathbf{d})^{\ast}_{\gd}(\mathbf{C})=\mathbf{d}^{\ast}_{\gd}\left(\mathbf{e}^{\ast}_{\gd}(\mathbf{C})\right);
$$
and, if, moreover, $f\colon\mathbf{C}\mor\mathbf{C}'$ is a homomorphism of partial $\Omega$-algebras and $\mathbf{C}'$ satisfies the same hypotheses, then $(\mathbf{e}\circ\mathbf{d})^{\ast}_{\gd}(f)=\mathbf{d}^{\ast}_{\gd}(\mathbf{e}^{\ast}_{\gd}(f))$.
\end{enumerate}
\end{proposition}
\begin{proof}
\textsf{(Inclusion of the cells)} Let $s\in S$ and $c$ a cell of sort $s$ of $\mathbf{d}^{\ast}_{\gd}(\mathbf{e}^{\ast}_{\gd}(\mathbf{C}))$, i.e., $c\in(C_{\Dom^{(\psi,e)}})_{\varphi(s)}$ and $\mathbf{e}^{\ast}_{\gd}(\mathbf{C})\models\Dom^{(\varphi,d)}_{s}[c]$. By Lemma~\ref{LTermComp}, $\mathbf{C}\models e^{\sharp}[\Dom^{(\varphi,d)}_{s}][c]$. Hence, by Definition~\ref{DHepComp}, $\mathbf{C}\models\Dom^{\mathbf{e}\circ\mathbf{d}}_{s}[c]$, i.e., $c$ is a cell of sort $s$ of $(\mathbf{e}\circ\mathbf{d})^{\ast}_{\gd}(\mathbf{C})$.

\textsf{(The inclusion is a homomorphism)} We prove, by well-founded induction along $\lhd$, that, for every $\sigma\in\Sigma_{\mathbf{s},s}$ and every $\mathbf{a}\in\Dom(\sigma^{\mathbf{d}^{\ast}_{\gd}(\mathbf{e}^{\ast}_{\gd}(\mathbf{C}))})$, the family $\mathbf{a}$ is in $\Dom(\sigma^{(\mathbf{e}\circ\mathbf{d})^{\ast}_{\gd}(\mathbf{C})})$ and
$$
\sigma^{(\mathbf{e}\circ\mathbf{d})^{\ast}_{\gd}(\mathbf{C})}(\mathbf{a})=\sigma^{\mathbf{d}^{\ast}_{\gd}(\mathbf{e}^{\ast}_{\gd}(\mathbf{C}))}(\mathbf{a}).
$$
Let us assume that this holds for every $\tau\lhd\sigma$ and let $\mathbf{a}\in\Dom(\sigma^{\mathbf{d}^{\ast}_{\gd}(\mathbf{e}^{\ast}_{\gd}(\mathbf{C}))})$. We check the conditions (D1), (D2) and (D3) of Proposition~\ref{PTransformedAlg} for $\mathbf{e}\circ\mathbf{d}$.

\textsf{(D1)} By the induction hypothesis, the inclusion is a homomorphism between the reducts of $\mathbf{d}^{\ast}_{\gd}(\mathbf{e}^{\ast}_{\gd}(\mathbf{C}))$ and $(\mathbf{e}\circ\mathbf{d})^{\ast}_{\gd}(\mathbf{C})$ to the operation symbols that occur in $\mathrm{def}^{(S,\Sigma)}(\sigma)$. Since $\mathbf{d}^{\ast}_{\gd}(\mathbf{e}^{\ast}_{\gd}(\mathbf{C}))\models\mathrm{def}^{(S,\Sigma)}(\sigma)[\mathbf{a}]$, by Remark~\ref{RHomTerm}, $(\mathbf{e}\circ\mathbf{d})^{\ast}_{\gd}(\mathbf{C})\models\mathrm{def}^{(S,\Sigma)}(\sigma)[\mathbf{a}]$.

\textsf{(D2)} By Proposition~\ref{PTransformedAlg}, $\mathbf{a}\in\Dom(d(\sigma)^{\mathbf{e}^{\ast}_{\gd}(\mathbf{C})})$. By Lemma~\ref{LTermComp}, $\mathbf{a}\in\Dom(e^{\sharp}(d(\sigma))^{\mathbf{C}})$ and $e^{\sharp}(d(\sigma))^{\mathbf{C}}(\mathbf{a})=d(\sigma)^{\mathbf{e}^{\ast}_{\gd}(\mathbf{C})}(\mathbf{a})$.

\textsf{(D3)} By Proposition~\ref{PTransformedAlg}, $d(\sigma)^{\mathbf{e}^{\ast}_{\gd}(\mathbf{C})}(\mathbf{a})=\sigma^{\mathbf{d}^{\ast}_{\gd}(\mathbf{e}^{\ast}_{\gd}(\mathbf{C}))}(\mathbf{a})$ is a cell of $\mathbf{d}^{\ast}_{\gd}(\mathbf{e}^{\ast}_{\gd}(\mathbf{C}))$, hence, by the inclusion of the cells, a cell of $(\mathbf{e}\circ\mathbf{d})^{\ast}_{\gd}(\mathbf{C})$.

Therefore, by Proposition~\ref{PTransformedAlg}, $\mathbf{a}\in\Dom(\sigma^{(\mathbf{e}\circ\mathbf{d})^{\ast}_{\gd}(\mathbf{C})})$ and the following chain of equalities holds
\allowdisplaybreaks
\begin{align*}
\sigma^{(\mathbf{e}\circ\mathbf{d})^{\ast}_{\gd}(\mathbf{C})}\left(\mathbf{a}\right)
&=
e^{\sharp}(d(\sigma))^{\mathbf{C}}\left(\mathbf{a}\right)
\tag{1}
\\&=
d(\sigma)^{\mathbf{e}^{\ast}_{\gd}(\mathbf{C})}\left(\mathbf{a}\right)
\tag{2}
\\&=
\sigma^{\mathbf{d}^{\ast}_{\gd}(\mathbf{e}^{\ast}_{\gd}(\mathbf{C}))}\left(\mathbf{a}\right).
\tag{3}
\end{align*}

The first equality follows from Proposition~\ref{PTransformedAlg} applied to $\mathbf{e}\circ\mathbf{d}$; the second equality follows from Lemma~\ref{LTermComp}; finally, the last equality follows from Proposition~\ref{PTransformedAlg} applied to $\mathbf{d}$.

\textsf{(Naturality)} By Proposition~\ref{PGenHepDerivorFunctor}, for every $s\in S$, both $((\mathbf{e}\circ\mathbf{d})^{\ast}_{\gd}(f)\circ\kappa^{\mathbf{d},\mathbf{e}}_{\mathbf{C}})_{s}$ and $(\kappa^{\mathbf{d},\mathbf{e}}_{\mathbf{C}'}\circ\mathbf{d}^{\ast}_{\gd}(\mathbf{e}^{\ast}_{\gd}(f)))_{s}$ send a cell $c$ of sort $s$ of $\mathbf{d}^{\ast}_{\gd}(\mathbf{e}^{\ast}_{\gd}(\mathbf{C}))$ to $f_{\psi(\varphi(s))}(c)$. This proves statement~(1).

\textsf{(Statement (2))} Let us assume that $\mathbf{d}^{\ast}_{\gd}(\mathbf{e}^{\ast}_{\gd}(\mathbf{C}))$ and $(\mathbf{e}\circ\mathbf{d})^{\ast}_{\gd}(\mathbf{C})$ have the same underlying $S$-sorted set and that $\mathbf{d}^{\ast}_{\gd}(\mathbf{e}^{\ast}_{\gd}(\mathbf{C}))$ is a $\sig{\Sigma}$-algebra. By statement~(1), for every $\sigma\in\Sigma_{\mathbf{s},s}$,
$$
\Dom\left(\sigma^{\mathbf{d}^{\ast}_{\gd}(\mathbf{e}^{\ast}_{\gd}(\mathbf{C}))}\right)\subseteq\Dom\left(\sigma^{(\mathbf{e}\circ\mathbf{d})^{\ast}_{\gd}(\mathbf{C})}\right)
$$
and both operations agree on $\Dom(\sigma^{\mathbf{d}^{\ast}_{\gd}(\mathbf{e}^{\ast}_{\gd}(\mathbf{C}))})$. We prove, by well-founded induction along $\lhd$, the reverse inclusion. Let us assume that every $\tau\lhd\sigma$ has the same interpretation in both partial $\Sigma$-algebras and let $\mathbf{a}\in\Dom(\sigma^{(\mathbf{e}\circ\mathbf{d})^{\ast}_{\gd}(\mathbf{C})})$. Then, by Proposition~\ref{PTransformedAlg}, $(\mathbf{e}\circ\mathbf{d})^{\ast}_{\gd}(\mathbf{C})\models\mathrm{def}^{(S,\Sigma)}(\sigma)[\mathbf{a}]$, hence, by the induction hypothesis and Definition~\ref{DHepSignature}, $\mathbf{d}^{\ast}_{\gd}(\mathbf{e}^{\ast}_{\gd}(\mathbf{C}))\models\mathrm{def}^{(S,\Sigma)}(\sigma)[\mathbf{a}]$. Since $\mathbf{d}^{\ast}_{\gd}(\mathbf{e}^{\ast}_{\gd}(\mathbf{C}))$ is a $\sig{\Sigma}$-algebra, by Definition~\ref{DEpAlgebra}, $\mathbf{a}\in\Dom(\sigma^{\mathbf{d}^{\ast}_{\gd}(\mathbf{e}^{\ast}_{\gd}(\mathbf{C}))})$. Therefore both partial $\Sigma$-algebras are equal, i.e., $\kappa^{\mathbf{d},\mathbf{e}}_{\mathbf{C}}$ is the identity. If $\mathbf{C}'$ satisfies the same hypotheses, then $\kappa^{\mathbf{d},\mathbf{e}}_{\mathbf{C}'}$ is also the identity and, by statement~(1), $(\mathbf{e}\circ\mathbf{d})^{\ast}_{\gd}(f)=\mathbf{d}^{\ast}_{\gd}(\mathbf{e}^{\ast}_{\gd}(f))$.

This completes the proof.
\end{proof}

By Lemma~\ref{LTermComp}, the hypothesis on the cells in Proposition~\ref{PHepComp}~(2) amounts to the following condition: for every $s\in S$ and every cell $c$ of sort $\varphi(s)$ of $\mathbf{C}$ relative to $\mathbf{e}$, if $\mathbf{C}\models e^{\sharp}[\Dom^{(\varphi,d)}_{s}][c]$, then $\mathbf{e}^{\ast}_{\gd}(\mathbf{C})\models\Dom^{(\varphi,d)}_{s}[c]$. It trivially holds when the domains of $\mathbf{d}$ are empty.

\subsection{The category \texorpdfstring{$\mathsf{hepAlg}_{\gd}$}{hepAlggd}}\label{SHepAlgCat}

In~\cite{CVST10b}, the category $\mathsf{Alg}_{\mathfrak{d}}$, whose objects are the pairs formed by a many-sorted signature and an algebra for it, is obtained by means of the Grothendieck construction, see Definition~\ref{DGConst}, for the split indexed category that sends a many-sorted signature $\sig{\Sigma}$ to $\mathsf{Alg}(\Sigma)$ and a derivor to the functor that it induces. Following this pattern, we next assemble the ep-algebras of all the hep-signatures into a single category. However, by Proposition~\ref{PHepComp}, the composite $\mathbf{d}^{\ast}_{\gd}\circ\mathbf{e}^{\ast}_{\gd}$ is, in general, only contained in $(\mathbf{e}\circ\mathbf{d})^{\ast}_{\gd}$, and the inclusions $\kappa^{\mathbf{d},\mathbf{e}}_{\mathbf{C}}$ need not be isomorphisms. Therefore the induced functors do not constitute a split, nor even a pseudo, indexed category, but a normal lax indexed category, see Definition~\ref{DLaxIndexed}, provided that we take as fibres the categories of weak ep-algebras, see Definition~\ref{DWeakEpAlgebra}, on which, by Proposition~\ref{PHepId}, the identity generalized hep-derivors induce identity functors. We begin by proving that the transformed algebras are weak ep-algebras.

\begin{lemma}\label{LHepWeak}
Let $\sig{\Sigma}=(S,\Sigma,\mathrm{def}^{(S,\Sigma)})$ and $\sig{\Lambda}=(T,\Lambda,\mathrm{def}^{(T,\Lambda)})$ be hep-signatures and $\mathbf{d}$ a generalized hep-derivor from $\sig{\Sigma}$ to $\sig{\Lambda}$. Then, for every partial $\Lambda$-algebra $\mathbf{B}$, $\mathbf{d}^{\ast}_{\gd}(\mathbf{B})$ is a weak $\sig{\Sigma}$-algebra. Consequently, the restriction of $\mathbf{d}^{\ast}_{\gd}$ to $\mathsf{wEPAlg}(\sig{\Lambda})$ corestricts to $\mathsf{wEPAlg}(\sig{\Sigma})$, see Definition~\ref{DCorestriction}, and we denote the corestriction by
$$
\mathbf{d}^{\ast}_{\gd}\colon\mathsf{wEPAlg}(\sig{\Lambda})\mor\mathsf{wEPAlg}(\sig{\Sigma}).
$$
\end{lemma}
\begin{proof}
Let $\mathbf{B}$ be a partial $\Lambda$-algebra, $(\mathbf{s},s)\in S^{\star}\times S$ and $\sigma\in\Sigma_{\mathbf{s},s}$. By condition (D1) of Proposition~\ref{PTransformedAlg}, every family in $\Dom(\sigma^{\mathbf{d}^{\ast}_{\gd}(\mathbf{B})})$ satisfies $\mathrm{def}^{(S,\Sigma)}(\sigma)$ in $\mathbf{d}^{\ast}_{\gd}(\mathbf{B})$. Hence, by Definition~\ref{DWeakEpAlgebra}, $\mathbf{d}^{\ast}_{\gd}(\mathbf{B})$ is a weak $\sig{\Sigma}$-algebra. The last statement follows from Definition~\ref{DCorestriction}, since $\mathsf{wEPAlg}(\sig{\Lambda})$ and $\mathsf{wEPAlg}(\sig{\Sigma})$ are full subcategories of $\mathsf{PAlg}(\Lambda)$ and of $\mathsf{PAlg}(\Sigma)$, respectively.
\end{proof}

In the following proposition we prove that the categories of weak ep-algebras, together with the functors induced by the generalized hep-derivors and the inclusions of Proposition~\ref{PHepComp}, constitute a normal lax indexed category.

\begin{proposition}\label{PHepLax}
Let $\mathrm{wEPAlg}$ stand for the assignment that sends every hep-signature $\sig{\Sigma}$ to the category $\mathsf{wEPAlg}(\sig{\Sigma})$ and every generalized hep-derivor $\mathbf{d}$ from $\sig{\Sigma}$ to $\sig{\Lambda}$ to the functor $\mathbf{d}^{\ast}_{\gd}\colon\mathsf{wEPAlg}(\sig{\Lambda})\mor\mathsf{wEPAlg}(\sig{\Sigma})$ of Lemma~\ref{LHepWeak}, and let $\kappa$ stand for the assignment that sends every pair of generalized hep-derivors $\mathbf{d}\colon\sig{\Sigma}\mor\sig{\Lambda}$ and $\mathbf{e}\colon\sig{\Lambda}\mor\sig{\Omega}$ to the family
$$
\kappa^{\mathbf{d},\mathbf{e}}=\left(\kappa^{\mathbf{d},\mathbf{e}}_{\mathbf{C}}\right)_{\mathbf{C}\in\mathrm{Ob}(\mathsf{wEPAlg}(\sig{\Omega}))}
$$
of the inclusions of Proposition~\ref{PHepComp}. Then $(\mathrm{wEPAlg},\kappa)$ is a normal lax $\mathsf{hepSig}_{\gd}$-indexed category, see Definition~\ref{DLaxIndexed}.
\end{proposition}
\begin{proof}
\textsf{(Identities)} By Proposition~\ref{PHepId}, for every hep-signature $\sig{\Sigma}$, $(\id_{\sig{\Sigma}})^{\ast}_{\gd}=\Id_{\mathsf{wEPAlg}(\sig{\Sigma})}$.

\textsf{(Naturality)} Let $\mathbf{d}\colon\sig{\Sigma}\mor\sig{\Lambda}$ and $\mathbf{e}\colon\sig{\Lambda}\mor\sig{\Omega}$ be generalized hep-derivors. By Lemma~\ref{LHepWeak}, for every weak $\sig{\Omega}$-algebra $\mathbf{C}$, both $\mathbf{d}^{\ast}_{\gd}(\mathbf{e}^{\ast}_{\gd}(\mathbf{C}))$ and $(\mathbf{e}\circ\mathbf{d})^{\ast}_{\gd}(\mathbf{C})$ are weak $\sig{\Sigma}$-algebras, so that, by Proposition~\ref{PHepComp}~(1), $\kappa^{\mathbf{d},\mathbf{e}}_{\mathbf{C}}$ is a morphism of the full subcategory $\mathsf{wEPAlg}(\sig{\Sigma})$ of $\mathsf{PAlg}(\Sigma)$. Moreover, again by Proposition~\ref{PHepComp}~(1), $\kappa^{\mathbf{d},\mathbf{e}}$ is a natural transformation from $\mathbf{d}^{\ast}_{\gd}\circ\mathbf{e}^{\ast}_{\gd}$ to $(\mathbf{e}\circ\mathbf{d})^{\ast}_{\gd}$.

Before checking the conditions (L1) and (L2) of Definition~\ref{DLaxIndexed}, let us note that a homomorphism of partial algebras is determined by its domain, its codomain and its underlying many-sorted mapping, and that, by Proposition~\ref{PHepComp}~(1) and Proposition~\ref{PGenHepDerivorFunctor}, the underlying many-sorted mapping of every $\kappa^{\mathbf{d},\mathbf{e}}_{\mathbf{C}}$, and of its image under every functor $\mathbf{g}^{\ast}_{\gd}$, which is a birestriction of it, is an inclusion.

\textsf{(L1)} Let $\mathbf{d}\colon\sig{\Sigma}\mor\sig{\Lambda}$ be a generalized hep-derivor and $\mathbf{B}$ a weak $\sig{\Lambda}$-algebra. By Proposition~\ref{PHepSigCat}, $\mathbf{d}\circ\id_{\sig{\Sigma}}=\mathbf{d}$ and, by Lemma~\ref{LHepWeak} and Proposition~\ref{PHepId}, $(\id_{\sig{\Sigma}})^{\ast}_{\gd}(\mathbf{d}^{\ast}_{\gd}(\mathbf{B}))=\mathbf{d}^{\ast}_{\gd}(\mathbf{B})$. Hence $\kappa^{\id_{\sig{\Sigma}},\mathbf{d}}_{\mathbf{B}}$ is the inclusion of the underlying $S$-sorted set of $\mathbf{d}^{\ast}_{\gd}(\mathbf{B})$ into itself, i.e., $\kappa^{\id_{\sig{\Sigma}},\mathbf{d}}_{\mathbf{B}}=\id_{\mathbf{d}^{\ast}_{\gd}(\mathbf{B})}$. Similarly, $\id_{\sig{\Lambda}}\circ\mathbf{d}=\mathbf{d}$ and, by Proposition~\ref{PHepId}, $\mathbf{d}^{\ast}_{\gd}((\id_{\sig{\Lambda}})^{\ast}_{\gd}(\mathbf{B}))=\mathbf{d}^{\ast}_{\gd}(\mathbf{B})$, so that $\kappa^{\mathbf{d},\id_{\sig{\Lambda}}}_{\mathbf{B}}=\id_{\mathbf{d}^{\ast}_{\gd}(\mathbf{B})}$. Therefore $\kappa^{\id_{\sig{\Sigma}},\mathbf{d}}=\id_{\mathbf{d}^{\ast}_{\gd}}=\kappa^{\mathbf{d},\id_{\sig{\Lambda}}}$.

\textsf{(L2)} Let $\mathbf{d}\colon\sig{\Sigma}\mor\sig{\Lambda}$, $\mathbf{e}\colon\sig{\Lambda}\mor\sig{\Omega}$ and $\mathbf{g}\colon\sig{\Omega}\mor\sig{\Xi}$ be generalized hep-derivors and $\mathbf{D}$ a weak $\sig{\Xi}$-algebra. By Proposition~\ref{PHepSigCat}, $(\mathbf{g}\circ\mathbf{e})\circ\mathbf{d}=\mathbf{g}\circ(\mathbf{e}\circ\mathbf{d})$. Therefore both
$$
\kappa^{\mathbf{d},\mathbf{g}\circ\mathbf{e}}_{\mathbf{D}}\circ\mathbf{d}^{\ast}_{\gd}\left(\kappa^{\mathbf{e},\mathbf{g}}_{\mathbf{D}}\right)
\qquad\text{and}\qquad
\kappa^{\mathbf{e}\circ\mathbf{d},\mathbf{g}}_{\mathbf{D}}\circ\kappa^{\mathbf{d},\mathbf{e}}_{\mathbf{g}^{\ast}_{\gd}(\mathbf{D})}
$$
are homomorphisms from $\mathbf{d}^{\ast}_{\gd}(\mathbf{e}^{\ast}_{\gd}(\mathbf{g}^{\ast}_{\gd}(\mathbf{D})))$ to $(\mathbf{g}\circ\mathbf{e}\circ\mathbf{d})^{\ast}_{\gd}(\mathbf{D})$ whose underlying $S$-sorted mappings are composites of inclusions, i.e., the inclusion of the underlying $S$-sorted set of the former into that of the latter. Hence both homomorphisms are equal.

This completes the proof.
\end{proof}

We are now in position to define the category of the ep-algebras of all the hep-signatures.

\begin{definition}\label{DHepAlgCat}
We let $\mathsf{hepAlg}_{\gd}$ stand for the full subcategory of $\int^{\mathsf{hepSig}_{\gd}}(\mathrm{wEPAlg},\kappa)$, see Definition~\ref{DLaxGConst} and Proposition~\ref{PHepLax}, whose objects are the ordered pairs $(\sig{\Sigma},\mathbf{A})$, where $\sig{\Sigma}$ is a hep-signature and $\mathbf{A}$ a $\sig{\Sigma}$-algebra. Explicitly, $\mathsf{hepAlg}_{\gd}$ is defined as follows:
\begin{enumerate}
\item its objects are the ordered pairs $(\sig{\Sigma},\mathbf{A})$, where $\sig{\Sigma}$ is a hep-signature and $\mathbf{A}$ a $\sig{\Sigma}$-algebra;
\item its morphisms from $(\sig{\Sigma},\mathbf{A})$ to $(\sig{\Lambda},\mathbf{B})$ are the ordered pairs $(\mathbf{d},f)$, where $\mathbf{d}$ is a generalized hep-derivor from $\sig{\Sigma}$ to $\sig{\Lambda}$ and $f\colon\mathbf{A}\mor\mathbf{d}^{\ast}_{\gd}(\mathbf{B})$ a homomorphism of partial $\Sigma$-algebras;
\item the composite of $(\mathbf{d},f)\colon(\sig{\Sigma},\mathbf{A})\mor(\sig{\Lambda},\mathbf{B})$ and $(\mathbf{e},g)\colon(\sig{\Lambda},\mathbf{B})\mor(\sig{\Omega},\mathbf{C})$ is
$$
(\mathbf{e},g)\circ(\mathbf{d},f)=\left(\mathbf{e}\circ\mathbf{d},\kappa^{\mathbf{d},\mathbf{e}}_{\mathbf{C}}\circ\mathbf{d}^{\ast}_{\gd}(g)\circ f\right),
$$
see Definition~\ref{DHepComp} and Proposition~\ref{PHepComp}; and
\item the identity at $(\sig{\Sigma},\mathbf{A})$ is $(\id_{\sig{\Sigma}},\id_{\mathbf{A}})$.
\end{enumerate}
\end{definition}

Let us note that, since every $\sig{\Sigma}$-algebra is a weak $\sig{\Sigma}$-algebra and $\mathsf{wEPAlg}(\sig{\Sigma})$ is a full subcategory of $\mathsf{PAlg}(\Sigma)$, the morphisms of $\int^{\mathsf{hepSig}_{\gd}}(\mathrm{wEPAlg},\kappa)$ between objects of $\mathsf{hepAlg}_{\gd}$ are exactly the pairs described in the second item of the above definition, so that both descriptions of $\mathsf{hepAlg}_{\gd}$ agree.

\begin{proposition}\label{PHepAlgCat}
$\mathsf{hepAlg}_{\gd}$ is a category. Moreover, the assignment that sends every object $(\sig{\Sigma},\mathbf{A})$ to $\sig{\Sigma}$ and every morphism $(\mathbf{d},f)$ to $\mathbf{d}$ is a functor from $\mathsf{hepAlg}_{\gd}$ to $\mathsf{hepSig}_{\gd}$.
\end{proposition}
\begin{proof}
By Propositions~\ref{PLaxGConst} and~\ref{PHepLax}, $\int^{\mathsf{hepSig}_{\gd}}(\mathrm{wEPAlg},\kappa)$ is a category and $\pi_{(\mathrm{wEPAlg},\kappa)}$ is a functor from it to $\mathsf{hepSig}_{\gd}$. Therefore its full subcategory $\mathsf{hepAlg}_{\gd}$ is a category, and the assignment of the statement, which is the restriction of $\pi_{(\mathrm{wEPAlg},\kappa)}$ to $\mathsf{hepAlg}_{\gd}$, is a functor.

This completes the proof.
\end{proof}

\begin{remark}\label{RHepAlgCat}
Let us consider the many-sorted signatures as hep-signatures with empty sets of domain conditions, see Remark~\ref{REpAlgebra}, and the derivors $(\varphi,d)$ between them as the generalized hep-derivors $(\varphi,d,\varnothing)$. For a many-sorted signature $(S,\Sigma)$, the weak $(S,\Sigma,\varnothing)$-algebras are all the partial $\Sigma$-algebras and the $(S,\Sigma,\varnothing)$-algebras are the $\Sigma$-algebras. Moreover, by Remark~\ref{RGenHepDerivorCases}~(1), the functors induced by the derivors send total algebras to total algebras and, since the domains of these generalized hep-derivors are empty, by the note after Proposition~\ref{PHepComp} and Proposition~\ref{PHepComp}~(2), the inclusions $\kappa^{\mathbf{d},\mathbf{e}}_{\mathbf{C}}$ are identities when $\mathbf{C}$ is a total algebra. Therefore the subcategory of $\mathsf{hepAlg}_{\gd}$ whose objects are the pairs $((S,\Sigma,\varnothing),\mathbf{A})$, with $\mathbf{A}$ a $\Sigma$-algebra, and whose morphisms are the pairs $((\varphi,d,\varnothing),f)$, is the category $\mathsf{Alg}_{\mathfrak{d}}$ of~\cite{CVST10b}.
\end{remark}

\section{Single-set presentation of higher-order categorial algebras}\label{SSingleSet}

In this section, we recall the definitions of single-set category, single-set $2$-category, single-set $n$-category, and many-sorted categorial algebra, together with their $\omega$-dimensional counterparts. These are the single-set presentations of the corresponding categorical structures, in which all the cells, whatever their dimension, form a single set and the dimension of a cell is read off from its sources and targets, as in~\cite{Str87}. Next we introduce the single-set $n$-categorial signature determined by a many-sorted signature and the $\QE$-variety of its models. We end this section by showing that the many-sorted $n$-categorial algebras are exactly the models of this $\QE$-variety, which is an ep-variety.

We begin by recalling the notion of single-set category.

\begin{definition}\label{DSingleSetCat}
A \emph{single-set category} is an ordered pair $\mathsf{C} = (C,\xi)$ where $C$ is a set and $\xi = (\circ,\src,\tgt)$ a structure of category on $C$, i.e., $\circ\colon C\times C\dmor C$ is a partial operation of \emph{composition}, and $\src$, $\tgt$ are unary operations on $C$ which assign, to every morphism in $C$, its \emph{source} and its \emph{target}, respectively, subject to satisfy the following conditions:
\begin{enumerate}
\item[(C1)] For every $f\in C$, 
\allowdisplaybreaks
\begin{align*}
\src(\src(f)) &= \src(f),&
 \src(\tgt(f)) &= \tgt(f),\\
\tgt(\src(f)) &= \src(f),&
 \tgt(\tgt(f)) &= \tgt(f).
\end{align*}
\item[(C2)] For every $f,g\in C$, 
\allowdisplaybreaks
\begin{align*}
g\circ f &\mbox{ is defined}&
 \mbox{if and only if}&&
  \src(g)& = \tgt(f).
\end{align*} 
\item[(C3)] For every $f,g\in C$, if $\src(g) = \tgt(f)$, then 
\allowdisplaybreaks
\begin{align*}
\src(g\circ f) &= \src(f),
&
\tgt(g\circ f) &= \tgt(g).
\end{align*}
\item[(C4)] For every $f\in C$, 
\allowdisplaybreaks
\begin{align*}
f\circ \src(f) &= f,
&
\tgt(f)\circ f &= f.
\end{align*}
\item[(C5)] For every $f,g,h\in C$, if $\src(h)=\tgt(g)$ and $\src(g)=\tgt(f)$, then 
\allowdisplaybreaks
\begin{align*}
h\circ (g\circ f) = (h\circ g)\circ f.
\end{align*}
\end{enumerate}

We will call the $f\in C$ such that $\src(f) = f$ or, what is equivalent, $\tgt(f) = f$,  \emph{identities} of $C$. 

Let $\mathsf{C}=(C,\xi)$ and $\mathsf{C}'=(C',\xi')$ be two single-set categories. A \emph{single-set functor} from $\mathsf{C}$ to $\mathsf{C}'$ is an ordered triple $(\mathsf{C},F,\mathsf{C}')$, denoted by
$F\colon \mathsf{C}\mor\mathsf{C}'$, where $F$ is a mapping from $C$ to $C'$ such that
\begin{itemize}
\item[(Ci)] If $f\in C$, then
\allowdisplaybreaks
\begin{align*}
F(\src(f))&=\src'(F(f)),
&
F(\tgt(f))&=\tgt'(F(f)).
\end{align*}
\item[(Cii)] If $f,g\in C$ are such that $\src(g)=\tgt(f)$, then
\allowdisplaybreaks
\begin{align*}
F(g\circ f)=F(g)\circ' F(f).
\end{align*}
\end{itemize}

We will say that a single-set category $\mathsf{C}$ is \emph{discrete} if every $f\in C$ is an identity (i.e., if it contains no non-identity morphisms).

Single-set functors can be composed. In fact, if $F$ is a single-set functor from $\mathsf{C}$ to $\mathsf{D}$ and $G$ is a single-set functor from $\mathsf{D}$ to $\mathsf{E}$, then $G\circ F$ is a single-set functor from $\mathsf{C}$ to $\mathsf{E}$, which we call the \emph{composite} of $F$ and $G$. This composition is associative. Moreover, for every single-set category $\mathsf{C}$, there is an \emph{identity} single-set functor $(\mathsf{C},\Id_{C},\mathsf{C})$, denoted by $\Id_{\mathsf{C}}\colon \mathsf{C}\mor\mathsf{C}$, where $\Id_{C}$ is the identity mapping  at $C$, which is the neutral element for the composition. We denote by $\mathsf{Cat}^{}$ the category consisting of single-set categories and single-set functors.
\end{definition}

We next recall the notion of single-set $2$-category, in which two structures of single-set category on the same set interact as in a $2$-category.

\begin{definition}\label{DSingleSet2Cat}
A \emph{single-set $2$-category} is an ordered pair $\mathsf{C} = (C,(\xi_{k})_{k\in 2})$
where $C$ is a set and, for every $k\in 2$, $\xi_{k}= (\circ^{k},\src^{k},\tgt^{k})$ a structure of single-set category on $C$, such that the following conditions are satisfied:
\begin{enumerate}
\item[(2C1)] For every $f\in C$,
\allowdisplaybreaks
\begin{align*}
\src^{1}(\src^{0}(f))&=\src^{0}(f)
=\src^{0}(\src^{1}(f))=\src^{0}(\tgt^{1}(f)),\\
\tgt^{1}(\tgt^{0}(f))&=\tgt^{0}(f)
=\tgt^{0}(\tgt^{1}(f))=\tgt^{0}(\src^{1}(f)).
\end{align*}
\item[(2C2)] For every $f,g\in C$, if $\src^{0}(g)=\tgt^{0}(f)$, then
\allowdisplaybreaks
\begin{align*}
\src^{1}(g{\circ^{0}}f)&=\src^{1}(g){\circ^{0}}\src^{1}(f),
&
\tgt^{1}(g{\circ^{0}}f)&=\tgt^{1}(g){\circ^{0}}\tgt^{1}(f).
\end{align*}
\item[(2C3)] For every $f,f',g,g'\in C$, if 
\allowdisplaybreaks
\begin{align*}
\src^{1}(g')&=\tgt^{1}(g);
&
\src^{1}(f')&=\tgt^{1}(f);\\
\src^{0}(g')&=\tgt^{0}(f');
&
\src^{0}(g)&=\tgt^{0}(f),
\end{align*}
then
\allowdisplaybreaks
\begin{align*}
(g'\circ^{0}f')\circ^{1}(g\circ^{0} f)
=
(g'\circ^{1}g)\circ^{0}(f'\circ^{1}f)
.
\end{align*}
\end{enumerate}

We identify the family $(\xi_{k})_{k\in 2}$ with the ordered pair $(\xi_{0},\xi_{1})$.

Let $\mathsf{C}=(C,(\xi_{k})_{k\in 2})$ and $\mathsf{C}'=(C',(\xi'_{k})_{k\in 2})$ be two single-set $2$-categories. A \emph{single-set $2$-functor} from $\mathsf{C}$ to $\mathsf{C}'$ is an ordered triple $(\mathsf{C},F,\mathsf{C}')$, denoted by $F\colon \mathsf{C}\mor\mathsf{C}'$, where $F$ is a mapping from $C$ to $C'$ such that, for every $k\in 2$, $F$ is a single-set functor from $(C,\xi_{k})$ to $(C',\xi'_{k})$. 

Single-set $2$-functors can be composed. In fact, if $F$ is a single-set $2$-functor from $\mathsf{C}$ to $\mathsf{D}$ and $G$ is a single-set $2$-functor from $\mathsf{D}$ to $\mathsf{E}$, then $G\circ F$ is a single-set $2$-functor from $\mathsf{C}$ to $\mathsf{E}$, which we call the \emph{composite} of $F$ and $G$. This composition is associative. Moreover, for every single-set $2$-category $\mathsf{C}$, there is an \emph{identity} single-set $2$-functor $(\mathsf{C},\Id_{C},\mathsf{C})$, denoted by $\Id_{\mathsf{C}}\colon \mathsf{C}\mor\mathsf{C}$, where $\Id_{C}$ is the identity mapping  at $C$, which is the neutral element for the composition. We denote by $\mathsf{2Cat}^{}$ the category consisting of single-set 
$2$-categories and single-set $2$-functors.
\end{definition}

We next recall the notion of single-set $n$-category, for $n\in\mathbb{N}$, in which every pair of the $n$ structures of single-set category forms a single-set $2$-category.

\begin{definition}\label{DSingleSetnCat} 
A \emph{single-set $n$-category}, for $n\in\mathbb{N}$, is an ordered pair $\mathsf{C} = (C,(\xi_{k})_{k\in n})$ where $C$ is a set and, for every $k\in n$, $\xi_{k}$ a structure of  single-set category on $C$ such that, for every $j,k\in n$ with $j<k$, the ordered pair $(C,(\xi_{i})_{i\in\{j,k\}})$ is a single-set $2$-category. For $j<k$, we identify $(\xi_{i})_{i\in \{j,k\}}$ with $(\xi_{j},\xi_{k})$. Thus a single-set $0$-category is just a set, a single-set $1$-category is a single-set category in the sense of Definition~\ref{DSingleSetCat}, and a single-set $2$-category is one in the sense of Definition~\ref{DSingleSet2Cat}.

For \(k\in n\) we will say that \(f\in C\) is a \emph{\(k\)-cell} if \(f=\src^{k}(f)\) or, equivalently, by item~(C1) in Definition~\ref{DSingleSetCat}, if \(f=\tgt^{k}(f)\). Then we denote by \(C_{k}\) the set of \(k\)-cells of \(\mathsf{C}\), i.e.,
\[
C_{k}=\{f\in C\mid f=\src^{k}(f)\}.
\]
Moreover, we let \(C_{n}\) stand for \(C\) and, by analogy to the above, the elements of \(C_{n}\) will be called \emph{\(n\)-cells}.

A \emph{single-set $n$-functor} $F$ from $\mathsf{C}$ to $\mathsf{C}'$, denoted by
$F\colon \mathsf{C}\mor\mathsf{C}'$, is a mapping from $C$ to $C'$ such that, for every $k\in n$, $F$ is a single-set functor from $(C, \xi_{k})$ to $(C', \xi'_{k})$. 

Single-set $n$-functors can be composed. In fact, if $F$ is a single-set $n$-functor from $\mathsf{C}$ to $\mathsf{D}$ and $G$ is a single-set $n$-functor from $\mathsf{D}$ to $\mathsf{E}$, then $G\circ F$ is a single-set $n$-functor from $\mathsf{C}$ to $\mathsf{E}$, which we call the \emph{composite} of $F$ and $G$. This composition is associative. Moreover, for every single-set $n$-category $\mathsf{C}$, there is an \emph{identity} single-set $n$-functor $(\mathsf{C},\Id_{C},\mathsf{C})$, denoted by $\Id_{\mathsf{C}}\colon \mathsf{C}\mor\mathsf{C}$, where $\Id_{C}$ is the identity mapping  at $C$, which is the neutral element for the composition. We denote by $\mathsf{nCat}^{}$ the category consisting of single-set $n$-categories and single-set $n$-functors.
\end{definition}

In the following definition we introduce the many-sorted $n$-categorial algebras, i.e., the many-sorted algebras whose underlying sorts are single-set $n$-categories and whose operations are single-set $n$-functors.

\begin{definition}\label{DManySortedCatAlg}
Let $S$ be a set of sorts and $n\in\mathbb{N}$.
We will call the objects of the functor category $\mathsf{nCat}^{S}$ \emph{$S$-sorted single-set $n$-categories}. Thus, an $S$-sorted single-set $n$-category is an $S$-indexed family $\mathsf{A}=(\mathsf{A}_{s})_{s\in S}$ such that, for every sort $s\in S$, $\mathsf{A}_{s}$ is a single-set $n$-category. Given two $S$-sorted single-set $n$-categories 
$\mathsf{A}$ and $\mathsf{B}$, we will call the morphisms in $\mathsf{nCat}^{S}$ from $\mathsf{A}$ to 
$\mathsf{B}$ \emph{$S$-sorted single-set $n$-functors}. Thus, an $S$-sorted single-set $n$-functor from $\mathsf{A}$ to $\mathsf{B}$ is an $S$-indexed family $F=(F_{s})_{s\in S}$ such that, for every sort $s\in S$, $F_{s}$ is a single-set $n$-functor from $\mathsf{A}_{s}$ to $\mathsf{B}_{s}$. We will denote by $\mathrm{nFunc}(\mathsf{A},\mathsf{B})$ the set of all $S$-sorted single-set $n$-functors from $\mathsf{A}$ to $\mathsf{B}$. We consider the expressions $F\in\mathrm{nFunc}(\mathsf{A},\mathsf{B})$ and $F\colon\mathsf{A}\mor\mathsf{B}$ as synonyms.

Let $\mathsf{A}$ be an $S$-sorted single-set $n$-category. The $S^{\star}\times S$-sorted set of finitary operations on $\mathsf{A}$ is $(\mathrm{nFunc}(\mathsf{A}_{\mathbf{s}},\mathsf{A}_{s}))_{(\mathbf{s},s)\in S^{\star}\times S}$, where, for every $\mathbf{s}\in S^{\star}$, $\mathsf{A}_{\mathbf{s}}=\prod_{j\in\bb{\mathbf{s}}}\mathsf{A}_{s_{j}}$, with $\bb{\mathbf{s}}$ denoting the length of the word $\mathbf{s}$ (if $\mathbf{s}=\lambda$, then $\mathsf{A}_{\lambda}$ is a final single-set $n$-category). 

Let $\sig{\Sigma}=(S,\Sigma)$ be a many-sorted signature. 
 A \emph{structure of many-sorted $n$-categorial $\Sigma$-algebra on} an $S$-sorted $n$-category $\mathsf{A}$ is a family $F = (F_{\mathbf{s},s})_{(\mathbf{s},s)\in S^{\star}\times S}$, where, for $(\mathbf{s},s)\in S^{\star}\times S$, $F_{\mathbf{s},s}$ is a mapping from $\Sigma_{\mathbf{s},s}$ to $\mathrm{nFunc}(\mathsf{A}_{\mathbf{s}},\mathsf{A}_{s})$ (if $(\mathbf{s},s) = (\lambda,s)$ and $\sigma\in \Sigma_{\lambda,s}$, then $F_{\lambda,s}(\sigma)$ picks out a cell of $\mathsf{A}_{s}$ which is a $k$-cell for every $k\in n$, since $\mathsf{A}_{\lambda}$ is a final single-set $n$-category). The finite products $\mathsf{A}_{\mathbf{s}}$ are taken componentwise: sources and targets are computed componentwise, and $(\mathsf{g}_{j})_{j}\circ^{k}(\mathsf{f}_{j})_{j}$ is defined if and only if every $\mathsf{g}_{j}\circ^{k}\mathsf{f}_{j}$ is defined, and then it is $(\mathsf{g}_{j}\circ^{k}\mathsf{f}_{j})_{j}$.  An \emph{$S$-sorted single-set $n$-categorial $\Sigma$-algebra} or, simply, a \emph{many-sorted $n$-categorial $\sig{\Sigma}$-algebra}, is a pair $(\mathsf{A},F)$, abbreviated to $\calg{A}$, where $\mathsf{A}=(\mathsf{A}_{s})_{s\in S}$ is an $S$-sorted single-set $n$-category and $F$ a structure of many-sorted $n$-categorial $\Sigma$-algebra on $\mathsf{A}$.  For a pair $(\mathbf{s},s)\in S^{\star}\times S$ and a formal operation $\sigma\in \Sigma_{\mathbf{s},s}$, in order to simplify the notation, the $n$-functor $F_{\mathbf{s},s}(\sigma)$ from $\mathsf{A}_{\mathbf{s}}$ to $\mathsf{A}_{s}$ will be written as $F_{\sigma}$, or simply as $\sigma^{\calg{A}}$. 

An \emph{$S$-sorted single-set $n$-categorial $\Sigma$-homomorphism} or, simply, a \emph{many-sorted $n$-categorial $\sig{\Sigma}$-ho\-mo\-mor\-phism}, from $\calg{A}$ to $\calg{B}$, where $\calg{B}=(\mathsf{B},G)$, is a triple $(\calg{A},F,\calg{B})$, abbreviated to $F\colon\calg{A}\mor\calg{B}$, where $F$ is an $S$-sorted single-set $n$-functor from $\mathsf{A}$ to $\mathsf{B}$ such that, for every $(\mathbf{s},s)\in S^{\star}\times S$, every $\sigma\in\Sigma_{\mathbf{s},s}$ and every family of cells $(\mathsf{a}_{j})_{j\in\bb{\mathbf{s}}}\in\mathsf{A}_{\mathbf{s}}$ (i.e., of $n$-cells, since $C_{n}=C$), we have that
$$
F_{s}\left(\sigma^{\calg{A}}
\left(\left(\mathsf{a}_{j}
\right)_{j\in\bb{\mathbf{s}}}\right)\right)
=
\sigma^{\calg{B}}
\left(\left(
F_{s_{j}}\left(
\mathsf{a}_{j}
\right)
\right)_{j\in\bb{\mathbf{s}}}
\right).
$$
We denote by $\mathsf{nCat}\mathsf{Alg}(\sig{\Sigma})$ the category of $n$-categorial  $\sig{\Sigma}$-algebras and $n$-categorial $\sig{\Sigma}$-homomorphisms.
\end{definition}

We next state the $\omega$-dimensional counterparts of these notions, which are used in Sections~\ref{STruncation} and~\ref{SLimitSingleSet}.

\begin{definition}\label{DSingleSetOmegaCat}
A \emph{single-set $\omega$-category} is an ordered pair $\mathsf{C} = (C,(\xi_{k})_{k\in\omega})$ where $C$ is a set and, for every $k\in\omega$, $\xi_{k}$ a structure of single-set category on $C$ such that, for every $j,k\in\omega$ with $j<k$, the ordered pair $(C,(\xi_{j},\xi_{k}))$ is a single-set $2$-category, and such that every $f\in C$ is a $k$-cell, i.e., $\src^{k}(f)=f$, for some $k\in\omega$. Note that, by~(2C1), a $j$-cell is a $k$-cell for every $k\geq j$.

A \emph{single-set $\omega$-functor} is a mapping $F\colon C\mor C'$ which is a single-set functor from $(C,\xi_{k})$ to $(C',\xi'_{k})$ for every $k\in\omega$. We denote by $\omega\mathsf{Cat}$ the category of single-set $\omega$-categories and single-set $\omega$-functors.
\end{definition}

We next introduce the $\omega$-dimensional counterpart of Definition~\ref{DManySortedCatAlg}.

\begin{definition}\label{DOmegaCatAlg}
Let $\sig{\Sigma}=(S,\Sigma)$ be a many-sorted signature. An \emph{$\omega$-categorial $\sig{\Sigma}$-algebra} is a pair $\calg{A}=(\mathsf{A},F)$ where $\mathsf{A}=(\mathsf{A}_{s})_{s\in S}$ is an $S$-sorted single-set $\omega$-category and $F$ assigns to every $\sigma\in\Sigma_{\mathbf{s},s}$ a single-set $\omega$-functor $\sigma^{\calg{A}}\colon\mathsf{A}_{\mathbf{s}}\mor\mathsf{A}_{s}$. Here $\mathsf{A}_{\mathbf{s}}$ is a single-set $\omega$-category: if each $\mathsf{a}_{j}$ is a $k_{j}$-cell, then $(\mathsf{a}_{j})_{j\in\bb{\mathbf{s}}}$ is a $k$-cell for $k=\max_{j}k_{j}$. An \emph{$\omega$-categorial $\sig{\Sigma}$-homomorphism} is an $S$-sorted single-set $\omega$-functor commuting with the operations, as in Definition~\ref{DManySortedCatAlg}. We denote by $\omega\mathsf{CatAlg}(\sig{\Sigma})$ the category of $\omega$-categorial $\sig{\Sigma}$-algebras and $\omega$-categorial $\sig{\Sigma}$-homomorphisms.
\end{definition}


In the single-set presentation, an $n$-categorial $\sig{\Sigma}$-algebra is regarded as a partial algebra on a single $S$-sorted set of cells, for a signature that enlarges $\Sigma$ with sources, targets and compositions of every dimension below $n$.

\begin{definition}\label{DSigmaCatSingleSet}
Let $\Sigma$ be an $S$-sorted signature and $n\in\mathbb{N}$. The \emph{single-set $n$-categorial signature determined by $\Sigma$}, denoted by $\Sigma^{(n)}$, is the $S$-sorted signature defined, for every $(\mathbf{s},s)\in S^{\star}\times S$, as follows:
$$
{\Sigma}^{(n)}_{\mathbf{s},s}=
\begin{cases}
{\Sigma}_{\mathbf{s},s},
& \text{if $\mathbf{s}\not\in \{(s),(s,s)\}$;}\\
{\Sigma}_{\mathbf{s},s}\amalg\{\src^{k}_{s}, \tgt^{k}_{s}\mid k\in n\},
& \text{if $\mathbf{s}= (s)$;}\\
{\Sigma}_{\mathbf{s},s}\amalg\{\circ^{k}_{s}\mid k\in n\},
& \text{if $\mathbf{s}= (s,s)$.}
\end{cases}
$$
That is, ${\Sigma}^{(n)}$ is the expansion of $\Sigma$ obtained by adding, for every $k\in n$ and every sort $s\in S$, two unary operation symbols $\src^{k}_{s}$ and $\tgt^{k}_{s}$, the \emph{$k$-source} and the \emph{$k$-target}, of arity $(s)$ and coarity $s$, and a binary operation symbol $\circ^{k}_{s}$, the \emph{$k$-composition}, of arity $(s,s)$ and coarity $s$. As in Definition~\ref{DSpecSingleSet}, we write $x\circ^{k}_{s}y$ for $\circ^{k}_{s}(x,y)$.

Moreover, we regard $\Sigma^{(n)}$ as the hep-signature, see Definition~\ref{DHepSignature}, $\sig{\Sigma}^{(n)}=(S,\Sigma^{(n)},\mathrm{def}^{(S,\Sigma^{(n)})})$ in which the $k$-compositions are the only operation symbols with a nonempty set of domain conditions: for every $k\in n$ and every $s\in S$,
$$
\mathrm{def}^{(S,\Sigma^{(n)})}(\circ^{k}_{s})=\{(\src^{k}_{s}(v^{s}_{0}),\tgt^{k}_{s}(v^{s}_{1}))\},
$$
where $v^{s}_{0}$ and $v^{s}_{1}$ are the variables of $\vs{(s,s)}$ that correspond to the first and to the second argument, respectively, so that $v^{s}_{0}\circ^{k}_{s}v^{s}_{1}$ is defined if and only if the $k$-source of $v^{s}_{0}$ is equal to the $k$-target of $v^{s}_{1}$; and $\mathrm{def}^{(S,\Sigma^{(n)})}(\sigma)=\varnothing$, $\mathrm{def}^{(S,\Sigma^{(n)})}(\src^{k}_{s})=\varnothing$ and $\mathrm{def}^{(S,\Sigma^{(n)})}(\tgt^{k}_{s})=\varnothing$, for every $\sigma\in\Sigma$, every $k\in n$ and every $s\in S$. The relation $\lhd$ of $\sig{\Sigma}^{(n)}$, see Definition~\ref{DHepSignature}, is well-founded, since the domain conditions of the compositions only involve sources and targets, which have empty sets of domain conditions; hence $\sig{\Sigma}^{(n)}$ is indeed a hep-signature.
\end{definition}

Note that $\Sigma^{(0)}=\Sigma$ and that, for every $n,m\in\mathbb{N}$ with $n\leq m$, $\Sigma^{(n)}\subseteq\Sigma^{(m)}$ and $\mathrm{def}^{(S,\Sigma^{(m)})}$ extends $\mathrm{def}^{(S,\Sigma^{(n)})}$.

In the following definition we introduce the $\omega$-dimensional counterpart of the single-set $n$-categorial signature.

\begin{definition}\label{DSigmaCatSingleSetOmega}
Let $\Sigma$ be an $S$-sorted signature. The \emph{single-set $\omega$-categorial signature determined by $\Sigma$} is the $S$-sorted signature $\Sigma^{(\omega)}=\bigcup_{n\in\mathbb{N}}\Sigma^{(n)}$, and $\sig{\Sigma}^{(\omega)}=(S,\Sigma^{(\omega)},\mathrm{def}^{(S,\Sigma^{(\omega)})})$ is the hep-signature in which $\mathrm{def}^{(S,\Sigma^{(\omega)})}$ is the union of the $\mathrm{def}^{(S,\Sigma^{(n)})}$, for $n\in\mathbb{N}$. Thus $\Sigma^{(\omega)}$ has the operation symbols $\src^{k}_{s}$, $\tgt^{k}_{s}$ and $\circ^{k}_{s}$ for every $k\in\omega$, with the same domain conditions as in Definition~\ref{DSigmaCatSingleSet}.
\end{definition}

We next present the $n$-categorial $\sig{\Sigma}$-algebras as a $\QE$-variety of partial $\Sigma^{(n)}$-algebras. The axioms (A$^{k}$1)--(A$^{k}$8), stated for every dimension $k\in n$, describe the $k$-dimensional structure, and the axioms (AB$^{l,k}$1)--(AB$^{l,k}$3), stated for every pair of dimensions $l<k$ in $n$, relate the $l$-dimensional and the $k$-dimensional structures. Informally, these equations state that the operations of $\Sigma$ are always defined, that, for every sort $s\in S$, sources, targets and compositions of sort $s$ of every dimension behave as in an $n$-category, and that the operations of $\Sigma$ are compatible with them.

\begin{definition}\label{DSpecSingleSet}
Let $\sig{\Sigma}=(S,\Sigma)$ be a many-sorted signature and $n\in\mathbb{N}\cup\{\omega\}$. We will denote by $(S,\Sigma^{(n)},V^{S},\mathcal{C}^{\sig{\Sigma}^{(n)}})$, written $\boldsymbol{\mathcal{C}}^{\sig{\Sigma}^{(n)}}$ for short, the specification in which $\Sigma^{(n)}$ is the single-set $n$-categorial signature introduced in Definition~\ref{DSigmaCatSingleSet}, $V^{S}$ a fixed $S$-sorted set with a countable infinity of variables in each coordinate, and $\mathcal{C}^{\sig{\Sigma}^{(n)}}$ the subset of $\mathrm{QEeqt}(\Sigma^{(n)})_{V^{S}}$, the set of all $\QE$-equations with variables in $V^{S}$, consisting of the following equations:

For every $(\mathbf{s},s)\in S^{\star}\times S$, every operation symbol $\sigma\in \Sigma_{\mathbf{s},s}$, and every family of variables $(x_{j})_{j\in \lvert\mathbf{s}\rvert}\in V^{S}_{\mathbf{s}}$, the operation $\sigma$ applied to the family $(x_{j})_{j\in \lvert\mathbf{s}\rvert}$ is always defined. Formally,
\allowdisplaybreaks
\begin{align*}
\sigma((x_{j})_{j\in \lvert\mathbf{s}\rvert})\eeq\sigma((x_{j})_{j\in \lvert\mathbf{s}\rvert}).\label{ASingleSetA0}\tag{A0}
\end{align*}

For every $k\in n$, the following equations describe the $k$-dimensional structure.

For every sort $s\in S$ and every variable $x\in V^{S}_{s}$, the $k$-source and $k$-target of $x$ are always defined. Formally,
\allowdisplaybreaks
\begin{align*}
\src^{k}_{s}(x)&\eeq\src^{k}_{s}(x);
&
\tgt^{k}_{s}(x)&\eeq\tgt^{k}_{s}(x).\label{ASingleSetAk1}\tag{A$^{k}$1}
\end{align*}

For every sort $s\in S$ and every variable $x\in V^{S}_{s}$, we have the following equations:
\allowdisplaybreaks
\begin{align*}
\src^{k}_{s}(\src^{k}_{s}(x))&\eeq\src^{k}_{s}(x);
&
\src^{k}_{s}(\tgt^{k}_{s}(x))&\eeq\tgt^{k}_{s}(x);\\
\tgt^{k}_{s}(\src^{k}_{s}(x))&\eeq\src^{k}_{s}(x);
&
\tgt^{k}_{s}(\tgt^{k}_{s}(x))&\eeq\tgt^{k}_{s}(x).\label{ASingleSetAk2}\tag{A$^{k}$2}
\end{align*}
In other words, $\src^{k}_{s}$ and $\tgt^{k}_{s}$ are right zeros. In particular, $\src^{k}_{s}$ and $\tgt^{k}_{s}$ are idempotent.

For every sort $s\in S$ and every pair of variables $x,y\in V^{S}_{s}$, $x\circ^{k}_{s} y$ is defined if and only if the $k$-target of $y$ is equal to the $k$-source of $x$. Formally,
\allowdisplaybreaks
\begin{align*}
x\circ^{k}_{s} y\eeq x\circ^{k}_{s} y \;\;
&\longrightarrow \;\;
\src^{k}_{s}(x)\eeq\tgt^{k}_{s}(y);\\
\src^{k}_{s}(x)\eeq\tgt^{k}_{s}(y) \;\;
&\longrightarrow \;\;
x\circ^{k}_{s} y\eeq x\circ^{k}_{s}y.\label{ASingleSetAk3}\tag{A$^{k}$3}
\end{align*}

For every sort $s\in S$ and every pair of variables $x,y\in V^{S}_{s}$, if $x\circ^{k}_{s} y$ is defined, then the $k$-source of $x\circ^{k}_{s}y$ is that of $y$ and the $k$-target of $x\circ^{k}_{s}y$ is that of $x$. Formally,
\allowdisplaybreaks
\begin{align*}
x\circ^{k}_{s}y\eeq x\circ^{k}_{s}y \;\;
&\longrightarrow \;\;
\src^{k}_{s}(x\circ^{k}_{s}y)\eeq\src^{k}_{s}(y);\\
x\circ^{k}_{s}y\eeq x\circ^{k}_{s}y \;\;
&\longrightarrow \;\;
\tgt^{k}_{s}(x\circ^{k}_{s}y)\eeq\tgt^{k}_{s}(x).\label{ASingleSetAk4}\tag{A$^{k}$4}
\end{align*}

For every sort $s\in S$ and every variable $x\in V^{S}_{s}$, the compositions $x\circ^{k}_{s}\src^{k}_{s}(x)$ and $\tgt^{k}_{s}(x)\circ^{k}_{s}x$ are always defined and are equal to $x$, i.e., $\src^{k}_{s}(x)$ is a right unit element for the $k$-composition with $x$ and $\tgt^{k}_{s}(x)$ is a left unit element for the $k$-composition with $x$. Formally,
\allowdisplaybreaks
\begin{align*}
x\circ^{k}_{s}\src^{k}_{s}(x)&\eeq x;
&
\tgt^{k}_{s}(x)\circ^{k}_{s}x&\eeq x.\label{ASingleSetAk5}\tag{A$^{k}$5}
\end{align*}

For every sort $s\in S$ and every triple of variables $x, y, z\in V^{S}_{s}$, if the $k$-compositions $x\circ^{k}_{s}y$ and $y\circ^{k}_{s}z$ are defined, then the $k$-compositions $x\circ^{k}_{s}(y\circ^{k}_{s}z)$ and $(x\circ^{k}_{s}y)\circ^{k}_{s}z$ are defined and they are equal, i.e., the $k$-composition, when defined, is associative. Formally,
\allowdisplaybreaks
\begin{align*}
(x\circ^{k}_{s}y\eeq x\circ^{k}_{s}y)
\wedge
(y\circ^{k}_{s}z\eeq y\circ^{k}_{s}z)\;\;
&\longrightarrow \;\;
(x\circ^{k}_{s}y)\circ^{k}_{s}z\eeq x\circ^{k}_{s}(y\circ^{k}_{s}z).\label{ASingleSetAk6}\tag{A$^{k}$6}
\end{align*}

For every $(\mathbf{s},s)\in S^{\star}\times S$, every operation symbol $\sigma\in \Sigma_{\mathbf{s},s}$, and every family of variables $(x_{j})_{j\in \lvert\mathbf{s}\rvert}\in V^{S}_{\mathbf{s}}$, the $k$-source of $\sigma((x_{j})_{j\in \lvert\mathbf{s}\rvert})$ is equal to $\sigma$ applied to the family $((\src^{k}_{s_{j}}(x_{j}))_{j\in \lvert\mathbf{s}\rvert})$, and the $k$-target of $\sigma((x_{j})_{j\in \lvert\mathbf{s}\rvert})$ is equal to $\sigma$ applied to the family $((\tgt^{k}_{s_{j}}(x_{j}))_{j\in \lvert\mathbf{s}\rvert})$. Formally,
\allowdisplaybreaks
\begin{align*}
\src^{k}_{s}(\sigma((x_{j})_{j\in \lvert\mathbf{s}\rvert}))&\eeq\sigma((\src^{k}_{s_{j}}(x_{j}))_{j\in \lvert\mathbf{s}\rvert});\\
\tgt^{k}_{s}(\sigma((x_{j})_{j\in \lvert\mathbf{s}\rvert}))&\eeq\sigma((\tgt^{k}_{s_{j}}(x_{j}))_{j\in \lvert\mathbf{s}\rvert}).\label{ASingleSetAk7}\tag{A$^{k}$7}
\end{align*}

For every $(\mathbf{s},s)\in S^{\star}\times S$, every operation symbol $\sigma\in \Sigma_{\mathbf{s},s}$, and every pair of families of variables $(x_{j})_{j\in \lvert\mathbf{s}\rvert}, (y_{j})_{j\in \lvert\mathbf{s}\rvert}\in V^{S}_{\mathbf{s}}$, if, for every $j\in \lvert\mathbf{s}\rvert$, the $k$-compositions $x_{j}\circ^{k}_{s_{j}}y_{j}$ are defined, then the $k$-composition $\sigma((x_{j})_{j\in \lvert\mathbf{s}\rvert})\circ^{k}_{s}\sigma((y_{j})_{j\in \lvert\mathbf{s}\rvert})$ is defined and it is equal to $\sigma$ applied to the family $(x_{j}\circ^{k}_{s_{j}}y_{j})_{j\in \lvert\mathbf{s}\rvert}$. Formally,
\allowdisplaybreaks
\begin{multline*}
\textstyle
\bigwedge_{j\in \lvert\mathbf{s}\rvert}
(x_{j}\circ^{k}_{s_{j}} y_{j}\eeq x_{j}\circ^{k}_{s_{j}} y_{j})
\;\;\longrightarrow\;\;\\
\sigma((x_{j}\circ^{k}_{s_{j}}y_{j})_{j\in \lvert\mathbf{s}\rvert})
\eeq
\sigma((x_{j})_{j\in \lvert\mathbf{s}\rvert})\circ^{k}_{s}\sigma((y_{j})_{j\in \lvert\mathbf{s}\rvert}).
\label{ASingleSetAk8}\tag{A$^{k}$8}
\end{multline*}

For every $l,k\in n$ with $l<k$, the following equations relate the $l$-dimensional and the $k$-dimensional structures.

For every sort $s\in S$ and every variable $x\in V^{S}_{s}$, the elements $\src^{k}_{s}(\src^{l}_{s}(x))$, $\src^{l}_{s}(\src^{k}_{s}(x))$ and $\src^{l}_{s}(\tgt^{k}_{s}(x))$ are always defined and are equal to $\src^{l}_{s}(x)$. Similarly, the elements $\tgt^{k}_{s}(\tgt^{l}_{s}(x))$, $\tgt^{l}_{s}(\tgt^{k}_{s}(x))$, and $\tgt^{l}_{s}(\src^{k}_{s}(x))$ are always defined and are equal to $\tgt^{l}_{s}(x)$. Formally,
\allowdisplaybreaks
\begin{align*}
\src^{k}_{s}(\src^{l}_{s}(x))&\eeq\src^{l}_{s}(x);
&
\tgt^{k}_{s}(\tgt^{l}_{s}(x))&\eeq\tgt^{l}_{s}(x);\\
\src^{l}_{s}(\src^{k}_{s}(x))&\eeq\src^{l}_{s}(x);
&
\tgt^{l}_{s}(\tgt^{k}_{s}(x))&\eeq\tgt^{l}_{s}(x);\\
\src^{l}_{s}(\tgt^{k}_{s}(x))&\eeq\src^{l}_{s}(x);
&
\tgt^{l}_{s}(\src^{k}_{s}(x))&\eeq\tgt^{l}_{s}(x).\label{ASingleSetABlk1}\tag{AB$^{l,k}$1}
\end{align*}

For every sort $s\in S$ and every pair of variables $x,y\in V^{S}_{s}$, if $x\circ^{l}_{s}y$ is defined, then the $k$-source of $x\circ^{l}_{s}y$ is the $l$-composition of the $k$-source of $x$ with the $k$-source of $y$ and the $k$-target of $x\circ^{l}_{s}y$ is the $l$-composition of the $k$-target of $x$ with the $k$-target of $y$. Formally,
\allowdisplaybreaks
\begin{align*}
x\circ^{l}_{s}y&\eeq x\circ^{l}_{s}y \;\;\longrightarrow\;\;
\src^{k}_{s}(x\circ^{l}_{s}y)\eeq\src^{k}_{s}(x)\circ^{l}_{s}\src^{k}_{s}(y);\\
x\circ^{l}_{s}y&\eeq x\circ^{l}_{s}y \;\;\longrightarrow\;\;
\tgt^{k}_{s}(x\circ^{l}_{s}y)\eeq\tgt^{k}_{s}(x)\circ^{l}_{s}\tgt^{k}_{s}(y).
\label{ASingleSetABlk2}\tag{AB$^{l,k}$2}
\end{align*}

For every sort $s\in S$ and every quadruple of variables $x,y,z,t\in V^{S}_{s}$, if the $l$-compositions $x\circ^{l}_{s}y$ and $z\circ^{l}_{s}t$ are defined and the $k$-compositions $x\circ^{k}_{s}z$ and $y\circ^{k}_{s}t$ are defined, the $k$-composition of $x\circ^{l}_{s}y$ with $z\circ^{l}_{s}t$ is equal to the $l$-composition of $x\circ^{k}_{s}z$ with $y\circ^{k}_{s}t$. Formally,
\allowdisplaybreaks
\begin{multline*}
(x\circ^{l}_{s}y\eeq x\circ^{l}_{s}y)
\wedge
(z\circ^{l}_{s}t\eeq z\circ^{l}_{s}t)
\wedge
(x\circ^{k}_{s}z\eeq x\circ^{k}_{s}z)
\wedge
(y\circ^{k}_{s}t\eeq y\circ^{k}_{s}t)
\\
\;\;\longrightarrow\;\;
(x\circ^{l}_{s}y)\circ^{k}_{s}(z\circ^{l}_{s}t)\eeq
(x\circ^{k}_{s}z)\circ^{l}_{s}(y\circ^{k}_{s}t).
\label{ASingleSetABlk3}\tag{AB$^{l,k}$3}
\end{multline*}

For $n\in\mathbb{N}$, a model of the axioms (A$^{k}$1)--(A$^{k}$6), for $k\in n$, and (AB$^{l,k}$1)--(AB$^{l,k}$3), for $l<k$ in $n$, is an $S$-sorted single-set $n$-category, see Definition~\ref{DManySortedCatAlg}, presented with all the cells of a given sort, whatever their dimension, in a single set; for $n=\omega$, see Remark~\ref{RSingleSetIdentOmega}.

Since $\sig{\Sigma}^{(n)}$ is a hep-signature, see Definitions~\ref{DSigmaCatSingleSet} and~\ref{DSigmaCatSingleSetOmega}, $(\sig{\Sigma}^{(n)},\mathcal{C}^{\sig{\Sigma}^{(n)}})$ is a hep-specification, see Definition~\ref{DEpSpec}, and we will work with the ep-variety $\mathcal{V}(\sig{\Sigma}^{(n)},\mathcal{C}^{\sig{\Sigma}^{(n)}})$ and with its canonically associated category $\mathsf{EPAlg}(\sig{\Sigma}^{(n)},\mathcal{C}^{\sig{\Sigma}^{(n)}})$, see Definition~\ref{DEpVariety}.
\end{definition}

In the following proposition we characterize the ep-algebras of type $\sig{\Sigma}^{(n)}$ in terms of the axioms of Definition~\ref{DSpecSingleSet}.

\begin{proposition}\label{PSingleSetEp}
Let $\sig{\Sigma}=(S,\Sigma)$ be a many-sorted signature, $n\in\mathbb{N}\cup\{\omega\}$ and $\mathbf{A}$ a partial $\Sigma^{(n)}$-algebra. Then $\mathbf{A}$ is a $\sig{\Sigma}^{(n)}$-algebra, see Definitions~\ref{DEpAlgebra} and~\ref{DSigmaCatSingleSet}, if, and only if, $\mathbf{A}$ satisfies the axioms (A0), (A$^{k}$1) and (A$^{k}$3), for every $k\in n$.
\end{proposition}
\begin{proof}
By Definition~\ref{DEpAlgebra}, $\mathbf{A}$ is a $\sig{\Sigma}^{(n)}$-algebra if, and only if, for every operation symbol $\tau$ of $\Sigma^{(n)}$, the domain of $\tau^{\mathbf{A}}$ is the set of the families that satisfy the domain conditions of $\tau$ in $\mathbf{A}$. We consider the different kinds of operation symbols of $\Sigma^{(n)}$, see Definition~\ref{DSigmaCatSingleSet}.

\textsf{(Operation symbols of $\Sigma$)} For every $(\mathbf{s},s)\in S^{\star}\times S$ and every $\sigma\in\Sigma_{\mathbf{s},s}$, the set $\mathrm{def}^{(S,\Sigma^{(n)})}(\sigma)$ is empty, so that $A_{\mathrm{def}^{(S,\Sigma^{(n)})}(\sigma)}=A_{\mathbf{s}}$. Therefore, the condition of Definition~\ref{DEpAlgebra} for $\sigma$ states that $\sigma^{\mathbf{A}}$ is total, i.e., it is the axiom (A0).

\textsf{(Sources and targets)} For every $k\in n$ and every $s\in S$, the sets $\mathrm{def}^{(S,\Sigma^{(n)})}(\src^{k}_{s})$ and $\mathrm{def}^{(S,\Sigma^{(n)})}(\tgt^{k}_{s})$ are empty. Therefore, the conditions of Definition~\ref{DEpAlgebra} for $\src^{k}_{s}$ and $\tgt^{k}_{s}$ state that $\src^{k\mathbf{A}}_{s}$ and $\tgt^{k\mathbf{A}}_{s}$ are total, i.e., they are the axiom (A$^{k}$1).

\textsf{(Compositions)} For every $k\in n$ and every $s\in S$, assuming that $\src^{k\mathbf{A}}_{s}$ and $\tgt^{k\mathbf{A}}_{s}$ are total, the following equality holds
$$
A_{\mathrm{def}^{(S,\Sigma^{(n)})}(\circ^{k}_{s})}
=
\left\lbrace(x,y)\in A_{s}\times A_{s}\,\middle|\,\src^{k\mathbf{A}}_{s}(x)=\tgt^{k\mathbf{A}}_{s}(y)\right\rbrace.
$$
Therefore, the condition of Definition~\ref{DEpAlgebra} for $\circ^{k}_{s}$ states that $x\circ^{k}_{s}y$ is defined if, and only if, $\src^{k}_{s}(x)=\tgt^{k}_{s}(y)$, i.e., it is the axiom (A$^{k}$3).

All in all, $\mathbf{A}$ is a $\sig{\Sigma}^{(n)}$-algebra if, and only if, it satisfies (A0), (A$^{k}$1) and (A$^{k}$3), for every $k\in n$.

This completes the proof.
\end{proof}

By Proposition~\ref{PSingleSetEp}, a partial $\Sigma^{(n)}$-algebra in which every axiom of $\mathcal{C}^{\sig{\Sigma}^{(n)}}$ holds is automatically a $\sig{\Sigma}^{(n)}$-algebra, since it satisfies, in particular, the axioms (A0), (A$^{k}$1) and (A$^{k}$3), for every $k\in n$. Hence the objects of $\mathsf{EPAlg}(\sig{\Sigma}^{(n)},\mathcal{C}^{\sig{\Sigma}^{(n)}})$ are exactly the models of the $\QE$-specification $\boldsymbol{\mathcal{C}}^{\sig{\Sigma}^{(n)}}$.

We next relate the many-sorted $n$-categorial $\sig{\Sigma}$-algebras with the $\sig{\Sigma}^{(n)}$-algebras. Let $\sig{\Sigma}=(S,\Sigma)$ be a many-sorted signature and $n\in\mathbb{N}\cup\{\omega\}$. In what follows, a pair $(\mathsf{A},F)$, where $\mathsf{A}=(\mathsf{A}_{s})_{s\in S}$ is an $S$-sorted single-set $n$-category, with $\mathsf{A}_{s}=(A_{s},(\circ^{k}_{s},\src^{k}_{s},\tgt^{k}_{s})_{k\in n})$, and $F$ assigns to every $\sigma\in\Sigma_{\mathbf{s},s}$ a mapping $F_{\sigma}\colon A_{\mathbf{s}}\mor A_{s}$, is identified with the partial $\Sigma^{(n)}$-algebra with underlying $S$-sorted set $A=(A_{s})_{s\in S}$ in which, for every $k\in n$ and every $s\in S$, the operation symbols $\src^{k}_{s}$, $\tgt^{k}_{s}$ and $\circ^{k}_{s}$ are interpreted as $\src^{k}_{s}$, $\tgt^{k}_{s}$ and $\circ^{k}_{s}$, respectively, and every $\sigma\in\Sigma$ is interpreted as $F_{\sigma}$.

\begin{remark}\label{RSingleSetIdent}
Let $\sig{\Sigma}=(S,\Sigma)$ be a many-sorted signature and $n\in\mathbb{N}$. Then, under the above identification,
$$
\mathsf{nCatAlg}(\sig{\Sigma})=\mathsf{EPAlg}(\sig{\Sigma}^{(n)},\mathcal{C}^{\sig{\Sigma}^{(n)}}).
$$
In fact, for every $k\in n$, the conditions (C1)--(C5) of Definition~\ref{DSingleSetCat} for the $k$-th structure of category are the axioms (A$^{k}$2)--(A$^{k}$6), the totality of $\src^{k}_{s}$ and $\tgt^{k}_{s}$ being the axiom (A$^{k}$1); for every $l<k$ in $n$, the conditions (2C1)--(2C3) of Definition~\ref{DSingleSet2Cat} are the axioms (AB$^{l,k}$1)--(AB$^{l,k}$3); and, since the products $\mathsf{A}_{\mathbf{s}}$ are taken componentwise, the operations of $\Sigma$ are total single-set $n$-functors if, and only if, the axioms (A0), (A$^{k}$7) and (A$^{k}$8) hold, for every $k\in n$. Moreover, the $S$-sorted single-set $n$-functors that commute with the operations of $\Sigma$, i.e., the $n$-categorial $\sig{\Sigma}$-homomorphisms, are exactly the homomorphisms of partial $\Sigma^{(n)}$-algebras, see Definition~\ref{DPartialAlg}, since the latter preserve the defined $k$-compositions. Finally, by Proposition~\ref{PSingleSetEp}, the models of $\boldsymbol{\mathcal{C}}^{\sig{\Sigma}^{(n)}}$ are $\sig{\Sigma}^{(n)}$-algebras.
\end{remark}

In dimension $\omega$ the identification works in the same way, except for the condition that every cell has a finite dimension.

\begin{remark}\label{RSingleSetIdentOmega}
Let $\sig{\Sigma}=(S,\Sigma)$ be a many-sorted signature. Then, under the above identification, and by the same arguments as in Remark~\ref{RSingleSetIdent}, applied for every $k\in\omega$ and every $l<k$ in $\omega$, $\omega\mathsf{CatAlg}(\sig{\Sigma})$ is the full subcategory of $\mathsf{EPAlg}(\sig{\Sigma}^{(\omega)},\mathcal{C}^{\sig{\Sigma}^{(\omega)}})$ whose objects are the $\sig{\Sigma}^{(\omega)}$-algebras in which every cell is a $k$-cell, i.e., satisfies $\src^{k}_{s}(x)=x$, for some $k\in\omega$, see Definitions~\ref{DSingleSetOmegaCat} and~\ref{DOmegaCatAlg}.
\end{remark}

One could try to add to $\mathcal{C}^{\sig{\Sigma}^{(\omega)}}$ some axioms expressing that every cell has a finite dimension, in order to present $\omega\mathsf{CatAlg}(\sig{\Sigma})$ as the category canonically associated to an ep-variety, as in the finite-dimensional case. In the following proposition we prove that this is impossible.

\begin{proposition}\label{PFiniteDimNotQE}
Let $\sig{\Sigma}=(S,\Sigma)$ be a many-sorted signature with $S\neq\varnothing$. Then there is no set $\mathcal{F}$ of $\QE$-equations for $\Sigma^{(\omega)}$ such that the objects of $\omega\mathsf{CatAlg}(\sig{\Sigma})$, under the above identification, are exactly the partial $\Sigma^{(\omega)}$-algebras in which every $\QE$-equation of $\mathcal{F}$ holds. In particular, $\omega\mathsf{CatAlg}(\sig{\Sigma})$ is a proper full subcategory of $\mathsf{EPAlg}(\sig{\Sigma}^{(\omega)},\mathcal{C}^{\sig{\Sigma}^{(\omega)}})$.
\end{proposition}
\begin{proof}
Let us assume, towards a contradiction, that such a set $\mathcal{F}$ exists. Then, by the note after Definition~\ref{DQEEq}, the class of the objects of $\omega\mathsf{CatAlg}(\sig{\Sigma})$ is closed under products of partial $\Sigma^{(\omega)}$-algebras, see Definition~\ref{DPAlgProd}. We next construct a family $(\calg{A}_{k})_{k\in\omega}$ of $\omega$-categorial $\sig{\Sigma}$-algebras whose product is not an $\omega$-categorial $\sig{\Sigma}$-algebra.

Let $s_{0}\in S$. For every $k\in\omega$, let $\mathsf{G}_{k}$ be a single-set $\omega$-category with a cell $g_{k}$ which is a $(k+1)$-cell but not a $k$-cell and such that every cell of $\mathsf{G}_{k}$ is a $(k+1)$-cell; for instance, the single-set $\omega$-category freely generated by a $(k+1)$-dimensional globe, whose cells are $g_{k}$ and its iterated $j$-sources and $j$-targets, for $j\leq k$. Let $\calg{A}_{k}$ be the $\omega$-categorial $\sig{\Sigma}$-algebra in which, for every $s\in S$, the single-set $\omega$-category of sort $s$ is $\mathsf{G}_{k}$, and in which every operation symbol $\sigma\in\Sigma_{\mathbf{s},s}$ is interpreted as the constant mapping with value the $0$-cell $u_{k}=\src^{0}(g_{k})$. Every such constant mapping is a single-set $\omega$-functor, since $u_{k}$ is a $j$-cell for every $j\in\omega$, so that $\src^{j}(u_{k})=u_{k}=\tgt^{j}(u_{k})$ and, by~(C4), $u_{k}\circ^{j}u_{k}=u_{k}$. Moreover, every cell of $\calg{A}_{k}$ is a $(k+1)$-cell, so that $\calg{A}_{k}$ is an object of $\omega\mathsf{CatAlg}(\sig{\Sigma})$, see Remark~\ref{RSingleSetIdentOmega}.

Let $\mathbf{A}=\prod_{k\in\omega}\calg{A}_{k}$ be the product of the corresponding partial $\Sigma^{(\omega)}$-algebras and $g=(g_{k})_{k\in\omega}$, which is an element of $A_{s_{0}}$. By Definition~\ref{DPAlgProd}, for every $j\in\omega$, $\src^{j\,\mathbf{A}}_{s_{0}}(g)=(\src^{j}(g_{k}))_{k\in\omega}$, which is different from $g$, since $\src^{j}(g_{j})\neq g_{j}$. Hence $g$ is not a $j$-cell of $\mathbf{A}$ for any $j\in\omega$, so that $\mathbf{A}$ is not an object of $\omega\mathsf{CatAlg}(\sig{\Sigma})$, a contradiction.

Finally, since $\mathcal{C}^{\sig{\Sigma}^{(\omega)}}$ is a set of $\QE$-equations, by the note after Definition~\ref{DQEEq}, $\mathbf{A}$ is an object of $\mathsf{EPAlg}(\sig{\Sigma}^{(\omega)},\mathcal{C}^{\sig{\Sigma}^{(\omega)}})$ which is not an object of $\omega\mathsf{CatAlg}(\sig{\Sigma})$.

This completes the proof.
\end{proof}

Let us note that the same argument shows that the condition that every cell has a finite dimension cannot be expressed by any set of axioms whose class of models is closed under products. Therefore, the presentation of the $\omega$-categorial $\sig{\Sigma}$-algebras cannot be unified with that of the $n$-categorial $\sig{\Sigma}$-algebras, and in Sections~\ref{SLimitSingleSet} and~\ref{SLimitGlobular} we present them, instead, by means of the projective limits of the finite-dimensional ones.


%The axiom (A0) states that the operations of $\Sigma$ are total, and the axioms (A$^{k}$7) and (A$^{k}$8), for $k\in n$, state that they preserve sources, targets and compositions of every dimension, i.e., that they are $n$-functors. As in Chapters~2 and~3, the morphisms of $\mathsf{PAlg}(\boldsymbol{\mathcal{C}}^{\sig{\Sigma}^{(n)}})$ are the homomorphisms of partial $\Sigma^{(n)}$-algebras, i.e., the $S$-sorted mappings $h$ such that, whenever an operation is defined on a family of elements, it is also defined on the image of that family under $h$, and $h$ sends the result of the former to the result of the latter. These are the appropriate morphisms here: closed homomorphisms, which also reflect definedness, would be too restrictive, since an $n$-functor may send cells that are not composable to cells that are. Consequently, we have the following result.
%
%\begin{proposition}\label{PCompactVar}
%Let $\sig{\Sigma}$ be a many-sorted signature and $n\in\mathbb{N}$. Consider the functor from $\mathsf{nCatAlg}(\sig{\Sigma})$ to $\mathsf{PAlg}(\boldsymbol{\mathcal{C}}^{\sig{\Sigma}^{(n)}})$ that sends an $n$-categorial $\sig{\Sigma}$-algebra $\calg{A}$ to the partial $\Sigma^{(n)}$-algebra whose underlying set of sort $s$ is the set of $n$-cells of $\mathsf{A}_{s}$, in which every $k$-cell, for $k\in n$, is identified with its identity $n$-cell, and whose operations are the operations of $\calg{A}$ and the $k$-sources, $k$-targets and $k$-compositions of the $n$-categories $\mathsf{A}_{s}$; and that sends an $n$-categorial $\sig{\Sigma}$-homomorphism to its action on $n$-cells. Then this functor is an equivalence of categories, and it is an isomorphism of categories if $n$-categories are taken in their single-set presentation.
%\end{proposition}
%
%The proof reduces to the classical correspondence between the globular and the single-set presentations of $n$-categories, see, e.g.,~\cite{Str87,abgmmm23}, together with two observations: under this correspondence, the $n$-functors from $\mathsf{A}_{\mathbf{s}}$ to $\mathsf{A}_{s}$ are exactly the total operations that satisfy (A$^{k}$7) and (A$^{k}$8), for $k\in n$; and the $S$-sorted $n$-functors between the underlying $S$-sorted $n$-categories are exactly the homomorphisms of partial $\Sigma^{(n)}$-algebras defined above.

\section{Globular presentation of higher-order categorial algebras}\label{SGlobular}

In this section, we present the higher-order categorial algebras in their globular form, in which the cells of every dimension form a separate sort. We follow the same pattern as in Section~\ref{SSingleSet}: after recalling the globular $\omega$-categories and the globular $n$-categories, we introduce the globular $n$-categorial algebras, the globular $n$-categorial signature determined by a many-sorted signature and the $\QE$-variety of its models. We end this section by showing that the globular $n$-categorial algebras are exactly the models of this $\QE$-variety, which is an ep-variety.

We begin by defining globular $\omega$-categories, which present $\omega$-categories with separate sorts for each dimension, together with sources, targets and identities between dimensions, as in~\cite{abgmmm23}. A globular $n$-category is then a globular $\omega$-category truncated at level $n$.

\begin{definition}\label{DOmegaCatGl}
An \emph{globular $\omega$-category} is an ordered pair $\mathsf{C} = ((C_{k})_{k\in\omega},\zeta)$ consisting of an $\omega$-sorted set $(C_{k})_{k\in\omega}$ and a $\coprod_{k\in \omega}k$-indexed family
\[
\zeta = (\zeta_{j,k})_{(j,k)\in \coprod_{k\in \omega}k},
\] 
where, for every $k\in \omega$ and every $j\in k$, $\zeta_{j,k} = (\circ^{(j,k)},\src^{(j,k)},\tgt^{(j,k)},\mathrm{i}^{(k,j)})$ with
\begin{enumerate}
\item $\src^{(j,k)}$ and $\tgt^{(j,k)}$ mappings from $C_{k}$ to $C_{j}$  
which send $f\in C_{k}$ to its \emph{$(j,k)$-source} and \emph{$(j,k)$-target}, respectively;
\item $\mathrm{i}^{(k,j)}$ a mapping from $C_{j}$ to $C_{k}$ 
which sends $f\in C_{j}$ to its \emph{$(k,j)$-identity}; and
\item $\circ^{(j,k)}$ a partial binary operation 
$$
\circ^{(j,k)}\colon C_{k}\times C_{k} \dmor C_{k},
$$ 
whose domain of definition is
$$
\Dom\left(\circ^{(j,k)}\right) = 
\left\lbrace
(f,g)\in C_{k}\times C_{k}
\,\middle|\, \src^{(j,k)}(g)=\tgt^{(j,k)}(f) \right\rbrace,
$$ 
i.e., the fibered product of $C_{k}$ and $C_{k}$ over $\src^{(j,k)}$ and $\tgt^{(j,k)}$, and
which sends a pair $(f,g)\in \Dom(\circ^{(j,k)})$ to its \emph{$j$-composition of $k$-cells}.
\end{enumerate}
These data are required to satisfy the following conditions for all $i,j,k\in \omega$ with $i<j<k$ (the conditions (MS2), (MS6) and (MS7) only involve the dimensions $j<k$, and are required for every such pair):
\begin{itemize}
\item[(MS1)] For every $f\in C_{k}$,
\begin{align*}
\src^{(i,j)}(\src^{(j,k)}(f))&=\src^{(i,k)}(f)=\src^{(i,j)}(\tgt^{(j,k)}(f)),\\
\tgt^{(i,j)}(\tgt^{(j,k)}(f))&=\tgt^{(i,k)}(f)=\tgt^{(i,j)}(\src^{(j,k)}(f)).
\end{align*}
\item[(MS2)] If $g,f\in C_{k}$ are such that $\src^{(j,k)}(g)=\tgt^{(j,k)}(f)$, then
\begin{align*}
\src^{(j,k)}(g\circ^{(j,k)}f)&=\src^{(j,k)}(f), &
\tgt^{(j,k)}(g\circ^{(j,k)}f)&=\tgt^{(j,k)}(g).
\end{align*}
\item[(MS3)] If $g,f\in C_{k}$ are such that $\src^{(i,k)}(g)=\tgt^{(i,k)}(f)$, then
\begin{align*}
\src^{(j,k)}(g\circ^{(i,k)}f)&=\src^{(j,k)}(g)\circ^{(i,j)}\src^{(j,k)}(f),\\
\tgt^{(j,k)}(g\circ^{(i,k)}f)&=\tgt^{(j,k)}(g)\circ^{(i,j)}\tgt^{(j,k)}(f).
\end{align*}
\item[(MS4)] If $f\in C_{i}$, then
$$
\mathrm{i}^{(k,j)}(\mathrm{i}^{(j,i)}(f))=\mathrm{i}^{(k,i)}(f).
$$
\item[(MS5)] If $f\in C_{i}$, then
$$
\src^{(i,j)}(\mathrm{i}^{(j,i)}(f))=f=
\tgt^{(i,j)}(\mathrm{i}^{(j,i)}(f)).
$$
\item[(MS6)] If $f\in C_{k}$, then
$$
f\circ^{(j,k)}
(\mathrm{i}^{(k,j)}(\src^{(j,k)}(f)))
=
f
=
(\mathrm{i}^{(k,j)}(\tgt^{(j,k)}(f)))
\circ^{(j,k)}
f.
$$
\item[(MS7)] If $h,g,f\in C_{k}$ are such that 
\begin{align*}
\src^{(j,k)}(h)&=\tgt^{(j,k)}(g); &
\src^{(j,k)}(g)&=\tgt^{(j,k)}(f),
\end{align*}
then
$$
h\circ^{(j,k)}(g\circ^{(j,k)}f)=
(h\circ^{(j,k)}g)\circ^{(j,k)}f.
$$
\item[(MS8)] If $g,f\in C_{j}$ are such that $\src^{(i,j)}(g)=\tgt^{(i,j)}(f)$, then
$$
\mathrm{i}^{(k,j)}(g\circ^{(i,j)}f)=\mathrm{i}^{(k,j)}(g)\circ^{(i,k)}\mathrm{i}^{(k,j)}(f).
$$
\item[(MS9)] If $g',g,f',f\in C_{k}$ are such that
\begin{align*}
\src^{(j,k)}(g')&=\tgt^{(j,k)}(g); &
\src^{(j,k)}(f')&=\tgt^{(j,k)}(f);\\
\src^{(i,k)}(g')&=\tgt^{(i,k)}(f'); &
\src^{(i,k)}(g)&=\tgt^{(i,k)}(f).
\end{align*}
then,
$$
(g'\circ^{(i,k)}f')\circ^{(j,k)}(g\circ^{(i,k)} f)
=
(g'\circ^{(j,k)}g)\circ^{(i,k)}(f'\circ^{(j,k)}f)
.
$$
\end{itemize}

An \emph{globular $\omega$-functor} $F$ from $\mathsf{C}$ to $\mathsf{C}'$ is a family $F=(F_{k})_{k\in\omega}$ of mappings $F_{k}\colon C_{k}\mor C'_{k}$ such that, for all $j,k\in\omega$ with $j<k$,
\begin{align*}
F_{j}\circ \src^{(j,k)} &= \src'^{(j,k)}\circ F_{k}, &
F_{j}\circ \tgt^{(j,k)} &= \tgt'^{(j,k)}\circ F_{k},\\
F_{k}\circ \mathrm{i}^{(k,j)} &= \mathrm{i}'^{(k,j)}\circ F_{j}, &
F_{k}(g\circ^{(j,k)}f) &= F_{k}(g)\circ'^{(j,k)} F_{k}(f),
\end{align*}
the last one whenever $g\circ^{(j,k)}f$ is defined. We denote by $\omega\mathsf{Cat}^{\gl}$ the category of globular $\omega$-categories and globular $\omega$-functors.
\end{definition}

For $n\in\omega$, an \emph{globular $n$-category} is an ordered pair $\mathsf{C} = ((C_{k})_{k\in n+1},\zeta)$ consisting of an $n+1$-sorted set $(C_{k})_{k\in n+1}$ and a structure of globular $n$-category 
\[
\zeta = (\zeta_{j,k})_{(j,k)\in \coprod_{k\in n+1}k}
\] 
on $(C_{k})_{k\in n+1}$, where, for every $k\in n+1$ and every $j\in k$, $\zeta_{j,k} = (\circ^{(j,k)},\src^{(j,k)},\tgt^{(j,k)},\mathrm{i}^{(k,j)})$ with the same data as above, required to satisfy the same conditions truncated to the $n$-th level.

The notion of globular $n$-functor between $n$-categories is entirely analogous to that of globular $\omega$-functor between $\omega$-categories, but truncated to the $n$-th level. A globular $n$-functor $F$ from $\mathsf{C}$ to $\mathsf{C}'$ is a family $F=(F_{k})_{k\in n+1}$ of mappings $F_{k}\colon C_{k}\to C'_{k}$ such that for all $j,k\in n+1$ with $j<k$, the following equations hold (the last one whenever $g\circ^{(j,k)}f$ is defined):
\begin{align*}
F_{j}\circ \src^{(j,k)} &= \src'^{(j,k)}\circ F_{k}, &
F_{j}\circ \tgt^{(j,k)} &= \tgt'^{(j,k)}\circ F_{k},\\
F_{k}\circ \mathrm{i}^{(k,j)} &= \mathrm{i}'^{(k,j)}\circ F_{j}, &
F_{k}(g\circ^{(j,k)}f) &= F_{k}(g)\circ'^{(j,k)} F_{k}(f).
\end{align*}
We denote by $\mathsf{nCat}^{\gl}$ the category consisting of globular $n$-categories and globular $n$-functors.


We next define the $S$-sorted globular $n$-categories and the $S$-sorted globular $n$-functors.

\begin{definition}\label{DSortednCatGl}
Let $S$ be a set of sorts and $n\in\mathbb{N}$.
We will call the objects of the functor category $(\mathsf{nCat}^{\gl})^{S}$ \emph{$S$-sorted globular $n$-categories}. Thus, an $S$-sorted globular $n$-category is an $S$-indexed family $\mathsf{A}=(\mathsf{A}_{s})_{s\in S}$ such that, for every sort $s\in S$, $\mathsf{A}_{s}$ is a globular $n$-category. Given two $S$-sorted globular $n$-categories 
$\mathsf{A}$ and $\mathsf{B}$, we will call the morphisms in $(\mathsf{nCat}^{\gl})^{S}$ from $\mathsf{A}$ to 
$\mathsf{B}$ \emph{$S$-sorted globular $n$-functors}. Thus, an $S$-sorted globular $n$-functor from $\mathsf{A}$ to $\mathsf{B}$ is an $S$-indexed family $F=(F_{s})_{s\in S}$ such that, for every sort $s\in S$, $F_{s}$ is a globular $n$-functor from $\mathsf{A}_{s}$ to $\mathsf{B}_{s}$. We will denote by $\mathrm{nFunc}^{\gl}(\mathsf{A},\mathsf{B})$ the set of all $S$-sorted globular $n$-functors from $\mathsf{A}$ to $\mathsf{B}$. We consider the expressions $F\in\mathrm{nFunc}^{\gl}(\mathsf{A},\mathsf{B})$ and $F\colon\mathsf{A}\mor\mathsf{B}$ as synonyms.
We recall that for $F\colon \mathsf{A}\mor\mathsf{B}$ and $G\colon \mathsf{B}\mor\mathsf{C}$, the \emph{composition} $G\circ F$ is $(G_{s}\circ F_{s})_{s\in S}$, and that for an $S$-sorted globular $n$-category $\mathsf{A}$, the \emph{identity} 
$S$-sorted globular $n$-functor at $\mathsf{A}$, that is, $\Id_{\mathsf{A}}$, is $(\Id_{\mathsf{A}_{s}})_{s\in S}$. 
\end{definition}

Next we define the corresponding notion of globular $n$-categorial algebra, following the style of the single-set $n$-categorial algebra (Definition~\ref{DManySortedCatAlg}).

\begin{definition}\label{DGlCatAlg}
Let $\mathsf{A}$ be an $S$-sorted globular $n$-category. The $S^{\star}\times S$-sorted set of finitary operations on $\mathsf{A}$ is $(\mathrm{nFunc}^{\gl}(\mathsf{A}_{\mathbf{s}},\mathsf{A}_{s}))_{(\mathbf{s},s)\in S^{\star}\times S}$, where, for every $\mathbf{s}\in S^{\star}$, $\mathsf{A}_{\mathbf{s}}=\prod_{j\in\bb{\mathbf{s}}}\mathsf{A}_{s_{j}}$, with $\bb{\mathbf{s}}$ denoting the length of the word $\mathbf{s}$ (if $\mathbf{s}=\lambda$, then $\mathsf{A}_{\lambda}$ is a final globular $n$-category). 

Let $\sig{\Sigma}=(S,\Sigma)$ be a many-sorted signature. 
A \emph{structure of globular $n$-categorial $\Sigma$-algebra on} an $S$-sorted globular $n$-category $\mathsf{A}$ is a family $F = (F_{\mathbf{s},s})_{(\mathbf{s},s)\in S^{\star}\times S}$, where, for $(\mathbf{s},s)\in S^{\star}\times S$, $F_{\mathbf{s},s}$ is a mapping from $\Sigma_{\mathbf{s},s}$ to $\mathrm{nFunc}^{\gl}(\mathsf{A}_{\mathbf{s}},\mathsf{A}_{s})$ (if $(\mathbf{s},s) = (\lambda,s)$ and $\sigma\in \Sigma_{\lambda,s}$, then $F_{\lambda,s}(\sigma)$ picks out a $0$-cell $u$ of $\mathsf{A}_{s}$ together with its identities $\mathrm{i}^{(k,0)}(u)$, $k\in n+1$, since $\mathsf{A}_{\lambda}$ is a final globular $n$-category, with a single cell in each dimension). The finite products $\mathsf{A}_{\mathbf{s}}$ are taken dimensionwise: the $k$-cells of $\mathsf{A}_{\mathbf{s}}$ are the families of $k$-cells, and all the operations are computed componentwise. An \emph{$S$-sorted globular $n$-categorial $\Sigma$-algebra} or, simply, an \emph{globular $n$-categorial $\sig{\Sigma}$-algebra}, is a pair $(\mathsf{A},F)$, abbreviated to $\calg{A}$, where $\mathsf{A}=(\mathsf{A}_{s})_{s\in S}$ is an $S$-sorted globular $n$-category and $F$ a structure of globular $n$-categorial $\Sigma$-algebra on $\mathsf{A}$. For a pair $(\mathbf{s},s)\in S^{\star}\times S$ and a formal operation $\sigma\in \Sigma_{\mathbf{s},s}$, in order to simplify the notation, the $n$-functor $F_{\mathbf{s},s}(\sigma)$ from $\mathsf{A}_{\mathbf{s}}$ to $\mathsf{A}_{s}$ will be written as $F_{\sigma}$, or simply as $\sigma^{\calg{A}}$. 

An \emph{$S$-sorted globular $n$-categorial $\Sigma$-homomorphism} or, simply, an \emph{globular $n$-categorial $\sig{\Sigma}$-homomorphism}, from $\calg{A}$ to $\calg{B}$, where $\calg{B}=(\mathsf{B},G)$, is a triple $(\calg{A},F,\calg{B})$, abbreviated to $F\colon\calg{A}\mor\calg{B}$, where $F$ is an $S$-sorted globular $n$-functor from $\mathsf{A}$ to $\mathsf{B}$ such that, for every $k\in n+1$, every $(\mathbf{s},s)\in S^{\star}\times S$, every $\sigma\in\Sigma_{\mathbf{s},s}$ and every family of $k$-cells $(\mathsf{a}_{j})_{j\in\bb{\mathbf{s}}}\in\mathsf{A}_{\mathbf{s}}$, we have that
$$
F_{s,k}\left(\sigma^{\calg{A}}_{k}
\left(\left(\mathsf{a}_{j}
\right)_{j\in\bb{\mathbf{s}}}\right)\right)
=
\sigma^{\calg{B}}_{k}
\left(\left(
F_{s_{j},k}\left(
\mathsf{a}_{j}
\right)
\right)_{j\in\bb{\mathbf{s}}}
\right),
$$
where $F_{s,k}$ is the component at $k$-cells of the globular $n$-functor $F_{s}$ and $\sigma^{\calg{A}}_{k}$ that of $\sigma^{\calg{A}}$. The equation is required for the $n$-cells too, the cells of top dimension.
We denote by $\mathsf{nCat}^{\gl}\mathsf{Alg}(\sig{\Sigma})$ the category of globular $n$-categorial  $\sig{\Sigma}$-algebras and globular $n$-categorial $\sig{\Sigma}$-homomorphisms.
\end{definition}

We next introduce the $\omega$-dimensional counterpart of Definition~\ref{DGlCatAlg}.

\begin{definition}\label{DOmegaGlCatAlg}
Let $\sig{\Sigma}=(S,\Sigma)$ be a many-sorted signature. An \emph{globular $\omega$-categorial $\sig{\Sigma}$-algebra} is a pair $\calg{A}=(\mathsf{A},F)$ where $\mathsf{A}=(\mathsf{A}_{s})_{s\in S}$ is an $S$-sorted globular $\omega$-category and $F$ assigns to every $\sigma\in\Sigma_{\mathbf{s},s}$ a globular $\omega$-functor $\sigma^{\calg{A}}\colon\mathsf{A}_{\mathbf{s}}\mor\mathsf{A}_{s}$. An \emph{globular $\omega$-categorial $\sig{\Sigma}$-homomorphism} is an $S$-sorted globular $\omega$-functor satisfying the equation of Definition~\ref{DGlCatAlg} for every $k\in\omega$. We denote by $\omega\mathsf{Cat}^{\gl}\mathsf{Alg}(\sig{\Sigma})$ the category of globular $\omega$-categorial $\sig{\Sigma}$-algebras and globular $\omega$-categorial $\sig{\Sigma}$-homomorphisms.
\end{definition}

In the globular presentation, a globular $n$-categorial $\sig{\Sigma}$-algebra is regarded as a partial algebra on an $(n+1)\times S$-sorted set of cells, for a signature that enlarges $\Sigma$ with a copy of every operation symbol for every dimension, and with sources, targets, identities and compositions between dimensions.

\begin{definition}\label{DSigmaCatGlobular}
Let $\Sigma$ be an $S$-sorted signature and $n\in\mathbb{N}$. For every $k\in n+1$ and every $\mathbf{s}\in S^{\star}$, we let $(k,\mathbf{s})$ stand for the word $((k,s_{j}))_{j\in\lvert\mathbf{s}\rvert}$ in $((n+1)\times S)^{\star}$; in particular, $(k,\lambda)=\lambda$. The \emph{globular $n$-categorial signature determined by $\Sigma$}, denoted by $\Sigma^{\gl(n)}$, is the $(n+1)\times S$-sorted signature defined, for every $(\mathbf{w},(k,s))\in((n+1)\times S)^{\star}\times((n+1)\times S)$, as follows:
$$
{\Sigma}^{\gl(n)}_{\mathbf{w},(k,s)}=
\begin{cases}
\{\sigma_{k}\mid\sigma\in{\Sigma}_{\mathbf{s},s}\},
& \text{if $\mathbf{w}=(k,\mathbf{s})$ with $\mathbf{s}\neq(s,s)$;}\\
\{\sigma_{k}\mid\sigma\in{\Sigma}_{(s,s),s}\}\amalg\{\circ^{(j,k)}_{s}\mid j\in k\},
& \text{if $\mathbf{w}=((k,s),(k,s))$;}\\
\{\src^{(k,m)}_{s}, \tgt^{(k,m)}_{s}\},
& \text{if $\mathbf{w}=((m,s))$ with $k<m$;}\\
\{\mathrm{i}^{(k,j)}_{s}\},
& \text{if $\mathbf{w}=((j,s))$ with $j<k$;}\\
\varnothing,
& \text{otherwise.}
\end{cases}
$$
That is, the sort $(k,s)$ of ${\Sigma}^{\gl(n)}$ stands for the $k$-cells of sort $s$, and ${\Sigma}^{\gl(n)}$ is obtained from $\Sigma$ by taking, for every $k\in n+1$ and every operation symbol $\sigma\in\Sigma_{\mathbf{s},s}$, a copy $\sigma_{k}$ of $\sigma$ of arity $(k,\mathbf{s})$ and coarity $(k,s)$, which acts on the $k$-cells, and by adding, for every $j,k\in n+1$ with $j<k$ and every sort $s\in S$, two unary operation symbols $\src^{(j,k)}_{s}$ and $\tgt^{(j,k)}_{s}$, the \emph{$(j,k)$-source} and the \emph{$(j,k)$-target}, of arity $((k,s))$ and coarity $(j,s)$, a unary operation symbol $\mathrm{i}^{(k,j)}_{s}$, the \emph{$(k,j)$-identity}, of arity $((j,s))$ and coarity $(k,s)$, and a binary operation symbol $\circ^{(j,k)}_{s}$, the \emph{$j$-composition of $k$-cells}, of arity $((k,s),(k,s))$ and coarity $(k,s)$. As in Definition~\ref{DSpecGlobular}, we write $x\circ^{(j,k)}_{s}y$ for $\circ^{(j,k)}_{s}(x,y)$.

Moreover, we regard $\Sigma^{\gl(n)}$ as the hep-signature, see Definition~\ref{DHepSignature}, $\sig{\Sigma}^{\gl(n)}=((n+1)\times S,\Sigma^{\gl(n)},\mathrm{def}^{((n+1)\times S,\Sigma^{\gl(n)})})$ in which the compositions are the only operation symbols with a nonempty set of domain conditions: for every $j,k\in n+1$ with $j<k$ and every $s\in S$,
$$
\mathrm{def}^{((n+1)\times S,\Sigma^{\gl(n)})}(\circ^{(j,k)}_{s})=\{(\src^{(j,k)}_{s}(v^{(k,s)}_{0}),\tgt^{(j,k)}_{s}(v^{(k,s)}_{1}))\},
$$
where $v^{(k,s)}_{0}$ and $v^{(k,s)}_{1}$ are the variables of $\vs{((k,s),(k,s))}$, so that $v^{(k,s)}_{0}\circ^{(j,k)}_{s}v^{(k,s)}_{1}$ is defined if and only if the $(j,k)$-source of $v^{(k,s)}_{0}$ is equal to the $(j,k)$-target of $v^{(k,s)}_{1}$; and $\mathrm{def}^{((n+1)\times S,\Sigma^{\gl(n)})}(\tau)=\varnothing$ for every other operation symbol $\tau$ of $\Sigma^{\gl(n)}$. The relation $\lhd$ of $\sig{\Sigma}^{\gl(n)}$, see Definition~\ref{DHepSignature}, is well-founded, since the domain conditions of the compositions only involve sources and targets, which have empty sets of domain conditions; hence $\sig{\Sigma}^{\gl(n)}$ is indeed a hep-signature.

Finally, the \emph{globular $\omega$-categorial signature determined by $\Sigma$}, $\Sigma^{\gl(\omega)}$, and the hep-signature $\sig{\Sigma}^{\gl(\omega)}$ are the $\omega\times S$-sorted signature and the hep-signature defined by the same clauses, with $\omega$ in place of $n+1$.
\end{definition}

In contrast with the single-set case, $\Sigma^{\gl(n)}$ is not a subset of $\Sigma^{\gl(m)}$ for $n\leq m$ in $\mathbb{N}\cup\{\omega\}$, since the sets of sorts are different; however, the inclusion of $(n+1)\times S$ into $(m+1)\times S$, or into $\omega\times S$ if $m=\omega$, induces a canonical embedding of $\Sigma^{\gl(n)}$ into $\Sigma^{\gl(m)}$, and $\Sigma^{\gl(\omega)}$ is the union of the images of the embeddings of the $\Sigma^{\gl(n)}$, for $n\in\mathbb{N}$. Note also that, for $n=0$, $\Sigma^{\gl(0)}$ is, up to the identification of $1\times S$ with $S$, the signature $\Sigma$.

The specification of the globular presentation follows the pattern of that of the single-set presentation, see Definition~\ref{DSpecSingleSet}. Now, however, the role played in the single-set case by a dimension $k\in n$ is played by a pair of dimensions $j<k$ in $n+1$, since the $j$-composition acts on the $k$-cells and the $(k,j)$-identities take the place of the $k$-sources and $k$-targets as unit elements; and the role played by a pair of dimensions $l<k$ in $n$ is played by a triple of dimensions $h<j<k$ in $n+1$. Informally, the equations state that the operations of $\Sigma$ are always defined, that, for every sort $s\in S$ and every pair $j<k$, the $j$-cells and the $k$-cells of sort $s$ are the objects and the morphisms of a category, that these categories fit together as in an $n$-category, and that the operations of $\Sigma$ are compatible with them. In what follows, we abbreviate $V^{(n+1)\times S}$ to $V$.

\begin{definition}\label{DSpecGlobular}
Let $\sig{\Sigma}=(S,\Sigma)$ be a many-sorted signature and $n\in\mathbb{N}$. We will denote by $((n+1)\times S,\Sigma^{\gl(n)},V^{(n+1)\times S},\mathcal{C}^{\sig{\Sigma}^{\gl(n)}})$, written $\boldsymbol{\mathcal{C}}^{\sig{\Sigma}^{\gl(n)}}$ for short, the specification in which $\Sigma^{\gl(n)}$ is the globular $n$-categorial signature introduced in Definition~\ref{DSigmaCatGlobular}, $V^{(n+1)\times S}$ a fixed $(n+1)\times S$-sorted set with a countable infinity of variables in each coordinate, and $\mathcal{C}^{\sig{\Sigma}^{\gl(n)}}$ the subset of $\mathrm{QEeqt}(\Sigma^{\gl(n)})_{V^{(n+1)\times S}}$, the set of all $\QE$-equations with variables in $V^{(n+1)\times S}$, consisting of the following equations:

For every $k\in n+1$, every $(\mathbf{s},s)\in S^{\star}\times S$, every operation symbol $\sigma\in \Sigma_{\mathbf{s},s}$, and every family of variables $(x_{p})_{p\in \lvert\mathbf{s}\rvert}\in V_{(k,\mathbf{s})}$, the operation $\sigma_{k}$ applied to the family $(x_{p})_{p\in \lvert\mathbf{s}\rvert}$ is always defined. Formally,
\allowdisplaybreaks
\begin{align*}
\sigma_{k}((x_{p})_{p\in \lvert\mathbf{s}\rvert})\eeq\sigma_{k}((x_{p})_{p\in \lvert\mathbf{s}\rvert}).\label{AGlobularA0}\tag{EA0}
\end{align*}

For every $j,k\in n+1$ with $j<k$, the following equations describe the $j$-composition of $k$-cells.

For every sort $s\in S$, every variable $x\in V_{(k,s)}$ and every variable $u\in V_{(j,s)}$, the $(j,k)$-source and $(j,k)$-target of $x$ and the $(k,j)$-identity of $u$ are always defined. Formally,
\allowdisplaybreaks
\begin{align*}
\src^{(j,k)}_{s}(x)&\eeq\src^{(j,k)}_{s}(x);
&
\tgt^{(j,k)}_{s}(x)&\eeq\tgt^{(j,k)}_{s}(x);
&
\mathrm{i}^{(k,j)}_{s}(u)&\eeq\mathrm{i}^{(k,j)}_{s}(u).\label{AGlobularAjk1}\tag{EA$^{j,k}$1}
\end{align*}

For every sort $s\in S$ and every variable $u\in V_{(j,s)}$, the $(j,k)$-source and the $(j,k)$-target of the $(k,j)$-identity of $u$ are equal to $u$. Formally,
\allowdisplaybreaks
\begin{align*}
\src^{(j,k)}_{s}(\mathrm{i}^{(k,j)}_{s}(u))&\eeq u;
&
\tgt^{(j,k)}_{s}(\mathrm{i}^{(k,j)}_{s}(u))&\eeq u.\label{AGlobularAjk2}\tag{EA$^{j,k}$2}
\end{align*}
In other words, $\mathrm{i}^{(k,j)}_{s}$ is a common section of $\src^{(j,k)}_{s}$ and $\tgt^{(j,k)}_{s}$. This is the globular counterpart of the fact that $\src^{k}_{s}$ and $\tgt^{k}_{s}$ are right zeros.

For every sort $s\in S$ and every pair of variables $x,y\in V_{(k,s)}$, $x\circ^{(j,k)}_{s} y$ is defined if and only if the $(j,k)$-target of $y$ is equal to the $(j,k)$-source of $x$. Formally,
\allowdisplaybreaks
\begin{align*}
x\circ^{(j,k)}_{s} y\eeq x\circ^{(j,k)}_{s} y \;\;
&\longrightarrow \;\;
\src^{(j,k)}_{s}(x)\eeq\tgt^{(j,k)}_{s}(y);\\
\src^{(j,k)}_{s}(x)\eeq\tgt^{(j,k)}_{s}(y) \;\;
&\longrightarrow \;\;
x\circ^{(j,k)}_{s} y\eeq x\circ^{(j,k)}_{s}y.\label{AGlobularAjk3}\tag{EA$^{j,k}$3}
\end{align*}

For every sort $s\in S$ and every pair of variables $x,y\in V_{(k,s)}$, if $x\circ^{(j,k)}_{s} y$ is defined, then the $(j,k)$-source of $x\circ^{(j,k)}_{s}y$ is that of $y$ and the $(j,k)$-target of $x\circ^{(j,k)}_{s}y$ is that of $x$. Formally,
\allowdisplaybreaks
\begin{align*}
x\circ^{(j,k)}_{s}y\eeq x\circ^{(j,k)}_{s}y \;\;
&\longrightarrow \;\;
\src^{(j,k)}_{s}(x\circ^{(j,k)}_{s}y)\eeq\src^{(j,k)}_{s}(y);\\
x\circ^{(j,k)}_{s}y\eeq x\circ^{(j,k)}_{s}y \;\;
&\longrightarrow \;\;
\tgt^{(j,k)}_{s}(x\circ^{(j,k)}_{s}y)\eeq\tgt^{(j,k)}_{s}(x).\label{AGlobularAjk4}\tag{EA$^{j,k}$4}
\end{align*}

For every sort $s\in S$ and every variable $x\in V_{(k,s)}$, the compositions $x\circ^{(j,k)}_{s}\mathrm{i}^{(k,j)}_{s}(\src^{(j,k)}_{s}(x))$ and $\mathrm{i}^{(k,j)}_{s}(\tgt^{(j,k)}_{s}(x))\circ^{(j,k)}_{s}x$ are always defined and are equal to $x$, i.e., $\mathrm{i}^{(k,j)}_{s}(\src^{(j,k)}_{s}(x))$ is a right unit element for the $j$-composition with $x$ and $\mathrm{i}^{(k,j)}_{s}(\tgt^{(j,k)}_{s}(x))$ is a left unit element for the $j$-composition with $x$. Formally,
\allowdisplaybreaks
\begin{align*}
x\circ^{(j,k)}_{s}\mathrm{i}^{(k,j)}_{s}(\src^{(j,k)}_{s}(x))&\eeq x;
&
\mathrm{i}^{(k,j)}_{s}(\tgt^{(j,k)}_{s}(x))\circ^{(j,k)}_{s}x&\eeq x.\label{AGlobularAjk5}\tag{EA$^{j,k}$5}
\end{align*}

For every sort $s\in S$ and every triple of variables $x, y, z\in V_{(k,s)}$, if the compositions $x\circ^{(j,k)}_{s}y$ and $y\circ^{(j,k)}_{s}z$ are defined, then the compositions $x\circ^{(j,k)}_{s}(y\circ^{(j,k)}_{s}z)$ and $(x\circ^{(j,k)}_{s}y)\circ^{(j,k)}_{s}z$ are defined and they are equal, i.e., the $j$-composition of $k$-cells, when defined, is associative. Formally,
\allowdisplaybreaks
\begin{multline*}
(x\circ^{(j,k)}_{s}y\eeq x\circ^{(j,k)}_{s}y)
\wedge
(y\circ^{(j,k)}_{s}z\eeq y\circ^{(j,k)}_{s}z)\;\;
\longrightarrow \;\;\\
(x\circ^{(j,k)}_{s}y)\circ^{(j,k)}_{s}z\eeq x\circ^{(j,k)}_{s}(y\circ^{(j,k)}_{s}z).\label{AGlobularAjk6}\tag{EA$^{j,k}$6}
\end{multline*}

For every $(\mathbf{s},s)\in S^{\star}\times S$, every operation symbol $\sigma\in \Sigma_{\mathbf{s},s}$, every family of variables $(x_{p})_{p\in \lvert\mathbf{s}\rvert}\in V_{(k,\mathbf{s})}$ and every family of variables $(u_{p})_{p\in \lvert\mathbf{s}\rvert}\in V_{(j,\mathbf{s})}$, the $(j,k)$-source of $\sigma_{k}((x_{p})_{p\in \lvert\mathbf{s}\rvert})$ is equal to $\sigma_{j}$ applied to the family $((\src^{(j,k)}_{s_{p}}(x_{p}))_{p\in \lvert\mathbf{s}\rvert})$, the $(j,k)$-target of $\sigma_{k}((x_{p})_{p\in \lvert\mathbf{s}\rvert})$ is equal to $\sigma_{j}$ applied to the family $((\tgt^{(j,k)}_{s_{p}}(x_{p}))_{p\in \lvert\mathbf{s}\rvert})$, and the $(k,j)$-identity of $\sigma_{j}((u_{p})_{p\in \lvert\mathbf{s}\rvert})$ is equal to $\sigma_{k}$ applied to the family $((\mathrm{i}^{(k,j)}_{s_{p}}(u_{p}))_{p\in \lvert\mathbf{s}\rvert})$. Formally,
\allowdisplaybreaks
\begin{align*}
\src^{(j,k)}_{s}(\sigma_{k}((x_{p})_{p\in \lvert\mathbf{s}\rvert}))&\eeq\sigma_{j}((\src^{(j,k)}_{s_{p}}(x_{p}))_{p\in \lvert\mathbf{s}\rvert});\\
\tgt^{(j,k)}_{s}(\sigma_{k}((x_{p})_{p\in \lvert\mathbf{s}\rvert}))&\eeq\sigma_{j}((\tgt^{(j,k)}_{s_{p}}(x_{p}))_{p\in \lvert\mathbf{s}\rvert});\\
\mathrm{i}^{(k,j)}_{s}(\sigma_{j}((u_{p})_{p\in \lvert\mathbf{s}\rvert}))&\eeq\sigma_{k}((\mathrm{i}^{(k,j)}_{s_{p}}(u_{p}))_{p\in \lvert\mathbf{s}\rvert}).\label{AGlobularAjk7}\tag{EA$^{j,k}$7}
\end{align*}

For every $(\mathbf{s},s)\in S^{\star}\times S$, every operation symbol $\sigma\in \Sigma_{\mathbf{s},s}$, and every pair of families of variables $(x_{p})_{p\in \lvert\mathbf{s}\rvert}, (y_{p})_{p\in \lvert\mathbf{s}\rvert}\in V_{(k,\mathbf{s})}$, if, for every $p\in \lvert\mathbf{s}\rvert$, the compositions $x_{p}\circ^{(j,k)}_{s_{p}}y_{p}$ are defined, then the composition $\sigma_{k}((x_{p})_{p\in \lvert\mathbf{s}\rvert})\circ^{(j,k)}_{s}\sigma_{k}((y_{p})_{p\in \lvert\mathbf{s}\rvert})$ is defined and it is equal to $\sigma_{k}$ applied to the family $(x_{p}\circ^{(j,k)}_{s_{p}}y_{p})_{p\in \lvert\mathbf{s}\rvert}$. Formally,
\allowdisplaybreaks
\begin{multline*}
\textstyle
\bigwedge_{p\in \lvert\mathbf{s}\rvert}
(x_{p}\circ^{(j,k)}_{s_{p}} y_{p}\eeq x_{p}\circ^{(j,k)}_{s_{p}} y_{p})
\;\;\longrightarrow\;\;\\
\sigma_{k}((x_{p}\circ^{(j,k)}_{s_{p}}y_{p})_{p\in \lvert\mathbf{s}\rvert})
\eeq
\sigma_{k}((x_{p})_{p\in \lvert\mathbf{s}\rvert})\circ^{(j,k)}_{s}\sigma_{k}((y_{p})_{p\in \lvert\mathbf{s}\rvert}).
\label{AGlobularAjk8}\tag{EA$^{j,k}$8}
\end{multline*}

For every $h,j,k\in n+1$ with $h<j<k$, the following equations relate the $h$-, $j$- and $k$-cells.

For every sort $s\in S$ and every variable $x\in V_{(k,s)}$, the $(h,j)$-source of both the $(j,k)$-source and the $(j,k)$-target of $x$ is equal to the $(h,k)$-source of $x$, and the $(h,j)$-target of both the $(j,k)$-source and the $(j,k)$-target of $x$ is equal to the $(h,k)$-target of $x$. Formally,
\allowdisplaybreaks
\begin{align*}
\src^{(h,j)}_{s}(\src^{(j,k)}_{s}(x))&\eeq\src^{(h,k)}_{s}(x);
&
\tgt^{(h,j)}_{s}(\tgt^{(j,k)}_{s}(x))&\eeq\tgt^{(h,k)}_{s}(x);\\
\src^{(h,j)}_{s}(\tgt^{(j,k)}_{s}(x))&\eeq\src^{(h,k)}_{s}(x);
&
\tgt^{(h,j)}_{s}(\src^{(j,k)}_{s}(x))&\eeq\tgt^{(h,k)}_{s}(x).\label{AGlobularAhjk1}\tag{EA$^{h,j,k}$1}
\end{align*}

For every sort $s\in S$ and every variable $u\in V_{(h,s)}$, the $(k,j)$-identity of the $(j,h)$-identity of $u$ is the $(k,h)$-identity of $u$. Formally,
\allowdisplaybreaks
\begin{align*}
\mathrm{i}^{(k,j)}_{s}(\mathrm{i}^{(j,h)}_{s}(u))&\eeq\mathrm{i}^{(k,h)}_{s}(u).\label{AGlobularAhjk2}\tag{EA$^{h,j,k}$2}
\end{align*}

For every sort $s\in S$ and every pair of variables $x,y\in V_{(k,s)}$, if $x\circ^{(h,k)}_{s}y$ is defined, then the $(j,k)$-source of $x\circ^{(h,k)}_{s}y$ is the $h$-composition of the $(j,k)$-source of $x$ with the $(j,k)$-source of $y$, and the $(j,k)$-target of $x\circ^{(h,k)}_{s}y$ is the $h$-composition of the $(j,k)$-target of $x$ with the $(j,k)$-target of $y$. Formally,
\allowdisplaybreaks
\begin{align*}
x\circ^{(h,k)}_{s}y&\eeq x\circ^{(h,k)}_{s}y \;\;\longrightarrow\;\;
\src^{(j,k)}_{s}(x\circ^{(h,k)}_{s}y)\eeq\src^{(j,k)}_{s}(x)\circ^{(h,j)}_{s}\src^{(j,k)}_{s}(y);\\
x\circ^{(h,k)}_{s}y&\eeq x\circ^{(h,k)}_{s}y \;\;\longrightarrow\;\;
\tgt^{(j,k)}_{s}(x\circ^{(h,k)}_{s}y)\eeq\tgt^{(j,k)}_{s}(x)\circ^{(h,j)}_{s}\tgt^{(j,k)}_{s}(y).
\label{AGlobularAhjk3}\tag{EA$^{h,j,k}$3}
\end{align*}

For every sort $s\in S$ and every quadruple of variables $x,y,z,t\in V_{(k,s)}$, if the compositions $x\circ^{(j,k)}_{s}y$, $z\circ^{(j,k)}_{s}t$, $x\circ^{(h,k)}_{s}z$ and $y\circ^{(h,k)}_{s}t$ are defined, the $(h,k)$-composition of $x\circ^{(j,k)}_{s}y$ with $z\circ^{(j,k)}_{s}t$ is equal to the $(j,k)$-composition of $x\circ^{(h,k)}_{s}z$ with $y\circ^{(h,k)}_{s}t$. Formally,
\allowdisplaybreaks
\begin{multline*}
(x\circ^{(j,k)}_{s}y\eeq x\circ^{(j,k)}_{s}y)
\wedge
(z\circ^{(j,k)}_{s}t\eeq z\circ^{(j,k)}_{s}t)
\wedge
(x\circ^{(h,k)}_{s}z\eeq x\circ^{(h,k)}_{s}z)
\wedge
(y\circ^{(h,k)}_{s}t\eeq y\circ^{(h,k)}_{s}t)
\\
\;\;\longrightarrow\;\;
(x\circ^{(j,k)}_{s}y)\circ^{(h,k)}_{s}(z\circ^{(j,k)}_{s}t)\eeq
(x\circ^{(h,k)}_{s}z)\circ^{(j,k)}_{s}(y\circ^{(h,k)}_{s}t).
\label{AGlobularAhjk4}\tag{EA$^{h,j,k}$4}
\end{multline*}

For every sort $s\in S$ and every pair of variables $u,v\in V_{(j,s)}$, if $u\circ^{(h,j)}_{s}v$ is defined, then the $(k,j)$-identity of $u\circ^{(h,j)}_{s}v$ is the $h$-composition of the $(k,j)$-identities of $u$ and $v$. Formally,
\allowdisplaybreaks
\begin{align*}
u\circ^{(h,j)}_{s}v\eeq u\circ^{(h,j)}_{s}v \;\;\longrightarrow\;\;
\mathrm{i}^{(k,j)}_{s}(u\circ^{(h,j)}_{s}v)\eeq\mathrm{i}^{(k,j)}_{s}(u)\circ^{(h,k)}_{s}\mathrm{i}^{(k,j)}_{s}(v).
\label{AGlobularAhjk5}\tag{EA$^{h,j,k}$5}
\end{align*}

The axioms (EA$^{h,j,k}$1)--(EA$^{h,j,k}$5) are the counterparts of the conditions (MS1), (MS4), (MS3), (MS9) and (MS8) of Definition~\ref{DOmegaCatGl}, and (EA$^{j,k}$1)--(EA$^{j,k}$6) those of (MS2), (MS5), (MS6) and (MS7). Hence a model of the axioms (EA$^{j,k}$1)--(EA$^{j,k}$6), for $j<k$ in $n+1$, and (EA$^{h,j,k}$1)--(EA$^{h,j,k}$5), for $h<j<k$ in $n+1$, is an $S$-sorted globular $n$-category, see Definition~\ref{DSortednCatGl}, presented in its globular form, i.e., with the cells of each dimension in a separate sort.

Since $\sig{\Sigma}^{\gl(n)}$ is a hep-signature, see Definition~\ref{DSigmaCatGlobular}, $(\sig{\Sigma}^{\gl(n)},\mathcal{C}^{\sig{\Sigma}^{\gl(n)}})$ is a hep-specification, see Definition~\ref{DEpSpec}, and we will work with the ep-variety $\mathcal{V}(\sig{\Sigma}^{\gl(n)},\mathcal{C}^{\sig{\Sigma}^{\gl(n)}})$ and with its canonically associated category $\mathsf{EPAlg}(\sig{\Sigma}^{\gl(n)},\mathcal{C}^{\sig{\Sigma}^{\gl(n)}})$, see Definition~\ref{DEpVariety}.
\end{definition}

%The axioms (EA0), (EA$^{j,k}$7) and (EA$^{j,k}$8) state, as in the compact case, that the operations of $\Sigma$ act layer by layer and compatibly with the $n$-categorial structure, i.e., that they are $n$-functors between $n$-categories in their globular presentation. The relation between the two presentations, however, does not come from these axioms, but from the correspondence between the globular and the single-set presentations of $n$-categories. A model of $\boldsymbol{\mathcal{C}}^{\sig{\Sigma}^{\mathrm{ext}(n)}}$ gives rise to a model of $\boldsymbol{\mathcal{C}}^{\sig{\Sigma}^{(n)}}$ by keeping only its $n$-cells, on which $k$-sources and $k$-targets are obtained by composing $(k,n)$-sources and $(k,n)$-targets with $(n,k)$-identities. Conversely, a model of $\boldsymbol{\mathcal{C}}^{\sig{\Sigma}^{(n)}}$ gives rise to a model of $\boldsymbol{\mathcal{C}}^{\sig{\Sigma}^{\mathrm{ext}(n)}}$ by taking, as $k$-cells of sort $s$, for $k\in n$, the cells $x$ with $\src^{k}_{s}(x)=x$. Going from the compact presentation to the extended one and back gives the original algebra, whereas going from the extended presentation to the compact one and back only gives an isomorphic algebra, in which every $k$-cell has been replaced by its identity $n$-cell. It is therefore to be expected that $\mathsf{PAlg}(\boldsymbol{\mathcal{C}}^{\sig{\Sigma}^{\mathrm{ext}(n)}})$ is equivalent to $\mathsf{PAlg}(\boldsymbol{\mathcal{C}}^{\sig{\Sigma}^{(n)}})$, and hence to $\mathsf{nCatAlg}(\sig{\Sigma})$, although the two canonical functors are not mutually inverse; proving this equivalence in a way that is functorial in $\sig{\Sigma}$ is part of Objective~II.2.

For $n=\omega$, we let $\boldsymbol{\mathcal{C}}^{\sig{\Sigma}^{\gl(\omega)}}=(\omega\times S,\Sigma^{\gl(\omega)},V^{\omega\times S},\mathcal{C}^{\sig{\Sigma}^{\gl(\omega)}})$ stand for the specification whose axioms are the equations (EA0), for every $k\in\omega$, (EA$^{j,k}$1)--(EA$^{j,k}$8), for every $j<k$ in $\omega$, and (EA$^{h,j,k}$1)--(EA$^{h,j,k}$5), for every $h<j<k$ in $\omega$, and by $(\sig{\Sigma}^{\gl(\omega)},\mathcal{C}^{\sig{\Sigma}^{\gl(\omega)}})$ the corresponding hep-specification.

In the following proposition we characterize, as in the single-set case, the ep-algebras of type $\sig{\Sigma}^{\gl(n)}$ in terms of the axioms of Definition~\ref{DSpecGlobular}.

\begin{proposition}\label{PGlobularEp}
Let $\sig{\Sigma}=(S,\Sigma)$ be a many-sorted signature, $n\in\mathbb{N}\cup\{\omega\}$ and $\mathbf{A}$ a partial $\Sigma^{\gl(n)}$-algebra. Then $\mathbf{A}$ is a $\sig{\Sigma}^{\gl(n)}$-algebra, see Definitions~\ref{DEpAlgebra} and~\ref{DSigmaCatGlobular}, if, and only if, $\mathbf{A}$ satisfies the axiom (EA0) and, for every pair of dimensions $j<k$, the axioms (EA$^{j,k}$1) and (EA$^{j,k}$3).
\end{proposition}
\begin{proof}
By Definition~\ref{DEpAlgebra}, $\mathbf{A}$ is a $\sig{\Sigma}^{\gl(n)}$-algebra if, and only if, for every operation symbol $\tau$ of $\Sigma^{\gl(n)}$, the domain of $\tau^{\mathbf{A}}$ is the set of the families that satisfy the domain conditions of $\tau$ in $\mathbf{A}$. We consider the different kinds of operation symbols of $\Sigma^{\gl(n)}$, see Definition~\ref{DSigmaCatGlobular}. Let $j<k$ be dimensions and $s\in S$.

\textsf{(Operation symbols $\sigma_{k}$)} The empty sets of domain conditions of the operation symbols $\sigma_{k}$ state that they are total, i.e., the axiom (EA0).

\textsf{(Sources, targets and identities)} The empty sets of domain conditions of the operation symbols $\src^{(j,k)}_{s}$, $\tgt^{(j,k)}_{s}$ and $\mathrm{i}^{(k,j)}_{s}$ state that they are total, i.e., the axiom (EA$^{j,k}$1).

\textsf{(Compositions)} Assuming that the $(j,k)$-sources and the $(j,k)$-targets are total, the domain condition of $\circ^{(j,k)}_{s}$ states that $x\circ^{(j,k)}_{s}y$ is defined if, and only if, $\src^{(j,k)}_{s}(x)=\tgt^{(j,k)}_{s}(y)$, i.e., the axiom (EA$^{j,k}$3).

This completes the proof.
\end{proof}

By Proposition~\ref{PGlobularEp}, a partial $\Sigma^{\gl(n)}$-algebra in which every axiom of $\mathcal{C}^{\sig{\Sigma}^{\gl(n)}}$ holds is automatically a $\sig{\Sigma}^{\gl(n)}$-algebra, since it satisfies, in particular, the axiom (EA0) and, for every pair of dimensions $j<k$, the axioms (EA$^{j,k}$1) and (EA$^{j,k}$3). Hence the objects of $\mathsf{EPAlg}(\sig{\Sigma}^{\gl(n)},\mathcal{C}^{\sig{\Sigma}^{\gl(n)}})$ are exactly the models of the $\QE$-specification $\boldsymbol{\mathcal{C}}^{\sig{\Sigma}^{\gl(n)}}$.

We next relate the globular $n$-categorial $\sig{\Sigma}$-algebras with the $\sig{\Sigma}^{\gl(n)}$-algebras. Let $\sig{\Sigma}=(S,\Sigma)$ be a many-sorted signature and $n\in\mathbb{N}\cup\{\omega\}$. In what follows, a pair $(\mathsf{A},F)$, where $\mathsf{A}=(\mathsf{A}_{s})_{s\in S}$ is an $S$-sorted globular $n$-category and $F$ assigns to every $\sigma\in\Sigma_{\mathbf{s},s}$ a family of mappings $(F_{\sigma,k})_{k}$, one for each dimension $k$, is identified with the partial $\Sigma^{\gl(n)}$-algebra whose underlying set of sort $(k,s)$ is the set of $k$-cells of $\mathsf{A}_{s}$, in which $\src^{(j,k)}_{s}$, $\tgt^{(j,k)}_{s}$, $\mathrm{i}^{(k,j)}_{s}$ and $\circ^{(j,k)}_{s}$ are interpreted as the corresponding operations of $\mathsf{A}_{s}$, and $\sigma_{k}$ is interpreted as $F_{\sigma,k}$, the component at $k$-cells of $F_{\sigma}$.

\begin{remark}\label{RGlobularIdent}
Let $\sig{\Sigma}=(S,\Sigma)$ be a many-sorted signature. Then, under the above identification, the following equalities hold
$$
\mathsf{nCat}^{\gl}\mathsf{Alg}(\sig{\Sigma})=\mathsf{EPAlg}(\sig{\Sigma}^{\gl(n)},\mathcal{C}^{\sig{\Sigma}^{\gl(n)}}),
\text{ for }n\in\mathbb{N},
\qquad\text{and}\qquad
\omega\mathsf{Cat}^{\gl}\mathsf{Alg}(\sig{\Sigma})=\mathsf{EPAlg}(\sig{\Sigma}^{\gl(\omega)},\mathcal{C}^{\sig{\Sigma}^{\gl(\omega)}}).
$$
In fact, in a globular $n$-category, the $(j,k)$-sources, $(j,k)$-targets and $(k,j)$-identities are total mappings, i.e., (EA$^{j,k}$1) holds, and the domain of $\circ^{(j,k)}$ is the fibered product of Definition~\ref{DOmegaCatGl}, i.e., (EA$^{j,k}$3) holds; as noted after the axioms of Definition~\ref{DSpecGlobular}, the conditions (MS1)--(MS9) of Definition~\ref{DOmegaCatGl} correspond to the remaining axioms on the globular structure; since the products $\mathsf{A}_{\mathbf{s}}$ are taken dimensionwise, a family of total mappings $(F_{\sigma,k})_{k}$ is a globular $n$-functor if, and only if, the axioms (EA0), (EA$^{j,k}$7) and (EA$^{j,k}$8) hold for $\sigma$; and the globular $n$-categorial $\sig{\Sigma}$-homomorphisms are exactly the homomorphisms of partial $\Sigma^{\gl(n)}$-algebras, see Definitions~\ref{DPartialAlg} and~\ref{DGlCatAlg}. Finally, by Proposition~\ref{PGlobularEp}, the models of $\boldsymbol{\mathcal{C}}^{\sig{\Sigma}^{\gl(n)}}$ are $\sig{\Sigma}^{\gl(n)}$-algebras. In contrast with the single-set presentation, see Proposition~\ref{PFiniteDimNotQE}, in the globular presentation every cell has a dimension, namely the first coordinate of its sort, so that no further condition is needed in dimension $\omega$. The correspondences between the conditions of Definition~\ref{DOmegaCatGl} and the axioms are the following:
\begin{align*}
&\text{(MS1)}\leftrightarrow\text{(EA$^{h,j,k}$1)};
&&\text{(MS2)}\leftrightarrow\text{(EA$^{j,k}$4)};
&&\text{(MS3)}\leftrightarrow\text{(EA$^{h,j,k}$3)};
\\
&\text{(MS4)}\leftrightarrow\text{(EA$^{h,j,k}$2)};
&&\text{(MS5)}\leftrightarrow\text{(EA$^{j,k}$2)};
&&\text{(MS6)}\leftrightarrow\text{(EA$^{j,k}$5)};
\\
&\text{(MS7)}\leftrightarrow\text{(EA$^{j,k}$6)};
&&\text{(MS8)}\leftrightarrow\text{(EA$^{h,j,k}$5)};
&&\text{(MS9)}\leftrightarrow\text{(EA$^{h,j,k}$4)}.
\end{align*}
\end{remark}

\section{Graduation and Unification of Higher-Order Categorial Algebras}\label{SGraduation}

In this section, we associate to a single-set $n$-categorial $\sig{\Sigma}$-algebra a globular $n$-categorial $\sig{\Sigma}$-algebra by means of the notion of graduation and, conversely, to a globular $n$-categorial $\sig{\Sigma}$-algebra a single-set $n$-categorial $\sig{\Sigma}$-algebra by means of the notion of unification. Both constructions are induced by generalized hep-derivors. We end this section by showing that graduation and unification are mutually quasi-inverse equivalences of categories.

By Remarks~\ref{RSingleSetIdent} and~\ref{RGlobularIdent}, for every many-sorted signature $\sig{\Sigma}=(S,\Sigma)$ and every $n\in\mathbb{N}$, the following equalities hold
$$
\mathsf{nCatAlg}(\sig{\Sigma})=\mathsf{EPAlg}(\sig{\Sigma}^{(n)},\mathcal{C}^{\sig{\Sigma}^{(n)}})
\qquad\text{and}\qquad
\mathsf{nCat}^{\gl}\mathsf{Alg}(\sig{\Sigma})=\mathsf{EPAlg}(\sig{\Sigma}^{\gl(n)},\mathcal{C}^{\sig{\Sigma}^{\gl(n)}}).
$$
Thus, the categories of higher-order categorial algebras are the categories canonically associated to the ep-varieties $\mathcal{V}(\sig{\Sigma}^{(n)},\mathcal{C}^{\sig{\Sigma}^{(n)}})$ and $\mathcal{V}(\sig{\Sigma}^{\gl(n)},\mathcal{C}^{\sig{\Sigma}^{\gl(n)}})$ of hep-specifications. The functors of this section are induced by generalized hep-derivors between the corresponding hep-signatures, i.e., by morphisms of $\mathsf{hepSig}_{\gd}$, see Proposition~\ref{PHepSigCat}. They are defined on all the partial algebras, and we prove, by means of Proposition~\ref{PTransformedEpAlg}, that their restrictions to the categories of higher-order categorial algebras corestrict to categories of higher-order categorial algebras, see Definition~\ref{DCorestriction} and Subsection~\ref{SEpRestriction}. By abuse of notation, we denote these corestrictions as the functors themselves.

For the generalized hep-derivors of this work we use the following notation, which we explain for the graduation. For a many-sorted signature $\sig{\Sigma}=(S,\Sigma)$ and $n\in\mathbb{N}$, the graduation functor $\Gd^{(n)}_{\sig{\Sigma}}$ is induced by the generalized hep-derivor
$$
\bGd^{(n)}_{\sig{\Sigma}}=\left(\varphi^{\Gd^{(n)}_{\sig{\Sigma}}},\Gd^{(n)}_{\sig{\Sigma}},\Dom^{\Gd^{(n)}_{\sig{\Sigma}}}\right),
$$
whose mapping on sorts is $\varphi^{\Gd^{(n)}_{\sig{\Sigma}}}$, whose assignment of terms to operation symbols is denoted, as the functor, by $\Gd^{(n)}_{\sig{\Sigma}}$, and whose domains are $\Dom^{\Gd^{(n)}_{\sig{\Sigma}}}$. Since the assignment of terms is applied to operation symbols and the functor to algebras and homomorphisms, this will not cause confusion. The same convention applies to the unification, discretization and truncation derivors $\bUf$, $\bDs$ and $\bTr$. For the compositions of these generalized hep-derivors, see Definition~\ref{DHepComp}, we keep the notation $\Dom^{\mathbf{e}\circ\mathbf{d}}$, e.g., $\Dom^{\bGd^{(n)}_{\sig{\Sigma}}\circ\bUf^{(n)}_{\sig{\Sigma}}}$. The dimensions of $\Sigma^{\gl(n)}$ are the elements of $n+1$. In this section and in Section~\ref{STruncation} we only consider finite dimensions; the $\omega$-dimensional functors are considered in Subsection~\ref{SDTOmega} and in Section~\ref{SOmegaFunctors}.

\subsection{Graduation of higher-order categorial algebras}\label{SGraduationSub}

Graduating a single-set $n$-categorial $\sig{\Sigma}$-algebra $\calg{A}$ consists in distributing its cells into sorts according to their dimension: for $k\in n$, the $k$-cells of sort $s$, i.e., the cells $x$ of sort $s$ with $\src^{k}_{s}(x)=x$, become the cells of sort $(k,s)$, and all the cells of sort $s$ become the cells of sort $(n,s)$. The $(j,k)$-source and the $(j,k)$-target of a $k$-cell are its $j$-source and its $j$-target, the $(k,j)$-identity of a $j$-cell is the cell itself, the $j$-composition of two $k$-cells is their $j$-composition, and $\sigma_{k}$ acts on $k$-cells as $\sigma$. All these data are given by a single generalized hep-derivor, which we define next.

\begin{definition}\label{DGradDerivor}
Let $\sig{\Sigma}=(S,\Sigma)$ be a many-sorted signature and $n\in\mathbb{N}$. The \emph{graduation derivor} is the generalized hep-derivor
$$
\bGd^{(n)}_{\sig{\Sigma}}=\left(\varphi^{\Gd^{(n)}_{\sig{\Sigma}}},\Gd^{(n)}_{\sig{\Sigma}},\Dom^{\Gd^{(n)}_{\sig{\Sigma}}}\right)
$$
from $\sig{\Sigma}^{\gl(n)}$ to $\sig{\Sigma}^{(n)}$ defined as follows:
\begin{enumerate}
\item for every dimension $k$ and every $s\in S$, $\varphi^{\Gd^{(n)}_{\sig{\Sigma}}}(k,s)=s$;
\item for every $\sigma\in\Sigma_{\mathbf{s},s}$, every pair of dimensions $j<k$ and every $s\in S$,
\begin{align*}
\Gd^{(n)}_{\sig{\Sigma}}(\sigma_{k})&=\sigma((v^{\mathbf{s}}_{p})_{p\in\bb{\mathbf{s}}}), &
\Gd^{(n)}_{\sig{\Sigma}}(\mathrm{i}^{(k,j)}_{s})&=v^{s}_{0},\\
\Gd^{(n)}_{\sig{\Sigma}}(\src^{(j,k)}_{s})&=\src^{j}_{s}(v^{s}_{0}), &
\Gd^{(n)}_{\sig{\Sigma}}(\tgt^{(j,k)}_{s})&=\tgt^{j}_{s}(v^{s}_{0}),\\
\Gd^{(n)}_{\sig{\Sigma}}(\circ^{(j,k)}_{s})&=v^{s}_{0}\circ^{j}_{s}v^{s}_{1};
\end{align*}
\item for every $s\in S$ and every $k\in n$,
$$
\Dom^{\Gd^{(n)}_{\sig{\Sigma}}}_{(k,s)}=\left\lbrace\left(v^{s}_{0},\src^{k}_{s}(v^{s}_{0})\right)\right\rbrace
\qquad\text{and}\qquad
\Dom^{\Gd^{(n)}_{\sig{\Sigma}}}_{(n,s)}=\varnothing,
$$
\end{enumerate}
\end{definition}

Let us note that $(\varphi^{\Gd^{(n)}_{\sig{\Sigma}}})^{\star}((k,\mathbf{s}))=\mathbf{s}$, so that $\Gd^{(n)}_{\sig{\Sigma}}(\sigma_{k})$ is a term with variables in $\vs{\mathbf{s}}$, and that $j\in n$ whenever $j<k$, so that $\src^{j}_{s}$, $\tgt^{j}_{s}$ and $\circ^{j}_{s}$ are operation symbols of $\Sigma^{(n)}$. Moreover, by Definition~\ref{DEEq}, for every partial $\Sigma^{(n)}$-algebra $\mathbf{B}$, every $s\in S$ and every $x\in B_{s}$, $\mathbf{B}\models\Dom^{\Gd^{(n)}_{\sig{\Sigma}}}_{(k,s)}[x]$ if, and only if, $x=\src^{k}_{s}(x)$, i.e., $x$ is a $k$-cell, for $k\in n$, and every element satisfies $\Dom^{\Gd^{(n)}_{\sig{\Sigma}}}_{(n,s)}$.

\begin{definition}\label{DGraduation}
Let $\sig{\Sigma}=(S,\Sigma)$ be a many-sorted signature and $n\in\mathbb{N}$. The \emph{graduation functor} $\Gd^{(n)}_{\sig{\Sigma}}$ is the functor $(\bGd^{(n)}_{\sig{\Sigma}})^{\ast}_{\gd}$ induced by the generalized hep-derivor $\bGd^{(n)}_{\sig{\Sigma}}$ from $\mathsf{PAlg}(\Sigma^{(n)})$ to $\mathsf{PAlg}(\Sigma^{\gl(n)})$, see Proposition~\ref{PGenHepDerivorFunctor}.
\end{definition}

In the following lemma we collect the facts about cells of a single-set $n$-categorial $\sig{\Sigma}$-algebra that we need to graduate it.

\begin{lemma}\label{LGradCells}
Let $\sig{\Sigma}=(S,\Sigma)$ be a many-sorted signature, $n\in\mathbb{N}$, $\mathbf{B}$ an object in $\mathsf{EPAlg}(\sig{\Sigma}^{(n)},\mathcal{C}^{\sig{\Sigma}^{(n)}})$, $s\in S$ and $j<k$ dimensions in $n+1$. Then the following statements hold:
\begin{enumerate}
\item if $k\in n$, $(\mathbf{s},s)\in S^{\star}\times S$, $\sigma\in\Sigma_{\mathbf{s},s}$ and $(x_{p})_{p\in\bb{\mathbf{s}}}$ is a family of $k$-cells, then $\sigma^{\mathbf{B}}((x_{p})_{p\in\bb{\mathbf{s}}})$ is a $k$-cell;
\item for every cell $x$ of sort $s$, $\src^{j}_{s}(x)$ and $\tgt^{j}_{s}(x)$ are $j$-cells;
\item if $k\in n$, then every $j$-cell of sort $s$ is a $k$-cell;
\item if $k\in n$ and $x$, $y$ are $k$-cells of sort $s$ such that $\src^{j}_{s}(x)=\tgt^{j}_{s}(y)$, then $x\circ^{j}_{s}y$ is defined and it is a $k$-cell.
\end{enumerate}
\end{lemma}
\begin{proof}
(1) The following chain of equalities holds
\allowdisplaybreaks
\begin{align*}
\src^{k}_{s}\left(\sigma\left((x_{p})_{p\in\bb{\mathbf{s}}}\right)\right)
&=
\sigma\left(\left(\src^{k}_{s_{p}}(x_{p})\right)_{p\in\bb{\mathbf{s}}}\right)
\tag{1}
\\&=
\sigma\left((x_{p})_{p\in\bb{\mathbf{s}}}\right).
\tag{2}
\end{align*}

The first equality follows from~(A$^{k}$7); finally, the last equality follows from the fact that every $x_{p}$ is a $k$-cell.

(2) The following chain of equalities holds
\allowdisplaybreaks
\begin{align*}
\src^{j}_{s}\left(\src^{j}_{s}(x)\right)
&=
\src^{j}_{s}(x),
&
\src^{j}_{s}\left(\tgt^{j}_{s}(x)\right)
&=
\tgt^{j}_{s}(x),
\end{align*}
by~(A$^{j}$2).

(3) Let $u$ be a $j$-cell of sort $s$. Then the following chain of equalities holds
\allowdisplaybreaks
\begin{align*}
\src^{k}_{s}(u)
&=
\src^{k}_{s}\left(\src^{j}_{s}(u)\right)
\tag{1}
\\&=
\src^{j}_{s}(u)
\tag{2}
\\&=
u.
\tag{3}
\end{align*}

The first equality follows from the fact that $u$ is a $j$-cell; the second equality follows from~(AB$^{j,k}$1); finally, the last equality follows, again, from the fact that $u$ is a $j$-cell.

(4) By~(A$^{j}$3), $x\circ^{j}_{s}y$ is defined. Moreover, the following chain of equalities holds
\allowdisplaybreaks
\begin{align*}
\src^{k}_{s}\left(x\circ^{j}_{s}y\right)
&=
\src^{k}_{s}(x)\circ^{j}_{s}\src^{k}_{s}(y)
\tag{1}
\\&=
x\circ^{j}_{s}y.
\tag{2}
\end{align*}

The first equality follows from~(AB$^{j,k}$2); finally, the last equality follows from the fact that $x$ and $y$ are $k$-cells.

This completes the proof.
\end{proof}

In the following lemma we describe the transformed algebra along the graduation derivor.

\begin{lemma}\label{LGradDescription}
Let $\sig{\Sigma}=(S,\Sigma)$ be a many-sorted signature, $n\in\mathbb{N}$ and $\mathbf{B}$ an object in $\mathsf{EPAlg}(\sig{\Sigma}^{(n)},\mathcal{C}^{\sig{\Sigma}^{(n)}})$. Then $\Gd^{(n)}_{\sig{\Sigma}}(\mathbf{B})$ is a $\sig{\Sigma}^{\gl(n)}$-algebra whose cells of sort $(k,s)$ are the $k$-cells of sort $s$ of $\mathbf{B}$, for $k\in n$, and all the cells of sort $s$ of $\mathbf{B}$, for $k=n$. Moreover, for every pair of dimensions $j<k$ and every $s\in S$, the following statements hold:
\begin{enumerate}
\item $\sigma_{k}$, $\src^{(j,k)}_{s}$, $\tgt^{(j,k)}_{s}$ and $\mathrm{i}^{(k,j)}_{s}$ are total in $\Gd^{(n)}_{\sig{\Sigma}}(\mathbf{B})$ and act as $\sigma^{\mathbf{B}}$, $\src^{j\,\mathbf{B}}_{s}$, $\tgt^{j\,\mathbf{B}}_{s}$ and the inclusion, respectively; and
\item for every pair of $k$-cells $x$, $y$ of sort $s$, $x\circ^{(j,k)}_{s}y$ is defined in $\Gd^{(n)}_{\sig{\Sigma}}(\mathbf{B})$ if, and only if, $\src^{j\,\mathbf{B}}_{s}(x)=\tgt^{j\,\mathbf{B}}_{s}(y)$, and, in this case, the $j$-composition of $x$ and $y$ in $\Gd^{(n)}_{\sig{\Sigma}}(\mathbf{B})$ is $x\circ^{j\,\mathbf{B}}_{s}y$.
\end{enumerate}
\end{lemma}
\begin{proof}
The description of the cells follows from Definitions~\ref{DDomainSet} and~\ref{DGradDerivor} and the note after Definition~\ref{DGradDerivor}. By Proposition~\ref{PTransformedEpAlg}, to prove that $\Gd^{(n)}_{\sig{\Sigma}}(\mathbf{B})$ is a $\sig{\Sigma}^{\gl(n)}$-algebra, we have to show that, for every operation symbol $\tau$ of $\Sigma^{\gl(n)}$ and every family $\mathbf{x}$ of cells of $\Gd^{(n)}_{\sig{\Sigma}}(\mathbf{B})$ that satisfies the domain conditions of $\tau$ in $\Gd^{(n)}_{\sig{\Sigma}}(\mathbf{B})$, the partial term operation $\Gd^{(n)}_{\sig{\Sigma}}(\tau)^{\mathbf{B}}$ is defined at $\mathbf{x}$ and takes a value in the domain of the coarity of $\tau$. We consider the different kinds of operation symbols of $\Sigma^{\gl(n)}$. Let $j<k$ be dimensions and $s\in S$.

\textsf{(Operation symbols $\sigma_{k}$)} The set of domain conditions of $\sigma_{k}$ is empty and the term $\sigma((v^{\mathbf{s}}_{p})_{p\in\bb{\mathbf{s}}})$ is defined everywhere in $\mathbf{B}$, by~(A0). If $k=n$, there is nothing more to prove. If $k\in n$, then, by Lemma~\ref{LGradCells}~(1), the value at a family of $k$-cells is a $k$-cell.

\textsf{(Sources and targets)} The sets of domain conditions of $\src^{(j,k)}_{s}$ and $\tgt^{(j,k)}_{s}$ are empty and the terms $\src^{j}_{s}(v^{s}_{0})$ and $\tgt^{j}_{s}(v^{s}_{0})$ are defined everywhere in $\mathbf{B}$, by~(A$^{j}$1). By Lemma~\ref{LGradCells}~(2), their values are $j$-cells.

\textsf{(Identities)} The set of domain conditions of $\mathrm{i}^{(k,j)}_{s}$ is empty and the term $v^{s}_{0}$ is defined everywhere. If $k=n$, there is nothing more to prove. If $k\in n$, then, by Lemma~\ref{LGradCells}~(3), every $j$-cell is a $k$-cell.

Therefore, in $\Gd^{(n)}_{\sig{\Sigma}}(\mathbf{B})$, $\sigma_{k}$, $\src^{(j,k)}_{s}$, $\tgt^{(j,k)}_{s}$ and $\mathrm{i}^{(k,j)}_{s}$ are total and act as $\sigma^{\mathbf{B}}$, $\src^{j\,\mathbf{B}}_{s}$, $\tgt^{j\,\mathbf{B}}_{s}$ and the inclusion, respectively.

\textsf{(Compositions)} By the above, two $k$-cells $x$, $y$ of sort $s$ satisfy the domain condition of $\circ^{(j,k)}_{s}$ in $\Gd^{(n)}_{\sig{\Sigma}}(\mathbf{B})$ if, and only if, $\src^{j}_{s}(x)=\tgt^{j}_{s}(y)$ in $\mathbf{B}$. In this case, by~(A$^{j}$3), $x\circ^{j}_{s}y$ is defined in $\mathbf{B}$. If $k=n$, there is nothing more to prove. If $k\in n$, then, by Lemma~\ref{LGradCells}~(4), $x\circ^{j}_{s}y$ is a $k$-cell.

All in all, by Proposition~\ref{PTransformedEpAlg}, $\Gd^{(n)}_{\sig{\Sigma}}(\mathbf{B})$ is a $\sig{\Sigma}^{\gl(n)}$-algebra and, by Proposition~\ref{PTransformedAlg}, its operations are as stated.

This completes the proof.
\end{proof}

We next prove that the graduation of a single-set $n$-categorial $\sig{\Sigma}$-algebra is a globular $n$-categorial $\sig{\Sigma}$-algebra.

\begin{proposition}\label{PGradDerivor}
Let $\sig{\Sigma}=(S,\Sigma)$ be a many-sorted signature and $n\in\mathbb{N}$. Then the restriction of $\Gd^{(n)}_{\sig{\Sigma}}$ to $\mathsf{EPAlg}(\sig{\Sigma}^{(n)},\mathcal{C}^{\sig{\Sigma}^{(n)}})$ corestricts to $\mathsf{EPAlg}(\sig{\Sigma}^{\gl(n)},\mathcal{C}^{\sig{\Sigma}^{\gl(n)}})$.
\end{proposition}
\begin{proof}
Let $\mathbf{B}$ be an object in $\mathsf{EPAlg}(\sig{\Sigma}^{(n)},\mathcal{C}^{\sig{\Sigma}^{(n)}})$. By Definitions~\ref{DCorestriction} and~\ref{DEpVariety}, we have to prove that $\Gd^{(n)}_{\sig{\Sigma}}(\mathbf{B})$ is a $\sig{\Sigma}^{\gl(n)}$-algebra in which every axiom of $\mathcal{C}^{\sig{\Sigma}^{\gl(n)}}$ holds.

\textsf{(Ep-algebra)} By Lemma~\ref{LGradDescription}, which rests on Proposition~\ref{PTransformedEpAlg}, $\Gd^{(n)}_{\sig{\Sigma}}(\mathbf{B})$ is a $\sig{\Sigma}^{\gl(n)}$-algebra.

\textsf{(Axioms)} By Lemma~\ref{LGradDescription}, the operations of $\Gd^{(n)}_{\sig{\Sigma}}(\mathbf{B})$ are the restrictions of those of $\mathbf{B}$, and $x\circ^{(j,k)}_{s}y$ is defined in $\Gd^{(n)}_{\sig{\Sigma}}(\mathbf{B})$ if, and only if, $x\circ^{j}_{s}y$ is defined in $\mathbf{B}$. Under this description, the axioms of $\boldsymbol{\mathcal{C}}^{\sig{\Sigma}^{\gl(n)}}$ become statements on $\mathbf{B}$, which we check in turn. In the following computations, the operation symbols of $\Sigma^{\gl(n)}$ are interpreted in $\Gd^{(n)}_{\sig{\Sigma}}(\mathbf{B})$ and those of $\Sigma^{(n)}$ in $\mathbf{B}$. Let $h<j<k$ be dimensions and $s\in S$.

\textsf{(EA0) and (EA$^{j,k}$1)} They hold by Lemma~\ref{LGradDescription}~(1).

\textsf{(EA$^{j,k}$2)} Let $u$ be a $j$-cell of sort $s$. Then the following chain of equalities holds
\allowdisplaybreaks
\begin{align*}
\src^{(j,k)}_{s}\left(\mathrm{i}^{(k,j)}_{s}(u)\right)
&=
\src^{j}_{s}(u)
\tag{1}
\\&=
u.
\tag{2}
\end{align*}

The first equality follows from Lemma~\ref{LGradDescription}~(1); finally, the last equality follows from the fact that $u$ is a $j$-cell. By a similar argument we also have that $\tgt^{(j,k)}_{s}(\mathrm{i}^{(k,j)}_{s}(u))=u$.

\textsf{(EA$^{j,k}$3)} It holds by Lemma~\ref{LGradDescription}~(2).

\textsf{(EA$^{j,k}$4)--(EA$^{j,k}$8)} By Lemma~\ref{LGradDescription}, they become the axioms (A$^{j}$4)--(A$^{j}$8) restricted to $k$-cells, except for the last equation of (EA$^{j,k}$7), which becomes trivial, since the identities are inclusions. For instance, for every $k$-cell $x$ of sort $s$, the following chain of equalities holds
\allowdisplaybreaks
\begin{align*}
x\circ^{(j,k)}_{s}\mathrm{i}^{(k,j)}_{s}\left(\src^{(j,k)}_{s}(x)\right)
&=
x\circ^{j}_{s}\src^{j}_{s}(x)
\tag{1}
\\&=
x.
\tag{2}
\end{align*}

The first equality follows from Lemma~\ref{LGradDescription}; finally, the last equality follows from~(A$^{j}$5).

\textsf{(EA$^{h,j,k}$1)} Let $x$ be a $k$-cell of sort $s$. Then the following chain of equalities holds
\allowdisplaybreaks
\begin{align*}
\src^{(h,j)}_{s}\left(\src^{(j,k)}_{s}(x)\right)
&=
\src^{h}_{s}\left(\src^{j}_{s}(x)\right)
\tag{1}
\\&=
\src^{h}_{s}(x)
\tag{2}
\\&=
\src^{(h,k)}_{s}(x).
\tag{3}
\end{align*}

The first equality follows from Lemma~\ref{LGradDescription}; the second equality follows from~(AB$^{h,j}$1); finally, the last equality follows from Lemma~\ref{LGradDescription}. By a similar argument, the remaining equations of (EA$^{h,j,k}$1) also hold.

\textsf{(EA$^{h,j,k}$2) and (EA$^{h,j,k}$5)} They become trivial, since the identities are inclusions.

\textsf{(EA$^{h,j,k}$3)} Let $x$, $y$ be $k$-cells of sort $s$ such that $x\circ^{(h,k)}_{s}y$ is defined. Then the following chain of equalities holds
\allowdisplaybreaks
\begin{align*}
\src^{(j,k)}_{s}\left(x\circ^{(h,k)}_{s}y\right)
&=
\src^{j}_{s}\left(x\circ^{h}_{s}y\right)
\tag{1}
\\&=
\src^{j}_{s}(x)\circ^{h}_{s}\src^{j}_{s}(y)
\tag{2}
\\&=
\src^{(j,k)}_{s}(x)\circ^{(h,j)}_{s}\src^{(j,k)}_{s}(y).
\tag{3}
\end{align*}

The first equality follows from Lemma~\ref{LGradDescription}; the second equality follows from~(AB$^{h,j}$2); finally, the last equality follows from Lemma~\ref{LGradDescription}, taking into account that, by Lemma~\ref{LGradCells}~(2), $\src^{j}_{s}(x)$ and $\src^{j}_{s}(y)$ are $j$-cells. By a similar argument we also have the equation for the targets.

\textsf{(EA$^{h,j,k}$4)} By Lemma~\ref{LGradDescription}, it becomes the axiom (AB$^{h,j}$3) restricted to $k$-cells, after exchanging the variables $y$ and $z$.

All in all, every axiom of $\boldsymbol{\mathcal{C}}^{\sig{\Sigma}^{\gl(n)}}$ holds in $\Gd^{(n)}_{\sig{\Sigma}}(\mathbf{B})$, so that $\Gd^{(n)}_{\sig{\Sigma}}(\mathbf{B})$ is an object in $\mathsf{EPAlg}(\sig{\Sigma}^{\gl(n)},\mathcal{C}^{\sig{\Sigma}^{\gl(n)}})$.

This completes the proof.
\end{proof}

\begin{remark}\label{RGradNotEp}
Let $\sig{\Sigma}=(S,\Sigma)$ be a many-sorted signature and $n\in\mathbb{N}$. The restriction of $\Gd^{(n)}_{\sig{\Sigma}}$ to $\mathsf{EPAlg}(\sig{\Sigma}^{(n)})$ does not corestrict to $\mathsf{EPAlg}(\sig{\Sigma}^{\gl(n)})$: in a $\sig{\Sigma}^{(n)}$-algebra in which $\src^{j}_{s}$ is not idempotent, the $j$-source of a cell need not be a $j$-cell, so that $\src^{(j,k)}_{s}$ need not be total in the transformed algebra. This is why we restrict $\Gd^{(n)}_{\sig{\Sigma}}$ to the category canonically associated to the ep-variety $\mathcal{V}(\sig{\Sigma}^{(n)},\mathcal{C}^{\sig{\Sigma}^{(n)}})$.
\end{remark}

Let $\sig{\Sigma}=(S,\Sigma)$ be a many-sorted signature and $n\in\mathbb{N}$. By Proposition~\ref{PGradDerivor} and Remarks~\ref{RSingleSetIdent} and~\ref{RGlobularIdent}, the restriction of $\Gd^{(n)}_{\sig{\Sigma}}$ to $\mathsf{nCatAlg}(\sig{\Sigma})$ corestricts to $\mathsf{nCat}^{\gl}\mathsf{Alg}(\sig{\Sigma})$, and, by abuse of notation, we denote the corestriction by
$$
\Gd^{(n)}_{\sig{\Sigma}}\colon\mathsf{nCatAlg}(\sig{\Sigma})\mor\mathsf{nCat}^{\gl}\mathsf{Alg}(\sig{\Sigma}).
$$
Explicitly, by Lemma~\ref{LGradDescription}, $\Gd^{(n)}_{\sig{\Sigma}}(\calg{A})$ is the graduated algebra described at the beginning of this subsection, and, by Proposition~\ref{PGenHepDerivorFunctor}, $\Gd^{(n)}_{\sig{\Sigma}}$ sends an $n$-categorial $\sig{\Sigma}$-homomorphism $F$ to the family of its restrictions to the $k$-cells.

\subsection{Unification of higher-order categorial algebras}\label{SUnificationSub}

Unifying a globular $n$-categorial $\sig{\Sigma}$-algebra $\calg{A}$, for $n\in\mathbb{N}$, consists in keeping only its $n$-cells, in which every $k$-cell has been replaced by its $(n,k)$-identity: the $k$-source of an $n$-cell is the $(n,k)$-identity of its $(k,n)$-source, the $k$-composition is the $k$-composition of $n$-cells, and $\sigma$ acts as $\sigma_{n}$. These data are given by the following generalized hep-derivor.

\begin{definition}\label{DUnifDerivor}
Let $\sig{\Sigma}=(S,\Sigma)$ be a many-sorted signature and $n\in\mathbb{N}$. The \emph{unification derivor} is the generalized hep-derivor
$$
\bUf^{(n)}_{\sig{\Sigma}}=\left(\varphi^{\Uf^{(n)}_{\sig{\Sigma}}},\Uf^{(n)}_{\sig{\Sigma}},\Dom^{\Uf^{(n)}_{\sig{\Sigma}}}\right)
$$
from $\sig{\Sigma}^{(n)}$ to $\sig{\Sigma}^{\gl(n)}$ defined as follows: $\varphi^{\Uf^{(n)}_{\sig{\Sigma}}}(s)=(n,s)$ and $\Dom^{\Uf^{(n)}_{\sig{\Sigma}}}_{s}=\varnothing$, for every $s\in S$, and, for every $\sigma\in\Sigma_{\mathbf{s},s}$, every $k\in n$ and every $s\in S$,
\begin{align*}
\Uf^{(n)}_{\sig{\Sigma}}(\sigma)&=\sigma_{n}((v^{(n,\mathbf{s})}_{p})_{p\in\bb{\mathbf{s}}}), &
\Uf^{(n)}_{\sig{\Sigma}}(\circ^{k}_{s})&=v^{(n,s)}_{0}\circ^{(k,n)}_{s}v^{(n,s)}_{1},\\
\Uf^{(n)}_{\sig{\Sigma}}(\src^{k}_{s})&=\mathrm{i}^{(n,k)}_{s}(\src^{(k,n)}_{s}(v^{(n,s)}_{0})), &
\Uf^{(n)}_{\sig{\Sigma}}(\tgt^{k}_{s})&=\mathrm{i}^{(n,k)}_{s}(\tgt^{(k,n)}_{s}(v^{(n,s)}_{0})).
\end{align*}
\end{definition}

\begin{definition}\label{DUnification}
Let $\sig{\Sigma}=(S,\Sigma)$ be a many-sorted signature and $n\in\mathbb{N}$. The \emph{unification functor} $\Uf^{(n)}_{\sig{\Sigma}}$ is the functor $(\bUf^{(n)}_{\sig{\Sigma}})^{\ast}_{\gd}$ induced by the generalized hep-derivor $\bUf^{(n)}_{\sig{\Sigma}}$ from $\mathsf{PAlg}(\Sigma^{\gl(n)})$ to $\mathsf{PAlg}(\Sigma^{(n)})$, see Proposition~\ref{PGenHepDerivorFunctor}.
\end{definition}

In the following lemma we prove that the identities of a globular $n$-categorial $\sig{\Sigma}$-algebra are injective.

\begin{lemma}\label{LUnifInj}
Let $\sig{\Sigma}=(S,\Sigma)$ be a many-sorted signature, $n\in\mathbb{N}$, $\mathbf{A}$ an object in $\mathsf{EPAlg}(\sig{\Sigma}^{\gl(n)},\mathcal{C}^{\sig{\Sigma}^{\gl(n)}})$, $j<k$ dimensions and $s\in S$. Then, for every $j$-cell $u$ of sort $s$, the following equalities hold
$$
\src^{(j,k)}_{s}\left(\mathrm{i}^{(k,j)}_{s}(u)\right)=u=\tgt^{(j,k)}_{s}\left(\mathrm{i}^{(k,j)}_{s}(u)\right).
$$
In particular, $\mathrm{i}^{(k,j)}_{s}$ is injective.
\end{lemma}
\begin{proof}
The equalities are the axiom (EA$^{j,k}$2). Hence $\mathrm{i}^{(k,j)}_{s}$ has a left inverse, so that it is injective.
\end{proof}

In the following lemma we describe the transformed algebra along the unification derivor.

\begin{lemma}\label{LUnifDescription}
Let $\sig{\Sigma}=(S,\Sigma)$ be a many-sorted signature, $n\in\mathbb{N}$ and $\mathbf{A}$ an object in $\mathsf{EPAlg}(\sig{\Sigma}^{\gl(n)},\mathcal{C}^{\sig{\Sigma}^{\gl(n)}})$. Then $\Uf^{(n)}_{\sig{\Sigma}}(\mathbf{A})$ is a $\sig{\Sigma}^{(n)}$-algebra whose cells of sort $s$ are the $n$-cells of sort $s$ of $\mathbf{A}$. Moreover, for every $k\in n$ and every $s\in S$, the following statements hold:
\begin{enumerate}
\item $\sigma$, $\src^{k}_{s}$ and $\tgt^{k}_{s}$ are total in $\Uf^{(n)}_{\sig{\Sigma}}(\mathbf{A})$ and act as $\sigma^{\mathbf{A}}_{n}$, $\mathrm{i}^{(n,k)\,\mathbf{A}}_{s}\circ\src^{(k,n)\,\mathbf{A}}_{s}$ and $\mathrm{i}^{(n,k)\,\mathbf{A}}_{s}\circ\tgt^{(k,n)\,\mathbf{A}}_{s}$, respectively; and
\item for every pair of $n$-cells $x$, $y$ of sort $s$, $x\circ^{k}_{s}y$ is defined in $\Uf^{(n)}_{\sig{\Sigma}}(\mathbf{A})$ if, and only if, $\src^{(k,n)\,\mathbf{A}}_{s}(x)=\tgt^{(k,n)\,\mathbf{A}}_{s}(y)$, and, in this case, the $k$-composition of $x$ and $y$ in $\Uf^{(n)}_{\sig{\Sigma}}(\mathbf{A})$ is $x\circ^{(k,n)\,\mathbf{A}}_{s}y$.
\end{enumerate}
\end{lemma}
\begin{proof}
Since the domains of $\bUf^{(n)}_{\sig{\Sigma}}$ are empty, the cells of sort $s$ of $\Uf^{(n)}_{\sig{\Sigma}}(\mathbf{A})$ are the $n$-cells of sort $s$ of $\mathbf{A}$. By Proposition~\ref{PTransformedEpAlg}, we check that, for every operation symbol of $\Sigma^{(n)}$, the corresponding partial term operation is defined on the families that satisfy its domain conditions in $\Uf^{(n)}_{\sig{\Sigma}}(\mathbf{A})$.

\textsf{(Operation symbols of $\Sigma$, sources and targets)} Their sets of domain conditions are empty, and the terms $\sigma_{n}((v^{(n,\mathbf{s})}_{p})_{p\in\bb{\mathbf{s}}})$, $\mathrm{i}^{(n,k)}_{s}(\src^{(k,n)}_{s}(v^{(n,s)}_{0}))$ and $\mathrm{i}^{(n,k)}_{s}(\tgt^{(k,n)}_{s}(v^{(n,s)}_{0}))$ are defined everywhere in $\mathbf{A}$, by~(EA0) and~(EA$^{k,n}$1). Moreover, all the $n$-cells are cells of $\Uf^{(n)}_{\sig{\Sigma}}(\mathbf{A})$.

\textsf{(Compositions)} By the above, two $n$-cells $x$, $y$ of sort $s$ satisfy the domain condition of $\circ^{k}_{s}$ in $\Uf^{(n)}_{\sig{\Sigma}}(\mathbf{A})$ if, and only if,
$$
\mathrm{i}^{(n,k)}_{s}\left(\src^{(k,n)}_{s}(x)\right)=\mathrm{i}^{(n,k)}_{s}\left(\tgt^{(k,n)}_{s}(y)\right),
$$
i.e., by Lemma~\ref{LUnifInj}, if, and only if, $\src^{(k,n)}_{s}(x)=\tgt^{(k,n)}_{s}(y)$. In this case, by~(EA$^{k,n}$3), $x\circ^{(k,n)}_{s}y$ is defined.

All in all, by Proposition~\ref{PTransformedEpAlg}, $\Uf^{(n)}_{\sig{\Sigma}}(\mathbf{A})$ is a $\sig{\Sigma}^{(n)}$-algebra and, by Proposition~\ref{PTransformedAlg}, its operations are as stated.

This completes the proof.
\end{proof}

We next prove that the unification of a globular $n$-categorial $\sig{\Sigma}$-algebra is a single-set $n$-categorial $\sig{\Sigma}$-algebra.

\begin{proposition}\label{PUnifDerivor}
Let $\sig{\Sigma}=(S,\Sigma)$ be a many-sorted signature and $n\in\mathbb{N}$. Then the restriction of $\Uf^{(n)}_{\sig{\Sigma}}$ to $\mathsf{EPAlg}(\sig{\Sigma}^{\gl(n)},\mathcal{C}^{\sig{\Sigma}^{\gl(n)}})$ corestricts to $\mathsf{EPAlg}(\sig{\Sigma}^{(n)},\mathcal{C}^{\sig{\Sigma}^{(n)}})$.
\end{proposition}
\begin{proof}
Let $\mathbf{A}$ be an object in $\mathsf{EPAlg}(\sig{\Sigma}^{\gl(n)},\mathcal{C}^{\sig{\Sigma}^{\gl(n)}})$. By Definitions~\ref{DCorestriction} and~\ref{DEpVariety}, we have to prove that $\Uf^{(n)}_{\sig{\Sigma}}(\mathbf{A})$ is a $\sig{\Sigma}^{(n)}$-algebra in which every axiom of $\mathcal{C}^{\sig{\Sigma}^{(n)}}$ holds.

\textsf{(Ep-algebra)} By Lemma~\ref{LUnifDescription}, which rests on Proposition~\ref{PTransformedEpAlg}, $\Uf^{(n)}_{\sig{\Sigma}}(\mathbf{A})$ is a $\sig{\Sigma}^{(n)}$-algebra.

\textsf{(Axioms)} By Lemma~\ref{LUnifDescription}, the axioms of $\boldsymbol{\mathcal{C}}^{\sig{\Sigma}^{(n)}}$ in $\Uf^{(n)}_{\sig{\Sigma}}(\mathbf{A})$ become statements on the $n$-cells of $\mathbf{A}$, which we call their translations and which we check in turn. In the following computations, all the operation symbols are interpreted in $\mathbf{A}$. Let $l<k$ in $n$ and $s\in S$.

\textsf{(A0), (A$^{k}$1) and (A$^{k}$3)} Their translations are consequences of (EA0), (EA$^{k,n}$1), and (EA$^{k,n}$3) together with Lemma~\ref{LUnifInj}, respectively.

\textsf{(A$^{k}$2)} Let $x$ be an $n$-cell of sort $s$. Then the following chain of equalities holds
\allowdisplaybreaks
\begin{align*}
\mathrm{i}^{(n,k)}_{s}\left(\src^{(k,n)}_{s}\left(
\mathrm{i}^{(n,k)}_{s}\left(\src^{(k,n)}_{s}(x)\right)
\right)\right)
&=
\mathrm{i}^{(n,k)}_{s}\left(\src^{(k,n)}_{s}(x)\right).
\tag{1}
\end{align*}

The equality follows from Lemma~\ref{LUnifInj}. By a similar argument, the remaining equations of (A$^{k}$2) also hold.

\textsf{(A$^{k}$4), (A$^{k}$5), (A$^{k}$6) and (A$^{k}$8)} Their translations are consequences of (EA$^{k,n}$4), (EA$^{k,n}$5), (EA$^{k,n}$6) and (EA$^{k,n}$8), respectively.

\textsf{(A$^{k}$7)} Let $(\mathbf{s},s)\in S^{\star}\times S$, $\sigma\in\Sigma_{\mathbf{s},s}$ and $(x_{p})_{p\in\bb{\mathbf{s}}}$ a family of $n$-cells. Then the following chain of equalities holds
\allowdisplaybreaks
\begin{align*}
\mathrm{i}^{(n,k)}_{s}\left(\src^{(k,n)}_{s}\left(\sigma_{n}\left((x_{p})_{p\in\bb{\mathbf{s}}}\right)\right)\right)
&=
\mathrm{i}^{(n,k)}_{s}\left(\sigma_{k}\left(\left(\src^{(k,n)}_{s_{p}}(x_{p})\right)_{p\in\bb{\mathbf{s}}}\right)\right)
\tag{1}
\\&=
\sigma_{n}\left(\left(\mathrm{i}^{(n,k)}_{s_{p}}\left(\src^{(k,n)}_{s_{p}}(x_{p})\right)\right)_{p\in\bb{\mathbf{s}}}\right).
\tag{2}
\end{align*}

The first equality follows from the first equation of~(EA$^{k,n}$7); finally, the last equality follows from the third equation of~(EA$^{k,n}$7). By a similar argument we also have the equation for the targets.

\textsf{(AB$^{l,k}$1)} Let $x$ be an $n$-cell of sort $s$ and $u=\src^{(l,n)}_{s}(x)$. Then the following chain of equalities holds
\allowdisplaybreaks
\begin{align*}
\mathrm{i}^{(n,k)}_{s}\left(\src^{(k,n)}_{s}\left(\mathrm{i}^{(n,l)}_{s}\left(\src^{(l,n)}_{s}(x)\right)\right)\right)
&=
\mathrm{i}^{(n,k)}_{s}\left(\src^{(k,n)}_{s}\left(\mathrm{i}^{(n,k)}_{s}\left(\mathrm{i}^{(k,l)}_{s}(u)\right)\right)\right)
\tag{1}
\\&=
\mathrm{i}^{(n,k)}_{s}\left(\mathrm{i}^{(k,l)}_{s}(u)\right)
\tag{2}
\\&=
\mathrm{i}^{(n,l)}_{s}\left(\src^{(l,n)}_{s}(x)\right).
\tag{3}
\end{align*}

The first equality follows from~(EA$^{l,k,n}$2); the second equality follows from Lemma~\ref{LUnifInj}; finally, the last equality follows from~(EA$^{l,k,n}$2).

Moreover, by~(EA$^{l,k,n}$1) and Lemma~\ref{LUnifInj}, for every $k$-cell $w$ of sort $s$, the following chain of equalities holds
\allowdisplaybreaks
\begin{align*}
\src^{(l,n)}_{s}\left(\mathrm{i}^{(n,k)}_{s}(w)\right)
&=
\src^{(l,k)}_{s}\left(\src^{(k,n)}_{s}\left(\mathrm{i}^{(n,k)}_{s}(w)\right)\right)
\tag{1}
\\&=
\src^{(l,k)}_{s}(w).
\tag{2}
\end{align*}

Hence, taking $w=\src^{(k,n)}_{s}(x)$, we have that
\allowdisplaybreaks
\begin{align*}
\mathrm{i}^{(n,l)}_{s}\left(\src^{(l,n)}_{s}\left(\mathrm{i}^{(n,k)}_{s}\left(\src^{(k,n)}_{s}(x)\right)\right)\right)
&=
\mathrm{i}^{(n,l)}_{s}\left(\src^{(l,k)}_{s}\left(\src^{(k,n)}_{s}(x)\right)\right)
\tag{1}
\\&=
\mathrm{i}^{(n,l)}_{s}\left(\src^{(l,n)}_{s}(x)\right).
\tag{2}
\end{align*}

The first equality follows from the above; finally, the last equality follows from~(EA$^{l,k,n}$1). By a similar argument, the remaining equations of (AB$^{l,k}$1) also hold.

\textsf{(AB$^{l,k}$2)} Let $x$, $y$ be $n$-cells of sort $s$ such that $x\circ^{(l,n)}_{s}y$ is defined. Then the following chain of equalities holds
\allowdisplaybreaks
\begin{align*}
\mathrm{i}^{(n,k)}_{s}\left(\src^{(k,n)}_{s}\left(x\circ^{(l,n)}_{s}y\right)\right)
&=
\mathrm{i}^{(n,k)}_{s}\left(\src^{(k,n)}_{s}(x)\circ^{(l,k)}_{s}\src^{(k,n)}_{s}(y)\right)
\tag{1}
\\&=
\mathrm{i}^{(n,k)}_{s}\left(\src^{(k,n)}_{s}(x)\right)\circ^{(l,n)}_{s}\mathrm{i}^{(n,k)}_{s}\left(\src^{(k,n)}_{s}(y)\right).
\tag{2}
\end{align*}

The first equality follows from~(EA$^{l,k,n}$3); finally, the last equality follows from~(EA$^{l,k,n}$5). By a similar argument we also have the equation for the targets.

\textsf{(AB$^{l,k}$3)} Its translation is the axiom (EA$^{l,k,n}$4), after exchanging the variables $y$ and $z$.

All in all, every axiom of $\boldsymbol{\mathcal{C}}^{\sig{\Sigma}^{(n)}}$ holds in $\Uf^{(n)}_{\sig{\Sigma}}(\mathbf{A})$, so that $\Uf^{(n)}_{\sig{\Sigma}}(\mathbf{A})$ is an object in $\mathsf{EPAlg}(\sig{\Sigma}^{(n)},\mathcal{C}^{\sig{\Sigma}^{(n)}})$.

This completes the proof.
\end{proof}

Let $\sig{\Sigma}=(S,\Sigma)$ be a many-sorted signature and $n\in\mathbb{N}$. By Proposition~\ref{PUnifDerivor} and Remarks~\ref{RSingleSetIdent} and~\ref{RGlobularIdent}, the restriction of $\Uf^{(n)}_{\sig{\Sigma}}$ to $\mathsf{nCat}^{\gl}\mathsf{Alg}(\sig{\Sigma})$ corestricts to $\mathsf{nCatAlg}(\sig{\Sigma})$, and, by abuse of notation, we denote the corestriction by
$$
\Uf^{(n)}_{\sig{\Sigma}}\colon\mathsf{nCat}^{\gl}\mathsf{Alg}(\sig{\Sigma})\mor\mathsf{nCatAlg}(\sig{\Sigma}).
$$
Explicitly, by Lemma~\ref{LUnifDescription}, the cells of sort $s$ of $\Uf^{(n)}_{\sig{\Sigma}}(\calg{A})$ are the $n$-cells of sort $s$ of $\calg{A}$, $\src^{k}_{s}=\mathrm{i}^{(n,k)}_{s}\circ\src^{(k,n)}_{s}$, $\tgt^{k}_{s}=\mathrm{i}^{(n,k)}_{s}\circ\tgt^{(k,n)}_{s}$, $\circ^{k}_{s}=\circ^{(k,n)}_{s}$ and $\sigma$ acts as $\sigma_{n}$; and, by Proposition~\ref{PGenHepDerivorFunctor}, $\Uf^{(n)}_{\sig{\Sigma}}(F)$ is the family of the components of $F$ at the $n$-cells.


\subsection{On the relationship between graduation and unification}\label{SGradUnif}

In this subsection, for a many-sorted signature $\sig{\Sigma}$ and $n\in\mathbb{N}$, we investigate the relation between the functors $\Gd^{(n)}_{\sig{\Sigma}}$ and $\Uf^{(n)}_{\sig{\Sigma}}$.
\begin{center}
	\begin{tikzpicture}
		[ACliment/.style={-{To [angle'=45, length=5.75pt, width=4pt, round]}}
		, scale=1,
		AClimentD/.style={dogble eqgal sign distance,
			-implies
			%{To [angle'=45, length=5.75pt, width=6pt, rognd]}
		}
		]

		\node[] (0R) at (0.5,0) [color=white] {};
		\node[] (1L) at (2.3,0)  [color=white] {};

		\node[] (C0) at (-0.5,0) [] {$\mathsf{nCatAlg}(\sig{\Sigma})$};
		\node[] (C1) at (3.5,0) [] {$\mathsf{nCat}^{\gl}\mathsf{Alg}(\sig{\Sigma})$};

		\draw[ACliment] (0R.north east) to node [above] {$\Gd^{(n)}_{\sig{\Sigma}}$} (1L.north west);
		\draw[ACliment] (1L.south west) to node [below] {$\Uf^{(n)}_{\sig{\Sigma}}$} (0R.south east);
	\end{tikzpicture}
\end{center}

Both functors are induced by generalized hep-derivors. Therefore we compute their composites by means of the compositions $\bGd^{(n)}_{\sig{\Sigma}}\circ\bUf^{(n)}_{\sig{\Sigma}}$ and $\bUf^{(n)}_{\sig{\Sigma}}\circ\bGd^{(n)}_{\sig{\Sigma}}$, see Definition~\ref{DHepComp}, and Proposition~\ref{PHepComp}. We start by proving that the process of graduation followed by that of unification does not provide anything new.

\begin{proposition}\label{PUnifGrad}
Let $\sig{\Sigma}=(S,\Sigma)$ be a many-sorted signature and $n\in\mathbb{N}$. Then the following equality holds in $\mathsf{hepSig}_{\gd}$
$$
\bGd^{(n)}_{\sig{\Sigma}}\circ\bUf^{(n)}_{\sig{\Sigma}}=\id_{\sig{\Sigma}^{(n)}}.
$$
Consequently, $\Uf^{(n)}_{\sig{\Sigma}}\circ\Gd^{(n)}_{\sig{\Sigma}}=\Id_{\mathsf{nCatAlg}(\sig{\Sigma})}$.
\end{proposition}
\begin{proof}
We prove the following two claims.

\begin{claim}\label{CUnifGradA}
The underlying derivor of $\bGd^{(n)}_{\sig{\Sigma}}\circ\bUf^{(n)}_{\sig{\Sigma}}$ is the identity derivor $(\id_{S},\eta^{\Sigma^{(n)}})$.
\end{claim}

By Definition~\ref{DHepComp}, the underlying derivor of $\bGd^{(n)}_{\sig{\Sigma}}\circ\bUf^{(n)}_{\sig{\Sigma}}$ is
$$
\left(\varphi^{\Gd^{(n)}_{\sig{\Sigma}}}\circ\varphi^{\Uf^{(n)}_{\sig{\Sigma}}},\left(\Gd^{(n)}_{\sig{\Sigma}}\right)^{\sharp}\circ\Uf^{(n)}_{\sig{\Sigma}}\right).
$$
Its mapping on sorts sends $s$ to $\varphi^{\Gd^{(n)}_{\sig{\Sigma}}}(n,s)=s$, i.e., it is $\id_{S}$. Let $k\in n$ and $s\in S$. Then the following chain of equalities holds
\allowdisplaybreaks
\begin{align*}
\left(\Gd^{(n)}_{\sig{\Sigma}}\right)^{\sharp}\left(\Uf^{(n)}_{\sig{\Sigma}}\left(\src^{k}_{s}\right)\right)
&=
\left(\Gd^{(n)}_{\sig{\Sigma}}\right)^{\sharp}\left(\mathrm{i}^{(n,k)}_{s}\left(\src^{(k,n)}_{s}(v^{(n,s)}_{0})\right)\right)
\tag{1}
\\&=
\src^{k}_{s}(v^{s}_{0}).
\tag{2}
\end{align*}

The first equality unravels Definition~\ref{DUnifDerivor}; finally, the last equality follows from~(Der1) and~(Der2), taking into account that, by Definition~\ref{DGradDerivor}, $\Gd^{(n)}_{\sig{\Sigma}}(\mathrm{i}^{(n,k)}_{s})$ is a variable and $\Gd^{(n)}_{\sig{\Sigma}}(\src^{(k,n)}_{s})=\src^{k}_{s}(v^{s}_{0})$. By a similar argument we also have that the term assigned to $\tgt^{k}_{s}$ is $\tgt^{k}_{s}(v^{s}_{0})$, the term assigned to $\circ^{k}_{s}$ is $v^{s}_{0}\circ^{k}_{s}v^{s}_{1}$ and the term assigned to $\sigma$ is $\sigma((v^{\mathbf{s}}_{p})_{p\in\bb{\mathbf{s}}})$, for every $\sigma\in\Sigma_{\mathbf{s},s}$. Thus $(\Gd^{(n)}_{\sig{\Sigma}})^{\sharp}\circ\Uf^{(n)}_{\sig{\Sigma}}=\eta^{\Sigma^{(n)}}$.

Claim~\ref{CUnifGradA} follows.

\begin{claim}\label{CUnifGradB}
The domains of $\bGd^{(n)}_{\sig{\Sigma}}\circ\bUf^{(n)}_{\sig{\Sigma}}$ are empty.
\end{claim}

Let $s\in S$. By Definitions~\ref{DHepComp}, \ref{DGradDerivor} and~\ref{DUnifDerivor}, the following chain of equalities holds
\allowdisplaybreaks
\begin{align*}
\Dom^{\bGd^{(n)}_{\sig{\Sigma}}\circ\bUf^{(n)}_{\sig{\Sigma}}}_{s}
&=
\Dom^{\Gd^{(n)}_{\sig{\Sigma}}}_{(n,s)}\cup\left(\Gd^{(n)}_{\sig{\Sigma}}\right)^{\sharp}\left[\Dom^{\Uf^{(n)}_{\sig{\Sigma}}}_{s}\right]
\tag{1}
\\&=
\varnothing\cup\left(\Gd^{(n)}_{\sig{\Sigma}}\right)^{\sharp}[\varnothing]
\tag{2}
\\&=
\varnothing.
\tag{3}
\end{align*}

The first equality unravels Definition~\ref{DHepComp}; the second equality follows from Definitions~\ref{DGradDerivor} and~\ref{DUnifDerivor}; finally, the last equality is clear.

Claim~\ref{CUnifGradB} follows.

By Claims~\ref{CUnifGradA} and~\ref{CUnifGradB} and Definition~\ref{DHepComp}, $\bGd^{(n)}_{\sig{\Sigma}}\circ\bUf^{(n)}_{\sig{\Sigma}}=\id_{\sig{\Sigma}^{(n)}}$.

Let $\calg{A}$ be an $n$-categorial $\sig{\Sigma}$-algebra. By Lemmas~\ref{LGradDescription} and~\ref{LUnifDescription}, the cells of sort $s$ of $\Uf^{(n)}_{\sig{\Sigma}}(\Gd^{(n)}_{\sig{\Sigma}}(\calg{A}))$ are the cells of sort $(n,s)$ of $\Gd^{(n)}_{\sig{\Sigma}}(\calg{A})$, i.e., all the cells of sort $s$ of $\calg{A}$, which, since the domains of $\id_{\sig{\Sigma}^{(n)}}$ are empty, are the cells of $(\id_{\sig{\Sigma}^{(n)}})^{\ast}_{\gd}(\calg{A})$. Moreover, by Propositions~\ref{PGradDerivor} and~\ref{PUnifDerivor}, $\Uf^{(n)}_{\sig{\Sigma}}(\Gd^{(n)}_{\sig{\Sigma}}(\calg{A}))$ is a $\sig{\Sigma}^{(n)}$-algebra. Therefore, the following chain of equalities holds
\allowdisplaybreaks
\begin{align*}
\Uf^{(n)}_{\sig{\Sigma}}\left(\Gd^{(n)}_{\sig{\Sigma}}(\calg{A})\right)
&=
\left(\bGd^{(n)}_{\sig{\Sigma}}\circ\bUf^{(n)}_{\sig{\Sigma}}\right)^{\ast}_{\gd}(\calg{A})
\tag{1}
\\&=
\left(\id_{\sig{\Sigma}^{(n)}}\right)^{\ast}_{\gd}(\calg{A})
\tag{2}
\\&=
\calg{A}.
\tag{3}
\end{align*}

The first equality follows from Proposition~\ref{PHepComp}~(2); the second equality follows from the above; finally, the last equality follows from Proposition~\ref{PHepId}. By the same propositions, the equality also holds for the $n$-categorial $\sig{\Sigma}$-homomorphisms. Since $\mathsf{nCatAlg}(\sig{\Sigma})=\mathsf{EPAlg}(\sig{\Sigma}^{(n)},\mathcal{C}^{\sig{\Sigma}^{(n)}})$, the stated equality of functors follows.

This completes the proof.
\end{proof}

In contrast to the previous result, the process of unification followed by that of graduation does not give rise to the identity functor on the category of globular $n$-categorial algebras, since the $k$-cells of $\Gd^{(n)}_{\sig{\Sigma}}(\Uf^{(n)}_{\sig{\Sigma}}(\calg{A}))$ are the $(n,k)$-identities of the $k$-cells of $\calg{A}$. However, we will show that both functors are naturally isomorphic.

\begin{proposition}\label{PGradUnif}
Let $\sig{\Sigma}=(S,\Sigma)$ be a many-sorted signature and $n\in\mathbb{N}$. Then $\Gd^{(n)}_{\sig{\Sigma}}\circ\Uf^{(n)}_{\sig{\Sigma}}\cong \Id_{\mathsf{nCat}^{\gl}\mathsf{Alg}(\sig{\Sigma})}$.
\end{proposition}
\begin{proof}
In this proof, for every $s\in S$, we let $\mathrm{i}^{(n,n)}_{s}$ and $\src^{(n,n)}_{s}$ stand for the identity mapping on the $n$-cells of sort $s$. Let $\calg{A}$ be a globular $n$-categorial $\sig{\Sigma}$-algebra. By Lemma~\ref{LUnifDescription}, $\src^{k}_{s}$ acts in $\Uf^{(n)}_{\sig{\Sigma}}(\calg{A})$ as $\mathrm{i}^{(n,k)}_{s}\circ\src^{(k,n)}_{s}$. Hence, by Definitions~\ref{DGradDerivor} and~\ref{DUnifDerivor}, an $n$-cell $x$ of sort $s$ satisfies $\Dom^{\Gd^{(n)}_{\sig{\Sigma}}}_{(k,s)}$ in $\Uf^{(n)}_{\sig{\Sigma}}(\calg{A})$ if, and only if, it satisfies $(\Uf^{(n)}_{\sig{\Sigma}})^{\sharp}[\Dom^{\Gd^{(n)}_{\sig{\Sigma}}}_{(k,s)}]=\{(v^{(n,s)}_{0},\mathrm{i}^{(n,k)}_{s}(\src^{(k,n)}_{s}(v^{(n,s)}_{0})))\}$ in $\calg{A}$. Thus, by the note after Proposition~\ref{PHepComp}, $\Gd^{(n)}_{\sig{\Sigma}}(\Uf^{(n)}_{\sig{\Sigma}}(\calg{A}))$ and $(\bUf^{(n)}_{\sig{\Sigma}}\circ\bGd^{(n)}_{\sig{\Sigma}})^{\ast}_{\gd}(\calg{A})$ have the same cells and, by Propositions~\ref{PGradDerivor} and~\ref{PUnifDerivor}, the former is a $\sig{\Sigma}^{\gl(n)}$-algebra. Therefore, by Proposition~\ref{PHepComp}~(2),
$$
\Gd^{(n)}_{\sig{\Sigma}}\left(\Uf^{(n)}_{\sig{\Sigma}}(\calg{A})\right)=\left(\bUf^{(n)}_{\sig{\Sigma}}\circ\bGd^{(n)}_{\sig{\Sigma}}\right)^{\ast}_{\gd}(\calg{A}).
$$
We prove the following claims.

\begin{claim}\label{CGradUnifA}
The underlying derivor of $\bUf^{(n)}_{\sig{\Sigma}}\circ\bGd^{(n)}_{\sig{\Sigma}}$ is
$$
\left(\varphi^{\Uf^{(n)}_{\sig{\Sigma}}}\circ\varphi^{\Gd^{(n)}_{\sig{\Sigma}}},\left(\Uf^{(n)}_{\sig{\Sigma}}\right)^{\sharp}\circ\Gd^{(n)}_{\sig{\Sigma}}\right),
$$
where $\varphi^{\Uf^{(n)}_{\sig{\Sigma}}}(\varphi^{\Gd^{(n)}_{\sig{\Sigma}}}(k,s))=(n,s)$ and, for every pair of dimensions $j<k$ and every $s\in S$,
\begin{align*}
\left(\Uf^{(n)}_{\sig{\Sigma}}\right)^{\sharp}\left(\Gd^{(n)}_{\sig{\Sigma}}(\sigma_{k})\right)&=\sigma_{n}((v^{(n,\mathbf{s})}_{p})_{p\in\bb{\mathbf{s}}}), &
\left(\Uf^{(n)}_{\sig{\Sigma}}\right)^{\sharp}\left(\Gd^{(n)}_{\sig{\Sigma}}(\mathrm{i}^{(k,j)}_{s})\right)&=v^{(n,s)}_{0},\\
\left(\Uf^{(n)}_{\sig{\Sigma}}\right)^{\sharp}\left(\Gd^{(n)}_{\sig{\Sigma}}(\src^{(j,k)}_{s})\right)&=\mathrm{i}^{(n,j)}_{s}(\src^{(j,n)}_{s}(v^{(n,s)}_{0})), &
\left(\Uf^{(n)}_{\sig{\Sigma}}\right)^{\sharp}\left(\Gd^{(n)}_{\sig{\Sigma}}(\circ^{(j,k)}_{s})\right)&=v^{(n,s)}_{0}\circ^{(j,n)}_{s}v^{(n,s)}_{1},
\end{align*}
and similarly for the targets.
\end{claim}

It follows from Definition~\ref{DHepComp}, (Der1) and~(Der2), and Definitions~\ref{DGradDerivor} and~\ref{DUnifDerivor}. For instance, the following chain of equalities holds
\allowdisplaybreaks
\begin{align*}
\left(\Uf^{(n)}_{\sig{\Sigma}}\right)^{\sharp}\left(\Gd^{(n)}_{\sig{\Sigma}}\left(\src^{(j,k)}_{s}\right)\right)
&=
\left(\Uf^{(n)}_{\sig{\Sigma}}\right)^{\sharp}\left(\src^{j}_{s}(v^{s}_{0})\right)
\tag{1}
\\&=
\mathrm{i}^{(n,j)}_{s}\left(\src^{(j,n)}_{s}(v^{(n,s)}_{0})\right).
\tag{2}
\end{align*}

The first equality unravels Definition~\ref{DGradDerivor}; finally, the last equality follows from~(Der1) and~(Der2) and Definition~\ref{DUnifDerivor}.

Claim~\ref{CGradUnifA} follows.

\begin{claim}\label{CGradUnifB}
For every globular $n$-categorial $\sig{\Sigma}$-algebra $\calg{A}$, the cells of sort $(k,s)$ of $\Gd^{(n)}_{\sig{\Sigma}}(\Uf^{(n)}_{\sig{\Sigma}}(\calg{A}))$ are the $n$-cells $x$ of sort $s$ of $\calg{A}$ such that $x=\mathrm{i}^{(n,k)}_{s}(\src^{(k,n)}_{s}(x))$, for $k\in n$, and all the $n$-cells of sort $s$ of $\calg{A}$, for $k=n$. Moreover, its operations are the restrictions of the partial term operations in $\calg{A}$ of the terms of Claim~\ref{CGradUnifA}.
\end{claim}

Let $\calg{A}$ be a globular $n$-categorial $\sig{\Sigma}$-algebra. By Lemmas~\ref{LGradDescription} and~\ref{LUnifDescription}, the cells of sort $(k,s)$ of $\Gd^{(n)}_{\sig{\Sigma}}(\Uf^{(n)}_{\sig{\Sigma}}(\calg{A}))$, for $k\in n$, are the $n$-cells $x$ of sort $s$ of $\calg{A}$ that are $k$-cells of $\Uf^{(n)}_{\sig{\Sigma}}(\calg{A})$, i.e., such that $x=\mathrm{i}^{(n,k)}_{s}(\src^{(k,n)}_{s}(x))$; and, for $k=n$, all the $n$-cells of sort $s$. Moreover, by the above equality and Proposition~\ref{PTransformedAlg}, applied to $\bUf^{(n)}_{\sig{\Sigma}}\circ\bGd^{(n)}_{\sig{\Sigma}}$, its operations are the restrictions of the partial term operations in $\calg{A}$ of the terms of Claim~\ref{CGradUnifA}.

Claim~\ref{CGradUnifB} follows.

Let $\calg{A}$ be a globular $n$-categorial $\sig{\Sigma}$-algebra. We let $\theta_{\calg{A}}$ stand for the family of mappings that sends, for every $k\in n+1$ and every $s\in S$, a $k$-cell $u$ of sort $s$ of $\calg{A}$ to $\mathrm{i}^{(n,k)}_{s}(u)$.

\begin{claim}\label{CGradUnifC}
$\theta_{\calg{A}}$ is a bijection from the cells of $\calg{A}$ onto the cells of $\Gd^{(n)}_{\sig{\Sigma}}(\Uf^{(n)}_{\sig{\Sigma}}(\calg{A}))$, whose inverse sends, for every $k\in n+1$ and every $s\in S$, a cell $x$ of sort $(k,s)$ to $\src^{(k,n)}_{s}(x)$.
\end{claim}

Let $k\in n$, $s\in S$ and $u$ a $k$-cell of sort $s$. Then, by Lemma~\ref{LUnifInj}, the following chain of equalities holds
\allowdisplaybreaks
\begin{align*}
\mathrm{i}^{(n,k)}_{s}\left(\src^{(k,n)}_{s}\left(\theta_{\calg{A}}(u)\right)\right)
&=
\mathrm{i}^{(n,k)}_{s}\left(\src^{(k,n)}_{s}\left(\mathrm{i}^{(n,k)}_{s}(u)\right)\right)
\tag{1}
\\&=
\mathrm{i}^{(n,k)}_{s}(u)
\tag{2}
\\&=
\theta_{\calg{A}}(u).
\tag{3}
\end{align*}

Thus, by Claim~\ref{CGradUnifB}, $\theta_{\calg{A}}(u)$ is a cell of sort $(k,s)$ of $\Gd^{(n)}_{\sig{\Sigma}}(\Uf^{(n)}_{\sig{\Sigma}}(\calg{A}))$. Moreover, by Lemma~\ref{LUnifInj}, $\src^{(k,n)}_{s}(\theta_{\calg{A}}(u))=u$, and, for every cell $x$ of sort $(k,s)$ of $\Gd^{(n)}_{\sig{\Sigma}}(\Uf^{(n)}_{\sig{\Sigma}}(\calg{A}))$, $\theta_{\calg{A}}(\src^{(k,n)}_{s}(x))=\mathrm{i}^{(n,k)}_{s}(\src^{(k,n)}_{s}(x))=x$. For $k=n$, both mappings are identities.

Claim~\ref{CGradUnifC} follows.

\begin{claim}\label{CGradUnifD}
$\theta_{\calg{A}}$ is a homomorphism of partial $\Sigma^{\gl(n)}$-algebras from $\calg{A}$ to $\Gd^{(n)}_{\sig{\Sigma}}(\Uf^{(n)}_{\sig{\Sigma}}(\calg{A}))$ that reflects definedness. Hence it is an isomorphism.
\end{claim}

By Claims~\ref{CGradUnifA} and~\ref{CGradUnifB}, the operations of $\Gd^{(n)}_{\sig{\Sigma}}(\Uf^{(n)}_{\sig{\Sigma}}(\calg{A}))$ are the restrictions of the partial term operations in $\calg{A}$ of the terms $(\Uf^{(n)}_{\sig{\Sigma}})^{\sharp}(\Gd^{(n)}_{\sig{\Sigma}}(\tau))$, for $\tau$ an operation symbol of $\Sigma^{\gl(n)}$. Let $j<k$ be dimensions and $s\in S$. We consider the different kinds of operation symbols of $\Sigma^{\gl(n)}$.

\textsf{(Operation symbols $\sigma_{k}$)} Let $(\mathbf{s},s)\in S^{\star}\times S$, $\sigma\in\Sigma_{\mathbf{s},s}$ and $(u_{p})_{p\in\bb{\mathbf{s}}}$ a family of $k$-cells. Then the following chain of equalities holds
\allowdisplaybreaks
\begin{align*}
\theta_{\calg{A}}\left(\sigma_{k}\left((u_{p})_{p\in\bb{\mathbf{s}}}\right)\right)
&=
\mathrm{i}^{(n,k)}_{s}\left(\sigma_{k}\left((u_{p})_{p\in\bb{\mathbf{s}}}\right)\right)
\tag{1}
\\&=
\sigma_{n}\left(\left(\mathrm{i}^{(n,k)}_{s_{p}}(u_{p})\right)_{p\in\bb{\mathbf{s}}}\right)
\tag{2}
\\&=
\sigma_{k}^{\Gd^{(n)}_{\sig{\Sigma}}(\Uf^{(n)}_{\sig{\Sigma}}(\calg{A}))}\left(\left(\theta_{\calg{A}}(u_{p})\right)_{p\in\bb{\mathbf{s}}}\right).
\tag{3}
\end{align*}

The first equality unravels the definition of $\theta_{\calg{A}}$; the second equality follows from the third equation of~(EA$^{k,n}$7); finally, the last equality follows from Claims~\ref{CGradUnifA} and~\ref{CGradUnifB}.

\textsf{(Sources and targets)} Let $u$ be a $k$-cell of sort $s$. Then the following chain of equalities holds
\allowdisplaybreaks
\begin{align*}
\theta_{\calg{A}}\left(\src^{(j,k)}_{s}(u)\right)
&=
\mathrm{i}^{(n,j)}_{s}\left(\src^{(j,k)}_{s}(u)\right)
\tag{1}
\\&=
\mathrm{i}^{(n,j)}_{s}\left(\src^{(j,k)}_{s}\left(\src^{(k,n)}_{s}\left(\mathrm{i}^{(n,k)}_{s}(u)\right)\right)\right)
\tag{2}
\\&=
\mathrm{i}^{(n,j)}_{s}\left(\src^{(j,n)}_{s}\left(\mathrm{i}^{(n,k)}_{s}(u)\right)\right)
\tag{3}
\\&=
\src^{(j,k)\Gd^{(n)}_{\sig{\Sigma}}(\Uf^{(n)}_{\sig{\Sigma}}(\calg{A}))}_{s}\left(\theta_{\calg{A}}(u)\right).
\tag{4}
\end{align*}

The first equality unravels the definition of $\theta_{\calg{A}}$; the second equality follows from Lemma~\ref{LUnifInj}; the third equality follows from~(EA$^{j,k,n}$1); finally, the last equality follows from Claims~\ref{CGradUnifA} and~\ref{CGradUnifB}. By a similar argument, $\theta_{\calg{A}}$ commutes with the $(j,k)$-targets.

\textsf{(Identities)} Let $w$ be a $j$-cell of sort $s$. Then the following chain of equalities holds
\allowdisplaybreaks
\begin{align*}
\theta_{\calg{A}}\left(\mathrm{i}^{(k,j)}_{s}(w)\right)
&=
\mathrm{i}^{(n,k)}_{s}\left(\mathrm{i}^{(k,j)}_{s}(w)\right)
\tag{1}
\\&=
\mathrm{i}^{(n,j)}_{s}(w)
\tag{2}
\\&=
\mathrm{i}^{(k,j)\Gd^{(n)}_{\sig{\Sigma}}(\Uf^{(n)}_{\sig{\Sigma}}(\calg{A}))}_{s}\left(\theta_{\calg{A}}(w)\right).
\tag{3}
\end{align*}

The first equality unravels the definition of $\theta_{\calg{A}}$; the second equality follows from~(EA$^{j,k,n}$2); finally, the last equality follows from Claims~\ref{CGradUnifA} and~\ref{CGradUnifB}, since $(\Uf^{(n)}_{\sig{\Sigma}})^{\sharp}(\Gd^{(n)}_{\sig{\Sigma}}(\mathrm{i}^{(k,j)}_{s}))$ is a variable.

\textsf{(Compositions)} Let $u$, $u'$ be $k$-cells of sort $s$. By the computation of the $(j,k)$-sources and the $(j,k)$-targets and by Lemma~\ref{LUnifInj}, $\theta_{\calg{A}}(u)\circ^{(j,k)}_{s}\theta_{\calg{A}}(u')$ is defined in $\Gd^{(n)}_{\sig{\Sigma}}(\Uf^{(n)}_{\sig{\Sigma}}(\calg{A}))$ if, and only if, $\mathrm{i}^{(n,j)}_{s}(\src^{(j,k)}_{s}(u))=\mathrm{i}^{(n,j)}_{s}(\tgt^{(j,k)}_{s}(u'))$, if, and only if, $\src^{(j,k)}_{s}(u)=\tgt^{(j,k)}_{s}(u')$, i.e., by~(EA$^{j,k}$3), if, and only if, $u\circ^{(j,k)}_{s}u'$ is defined. In this case, the following chain of equalities holds
\allowdisplaybreaks
\begin{align*}
\theta_{\calg{A}}\left(u\circ^{(j,k)}_{s}u'\right)
&=
\mathrm{i}^{(n,k)}_{s}\left(u\circ^{(j,k)}_{s}u'\right)
\tag{1}
\\&=
\mathrm{i}^{(n,k)}_{s}(u)\circ^{(j,n)}_{s}\mathrm{i}^{(n,k)}_{s}(u')
\tag{2}
\\&=
\theta_{\calg{A}}(u)\circ^{(j,k)\Gd^{(n)}_{\sig{\Sigma}}(\Uf^{(n)}_{\sig{\Sigma}}(\calg{A}))}_{s}\theta_{\calg{A}}(u').
\tag{3}
\end{align*}

The first equality unravels the definition of $\theta_{\calg{A}}$; the second equality follows from~(EA$^{j,k,n}$5); finally, the last equality follows from Claims~\ref{CGradUnifA} and~\ref{CGradUnifB}.

Therefore, $\theta_{\calg{A}}$ is a homomorphism that reflects definedness and, by Claim~\ref{CGradUnifC}, it is bijective. Hence it is an isomorphism.

Claim~\ref{CGradUnifD} follows.

\begin{claim}\label{CGradUnifE}
The family $\theta=(\theta_{\calg{A}})_{\calg{A}\in\mathrm{Ob}(\mathsf{nCat}^{\gl}\mathsf{Alg}(\sig{\Sigma}))}$ is a natural transformation from $\Id_{\mathsf{nCat}^{\gl}\mathsf{Alg}(\sig{\Sigma})}$ to $\Gd^{(n)}_{\sig{\Sigma}}\circ\Uf^{(n)}_{\sig{\Sigma}}$.
\end{claim}

Let $F\colon\calg{A}\mor\calg{A}'$ be a globular $n$-categorial $\sig{\Sigma}$-homomorphism. By Claim~\ref{CGradUnifB} and Proposition~\ref{PGenHepDerivorFunctor}, $\Gd^{(n)}_{\sig{\Sigma}}(\Uf^{(n)}_{\sig{\Sigma}}(F))$ acts on the cells of sort $(k,s)$ as $F_{(n,s)}$. Then, for every $k$-cell $u$ of sort $s$, the following chain of equalities holds
\allowdisplaybreaks
\begin{align*}
\Gd^{(n)}_{\sig{\Sigma}}\left(\Uf^{(n)}_{\sig{\Sigma}}(F)\right)_{(k,s)}\left(\theta_{\calg{A}}(u)\right)
&=
F_{(n,s)}\left(\mathrm{i}^{(n,k)}_{s}(u)\right)
\tag{1}
\\&=
\mathrm{i}^{(n,k)}_{s}\left(F_{(k,s)}(u)\right)
\tag{2}
\\&=
\theta_{\calg{A}'}\left(F_{(k,s)}(u)\right).
\tag{3}
\end{align*}

The first equality unravels the definitions of $\Gd^{(n)}_{\sig{\Sigma}}(\Uf^{(n)}_{\sig{\Sigma}}(F))$ and of $\theta_{\calg{A}}$; the second equality follows from the fact that $F$ commutes with the identities, see Remark~\ref{RGlobularIdent}; finally, the last equality unravels the definition of $\theta_{\calg{A}'}$.

Claim~\ref{CGradUnifE} follows.

All in all, by Claims~\ref{CGradUnifD} and~\ref{CGradUnifE}, $\theta$ is a natural isomorphism from $\Id_{\mathsf{nCat}^{\gl}\mathsf{Alg}(\sig{\Sigma})}$ to $\Gd^{(n)}_{\sig{\Sigma}}\circ\Uf^{(n)}_{\sig{\Sigma}}$.

This completes the proof.
\end{proof}

From Propositions~\ref{PUnifGrad} and~\ref{PGradUnif} we obtain the following corollary.

\begin{corollary}\label{CGdUfEquiv}
Let $\sig{\Sigma}=(S,\Sigma)$ be a many-sorted signature and $n\in\mathbb{N}$. Then $\Gd^{(n)}_{\sig{\Sigma}}$ and $\Uf^{(n)}_{\sig{\Sigma}}$ are mutually quasi-inverse equivalences of categories between $\mathsf{nCatAlg}(\sig{\Sigma})$ and $\mathsf{nCat}^{\gl}\mathsf{Alg}(\sig{\Sigma})$.
\end{corollary}

We will use the following, more precise, form of Corollary~\ref{CGdUfEquiv} in Section~\ref{SAlgSquare}.

\begin{proposition}\label{PGdUfAdj}
Let $\sig{\Sigma}=(S,\Sigma)$ be a many-sorted signature, $n\in\mathbb{N}$ and $\theta=(\theta_{\calg{A}})_{\calg{A}\in\mathrm{Ob}(\mathsf{nCat}^{\gl}\mathsf{Alg}(\sig{\Sigma}))}$ the natural isomorphism from $\Id_{\mathsf{nCat}^{\gl}\mathsf{Alg}(\sig{\Sigma})}$ to $\Gd^{(n)}_{\sig{\Sigma}}\circ\Uf^{(n)}_{\sig{\Sigma}}$ of the proof of Proposition~\ref{PGradUnif}. Then $(\Uf^{(n)}_{\sig{\Sigma}},\Gd^{(n)}_{\sig{\Sigma}},\id,\theta^{-1})$ is an adjunction from $\mathsf{nCatAlg}(\sig{\Sigma})$ to $\mathsf{nCat}^{\gl}\mathsf{Alg}(\sig{\Sigma})$, see Definition~\ref{DAdj}, i.e., $\Gd^{(n)}_{\sig{\Sigma}}\dashv\Uf^{(n)}_{\sig{\Sigma}}$, with unit the identity and counit $\theta^{-1}=(\theta^{-1}_{\calg{A}})_{\calg{A}\in\mathrm{Ob}(\mathsf{nCat}^{\gl}\mathsf{Alg}(\sig{\Sigma}))}$. Thus it is an adjoint equivalence.
\end{proposition}
\begin{proof}
By Proposition~\ref{PUnifGrad}, $\Uf^{(n)}_{\sig{\Sigma}}\circ\Gd^{(n)}_{\sig{\Sigma}}=\Id_{\mathsf{nCatAlg}(\sig{\Sigma})}$, so that the identity $\id=(\id_{\calg{B}})_{\calg{B}\in\mathrm{Ob}(\mathsf{nCatAlg}(\sig{\Sigma}))}$ is a natural transformation from $\Id_{\mathsf{nCatAlg}(\sig{\Sigma})}$ to $\Uf^{(n)}_{\sig{\Sigma}}\circ\Gd^{(n)}_{\sig{\Sigma}}$; and, by Proposition~\ref{PGradUnif}, $\theta^{-1}$ is a natural transformation from $\Gd^{(n)}_{\sig{\Sigma}}\circ\Uf^{(n)}_{\sig{\Sigma}}$ to $\Id_{\mathsf{nCat}^{\gl}\mathsf{Alg}(\sig{\Sigma})}$. By Definition~\ref{DAdj}, we have to prove the triangular identities
$$
\left(\Uf^{(n)}_{\sig{\Sigma}}\ast\theta^{-1}\right)\circ\left(\id\ast\Uf^{(n)}_{\sig{\Sigma}}\right)=\id_{\Uf^{(n)}_{\sig{\Sigma}}}
\qquad\text{and}\qquad
\left(\theta^{-1}\ast\Gd^{(n)}_{\sig{\Sigma}}\right)\circ\left(\Gd^{(n)}_{\sig{\Sigma}}\ast\id\right)=\id_{\Gd^{(n)}_{\sig{\Sigma}}}.
$$
Let us recall that, by Claim~\ref{CGradUnifC} in the proof of Proposition~\ref{PGradUnif}, for every globular $n$-categorial $\sig{\Sigma}$-algebra $\calg{A}$, every $k\in n+1$ and every $s\in S$, the component of $\theta^{-1}_{\calg{A}}$ at the sort $(k,s)$ sends a cell $x$ of sort $(k,s)$ of $\Gd^{(n)}_{\sig{\Sigma}}(\Uf^{(n)}_{\sig{\Sigma}}(\calg{A}))$ to $\src^{(k,n)}_{s}(x)$, where $\src^{(n,n)}_{s}$ is the identity mapping on the $n$-cells of sort $s$.

\textsf{(First triangular identity)} Let $\calg{A}$ be a globular $n$-categorial $\sig{\Sigma}$-algebra. Since $\Uf^{(n)}_{\sig{\Sigma}}(\calg{A})$ is an $n$-categorial $\sig{\Sigma}$-algebra, by Proposition~\ref{PUnifGrad}, the domain $\Uf^{(n)}_{\sig{\Sigma}}(\Gd^{(n)}_{\sig{\Sigma}}(\Uf^{(n)}_{\sig{\Sigma}}(\calg{A})))$ of $\Uf^{(n)}_{\sig{\Sigma}}(\theta^{-1}_{\calg{A}})$ is $\Uf^{(n)}_{\sig{\Sigma}}(\calg{A})$, whose cells of sort $s$ are, by Lemma~\ref{LUnifDescription}, the $n$-cells of sort $s$ of $\calg{A}$. Let $s\in S$ and $x$ an $n$-cell of sort $s$ of $\calg{A}$. Then the following chain of equalities holds
\allowdisplaybreaks
\begin{align*}
\left(\left(\left(\Uf^{(n)}_{\sig{\Sigma}}\ast\theta^{-1}\right)\circ\left(\id\ast\Uf^{(n)}_{\sig{\Sigma}}\right)\right)_{\calg{A}}\right)_{s}(x)
&=
\left(\Uf^{(n)}_{\sig{\Sigma}}\left(\theta^{-1}_{\calg{A}}\right)\circ\id_{\Uf^{(n)}_{\sig{\Sigma}}(\calg{A})}\right)_{s}(x)
\tag{1}
\\&=
\Uf^{(n)}_{\sig{\Sigma}}\left(\theta^{-1}_{\calg{A}}\right)_{s}(x)
\tag{2}
\\&=
\left(\theta^{-1}_{\calg{A}}\right)_{(n,s)}(x)
\tag{3}
\\&=
\src^{(n,n)}_{s}(x)
\tag{4}
\\&=
x
\tag{5}
\\&=
\left(\id_{\Uf^{(n)}_{\sig{\Sigma}}}\right)_{\calg{A},s}(x).
\tag{6}
\end{align*}

The first equality unravels the definitions of the horizontal and vertical composition of natural transformations; the second equality follows from the fact that $\id_{\Uf^{(n)}_{\sig{\Sigma}}(\calg{A})}$ is the identity homomorphism; the third equality follows from Proposition~\ref{PGenHepDerivorFunctor}, since $\varphi^{\Uf^{(n)}_{\sig{\Sigma}}}(s)=(n,s)$, see Definition~\ref{DUnifDerivor}, so that $\Uf^{(n)}_{\sig{\Sigma}}(\theta^{-1}_{\calg{A}})_{s}$ is the birestriction of $(\theta^{-1}_{\calg{A}})_{(n,s)}$ to the cells; the fourth equality follows from Claim~\ref{CGradUnifC}; the fifth equality follows from the fact that $\src^{(n,n)}_{s}$ is the identity mapping on the $n$-cells of sort $s$; finally, the last equality unravels the definition of the identity natural transformation at $\Uf^{(n)}_{\sig{\Sigma}}$.

Since a homomorphism is determined by its domain, its codomain and its underlying many-sorted mapping, the first triangular identity holds.

\textsf{(Second triangular identity)} Let $\calg{B}$ be an $n$-categorial $\sig{\Sigma}$-algebra. By Proposition~\ref{PUnifGrad}, the codomain $\Gd^{(n)}_{\sig{\Sigma}}(\Uf^{(n)}_{\sig{\Sigma}}(\Gd^{(n)}_{\sig{\Sigma}}(\calg{B})))$ of $\theta_{\Gd^{(n)}_{\sig{\Sigma}}(\calg{B})}$ is $\Gd^{(n)}_{\sig{\Sigma}}(\calg{B})$. We first prove that $\theta_{\Gd^{(n)}_{\sig{\Sigma}}(\calg{B})}$ is the identity homomorphism of $\Gd^{(n)}_{\sig{\Sigma}}(\calg{B})$. Let $k\in n+1$, $s\in S$ and $u$ a cell of sort $(k,s)$ of $\Gd^{(n)}_{\sig{\Sigma}}(\calg{B})$. If $k=n$, then, by definition, $(\theta_{\Gd^{(n)}_{\sig{\Sigma}}(\calg{B})})_{(n,s)}(u)=\mathrm{i}^{(n,n)}_{s}(u)=u$. If $k\in n$, then the following chain of equalities holds
\allowdisplaybreaks
\begin{align*}
\left(\theta_{\Gd^{(n)}_{\sig{\Sigma}}(\calg{B})}\right)_{(k,s)}(u)
&=
\mathrm{i}^{(n,k)\,\Gd^{(n)}_{\sig{\Sigma}}(\calg{B})}_{s}(u)
\tag{1}
\\&=
u.
\tag{2}
\end{align*}

The first equality unravels the definition of $\theta_{\Gd^{(n)}_{\sig{\Sigma}}(\calg{B})}$, see the proof of Proposition~\ref{PGradUnif}; finally, the last equality follows from Lemma~\ref{LGradDescription}~(1), since the $(n,k)$-identities of $\Gd^{(n)}_{\sig{\Sigma}}(\calg{B})$ act as inclusions.

Hence $\theta_{\Gd^{(n)}_{\sig{\Sigma}}(\calg{B})}=\id_{\Gd^{(n)}_{\sig{\Sigma}}(\calg{B})}$ and, consequently, $\theta^{-1}_{\Gd^{(n)}_{\sig{\Sigma}}(\calg{B})}=\id_{\Gd^{(n)}_{\sig{\Sigma}}(\calg{B})}$. Therefore the following chain of equalities holds
\allowdisplaybreaks
\begin{align*}
\left(\left(\theta^{-1}\ast\Gd^{(n)}_{\sig{\Sigma}}\right)\circ\left(\Gd^{(n)}_{\sig{\Sigma}}\ast\id\right)\right)_{\calg{B}}
&=
\theta^{-1}_{\Gd^{(n)}_{\sig{\Sigma}}(\calg{B})}\circ\Gd^{(n)}_{\sig{\Sigma}}\left(\id_{\calg{B}}\right)
\tag{1}
\\&=
\id_{\Gd^{(n)}_{\sig{\Sigma}}(\calg{B})}\circ\id_{\Gd^{(n)}_{\sig{\Sigma}}(\calg{B})}
\tag{2}
\\&=
\left(\id_{\Gd^{(n)}_{\sig{\Sigma}}}\right)_{\calg{B}}.
\tag{3}
\end{align*}

The first equality unravels the definitions of the horizontal and vertical composition of natural transformations; the second equality follows from the above and from the fact that $\Gd^{(n)}_{\sig{\Sigma}}$ preserves identities; finally, the last equality unravels the definition of the identity natural transformation at $\Gd^{(n)}_{\sig{\Sigma}}$.

Thus the second triangular identity holds. Finally, since the unit and the counit are natural isomorphisms, the adjunction is an adjoint equivalence.

This completes the proof.
\end{proof}

\section{Discretization and Truncation of Higher-Order Categorial Algebras}\label{STruncation}

In this section we move between dimensions. Let $\sig{\Sigma}=(S,\Sigma)$ be a many-sorted signature and $m,n\in\mathbb{N}$ with $m\leq n$. The discretization functor regards an $m$-categorial $\sig{\Sigma}$-algebra as an $n$-categorial $\sig{\Sigma}$-algebra whose cells of dimension greater than $m$ are identities, and the truncation functor forgets the cells of dimension greater than $m$ of an $n$-categorial $\sig{\Sigma}$-algebra. As in Section~\ref{SGraduation}, both functors are induced by generalized hep-derivors, in the single-set and in the globular presentations. We begin by defining the discretization functors and the truncation functors. Next we prove that discretization is left adjoint to truncation and that these functors compose as expected. We end this section by showing that they are compatible with graduation and unification. We use the notation for generalized hep-derivors introduced at the beginning of Section~\ref{SGraduation}.

\subsection{Discretization of higher-order categorial algebras}\label{SDiscretization}

We begin by defining the single-set discretization derivor.

\begin{definition}\label{DDiscDerivor}
Let $\sig{\Sigma}=(S,\Sigma)$ be a many-sorted signature and $m,n\in\mathbb{N}$ with $m<n$. The \emph{single-set discretization derivor} is the generalized hep-derivor
$$
\bDs^{(n,m)}_{\sig{\Sigma}}=\left(\varphi^{\Ds^{(n,m)}_{\sig{\Sigma}}},\Ds^{(n,m)}_{\sig{\Sigma}},\Dom^{\Ds^{(n,m)}_{\sig{\Sigma}}}\right)
$$
from $\sig{\Sigma}^{(n)}$ to $\sig{\Sigma}^{(m)}$ defined as follows:
\begin{enumerate}
\item the mapping on sorts is $\varphi^{\Ds^{(n,m)}_{\sig{\Sigma}}}=\id_{S}$;
\item for every $(\mathbf{s},s)\in S^{\star}\times S$ and every $\sigma\in\Sigma_{\mathbf{s},s}$,
$$
\Ds^{(n,m)}_{\sig{\Sigma}}(\sigma)=\sigma((v^{\mathbf{s}}_{p})_{p\in\bb{\mathbf{s}}});
$$
\item for every $k\in m$ and every $s\in S$,
$$
\Ds^{(n,m)}_{\sig{\Sigma}}(\src^{k}_{s})=\src^{k}_{s}(v^{s}_{0}),\quad
\Ds^{(n,m)}_{\sig{\Sigma}}(\tgt^{k}_{s})=\tgt^{k}_{s}(v^{s}_{0}),\quad
\Ds^{(n,m)}_{\sig{\Sigma}}(\circ^{k}_{s})=v^{s}_{0}\circ^{k}_{s}v^{s}_{1};
$$
\item for every $k\in n$ with $k\geq m$ and every $s\in S$,
$$
\Ds^{(n,m)}_{\sig{\Sigma}}(\src^{k}_{s})=\Ds^{(n,m)}_{\sig{\Sigma}}(\tgt^{k}_{s})=\Ds^{(n,m)}_{\sig{\Sigma}}(\circ^{k}_{s})=v^{s}_{0};
$$
and
\item for every $s\in S$, the domain is $\Dom^{\Ds^{(n,m)}_{\sig{\Sigma}}}_{s}=\varnothing$.
\end{enumerate}
\end{definition}

Let us note that, for $k\geq m$, $\Ds^{(n,m)}_{\sig{\Sigma}}(\circ^{k}_{s})$ is total but, in the transformed algebra, $\src^{k}_{s}$ and $\tgt^{k}_{s}$ act as identities, so that two cells satisfy the domain condition $\src^{k}_{s}(v^{s}_{0})=\tgt^{k}_{s}(v^{s}_{1})$ of $\circ^{k}_{s}$ if, and only if, they are equal, see Proposition~\ref{PTransformedAlg}. Thus the $k$-composition of two cells is defined if, and only if, they are equal. No term of $\Sigma^{(m)}$ could express this condition: in a monoid, regarded as a single-set category with a single identity, every term operation is total.

\begin{definition}\label{DDiscretization}
Let $\sig{\Sigma}=(S,\Sigma)$ be a many-sorted signature and $m,n\in\mathbb{N}$ with $m\leq n$. The \emph{discretization functor} $\Ds^{(n,m)}_{\sig{\Sigma}}$ is the functor $(\bDs^{(n,m)}_{\sig{\Sigma}})^{\ast}_{\gd}$ induced by the generalized hep-derivor $\bDs^{(n,m)}_{\sig{\Sigma}}$, see Definitions~\ref{DDiscDerivor} and~\ref{DDiscGlDerivor}, from $\mathsf{PAlg}(\Sigma^{(m)})$ to $\mathsf{PAlg}(\Sigma^{(n)})$, see Proposition~\ref{PGenHepDerivorFunctor}.
\end{definition}

In the following lemma we describe the transformed algebra along the single-set discretization derivor.

\begin{lemma}\label{LDiscDescription}
Let $\sig{\Sigma}=(S,\Sigma)$ be a many-sorted signature, $m,n\in\mathbb{N}$ with $m<n$ and $\mathbf{B}$ an object in $\mathsf{EPAlg}(\sig{\Sigma}^{(m)},\mathcal{C}^{\sig{\Sigma}^{(m)}})$. Then $\Ds^{(n,m)}_{\sig{\Sigma}}(\mathbf{B})$ is a $\sig{\Sigma}^{(n)}$-algebra whose cells are all the cells of $\mathbf{B}$. Moreover, for every $s\in S$, the following statements hold:
\begin{enumerate}
\item the operations of $\Sigma$ and, for $k<m$, the $k$-sources, the $k$-targets and the $k$-compositions of $\Ds^{(n,m)}_{\sig{\Sigma}}(\mathbf{B})$ are those of $\mathbf{B}$; and
\item for $m\leq k<n$, $\src^{k}_{s}$ and $\tgt^{k}_{s}$ are identities in $\Ds^{(n,m)}_{\sig{\Sigma}}(\mathbf{B})$, and $x\circ^{k}_{s}y$ is defined in $\Ds^{(n,m)}_{\sig{\Sigma}}(\mathbf{B})$ if, and only if, $x=y$, and, in this case, the $k$-composition of $x$ and $y$ in $\Ds^{(n,m)}_{\sig{\Sigma}}(\mathbf{B})$ is $x$.
\end{enumerate}
\end{lemma}
\begin{proof}
Since the domains of $\bDs^{(n,m)}_{\sig{\Sigma}}$ are empty, the cells of $\Ds^{(n,m)}_{\sig{\Sigma}}(\mathbf{B})$ are all the cells of $\mathbf{B}$. By Proposition~\ref{PTransformedEpAlg}, we check that, for every operation symbol of $\Sigma^{(n)}$, the corresponding partial term operation is defined on the families that satisfy its domain conditions in $\Ds^{(n,m)}_{\sig{\Sigma}}(\mathbf{B})$.

\textsf{(Operation symbols with empty domain conditions)} They are sent to total derived operations, which act as $\sigma^{\mathbf{B}}$, $\src^{k\,\mathbf{B}}_{s}$ and $\tgt^{k\,\mathbf{B}}_{s}$, for $k<m$, by~(A0) and~(A$^{k}$1), and as identities, for $k\geq m$.

\textsf{(Compositions of dimension $k<m$)} By the above, two cells $x$, $y$ satisfy the domain condition of $\circ^{k}_{s}$ in $\Ds^{(n,m)}_{\sig{\Sigma}}(\mathbf{B})$ if, and only if, $\src^{k\,\mathbf{B}}_{s}(x)=\tgt^{k\,\mathbf{B}}_{s}(y)$, in which case, by~(A$^{k}$3), $x\circ^{k}_{s}y$ is defined in $\mathbf{B}$.

\textsf{(Compositions of dimension $k\geq m$)} By the above, two cells $x$, $y$ satisfy the domain condition of $\circ^{k}_{s}$ in $\Ds^{(n,m)}_{\sig{\Sigma}}(\mathbf{B})$ if, and only if, $x=y$, in which case the term $\Ds^{(n,m)}_{\sig{\Sigma}}(\circ^{k}_{s})=v^{s}_{0}$ is defined, with value $x$.

In all cases, the values are cells of $\mathbf{B}$, i.e., cells of $\Ds^{(n,m)}_{\sig{\Sigma}}(\mathbf{B})$. All in all, by Proposition~\ref{PTransformedEpAlg}, $\Ds^{(n,m)}_{\sig{\Sigma}}(\mathbf{B})$ is a $\sig{\Sigma}^{(n)}$-algebra and, by Proposition~\ref{PTransformedAlg}, its operations are as stated.

This completes the proof.
\end{proof}

We next prove that the discretization of a single-set $m$-categorial $\sig{\Sigma}$-algebra is a single-set $n$-categorial $\sig{\Sigma}$-algebra.

\begin{proposition}\label{PDiscDerivor}
Let $\sig{\Sigma}=(S,\Sigma)$ be a many-sorted signature and $m,n\in\mathbb{N}$ with $m\leq n$. Then the restriction of $\Ds^{(n,m)}_{\sig{\Sigma}}$ to $\mathsf{EPAlg}(\sig{\Sigma}^{(m)},\mathcal{C}^{\sig{\Sigma}^{(m)}})$ corestricts to $\mathsf{EPAlg}(\sig{\Sigma}^{(n)},\mathcal{C}^{\sig{\Sigma}^{(n)}})$.
\end{proposition}
\begin{proof}
For $m=n$ this follows from Proposition~\ref{PHepId}. Let $m<n$ and $\mathbf{B}$ an object in $\mathsf{EPAlg}(\sig{\Sigma}^{(m)},\mathcal{C}^{\sig{\Sigma}^{(m)}})$. By Definitions~\ref{DCorestriction} and~\ref{DEpVariety}, we have to prove that $\Ds^{(n,m)}_{\sig{\Sigma}}(\mathbf{B})$ is a $\sig{\Sigma}^{(n)}$-algebra in which every axiom of $\mathcal{C}^{\sig{\Sigma}^{(n)}}$ holds.

\textsf{(Ep-algebra)} By Lemma~\ref{LDiscDescription}, which rests on Proposition~\ref{PTransformedEpAlg}, $\Ds^{(n,m)}_{\sig{\Sigma}}(\mathbf{B})$ is a $\sig{\Sigma}^{(n)}$-algebra.

\textsf{(Axioms)} By Lemma~\ref{LDiscDescription}, the axioms of $\boldsymbol{\mathcal{C}}^{\sig{\Sigma}^{(n)}}$ in $\Ds^{(n,m)}_{\sig{\Sigma}}(\mathbf{B})$ become statements on $\mathbf{B}$, which we check according to the dimensions involved. In the following computations, the operation symbols are interpreted in $\Ds^{(n,m)}_{\sig{\Sigma}}(\mathbf{B})$.

\textsf{(Dimensions smaller than $m$)} The axioms of $\boldsymbol{\mathcal{C}}^{\sig{\Sigma}^{(n)}}$ in which all the dimensions are smaller than $m$ are axioms of $\boldsymbol{\mathcal{C}}^{\sig{\Sigma}^{(m)}}$, which hold in $\mathbf{B}$.

\textsf{(Axioms (A$^{k}$1)--(A$^{k}$8) for $k\geq m$)} They become trivial. For instance, (A$^{k}$5) becomes $x=x$, (A$^{k}$3) becomes $x=y\leftrightarrow x=y$, and (A$^{k}$8) becomes the statement that, if $x_{p}=y_{p}$ for every $p$, then $\sigma((x_{p})_{p})=\sigma((y_{p})_{p})$.

\textsf{(Axioms (AB$^{l,k}$1)--(AB$^{l,k}$3) for $l<m\leq k$)} They become trivial. For instance, for $x$, $y$ such that $x\circ^{l}_{s}y$ is defined, the following chain of equalities holds
\allowdisplaybreaks
\begin{align*}
\src^{k}_{s}\left(x\circ^{l}_{s}y\right)
&=
x\circ^{l}_{s}y
\tag{1}
\\&=
\src^{k}_{s}(x)\circ^{l}_{s}\src^{k}_{s}(y).
\tag{2}
\end{align*}

Both equalities follow from Lemma~\ref{LDiscDescription}, since $\src^{k}_{s}$ acts as the identity. Moreover, the premises of (AB$^{l,k}$3) become that $x\circ^{l}_{s}y$ and $z\circ^{l}_{s}t$ are defined, $x=z$ and $y=t$, and then its conclusion becomes $x\circ^{l}_{s}y=x\circ^{l}_{s}y$.

\textsf{(Axioms (AB$^{l,k}$1)--(AB$^{l,k}$3) for $m\leq l<k$)} They become trivial, since all the premises become equalities between the variables involved.

All in all, every axiom of $\boldsymbol{\mathcal{C}}^{\sig{\Sigma}^{(n)}}$ holds in $\Ds^{(n,m)}_{\sig{\Sigma}}(\mathbf{B})$, so that $\Ds^{(n,m)}_{\sig{\Sigma}}(\mathbf{B})$ is an object in $\mathsf{EPAlg}(\sig{\Sigma}^{(n)},\mathcal{C}^{\sig{\Sigma}^{(n)}})$.

This completes the proof.
\end{proof}

Let $\sig{\Sigma}=(S,\Sigma)$ be a many-sorted signature and $m,n\in\mathbb{N}$ with $m\leq n$. By Remark~\ref{RSingleSetIdent} and Proposition~\ref{PDiscDerivor}, the restriction of $\Ds^{(n,m)}_{\sig{\Sigma}}$ to $\mathsf{mCatAlg}(\sig{\Sigma})$ corestricts to $\mathsf{nCatAlg}(\sig{\Sigma})$, and, by abuse of notation, we denote the corestriction by
$$
\Ds^{(n,m)}_{\sig{\Sigma}}\colon\mathsf{mCatAlg}(\sig{\Sigma})\mor\mathsf{nCatAlg}(\sig{\Sigma}).
$$
By Lemma~\ref{LDiscDescription}, it sends an $m$-categorial $\sig{\Sigma}$-algebra $\calg{A}$ to the $n$-categorial $\sig{\Sigma}$-algebra with the same cells, the same structure in the dimensions $k<m$, and in which, for $k\geq m$, $\src^{k}_{s}$ and $\tgt^{k}_{s}$ are identities, so that all the cells are $k$-cells, and $x\circ^{k}_{s}y$ is defined if, and only if, $x=y$, and then it is $x$; and it sends a homomorphism to itself. By Proposition~\ref{PHepId}, $\Ds^{(n,n)}_{\sig{\Sigma}}$ is the identity.

In the globular presentation, a cell of dimension $k$ of the discretization is a cell of dimension $k$ of the given algebra if $k\leq m$, and a cell of dimension $m$ if $k>m$. Accordingly, every dimension $k\in n+1$ of $\Sigma^{\gl(n)}$ is sent to the dimension $\min\{k,m\}\in m+1$ of $\Sigma^{\gl(m)}$. We next define the globular discretization derivor.

\begin{definition}\label{DDiscGlDerivor}
Let $\sig{\Sigma}=(S,\Sigma)$ be a many-sorted signature and $m,n\in\mathbb{N}$ with $m<n$. The \emph{globular discretization derivor} is the generalized hep-derivor
$$
\bDs^{\gl(n,m)}_{\sig{\Sigma}}=\left(\varphi^{\Ds^{\gl(n,m)}_{\sig{\Sigma}}},\Ds^{\gl(n,m)}_{\sig{\Sigma}},\Dom^{\Ds^{\gl(n,m)}_{\sig{\Sigma}}}\right)
$$
from $\sig{\Sigma}^{\gl(n)}$ to $\sig{\Sigma}^{\gl(m)}$ defined as follows:
\begin{enumerate}
\item for every dimension $k\in n+1$ and every $s\in S$, the mapping on sorts sends $(k,s)$ to
$$
\varphi^{\Ds^{\gl(n,m)}_{\sig{\Sigma}}}(k,s)=(\min\{k,m\},s);
$$
\item for every $(\mathbf{s},s)\in S^{\star}\times S$, every $\sigma\in\Sigma_{\mathbf{s},s}$ and every $k\in n+1$,
$$
\Ds^{\gl(n,m)}_{\sig{\Sigma}}(\sigma_{k})=\sigma_{\min\{k,m\}}\left(\left(v^{(\min\{k,m\},\mathbf{s})}_{p}\right)_{p\in\bb{\mathbf{s}}}\right);
$$
\item for every pair of dimensions $j<k$ in $n+1$ with $j<m$, so that $j<\min\{k,m\}$, and every $s\in S$,
\begin{align*}
\Ds^{\gl(n,m)}_{\sig{\Sigma}}(\src^{(j,k)}_{s})&=\src^{(j,\min\{k,m\})}_{s}\left(v^{(\min\{k,m\},s)}_{0}\right),\\
\Ds^{\gl(n,m)}_{\sig{\Sigma}}(\tgt^{(j,k)}_{s})&=\tgt^{(j,\min\{k,m\})}_{s}\left(v^{(\min\{k,m\},s)}_{0}\right),\\
\Ds^{\gl(n,m)}_{\sig{\Sigma}}(\mathrm{i}^{(k,j)}_{s})&=\mathrm{i}^{(\min\{k,m\},j)}_{s}\left(v^{(j,s)}_{0}\right),\\
\Ds^{\gl(n,m)}_{\sig{\Sigma}}(\circ^{(j,k)}_{s})&=v^{(\min\{k,m\},s)}_{0}\circ^{(j,\min\{k,m\})}_{s}v^{(\min\{k,m\},s)}_{1};
\end{align*}
\item for every pair of dimensions $j<k$ in $n+1$ with $j\geq m$, so that $\min\{j,m\}=\min\{k,m\}=m$, and every $s\in S$,
$$
\Ds^{\gl(n,m)}_{\sig{\Sigma}}(\src^{(j,k)}_{s})=
\Ds^{\gl(n,m)}_{\sig{\Sigma}}(\tgt^{(j,k)}_{s})=
\Ds^{\gl(n,m)}_{\sig{\Sigma}}(\mathrm{i}^{(k,j)}_{s})=
\Ds^{\gl(n,m)}_{\sig{\Sigma}}(\circ^{(j,k)}_{s})=v^{(m,s)}_{0};
$$
and
\item for every dimension $k\in n+1$ and every $s\in S$, the domain is $\Dom^{\Ds^{\gl(n,m)}_{\sig{\Sigma}}}_{(k,s)}=\varnothing$.
\end{enumerate}
For $m=n$, we let $\bDs^{(n,n)}_{\sig{\Sigma}}$ and $\bDs^{\gl(n,n)}_{\sig{\Sigma}}$ be the identities $\id_{\sig{\Sigma}^{(n)}}$ and $\id_{\sig{\Sigma}^{\gl(n)}}$, see Definition~\ref{DHepComp}.
\end{definition}

Let us note that, on the dimensions $k\leq m$, $\bDs^{\gl(n,m)}_{\sig{\Sigma}}$ is the identity, and that every operation symbol that involves a dimension $j\geq m$ is sent to a variable of sort $(m,s)$.

\begin{definition}\label{DDiscretizationGl}
Let $\sig{\Sigma}=(S,\Sigma)$ be a many-sorted signature and $m,n\in\mathbb{N}$ with $m\leq n$. The \emph{globular discretization functor} $\Ds^{\gl(n,m)}_{\sig{\Sigma}}$ is the functor $(\bDs^{\gl(n,m)}_{\sig{\Sigma}})^{\ast}_{\gd}$ induced by the generalized hep-derivor $\bDs^{\gl(n,m)}_{\sig{\Sigma}}$ from $\mathsf{PAlg}(\Sigma^{\gl(m)})$ to $\mathsf{PAlg}(\Sigma^{\gl(n)})$, see Proposition~\ref{PGenHepDerivorFunctor}.
\end{definition}

In the following lemma we describe the transformed algebra along the globular discretization derivor.

\begin{lemma}\label{LDiscGlDescription}
Let $\sig{\Sigma}=(S,\Sigma)$ be a many-sorted signature, $m,n\in\mathbb{N}$ with $m<n$ and $\mathbf{B}$ an object in $\mathsf{EPAlg}(\sig{\Sigma}^{\gl(m)},\mathcal{C}^{\sig{\Sigma}^{\gl(m)}})$. Then $\Ds^{\gl(n,m)}_{\sig{\Sigma}}(\mathbf{B})$ is a $\sig{\Sigma}^{\gl(n)}$-algebra whose cells of sort $(k,s)$ are the $\min\{k,m\}$-cells of sort $s$ of $\mathbf{B}$. Moreover, for every pair of dimensions $j<k$ and every $s\in S$, the following statements hold:
\begin{enumerate}
\item $\sigma_{k}$ acts in $\Ds^{\gl(n,m)}_{\sig{\Sigma}}(\mathbf{B})$ as $\sigma^{\mathbf{B}}_{\min\{k,m\}}$;
\item if $j<m$, then the $(j,k)$-sources, $(j,k)$-targets, $(k,j)$-identities and $j$-compositions of $k$-cells of $\Ds^{\gl(n,m)}_{\sig{\Sigma}}(\mathbf{B})$ are the $(j,\min\{k,m\})$-sources, $(j,\min\{k,m\})$-targets, $(\min\{k,m\},j)$-identities and $j$-compositions of $\min\{k,m\}$-cells of $\mathbf{B}$; and
\item if $j\geq m$, then the $(j,k)$-sources, $(j,k)$-targets and $(k,j)$-identities of $\Ds^{\gl(n,m)}_{\sig{\Sigma}}(\mathbf{B})$ are identities, and the $j$-composition of two $k$-cells is defined if, and only if, they are equal, and then it is the common cell.
\end{enumerate}
\end{lemma}
\begin{proof}
Since the domains of $\bDs^{\gl(n,m)}_{\sig{\Sigma}}$ are empty, the cells of sort $(k,s)$ of $\Ds^{\gl(n,m)}_{\sig{\Sigma}}(\mathbf{B})$ are the $\min\{k,m\}$-cells of sort $s$ of $\mathbf{B}$. As in the proof of Lemma~\ref{LDiscDescription}, the operation symbols with empty domain conditions act through total derived operations. Let $j<k$ be dimensions, $s\in S$ and $x$, $y$ cells of sort $(k,s)$ of $\Ds^{\gl(n,m)}_{\sig{\Sigma}}(\mathbf{B})$.

\textsf{(Case $j<m$)} The cells $x$, $y$ satisfy the domain condition of $\circ^{(j,k)}_{s}$ in $\Ds^{\gl(n,m)}_{\sig{\Sigma}}(\mathbf{B})$ if, and only if, $\src^{(j,\min\{k,m\})}_{s}(x)=\tgt^{(j,\min\{k,m\})}_{s}(y)$ in $\mathbf{B}$, in which case, by~(EA$^{j,\min\{k,m\}}$3), $x\circ^{(j,\min\{k,m\})}_{s}y$ is defined in $\mathbf{B}$.

\textsf{(Case $j\geq m$)} The cells $x$, $y$ satisfy the domain condition of $\circ^{(j,k)}_{s}$ in $\Ds^{\gl(n,m)}_{\sig{\Sigma}}(\mathbf{B})$ if, and only if, $x=y$, in which case the term $\Ds^{\gl(n,m)}_{\sig{\Sigma}}(\circ^{(j,k)}_{s})=v^{(m,s)}_{0}$ is defined, with value $x$.

All in all, by Proposition~\ref{PTransformedEpAlg}, $\Ds^{\gl(n,m)}_{\sig{\Sigma}}(\mathbf{B})$ is a $\sig{\Sigma}^{\gl(n)}$-algebra and, by Proposition~\ref{PTransformedAlg}, its operations are as stated.

This completes the proof.
\end{proof}

We next prove that the globular discretization of a globular $m$-categorial $\sig{\Sigma}$-algebra is a globular $n$-categorial $\sig{\Sigma}$-algebra.

\begin{proposition}\label{PDiscGlDerivor}
Let $\sig{\Sigma}=(S,\Sigma)$ be a many-sorted signature and $m,n\in\mathbb{N}$ with $m\leq n$. Then the restriction of $\Ds^{\gl(n,m)}_{\sig{\Sigma}}$ to $\mathsf{EPAlg}(\sig{\Sigma}^{\gl(m)},\mathcal{C}^{\sig{\Sigma}^{\gl(m)}})$ corestricts to $\mathsf{EPAlg}(\sig{\Sigma}^{\gl(n)},\mathcal{C}^{\sig{\Sigma}^{\gl(n)}})$.
\end{proposition}
\begin{proof}
For $m=n$ this follows from Proposition~\ref{PHepId}. Let $m<n$ and $\mathbf{B}$ an object in $\mathsf{EPAlg}(\sig{\Sigma}^{\gl(m)},\mathcal{C}^{\sig{\Sigma}^{\gl(m)}})$. By Definitions~\ref{DCorestriction} and~\ref{DEpVariety}, we have to prove that $\Ds^{\gl(n,m)}_{\sig{\Sigma}}(\mathbf{B})$ is a $\sig{\Sigma}^{\gl(n)}$-algebra in which every axiom of $\mathcal{C}^{\sig{\Sigma}^{\gl(n)}}$ holds.

\textsf{(Ep-algebra)} By Lemma~\ref{LDiscGlDescription}, which rests on Proposition~\ref{PTransformedEpAlg}, $\Ds^{\gl(n,m)}_{\sig{\Sigma}}(\mathbf{B})$ is a $\sig{\Sigma}^{\gl(n)}$-algebra.

\textsf{(Axioms)} By Lemma~\ref{LDiscGlDescription}, the axioms of $\boldsymbol{\mathcal{C}}^{\sig{\Sigma}^{\gl(n)}}$ in $\Ds^{\gl(n,m)}_{\sig{\Sigma}}(\mathbf{B})$ become statements on $\mathbf{B}$, which we check according to the dimensions involved.

\textsf{(Dimensions at most $m$)} The axioms in which all the dimensions are at most $m$ become axioms of $\boldsymbol{\mathcal{C}}^{\sig{\Sigma}^{\gl(m)}}$.

\textsf{(Pairs $j<k$)} For $j<m<k$, the axioms (EA$^{j,k}$1)--(EA$^{j,k}$8) become (EA$^{j,m}$1)--(EA$^{j,m}$8). For $m\leq j<k$, they become trivial, as in the proof of Proposition~\ref{PDiscDerivor}.

\textsf{(Triples $h<j<k$)} For $j<m<k$, the axioms (EA$^{h,j,k}$1)--(EA$^{h,j,k}$5) become (EA$^{h,j,m}$1)--(EA$^{h,j,m}$5). For $h<m\leq j<k$, the $(j,k)$-sources, targets and identities become identities and the $j$-compositions of $k$-cells become trivial, so that the axioms become trivial: for instance, (EA$^{h,j,k}$3) becomes $x\circ^{(h,m)}_{s}y=x\circ^{(h,m)}_{s}y$, and (EA$^{h,j,k}$4), whose premises become $x=y$, $z=t$ and $x\circ^{(h,m)}_{s}z$ defined, becomes $x\circ^{(h,m)}_{s}z=x\circ^{(h,m)}_{s}z$. For $m\leq h$, all the axioms become trivial.

All in all, every axiom of $\boldsymbol{\mathcal{C}}^{\sig{\Sigma}^{\gl(n)}}$ holds in $\Ds^{\gl(n,m)}_{\sig{\Sigma}}(\mathbf{B})$, so that $\Ds^{\gl(n,m)}_{\sig{\Sigma}}(\mathbf{B})$ is an object in $\mathsf{EPAlg}(\sig{\Sigma}^{\gl(n)},\mathcal{C}^{\sig{\Sigma}^{\gl(n)}})$.

This completes the proof.
\end{proof}

Let $\sig{\Sigma}=(S,\Sigma)$ be a many-sorted signature and $m,n\in\mathbb{N}$ with $m\leq n$. By Remark~\ref{RGlobularIdent} and Proposition~\ref{PDiscGlDerivor}, the restriction of $\Ds^{\gl(n,m)}_{\sig{\Sigma}}$ to $\mathsf{mCat}^{\gl}\mathsf{Alg}(\sig{\Sigma})$ corestricts to $\mathsf{nCat}^{\gl}\mathsf{Alg}(\sig{\Sigma})$, and, by abuse of notation, we denote the corestriction by
$$
\Ds^{\gl(n,m)}_{\sig{\Sigma}}\colon\mathsf{mCat}^{\gl}\mathsf{Alg}(\sig{\Sigma})\mor\mathsf{nCat}^{\gl}\mathsf{Alg}(\sig{\Sigma}).
$$
By Lemma~\ref{LDiscGlDescription}, it sends a globular $m$-categorial $\sig{\Sigma}$-algebra $\calg{A}$ to the globular $n$-categorial $\sig{\Sigma}$-algebra whose $k$-cells are the $k$-cells of $\calg{A}$, if $k\leq m$, and the $m$-cells of $\calg{A}$, if $k>m$, in which the structure in the dimensions $k\leq m$ is that of $\calg{A}$, the $(j,k)$-sources, targets and identities with $j\geq m$ are identities, and the $j$-composition of two $k$-cells, for $j\geq m$, is defined if, and only if, they are equal, and then it is the common cell. By Proposition~\ref{PHepId}, $\Ds^{\gl(n,n)}_{\sig{\Sigma}}$ is the identity.

\subsection{Truncation of higher-order categorial algebras}\label{STruncationAlg}

We next define the truncation derivors, in the single-set and in the globular presentations.

\begin{definition}\label{DTruncDerivor}
Let $\sig{\Sigma}=(S,\Sigma)$ be a many-sorted signature and $m,n\in\mathbb{N}$ with $m<n$.
\begin{enumerate}
\item The \emph{single-set truncation derivor} is the generalized hep-derivor
$$
\bTr^{(m,n)}_{\sig{\Sigma}}=\left(\varphi^{\Tr^{(m,n)}_{\sig{\Sigma}}},\Tr^{(m,n)}_{\sig{\Sigma}},\Dom^{\Tr^{(m,n)}_{\sig{\Sigma}}}\right)
$$
from $\sig{\Sigma}^{(m)}$ to $\sig{\Sigma}^{(n)}$ defined as follows:
\begin{enumerate}
\item the mapping on sorts is $\varphi^{\Tr^{(m,n)}_{\sig{\Sigma}}}=\id_{S}$;
\item the assignment of terms is the inclusion, i.e., for every operation symbol $\tau$ of $\Sigma^{(m)}$ of arity $\mathbf{s}$,
$$
\Tr^{(m,n)}_{\sig{\Sigma}}(\tau)=\tau((v^{\mathbf{s}}_{p})_{p\in\bb{\mathbf{s}}});
$$
and
\item for every $s\in S$, the domain is
$$
\Dom^{\Tr^{(m,n)}_{\sig{\Sigma}}}_{s}=\left\lbrace\left(v^{s}_{0},\src^{m}_{s}(v^{s}_{0})\right)\right\rbrace.
$$
\end{enumerate}
\item The \emph{globular truncation derivor} is the generalized hep-derivor
$$
\bTr^{\gl(m,n)}_{\sig{\Sigma}}=\left(\varphi^{\Tr^{\gl(m,n)}_{\sig{\Sigma}}},\Tr^{\gl(m,n)}_{\sig{\Sigma}},\Dom^{\Tr^{\gl(m,n)}_{\sig{\Sigma}}}\right)
$$
from $\sig{\Sigma}^{\gl(m)}$ to $\sig{\Sigma}^{\gl(n)}$ determined by the canonical embedding of $\Sigma^{\gl(m)}$ into $\Sigma^{\gl(n)}$, see the note after Definition~\ref{DSigmaCatGlobular}, i.e., defined as follows:
\begin{enumerate}
\item for every $(k,s)\in(m+1)\times S$, the mapping on sorts sends $(k,s)$ to $\varphi^{\Tr^{\gl(m,n)}_{\sig{\Sigma}}}(k,s)=(k,s)$;
\item for every operation symbol $\tau$ of $\Sigma^{\gl(m)}$ of arity $\mathbf{w}$,
$$
\Tr^{\gl(m,n)}_{\sig{\Sigma}}(\tau)=\tau((v^{\mathbf{w}}_{p})_{p\in\bb{\mathbf{w}}});
$$
and
\item for every $(k,s)\in(m+1)\times S$, the domain is $\Dom^{\Tr^{\gl(m,n)}_{\sig{\Sigma}}}_{(k,s)}=\varnothing$.
\end{enumerate}
\end{enumerate}
For $m=n$, we let $\bTr^{(n,n)}_{\sig{\Sigma}}$ and $\bTr^{\gl(n,n)}_{\sig{\Sigma}}$ be the identities $\id_{\sig{\Sigma}^{(n)}}$ and $\id_{\sig{\Sigma}^{\gl(n)}}$, see Definition~\ref{DHepComp}.
\end{definition}

Let us note that, by Definition~\ref{DDomainSet}, for every partial $\Sigma^{(n)}$-algebra $\mathbf{B}$, the cells of sort $s$ of $\mathbf{B}$ relative to $\bTr^{(m,n)}_{\sig{\Sigma}}$ are its $m$-cells, i.e., the elements $x$ of $B_{s}$ such that $x=\src^{m}_{s}(x)$, and that, for every partial $\Sigma^{\gl(n)}$-algebra $\mathbf{B}$, the cells of sort $(k,s)$ of $\mathbf{B}$ relative to $\bTr^{\gl(m,n)}_{\sig{\Sigma}}$ are all the elements of $B_{(k,s)}$.

\begin{definition}\label{DTruncation}
Let $\sig{\Sigma}=(S,\Sigma)$ be a many-sorted signature and $m,n\in\mathbb{N}$ with $m\leq n$. The \emph{truncation functor} $\Tr^{(m,n)}_{\sig{\Sigma}}$ is the functor $(\bTr^{(m,n)}_{\sig{\Sigma}})^{\ast}_{\gd}$ induced by the generalized hep-derivor $\bTr^{(m,n)}_{\sig{\Sigma}}$ from $\mathsf{PAlg}(\Sigma^{(n)})$ to $\mathsf{PAlg}(\Sigma^{(m)})$, see Proposition~\ref{PGenHepDerivorFunctor}.
\end{definition}

\begin{definition}\label{DTruncationGl}
Let $\sig{\Sigma}=(S,\Sigma)$ be a many-sorted signature and $m,n\in\mathbb{N}$ with $m\leq n$. The \emph{globular truncation functor} $\Tr^{\gl(m,n)}_{\sig{\Sigma}}$ is the functor $(\bTr^{\gl(m,n)}_{\sig{\Sigma}})^{\ast}_{\gd}$ induced by the generalized hep-derivor $\bTr^{\gl(m,n)}_{\sig{\Sigma}}$ from $\mathsf{PAlg}(\Sigma^{\gl(n)})$ to $\mathsf{PAlg}(\Sigma^{\gl(m)})$, see Proposition~\ref{PGenHepDerivorFunctor}.
\end{definition}

In the following lemma we prove that the $m$-cells of a single-set $n$-categorial $\sig{\Sigma}$-algebra are closed under the operations of $\Sigma^{(m)}$.

\begin{lemma}\label{LTruncCells}
Let $\sig{\Sigma}=(S,\Sigma)$ be a many-sorted signature and $m,n\in\mathbb{N}$ with $m<n$, $\mathbf{B}$ an object in $\mathsf{EPAlg}(\sig{\Sigma}^{(n)},\mathcal{C}^{\sig{\Sigma}^{(n)}})$, $s\in S$ and $k<m$. Then the following statements hold:
\begin{enumerate}
\item if $(\mathbf{s},s)\in S^{\star}\times S$, $\sigma\in\Sigma_{\mathbf{s},s}$ and $(x_{p})_{p\in\bb{\mathbf{s}}}$ is a family of $m$-cells, then $\sigma((x_{p})_{p\in\bb{\mathbf{s}}})$ is an $m$-cell;
\item for every cell $x$ of sort $s$, $\src^{k}_{s}(x)$ and $\tgt^{k}_{s}(x)$ are $m$-cells;
\item if $x$, $y$ are $m$-cells of sort $s$ such that $\src^{k}_{s}(x)=\tgt^{k}_{s}(y)$, then $x\circ^{k}_{s}y$ is defined and it is an $m$-cell.
\end{enumerate}
\end{lemma}
\begin{proof}
(1) It follows from Lemma~\ref{LGradCells}~(1), with $m$ in place of $k$.

(2) By Lemma~\ref{LGradCells}~(2), $\src^{k}_{s}(x)$ and $\tgt^{k}_{s}(x)$ are $k$-cells and, by Lemma~\ref{LGradCells}~(3), they are $m$-cells.

(3) It follows from Lemma~\ref{LGradCells}~(4), with $k$ and $m$ in place of $j$ and $k$.
\end{proof}

We next prove that the truncation of a single-set $n$-categorial $\sig{\Sigma}$-algebra is a single-set $m$-categorial $\sig{\Sigma}$-algebra.

\begin{proposition}\label{PTruncDerivor}
Let $\sig{\Sigma}=(S,\Sigma)$ be a many-sorted signature and $m,n\in\mathbb{N}$ with $m\leq n$. Then the restriction of $\Tr^{(m,n)}_{\sig{\Sigma}}$ to $\mathsf{EPAlg}(\sig{\Sigma}^{(n)},\mathcal{C}^{\sig{\Sigma}^{(n)}})$ corestricts to $\mathsf{EPAlg}(\sig{\Sigma}^{(m)},\mathcal{C}^{\sig{\Sigma}^{(m)}})$.
\end{proposition}
\begin{proof}
For $m=n$ this follows from Proposition~\ref{PHepId}. Let $m<n$ and $\mathbf{B}$ an object in $\mathsf{EPAlg}(\sig{\Sigma}^{(n)},\mathcal{C}^{\sig{\Sigma}^{(n)}})$. By Definition~\ref{DTruncDerivor}, the note after it and Proposition~\ref{PTransformedAlg}, the cells of $\Tr^{(m,n)}_{\sig{\Sigma}}(\mathbf{B})$ are the $m$-cells of $\mathbf{B}$, and its operations are restrictions of those of $\mathbf{B}$. By Definitions~\ref{DCorestriction} and~\ref{DEpVariety}, we have to prove that $\Tr^{(m,n)}_{\sig{\Sigma}}(\mathbf{B})$ is a $\sig{\Sigma}^{(m)}$-algebra in which every axiom of $\mathcal{C}^{\sig{\Sigma}^{(m)}}$ holds.

\textsf{(Ep-algebra)} By Lemma~\ref{LTruncCells}, the $m$-cells are closed under the operations of $\Sigma^{(m)}$, which are defined on the families of $m$-cells that satisfy their domain conditions, by~(A0), (A$^{k}$1) and~(A$^{k}$3). Hence, by Proposition~\ref{PTransformedEpAlg}, $\Tr^{(m,n)}_{\sig{\Sigma}}(\mathbf{B})$ is a $\sig{\Sigma}^{(m)}$-algebra.

\textsf{(Axioms)} By the above, $\Tr^{(m,n)}_{\sig{\Sigma}}(\mathbf{B})$ is the relative subalgebra of $\mathbf{B}$ on the $m$-cells, in which the values of terms are their values in $\mathbf{B}$. Therefore, every axiom of $\boldsymbol{\mathcal{C}}^{\sig{\Sigma}^{(m)}}$, which is an axiom of $\boldsymbol{\mathcal{C}}^{\sig{\Sigma}^{(n)}}$, holds in $\Tr^{(m,n)}_{\sig{\Sigma}}(\mathbf{B})$.

This completes the proof.
\end{proof}

We next prove the analogous result for the globular presentation.

\begin{proposition}\label{PTruncGlDerivor}
Let $\sig{\Sigma}=(S,\Sigma)$ be a many-sorted signature and $m,n\in\mathbb{N}$ with $m\leq n$. Then the restriction of $\Tr^{\gl(m,n)}_{\sig{\Sigma}}$ to $\mathsf{EPAlg}(\sig{\Sigma}^{\gl(n)},\mathcal{C}^{\sig{\Sigma}^{\gl(n)}})$ corestricts to $\mathsf{EPAlg}(\sig{\Sigma}^{\gl(m)},\mathcal{C}^{\sig{\Sigma}^{\gl(m)}})$.
\end{proposition}
\begin{proof}
For $m=n$ this follows from Proposition~\ref{PHepId}. Let $m<n$ and $\mathbf{B}$ an object in $\mathsf{EPAlg}(\sig{\Sigma}^{\gl(n)},\mathcal{C}^{\sig{\Sigma}^{\gl(n)}})$. By Definitions~\ref{DCorestriction} and~\ref{DEpVariety}, we have to prove that $\Tr^{\gl(m,n)}_{\sig{\Sigma}}(\mathbf{B})$ is a $\sig{\Sigma}^{\gl(m)}$-algebra in which every axiom of $\mathcal{C}^{\sig{\Sigma}^{\gl(m)}}$ holds.

\textsf{(Ep-algebra)} By the note after Definition~\ref{DTruncDerivor}, all the cells of dimension at most $m$ of $\mathbf{B}$ are cells of the transformed algebra, and every operation symbol of $\Sigma^{\gl(m)}$ is sent to itself. Hence, by Proposition~\ref{PTransformedEpAlg}, since $\mathbf{B}$ is a $\sig{\Sigma}^{\gl(n)}$-algebra and the domain conditions of the operation symbols of $\Sigma^{\gl(m)}$ are the same in $\sig{\Sigma}^{\gl(m)}$ and in $\sig{\Sigma}^{\gl(n)}$, $\Tr^{\gl(m,n)}_{\sig{\Sigma}}(\mathbf{B})$ is a $\sig{\Sigma}^{\gl(m)}$-algebra.

\textsf{(Axioms)} The axioms of $\boldsymbol{\mathcal{C}}^{\sig{\Sigma}^{\gl(m)}}$ are axioms of $\boldsymbol{\mathcal{C}}^{\sig{\Sigma}^{\gl(n)}}$, which hold in $\mathbf{B}$, and the values of terms in the transformed algebra are their values in $\mathbf{B}$.

This completes the proof.
\end{proof}

Let $\sig{\Sigma}=(S,\Sigma)$ be a many-sorted signature and $m,n\in\mathbb{N}$ with $m\leq n$. By Remark~\ref{RSingleSetIdent} and Proposition~\ref{PTruncDerivor}, the restriction of $\Tr^{(m,n)}_{\sig{\Sigma}}$ to $\mathsf{nCatAlg}(\sig{\Sigma})$ corestricts to $\mathsf{mCatAlg}(\sig{\Sigma})$, and, by abuse of notation, we denote the corestriction by
$$
\Tr^{(m,n)}_{\sig{\Sigma}}\colon\mathsf{nCatAlg}(\sig{\Sigma})\mor\mathsf{mCatAlg}(\sig{\Sigma}).
$$
It sends an $n$-categorial $\sig{\Sigma}$-algebra $\calg{A}$ to the $m$-categorial $\sig{\Sigma}$-algebra formed by the $m$-cells of $\calg{A}$, i.e., by~(AB$^{m,k}$1), by the cells that are $k$-cells for every $k\in n$ with $k\geq m$, with the restrictions of the operations of $\Sigma^{(m)}$, and a homomorphism to its restriction. By Proposition~\ref{PHepId}, $\Tr^{(n,n)}_{\sig{\Sigma}}$ is the identity.

Similarly, by Remark~\ref{RGlobularIdent} and Proposition~\ref{PTruncGlDerivor}, the restriction of $\Tr^{\gl(m,n)}_{\sig{\Sigma}}$ to $\mathsf{nCat}^{\gl}\mathsf{Alg}(\sig{\Sigma})$ corestricts to $\mathsf{mCat}^{\gl}\mathsf{Alg}(\sig{\Sigma})$, and, by abuse of notation, we denote the corestriction by
$$
\Tr^{\gl(m,n)}_{\sig{\Sigma}}\colon\mathsf{nCat}^{\gl}\mathsf{Alg}(\sig{\Sigma})\mor\mathsf{mCat}^{\gl}\mathsf{Alg}(\sig{\Sigma}).
$$
It sends a globular $n$-categorial $\sig{\Sigma}$-algebra $\calg{A}$ to the globular $m$-categorial $\sig{\Sigma}$-algebra formed by its cells of dimension at most $m$, and a globular $n$-categorial $\sig{\Sigma}$-homomorphism $F$ to $(F_{(k,s)})_{(k,s)\in(m+1)\times S}$. By Proposition~\ref{PHepId}, $\Tr^{\gl(n,n)}_{\sig{\Sigma}}$ is the identity.

\subsection{On the relationship between discretization and truncation}\label{SDiscTrunc}

In this subsection we prove that discretization is left adjoint to truncation and that these functors compose as expected. As in Subsection~\ref{SGradUnif}, we compute the composites of these functors by composing the corresponding generalized hep-derivors, by means of Proposition~\ref{PHepComp} and Lemma~\ref{LSameCells}. We first consider the single-set presentation.

\begin{proposition}\label{PTruncDisc}
Let $\sig{\Sigma}=(S,\Sigma)$ be a many-sorted signature and $m,n\in\mathbb{N}$ with $m\leq n$. Then the following equality holds
$$
\Tr^{(m,n)}_{\sig{\Sigma}}\circ\Ds^{(n,m)}_{\sig{\Sigma}}=\Id_{\mathsf{mCatAlg}(\sig{\Sigma})}.
$$
\end{proposition}
\begin{proof}
For $m=n$, both functors are identities, by Proposition~\ref{PHepId}. Let $m<n$. We prove the following two claims.

\begin{claim}\label{CTruncDiscA}
The underlying derivor of $\bDs^{(n,m)}_{\sig{\Sigma}}\circ\bTr^{(m,n)}_{\sig{\Sigma}}$ is the identity derivor $(\id_{S},\eta^{\Sigma^{(m)}})$.
\end{claim}

By Definitions~\ref{DDiscDerivor} and~\ref{DTruncDerivor}, both mappings on sorts are $\id_{S}$. Since $\Tr^{(m,n)}_{\sig{\Sigma}}$ sends every operation symbol $\tau$ of $\Sigma^{(m)}$ of arity $\mathbf{s}$ to $\tau((v^{\mathbf{s}}_{p})_{p\in\bb{\mathbf{s}}})$ and $\Ds^{(n,m)}_{\sig{\Sigma}}$ sends every operation symbol of $\Sigma^{(m)}\subseteq\Sigma^{(n)}$ to itself, applied to variables, it follows from~(Der1) and~(Der2) that
$$
\left(\Ds^{(n,m)}_{\sig{\Sigma}}\right)^{\sharp}\circ \Tr^{(m,n)}_{\sig{\Sigma}}=\eta^{\Sigma^{(m)}},
$$
see~(Der3).

Claim~\ref{CTruncDiscA} follows.

\begin{claim}\label{CTruncDiscB}
For every $m$-categorial $\sig{\Sigma}$-algebra $\calg{A}$, $\Tr^{(m,n)}_{\sig{\Sigma}}(\Ds^{(n,m)}_{\sig{\Sigma}}(\calg{A}))=(\bDs^{(n,m)}_{\sig{\Sigma}}\circ\bTr^{(m,n)}_{\sig{\Sigma}})^{\ast}_{\gd}(\calg{A})$, and the cells of this partial $\Sigma^{(m)}$-algebra are all the cells of $\calg{A}$.
\end{claim}

Let $\calg{A}$ be an $m$-categorial $\sig{\Sigma}$-algebra. By Lemma~\ref{LDiscDescription}, $\src^{m}_{s}$ acts as the identity in $\Ds^{(n,m)}_{\sig{\Sigma}}(\calg{A})$, so that, by the note after Definition~\ref{DTruncDerivor}, every cell of $\calg{A}$ is a cell of $\Tr^{(m,n)}_{\sig{\Sigma}}(\Ds^{(n,m)}_{\sig{\Sigma}}(\calg{A}))$. On the other hand, by Definitions~\ref{DHepComp}, \ref{DDiscDerivor} and~\ref{DTruncDerivor}, for every $s\in S$,
$$
\Dom^{\bDs^{(n,m)}_{\sig{\Sigma}}\circ\bTr^{(m,n)}_{\sig{\Sigma}}}_{s}=\varnothing\cup\left(\Ds^{(n,m)}_{\sig{\Sigma}}\right)^{\sharp}\left[\left\lbrace\left(v^{s}_{0},\src^{m}_{s}(v^{s}_{0})\right)\right\rbrace\right]=\left\lbrace\left(v^{s}_{0},v^{s}_{0}\right)\right\rbrace,
$$
which is satisfied by every cell. Hence both partial algebras have the same cells. Moreover, by Propositions~\ref{PDiscDerivor} and~\ref{PTruncDerivor}, $\Tr^{(m,n)}_{\sig{\Sigma}}(\Ds^{(n,m)}_{\sig{\Sigma}}(\calg{A}))$ is a $\sig{\Sigma}^{(m)}$-algebra. Therefore, by Proposition~\ref{PHepComp}~(2), both partial algebras are equal.

Claim~\ref{CTruncDiscB} follows.

Let $\calg{A}$ be an $m$-categorial $\sig{\Sigma}$-algebra. By Claim~\ref{CTruncDiscA}, $\bDs^{(n,m)}_{\sig{\Sigma}}\circ\bTr^{(m,n)}_{\sig{\Sigma}}$ and $\id_{\sig{\Sigma}^{(m)}}$ have the same underlying derivor and, by Claim~\ref{CTruncDiscB}, they determine the same cells of $\calg{A}$, since the domains of $\id_{\sig{\Sigma}^{(m)}}$ are empty. Therefore the following chain of equalities holds
\allowdisplaybreaks
\begin{align*}
\Tr^{(m,n)}_{\sig{\Sigma}}\left(\Ds^{(n,m)}_{\sig{\Sigma}}(\calg{A})\right)
&=
\left(\bDs^{(n,m)}_{\sig{\Sigma}}\circ\bTr^{(m,n)}_{\sig{\Sigma}}\right)^{\ast}_{\gd}(\calg{A})
\tag{1}
\\&=
\left(\id_{\sig{\Sigma}^{(m)}}\right)^{\ast}_{\gd}(\calg{A})
\tag{2}
\\&=
\calg{A}.
\tag{3}
\end{align*}

The first equality follows from Claim~\ref{CTruncDiscB}; the second equality follows from Lemma~\ref{LSameCells}; finally, the last equality follows from Proposition~\ref{PHepId}. By the same results, the equality also holds for the $m$-categorial $\sig{\Sigma}$-homomorphisms.

This completes the proof.
\end{proof}

In the following lemma we introduce the counit of the adjunction between discretization and truncation.

\begin{lemma}\label{LDiscTruncCounit}
Let $\sig{\Sigma}=(S,\Sigma)$ be a many-sorted signature and $m,n\in\mathbb{N}$ with $m\leq n$. For every $n$-categorial $\sig{\Sigma}$-algebra $\calg{B}$, the inclusion $\varepsilon_{\calg{B}}$ of the $m$-cells of $\calg{B}$ into $\calg{B}$ is an $n$-categorial $\sig{\Sigma}$-homomorphism from $\Ds^{(n,m)}_{\sig{\Sigma}}(\Tr^{(m,n)}_{\sig{\Sigma}}(\calg{B}))$ to $\calg{B}$. Moreover, the family $\varepsilon=(\varepsilon_{\calg{B}})_{\calg{B}\in\mathrm{Ob}(\mathsf{nCatAlg}(\sig{\Sigma}))}$ is a natural transformation from $\Ds^{(n,m)}_{\sig{\Sigma}}\circ\Tr^{(m,n)}_{\sig{\Sigma}}$ to $\Id_{\mathsf{nCatAlg}(\sig{\Sigma})}$.
\end{lemma}
\begin{proof}
For $m=n$ there is nothing to prove. Let $m<n$ and $\calg{B}$ an $n$-categorial $\sig{\Sigma}$-algebra. By Remark~\ref{RSingleSetIdent}, it suffices to prove that $\varepsilon_{\calg{B}}$ is a homomorphism of partial $\Sigma^{(n)}$-algebras.

\textsf{(Operations of $\Sigma$ and of dimension $k<m$)} The inclusion commutes with them, since, by Lemma~\ref{LDiscDescription} and the description of $\Tr^{(m,n)}_{\sig{\Sigma}}$, the operations of $\Ds^{(n,m)}_{\sig{\Sigma}}(\Tr^{(m,n)}_{\sig{\Sigma}}(\calg{B}))$ of these kinds are restrictions of those of $\calg{B}$.

\textsf{(Sources and targets of dimension $k\geq m$)} Let $x$ be an $m$-cell of sort $s$. Then, by Lemma~\ref{LDiscDescription}, the $k$-source of $x$ in $\Ds^{(n,m)}_{\sig{\Sigma}}(\Tr^{(m,n)}_{\sig{\Sigma}}(\calg{B}))$ is $x$ and, by Lemma~\ref{LGradCells}~(3), $\src^{k}_{s}(x)=x$ in $\calg{B}$. By a similar argument for the targets, $\varepsilon_{\calg{B}}$ commutes with the $k$-sources and the $k$-targets.

\textsf{(Compositions of dimension $k\geq m$)} Let $x$ be an $m$-cell of sort $s$. By Lemma~\ref{LDiscDescription}, $x\circ^{k}_{s}x$ is defined in $\Ds^{(n,m)}_{\sig{\Sigma}}(\Tr^{(m,n)}_{\sig{\Sigma}}(\calg{B}))$ and equal to $x$. In $\calg{B}$, the following chain of equalities holds
\allowdisplaybreaks
\begin{align*}
x\circ^{k}_{s}x
&=
x\circ^{k}_{s}\src^{k}_{s}(x)
\tag{1}
\\&=
x.
\tag{2}
\end{align*}

The first equality follows from the fact that $x$ is a $k$-cell; finally, the last equality follows from~(A$^{k}$5). Hence $\varepsilon_{\calg{B}}$ preserves the defined $k$-compositions.

\textsf{(Naturality)} Let $H\colon\calg{B}\mor\calg{B}'$ be an $n$-categorial $\sig{\Sigma}$-homomorphism. Since $\Ds^{(n,m)}_{\sig{\Sigma}}(\Tr^{(m,n)}_{\sig{\Sigma}}(H))$ is the restriction of $H$ to the $m$-cells, we have that $H\circ\varepsilon_{\calg{B}}=\varepsilon_{\calg{B}'}\circ\Ds^{(n,m)}_{\sig{\Sigma}}(\Tr^{(m,n)}_{\sig{\Sigma}}(H))$.

This completes the proof.
\end{proof}

We are now in position to prove that discretization is left adjoint to truncation.

\begin{proposition}\label{PDiscTruncAdj}
Let $\sig{\Sigma}=(S,\Sigma)$ be a many-sorted signature and $m,n\in\mathbb{N}$ with $m\leq n$. Then $\Ds^{(n,m)}_{\sig{\Sigma}}$ is left adjoint to $\Tr^{(m,n)}_{\sig{\Sigma}}$. More precisely, see Definition~\ref{DAdj},
$$
(\Tr^{(m,n)}_{\sig{\Sigma}},\Ds^{(n,m)}_{\sig{\Sigma}},\id,\varepsilon)
$$
is an adjunction from $\mathsf{mCatAlg}(\sig{\Sigma})$ to $\mathsf{nCatAlg}(\sig{\Sigma})$ whose unit is the identity and whose counit is the natural transformation $\varepsilon$ of Lemma~\ref{LDiscTruncCounit}.
\end{proposition}
\begin{proof}
By Proposition~\ref{PTruncDisc}, $\Tr^{(m,n)}_{\sig{\Sigma}}\circ\Ds^{(n,m)}_{\sig{\Sigma}}=\Id_{\mathsf{mCatAlg}(\sig{\Sigma})}$, so that the identity $\id=(\id_{\calg{A}})_{\calg{A}\in\mathrm{Ob}(\mathsf{mCatAlg}(\sig{\Sigma}))}$ is a natural transformation from $\Id_{\mathsf{mCatAlg}(\sig{\Sigma})}$ to $\Tr^{(m,n)}_{\sig{\Sigma}}\circ\Ds^{(n,m)}_{\sig{\Sigma}}$; and, by Lemma~\ref{LDiscTruncCounit}, $\varepsilon$ is a natural transformation from $\Ds^{(n,m)}_{\sig{\Sigma}}\circ\Tr^{(m,n)}_{\sig{\Sigma}}$ to $\Id_{\mathsf{nCatAlg}(\sig{\Sigma})}$. By Definition~\ref{DAdj}, we have to prove the triangular identities
$$
\left(\Tr^{(m,n)}_{\sig{\Sigma}}\ast\varepsilon\right)\circ\left(\id\ast\Tr^{(m,n)}_{\sig{\Sigma}}\right)=\id_{\Tr^{(m,n)}_{\sig{\Sigma}}}
\qquad\text{and}\qquad
\left(\varepsilon\ast\Ds^{(n,m)}_{\sig{\Sigma}}\right)\circ\left(\Ds^{(n,m)}_{\sig{\Sigma}}\ast\id\right)=\id_{\Ds^{(n,m)}_{\sig{\Sigma}}}.
$$
If $m=n$, then all the functors involved are identities and $\varepsilon$ is the identity, so that both identities hold. Let us assume that $m<n$. We recall that a homomorphism is determined by its domain, its codomain and its underlying many-sorted mapping.

\textsf{(First triangular identity)} Let $\calg{B}$ be an $n$-categorial $\sig{\Sigma}$-algebra. By Proposition~\ref{PTruncDisc}, the domain $\Tr^{(m,n)}_{\sig{\Sigma}}(\Ds^{(n,m)}_{\sig{\Sigma}}(\Tr^{(m,n)}_{\sig{\Sigma}}(\calg{B})))$ of $\Tr^{(m,n)}_{\sig{\Sigma}}(\varepsilon_{\calg{B}})$ is $\Tr^{(m,n)}_{\sig{\Sigma}}(\calg{B})$, whose cells of sort $s$ are the $m$-cells of sort $s$ of $\calg{B}$. Let $s\in S$ and $x$ an $m$-cell of sort $s$ of $\calg{B}$. Then the following chain of equalities holds
\allowdisplaybreaks
\begin{align*}
\left(\left(\left(\Tr^{(m,n)}_{\sig{\Sigma}}\ast\varepsilon\right)\circ\left(\id\ast\Tr^{(m,n)}_{\sig{\Sigma}}\right)\right)_{\calg{B}}\right)_{s}(x)
&=
\left(\Tr^{(m,n)}_{\sig{\Sigma}}\left(\varepsilon_{\calg{B}}\right)\circ\id_{\Tr^{(m,n)}_{\sig{\Sigma}}(\calg{B})}\right)_{s}(x)
\tag{1}
\\&=
\Tr^{(m,n)}_{\sig{\Sigma}}\left(\varepsilon_{\calg{B}}\right)_{s}(x)
\tag{2}
\\&=
\left(\varepsilon_{\calg{B}}\right)_{s}(x)
\tag{3}
\\&=
x
\tag{4}
\\&=
\left(\id_{\Tr^{(m,n)}_{\sig{\Sigma}}}\right)_{\calg{B},s}(x).
\tag{5}
\end{align*}

The first equality unravels the definitions of the horizontal and vertical composition of natural transformations; the second equality follows from the fact that $\id_{\Tr^{(m,n)}_{\sig{\Sigma}}(\calg{B})}$ is the identity homomorphism; the third equality follows from Proposition~\ref{PGenHepDerivorFunctor}, since $\varphi^{\Tr^{(m,n)}_{\sig{\Sigma}}}=\id_{S}$, so that $\Tr^{(m,n)}_{\sig{\Sigma}}(\varepsilon_{\calg{B}})_{s}$ is the birestriction of $(\varepsilon_{\calg{B}})_{s}$ to the $m$-cells; the fourth equality follows from Lemma~\ref{LDiscTruncCounit}, since $\varepsilon_{\calg{B}}$ is an inclusion; finally, the last equality unravels the definition of the identity natural transformation at $\Tr^{(m,n)}_{\sig{\Sigma}}$.

Hence the first triangular identity holds.

\textsf{(Second triangular identity)} Let $\calg{A}$ be an $m$-categorial $\sig{\Sigma}$-algebra. By Proposition~\ref{PTruncDisc}, the domain $\Ds^{(n,m)}_{\sig{\Sigma}}(\Tr^{(m,n)}_{\sig{\Sigma}}(\Ds^{(n,m)}_{\sig{\Sigma}}(\calg{A})))$ of $\varepsilon_{\Ds^{(n,m)}_{\sig{\Sigma}}(\calg{A})}$ is $\Ds^{(n,m)}_{\sig{\Sigma}}(\calg{A})$. Let $s\in S$ and $x$ a cell of sort $s$ of $\Ds^{(n,m)}_{\sig{\Sigma}}(\calg{A})$. By Lemma~\ref{LDiscDescription}~(2), $\src^{m}_{s}$ acts as the identity in $\Ds^{(n,m)}_{\sig{\Sigma}}(\calg{A})$, so that $x$ is an $m$-cell of $\Ds^{(n,m)}_{\sig{\Sigma}}(\calg{A})$. Then the following chain of equalities holds
\allowdisplaybreaks
\begin{align*}
\left(\left(\left(\varepsilon\ast\Ds^{(n,m)}_{\sig{\Sigma}}\right)\circ\left(\Ds^{(n,m)}_{\sig{\Sigma}}\ast\id\right)\right)_{\calg{A}}\right)_{s}(x)
&=
\left(\varepsilon_{\Ds^{(n,m)}_{\sig{\Sigma}}(\calg{A})}\circ\Ds^{(n,m)}_{\sig{\Sigma}}\left(\id_{\calg{A}}\right)\right)_{s}(x)
\tag{1}
\\&=
\left(\varepsilon_{\Ds^{(n,m)}_{\sig{\Sigma}}(\calg{A})}\right)_{s}(x)
\tag{2}
\\&=
x
\tag{3}
\\&=
\left(\id_{\Ds^{(n,m)}_{\sig{\Sigma}}}\right)_{\calg{A},s}(x).
\tag{4}
\end{align*}

The first equality unravels the definitions of the horizontal and vertical composition of natural transformations; the second equality follows from the fact that $\Ds^{(n,m)}_{\sig{\Sigma}}$ preserves identities; the third equality follows from Lemma~\ref{LDiscTruncCounit}, since $\varepsilon_{\Ds^{(n,m)}_{\sig{\Sigma}}(\calg{A})}$ is the inclusion of the $m$-cells of $\Ds^{(n,m)}_{\sig{\Sigma}}(\calg{A})$ and $x$ is one of them; finally, the last equality unravels the definition of the identity natural transformation at $\Ds^{(n,m)}_{\sig{\Sigma}}$.

Hence the second triangular identity holds.

This completes the proof.
\end{proof}

We next consider the globular presentation.

\begin{proposition}\label{PTruncDiscGl}
Let $\sig{\Sigma}=(S,\Sigma)$ be a many-sorted signature and $m,n\in\mathbb{N}$ with $m\leq n$. Then the following equality holds in $\mathsf{hepSig}_{\gd}$
$$
\bDs^{\gl(n,m)}_{\sig{\Sigma}}\circ\bTr^{\gl(m,n)}_{\sig{\Sigma}}=\id_{\sig{\Sigma}^{\gl(m)}}.
$$
Consequently, $\Tr^{\gl(m,n)}_{\sig{\Sigma}}\circ\Ds^{\gl(n,m)}_{\sig{\Sigma}}=\Id_{\mathsf{mCat}^{\gl}\mathsf{Alg}(\sig{\Sigma})}$.
\end{proposition}
\begin{proof}
For $m=n$ both generalized hep-derivors are identities. Let $m<n$. We compare the underlying derivors and the domains.

\textsf{(Underlying derivors)} By Definitions~\ref{DDiscGlDerivor} and~\ref{DTruncDerivor}, for every $(k,s)\in(m+1)\times S$, the following chain of equalities holds
\allowdisplaybreaks
\begin{align*}
\varphi^{\Ds^{\gl(n,m)}_{\sig{\Sigma}}}\left(\varphi^{\Tr^{\gl(m,n)}_{\sig{\Sigma}}}(k,s)\right)
&=
\varphi^{\Ds^{\gl(n,m)}_{\sig{\Sigma}}}(k,s)
\tag{1}
\\&=
(\min\{k,m\},s)
\tag{2}
\\&=
(k,s).
\tag{3}
\end{align*}

The first equality unravels Definition~\ref{DTruncDerivor}; the second equality unravels Definition~\ref{DDiscGlDerivor}; finally, the last equality follows from the fact that $k\leq m$. Moreover, $\Tr^{\gl(m,n)}_{\sig{\Sigma}}$ sends every operation symbol $\tau$ of $\Sigma^{\gl(m)}$ to itself, applied to variables, and, by Definition~\ref{DDiscGlDerivor}, since all the dimensions of $\tau$ are at most $m$, $\Ds^{\gl(n,m)}_{\sig{\Sigma}}$ sends $\tau$ to itself, applied to variables. Hence, by~(Der1) and~(Der2), $(\Ds^{\gl(n,m)}_{\sig{\Sigma}})^{\sharp}\circ\Tr^{\gl(m,n)}_{\sig{\Sigma}}=\eta^{\Sigma^{\gl(m)}}$, i.e., the underlying derivor of $\bDs^{\gl(n,m)}_{\sig{\Sigma}}\circ\bTr^{\gl(m,n)}_{\sig{\Sigma}}$ is the identity derivor at $\Sigma^{\gl(m)}$, see~(Der3).

\textsf{(Domains)} By Definitions~\ref{DHepComp}, \ref{DDiscGlDerivor} and~\ref{DTruncDerivor}, for every $(k,s)\in(m+1)\times S$,
$$
\Dom^{\bDs^{\gl(n,m)}_{\sig{\Sigma}}\circ\bTr^{\gl(m,n)}_{\sig{\Sigma}}}_{(k,s)}=\varnothing\cup\left(\Ds^{\gl(n,m)}_{\sig{\Sigma}}\right)^{\sharp}[\varnothing]=\varnothing.
$$
Hence the stated equality holds.

\textsf{(Functors)} Let $\calg{A}$ be a globular $m$-categorial $\sig{\Sigma}$-algebra. Since the domains of $\bTr^{\gl(m,n)}_{\sig{\Sigma}}$ are empty, by the note after Proposition~\ref{PHepComp}, $\Tr^{\gl(m,n)}_{\sig{\Sigma}}(\Ds^{\gl(n,m)}_{\sig{\Sigma}}(\calg{A}))$ and $(\bDs^{\gl(n,m)}_{\sig{\Sigma}}\circ\bTr^{\gl(m,n)}_{\sig{\Sigma}})^{\ast}_{\gd}(\calg{A})$ have the same cells and, by Propositions~\ref{PDiscGlDerivor} and~\ref{PTruncGlDerivor}, the former is a $\sig{\Sigma}^{\gl(m)}$-algebra. Therefore, by Proposition~\ref{PHepComp}~(2), the stated equality and Proposition~\ref{PHepId}, $\Tr^{\gl(m,n)}_{\sig{\Sigma}}(\Ds^{\gl(n,m)}_{\sig{\Sigma}}(\calg{A}))=(\id_{\sig{\Sigma}^{\gl(m)}})^{\ast}_{\gd}(\calg{A})=\calg{A}$, and similarly for the homomorphisms.

This completes the proof.
\end{proof}

In the following lemma we introduce the counit of the adjunction between globular discretization and globular truncation.

\begin{lemma}\label{LDiscTruncCounitGl}
Let $\sig{\Sigma}=(S,\Sigma)$ be a many-sorted signature and $m,n\in\mathbb{N}$ with $m\leq n$. For every globular $n$-categorial $\sig{\Sigma}$-algebra $\calg{B}$, let $\varepsilon^{\gl}_{\calg{B}}$ be the family of mappings whose component at the sort $(k,s)$ is the identity, if $k\leq m$, and $\mathrm{i}^{(k,m)}_{s}$, if $k>m$. Then $\varepsilon^{\gl}_{\calg{B}}$ is a globular $n$-categorial $\sig{\Sigma}$-homomorphism from $\Ds^{\gl(n,m)}_{\sig{\Sigma}}(\Tr^{\gl(m,n)}_{\sig{\Sigma}}(\calg{B}))$ to $\calg{B}$, and the family $\varepsilon^{\gl}=(\varepsilon^{\gl}_{\calg{B}})_{\calg{B}\in\mathrm{Ob}(\mathsf{nCat}^{\gl}\mathsf{Alg}(\sig{\Sigma}))}$ is a natural transformation from $\Ds^{\gl(n,m)}_{\sig{\Sigma}}\circ\Tr^{\gl(m,n)}_{\sig{\Sigma}}$ to $\Id_{\mathsf{nCat}^{\gl}\mathsf{Alg}(\sig{\Sigma})}$.
\end{lemma}
\begin{proof}
For $m=n$ there is nothing to prove. Let $m<n$ and $\calg{B}$ a globular $n$-categorial $\sig{\Sigma}$-algebra. By Remark~\ref{RGlobularIdent}, it suffices to prove that $\varepsilon^{\gl}_{\calg{B}}$ is a homomorphism of partial $\Sigma^{\gl(n)}$-algebras. On the dimensions $k\leq m$ it is the identity, and the structure of $\Ds^{\gl(n,m)}_{\sig{\Sigma}}(\Tr^{\gl(m,n)}_{\sig{\Sigma}}(\calg{B}))$ is that of $\calg{B}$, by Lemma~\ref{LDiscGlDescription}. Let $j<k$ with $k>m$ and $s\in S$.

\textsf{(Operation symbols $\sigma_{k}$)} By the third equation of~(EA$^{m,k}$7), $\mathrm{i}^{(k,m)}_{s}$ commutes with the operations $\sigma$.

\textsf{(Case $j<m$)} Let $x$ be an $m$-cell of sort $s$. Then the following chain of equalities holds
\allowdisplaybreaks
\begin{align*}
\src^{(j,k)}_{s}\left(\mathrm{i}^{(k,m)}_{s}(x)\right)
&=
\src^{(j,m)}_{s}\left(\src^{(m,k)}_{s}\left(\mathrm{i}^{(k,m)}_{s}(x)\right)\right)
\tag{1}
\\&=
\src^{(j,m)}_{s}(x).
\tag{2}
\end{align*}

The first equality follows from~(EA$^{j,m,k}$1); finally, the last equality follows from Lemma~\ref{LUnifInj}. By a similar argument, $\varepsilon^{\gl}_{\calg{B}}$ commutes with the targets. Moreover, for every $j$-cell $u$, by~(EA$^{j,m,k}$2), $\mathrm{i}^{(k,m)}_{s}(\mathrm{i}^{(m,j)}_{s}(u))=\mathrm{i}^{(k,j)}_{s}(u)$, and, for every pair of $m$-cells $x$, $y$ such that $x\circ^{(j,m)}_{s}y$ is defined, by~(EA$^{j,m,k}$5), $\mathrm{i}^{(k,m)}_{s}(x\circ^{(j,m)}_{s}y)=\mathrm{i}^{(k,m)}_{s}(x)\circ^{(j,k)}_{s}\mathrm{i}^{(k,m)}_{s}(y)$.

\textsf{(Case $j\geq m$)} By~(EA$^{m,j,k}$2), $\mathrm{i}^{(k,m)}_{s}=\mathrm{i}^{(k,j)}_{s}\circ\mathrm{i}^{(j,m)}_{s}$. Hence, for every $m$-cell $x$ of sort $s$, by Lemma~\ref{LUnifInj}, $\src^{(j,k)}_{s}(\mathrm{i}^{(k,m)}_{s}(x))=\mathrm{i}^{(j,m)}_{s}(x)$, and similarly for the targets; moreover, $\mathrm{i}^{(k,j)}_{s}(\mathrm{i}^{(j,m)}_{s}(x))=\mathrm{i}^{(k,m)}_{s}(x)$. Finally, by~(EA$^{j,k}$5), the following chain of equalities holds
\allowdisplaybreaks
\begin{align*}
\mathrm{i}^{(k,m)}_{s}(x)\circ^{(j,k)}_{s}\mathrm{i}^{(k,m)}_{s}(x)
&=
\mathrm{i}^{(k,m)}_{s}(x)\circ^{(j,k)}_{s}\mathrm{i}^{(k,j)}_{s}\left(\src^{(j,k)}_{s}\left(\mathrm{i}^{(k,m)}_{s}(x)\right)\right)
\tag{1}
\\&=
\mathrm{i}^{(k,m)}_{s}(x).
\tag{2}
\end{align*}

The first equality follows from the above computation of the $(j,k)$-source; finally, the last equality follows from~(EA$^{j,k}$5).

\textsf{(Naturality)} It follows from the fact that the globular $n$-categorial $\sig{\Sigma}$-homomorphisms commute with the identities.

This completes the proof.
\end{proof}

\begin{proposition}\label{PDiscTruncAdjGl}
Let $\sig{\Sigma}=(S,\Sigma)$ be a many-sorted signature and $m,n\in\mathbb{N}$ with $m\leq n$. Then $\Ds^{\gl(n,m)}_{\sig{\Sigma}}$ is left adjoint to $\Tr^{\gl(m,n)}_{\sig{\Sigma}}$. More precisely,
$$
(\Tr^{\gl(m,n)}_{\sig{\Sigma}},\Ds^{\gl(n,m)}_{\sig{\Sigma}},\id,\varepsilon^{\gl})
$$
is an adjunction from $\mathsf{mCat}^{\gl}\mathsf{Alg}(\sig{\Sigma})$ to $\mathsf{nCat}^{\gl}\mathsf{Alg}(\sig{\Sigma})$ whose unit is the identity and whose counit is the natural transformation $\varepsilon^{\gl}$ of Lemma~\ref{LDiscTruncCounitGl}.
\end{proposition}
\begin{proof}
By Proposition~\ref{PTruncDiscGl}, $\Tr^{\gl(m,n)}_{\sig{\Sigma}}\circ\Ds^{\gl(n,m)}_{\sig{\Sigma}}=\Id_{\mathsf{mCat}^{\gl}\mathsf{Alg}(\sig{\Sigma})}$, so that the identity $\id=(\id_{\calg{A}})_{\calg{A}\in\mathrm{Ob}(\mathsf{mCat}^{\gl}\mathsf{Alg}(\sig{\Sigma}))}$ is a natural transformation from $\Id_{\mathsf{mCat}^{\gl}\mathsf{Alg}(\sig{\Sigma})}$ to $\Tr^{\gl(m,n)}_{\sig{\Sigma}}\circ\Ds^{\gl(n,m)}_{\sig{\Sigma}}$; and, by Lemma~\ref{LDiscTruncCounitGl}, $\varepsilon^{\gl}$ is a natural transformation from $\Ds^{\gl(n,m)}_{\sig{\Sigma}}\circ\Tr^{\gl(m,n)}_{\sig{\Sigma}}$ to $\Id_{\mathsf{nCat}^{\gl}\mathsf{Alg}(\sig{\Sigma})}$. By Definition~\ref{DAdj}, we have to prove the triangular identities
$$
\left(\Tr^{\gl(m,n)}_{\sig{\Sigma}}\ast\varepsilon^{\gl}\right)\circ\left(\id\ast\Tr^{\gl(m,n)}_{\sig{\Sigma}}\right)=\id_{\Tr^{\gl(m,n)}_{\sig{\Sigma}}}
\qquad\text{and}\qquad
\left(\varepsilon^{\gl}\ast\Ds^{\gl(n,m)}_{\sig{\Sigma}}\right)\circ\left(\Ds^{\gl(n,m)}_{\sig{\Sigma}}\ast\id\right)=\id_{\Ds^{\gl(n,m)}_{\sig{\Sigma}}}.
$$
If $m=n$, then all the functors involved are identities and $\varepsilon^{\gl}$ is the identity, so that both identities hold. Let us assume that $m<n$. We recall that a homomorphism is determined by its domain, its codomain and its underlying many-sorted mapping.

\textsf{(First triangular identity)} Let $\calg{B}$ be a globular $n$-categorial $\sig{\Sigma}$-algebra. By Proposition~\ref{PTruncDiscGl}, the domain of $\Tr^{\gl(m,n)}_{\sig{\Sigma}}(\varepsilon^{\gl}_{\calg{B}})$ is $\Tr^{\gl(m,n)}_{\sig{\Sigma}}(\calg{B})$. Let $(k,s)\in(m+1)\times S$ and $x$ a cell of sort $(k,s)$ of $\Tr^{\gl(m,n)}_{\sig{\Sigma}}(\calg{B})$, i.e., a $k$-cell of sort $s$ of $\calg{B}$. Then the following chain of equalities holds
\allowdisplaybreaks
\begin{align*}
\left(\left(\left(\Tr^{\gl(m,n)}_{\sig{\Sigma}}\ast\varepsilon^{\gl}\right)\circ\left(\id\ast\Tr^{\gl(m,n)}_{\sig{\Sigma}}\right)\right)_{\calg{B}}\right)_{(k,s)}(x)
&=
\left(\Tr^{\gl(m,n)}_{\sig{\Sigma}}\left(\varepsilon^{\gl}_{\calg{B}}\right)\circ\id_{\Tr^{\gl(m,n)}_{\sig{\Sigma}}(\calg{B})}\right)_{(k,s)}(x)
\tag{1}
\\&=
\left(\varepsilon^{\gl}_{\calg{B}}\right)_{(k,s)}(x)
\tag{2}
\\&=
x
\tag{3}
\\&=
\left(\id_{\Tr^{\gl(m,n)}_{\sig{\Sigma}}}\right)_{\calg{B},(k,s)}(x).
\tag{4}
\end{align*}

The first equality unravels the definitions of the horizontal and vertical composition of natural transformations; the second equality follows from the fact that $\id_{\Tr^{\gl(m,n)}_{\sig{\Sigma}}(\calg{B})}$ is the identity homomorphism and from Proposition~\ref{PGenHepDerivorFunctor}, since $\varphi^{\Tr^{\gl(m,n)}_{\sig{\Sigma}}}(k,s)=(k,s)$ and the domains of $\bTr^{\gl(m,n)}_{\sig{\Sigma}}$ are empty; the third equality follows from Lemma~\ref{LDiscTruncCounitGl}, since $k\leq m$; finally, the last equality unravels the definition of the identity natural transformation at $\Tr^{\gl(m,n)}_{\sig{\Sigma}}$.

Hence the first triangular identity holds.

\textsf{(Second triangular identity)} Let $\calg{A}$ be a globular $m$-categorial $\sig{\Sigma}$-algebra. By Proposition~\ref{PTruncDiscGl}, the domain of $\varepsilon^{\gl}_{\Ds^{\gl(n,m)}_{\sig{\Sigma}}(\calg{A})}$ is $\Ds^{\gl(n,m)}_{\sig{\Sigma}}(\calg{A})$. Let $(k,s)\in(n+1)\times S$ and $x$ a cell of sort $(k,s)$ of $\Ds^{\gl(n,m)}_{\sig{\Sigma}}(\calg{A})$. If $k\leq m$, then, by Lemma~\ref{LDiscTruncCounitGl}, $(\varepsilon^{\gl}_{\Ds^{\gl(n,m)}_{\sig{\Sigma}}(\calg{A})})_{(k,s)}(x)=x$. If $k>m$, then the following chain of equalities holds
\allowdisplaybreaks
\begin{align*}
\left(\varepsilon^{\gl}_{\Ds^{\gl(n,m)}_{\sig{\Sigma}}(\calg{A})}\right)_{(k,s)}(x)
&=
\mathrm{i}^{(k,m)\,\Ds^{\gl(n,m)}_{\sig{\Sigma}}(\calg{A})}_{s}(x)
\tag{1}
\\&=
x.
\tag{2}
\end{align*}

The first equality unravels the definition of $\varepsilon^{\gl}$, see Lemma~\ref{LDiscTruncCounitGl}; finally, the last equality follows from Lemma~\ref{LDiscGlDescription}~(3), since $m\geq m$, so that the $(k,m)$-identities of $\Ds^{\gl(n,m)}_{\sig{\Sigma}}(\calg{A})$ are identity mappings.

Hence $\varepsilon^{\gl}_{\Ds^{\gl(n,m)}_{\sig{\Sigma}}(\calg{A})}=\id_{\Ds^{\gl(n,m)}_{\sig{\Sigma}}(\calg{A})}$, and the following chain of equalities holds
\allowdisplaybreaks
\begin{align*}
\left(\left(\varepsilon^{\gl}\ast\Ds^{\gl(n,m)}_{\sig{\Sigma}}\right)\circ\left(\Ds^{\gl(n,m)}_{\sig{\Sigma}}\ast\id\right)\right)_{\calg{A}}
&=
\varepsilon^{\gl}_{\Ds^{\gl(n,m)}_{\sig{\Sigma}}(\calg{A})}\circ\Ds^{\gl(n,m)}_{\sig{\Sigma}}\left(\id_{\calg{A}}\right)
\tag{1}
\\&=
\id_{\Ds^{\gl(n,m)}_{\sig{\Sigma}}(\calg{A})}\circ\id_{\Ds^{\gl(n,m)}_{\sig{\Sigma}}(\calg{A})}
\tag{2}
\\&=
\left(\id_{\Ds^{\gl(n,m)}_{\sig{\Sigma}}}\right)_{\calg{A}}.
\tag{3}
\end{align*}

The first equality unravels the definitions of the horizontal and vertical composition of natural transformations; the second equality follows from the above and from the fact that $\Ds^{\gl(n,m)}_{\sig{\Sigma}}$ preserves identities; finally, the last equality unravels the definition of the identity natural transformation at $\Ds^{\gl(n,m)}_{\sig{\Sigma}}$.

Hence the second triangular identity holds.

This completes the proof.
\end{proof}

Finally, the truncation functors compose as expected, and so do the discretization functors; this is the fact used in Sections~\ref{SLimitSingleSet} and~\ref{SLimitGlobular} to form the projective systems of categorial algebras.

\begin{proposition}\label{PTruncComp}
Let $\sig{\Sigma}=(S,\Sigma)$ be a many-sorted signature and $m,n,p\in\mathbb{N}$ with $m\leq n\leq p$. Then the following equality holds in $\mathsf{hepSig}_{\gd}$
$$
\bTr^{\gl(n,p)}_{\sig{\Sigma}}\circ\bTr^{\gl(m,n)}_{\sig{\Sigma}}=\bTr^{\gl(m,p)}_{\sig{\Sigma}},
$$
and the following equalities of functors hold
$$
\Tr^{(m,n)}_{\sig{\Sigma}}\circ \Tr^{(n,p)}_{\sig{\Sigma}}=\Tr^{(m,p)}_{\sig{\Sigma}}
\qquad\text{and}\qquad
\Tr^{\gl(m,n)}_{\sig{\Sigma}}\circ \Tr^{\gl(n,p)}_{\sig{\Sigma}}=\Tr^{\gl(m,p)}_{\sig{\Sigma}}.
$$
\end{proposition}
\begin{proof}
If two of the dimensions are equal, the statement follows from Definition~\ref{DTruncDerivor} and Propositions~\ref{PHepSigCat} and~\ref{PHepId}, since $\bTr^{(n,n)}_{\sig{\Sigma}}$ and $\bTr^{\gl(n,n)}_{\sig{\Sigma}}$ are identities. Let $m<n<p$. We consider first the single-set case and then the globular case.

\textsf{(Single-set case: underlying derivors)} The mappings on sorts of $\bTr^{(m,n)}_{\sig{\Sigma}}$, $\bTr^{(n,p)}_{\sig{\Sigma}}$ and $\bTr^{(m,p)}_{\sig{\Sigma}}$ are $\id_{S}$. Moreover, for every operation symbol $\tau$ of $\Sigma^{(m)}$ of arity $\mathbf{s}$, the following chain of equalities holds
\allowdisplaybreaks
\begin{align*}
\left(\Tr^{(n,p)}_{\sig{\Sigma}}\right)^{\sharp}\left(\Tr^{(m,n)}_{\sig{\Sigma}}(\tau)\right)
&=
\left(\Tr^{(n,p)}_{\sig{\Sigma}}\right)^{\sharp}\left(\tau\left((v^{\mathbf{s}}_{q})_{q\in\bb{\mathbf{s}}}\right)\right)
\tag{1}
\\&=
\tau\left((v^{\mathbf{s}}_{q})_{q\in\bb{\mathbf{s}}}\right)
\tag{2}
\\&=
\Tr^{(m,p)}_{\sig{\Sigma}}(\tau).
\tag{3}
\end{align*}

The first equality unravels Definition~\ref{DTruncDerivor}; the second equality follows from~(Der1) and~(Der2), taking into account that $\Tr^{(n,p)}_{\sig{\Sigma}}(\tau)=\tau((v^{\mathbf{s}}_{q})_{q\in\bb{\mathbf{s}}})$; finally, the last equality unravels Definition~\ref{DTruncDerivor}. Thus the underlying derivor of $\bTr^{(n,p)}_{\sig{\Sigma}}\circ\bTr^{(m,n)}_{\sig{\Sigma}}$ is that of $\bTr^{(m,p)}_{\sig{\Sigma}}$.

\textsf{(Single-set case: cells)} Let $\calg{B}$ be a $p$-categorial $\sig{\Sigma}$-algebra. By Definitions~\ref{DHepComp} and~\ref{DTruncDerivor}, for every $s\in S$,
$$
\Dom^{\bTr^{(n,p)}_{\sig{\Sigma}}\circ\bTr^{(m,n)}_{\sig{\Sigma}}}_{s}=\left\lbrace\left(v^{s}_{0},\src^{n}_{s}(v^{s}_{0})\right),\left(v^{s}_{0},\src^{m}_{s}(v^{s}_{0})\right)\right\rbrace,
$$
since $\Tr^{(n,p)}_{\sig{\Sigma}}$ sends $\src^{m}_{s}$ to itself. Hence the cells of $(\bTr^{(n,p)}_{\sig{\Sigma}}\circ\bTr^{(m,n)}_{\sig{\Sigma}})^{\ast}_{\gd}(\calg{B})$ are the cells of $\calg{B}$ that are $m$-cells and $n$-cells. By the note after Definition~\ref{DTruncDerivor}, the cells of $\Tr^{(m,n)}_{\sig{\Sigma}}(\Tr^{(n,p)}_{\sig{\Sigma}}(\calg{B}))$ are the $m$-cells among the $n$-cells of $\calg{B}$, since $\src^{m}_{s}$ acts in $\Tr^{(n,p)}_{\sig{\Sigma}}(\calg{B})$ as in $\calg{B}$. By Lemma~\ref{LGradCells}~(3), in both cases these are the $m$-cells of $\calg{B}$, which are the cells of $\Tr^{(m,p)}_{\sig{\Sigma}}(\calg{B})$. Moreover, by Proposition~\ref{PTruncDerivor}, $\Tr^{(m,n)}_{\sig{\Sigma}}(\Tr^{(n,p)}_{\sig{\Sigma}}(\calg{B}))$ is a $\sig{\Sigma}^{(m)}$-algebra.

Hence, by Proposition~\ref{PHepComp}~(2), Lemma~\ref{LSameCells} and the comparison of the underlying derivors, the following chain of equalities holds
$$
\Tr^{(m,n)}_{\sig{\Sigma}}\left(\Tr^{(n,p)}_{\sig{\Sigma}}(\calg{B})\right)
=
\left(\bTr^{(n,p)}_{\sig{\Sigma}}\circ\bTr^{(m,n)}_{\sig{\Sigma}}\right)^{\ast}_{\gd}(\calg{B})
=
\Tr^{(m,p)}_{\sig{\Sigma}}(\calg{B}),
$$
and, by the same results, the equality also holds for the $p$-categorial $\sig{\Sigma}$-homomorphisms.

\textsf{(Globular case: underlying derivors)} By Definition~\ref{DTruncDerivor}, the underlying derivors of $\bTr^{\gl(m,n)}_{\sig{\Sigma}}$, $\bTr^{\gl(n,p)}_{\sig{\Sigma}}$ and $\bTr^{\gl(m,p)}_{\sig{\Sigma}}$ are determined by the canonical embeddings of $\Sigma^{\gl(m)}$ into $\Sigma^{\gl(n)}$, of $\Sigma^{\gl(n)}$ into $\Sigma^{\gl(p)}$, and of $\Sigma^{\gl(m)}$ into $\Sigma^{\gl(p)}$, respectively. Their mappings on sorts send $(k,s)$ to $(k,s)$. Moreover, for every operation symbol $\tau$ of $\Sigma^{\gl(m)}$ of arity $\mathbf{w}$, the following chain of equalities holds
\allowdisplaybreaks
\begin{align*}
\left(\Tr^{\gl(n,p)}_{\sig{\Sigma}}\right)^{\sharp}\left(\Tr^{\gl(m,n)}_{\sig{\Sigma}}(\tau)\right)
&=
\left(\Tr^{\gl(n,p)}_{\sig{\Sigma}}\right)^{\sharp}\left(\tau\left((v^{\mathbf{w}}_{q})_{q\in\bb{\mathbf{w}}}\right)\right)
\tag{1}
\\&=
\tau\left((v^{\mathbf{w}}_{q})_{q\in\bb{\mathbf{w}}}\right)
\tag{2}
\\&=
\Tr^{\gl(m,p)}_{\sig{\Sigma}}(\tau).
\tag{3}
\end{align*}

The first equality unravels Definition~\ref{DTruncDerivor}; the second equality follows from~(Der1) and~(Der2), taking into account that $\Tr^{\gl(n,p)}_{\sig{\Sigma}}(\tau)=\tau((v^{\mathbf{w}}_{q})_{q\in\bb{\mathbf{w}}})$; finally, the last equality unravels Definition~\ref{DTruncDerivor}. Thus the underlying derivor of $\bTr^{\gl(n,p)}_{\sig{\Sigma}}\circ\bTr^{\gl(m,n)}_{\sig{\Sigma}}$ is that of $\bTr^{\gl(m,p)}_{\sig{\Sigma}}$.

\textsf{(Globular case: domains)} By Definitions~\ref{DHepComp} and~\ref{DTruncDerivor}, the domains of $\bTr^{\gl(n,p)}_{\sig{\Sigma}}\circ\bTr^{\gl(m,n)}_{\sig{\Sigma}}$ are $\varnothing\cup(\Tr^{\gl(n,p)}_{\sig{\Sigma}})^{\sharp}[\varnothing]=\varnothing$, which are the domains of $\bTr^{\gl(m,p)}_{\sig{\Sigma}}$. Hence $\bTr^{\gl(n,p)}_{\sig{\Sigma}}\circ\bTr^{\gl(m,n)}_{\sig{\Sigma}}=\bTr^{\gl(m,p)}_{\sig{\Sigma}}$.

\textsf{(Globular case: functors)} Since the domains of $\bTr^{\gl(m,n)}_{\sig{\Sigma}}$ are empty, by the note after Proposition~\ref{PHepComp}, Proposition~\ref{PTruncGlDerivor} and Proposition~\ref{PHepComp}~(2), for every globular $p$-categorial $\sig{\Sigma}$-algebra $\calg{B}$, $\Tr^{\gl(m,n)}_{\sig{\Sigma}}(\Tr^{\gl(n,p)}_{\sig{\Sigma}}(\calg{B}))=(\bTr^{\gl(n,p)}_{\sig{\Sigma}}\circ\bTr^{\gl(m,n)}_{\sig{\Sigma}})^{\ast}_{\gd}(\calg{B})=\Tr^{\gl(m,p)}_{\sig{\Sigma}}(\calg{B})$, and similarly for the homomorphisms.

This completes the proof.
\end{proof}

\begin{proposition}\label{PDiscComp}
Let $\sig{\Sigma}=(S,\Sigma)$ be a many-sorted signature and $m,n,p\in\mathbb{N}$ with $m\leq n\leq p$. Then the following equalities hold in $\mathsf{hepSig}_{\gd}$
$$
\bDs^{(n,m)}_{\sig{\Sigma}}\circ\bDs^{(p,n)}_{\sig{\Sigma}}=\bDs^{(p,m)}_{\sig{\Sigma}}
\qquad\text{and}\qquad
\bDs^{\gl(n,m)}_{\sig{\Sigma}}\circ\bDs^{\gl(p,n)}_{\sig{\Sigma}}=\bDs^{\gl(p,m)}_{\sig{\Sigma}}.
$$
Consequently, the following equalities of functors hold
$$
\Ds^{(p,n)}_{\sig{\Sigma}}\circ\Ds^{(n,m)}_{\sig{\Sigma}}=\Ds^{(p,m)}_{\sig{\Sigma}}
\qquad\text{and}\qquad
\Ds^{\gl(p,n)}_{\sig{\Sigma}}\circ\Ds^{\gl(n,m)}_{\sig{\Sigma}}=\Ds^{\gl(p,m)}_{\sig{\Sigma}}.
$$
\end{proposition}
\begin{proof}
If two of the dimensions are equal, the statement follows from Definition~\ref{DDiscGlDerivor} and Propositions~\ref{PHepSigCat} and~\ref{PHepId}. Let $m<n<p$. We compare the underlying derivors and the domains. We consider first the single-set case and then the globular case.

\textsf{(Single-set case: underlying derivors)} The mappings on sorts are $\id_{S}$. By~(Der1) and~(Der2) and Definition~\ref{DDiscDerivor}, $(\Ds^{(n,m)}_{\sig{\Sigma}})^{\sharp}\circ\Ds^{(p,n)}_{\sig{\Sigma}}$ is $\Ds^{(p,m)}_{\sig{\Sigma}}$: the operation symbols of $\Sigma$ and those of dimension $k<m$ are sent to themselves, applied to variables, since both $\Ds^{(p,n)}_{\sig{\Sigma}}$ and $\Ds^{(n,m)}_{\sig{\Sigma}}$ do so; those of dimension $m\leq k<n$ are sent by $\Ds^{(p,n)}_{\sig{\Sigma}}$ to themselves, applied to variables, and then by $\Ds^{(n,m)}_{\sig{\Sigma}}$ to $v^{s}_{0}$; and those of dimension $k\geq n$ are sent by $\Ds^{(p,n)}_{\sig{\Sigma}}$ to $v^{s}_{0}$, which, by~(Der1), $(\Ds^{(n,m)}_{\sig{\Sigma}})^{\sharp}$ sends to $v^{s}_{0}$.

\textsf{(Single-set case: domains)} By Definitions~\ref{DHepComp} and~\ref{DDiscDerivor}, the domains of $\bDs^{(n,m)}_{\sig{\Sigma}}\circ\bDs^{(p,n)}_{\sig{\Sigma}}$ are $\varnothing\cup(\Ds^{(n,m)}_{\sig{\Sigma}})^{\sharp}[\varnothing]=\varnothing$, which are the domains of $\bDs^{(p,m)}_{\sig{\Sigma}}$. Hence $\bDs^{(n,m)}_{\sig{\Sigma}}\circ\bDs^{(p,n)}_{\sig{\Sigma}}=\bDs^{(p,m)}_{\sig{\Sigma}}$.

\textsf{(Globular case: mappings on sorts)} The mappings on sorts of $\bDs^{\gl(p,n)}_{\sig{\Sigma}}$ and $\bDs^{\gl(n,m)}_{\sig{\Sigma}}$ send $(k,s)$ to $(\min\{k,n\},s)$ and to $(\min\{k,m\},s)$, respectively. Since $m<n$, their composite sends $(k,s)$ to $(\min\{\min\{k,n\},m\},s)=(\min\{k,m\},s)$, i.e., it is the mapping on sorts of $\bDs^{\gl(p,m)}_{\sig{\Sigma}}$.

\textsf{(Globular case: terms)} Let $j<k$ be dimensions in $p+1$ and $s\in S$. We compute $(\Ds^{\gl(n,m)}_{\sig{\Sigma}})^{\sharp}\circ\Ds^{\gl(p,n)}_{\sig{\Sigma}}$ by~(Der1) and~(Der2) and Definition~\ref{DDiscGlDerivor}, distinguishing the position of $j$.
\begin{enumerate}
\item For every $\sigma\in\Sigma_{\mathbf{s},s}$, $\Ds^{\gl(p,n)}_{\sig{\Sigma}}$ sends $\sigma_{k}$ to $\sigma_{\min\{k,n\}}$, applied to variables, and $\Ds^{\gl(n,m)}_{\sig{\Sigma}}$ sends $\sigma_{\min\{k,n\}}$ to $\sigma_{\min\{k,m\}}$, applied to variables. Thus $(\Ds^{\gl(n,m)}_{\sig{\Sigma}})^{\sharp}(\Ds^{\gl(p,n)}_{\sig{\Sigma}}(\sigma_{k}))=\Ds^{\gl(p,m)}_{\sig{\Sigma}}(\sigma_{k})$.
\item If $j<m$, then $j<n$, and $\Ds^{\gl(p,n)}_{\sig{\Sigma}}$ sends $\src^{(j,k)}_{s}$, $\tgt^{(j,k)}_{s}$, $\mathrm{i}^{(k,j)}_{s}$ and $\circ^{(j,k)}_{s}$ to the corresponding operation symbols of dimensions $j$ and $\min\{k,n\}$, applied to variables, which $\Ds^{\gl(n,m)}_{\sig{\Sigma}}$ sends to the corresponding operation symbols of dimensions $j$ and $\min\{k,m\}$, applied to variables. For instance, the following chain of equalities holds
\allowdisplaybreaks
\begin{align*}
\left(\Ds^{\gl(n,m)}_{\sig{\Sigma}}\right)^{\sharp}\left(\Ds^{\gl(p,n)}_{\sig{\Sigma}}\left(\src^{(j,k)}_{s}\right)\right)
&=
\left(\Ds^{\gl(n,m)}_{\sig{\Sigma}}\right)^{\sharp}\left(\src^{(j,\min\{k,n\})}_{s}\left(v^{(\min\{k,n\},s)}_{0}\right)\right)
\tag{1}
\\&=
\src^{(j,\min\{k,m\})}_{s}\left(v^{(\min\{k,m\},s)}_{0}\right)
\tag{2}
\\&=
\Ds^{\gl(p,m)}_{\sig{\Sigma}}\left(\src^{(j,k)}_{s}\right).
\tag{3}
\end{align*}
The first equality unravels Definition~\ref{DDiscGlDerivor} for $\bDs^{\gl(p,n)}_{\sig{\Sigma}}$; the second equality follows from~(Der1) and~(Der2) and Definition~\ref{DDiscGlDerivor} for $\bDs^{\gl(n,m)}_{\sig{\Sigma}}$; finally, the last equality unravels Definition~\ref{DDiscGlDerivor} for $\bDs^{\gl(p,m)}_{\sig{\Sigma}}$. By a similar argument, the same holds for $\tgt^{(j,k)}_{s}$, $\mathrm{i}^{(k,j)}_{s}$ and $\circ^{(j,k)}_{s}$.
\item If $m\leq j<n$, then $\Ds^{\gl(p,n)}_{\sig{\Sigma}}$ sends $\src^{(j,k)}_{s}$, $\tgt^{(j,k)}_{s}$, $\mathrm{i}^{(k,j)}_{s}$ and $\circ^{(j,k)}_{s}$ to the corresponding operation symbols of dimensions $j$ and $\min\{k,n\}$, applied to variables, which, since $j\geq m$, $\Ds^{\gl(n,m)}_{\sig{\Sigma}}$ sends to variables of sort $(m,s)$. These are the values of $\Ds^{\gl(p,m)}_{\sig{\Sigma}}$, since $j\geq m$.
\item If $j\geq n$, then $\Ds^{\gl(p,n)}_{\sig{\Sigma}}$ sends these operation symbols to variables of sort $(n,s)$, which, by~(Der1), $(\Ds^{\gl(n,m)}_{\sig{\Sigma}})^{\sharp}$ sends to variables of sort $(m,s)$. These are the values of $\Ds^{\gl(p,m)}_{\sig{\Sigma}}$, since $j\geq m$.
\end{enumerate}
Thus the underlying derivor of $\bDs^{\gl(n,m)}_{\sig{\Sigma}}\circ\bDs^{\gl(p,n)}_{\sig{\Sigma}}$ is that of $\bDs^{\gl(p,m)}_{\sig{\Sigma}}$.

\textsf{(Globular case: domains)} By Definitions~\ref{DHepComp} and~\ref{DDiscGlDerivor}, the domains of $\bDs^{\gl(n,m)}_{\sig{\Sigma}}\circ\bDs^{\gl(p,n)}_{\sig{\Sigma}}$ are $\varnothing\cup(\Ds^{\gl(n,m)}_{\sig{\Sigma}})^{\sharp}[\varnothing]=\varnothing$, which are the domains of $\bDs^{\gl(p,m)}_{\sig{\Sigma}}$. Hence $\bDs^{\gl(n,m)}_{\sig{\Sigma}}\circ\bDs^{\gl(p,n)}_{\sig{\Sigma}}=\bDs^{\gl(p,m)}_{\sig{\Sigma}}$.

\textsf{(Functors)} Since the domains of the discretization derivors are empty, by the note after Proposition~\ref{PHepComp}, Propositions~\ref{PDiscDerivor} and~\ref{PDiscGlDerivor} and Proposition~\ref{PHepComp}~(2), the functors induced by the above compositions are, on categorial algebras, the composites of the induced functors. Hence the stated equalities of functors follow from the equalities of generalized hep-derivors.

This completes the proof.
\end{proof}

The composite of a discretization followed by a truncation depends on the relative position of the dimensions.

\begin{proposition}\label{PTruncDiscComp}
Let $\sig{\Sigma}=(S,\Sigma)$ be a many-sorted signature and $m,n,p\in\mathbb{N}$ with $m\leq p$ and $n\leq p$. Then the following equality holds
$$
\Tr^{(n,p)}_{\sig{\Sigma}}\circ\Ds^{(p,m)}_{\sig{\Sigma}}=
\begin{cases}
\Ds^{(n,m)}_{\sig{\Sigma}}, & \text{if $m\leq n$;}\\
\Tr^{(n,m)}_{\sig{\Sigma}}, & \text{if $n\leq m$,}
\end{cases}
$$
and the same holds for the globular truncation and discretization functors.
\end{proposition}

\begin{proof}
If $n=p$, the truncation derivor is an identity, and if $m=p$, so is the discretization derivor; in both cases the statement follows from Proposition~\ref{PHepId}. Let $m<p$ and $n<p$. We compare the underlying derivors and the cells, and we conclude by Proposition~\ref{PHepComp}~(2) and Lemma~\ref{LSameCells}. We consider first the single-set case and then the globular case.

\textsf{(Single-set case: underlying derivors)} The mappings on sorts are $\id_{S}$. Since $\Tr^{(n,p)}_{\sig{\Sigma}}$ sends every operation symbol $\tau$ of $\Sigma^{(n)}$ to itself, applied to variables, by~(Der1) and~(Der2), $(\Ds^{(p,m)}_{\sig{\Sigma}})^{\sharp}(\Tr^{(n,p)}_{\sig{\Sigma}}(\tau))=\Ds^{(p,m)}_{\sig{\Sigma}}(\tau)$. By Definition~\ref{DDiscDerivor}, this term is $\tau$, applied to variables, if $\tau$ is an operation symbol of $\Sigma$ or of dimension $k<m$, and $v^{s}_{0}$, if $\tau$ is of dimension $k\geq m$. If $m\leq n$, this is $\Ds^{(n,m)}_{\sig{\Sigma}}(\tau)$; if $n\leq m$, then all the dimensions of $\Sigma^{(n)}$ are smaller than $m$, and this is $\Tr^{(n,m)}_{\sig{\Sigma}}(\tau)$.

\textsf{(Single-set case: cells)} Let $\calg{A}$ be an $m$-categorial $\sig{\Sigma}$-algebra. By the note after Definition~\ref{DTruncDerivor}, the cells of $\Tr^{(n,p)}_{\sig{\Sigma}}(\Ds^{(p,m)}_{\sig{\Sigma}}(\calg{A}))$ are the cells $x$ of $\calg{A}$ such that $\src^{n}_{s}(x)=x$ in $\Ds^{(p,m)}_{\sig{\Sigma}}(\calg{A})$. If $m\leq n$, then, by Lemma~\ref{LDiscDescription}, $\src^{n}_{s}$ acts as the identity, so that they are all the cells of $\calg{A}$, which are the cells of $\Ds^{(n,m)}_{\sig{\Sigma}}(\calg{A})$. If $n<m$, then $\src^{n}_{s}$ acts as in $\calg{A}$, so that they are the $n$-cells of $\calg{A}$, which are the cells of $\Tr^{(n,m)}_{\sig{\Sigma}}(\calg{A})$. On the other hand, by Definitions~\ref{DHepComp}, \ref{DDiscDerivor} and~\ref{DTruncDerivor}, the domain of sort $s$ of $\bDs^{(p,m)}_{\sig{\Sigma}}\circ\bTr^{(n,p)}_{\sig{\Sigma}}$ is $(\Ds^{(p,m)}_{\sig{\Sigma}})^{\sharp}[\{(v^{s}_{0},\src^{n}_{s}(v^{s}_{0}))\}]$, i.e., $\{(v^{s}_{0},v^{s}_{0})\}$, if $m\leq n$, and $\{(v^{s}_{0},\src^{n}_{s}(v^{s}_{0}))\}$, if $n<m$, so that the cells of $(\bDs^{(p,m)}_{\sig{\Sigma}}\circ\bTr^{(n,p)}_{\sig{\Sigma}})^{\ast}_{\gd}(\calg{A})$ are the same. Moreover, by Propositions~\ref{PDiscDerivor} and~\ref{PTruncDerivor}, $\Tr^{(n,p)}_{\sig{\Sigma}}(\Ds^{(p,m)}_{\sig{\Sigma}}(\calg{A}))$ is a $\sig{\Sigma}^{(n)}$-algebra. Thus, by Proposition~\ref{PHepComp}~(2), Lemma~\ref{LSameCells} and the comparison of the underlying derivors, the single-set equality holds, on objects and, by the same results, on morphisms.

\textsf{(Globular case: mappings on sorts)} The composite of the mappings on sorts of $\bTr^{\gl(n,p)}_{\sig{\Sigma}}$ and $\bDs^{\gl(p,m)}_{\sig{\Sigma}}$ sends $(k,s)$, for $k\leq n$, to $(\min\{k,m\},s)$. If $m\leq n$, it is the mapping on sorts of $\bDs^{\gl(n,m)}_{\sig{\Sigma}}$; if $n\leq m$, then $\min\{k,m\}=k$, and it is the mapping on sorts of $\bTr^{\gl(n,m)}_{\sig{\Sigma}}$.

\textsf{(Globular case: terms)} Since $\Tr^{\gl(n,p)}_{\sig{\Sigma}}$ sends every operation symbol $\tau$ of $\Sigma^{\gl(n)}$ to itself, applied to variables, by~(Der1) and~(Der2), $(\Ds^{\gl(p,m)}_{\sig{\Sigma}})^{\sharp}(\Tr^{\gl(n,p)}_{\sig{\Sigma}}(\tau))=\Ds^{\gl(p,m)}_{\sig{\Sigma}}(\tau)$. If $m\leq n$, then, by Definition~\ref{DDiscGlDerivor}, the clauses that define $\Ds^{\gl(p,m)}_{\sig{\Sigma}}(\tau)$, for the operation symbols of $\Sigma^{\gl(n)}$, are those that define $\Ds^{\gl(n,m)}_{\sig{\Sigma}}(\tau)$, since both only depend on the dimensions involved and on $m$. If $n\leq m$, then all the dimensions of $\Sigma^{\gl(n)}$ are at most $m$, so that, by Definition~\ref{DDiscGlDerivor}, $\Ds^{\gl(p,m)}_{\sig{\Sigma}}(\tau)$ is $\tau$, applied to variables, i.e., $\Tr^{\gl(n,m)}_{\sig{\Sigma}}(\tau)$.

\textsf{(Globular case: cells)} Let $\calg{A}$ be a globular $m$-categorial $\sig{\Sigma}$-algebra. Since the domains of the globular truncation and discretization derivors are empty, by the note after Proposition~\ref{PHepComp} and Propositions~\ref{PDiscGlDerivor}, \ref{PTruncGlDerivor} and~\ref{PHepComp}~(2), $\Tr^{\gl(n,p)}_{\sig{\Sigma}}(\Ds^{\gl(p,m)}_{\sig{\Sigma}}(\calg{A}))=(\bDs^{\gl(p,m)}_{\sig{\Sigma}}\circ\bTr^{\gl(n,p)}_{\sig{\Sigma}})^{\ast}_{\gd}(\calg{A})$. By the note after Definition~\ref{DTruncDerivor} and Lemma~\ref{LDiscGlDescription}, its cells of sort $(k,s)$, for $k\leq n$, are the $\min\{k,m\}$-cells of sort $s$ of $\calg{A}$. If $m\leq n$, they are the cells of sort $(k,s)$ of $\Ds^{\gl(n,m)}_{\sig{\Sigma}}(\calg{A})$; if $n\leq m$, they are the $k$-cells of $\calg{A}$, i.e., the cells of sort $(k,s)$ of $\Tr^{\gl(n,m)}_{\sig{\Sigma}}(\calg{A})$. Thus, by Lemma~\ref{LSameCells} and the comparison of the underlying derivors, the globular equality holds, on objects and on morphisms.

This completes the proof.
\end{proof}

\begin{remark}\label{RTruncDiscCompEqual}
Let $\sig{\Sigma}=(S,\Sigma)$ be a many-sorted signature and $n,p\in\mathbb{N}$ with $n\leq p$. In Proposition~\ref{PTruncDiscComp}, the two cases overlap when $m=n$. In this case both expressions agree, since, by Definitions~\ref{DDiscGlDerivor} and~\ref{DTruncDerivor} and Proposition~\ref{PHepId}, $\Ds^{(n,n)}_{\sig{\Sigma}}=\Tr^{(n,n)}_{\sig{\Sigma}}=\Id_{\mathsf{nCatAlg}(\sig{\Sigma})}$ and $\Ds^{\gl(n,n)}_{\sig{\Sigma}}=\Tr^{\gl(n,n)}_{\sig{\Sigma}}=\Id_{\mathsf{nCat}^{\gl}\mathsf{Alg}(\sig{\Sigma})}$. Thus, for $m=n$, Proposition~\ref{PTruncDiscComp} states that
$$
\Tr^{(n,p)}_{\sig{\Sigma}}\circ\Ds^{(p,n)}_{\sig{\Sigma}}=\Id_{\mathsf{nCatAlg}(\sig{\Sigma})}
\qquad\text{and}\qquad
\Tr^{\gl(n,p)}_{\sig{\Sigma}}\circ\Ds^{\gl(p,n)}_{\sig{\Sigma}}=\Id_{\mathsf{nCat}^{\gl}\mathsf{Alg}(\sig{\Sigma})},
$$
i.e., it recovers Propositions~\ref{PTruncDisc} and~\ref{PTruncDiscGl}, with $n$ and $p$ in place of $m$ and $n$, respectively. Moreover, when $n=p$, all the functors involved are identities.
\end{remark}

\subsection{Compatibility with graduation and unification}\label{SCompatGdUf}

In this subsection we prove that the functors of this section commute with graduation and unification, strictly or up to a natural isomorphism. These facts are used in Sections~\ref{SOmegaFunctors} and~\ref{SAlgSquare}. All the functors involved act on homomorphisms by restriction to the corresponding cells, so that, in order to prove the equalities between them, it suffices to compare their values on objects. Moreover, for $m=n$ all the statements are trivial, since all the truncation and discretization functors involved are identities.

We begin with the compatibility of truncation with graduation.

\begin{proposition}\label{PCompatGdTr}
Let $\sig{\Sigma}=(S,\Sigma)$ be a many-sorted signature and $m,n\in\mathbb{N}$ with $m\leq n$. Then the following equality holds
$$
\Tr^{\gl(m,n)}_{\sig{\Sigma}}\circ\Gd^{(n)}_{\sig{\Sigma}}=\Gd^{(m)}_{\sig{\Sigma}}\circ \Tr^{(m,n)}_{\sig{\Sigma}}.
$$
\end{proposition}
\begin{proof}
Let $\calg{A}$ be an $n$-categorial $\sig{\Sigma}$-algebra. By Lemma~\ref{LGradDescription}, both globular $m$-categorial algebras have as cells of sort $(k,s)$ the $k$-cells of sort $s$ of $\calg{A}$, for $k\leq m$, because, by Lemma~\ref{LGradCells}~(3), every $k$-cell is an $m$-cell; and, in both of them, the operations are the restrictions of those of $\calg{A}$ given in Definition~\ref{DGradDerivor}.
\end{proof}

We next consider the compatibility of discretization with graduation.

\begin{proposition}\label{PCompatGdDs}
Let $\sig{\Sigma}=(S,\Sigma)$ be a many-sorted signature and $m,n\in\mathbb{N}$ with $m\leq n$. Then the following equality holds
$$
\Gd^{(n)}_{\sig{\Sigma}}\circ\Ds^{(n,m)}_{\sig{\Sigma}}=\Ds^{\gl(n,m)}_{\sig{\Sigma}}\circ\Gd^{(m)}_{\sig{\Sigma}}.
$$
\end{proposition}
\begin{proof}
Let $\calg{A}$ be an $m$-categorial $\sig{\Sigma}$-algebra. By Lemmas~\ref{LGradDescription}, \ref{LDiscDescription} and~\ref{LDiscGlDescription}, both globular $n$-categorial algebras have as cells of sort $(k,s)$ the $k$-cells of $\calg{A}$, for $k<m$, and all the cells of sort $s$ of $\calg{A}$, for $m\leq k\leq n$. In both of them, for every pair of dimensions $j<k$,
\begin{enumerate}
\item if $j<m$, the $(j,k)$-sources, $(j,k)$-targets and $j$-compositions are those of $\calg{A}$, and the $(k,j)$-identities are inclusions;
\item if $j\geq m$, the $(j,k)$-sources, $(j,k)$-targets and $(k,j)$-identities are identities and the $j$-composition of two cells is defined if, and only if, they are equal; and
\item $\sigma_{k}$ acts as $\sigma$.
\end{enumerate}
Hence both algebras are equal.
\end{proof}

We next consider the compatibility of discretization with unification.

\begin{proposition}\label{PCompatUfDs}
Let $\sig{\Sigma}=(S,\Sigma)$ be a many-sorted signature and $m,n\in\mathbb{N}$ with $m\leq n$. Then the following equality holds
$$
\Uf^{(n)}_{\sig{\Sigma}}\circ\Ds^{\gl(n,m)}_{\sig{\Sigma}}=\Ds^{(n,m)}_{\sig{\Sigma}}\circ\Uf^{(m)}_{\sig{\Sigma}}.
$$
\end{proposition}
\begin{proof}
Let $\calg{B}$ be a globular $m$-categorial $\sig{\Sigma}$-algebra. By Lemmas~\ref{LUnifDescription}, \ref{LDiscDescription} and~\ref{LDiscGlDescription}, both $n$-categorial algebras have as cells of sort $s$ the $m$-cells of sort $s$ of $\calg{B}$. In both of them,
\begin{enumerate}
\item for $k<m$, $\src^{k}_{s}=\mathrm{i}^{(m,k)}_{s}\circ\src^{(k,m)}_{s}$, $\tgt^{k}_{s}=\mathrm{i}^{(m,k)}_{s}\circ\tgt^{(k,m)}_{s}$ and $\circ^{k}_{s}=\circ^{(k,m)}_{s}$;
\item for $m\leq k<n$, $\src^{k}_{s}$ and $\tgt^{k}_{s}$ are identities and $x\circ^{k}_{s}y$ is defined if, and only if, $x=y$; and
\item $\sigma$ acts as $\sigma_{m}$.
\end{enumerate}
Hence both algebras are equal.
\end{proof}

Finally, truncation commutes with unification only up to a natural isomorphism.

\begin{proposition}\label{PCompatUfTr}
Let $\sig{\Sigma}=(S,\Sigma)$ be a many-sorted signature and $m,n\in\mathbb{N}$ with $m\leq n$. For every globular $n$-categorial $\sig{\Sigma}$-algebra $\calg{B}$, let $\beta^{(m,n)}_{\calg{B}}$ stand for the family $(\mathrm{i}^{(n,m)}_{s})_{s\in S}$, where $\mathrm{i}^{(n,n)}_{s}$ is the identity. Then $\beta^{(m,n)}_{\calg{B}}$ is an isomorphism from $\Uf^{(m)}_{\sig{\Sigma}}(\Tr^{\gl(m,n)}_{\sig{\Sigma}}(\calg{B}))$ to $\Tr^{(m,n)}_{\sig{\Sigma}}(\Uf^{(n)}_{\sig{\Sigma}}(\calg{B}))$, and $\beta^{(m,n)}=(\beta^{(m,n)}_{\calg{B}})_{\calg{B}\in\mathrm{Ob}(\mathsf{nCat}^{\gl}\mathsf{Alg}(\sig{\Sigma}))}$ is a natural isomorphism from $\Uf^{(m)}_{\sig{\Sigma}}\circ \Tr^{\gl(m,n)}_{\sig{\Sigma}}$ to $\Tr^{(m,n)}_{\sig{\Sigma}}\circ\Uf^{(n)}_{\sig{\Sigma}}$. Thus
$$
\Uf^{(m)}_{\sig{\Sigma}}\circ \Tr^{\gl(m,n)}_{\sig{\Sigma}}\cong \Tr^{(m,n)}_{\sig{\Sigma}}\circ\Uf^{(n)}_{\sig{\Sigma}}.
$$
\end{proposition}
\begin{proof}
For $m=n$ there is nothing to prove. Let $m<n$. We prove the following claims.

\begin{claim}\label{CCompatUfTrA}
$\beta^{(m,n)}_{\calg{B}}$ is a bijection.
\end{claim}

By Lemma~\ref{LUnifDescription}, the cells of sort $s$ of $\Tr^{(m,n)}_{\sig{\Sigma}}(\Uf^{(n)}_{\sig{\Sigma}}(\calg{B}))$ are the $n$-cells $x$ with $x=\mathrm{i}^{(n,m)}_{s}(\src^{(m,n)}_{s}(x))$, i.e., by Lemma~\ref{LUnifInj}, the elements of the image of $\mathrm{i}^{(n,m)}_{s}$, which is injective.

Claim~\ref{CCompatUfTrA} follows.

\begin{claim}\label{CCompatUfTrB}
$\beta^{(m,n)}_{\calg{B}}$ is a homomorphism that reflects definedness.
\end{claim}

\textsf{(Operations of $\Sigma$)} $\beta^{(m,n)}_{\calg{B}}$ commutes with them, by the third equation of~(EA$^{m,n}$7).

\textsf{(Sources and targets)} Let $k<m$ and $u$ an $m$-cell of sort $s$. Then the following chain of equalities holds
\allowdisplaybreaks
\begin{align*}
\mathrm{i}^{(n,k)}_{s}\left(\src^{(k,n)}_{s}\left(\mathrm{i}^{(n,m)}_{s}(u)\right)\right)
&=
\mathrm{i}^{(n,k)}_{s}\left(\src^{(k,m)}_{s}(u)\right)
\tag{1}
\\&=
\mathrm{i}^{(n,m)}_{s}\left(\mathrm{i}^{(m,k)}_{s}\left(\src^{(k,m)}_{s}(u)\right)\right).
\tag{2}
\end{align*}

The first equality follows from~(EA$^{k,m,n}$1) and Lemma~\ref{LUnifInj}; finally, the last equality follows from~(EA$^{k,m,n}$2). Thus, by Lemma~\ref{LUnifDescription}, $\beta^{(m,n)}_{\calg{B}}$ commutes with the $k$-sources. By a similar argument, it commutes with the $k$-targets.

\textsf{(Compositions)} Let $k<m$ and $u$, $u'$ $m$-cells of sort $s$. Since $\mathrm{i}^{(n,m)}_{s}$ is injective, by the above computation, $\beta^{(m,n)}_{\calg{B}}(u)\circ^{k}_{s}\beta^{(m,n)}_{\calg{B}}(u')$ is defined if, and only if, $u\circ^{k}_{s}u'$ is defined. In this case, by~(EA$^{k,m,n}$5), the following equality holds
$$
\mathrm{i}^{(n,m)}_{s}\left(u\circ^{(k,m)}_{s}u'\right)=\mathrm{i}^{(n,m)}_{s}(u)\circ^{(k,n)}_{s}\mathrm{i}^{(n,m)}_{s}(u').
$$

Claim~\ref{CCompatUfTrB} follows.

By Claims~\ref{CCompatUfTrA} and~\ref{CCompatUfTrB}, $\beta^{(m,n)}_{\calg{B}}$ is a bijective homomorphism that reflects definedness, i.e., an isomorphism. Finally, it is natural in $\calg{B}$, because the globular $n$-categorial $\sig{\Sigma}$-homomorphisms commute with the identities.

This completes the proof.
\end{proof}

\subsection{Truncation and discretization in dimension \texorpdfstring{$\omega$}{omega}}\label{SDTOmega}

In this subsection we define the truncation functors from the categories of $\omega$-categorial algebras to the categories of $n$-categorial algebras, and the discretization functors in the opposite direction, in the single-set and in the globular presentations. As in the finite-dimensional case, they are induced by generalized hep-derivors, which are defined by the same clauses as in Definitions~\ref{DDiscDerivor}, \ref{DDiscGlDerivor} and~\ref{DTruncDerivor}, with $\omega$ in place of the greatest dimension. In Sections~\ref{SLimitSingleSet} and~\ref{SLimitGlobular} we prove that the truncation functors are the projections of projective limits.

Let us recall that, by Remarks~\ref{RSingleSetIdentOmega} and~\ref{RGlobularIdent}, for every many-sorted signature $\sig{\Sigma}=(S,\Sigma)$, $\omega\mathsf{CatAlg}(\sig{\Sigma})$ is the full subcategory of $\mathsf{EPAlg}(\sig{\Sigma}^{(\omega)},\mathcal{C}^{\sig{\Sigma}^{(\omega)}})$ whose objects are the $\sig{\Sigma}^{(\omega)}$-algebras in which every cell is a $k$-cell for some $k\in\omega$, and $\omega\mathsf{Cat}^{\gl}\mathsf{Alg}(\sig{\Sigma})=\mathsf{EPAlg}(\sig{\Sigma}^{\gl(\omega)},\mathcal{C}^{\sig{\Sigma}^{\gl(\omega)}})$.

We begin by defining the truncation derivors in dimension $\omega$.

\begin{definition}\label{DTruncOmegaDerivor}
Let $\sig{\Sigma}=(S,\Sigma)$ be a many-sorted signature and $n\in\mathbb{N}$.
\begin{enumerate}
\item The \emph{single-set truncation derivor} in dimension $\omega$ is the generalized hep-derivor
$$
\bTr^{(n,\omega)}_{\sig{\Sigma}}=\left(\varphi^{\Tr^{(n,\omega)}_{\sig{\Sigma}}},\Tr^{(n,\omega)}_{\sig{\Sigma}},\Dom^{\Tr^{(n,\omega)}_{\sig{\Sigma}}}\right)
$$
from $\sig{\Sigma}^{(n)}$ to $\sig{\Sigma}^{(\omega)}$ defined as follows:
\begin{enumerate}
\item the mapping on sorts is $\varphi^{\Tr^{(n,\omega)}_{\sig{\Sigma}}}=\id_{S}$;
\item the assignment of terms is the inclusion, i.e., for every operation symbol $\tau$ of $\Sigma^{(n)}$ of arity $\mathbf{s}$,
$$
\Tr^{(n,\omega)}_{\sig{\Sigma}}(\tau)=\tau((v^{\mathbf{s}}_{p})_{p\in\bb{\mathbf{s}}});
$$
and
\item for every $s\in S$, the domain is
$$
\Dom^{\Tr^{(n,\omega)}_{\sig{\Sigma}}}_{s}=\left\lbrace\left(v^{s}_{0},\src^{n}_{s}(v^{s}_{0})\right)\right\rbrace.
$$
\end{enumerate}
\item The \emph{globular truncation derivor} in dimension $\omega$ is the generalized hep-derivor
$$
\bTr^{\gl(n,\omega)}_{\sig{\Sigma}}=\left(\varphi^{\Tr^{\gl(n,\omega)}_{\sig{\Sigma}}},\Tr^{\gl(n,\omega)}_{\sig{\Sigma}},\Dom^{\Tr^{\gl(n,\omega)}_{\sig{\Sigma}}}\right)
$$
from $\sig{\Sigma}^{\gl(n)}$ to $\sig{\Sigma}^{\gl(\omega)}$ determined by the canonical embedding of $\Sigma^{\gl(n)}$ into $\Sigma^{\gl(\omega)}$, see the note after Definition~\ref{DSigmaCatGlobular}, i.e., defined as follows:
\begin{enumerate}
\item for every $(k,s)\in(n+1)\times S$, the mapping on sorts sends $(k,s)$ to $\varphi^{\Tr^{\gl(n,\omega)}_{\sig{\Sigma}}}(k,s)=(k,s)$;
\item for every operation symbol $\tau$ of $\Sigma^{\gl(n)}$ of arity $\mathbf{w}$,
$$
\Tr^{\gl(n,\omega)}_{\sig{\Sigma}}(\tau)=\tau((v^{\mathbf{w}}_{p})_{p\in\bb{\mathbf{w}}});
$$
and
\item for every $(k,s)\in(n+1)\times S$, the domain is $\Dom^{\Tr^{\gl(n,\omega)}_{\sig{\Sigma}}}_{(k,s)}=\varnothing$.
\end{enumerate}
\end{enumerate}
\end{definition}

As in the note after Definition~\ref{DTruncDerivor}, the cells of a partial $\Sigma^{(\omega)}$-algebra $\mathbf{B}$ relative to $\bTr^{(n,\omega)}_{\sig{\Sigma}}$ are its $n$-cells, and the cells of sort $(k,s)$ of a partial $\Sigma^{\gl(\omega)}$-algebra $\mathbf{B}$ relative to $\bTr^{\gl(n,\omega)}_{\sig{\Sigma}}$ are all the elements of $B_{(k,s)}$.

In the following proposition we prove that the truncation derivors in dimension $\omega$ send categorial algebras to categorial algebras.

\begin{proposition}\label{PTruncOmegaDerivor}
Let $\sig{\Sigma}=(S,\Sigma)$ be a many-sorted signature and $n\in\mathbb{N}$. Then the restriction of $(\bTr^{(n,\omega)}_{\sig{\Sigma}})^{\ast}_{\gd}$ to $\mathsf{EPAlg}(\sig{\Sigma}^{(\omega)},\mathcal{C}^{\sig{\Sigma}^{(\omega)}})$ corestricts to $\mathsf{EPAlg}(\sig{\Sigma}^{(n)},\mathcal{C}^{\sig{\Sigma}^{(n)}})$, and the restriction of $(\bTr^{\gl(n,\omega)}_{\sig{\Sigma}})^{\ast}_{\gd}$ to $\mathsf{EPAlg}(\sig{\Sigma}^{\gl(\omega)},\mathcal{C}^{\sig{\Sigma}^{\gl(\omega)}})$ corestricts to $\mathsf{EPAlg}(\sig{\Sigma}^{\gl(n)},\mathcal{C}^{\sig{\Sigma}^{\gl(n)}})$, see Definition~\ref{DCorestriction}.
\end{proposition}
\begin{proof}
By Definitions~\ref{DCorestriction} and~\ref{DEpVariety}, in each case we have to prove that the transformed algebras are ep-algebras in which the axioms of the corresponding ep-variety hold.

\textsf{(Single-set case: ep-algebra)} Let $\mathbf{B}$ be an object in $\mathsf{EPAlg}(\sig{\Sigma}^{(\omega)},\mathcal{C}^{\sig{\Sigma}^{(\omega)}})$. By the above note and Proposition~\ref{PTransformedAlg}, the cells of $(\bTr^{(n,\omega)}_{\sig{\Sigma}})^{\ast}_{\gd}(\mathbf{B})$ are the $n$-cells of $\mathbf{B}$, and its operations are restrictions of those of $\mathbf{B}$. The proof of Lemma~\ref{LGradCells} only uses, for dimensions $j<k$, the axioms (A$^{k}$7), (A$^{j}$2), (A$^{j}$3), (AB$^{j,k}$1) and (AB$^{j,k}$2), which are axioms of $\boldsymbol{\mathcal{C}}^{\sig{\Sigma}^{(\omega)}}$ for every $j<k$ in $\omega$. Hence, as in the proof of Lemma~\ref{LTruncCells}, the following statements hold, for every $s\in S$ and every $k<n$:
\begin{enumerate}
\item for every $\sigma\in\Sigma_{\mathbf{s},s}$ and every family $(x_{p})_{p\in\bb{\mathbf{s}}}$ of $n$-cells, $\sigma((x_{p})_{p\in\bb{\mathbf{s}}})$ is an $n$-cell;
\item for every cell $x$ of sort $s$, $\src^{k}_{s}(x)$ and $\tgt^{k}_{s}(x)$ are $n$-cells; and
\item for every pair of $n$-cells $x$, $y$ of sort $s$ such that $\src^{k}_{s}(x)=\tgt^{k}_{s}(y)$, $x\circ^{k}_{s}y$ is defined and it is an $n$-cell.
\end{enumerate}
Therefore, the $n$-cells are closed under the operations of $\Sigma^{(n)}$, which are defined on the families of $n$-cells that satisfy their domain conditions, by~(A0), (A$^{k}$1) and~(A$^{k}$3). Thus, by Proposition~\ref{PTransformedEpAlg}, $(\bTr^{(n,\omega)}_{\sig{\Sigma}})^{\ast}_{\gd}(\mathbf{B})$ is a $\sig{\Sigma}^{(n)}$-algebra.

\textsf{(Single-set case: axioms)} By the above, $(\bTr^{(n,\omega)}_{\sig{\Sigma}})^{\ast}_{\gd}(\mathbf{B})$ is the relative subalgebra of $\mathbf{B}$ on the $n$-cells, in which the values of terms are their values in $\mathbf{B}$. Therefore, every axiom of $\boldsymbol{\mathcal{C}}^{\sig{\Sigma}^{(n)}}$, which is an axiom of $\boldsymbol{\mathcal{C}}^{\sig{\Sigma}^{(\omega)}}$, holds in $(\bTr^{(n,\omega)}_{\sig{\Sigma}})^{\ast}_{\gd}(\mathbf{B})$.

\textsf{(Globular case: ep-algebra)} Let $\mathbf{B}$ be an object in $\mathsf{EPAlg}(\sig{\Sigma}^{\gl(\omega)},\mathcal{C}^{\sig{\Sigma}^{\gl(\omega)}})$. By the above note, all the cells of dimension at most $n$ of $\mathbf{B}$ are cells of the transformed algebra, and every operation symbol of $\Sigma^{\gl(n)}$ is sent to itself, applied to variables. Hence the transformed algebra is the reduct of $\mathbf{B}$ to the sorts of dimension at most $n$ and to the operation symbols of $\Sigma^{\gl(n)}$, which is a $\sig{\Sigma}^{\gl(n)}$-algebra, since the domain conditions of the operation symbols of $\Sigma^{\gl(n)}$ are the same in $\sig{\Sigma}^{\gl(n)}$ and in $\sig{\Sigma}^{\gl(\omega)}$.

\textsf{(Globular case: axioms)} The axioms of $\boldsymbol{\mathcal{C}}^{\sig{\Sigma}^{\gl(n)}}$ are axioms of $\boldsymbol{\mathcal{C}}^{\sig{\Sigma}^{\gl(\omega)}}$, which hold in $\mathbf{B}$, and the values of terms in the transformed algebra are their values in $\mathbf{B}$.

This completes the proof.
\end{proof}

We are now in position to define the truncation functors in dimension $\omega$.

\begin{definition}\label{DTruncationOmega}
Let $\sig{\Sigma}=(S,\Sigma)$ be a many-sorted signature and $n\in\mathbb{N}$. The \emph{truncation functor}
$$
\Tr^{(n,\omega)}_{\sig{\Sigma}}\colon\omega\mathsf{CatAlg}(\sig{\Sigma})\mor\mathsf{nCatAlg}(\sig{\Sigma})
$$
is the restriction to $\omega\mathsf{CatAlg}(\sig{\Sigma})$ of the functor $(\bTr^{(n,\omega)}_{\sig{\Sigma}})^{\ast}_{\gd}$ from $\mathsf{EPAlg}(\sig{\Sigma}^{(\omega)},\mathcal{C}^{\sig{\Sigma}^{(\omega)}})$ to $\mathsf{EPAlg}(\sig{\Sigma}^{(n)},\mathcal{C}^{\sig{\Sigma}^{(n)}})=\mathsf{nCatAlg}(\sig{\Sigma})$, see Proposition~\ref{PTruncOmegaDerivor}.
\end{definition}

By the proof of Proposition~\ref{PTruncOmegaDerivor}, $\Tr^{(n,\omega)}_{\sig{\Sigma}}$ sends an $\omega$-categorial $\sig{\Sigma}$-algebra $\calg{A}$ to the $n$-categorial $\sig{\Sigma}$-algebra whose cells of sort $s$ are the $n$-cells of sort $s$ of $\calg{A}$, in which the operations of $\Sigma$ and the $k$-sources, $k$-targets and $k$-compositions, for $k<n$, are the restrictions of those of $\calg{A}$; and, by Proposition~\ref{PGenHepDerivorFunctor}, it sends an $\omega$-categorial $\sig{\Sigma}$-homomorphism to its restriction to the $n$-cells.

\begin{definition}\label{DTruncationOmegaGl}
Let $\sig{\Sigma}=(S,\Sigma)$ be a many-sorted signature and $n\in\mathbb{N}$. The \emph{globular truncation functor}
$$
\Tr^{\gl(n,\omega)}_{\sig{\Sigma}}\colon\omega\mathsf{Cat}^{\gl}\mathsf{Alg}(\sig{\Sigma})\mor\mathsf{nCat}^{\gl}\mathsf{Alg}(\sig{\Sigma})
$$
is the functor $(\bTr^{\gl(n,\omega)}_{\sig{\Sigma}})^{\ast}_{\gd}$ from $\mathsf{EPAlg}(\sig{\Sigma}^{\gl(\omega)},\mathcal{C}^{\sig{\Sigma}^{\gl(\omega)}})=\omega\mathsf{Cat}^{\gl}\mathsf{Alg}(\sig{\Sigma})$ to $\mathsf{EPAlg}(\sig{\Sigma}^{\gl(n)},\mathcal{C}^{\sig{\Sigma}^{\gl(n)}})=\mathsf{nCat}^{\gl}\mathsf{Alg}(\sig{\Sigma})$, see Proposition~\ref{PTruncOmegaDerivor}.
\end{definition}

By the proof of Proposition~\ref{PTruncOmegaDerivor}, $\Tr^{\gl(n,\omega)}_{\sig{\Sigma}}$ sends a globular $\omega$-categorial $\sig{\Sigma}$-algebra $\calg{A}$ to the globular $n$-categorial $\sig{\Sigma}$-algebra whose $k$-cells of sort $s$, for $k\leq n$, are the $k$-cells of sort $s$ of $\calg{A}$, in which the operations of $\Sigma^{\gl(n)}$ are those of $\calg{A}$; and, by Proposition~\ref{PGenHepDerivorFunctor}, it sends a globular $\omega$-categorial $\sig{\Sigma}$-homomorphism $F$ to $(F_{(k,s)})_{(k,s)\in(n+1)\times S}$.

We next define the discretization derivors in dimension $\omega$.

\begin{definition}\label{DDiscOmegaDerivor}
Let $\sig{\Sigma}=(S,\Sigma)$ be a many-sorted signature and $n\in\mathbb{N}$.
\begin{enumerate}
\item The \emph{single-set discretization derivor} in dimension $\omega$ is the generalized hep-derivor
$$
\bDs^{(\omega,n)}_{\sig{\Sigma}}=\left(\varphi^{\Ds^{(\omega,n)}_{\sig{\Sigma}}},\Ds^{(\omega,n)}_{\sig{\Sigma}},\Dom^{\Ds^{(\omega,n)}_{\sig{\Sigma}}}\right)
$$
from $\sig{\Sigma}^{(\omega)}$ to $\sig{\Sigma}^{(n)}$ defined as follows:
\begin{enumerate}
\item the mapping on sorts is $\varphi^{\Ds^{(\omega,n)}_{\sig{\Sigma}}}=\id_{S}$;
\item for every $(\mathbf{s},s)\in S^{\star}\times S$ and every $\sigma\in\Sigma_{\mathbf{s},s}$,
$$
\Ds^{(\omega,n)}_{\sig{\Sigma}}(\sigma)=\sigma((v^{\mathbf{s}}_{p})_{p\in\bb{\mathbf{s}}});
$$
\item for every $k\in n$ and every $s\in S$,
$$
\Ds^{(\omega,n)}_{\sig{\Sigma}}(\src^{k}_{s})=\src^{k}_{s}(v^{s}_{0}),\quad
\Ds^{(\omega,n)}_{\sig{\Sigma}}(\tgt^{k}_{s})=\tgt^{k}_{s}(v^{s}_{0}),\quad
\Ds^{(\omega,n)}_{\sig{\Sigma}}(\circ^{k}_{s})=v^{s}_{0}\circ^{k}_{s}v^{s}_{1};
$$
\item for every $k\in\omega$ with $k\geq n$ and every $s\in S$,
$$
\Ds^{(\omega,n)}_{\sig{\Sigma}}(\src^{k}_{s})=\Ds^{(\omega,n)}_{\sig{\Sigma}}(\tgt^{k}_{s})=\Ds^{(\omega,n)}_{\sig{\Sigma}}(\circ^{k}_{s})=v^{s}_{0};
$$
and
\item for every $s\in S$, the domain is $\Dom^{\Ds^{(\omega,n)}_{\sig{\Sigma}}}_{s}=\varnothing$.
\end{enumerate}
\item The \emph{globular discretization derivor} in dimension $\omega$ is the generalized hep-derivor
$$
\bDs^{\gl(\omega,n)}_{\sig{\Sigma}}=\left(\varphi^{\Ds^{\gl(\omega,n)}_{\sig{\Sigma}}},\Ds^{\gl(\omega,n)}_{\sig{\Sigma}},\Dom^{\Ds^{\gl(\omega,n)}_{\sig{\Sigma}}}\right)
$$
from $\sig{\Sigma}^{\gl(\omega)}$ to $\sig{\Sigma}^{\gl(n)}$ defined by the clauses of Definition~\ref{DDiscGlDerivor}, with $\omega$ in place of $n+1$ and $n$ in place of $m$, i.e., defined as follows:
\begin{enumerate}
\item for every $k\in\omega$ and every $s\in S$, the mapping on sorts sends $(k,s)$ to
$$
\varphi^{\Ds^{\gl(\omega,n)}_{\sig{\Sigma}}}(k,s)=(\min\{k,n\},s);
$$
\item for every $(\mathbf{s},s)\in S^{\star}\times S$, every $\sigma\in\Sigma_{\mathbf{s},s}$ and every $k\in\omega$,
$$
\Ds^{\gl(\omega,n)}_{\sig{\Sigma}}(\sigma_{k})=\sigma_{\min\{k,n\}}\left(\left(v^{(\min\{k,n\},\mathbf{s})}_{p}\right)_{p\in\bb{\mathbf{s}}}\right);
$$
\item for every pair of dimensions $j<k$ in $\omega$ with $j<n$ and every $s\in S$,
\begin{align*}
\Ds^{\gl(\omega,n)}_{\sig{\Sigma}}(\src^{(j,k)}_{s})&=\src^{(j,\min\{k,n\})}_{s}\left(v^{(\min\{k,n\},s)}_{0}\right),\\
\Ds^{\gl(\omega,n)}_{\sig{\Sigma}}(\tgt^{(j,k)}_{s})&=\tgt^{(j,\min\{k,n\})}_{s}\left(v^{(\min\{k,n\},s)}_{0}\right),\\
\Ds^{\gl(\omega,n)}_{\sig{\Sigma}}(\mathrm{i}^{(k,j)}_{s})&=\mathrm{i}^{(\min\{k,n\},j)}_{s}\left(v^{(j,s)}_{0}\right),\\
\Ds^{\gl(\omega,n)}_{\sig{\Sigma}}(\circ^{(j,k)}_{s})&=v^{(\min\{k,n\},s)}_{0}\circ^{(j,\min\{k,n\})}_{s}v^{(\min\{k,n\},s)}_{1};
\end{align*}
\item for every pair of dimensions $j<k$ in $\omega$ with $j\geq n$ and every $s\in S$,
$$
\Ds^{\gl(\omega,n)}_{\sig{\Sigma}}(\src^{(j,k)}_{s})=
\Ds^{\gl(\omega,n)}_{\sig{\Sigma}}(\tgt^{(j,k)}_{s})=
\Ds^{\gl(\omega,n)}_{\sig{\Sigma}}(\mathrm{i}^{(k,j)}_{s})=
\Ds^{\gl(\omega,n)}_{\sig{\Sigma}}(\circ^{(j,k)}_{s})=v^{(n,s)}_{0};
$$
and
\item for every $k\in\omega$ and every $s\in S$, the domain is $\Dom^{\Ds^{\gl(\omega,n)}_{\sig{\Sigma}}}_{(k,s)}=\varnothing$.
\end{enumerate}
\end{enumerate}
\end{definition}

Let $\sig{\Sigma}=(S,\Sigma)$ be a many-sorted signature and $n\in\mathbb{N}$. Let us note that, for every $p\in\mathbb{N}$ with $n<p$, the restriction of the underlying derivor of $\bDs^{(\omega,n)}_{\sig{\Sigma}}$ to $\Sigma^{(p)}\subseteq\Sigma^{(\omega)}$ is the underlying derivor of $\bDs^{(p,n)}_{\sig{\Sigma}}$, and that the restriction of the underlying derivor of $\bDs^{\gl(\omega,n)}_{\sig{\Sigma}}$ to the image of $\Sigma^{\gl(p)}$ in $\Sigma^{\gl(\omega)}$ is the underlying derivor of $\bDs^{\gl(p,n)}_{\sig{\Sigma}}$, since, in both cases, the defining clauses coincide, see Definitions~\ref{DDiscDerivor} and~\ref{DDiscGlDerivor}. Moreover, all of them have the same domains. Therefore, by Proposition~\ref{PTransformedAlg}, for every object $\mathbf{B}$ in $\mathsf{EPAlg}(\sig{\Sigma}^{(n)},\mathcal{C}^{\sig{\Sigma}^{(n)}})$, the reduct of $(\bDs^{(\omega,n)}_{\sig{\Sigma}})^{\ast}_{\gd}(\mathbf{B})$ to $\Sigma^{(p)}$ is $\Ds^{(p,n)}_{\sig{\Sigma}}(\mathbf{B})$, and, for every object $\mathbf{B}$ in $\mathsf{EPAlg}(\sig{\Sigma}^{\gl(n)},\mathcal{C}^{\sig{\Sigma}^{\gl(n)}})$, the reduct of $(\bDs^{\gl(\omega,n)}_{\sig{\Sigma}})^{\ast}_{\gd}(\mathbf{B})$ to the sorts of dimension at most $p$ and to the operation symbols of $\Sigma^{\gl(p)}$ is $\Ds^{\gl(p,n)}_{\sig{\Sigma}}(\mathbf{B})$. In fact, by Proposition~\ref{PTransformedAlg}, in both algebras the cells are the same and every operation symbol is interpreted by means of the same term and the same domain conditions.

In the following proposition we prove that the discretization derivors in dimension $\omega$ send categorial algebras to categorial algebras.

\begin{proposition}\label{PDiscOmegaDerivor}
Let $\sig{\Sigma}=(S,\Sigma)$ be a many-sorted signature and $n\in\mathbb{N}$. Then the restriction of $(\bDs^{(\omega,n)}_{\sig{\Sigma}})^{\ast}_{\gd}$ to $\mathsf{EPAlg}(\sig{\Sigma}^{(n)},\mathcal{C}^{\sig{\Sigma}^{(n)}})$ corestricts to $\mathsf{EPAlg}(\sig{\Sigma}^{(\omega)},\mathcal{C}^{\sig{\Sigma}^{(\omega)}})$, and the restriction of $(\bDs^{\gl(\omega,n)}_{\sig{\Sigma}})^{\ast}_{\gd}$ to $\mathsf{EPAlg}(\sig{\Sigma}^{\gl(n)},\mathcal{C}^{\sig{\Sigma}^{\gl(n)}})$ corestricts to $\mathsf{EPAlg}(\sig{\Sigma}^{\gl(\omega)},\mathcal{C}^{\sig{\Sigma}^{\gl(\omega)}})$, see Definition~\ref{DCorestriction}. Moreover, for every object $\mathbf{B}$ in $\mathsf{EPAlg}(\sig{\Sigma}^{(n)},\mathcal{C}^{\sig{\Sigma}^{(n)}})$, every cell of $(\bDs^{(\omega,n)}_{\sig{\Sigma}})^{\ast}_{\gd}(\mathbf{B})$ is an $n$-cell, so that $(\bDs^{(\omega,n)}_{\sig{\Sigma}})^{\ast}_{\gd}(\mathbf{B})$ is an object in $\omega\mathsf{CatAlg}(\sig{\Sigma})$.
\end{proposition}
\begin{proof}
By Definitions~\ref{DCorestriction} and~\ref{DEpVariety}, in each case we have to prove that the transformed algebras are ep-algebras in which the axioms of the corresponding ep-variety hold.

\textsf{(Single-set case: ep-algebra)} Let $\mathbf{B}$ be an object in $\mathsf{EPAlg}(\sig{\Sigma}^{(n)},\mathcal{C}^{\sig{\Sigma}^{(n)}})$. Let $\tau$ be an operation symbol of $\Sigma^{(\omega)}$ and $p\in\mathbb{N}$ with $n<p$ such that $\tau$ is an operation symbol of $\Sigma^{(p)}$. By the above note, the interpretation of $\tau$ in $(\bDs^{(\omega,n)}_{\sig{\Sigma}})^{\ast}_{\gd}(\mathbf{B})$ is its interpretation in $\Ds^{(p,n)}_{\sig{\Sigma}}(\mathbf{B})$, which, by Lemma~\ref{LDiscDescription}, is a $\sig{\Sigma}^{(p)}$-algebra. Since the domain conditions of $\tau$ are the same in $\sig{\Sigma}^{(p)}$ and in $\sig{\Sigma}^{(\omega)}$, the domain of the interpretation of $\tau$ in $(\bDs^{(\omega,n)}_{\sig{\Sigma}})^{\ast}_{\gd}(\mathbf{B})$ is the set of the families that satisfy $\mathrm{def}^{(S,\Sigma^{(\omega)})}(\tau)$ in $(\bDs^{(\omega,n)}_{\sig{\Sigma}})^{\ast}_{\gd}(\mathbf{B})$. Thus $(\bDs^{(\omega,n)}_{\sig{\Sigma}})^{\ast}_{\gd}(\mathbf{B})$ is a $\sig{\Sigma}^{(\omega)}$-algebra, see Definition~\ref{DEpAlgebra} and Proposition~\ref{PTransformedEpAlg}.

\textsf{(Single-set case: axioms)} Every axiom of $\boldsymbol{\mathcal{C}}^{\sig{\Sigma}^{(\omega)}}$ only involves finitely many dimensions. Let $p\in\mathbb{N}$ with $n<p$ be greater than all of them. Then the axiom is an axiom of $\boldsymbol{\mathcal{C}}^{\sig{\Sigma}^{(p)}}$ and, by the above note, it holds in $(\bDs^{(\omega,n)}_{\sig{\Sigma}})^{\ast}_{\gd}(\mathbf{B})$ if, and only if, it holds in $\Ds^{(p,n)}_{\sig{\Sigma}}(\mathbf{B})$, which, by Proposition~\ref{PDiscDerivor}, is an object in $\mathsf{EPAlg}(\sig{\Sigma}^{(p)},\mathcal{C}^{\sig{\Sigma}^{(p)}})$. Hence every axiom of $\boldsymbol{\mathcal{C}}^{\sig{\Sigma}^{(\omega)}}$ holds in $(\bDs^{(\omega,n)}_{\sig{\Sigma}})^{\ast}_{\gd}(\mathbf{B})$.

\textsf{(Single-set case: $n$-cells)} By the above note and Lemma~\ref{LDiscDescription}, applied to any $p>n$, $\src^{n}_{s}$ acts as the identity in $(\bDs^{(\omega,n)}_{\sig{\Sigma}})^{\ast}_{\gd}(\mathbf{B})$, so that every cell of $(\bDs^{(\omega,n)}_{\sig{\Sigma}})^{\ast}_{\gd}(\mathbf{B})$ is an $n$-cell. Therefore, by Remark~\ref{RSingleSetIdentOmega}, $(\bDs^{(\omega,n)}_{\sig{\Sigma}})^{\ast}_{\gd}(\mathbf{B})$ is an object in $\omega\mathsf{CatAlg}(\sig{\Sigma})$.

\textsf{(Globular case: ep-algebra)} Let $\mathbf{B}$ be an object in $\mathsf{EPAlg}(\sig{\Sigma}^{\gl(n)},\mathcal{C}^{\sig{\Sigma}^{\gl(n)}})$. Let $\tau$ be an operation symbol of $\Sigma^{\gl(\omega)}$ and $p\in\mathbb{N}$ with $n<p$ such that $\tau$ is in the image of $\Sigma^{\gl(p)}$. By the above note, the interpretation of $\tau$ in $(\bDs^{\gl(\omega,n)}_{\sig{\Sigma}})^{\ast}_{\gd}(\mathbf{B})$ is its interpretation in $\Ds^{\gl(p,n)}_{\sig{\Sigma}}(\mathbf{B})$, which, by Lemma~\ref{LDiscGlDescription}, is a $\sig{\Sigma}^{\gl(p)}$-algebra. Since the domain conditions of $\tau$ are the same in $\sig{\Sigma}^{\gl(p)}$ and in $\sig{\Sigma}^{\gl(\omega)}$, $(\bDs^{\gl(\omega,n)}_{\sig{\Sigma}})^{\ast}_{\gd}(\mathbf{B})$ is a $\sig{\Sigma}^{\gl(\omega)}$-algebra, see Definition~\ref{DEpAlgebra} and Proposition~\ref{PTransformedEpAlg}.

\textsf{(Globular case: axioms)} Every axiom of $\boldsymbol{\mathcal{C}}^{\sig{\Sigma}^{\gl(\omega)}}$ only involves finitely many dimensions. Let $p\in\mathbb{N}$ with $n<p$ be greater than all of them. Then the axiom is, up to the canonical embedding, an axiom of $\boldsymbol{\mathcal{C}}^{\sig{\Sigma}^{\gl(p)}}$ and, by the above note, it holds in $(\bDs^{\gl(\omega,n)}_{\sig{\Sigma}})^{\ast}_{\gd}(\mathbf{B})$ if, and only if, it holds in $\Ds^{\gl(p,n)}_{\sig{\Sigma}}(\mathbf{B})$, which, by Proposition~\ref{PDiscGlDerivor}, is an object in $\mathsf{EPAlg}(\sig{\Sigma}^{\gl(p)},\mathcal{C}^{\sig{\Sigma}^{\gl(p)}})$. Hence every axiom of $\boldsymbol{\mathcal{C}}^{\sig{\Sigma}^{\gl(\omega)}}$ holds in $(\bDs^{\gl(\omega,n)}_{\sig{\Sigma}})^{\ast}_{\gd}(\mathbf{B})$.

This completes the proof.
\end{proof}

We are now in position to define the discretization functors in dimension $\omega$.

\begin{definition}\label{DDiscretizationOmega}
Let $\sig{\Sigma}=(S,\Sigma)$ be a many-sorted signature and $n\in\mathbb{N}$. The \emph{discretization functor}
$$
\Ds^{(\omega,n)}_{\sig{\Sigma}}\colon\mathsf{nCatAlg}(\sig{\Sigma})\mor\omega\mathsf{CatAlg}(\sig{\Sigma})
$$
is the corestriction to $\omega\mathsf{CatAlg}(\sig{\Sigma})$ of the functor $(\bDs^{(\omega,n)}_{\sig{\Sigma}})^{\ast}_{\gd}$ from $\mathsf{EPAlg}(\sig{\Sigma}^{(n)},\mathcal{C}^{\sig{\Sigma}^{(n)}})=\mathsf{nCatAlg}(\sig{\Sigma})$ to $\mathsf{EPAlg}(\sig{\Sigma}^{(\omega)},\mathcal{C}^{\sig{\Sigma}^{(\omega)}})$, see Proposition~\ref{PDiscOmegaDerivor}.
\end{definition}

By the above note, Lemma~\ref{LDiscDescription} and Proposition~\ref{PGenHepDerivorFunctor}, $\Ds^{(\omega,n)}_{\sig{\Sigma}}$ sends an $n$-categorial $\sig{\Sigma}$-algebra $\calg{A}$ to the $\omega$-categorial $\sig{\Sigma}$-algebra whose cells of sort $s$ are the cells of sort $s$ of $\calg{A}$, in which the operations of $\Sigma$ and the $k$-sources, $k$-targets and $k$-compositions, for $k<n$, are those of $\calg{A}$, whereas, for $k\geq n$, $\src^{k}_{s}$ and $\tgt^{k}_{s}$ are identities and $x\circ^{k}_{s}y$ is defined if, and only if, $x=y$, and then it is $x$; and it sends an $n$-categorial $\sig{\Sigma}$-homomorphism to itself.

\begin{definition}\label{DDiscretizationOmegaGl}
Let $\sig{\Sigma}=(S,\Sigma)$ be a many-sorted signature and $n\in\mathbb{N}$. The \emph{globular discretization functor}
$$
\Ds^{\gl(\omega,n)}_{\sig{\Sigma}}\colon\mathsf{nCat}^{\gl}\mathsf{Alg}(\sig{\Sigma})\mor\omega\mathsf{Cat}^{\gl}\mathsf{Alg}(\sig{\Sigma})
$$
is the functor $(\bDs^{\gl(\omega,n)}_{\sig{\Sigma}})^{\ast}_{\gd}$ from $\mathsf{EPAlg}(\sig{\Sigma}^{\gl(n)},\mathcal{C}^{\sig{\Sigma}^{\gl(n)}})=\mathsf{nCat}^{\gl}\mathsf{Alg}(\sig{\Sigma})$ to $\mathsf{EPAlg}(\sig{\Sigma}^{\gl(\omega)},\mathcal{C}^{\sig{\Sigma}^{\gl(\omega)}})=\omega\mathsf{Cat}^{\gl}\mathsf{Alg}(\sig{\Sigma})$, see Proposition~\ref{PDiscOmegaDerivor}.
\end{definition}

By the above note, Lemma~\ref{LDiscGlDescription} and Proposition~\ref{PGenHepDerivorFunctor}, $\Ds^{\gl(\omega,n)}_{\sig{\Sigma}}$ sends a globular $n$-categorial $\sig{\Sigma}$-algebra $\calg{A}$ to the globular $\omega$-categorial $\sig{\Sigma}$-algebra whose $k$-cells of sort $s$, for $k\in\omega$, are the $\min\{k,n\}$-cells of sort $s$ of $\calg{A}$, in which, for every pair of dimensions $j<k$, $\sigma_{k}$ is $\sigma_{\min\{k,n\}}$, and the $(j,k)$-sources, $(j,k)$-targets, $(k,j)$-identities and $j$-compositions of $k$-cells are, for $j<n$, the $(j,\min\{k,n\})$-sources, $(j,\min\{k,n\})$-targets, $(\min\{k,n\},j)$-identities and $j$-compositions of $\min\{k,n\}$-cells of $\calg{A}$, whereas, for $j\geq n$, the $(j,k)$-sources, $(j,k)$-targets and $(k,j)$-identities are identities and the $j$-composition of two $k$-cells is defined if, and only if, they are equal, and then it is the common cell; and it sends a globular $n$-categorial $\sig{\Sigma}$-homomorphism $F$ to the family whose component at $(k,s)$ is $F_{(\min\{k,n\},s)}$.

In the following remark we collect the relations between these functors and the finite-dimensional ones.

\begin{remark}\label{RDiscOmegaTrunc}
Let $\sig{\Sigma}=(S,\Sigma)$ be a many-sorted signature and $m\in\mathbb{N}$. For $n\in\mathbb{N}$, let $F_{n}=\Ds^{(n,m)}_{\sig{\Sigma}}$, if $m\leq n$, and $F_{n}=\Tr^{(n,m)}_{\sig{\Sigma}}$, if $n\leq m$, and let $F^{\gl}_{n}$ be defined in the same way with the globular functors. Then the following statements hold:
\begin{enumerate}
\item for every $n\in\mathbb{N}$ with $m\leq n$, $\Tr^{(m,n)}_{\sig{\Sigma}}\circ\Tr^{(n,\omega)}_{\sig{\Sigma}}=\Tr^{(m,\omega)}_{\sig{\Sigma}}$ and $\Tr^{\gl(m,n)}_{\sig{\Sigma}}\circ\Tr^{\gl(n,\omega)}_{\sig{\Sigma}}=\Tr^{\gl(m,\omega)}_{\sig{\Sigma}}$;
\item for every $n\in\mathbb{N}$, $\Tr^{(n,\omega)}_{\sig{\Sigma}}\circ\Ds^{(\omega,m)}_{\sig{\Sigma}}=F_{n}$ and $\Tr^{\gl(n,\omega)}_{\sig{\Sigma}}\circ\Ds^{\gl(\omega,m)}_{\sig{\Sigma}}=F^{\gl}_{n}$; in particular, $\Tr^{(m,\omega)}_{\sig{\Sigma}}\circ\Ds^{(\omega,m)}_{\sig{\Sigma}}$ and $\Tr^{\gl(m,\omega)}_{\sig{\Sigma}}\circ\Ds^{\gl(\omega,m)}_{\sig{\Sigma}}$ are identities.
\end{enumerate}
In fact, (1) follows from the fact that, by~(AB$^{m,n}$1), the $m$-cells are $n$-cells; and (2) follows from the explicit descriptions of the functors involved: if $m\leq n$, then the $n$-cells of $\Ds^{(\omega,m)}_{\sig{\Sigma}}(\calg{A})$ are all the cells of $\calg{A}$, and, if $n<m$, then they are the $n$-cells of $\calg{A}$.
\end{remark}

\section{$\omega\mathsf{CatAlg}(\sig{\Sigma})$ as a projective limit}\label{SLimitSingleSet}

In this section we prove that the category $\omega\mathsf{CatAlg}(\sig{\Sigma})$ of $S$-sorted $\omega$-categorial $\Sigma$-algebras and $\omega$-categorial $\Sigma$-homomorphisms is the projective limit of the projective system $(\mathsf{nCatAlg}(\sig{\Sigma}),\Tr^{(\bigcdot)}_{\sig{\Sigma}})$, where the functors $\Tr^{(m,n)}_{\sig{\Sigma}}$ are the truncation functors defined in Definition~\ref{DTruncation}, and the projections of the limit are the truncation functors $\Tr^{(n,\omega)}_{\sig{\Sigma}}$ from $\omega\mathsf{CatAlg}(\sig{\Sigma})$ to $\mathsf{nCatAlg}(\sig{\Sigma})$ of Definition~\ref{DTruncationOmega}. We begin by introducing the projective system.

\begin{remark}\label{RTruncProjSys}
Let $\sig{\Sigma}=(S,\Sigma)$ be a many-sorted signature. We let $(\bigcdot)\mathsf{CatAlg}(\sig{\Sigma})$ stand for the family of categories $(\mathsf{nCatAlg}(\sig{\Sigma}))_{n\in\omega}$ and we let $\Tr^{(\bigcdot)}_{\sig{\Sigma}}$ stand for the family of functors $(\Tr^{(m,n)}_{\sig{\Sigma}})_{(m,n)\in\leq}$. Let us note that $\Tr^{(n,n)}_{\sig{\Sigma}}=\Id_{\mathsf{nCatAlg}(\sig{\Sigma})}$ and that, for every $m\leq n\leq p$, by Proposition~\ref{PTruncComp},
\[
\Tr^{(m,n)}_{\sig{\Sigma}}\circ \Tr^{(n,p)}_{\sig{\Sigma}}=\Tr^{(m,p)}_{\sig{\Sigma}}.
\]
Therefore, we obtain the projective system $((\bigcdot)\mathsf{CatAlg}(\sig{\Sigma}),\Tr^{(\bigcdot)}_{\sig{\Sigma}})$ in $\mathsf{Cat}$ and its projective limit $\varprojlim ((\bigcdot)\mathsf{CatAlg}(\sig{\Sigma}),\Tr^{(\bigcdot)}_{\sig{\Sigma}})$.
\end{remark}

Our next goal is to prove that such limit is isomorphic to $(\omega\mathsf{CatAlg}(\sig{\Sigma}),\Tr^{(\bigcdot,\omega)}_{\sig{\Sigma}})$ with $\Tr^{(\bigcdot,\omega)}_{\sig{\Sigma}}$ being the family of functors $(\Tr^{(n,\omega)}_{\sig{\Sigma}})_{n\in\omega}$. That the composition of functors $\Tr^{(m,n)}_{\sig{\Sigma}}\circ \Tr^{(n,\omega)}_{\sig{\Sigma}}$ is equal to the functor $\Tr^{(m,\omega)}_{\sig{\Sigma}}$ has been noted in Remark~\ref{RDiscOmegaTrunc}. Thus, all that remains to be proven is that $(\omega\mathsf{CatAlg}(\sig{\Sigma}),\Tr^{(\bigcdot,\omega)}_{\sig{\Sigma}})$ satisfies the universal property.

\begin{proposition}\label{PUnivPropCatAlg}
Let $\sig{\Sigma}=(S,\Sigma)$ be a many-sorted signature, $\mathsf{L}$ a category and, for every $n\in\omega$, let $F_{n}$ be a functor from $\mathsf{L}$ to $\mathsf{nCatAlg}(\sig{\Sigma})$. If, for every $m,n\in\omega$ with $m\leq n$, $F_{m} = \Tr^{(m,n)}_{\sig{\Sigma}} \circ F_{n}$, then there exists a unique functor $\langle F_{k}\rangle_{k\in\omega}$ from $\mathsf{L}$ to the category $\omega\mathsf{CatAlg}(\sig{\Sigma})$ such that, for every $n\in\omega$, $F_{n} = \Tr^{(n,\omega)}_{\sig{\Sigma}} \circ \langle F_{k}\rangle_{k\in\omega}$. The diagram in Figure~\ref{FUnivPropCatAlg} depicts this situation.
\begin{figure}
\[
\xymatrix{
& \mathsf{L}\ar@{-->}[d]^{\exists!\langle F_{k}\rangle_{k\in\omega}}\ar@/_30pt/[ddl]_{F_{m}}\ar@/^30pt/[ddr]^{F_{n}} &
\\
& \omega\mathsf{CatAlg}(\sig{\Sigma})\ar[dr]^{\Tr^{(n,\omega)}_{\sig{\Sigma}}}\ar[dl]_{\Tr^{(m,\omega)}_{\sig{\Sigma}}} &
\\
\mathsf{mCatAlg}(\sig{\Sigma}) && \mathsf{nCatAlg}(\sig{\Sigma})\ar[ll]^{\Tr^{(m,n)}_{\sig{\Sigma}}}
}
\]
\caption{Universal Property of $\omega\mathsf{CatAlg}(\sig{\Sigma})$.}
\label{FUnivPropCatAlg}
\end{figure}
\end{proposition}

\begin{proof}
Before proceeding any further let us fix some terminology. We will say that a family $(\calg{A}_{n})_{n\in\omega}$, with $\calg{A}_{n}$ an $n$-categorial $\sig{\Sigma}$-algebra, is \emph{compatible} if $\Tr^{(m,n)}_{\sig{\Sigma}}(\calg{A}_{n})=\calg{A}_{m}$ for all $m\leq n$, i.e., if, for $m\leq n$, the cells of $\calg{A}_{m}$ are the $m$-cells of $\calg{A}_{n}$ and its operations are the restrictions of those of $\calg{A}_{n}$. Similarly, a family of homomorphisms $(H_{n}\colon\calg{A}_{n}\mor\calg{A}'_{n})_{n\in\omega}$ between compatible families is \emph{compatible} if $\Tr^{(m,n)}_{\sig{\Sigma}}(H_{n})=H_{m}$ for $m\leq n$, i.e., if $H_{m}$ is the restriction of $H_{n}$.

In order to define the functor $\langle F_{k}\rangle_{k\in\omega}$, we will prove the following two claims.

\begin{claim}\label{CUnivPropCatAlgA}
For every compatible family $(\calg{A}_{n})_{n\in\omega}$ there exists a unique $\omega$-categorial $\sig{\Sigma}$-algebra $\varprojlim_{n}\calg{A}_{n}$ such that, for every $n\in\omega$, $\Tr^{(n,\omega)}_{\sig{\Sigma}}(\varprojlim_{m}\calg{A}_{m})=\calg{A}_{n}$.
\end{claim}

Let $(\calg{A}_{n})_{n\in\omega}$ be a compatible family. Then the cells of every sort form an increasing chain of sets. We let $\varprojlim_{n}\calg{A}_{n}$ be the partial $\Sigma^{(\omega)}$-algebra whose cells of sort $s$ are those of the union of the $\calg{A}_{n}$ and in which, for finitely many cells and a dimension $k$, the operations of $\Sigma$ and the $k$-sources, $k$-targets and $k$-compositions are computed in any $\calg{A}_{n}$ with $n>k$ that contains the cells.

\textsf{(Well-definedness)} This does not depend on $n$: if $n\leq n'$, the operations of $\calg{A}_{n}$ of dimension $k<n$ are the restrictions of those of $\calg{A}_{n'}$, with the restricted domains.

\textsf{(Axioms)} Every axiom of $\boldsymbol{\mathcal{C}}^{\sig{\Sigma}^{(\omega)}}$ involves finitely many cells and dimensions. Since the operations of $\calg{A}_{n}$ of dimension $k<n$ take their values in $\calg{A}_{n}$, it holds in $\varprojlim_{n}\calg{A}_{n}$ because it holds in $\calg{A}_{n}$ for $n$ large enough.

\textsf{(Finite dimension)} Every cell $x$ of $\calg{A}_{n}$ is an $n$-cell of the limit, because $\calg{A}_{n}=\Tr^{(n,n+1)}_{\sig{\Sigma}}(\calg{A}_{n+1})$, so that $\src^{n}_{s}(x)=x$ in $\calg{A}_{n+1}$. By Remark~\ref{RSingleSetIdentOmega}, $\varprojlim_{n}\calg{A}_{n}$ is an $\omega$-categorial $\sig{\Sigma}$-algebra.

\textsf{(Truncations)} The $n$-cells of the limit are the cells of $\calg{A}_{n}$: an $n$-cell of $\calg{A}_{p}$, with $p\geq n$, is a cell of $\Tr^{(n,p)}_{\sig{\Sigma}}(\calg{A}_{p})=\calg{A}_{n}$. Hence $\Tr^{(n,\omega)}_{\sig{\Sigma}}(\varprojlim_{m}\calg{A}_{m})=\calg{A}_{n}$.

\textsf{(Uniqueness)} Every $\omega$-categorial $\sig{\Sigma}$-algebra $\calg{B}$ is the limit of its truncations, since every cell of $\calg{B}$ is an $n$-cell for some $n$ and the operations of $\Tr^{(n,\omega)}_{\sig{\Sigma}}(\calg{B})$ are restrictions of those of $\calg{B}$. Hence, if $\Tr^{(n,\omega)}_{\sig{\Sigma}}(\calg{B})=\calg{A}_{n}$ for every $n$, then $\calg{B}=\varprojlim_{n}\calg{A}_{n}$.

Claim~\ref{CUnivPropCatAlgA} follows.

\begin{claim}\label{CUnivPropCatAlgB}
For every compatible family of homomorphisms $(H_{n}\colon\calg{A}_{n}\mor\calg{A}'_{n})_{n\in\omega}$ there exists a unique $\omega$-categorial $\sig{\Sigma}$-homomorphism $\varprojlim_{n}H_{n}$ from $\varprojlim_{n}\calg{A}_{n}$ to $\varprojlim_{n}\calg{A}'_{n}$ such that, for every $n\in\omega$, $\Tr^{(n,\omega)}_{\sig{\Sigma}}(\varprojlim_{m}H_{m})=H_{n}$.
\end{claim}

We let $\varprojlim_{n}H_{n}$ be the mapping that acts on the cells of $\calg{A}_{n}$ as $H_{n}$, which is well defined by compatibility. It is a homomorphism, because every operation of the limit is computed at a finite level, where $H_{n}$ is a homomorphism. By construction, its truncations are the $H_{n}$, and it is the unique mapping with this property, since every cell of the limit belongs to some $\calg{A}_{n}$.

Claim~\ref{CUnivPropCatAlgB} follows.

For an object $X$ of $\mathsf{L}$, the family $(F_{n}(X))_{n\in\omega}$ is compatible, and, for a morphism $f$ of $\mathsf{L}$, so is $(F_{n}(f))_{n\in\omega}$. We let
$$
\langle F_{k}\rangle_{k\in\omega}(X)=\varprojlim_{n}F_{n}(X)
\qquad\text{and}\qquad
\langle F_{k}\rangle_{k\in\omega}(f)=\varprojlim_{n}F_{n}(f).
$$

\textsf{(Functoriality)} By the uniqueness in Claim~\ref{CUnivPropCatAlgB}, $\langle F_{k}\rangle_{k\in\omega}$ preserves identities and composites, since every $F_{n}$ does.

\textsf{(Factorization)} By Claims~\ref{CUnivPropCatAlgA} and~\ref{CUnivPropCatAlgB}, $\Tr^{(n,\omega)}_{\sig{\Sigma}}\circ\langle F_{k}\rangle_{k\in\omega}=F_{n}$, for every $n\in\omega$.

\textsf{(Uniqueness)} If $G\colon\mathsf{L}\mor\omega\mathsf{CatAlg}(\sig{\Sigma})$ also satisfies $\Tr^{(n,\omega)}_{\sig{\Sigma}}\circ G=F_{n}$ for every $n$, then, by the uniqueness in Claims~\ref{CUnivPropCatAlgA} and~\ref{CUnivPropCatAlgB}, $G(X)=\varprojlim_{n}F_{n}(X)$ and $G(f)=\varprojlim_{n}F_{n}(f)$, so that $G=\langle F_{k}\rangle_{k\in\omega}$.

This completes the proof.
\end{proof}

We will use the universal property in the following form.

\begin{corollary}\label{CLimitNat}
Let $\sig{\Sigma}=(S,\Sigma)$ be a many-sorted signature. Then the following statements hold:
\begin{enumerate}
\item For $\omega$-categorial $\sig{\Sigma}$-algebras $\calg{A}$ and $\calg{B}$, the mapping that sends $H$ to $(\Tr^{(n,\omega)}_{\sig{\Sigma}}(H))_{n\in\omega}$ is a bijection from the set of homomorphisms from $\calg{A}$ to $\calg{B}$ to the set of families $(H_{n}\colon \Tr^{(n,\omega)}_{\sig{\Sigma}}(\calg{A})\mor \Tr^{(n,\omega)}_{\sig{\Sigma}}(\calg{B}))_{n\in\omega}$ such that $\Tr^{(m,n)}_{\sig{\Sigma}}(H_{n})=H_{m}$ for $m\leq n$.
\item Let $F,G\colon\mathsf{L}\mor\omega\mathsf{CatAlg}(\sig{\Sigma})$ be functors and, for every $n\in\omega$, $\varphi^{n}$ a natural transformation from $\Tr^{(n,\omega)}_{\sig{\Sigma}}\circ F$ to $\Tr^{(n,\omega)}_{\sig{\Sigma}}\circ G$ such that $\Tr^{(m,n)}_{\sig{\Sigma}}\ast\varphi^{n}=\varphi^{m}$ for $m\leq n$. Then there exists a unique natural transformation $\varphi$ from $F$ to $G$ such that $\Tr^{(n,\omega)}_{\sig{\Sigma}}\ast\varphi=\varphi^{n}$ for every $n$, and $\varphi$ is a natural isomorphism if every $\varphi^{n}$ is.
\end{enumerate}
\end{corollary}
\begin{proof}
(1) It follows from Claim~\ref{CUnivPropCatAlgB} in the proof of Proposition~\ref{PUnivPropCatAlg}.

(2) By~(1), for every object $X$ of $\mathsf{L}$ there exists a unique homomorphism $\varphi_{X}$ with $\Tr^{(n,\omega)}_{\sig{\Sigma}}(\varphi_{X})=\varphi^{n}_{X}$ for every $n$.

\textsf{(Naturality)} For a morphism $f\colon X\mor Y$, the homomorphisms $G(f)\circ\varphi_{X}$ and $\varphi_{Y}\circ F(f)$ have the same truncations, by the naturality of every $\varphi^{n}$, so that, by~(1), they are equal.

\textsf{(Isomorphism)} If every $\varphi^{n}_{X}$ is an isomorphism, the inverses form a family as in~(1), whose image under the bijection of~(1) is the inverse of $\varphi_{X}$.

This completes the proof.
\end{proof}


\section{$\omega\mathsf{Cat}^{\gl}\mathsf{Alg}(\sig{\Sigma})$ as a projective limit}\label{SLimitGlobular}

In this section we prove that the category $\omega\mathsf{Cat}^{\gl}\mathsf{Alg}(\sig{\Sigma})$ of $S$-sorted globular $\omega$-categorial $\Sigma$-algebras and globular $\omega$-categorial $\Sigma$-homomorphisms is the projective limit of the projective system $(\mathsf{nCat}^{\gl}\mathsf{Alg}(\sig{\Sigma}),\Tr^{\gl(\bigcdot)}_{\sig{\Sigma}})$, where the functors $\Tr^{\gl(m,n)}_{\sig{\Sigma}}$ are the truncation functors defined in Definition~\ref{DTruncationGl}, and the projections of the limit are the truncation functors $\Tr^{\gl(n,\omega)}_{\sig{\Sigma}}$ of Definition~\ref{DTruncationOmegaGl}. We follow the same pattern as in Section~\ref{SLimitSingleSet}.

\begin{remark}\label{RTruncGlProjSys}
	Let $\sig{\Sigma}=(S,\Sigma)$ be a many-sorted signature. We let $(\bigcdot)\mathsf{Cat}^{\gl}\mathsf{Alg}(\sig{\Sigma})$ stand for the family of categories $(\mathsf{nCat}^{\gl}\mathsf{Alg}(\sig{\Sigma}))_{n\in\omega}$ and we let $\Tr^{\gl(\bigcdot)}_{\sig{\Sigma}}$ stand for the family of functors $(\Tr^{\gl(m,n)}_{\sig{\Sigma}})_{(m,n)\in\leq}$. Let us note that $\Tr^{\gl(n,n)}_{\sig{\Sigma}}=\Id_{\mathsf{nCat}^{\gl}\mathsf{Alg}(\sig{\Sigma})}$ and that, for every $m\leq n\leq p$, by Proposition~\ref{PTruncComp},
	\[
	\Tr^{\gl(m,n)}_{\sig{\Sigma}}\circ \Tr^{\gl(n,p)}_{\sig{\Sigma}}=\Tr^{\gl(m,p)}_{\sig{\Sigma}}.
	\]
	Therefore, we obtain the projective system $((\bigcdot)\mathsf{Cat}^{\gl}\mathsf{Alg}(\sig{\Sigma}),\Tr^{\gl(\bigcdot)}_{\sig{\Sigma}})$ in $\mathsf{Cat}$ and its projective limit $\varprojlim ((\bigcdot)\mathsf{Cat}^{\gl}\mathsf{Alg}(\sig{\Sigma}),\Tr^{\gl(\bigcdot)}_{\sig{\Sigma}})$.
\end{remark}

Our next goal is to prove that such limit is isomorphic to $(\omega\mathsf{Cat}^{\gl}\mathsf{Alg}(\sig{\Sigma}),\Tr^{\gl(\bigcdot,\omega)}_{\sig{\Sigma}})$ with $\Tr^{\gl(\bigcdot,\omega)}_{\sig{\Sigma}}$ being the family of functors $(\Tr^{\gl(n,\omega)}_{\sig{\Sigma}})_{n\in\omega}$. That the composition of functors $\Tr^{\gl(m,n)}_{\sig{\Sigma}}\circ \Tr^{\gl(n,\omega)}_{\sig{\Sigma}}$ is equal to the functor $\Tr^{\gl(m,\omega)}_{\sig{\Sigma}}$ has been noted in Remark~\ref{RDiscOmegaTrunc}. Thus, all that remains to be proven is that $(\omega\mathsf{Cat}^{\gl}\mathsf{Alg}(\sig{\Sigma}),\Tr^{\gl(\bigcdot,\omega)}_{\sig{\Sigma}})$ satisfies the universal property.

\begin{proposition}\label{PUnivPropGlCatAlg}
	Let $\sig{\Sigma}=(S,\Sigma)$ be a many-sorted signature, $\mathsf{L}$ a category and, for every $n\in\omega$, let $F_{n}$ be a functor from $\mathsf{L}$ to $\mathsf{nCat}^{\gl}\mathsf{Alg}(\sig{\Sigma})$. If, for every $m,n\in\omega$ with $m\leq n$, $F_{m} = \Tr^{\gl(m,n)}_{\sig{\Sigma}} \circ F_{n}$, then there exists a unique functor $\langle F_{k}\rangle_{k\in\omega}$ from $\mathsf{L}$ to the category $\omega\mathsf{Cat}^{\gl}\mathsf{Alg}(\sig{\Sigma})$ such that, for every $n\in\omega$, $F_{n} = \Tr^{\gl(n,\omega)}_{\sig{\Sigma}} \circ \langle F_{k}\rangle_{k\in\omega}$. The diagram in Figure~\ref{FUnivPropGlCatAlg} depicts this situation.
	\begin{figure}
		\[
		\xymatrix{
			& \mathsf{L}\ar@{-->}[d]^{\exists!\langle F_{k}\rangle_{k\in\omega}}\ar@/_30pt/[ddl]_{F_{m}}\ar@/^30pt/[ddr]^{F_{n}} &
			\\
			& \omega\mathsf{Cat}^{\gl}\mathsf{Alg}(\sig{\Sigma})\ar[dr]^{\Tr^{\gl(n,\omega)}_{\sig{\Sigma}}}\ar[dl]_{\Tr^{\gl(m,\omega)}_{\sig{\Sigma}}} &
			\\
			\mathsf{mCat}^{\gl}\mathsf{Alg}(\sig{\Sigma}) && \mathsf{nCat}^{\gl}\mathsf{Alg}(\sig{\Sigma})\ar[ll]^{\Tr^{\gl(m,n)}_{\sig{\Sigma}}}
		}
		\]
		\caption{Universal Property of $\omega\mathsf{Cat}^{\gl}\mathsf{Alg}(\sig{\Sigma})$.}
		\label{FUnivPropGlCatAlg}
	\end{figure}
\end{proposition}

\begin{proof}
As in the proof of Proposition~\ref{PUnivPropCatAlg}, we will say that a family $(\calg{A}_{n})_{n\in\omega}$, with $\calg{A}_{n}$ a globular $n$-categorial $\sig{\Sigma}$-algebra, is \emph{compatible} if $\Tr^{\gl(m,n)}_{\sig{\Sigma}}(\calg{A}_{n})=\calg{A}_{m}$ for all $m\leq n$, and similarly for families of homomorphisms. We prove the following two claims.

\begin{claim}\label{CUnivPropGlCatAlgA}
For every compatible family $(\calg{A}_{n})_{n\in\omega}$ there exists a unique globular $\omega$-categorial $\sig{\Sigma}$-algebra $\varprojlim_{n}\calg{A}_{n}$ such that, for every $n\in\omega$, $\Tr^{\gl(n,\omega)}_{\sig{\Sigma}}(\varprojlim_{m}\calg{A}_{m})=\calg{A}_{n}$.
\end{claim}

We let $\varprojlim_{n}\calg{A}_{n}$ be the globular $\omega$-categorial $\sig{\Sigma}$-algebra whose $k$-cells of sort $s$ are the $k$-cells of sort $s$ of $\calg{A}_{k}$, its cells of top dimension. For $j<k$, the $(j,k)$-source of a $k$-cell $x$ is its $(j,k)$-source in $\calg{A}_{k}$, which is a $j$-cell of $\calg{A}_{k}$ and hence, by compatibility, a $j$-cell of $\calg{A}_{j}$; the targets, identities and compositions are defined in the same way, and every $\sigma\in\Sigma_{\mathbf{s},s}$ acts on $k$-cells as $\sigma^{\calg{A}_{k}}$ does. Each of the conditions of Definition~\ref{DOmegaCatGl} involves finitely many dimensions $i<j<k$, and holds because it holds in $\calg{A}_{k}$. By construction, $\Tr^{\gl(n,\omega)}_{\sig{\Sigma}}(\varprojlim_{m}\calg{A}_{m})=\calg{A}_{n}$, and every globular $\omega$-categorial $\sig{\Sigma}$-algebra $\calg{B}$ is the limit of its truncations, i.e., $\calg{B}=\varprojlim_{n}\Tr^{\gl(n,\omega)}_{\sig{\Sigma}}(\calg{B})$, which gives the uniqueness.

Claim~\ref{CUnivPropGlCatAlgA} follows.

\begin{claim}\label{CUnivPropGlCatAlgB}
For every compatible family of homomorphisms $(H_{n}\colon\calg{A}_{n}\mor\calg{A}'_{n})_{n\in\omega}$ there exists a unique globular $\omega$-categorial $\sig{\Sigma}$-homomorphism $\varprojlim_{n}H_{n}$ such that, for every $n\in\omega$, $\Tr^{\gl(n,\omega)}_{\sig{\Sigma}}(\varprojlim_{m}H_{m})=H_{n}$.
\end{claim}

We let $\varprojlim_{n}H_{n}$ be the family that acts on $k$-cells as $H_{k}$. It is a homomorphism, because every condition involves finitely many dimensions, and it is clearly the unique one with the stated truncations.

Claim~\ref{CUnivPropGlCatAlgB} follows.

For an object $X$ of $\mathsf{L}$, the family $(F_{n}(X))_{n\in\omega}$ is compatible, and, for a morphism $f$ of $\mathsf{L}$, so is $(F_{n}(f))_{n\in\omega}$. We let
$$
\langle F_{k}\rangle_{k\in\omega}(X)=\varprojlim_{n}F_{n}(X)
\qquad\text{and}\qquad
\langle F_{k}\rangle_{k\in\omega}(f)=\varprojlim_{n}F_{n}(f).
$$

\textsf{(Functoriality)} By the uniqueness in Claim~\ref{CUnivPropGlCatAlgB}, $\langle F_{k}\rangle_{k\in\omega}$ preserves identities and composites, since every $F_{n}$ does.

\textsf{(Factorization)} By Claims~\ref{CUnivPropGlCatAlgA} and~\ref{CUnivPropGlCatAlgB}, $\Tr^{\gl(n,\omega)}_{\sig{\Sigma}}\circ\langle F_{k}\rangle_{k\in\omega}=F_{n}$, for every $n\in\omega$.

\textsf{(Uniqueness)} If $G\colon\mathsf{L}\mor\omega\mathsf{Cat}^{\gl}\mathsf{Alg}(\sig{\Sigma})$ also satisfies $\Tr^{\gl(n,\omega)}_{\sig{\Sigma}}\circ G=F_{n}$ for every $n$, then, by the uniqueness in Claims~\ref{CUnivPropGlCatAlgA} and~\ref{CUnivPropGlCatAlgB}, $G(X)=\varprojlim_{n}F_{n}(X)$ and $G(f)=\varprojlim_{n}F_{n}(f)$, so that $G=\langle F_{k}\rangle_{k\in\omega}$.

This completes the proof.
\end{proof}

\begin{corollary}\label{CLimitNatGl}
Let $\sig{\Sigma}=(S,\Sigma)$ be a many-sorted signature. Then the statements of Corollary~\ref{CLimitNat} hold for $\omega\mathsf{Cat}^{\gl}\mathsf{Alg}(\sig{\Sigma})$ and the truncation functors $\Tr^{\gl(n,\omega)}_{\sig{\Sigma}}$ and $\Tr^{\gl(m,n)}_{\sig{\Sigma}}$.
\end{corollary}
\begin{proof}
(1) It follows from Claim~\ref{CUnivPropGlCatAlgB} in the proof of Proposition~\ref{PUnivPropGlCatAlg}.

(2) Let $F,G\colon\mathsf{L}\mor\omega\mathsf{Cat}^{\gl}\mathsf{Alg}(\sig{\Sigma})$ be functors and, for every $n\in\omega$, $\varphi^{n}$ a natural transformation from $\Tr^{\gl(n,\omega)}_{\sig{\Sigma}}\circ F$ to $\Tr^{\gl(n,\omega)}_{\sig{\Sigma}}\circ G$ such that $\Tr^{\gl(m,n)}_{\sig{\Sigma}}\ast\varphi^{n}=\varphi^{m}$ for $m\leq n$. By~(1), for every object $X$ of $\mathsf{L}$ there exists a unique homomorphism $\varphi_{X}$ with $\Tr^{\gl(n,\omega)}_{\sig{\Sigma}}(\varphi_{X})=\varphi^{n}_{X}$ for every $n$.

\textsf{(Naturality)} For a morphism $f\colon X\mor Y$, the homomorphisms $G(f)\circ\varphi_{X}$ and $\varphi_{Y}\circ F(f)$ have the same truncations, by the naturality of every $\varphi^{n}$, so that, by~(1), they are equal.

\textsf{(Isomorphism)} If every $\varphi^{n}_{X}$ is an isomorphism, the inverses form a family as in~(1), whose image under the bijection of~(1) is the inverse of $\varphi_{X}$.

This completes the proof.
\end{proof}

\section{Graduation, unification, discretization and truncation in dimension \texorpdfstring{$\omega$}{omega}}\label{SOmegaFunctors}

In this section we obtain the $\omega$-dimensional versions of the functors of Sections~\ref{SGraduation} and~\ref{STruncation}. The truncation and discretization functors in dimension $\omega$ have been defined in Subsection~\ref{SDTOmega}, and, by Sections~\ref{SLimitSingleSet} and~\ref{SLimitGlobular}, the truncation functors are the projections of projective limits. We begin by obtaining the graduation functor, by the universal property of Proposition~\ref{PUnivPropGlCatAlg}, from a compatible family of the finite-dimensional ones. Next we obtain the unification functor as an inductive limit. We end this section by proving that graduation and unification are mutually quasi-inverse equivalences and that discretization is left adjoint to truncation also in dimension $\omega$.

\subsection{Graduation and unification}\label{SGdUfOmega}

In the following lemma we prove that the finite-dimensional graduation functors give rise to a compatible family.

\begin{lemma}\label{LGdOmegaCompatible}
Let $\sig{\Sigma}=(S,\Sigma)$ be a many-sorted signature. Then, for every $m,n\in\mathbb{N}$ with $m\leq n$, the following equality holds
$$
\Tr^{\gl(m,n)}_{\sig{\Sigma}}\circ\Gd^{(n)}_{\sig{\Sigma}}\circ \Tr^{(n,\omega)}_{\sig{\Sigma}}=\Gd^{(m)}_{\sig{\Sigma}}\circ \Tr^{(m,\omega)}_{\sig{\Sigma}}.
$$
\end{lemma}
\begin{proof}
The following chain of equalities holds
\allowdisplaybreaks
\begin{align*}
\Tr^{\gl(m,n)}_{\sig{\Sigma}}\circ\Gd^{(n)}_{\sig{\Sigma}}\circ \Tr^{(n,\omega)}_{\sig{\Sigma}}
&=
\Gd^{(m)}_{\sig{\Sigma}}\circ \Tr^{(m,n)}_{\sig{\Sigma}}\circ \Tr^{(n,\omega)}_{\sig{\Sigma}}
\tag{1}
\\&=
\Gd^{(m)}_{\sig{\Sigma}}\circ \Tr^{(m,\omega)}_{\sig{\Sigma}}.
\tag{2}
\end{align*}

The first equality follows from Proposition~\ref{PCompatGdTr}; finally, the last equality follows from Remark~\ref{RDiscOmegaTrunc}.
\end{proof}

We are now in position to define the graduation functor in dimension $\omega$.

\begin{definition}\label{DGdOmega}
Let $\sig{\Sigma}=(S,\Sigma)$ be a many-sorted signature. The \emph{graduation functor} $\Gd^{(\omega)}_{\sig{\Sigma}}\colon\omega\mathsf{CatAlg}(\sig{\Sigma})\mor\omega\mathsf{Cat}^{\gl}\mathsf{Alg}(\sig{\Sigma})$ is the unique functor such that, for every $n\in\mathbb{N}$,
$$
\Tr^{\gl(n,\omega)}_{\sig{\Sigma}}\circ\Gd^{(\omega)}_{\sig{\Sigma}}=\Gd^{(n)}_{\sig{\Sigma}}\circ \Tr^{(n,\omega)}_{\sig{\Sigma}}.
$$
\end{definition}

Let us note that the functor exists and is unique by Proposition~\ref{PUnivPropGlCatAlg}, since, by Lemma~\ref{LGdOmegaCompatible}, the family $(\Gd^{(n)}_{\sig{\Sigma}}\circ \Tr^{(n,\omega)}_{\sig{\Sigma}})_{n\in\omega}$ is compatible. Explicitly, by Claim~\ref{CUnivPropGlCatAlgA} in the proof of Proposition~\ref{PUnivPropGlCatAlg}, the cells of sort $(k,s)$ of $\Gd^{(\omega)}_{\sig{\Sigma}}(\calg{A})$ are the $k$-cells of sort $s$ of $\calg{A}$, the $(j,k)$-sources, $(j,k)$-targets and $j$-compositions of $k$-cells are the $j$-sources, $j$-targets and $j$-compositions of $\calg{A}$, the $(k,j)$-identities are inclusions, and $\sigma_{k}$ acts on $k$-cells as $\sigma$.

We next consider unification. In contrast with graduation, it cannot be obtained in the same way.

\begin{remark}\label{RUfOmegaNotDerivor}
Let $\sig{\Sigma}=(S,\Sigma)$ be a many-sorted signature. The unification functor cannot be obtained as the functor induced by a generalized hep-derivor. In fact, for every generalized hep-derivor $\mathbf{d}=(\varphi,d,\Dom^{(\varphi,d)})$ from $\sig{\Sigma}^{(\omega)}$ to $\sig{\Sigma}^{\gl(\omega)}$, the cells of sort $s$ of $\mathbf{d}^{\ast}_{\gd}(\calg{A})$ are cells of a single sort $\varphi(s)=(k,s')$ of $\calg{A}$. For $S=\{s\}$, $\Sigma=\varnothing$ and $\varphi(s)=(k,s)$, let $\calg{A}$ be the globular $\omega$-category with exactly one cell in every dimension up to $k$, whose $(k+1)$-cells form the two-element monoid $\{e,f\}$, with $f\circ f=f$, for all the compositions, and whose cells of dimension greater than $k+1$ are the identities of the $(k+1)$-cells. Then $\mathbf{d}^{\ast}_{\gd}(\calg{A})$ has at most one cell, whereas its unification must have two distinct cells, which correspond to $e$ and $f$.

Nor can we apply directly the universal property of Proposition~\ref{PUnivPropCatAlg} to the family $(\Uf^{(n)}_{\sig{\Sigma}}\circ \Tr^{\gl(n,\omega)}_{\sig{\Sigma}})_{n\in\omega}$, since, by Proposition~\ref{PCompatUfTr}, it is only compatible up to the natural isomorphisms $\beta^{(m,n)}$: the $m$-cells of $\Uf^{(n)}_{\sig{\Sigma}}(\calg{B})$ are not the $m$-cells of $\calg{B}$, but their $(n,m)$-identities. These isomorphisms, however, are the inclusions of an inductive system, and the unification of a globular $\omega$-categorial $\sig{\Sigma}$-algebra is its inductive limit, as for $\omega$-categories.
\end{remark}

We begin by defining the $\omega$-cells of a globular $\omega$-categorial $\sig{\Sigma}$-algebra.

\begin{definition}\label{DUfOmegaCells}
Let $\sig{\Sigma}=(S,\Sigma)$ be a many-sorted signature, $\calg{B}$ a globular $\omega$-categorial $\sig{\Sigma}$-algebra and $s\in S$. The set of \emph{$\omega$-cells} of sort $s$ of $\calg{B}$ is $\coprod_{k\in\omega}B_{(k,s)}$, whose elements are the pairs $(x,k)$ with $k\in\omega$ and $x\in B_{(k,s)}$. We let $\Theta_{s}$ be the binary relation on $\coprod_{k\in\omega}B_{(k,s)}$ defined by $((x,j),(y,k))\in\Theta_{s}$ if
\begin{enumerate}
\item $j<k$ and $y=\mathrm{i}^{(k,j)}_{s}(x)$; or
\item $j=k$ and $x=y$; or
\item $j>k$ and $x=\mathrm{i}^{(j,k)}_{s}(y)$.
\end{enumerate}
By~(EA$^{h,j,k}$2), $\Theta_{s}$ is an equivalence relation. We let $[B]_{s}$ stand for the quotient set $\coprod_{k\in\omega}B_{(k,s)}/\Theta_{s}$, $[x]_{k}$ for the class of $(x,k)$, and $\mathrm{in}^{[k]}_{s}$ for the mapping from $B_{(k,s)}$ to $[B]_{s}$ that sends $x$ to $[x]_{k}$.
\end{definition}

In the following lemma we collect the properties of the $\omega$-cells that we need.

\begin{lemma}\label{LUfOmegaCells}
Let $\sig{\Sigma}=(S,\Sigma)$ be a many-sorted signature, $\calg{B}$ a globular $\omega$-categorial $\sig{\Sigma}$-algebra and $s\in S$. Then the following statements hold:
\begin{enumerate}
\item $[B]_{s}$, together with the mappings $\mathrm{in}^{[k]}_{s}$, is the inductive limit of the inductive system of sets $(B_{(k,s)},\mathrm{i}^{(k+1,k)}_{s})_{k\in\omega}$, i.e., for every set $L$ and every family of mappings $f_{k}\colon B_{(k,s)}\mor L$ such that $f_{j}=f_{k}\circ\mathrm{i}^{(k,j)}_{s}$ for every $j<k$, there exists a unique mapping $[f_{k}]_{k\in\omega}\colon[B]_{s}\mor L$ such that $[f_{k}]_{k\in\omega}\circ\mathrm{in}^{[j]}_{s}=f_{j}$ for every $j\in\omega$;
\item every $\mathrm{in}^{[k]}_{s}$ is injective; and
\item $[B]_{s}$ is the increasing union of the images of the $\mathrm{in}^{[k]}_{s}$; in particular, every finite family of $\omega$-cell classes has representatives in a common $B_{(p,s)}$.
\end{enumerate}
\end{lemma}
\begin{proof}
(1) It is the usual construction of the inductive limit of an inductive system of sets.

(2) By Lemma~\ref{LUnifInj}, every $\mathrm{i}^{(k,j)}_{s}$ is injective, so that two elements of $B_{(k,s)}$ are $\Theta_{s}$-related if, and only if, they are equal.

(3) By Definition~\ref{DUfOmegaCells}, $[x]_{j}=[\mathrm{i}^{(k,j)}_{s}(x)]_{k}$ for every $j<k$, so that the image of $\mathrm{in}^{[j]}_{s}$ is contained in that of $\mathrm{in}^{[k]}_{s}$.
\end{proof}

Let $\sig{\Sigma}=(S,\Sigma)$ be a many-sorted signature and $\calg{B}$ a globular $\omega$-categorial $\sig{\Sigma}$-algebra. For every $n\in\omega$, $\Uf^{(n)}_{\sig{\Sigma}}(\Tr^{\gl(n,\omega)}_{\sig{\Sigma}}(\calg{B}))$ is an $n$-categorial $\sig{\Sigma}$-algebra whose cells of sort $s$ are the elements of $B_{(n,s)}$. By Proposition~\ref{PCompatUfTr} and Remark~\ref{RDiscOmegaTrunc}, for $n\leq n'$, $\mathrm{i}^{(n',n)}=(\mathrm{i}^{(n',n)}_{s})_{s\in S}$ is an isomorphism
$$
\mathrm{i}^{(n',n)}\colon\Uf^{(n)}_{\sig{\Sigma}}\left(\Tr^{\gl(n,\omega)}_{\sig{\Sigma}}(\calg{B})\right)\mor\Tr^{(n,n')}_{\sig{\Sigma}}\left(\Uf^{(n')}_{\sig{\Sigma}}\left(\Tr^{\gl(n',\omega)}_{\sig{\Sigma}}(\calg{B})\right)\right);
$$
that is, $\Uf^{(n)}_{\sig{\Sigma}}(\Tr^{\gl(n,\omega)}_{\sig{\Sigma}}(\calg{B}))$ is, up to $\mathrm{i}^{(n',n)}$, the algebra of the $n$-cells of $\Uf^{(n')}_{\sig{\Sigma}}(\Tr^{\gl(n',\omega)}_{\sig{\Sigma}}(\calg{B}))$. We next define the unification of $\calg{B}$.

\begin{definition}\label{DUfOmega}
Let $\sig{\Sigma}=(S,\Sigma)$ be a many-sorted signature and $\calg{B}$ a globular $\omega$-categorial $\sig{\Sigma}$-algebra. The \emph{unification} of $\calg{B}$ is the partial $\Sigma^{(\omega)}$-algebra $\Uf^{(\omega)}_{\sig{\Sigma}}(\calg{B})$ with underlying $S$-sorted set $[B]=([B]_{s})_{s\in S}$ in which, for every operation symbol $\tau$ of $\Sigma^{(\omega)}$, of arity $\mathbf{s}$ and coarity $s$, and every family $([x_{j}]_{p})_{j\in\bb{\mathbf{s}}}$ of $\omega$-cell classes represented in a common level $p$ such that $\tau$ is an operation symbol of $\Sigma^{(p)}$, $([x_{j}]_{p})_{j\in\bb{\mathbf{s}}}\in\Dom(\tau^{\Uf^{(\omega)}_{\sig{\Sigma}}(\calg{B})})$ if, and only if, $(x_{j})_{j\in\bb{\mathbf{s}}}\in\Dom(\tau^{\Uf^{(p)}_{\sig{\Sigma}}(\Tr^{\gl(p,\omega)}_{\sig{\Sigma}}(\calg{B}))})$, and then
$$
\tau^{\Uf^{(\omega)}_{\sig{\Sigma}}(\calg{B})}(([x_{j}]_{p})_{j\in\bb{\mathbf{s}}})=[\tau^{\Uf^{(p)}_{\sig{\Sigma}}(\Tr^{\gl(p,\omega)}_{\sig{\Sigma}}(\calg{B}))}((x_{j})_{j\in\bb{\mathbf{s}}})]_{p}.
$$
For a globular $\omega$-categorial $\sig{\Sigma}$-homomorphism $F\colon\calg{B}\mor\calg{B}'$, $\Uf^{(\omega)}_{\sig{\Sigma}}(F)$ is the $S$-sorted mapping from $[B]$ to $[B']$ that sends $[x]_{k}$ to $[F_{(k,s)}(x)]_{k}$.
\end{definition}

In the following lemma we prove that the above definition is correct.

\begin{lemma}\label{LUfOmegaWellDef}
Let $\sig{\Sigma}=(S,\Sigma)$ be a many-sorted signature, $\calg{B}$ a globular $\omega$-categorial $\sig{\Sigma}$-algebra and $F$ a globular $\omega$-categorial $\sig{\Sigma}$-homomorphism. Then $\Uf^{(\omega)}_{\sig{\Sigma}}(\calg{B})$ and $\Uf^{(\omega)}_{\sig{\Sigma}}(F)$ are well defined.
\end{lemma}
\begin{proof}
\textsf{(Independence of the level)} Let $p\leq p'$. Then $\mathrm{in}^{[p]}_{s}=\mathrm{in}^{[p']}_{s}\circ\mathrm{i}^{(p',p)}_{s}$, and, by Proposition~\ref{PCompatUfTr}, $\mathrm{i}^{(p',p)}$ is an isomorphism from $\Uf^{(p)}_{\sig{\Sigma}}(\Tr^{\gl(p,\omega)}_{\sig{\Sigma}}(\calg{B}))$ onto the algebra of the $p$-cells of $\Uf^{(p')}_{\sig{\Sigma}}(\Tr^{\gl(p',\omega)}_{\sig{\Sigma}}(\calg{B}))$, whose operations are the restrictions of those of $\Uf^{(p')}_{\sig{\Sigma}}(\Tr^{\gl(p',\omega)}_{\sig{\Sigma}}(\calg{B}))$. Hence computing at the level $p$ or at the level $p'$ gives the same domain and the same values.

\textsf{(Homomorphisms)} Since $F$ commutes with the identities, the family of mappings $(\mathrm{in}^{[k]}_{s}\circ F_{(k,s)})_{k\in\omega}$ satisfies the compatibility condition of Lemma~\ref{LUfOmegaCells}~(1), so that $\Uf^{(\omega)}_{\sig{\Sigma}}(F)$ is well defined.

This completes the proof.
\end{proof}

In other words, $\Uf^{(\omega)}_{\sig{\Sigma}}(\calg{B})$ is the inductive limit of the inductive system of categorial algebras $(\Uf^{(n)}_{\sig{\Sigma}}(\Tr^{\gl(n,\omega)}_{\sig{\Sigma}}(\calg{B})),\mathrm{i}^{(n+1,n)})_{n\in\omega}$. In the following proposition we prove that it is an $\omega$-categorial $\sig{\Sigma}$-algebra.

\begin{proposition}\label{PUfOmega}
Let $\sig{\Sigma}=(S,\Sigma)$ be a many-sorted signature and $\calg{B}$ a globular $\omega$-categorial $\sig{\Sigma}$-algebra. Then $\Uf^{(\omega)}_{\sig{\Sigma}}(\calg{B})$ is an $\omega$-categorial $\sig{\Sigma}$-algebra whose $k$-cells of sort $s$ are the classes $[x]_{k}$, with $x\in B_{(k,s)}$, and, for every $n\in\omega$, $\mathrm{in}^{[n]}=(\mathrm{in}^{[n]}_{s})_{s\in S}$ is an isomorphism, natural in $\calg{B}$, from $\Uf^{(n)}_{\sig{\Sigma}}(\Tr^{\gl(n,\omega)}_{\sig{\Sigma}}(\calg{B}))$ onto $\Tr^{(n,\omega)}_{\sig{\Sigma}}(\Uf^{(\omega)}_{\sig{\Sigma}}(\calg{B}))$. Moreover, $\Uf^{(\omega)}_{\sig{\Sigma}}$ is a functor from $\omega\mathsf{Cat}^{\gl}\mathsf{Alg}(\sig{\Sigma})$ to $\omega\mathsf{CatAlg}(\sig{\Sigma})$, the \emph{unification functor}.
\end{proposition}
\begin{proof}
We prove the following claims.

\begin{claim}\label{CUfOmegaA}
Every axiom of $\boldsymbol{\mathcal{C}}^{\sig{\Sigma}^{(\omega)}}$ holds in $\Uf^{(\omega)}_{\sig{\Sigma}}(\calg{B})$.
\end{claim}

Every axiom of $\boldsymbol{\mathcal{C}}^{\sig{\Sigma}^{(\omega)}}$ involves finitely many dimensions and finitely many variables. Given a valuation of its variables in $[B]$, by Lemma~\ref{LUfOmegaCells}~(3), we choose a level $p$ greater than all these dimensions at which all the values of the variables are represented. Since $\mathrm{in}^{[p]}$ is injective and, by Definition~\ref{DUfOmega}, it preserves and reflects the definedness and the values of the operations of $\Sigma^{(p)}$, the axiom holds in $\Uf^{(\omega)}_{\sig{\Sigma}}(\calg{B})$ under the given valuation, because it is an axiom of $\boldsymbol{\mathcal{C}}^{\sig{\Sigma}^{(p)}}$, which holds in the $p$-categorial $\sig{\Sigma}$-algebra $\Uf^{(p)}_{\sig{\Sigma}}(\Tr^{\gl(p,\omega)}_{\sig{\Sigma}}(\calg{B}))$.

Claim~\ref{CUfOmegaA} follows.

\begin{claim}\label{CUfOmegaB}
The $k$-cells of sort $s$ of $\Uf^{(\omega)}_{\sig{\Sigma}}(\calg{B})$ are the classes $[x]_{k}$, with $x\in B_{(k,s)}$. In particular, $\Uf^{(\omega)}_{\sig{\Sigma}}(\calg{B})$ is an $\omega$-categorial $\sig{\Sigma}$-algebra.
\end{claim}

Let $x\in B_{(k,s)}$ and $p>k$. In $\Uf^{(p)}_{\sig{\Sigma}}(\Tr^{\gl(p,\omega)}_{\sig{\Sigma}}(\calg{B}))$, the following chain of equalities holds
\allowdisplaybreaks
\begin{align*}
\src^{k}_{s}\left(\mathrm{i}^{(p,k)}_{s}(x)\right)
&=
\mathrm{i}^{(p,k)}_{s}\left(\src^{(k,p)}_{s}\left(\mathrm{i}^{(p,k)}_{s}(x)\right)\right)
\tag{1}
\\&=
\mathrm{i}^{(p,k)}_{s}(x).
\tag{2}
\end{align*}

The first equality follows from Lemma~\ref{LUnifDescription}; finally, the last equality follows from Lemma~\ref{LUnifInj}. Thus $[x]_{k}=[\mathrm{i}^{(p,k)}_{s}(x)]_{p}$ is a $k$-cell. Conversely, if $[y]_{p}$, with $p>k$, is a $k$-cell, then $y=\mathrm{i}^{(p,k)}_{s}(\src^{(k,p)}_{s}(y))$, so that $[y]_{p}=[\src^{(k,p)}_{s}(y)]_{k}$. Hence every $\omega$-cell class is a $k$-cell for some $k$, and, by Claim~\ref{CUfOmegaA} and Remark~\ref{RSingleSetIdentOmega}, $\Uf^{(\omega)}_{\sig{\Sigma}}(\calg{B})$ is an $\omega$-categorial $\sig{\Sigma}$-algebra.

Claim~\ref{CUfOmegaB} follows.

\begin{claim}\label{CUfOmegaC}
For every $n\in\omega$, $\mathrm{in}^{[n]}$ is an isomorphism from $\Uf^{(n)}_{\sig{\Sigma}}(\Tr^{\gl(n,\omega)}_{\sig{\Sigma}}(\calg{B}))$ onto $\Tr^{(n,\omega)}_{\sig{\Sigma}}(\Uf^{(\omega)}_{\sig{\Sigma}}(\calg{B}))$, natural in $\calg{B}$.
\end{claim}

By Claim~\ref{CUfOmegaB} and Lemma~\ref{LUfOmegaCells}~(2), $\mathrm{in}^{[n]}$ is a bijection from the cells of $\Uf^{(n)}_{\sig{\Sigma}}(\Tr^{\gl(n,\omega)}_{\sig{\Sigma}}(\calg{B}))$ onto the $n$-cells of $\Uf^{(\omega)}_{\sig{\Sigma}}(\calg{B})$, i.e., onto the cells of $\Tr^{(n,\omega)}_{\sig{\Sigma}}(\Uf^{(\omega)}_{\sig{\Sigma}}(\calg{B}))$. By Definition~\ref{DUfOmega}, it preserves and reflects the definedness and the values of the operations of $\Sigma^{(n)}$, so that it is an isomorphism. It is natural in $\calg{B}$ by the definition of $\Uf^{(\omega)}_{\sig{\Sigma}}(F)$.

Claim~\ref{CUfOmegaC} follows.

\begin{claim}\label{CUfOmegaD}
$\Uf^{(\omega)}_{\sig{\Sigma}}$ is a functor.
\end{claim}

Let $F\colon\calg{B}\mor\calg{B}'$ be a globular $\omega$-categorial $\sig{\Sigma}$-homomorphism. For every level $p$, the following equality holds
$$
\Uf^{(\omega)}_{\sig{\Sigma}}(F)\circ\mathrm{in}^{[p]}=\mathrm{in}^{[p]}\circ\Uf^{(p)}_{\sig{\Sigma}}\left(\Tr^{\gl(p,\omega)}_{\sig{\Sigma}}(F)\right),
$$
so that, since $\Uf^{(p)}_{\sig{\Sigma}}(\Tr^{\gl(p,\omega)}_{\sig{\Sigma}}(F))$ is a homomorphism, $\Uf^{(\omega)}_{\sig{\Sigma}}(F)$ commutes with the operations of $\Sigma^{(p)}$. Thus it is a homomorphism. Moreover, $\Uf^{(\omega)}_{\sig{\Sigma}}$ preserves identities and composites, since it acts componentwise on representatives.

Claim~\ref{CUfOmegaD} follows.

This completes the proof.
\end{proof}

Explicitly, the $k$-source of a $p$-cell class $[y]_{p}$, with $k<p$, is $[\src^{(k,p)}_{s}(y)]_{k}$; the $k$-composition of $[x]_{p}$ and $[y]_{p}$ is $[x\circ^{(k,p)}_{s}y]_{p}$, defined if, and only if, $\src^{(k,p)}_{s}(x)=\tgt^{(k,p)}_{s}(y)$; for $k\geq p$, $[y]_{p}$ is a $k$-cell, so that its $k$-source and $k$-target are $[y]_{p}$ itself; and $\sigma$ acts on classes represented at the level $p$ as $\sigma_{p}$.

\subsection{Relationships}\label{SRelOmega}

In this subsection we prove that the $\omega$-dimensional functors are related as the finite-dimensional ones. We begin with graduation and unification.

\begin{proposition}\label{POmegaGdUf}
Let $\sig{\Sigma}=(S,\Sigma)$ be a many-sorted signature. Then the following natural isomorphisms hold
$$
\Uf^{(\omega)}_{\sig{\Sigma}}\circ\Gd^{(\omega)}_{\sig{\Sigma}}\cong\Id_{\omega\mathsf{CatAlg}(\sig{\Sigma})}
\qquad\text{and}\qquad
\Gd^{(\omega)}_{\sig{\Sigma}}\circ\Uf^{(\omega)}_{\sig{\Sigma}}\cong\Id_{\omega\mathsf{Cat}^{\gl}\mathsf{Alg}(\sig{\Sigma})}.
$$
Hence $\Gd^{(\omega)}_{\sig{\Sigma}}$ and $\Uf^{(\omega)}_{\sig{\Sigma}}$ are mutually quasi-inverse equivalences of categories, and there exists an adjoint equivalence in which $\Gd^{(\omega)}_{\sig{\Sigma}}$ is left adjoint to $\Uf^{(\omega)}_{\sig{\Sigma}}$.
\end{proposition}
\begin{proof}
We prove the following claims.

\begin{claim}\label{COmegaGdUfA}
$\Uf^{(\omega)}_{\sig{\Sigma}}\circ\Gd^{(\omega)}_{\sig{\Sigma}}\cong\Id_{\omega\mathsf{CatAlg}(\sig{\Sigma})}$.
\end{claim}

Let $\calg{A}$ be an $\omega$-categorial $\sig{\Sigma}$-algebra. Since the identities of $\Gd^{(\omega)}_{\sig{\Sigma}}(\calg{A})$ are inclusions, two $\omega$-cells $(x,j)$ and $(y,k)$ of $\Gd^{(\omega)}_{\sig{\Sigma}}(\calg{A})$ are $\Theta_{s}$-related if, and only if, $x=y$. Hence the mapping $\zeta_{\calg{A},s}$ that sends $[x]_{k}$ to $x$ is a bijection from the set of the cells of sort $s$ of $\Uf^{(\omega)}_{\sig{\Sigma}}(\Gd^{(\omega)}_{\sig{\Sigma}}(\calg{A}))$ onto the set of the cells of sort $s$ of $\calg{A}$, all of which are $k$-cells for some $k$. For every level $p$, the following chain of equalities holds
\allowdisplaybreaks
\begin{align*}
\Uf^{(p)}_{\sig{\Sigma}}\left(\Tr^{\gl(p,\omega)}_{\sig{\Sigma}}\left(\Gd^{(\omega)}_{\sig{\Sigma}}(\calg{A})\right)\right)
&=
\Uf^{(p)}_{\sig{\Sigma}}\left(\Gd^{(p)}_{\sig{\Sigma}}\left(\Tr^{(p,\omega)}_{\sig{\Sigma}}(\calg{A})\right)\right)
\tag{1}
\\&=
\Tr^{(p,\omega)}_{\sig{\Sigma}}(\calg{A}).
\tag{2}
\end{align*}

The first equality follows from Definition~\ref{DGdOmega}; finally, the last equality follows from Proposition~\ref{PUnifGrad}. Therefore, by Definition~\ref{DUfOmega}, $\zeta_{\calg{A}}=(\zeta_{\calg{A},s})_{s\in S}$ preserves and reflects the definedness and the values of the operations of $\Sigma^{(p)}$. Hence it is an isomorphism from $\Uf^{(\omega)}_{\sig{\Sigma}}(\Gd^{(\omega)}_{\sig{\Sigma}}(\calg{A}))$ onto $\calg{A}$, and $\zeta=(\zeta_{\calg{A}})_{\calg{A}\in\mathrm{Ob}(\omega\mathsf{CatAlg}(\sig{\Sigma}))}$ is a natural transformation from $\Uf^{(\omega)}_{\sig{\Sigma}}\circ\Gd^{(\omega)}_{\sig{\Sigma}}$ to $\Id_{\omega\mathsf{CatAlg}(\sig{\Sigma})}$, since $\Gd^{(\omega)}_{\sig{\Sigma}}$ and $\Uf^{(\omega)}_{\sig{\Sigma}}$ act on homomorphisms by restriction and on representatives, respectively.

Claim~\ref{COmegaGdUfA} follows.

\begin{claim}\label{COmegaGdUfB}
$\Gd^{(\omega)}_{\sig{\Sigma}}\circ\Uf^{(\omega)}_{\sig{\Sigma}}\cong\Id_{\omega\mathsf{Cat}^{\gl}\mathsf{Alg}(\sig{\Sigma})}$.
\end{claim}

Let $\calg{B}$ be a globular $\omega$-categorial $\sig{\Sigma}$-algebra. By Proposition~\ref{PUfOmega}, the cells of sort $(k,s)$ of $\Gd^{(\omega)}_{\sig{\Sigma}}(\Uf^{(\omega)}_{\sig{\Sigma}}(\calg{B}))$ are the classes $[x]_{k}$, with $x\in B_{(k,s)}$, so that the family $\theta_{\calg{B}}=(\mathrm{in}^{[k]}_{s})_{(k,s)\in\omega\times S}$ is a bijection from the cells of $\calg{B}$ onto those of $\Gd^{(\omega)}_{\sig{\Sigma}}(\Uf^{(\omega)}_{\sig{\Sigma}}(\calg{B}))$. It is a homomorphism that reflects definedness. In fact, for $j<k$,
\begin{enumerate}
\item the $(j,k)$-source of $[x]_{k}$ is its $j$-source, which, by the explicit description after Proposition~\ref{PUfOmega}, is $[\src^{(j,k)}_{s}(x)]_{j}$, and similarly for the targets;
\item the $(k,j)$-identity of $[u]_{j}$ is $[u]_{j}=[\mathrm{i}^{(k,j)}_{s}(u)]_{k}$;
\item the $j$-composition of $[x]_{k}$ and $[y]_{k}$ is $[x\circ^{(j,k)}_{s}y]_{k}$, defined if, and only if, $x\circ^{(j,k)}_{s}y$ is defined; and
\item $\sigma_{k}$ acts on $k$-cells as $\sigma$, i.e., as $\sigma_{k}$ on representatives.
\end{enumerate}
Hence $\theta_{\calg{B}}$ is an isomorphism from $\calg{B}$ onto $\Gd^{(\omega)}_{\sig{\Sigma}}(\Uf^{(\omega)}_{\sig{\Sigma}}(\calg{B}))$ and, by Definition~\ref{DUfOmega}, $\theta=(\theta_{\calg{B}})_{\calg{B}\in\mathrm{Ob}(\omega\mathsf{Cat}^{\gl}\mathsf{Alg}(\sig{\Sigma}))}$ is a natural isomorphism from $\Id_{\omega\mathsf{Cat}^{\gl}\mathsf{Alg}(\sig{\Sigma})}$ to $\Gd^{(\omega)}_{\sig{\Sigma}}\circ\Uf^{(\omega)}_{\sig{\Sigma}}$.

Claim~\ref{COmegaGdUfB} follows.

All in all, by Claims~\ref{COmegaGdUfA} and~\ref{COmegaGdUfB}, $\Gd^{(\omega)}_{\sig{\Sigma}}$ and $\Uf^{(\omega)}_{\sig{\Sigma}}$ are mutually quasi-inverse equivalences of categories. Finally, every equivalence of categories can be completed to an adjoint equivalence, see~\cite{sM98}, Theorem~IV.4.1.

This completes the proof.
\end{proof}

We next consider discretization and truncation.

\begin{proposition}\label{POmegaDiscTrunc}
Let $\sig{\Sigma}=(S,\Sigma)$ be a many-sorted signature and $m\in\mathbb{N}$. Then $\Tr^{(m,\omega)}_{\sig{\Sigma}}\circ\Ds^{(\omega,m)}_{\sig{\Sigma}}=\Id_{\mathsf{mCatAlg}(\sig{\Sigma})}$ and $\Ds^{(\omega,m)}_{\sig{\Sigma}}$ is left adjoint to $\Tr^{(m,\omega)}_{\sig{\Sigma}}$, with the identity as unit and, as counit, the natural transformation $\varepsilon$ whose component at an $\omega$-categorial $\sig{\Sigma}$-algebra $\calg{B}$ is the inclusion of the $m$-cells. The same holds for $\Ds^{\gl(\omega,m)}_{\sig{\Sigma}}$ and $\Tr^{\gl(m,\omega)}_{\sig{\Sigma}}$, the component of the counit at the sort $(k,s)$ being the identity, if $k\leq m$, and $\mathrm{i}^{(k,m)}_{s}$, if $k>m$.
\end{proposition}
\begin{proof}
We consider first the single-set case and then the globular case.

\textsf{(Single-set case: unit)} The equation $\Tr^{(m,\omega)}_{\sig{\Sigma}}\circ\Ds^{(\omega,m)}_{\sig{\Sigma}}=\Id_{\mathsf{mCatAlg}(\sig{\Sigma})}$ is Remark~\ref{RDiscOmegaTrunc}~(2) for $n=m$.

\textsf{(Single-set case: counit)} Let $\calg{B}$ be an $\omega$-categorial $\sig{\Sigma}$-algebra and $n\in\mathbb{N}$. By Remark~\ref{RDiscOmegaTrunc}, the following equality holds
$$
\Tr^{(n,\omega)}_{\sig{\Sigma}}\left(\Ds^{(\omega,m)}_{\sig{\Sigma}}\left(\Tr^{(m,\omega)}_{\sig{\Sigma}}(\calg{B})\right)\right)
=
\begin{cases}
\Ds^{(n,m)}_{\sig{\Sigma}}\left(\Tr^{(m,n)}_{\sig{\Sigma}}\left(\Tr^{(n,\omega)}_{\sig{\Sigma}}(\calg{B})\right)\right), & \text{if $m\leq n$;}\\
\Tr^{(n,\omega)}_{\sig{\Sigma}}(\calg{B}), & \text{if $n<m$.}
\end{cases}
$$
Let $\varepsilon^{n}_{\calg{B}}$ be the component of the counit of Proposition~\ref{PDiscTruncAdj} at $\Tr^{(n,\omega)}_{\sig{\Sigma}}(\calg{B})$, if $m\leq n$, and the identity, if $n<m$. All of them are inclusions of $m$-cells, or identities, so that they satisfy the compatibility condition of Corollary~\ref{CLimitNat}, and they induce a natural transformation $\varepsilon=(\varepsilon_{\calg{B}})_{\calg{B}\in\mathrm{Ob}(\omega\mathsf{CatAlg}(\sig{\Sigma}))}$ from $\Ds^{(\omega,m)}_{\sig{\Sigma}}\circ\Tr^{(m,\omega)}_{\sig{\Sigma}}$ to $\Id_{\omega\mathsf{CatAlg}(\sig{\Sigma})}$, whose component $\varepsilon_{\calg{B}}$ at $\calg{B}$ is the inclusion of the $m$-cells of $\calg{B}$ into $\calg{B}$.

\textsf{(Single-set case: triangular identities)} By Definition~\ref{DAdj}, we have to prove the triangular identities
$$
\left(\Tr^{(m,\omega)}_{\sig{\Sigma}}\ast\varepsilon\right)\circ\left(\id\ast\Tr^{(m,\omega)}_{\sig{\Sigma}}\right)=\id_{\Tr^{(m,\omega)}_{\sig{\Sigma}}}
\qquad\text{and}\qquad
\left(\varepsilon\ast\Ds^{(\omega,m)}_{\sig{\Sigma}}\right)\circ\left(\Ds^{(\omega,m)}_{\sig{\Sigma}}\ast\id\right)=\id_{\Ds^{(\omega,m)}_{\sig{\Sigma}}}.
$$
We recall that a homomorphism is determined by its domain, its codomain and its underlying many-sorted mapping.

Let $\calg{B}$ be an $\omega$-categorial $\sig{\Sigma}$-algebra. By the equation of the unit, the domain of $\Tr^{(m,\omega)}_{\sig{\Sigma}}(\varepsilon_{\calg{B}})$ is $\Tr^{(m,\omega)}_{\sig{\Sigma}}(\calg{B})$, whose cells of sort $s$ are the $m$-cells of sort $s$ of $\calg{B}$. Let $s\in S$ and $x$ an $m$-cell of sort $s$ of $\calg{B}$. Then the following chain of equalities holds
\allowdisplaybreaks
\begin{align*}
\left(\left(\left(\Tr^{(m,\omega)}_{\sig{\Sigma}}\ast\varepsilon\right)\circ\left(\id\ast\Tr^{(m,\omega)}_{\sig{\Sigma}}\right)\right)_{\calg{B}}\right)_{s}(x)
&=
\left(\Tr^{(m,\omega)}_{\sig{\Sigma}}\left(\varepsilon_{\calg{B}}\right)\circ\id_{\Tr^{(m,\omega)}_{\sig{\Sigma}}(\calg{B})}\right)_{s}(x)
\tag{1}
\\&=
\left(\varepsilon_{\calg{B}}\right)_{s}(x)
\tag{2}
\\&=
x
\tag{3}
\\&=
\left(\id_{\Tr^{(m,\omega)}_{\sig{\Sigma}}}\right)_{\calg{B},s}(x).
\tag{4}
\end{align*}

The first equality unravels the definitions of the horizontal and vertical composition of natural transformations; the second equality follows from the fact that $\id_{\Tr^{(m,\omega)}_{\sig{\Sigma}}(\calg{B})}$ is the identity homomorphism and from Proposition~\ref{PGenHepDerivorFunctor}, since $\varphi^{\Tr^{(m,\omega)}_{\sig{\Sigma}}}=\id_{S}$; the third equality follows from the fact that $\varepsilon_{\calg{B}}$ is an inclusion; finally, the last equality unravels the definition of the identity natural transformation at $\Tr^{(m,\omega)}_{\sig{\Sigma}}$. Hence the first triangular identity holds.

Let $\calg{A}$ be an $m$-categorial $\sig{\Sigma}$-algebra. By the equation of the unit, the domain of $\varepsilon_{\Ds^{(\omega,m)}_{\sig{\Sigma}}(\calg{A})}$ is $\Ds^{(\omega,m)}_{\sig{\Sigma}}(\calg{A})$, and, by Proposition~\ref{PDiscOmegaDerivor}, every cell of $\Ds^{(\omega,m)}_{\sig{\Sigma}}(\calg{A})$ is an $m$-cell. Let $s\in S$ and $x$ a cell of sort $s$ of $\Ds^{(\omega,m)}_{\sig{\Sigma}}(\calg{A})$. Then the following chain of equalities holds
\allowdisplaybreaks
\begin{align*}
\left(\left(\left(\varepsilon\ast\Ds^{(\omega,m)}_{\sig{\Sigma}}\right)\circ\left(\Ds^{(\omega,m)}_{\sig{\Sigma}}\ast\id\right)\right)_{\calg{A}}\right)_{s}(x)
&=
\left(\varepsilon_{\Ds^{(\omega,m)}_{\sig{\Sigma}}(\calg{A})}\circ\Ds^{(\omega,m)}_{\sig{\Sigma}}\left(\id_{\calg{A}}\right)\right)_{s}(x)
\tag{1}
\\&=
\left(\varepsilon_{\Ds^{(\omega,m)}_{\sig{\Sigma}}(\calg{A})}\right)_{s}(x)
\tag{2}
\\&=
x
\tag{3}
\\&=
\left(\id_{\Ds^{(\omega,m)}_{\sig{\Sigma}}}\right)_{\calg{A},s}(x).
\tag{4}
\end{align*}

The first equality unravels the definitions of the horizontal and vertical composition of natural transformations; the second equality follows from the fact that $\Ds^{(\omega,m)}_{\sig{\Sigma}}$ preserves identities; the third equality follows from the fact that $\varepsilon_{\Ds^{(\omega,m)}_{\sig{\Sigma}}(\calg{A})}$ is the inclusion of the $m$-cells of $\Ds^{(\omega,m)}_{\sig{\Sigma}}(\calg{A})$ and $x$ is one of them; finally, the last equality unravels the definition of the identity natural transformation at $\Ds^{(\omega,m)}_{\sig{\Sigma}}$. Hence the second triangular identity holds.

\textsf{(Globular case: unit)} The equation $\Tr^{\gl(m,\omega)}_{\sig{\Sigma}}\circ\Ds^{\gl(\omega,m)}_{\sig{\Sigma}}=\Id_{\mathsf{mCat}^{\gl}\mathsf{Alg}(\sig{\Sigma})}$ is Remark~\ref{RDiscOmegaTrunc}~(2) for $n=m$.

\textsf{(Globular case: counit)} Let $\calg{B}$ be a globular $\omega$-categorial $\sig{\Sigma}$-algebra and $n\in\mathbb{N}$. By Remark~\ref{RDiscOmegaTrunc}, the following equality holds
$$
\Tr^{\gl(n,\omega)}_{\sig{\Sigma}}\left(\Ds^{\gl(\omega,m)}_{\sig{\Sigma}}\left(\Tr^{\gl(m,\omega)}_{\sig{\Sigma}}(\calg{B})\right)\right)
=
\begin{cases}
\Ds^{\gl(n,m)}_{\sig{\Sigma}}\left(\Tr^{\gl(m,n)}_{\sig{\Sigma}}\left(\Tr^{\gl(n,\omega)}_{\sig{\Sigma}}(\calg{B})\right)\right), & \text{if $m\leq n$;}\\
\Tr^{\gl(n,\omega)}_{\sig{\Sigma}}(\calg{B}), & \text{if $n<m$.}
\end{cases}
$$
Let $\varepsilon^{\gl,n}_{\calg{B}}$ be the component of the counit of Proposition~\ref{PDiscTruncAdjGl} at $\Tr^{\gl(n,\omega)}_{\sig{\Sigma}}(\calg{B})$, if $m\leq n$, and the identity, if $n<m$. By Lemma~\ref{LDiscTruncCounitGl}, the component of $\varepsilon^{\gl,n}_{\calg{B}}$ at the sort $(k,s)$, for $k\leq n$, is the identity, if $k\leq m$, and $\mathrm{i}^{(k,m)}_{s}$, if $k>m$. Hence, for $n\leq n'$, the components of $\varepsilon^{\gl,n}_{\calg{B}}$ are the components of $\varepsilon^{\gl,n'}_{\calg{B}}$ at the sorts of dimension at most $n$, i.e., $\Tr^{\gl(n,n')}_{\sig{\Sigma}}(\varepsilon^{\gl,n'}_{\calg{B}})=\varepsilon^{\gl,n}_{\calg{B}}$. Thus the family satisfies the compatibility condition of Corollary~\ref{CLimitNatGl}, and it induces a natural transformation $\varepsilon^{\gl}=(\varepsilon^{\gl}_{\calg{B}})_{\calg{B}\in\mathrm{Ob}(\omega\mathsf{Cat}^{\gl}\mathsf{Alg}(\sig{\Sigma}))}$ from $\Ds^{\gl(\omega,m)}_{\sig{\Sigma}}\circ\Tr^{\gl(m,\omega)}_{\sig{\Sigma}}$ to $\Id_{\omega\mathsf{Cat}^{\gl}\mathsf{Alg}(\sig{\Sigma})}$, whose component at $\calg{B}$ and at the sort $(k,s)$ is the identity, if $k\leq m$, and $\mathrm{i}^{(k,m)}_{s}$, if $k>m$.

\textsf{(Globular case: triangular identities)} By Definition~\ref{DAdj}, we have to prove the triangular identities
$$
\left(\Tr^{\gl(m,\omega)}_{\sig{\Sigma}}\ast\varepsilon^{\gl}\right)\circ\left(\id\ast\Tr^{\gl(m,\omega)}_{\sig{\Sigma}}\right)=\id_{\Tr^{\gl(m,\omega)}_{\sig{\Sigma}}}
\quad\text{and}\quad
\left(\varepsilon^{\gl}\ast\Ds^{\gl(\omega,m)}_{\sig{\Sigma}}\right)\circ\left(\Ds^{\gl(\omega,m)}_{\sig{\Sigma}}\ast\id\right)=\id_{\Ds^{\gl(\omega,m)}_{\sig{\Sigma}}}.
$$

Let $\calg{B}$ be a globular $\omega$-categorial $\sig{\Sigma}$-algebra. By the equation of the unit, the domain of $\Tr^{\gl(m,\omega)}_{\sig{\Sigma}}(\varepsilon^{\gl}_{\calg{B}})$ is $\Tr^{\gl(m,\omega)}_{\sig{\Sigma}}(\calg{B})$. Let $(k,s)\in(m+1)\times S$ and $x$ a $k$-cell of sort $s$ of $\calg{B}$. Then the following chain of equalities holds
\allowdisplaybreaks
\begin{align*}
\left(\left(\left(\Tr^{\gl(m,\omega)}_{\sig{\Sigma}}\ast\varepsilon^{\gl}\right)\circ\left(\id\ast\Tr^{\gl(m,\omega)}_{\sig{\Sigma}}\right)\right)_{\calg{B}}\right)_{(k,s)}(x)
&=
\left(\varepsilon^{\gl}_{\calg{B}}\right)_{(k,s)}(x)
\tag{1}
\\&=
x
\tag{2}
\\&=
\left(\id_{\Tr^{\gl(m,\omega)}_{\sig{\Sigma}}}\right)_{\calg{B},(k,s)}(x).
\tag{3}
\end{align*}

The first equality unravels the definitions of the horizontal and vertical composition of natural transformations, taking into account Proposition~\ref{PGenHepDerivorFunctor} and the fact that $\varphi^{\Tr^{\gl(m,\omega)}_{\sig{\Sigma}}}(k,s)=(k,s)$; the second equality follows from the description of $\varepsilon^{\gl}_{\calg{B}}$, since $k\leq m$; finally, the last equality unravels the definition of the identity natural transformation at $\Tr^{\gl(m,\omega)}_{\sig{\Sigma}}$. Hence the first triangular identity holds.

Let $\calg{A}$ be a globular $m$-categorial $\sig{\Sigma}$-algebra. By the equation of the unit, the domain of $\varepsilon^{\gl}_{\Ds^{\gl(\omega,m)}_{\sig{\Sigma}}(\calg{A})}$ is $\Ds^{\gl(\omega,m)}_{\sig{\Sigma}}(\calg{A})$. Let $(k,s)\in\omega\times S$ and $x$ a cell of sort $(k,s)$ of $\Ds^{\gl(\omega,m)}_{\sig{\Sigma}}(\calg{A})$. If $k\leq m$, then $(\varepsilon^{\gl}_{\Ds^{\gl(\omega,m)}_{\sig{\Sigma}}(\calg{A})})_{(k,s)}(x)=x$. If $k>m$, then, by the explicit description of $\Ds^{\gl(\omega,m)}_{\sig{\Sigma}}$ after Definition~\ref{DDiscretizationOmegaGl}, the $(k,m)$-identities of $\Ds^{\gl(\omega,m)}_{\sig{\Sigma}}(\calg{A})$ are identity mappings, so that
$$
\left(\varepsilon^{\gl}_{\Ds^{\gl(\omega,m)}_{\sig{\Sigma}}(\calg{A})}\right)_{(k,s)}(x)
=
\mathrm{i}^{(k,m)\,\Ds^{\gl(\omega,m)}_{\sig{\Sigma}}(\calg{A})}_{s}(x)
=
x.
$$
Hence $\varepsilon^{\gl}_{\Ds^{\gl(\omega,m)}_{\sig{\Sigma}}(\calg{A})}=\id_{\Ds^{\gl(\omega,m)}_{\sig{\Sigma}}(\calg{A})}$ and the following chain of equalities holds
\allowdisplaybreaks
\begin{align*}
\left(\left(\varepsilon^{\gl}\ast\Ds^{\gl(\omega,m)}_{\sig{\Sigma}}\right)\circ\left(\Ds^{\gl(\omega,m)}_{\sig{\Sigma}}\ast\id\right)\right)_{\calg{A}}
&=
\varepsilon^{\gl}_{\Ds^{\gl(\omega,m)}_{\sig{\Sigma}}(\calg{A})}\circ\Ds^{\gl(\omega,m)}_{\sig{\Sigma}}\left(\id_{\calg{A}}\right)
\tag{1}
\\&=
\id_{\Ds^{\gl(\omega,m)}_{\sig{\Sigma}}(\calg{A})}
\tag{2}
\\&=
\left(\id_{\Ds^{\gl(\omega,m)}_{\sig{\Sigma}}}\right)_{\calg{A}}.
\tag{3}
\end{align*}

The first equality unravels the definitions of the horizontal and vertical composition of natural transformations; the second equality follows from the above and from the fact that $\Ds^{\gl(\omega,m)}_{\sig{\Sigma}}$ preserves identities; finally, the last equality unravels the definition of the identity natural transformation at $\Ds^{\gl(\omega,m)}_{\sig{\Sigma}}$. Hence the second triangular identity holds.

This completes the proof.
\end{proof}

\section{An algebraic square}\label{SAlgSquare}

In this section, we show that graduation, unification, discretization and truncation fit together in algebraic squares, both between two finite dimensions and between a finite dimension and $\omega$. We begin by recalling that four adjunctions form an \emph{algebraic square} when there exist natural isomorphisms between the appropriate composites, as described in the following definition, see~\cite{CVCLRM26}.

\begin{definition}\label{DAlgSquare}
Let $\mathsf{C}$, $\mathsf{C}'$, $\mathsf{D}$ and $\mathsf{D}'$ be four categories. Moreover let
\begin{enumerate}
    \item
    $(G, F, \eta, \varepsilon)$ be an adjunction with $F\colon\mathsf{C}\mor\mathsf{D}$ left adjoint to $G\colon\mathsf{D}\mor\mathsf{C}$.

    \item
    $(G', F', \eta', \varepsilon')$ be an adjunction with $F'\colon\mathsf{C}'\mor\mathsf{D}'$ left adjoint to $G'\colon\mathsf{D}'\mor\mathsf{C}'$.

    \item
    $(K, J, \eta^{\mathsf{C}}, \varepsilon^{\mathsf{C}})$ be an adjunction with $J\colon\mathsf{C}\mor\mathsf{C}'$ left adjoint to $K\colon\mathsf{C}'\mor\mathsf{C}$.

    \item
    $(I, H, \eta^{\mathsf{D}}, \varepsilon^{\mathsf{D}})$ be an adjunction with $H\colon\mathsf{D}\mor\mathsf{D}'$ left adjoint to $I\colon\mathsf{D}'\mor\mathsf{D}$.
\end{enumerate}
We will say that these adjunctions form an \emph{algebraic square} if there exists a pair of natural isomorphisms $\alpha$ from $F' \circ J$ to $H \circ F$ and $\beta$ from $G \circ I$ to $K \circ G'$.

\begin{center}
\begin{tikzpicture}
    [ACliment/.style={-{To [angle'=45, length=5.75pt, width=4pt, round]}},
    RHACliment/.style={right hook-{To [angle'=45, length=5.75pt, width=4pt, round]}, font=\scriptsize},
    scale=1
    ]
    
    \node[] (C') at (0,0) {\textcolor{white}{$AAAA$}}; %Smu
    \node[] () at (0,0) {$\mathsf{C}'$};
    \node[] (D') at (3,0) {\textcolor{white}{$AAAA$}}; %Amu
    \node[] () at (3,0) {$\mathsf{D}'$};
    \node[] () at (1.5,0) {\rotatebox{90}{$\vdash$}};
    
    \draw[ACliment]  (C'.40) to node [above] {$F'$} (D'.140);
    \draw[ACliment]  (D'.220) to node [below] {$G'$} (C'.320);
    
    
    \node[] (C) at (0,3) {\textcolor{white}{$AAAA$}}; %Sz
    \node[] () at (0,3) {$\mathsf{C}$};
    \node[] (D) at (3,3) {\textcolor{white}{$AAAA$}}; %Az
    \node[] () at (3,3) {$\mathsf{D}$};
    \node[] () at (1.5,3) {\rotatebox{90}{$\vdash$}};
    \node[] () at (0,1.5) {\rotatebox{0}{$\vdash$}};
    \node[] () at (3,1.5) {\rotatebox{0}{$\vdash$}};

    \node[] () at (1.5,1.5) {$(\alpha, \beta)$};
    
    \draw[ACliment]  (C.40) to node [above] {$F$} (D.140);
    \draw[ACliment]  (D.220) to node [below] {$G$} (C.320);
    
    \draw[ACliment]  (C.310) to node [right] {$J$} (C'.50);
    \draw[ACliment]  (C'.140) to node [left] {$K$} (C.220);
    
    \draw[ACliment]  (D.310) to node [right] {$H$} (D'.50);
    \draw[ACliment]  (D'.140) to node [left] {$I$} (D.220);
\end{tikzpicture}
\end{center}
\end{definition}

We next show that graduation, unification, discretization and truncation form algebraic squares, both between two finite dimensions $m\leq n$ and between a finite dimension $m$ and $\omega$. In both cases the horizontal adjunctions relate the single-set and the globular presentations and the vertical ones relate the dimensions.

\begin{proposition}\label{PAlgSquareFinite}
Let $\sig{\Sigma}=(S,\Sigma)$ be a many-sorted signature and $m,n\in\mathbb{N}$ with $m\leq n$, and let
$$
\mathsf{C}=\mathsf{mCatAlg}(\sig{\Sigma}),\quad
\mathsf{D}=\mathsf{mCat}^{\gl}\mathsf{Alg}(\sig{\Sigma}),\quad
\mathsf{C}'=\mathsf{nCatAlg}(\sig{\Sigma}),\quad
\mathsf{D}'=\mathsf{nCat}^{\gl}\mathsf{Alg}(\sig{\Sigma}).
$$
Then the adjunctions
\begin{enumerate}
\item $\Gd^{(m)}_{\sig{\Sigma}}\dashv\Uf^{(m)}_{\sig{\Sigma}}$ and $\Gd^{(n)}_{\sig{\Sigma}}\dashv\Uf^{(n)}_{\sig{\Sigma}}$ of Proposition~\ref{PGdUfAdj}, and
\item $\Ds^{(n,m)}_{\sig{\Sigma}}\dashv \Tr^{(m,n)}_{\sig{\Sigma}}$ and $\Ds^{\gl(n,m)}_{\sig{\Sigma}}\dashv \Tr^{\gl(m,n)}_{\sig{\Sigma}}$ of Propositions~\ref{PDiscTruncAdj} and~\ref{PDiscTruncAdjGl}
\end{enumerate}
form an algebraic square, with $F=\Gd^{(m)}_{\sig{\Sigma}}$, $G=\Uf^{(m)}_{\sig{\Sigma}}$, $F'=\Gd^{(n)}_{\sig{\Sigma}}$, $G'=\Uf^{(n)}_{\sig{\Sigma}}$, $J=\Ds^{(n,m)}_{\sig{\Sigma}}$, $K=\Tr^{(m,n)}_{\sig{\Sigma}}$, $H=\Ds^{\gl(n,m)}_{\sig{\Sigma}}$ and $I=\Tr^{\gl(m,n)}_{\sig{\Sigma}}$, in which $\alpha$ is the identity of $\Gd^{(n)}_{\sig{\Sigma}}\circ\Ds^{(n,m)}_{\sig{\Sigma}}=\Ds^{\gl(n,m)}_{\sig{\Sigma}}\circ\Gd^{(m)}_{\sig{\Sigma}}$ and $\beta=\beta^{(m,n)}$ is the natural isomorphism from $\Uf^{(m)}_{\sig{\Sigma}}\circ \Tr^{\gl(m,n)}_{\sig{\Sigma}}$ to $\Tr^{(m,n)}_{\sig{\Sigma}}\circ\Uf^{(n)}_{\sig{\Sigma}}$ given by the $(n,m)$-identities.

\begin{center}
\begin{tikzpicture}
    [ACliment/.style={-{To [angle'=45, length=5.75pt, width=4pt, round]}},
    scale=1
    ]
    \node (C) at (0,3.5) {$\mathsf{mCatAlg}(\sig{\Sigma})$};
    \node (D) at (7,3.5) {$\mathsf{mCat}^{\gl}\mathsf{Alg}(\sig{\Sigma})$};
    \node (C') at (0,0) {$\mathsf{nCatAlg}(\sig{\Sigma})$};
    \node (D') at (7,0) {$\mathsf{nCat}^{\gl}\mathsf{Alg}(\sig{\Sigma})$};
    \draw[ACliment] ([yshift=5pt]C.east) to node [above] {$\Gd^{(m)}_{\sig{\Sigma}}$} ([yshift=5pt]D.west);
    \draw[ACliment] ([yshift=-5pt]D.west) to node [below] {$\Uf^{(m)}_{\sig{\Sigma}}$} ([yshift=-5pt]C.east);
    \draw[ACliment] ([yshift=5pt]C'.east) to node [above] {$\Gd^{(n)}_{\sig{\Sigma}}$} ([yshift=5pt]D'.west);
    \draw[ACliment] ([yshift=-5pt]D'.west) to node [below] {$\Uf^{(n)}_{\sig{\Sigma}}$} ([yshift=-5pt]C'.east);
    \draw[ACliment] ([xshift=8pt]C.south) to node [right] {$\Ds^{(n,m)}_{\sig{\Sigma}}$} ([xshift=8pt]C'.north);
    \draw[ACliment] ([xshift=-8pt]C'.north) to node [left] {$\Tr^{(m,n)}_{\sig{\Sigma}}$} ([xshift=-8pt]C.south);
    \draw[ACliment] ([xshift=8pt]D.south) to node [right] {$\Ds^{\gl(n,m)}_{\sig{\Sigma}}$} ([xshift=8pt]D'.north);
    \draw[ACliment] ([xshift=-8pt]D'.north) to node [left] {$\Tr^{\gl(m,n)}_{\sig{\Sigma}}$} ([xshift=-8pt]D.south);
    \node at ($(C.east)!0.5!(D.west)$) {\rotatebox{90}{$\vdash$}};
    \node at ($(C'.east)!0.5!(D'.west)$) {\rotatebox{90}{$\vdash$}};
    \node at ($(C.south)!0.5!(C'.north)$) {$\vdash$};
    \node at ($(D.south)!0.5!(D'.north)$) {$\vdash$};
    \node at (3.5,1.75) {$(\id,\beta^{(m,n)})$};
\end{tikzpicture}
\end{center}
\end{proposition}
\begin{proof}
By Definition~\ref{DAlgSquare}, we must check that the four pairs of functors involved are adjunctions and that there exist the natural isomorphisms $\alpha$ and $\beta$.

\textsf{(Adjunctions)} The four adjunctions are those of Propositions~\ref{PGdUfAdj}, \ref{PDiscTruncAdj} and~\ref{PDiscTruncAdjGl}.

\textsf{(The natural isomorphism $\alpha$)} By Proposition~\ref{PCompatGdDs}, the following equality holds
$$
F'\circ J=\Gd^{(n)}_{\sig{\Sigma}}\circ\Ds^{(n,m)}_{\sig{\Sigma}}=\Ds^{\gl(n,m)}_{\sig{\Sigma}}\circ\Gd^{(m)}_{\sig{\Sigma}}=H\circ F,
$$
so that we can take $\alpha$ to be the identity.

\textsf{(The natural isomorphism $\beta$)} By Proposition~\ref{PCompatUfTr}, $\beta^{(m,n)}$ is a natural isomorphism from $G\circ I=\Uf^{(m)}_{\sig{\Sigma}}\circ \Tr^{\gl(m,n)}_{\sig{\Sigma}}$ to $K\circ G'=\Tr^{(m,n)}_{\sig{\Sigma}}\circ\Uf^{(n)}_{\sig{\Sigma}}$.

This completes the proof.
\end{proof}

We next consider the algebraic square between a finite dimension and $\omega$.

\begin{proposition}\label{PAlgSquareOmega}
Let $\sig{\Sigma}=(S,\Sigma)$ be a many-sorted signature and $m\in\mathbb{N}$, and let
$$
\mathsf{C}=\mathsf{mCatAlg}(\sig{\Sigma}),\quad
\mathsf{D}=\mathsf{mCat}^{\gl}\mathsf{Alg}(\sig{\Sigma}),\quad
\mathsf{C}'=\omega\mathsf{CatAlg}(\sig{\Sigma}),\quad
\mathsf{D}'=\omega\mathsf{Cat}^{\gl}\mathsf{Alg}(\sig{\Sigma}).
$$
Then the adjunctions
\begin{enumerate}
\item $\Gd^{(m)}_{\sig{\Sigma}}\dashv\Uf^{(m)}_{\sig{\Sigma}}$ of Proposition~\ref{PGdUfAdj} and the adjoint equivalence $\Gd^{(\omega)}_{\sig{\Sigma}}\dashv\Uf^{(\omega)}_{\sig{\Sigma}}$ of Proposition~\ref{POmegaGdUf}, and
\item $\Ds^{(\omega,m)}_{\sig{\Sigma}}\dashv \Tr^{(m,\omega)}_{\sig{\Sigma}}$ and $\Ds^{\gl(\omega,m)}_{\sig{\Sigma}}\dashv \Tr^{\gl(m,\omega)}_{\sig{\Sigma}}$ of Proposition~\ref{POmegaDiscTrunc}
\end{enumerate}
form an algebraic square, with $F=\Gd^{(m)}_{\sig{\Sigma}}$, $G=\Uf^{(m)}_{\sig{\Sigma}}$, $F'=\Gd^{(\omega)}_{\sig{\Sigma}}$, $G'=\Uf^{(\omega)}_{\sig{\Sigma}}$, $J=\Ds^{(\omega,m)}_{\sig{\Sigma}}$, $K=\Tr^{(m,\omega)}_{\sig{\Sigma}}$, $H=\Ds^{\gl(\omega,m)}_{\sig{\Sigma}}$ and $I=\Tr^{\gl(m,\omega)}_{\sig{\Sigma}}$, in which $\alpha$ is the identity of $\Gd^{(\omega)}_{\sig{\Sigma}}\circ\Ds^{(\omega,m)}_{\sig{\Sigma}}=\Ds^{\gl(\omega,m)}_{\sig{\Sigma}}\circ\Gd^{(m)}_{\sig{\Sigma}}$ and $\beta$ is the natural isomorphism $\mathrm{in}^{[m]}$ of Proposition~\ref{PUfOmega}, from $\Uf^{(m)}_{\sig{\Sigma}}\circ \Tr^{\gl(m,\omega)}_{\sig{\Sigma}}$ to $\Tr^{(m,\omega)}_{\sig{\Sigma}}\circ\Uf^{(\omega)}_{\sig{\Sigma}}$, which sends every $m$-cell $x$ to its class $[x]_{m}$.
\end{proposition}
\begin{proof}
By Definition~\ref{DAlgSquare}, we must check that the four pairs of functors involved are adjunctions and that there exist the natural isomorphisms $\alpha$ and $\beta$.

\textsf{(Adjunctions)} The four adjunctions are those of Propositions~\ref{PGdUfAdj}, \ref{POmegaGdUf} and~\ref{POmegaDiscTrunc}.

\textsf{(The natural isomorphism $\alpha$)} For every $p\in\mathbb{N}$, by the definitions of the functors involved,
\begin{gather*}
\Tr^{\gl(p,\omega)}_{\sig{\Sigma}}\circ\Gd^{(\omega)}_{\sig{\Sigma}}\circ\Ds^{(\omega,m)}_{\sig{\Sigma}}=\Gd^{(p)}_{\sig{\Sigma}}\circ F_{p},\\
\Tr^{\gl(p,\omega)}_{\sig{\Sigma}}\circ\Ds^{\gl(\omega,m)}_{\sig{\Sigma}}\circ\Gd^{(m)}_{\sig{\Sigma}}=F^{\gl}_{p}\circ\Gd^{(m)}_{\sig{\Sigma}},
\end{gather*}
with $F_{p}$ and $F^{\gl}_{p}$ as in Remark~\ref{RDiscOmegaTrunc}. If $m\leq p$, $\Gd^{(p)}_{\sig{\Sigma}}\circ\Ds^{(p,m)}_{\sig{\Sigma}}=\Ds^{\gl(p,m)}_{\sig{\Sigma}}\circ\Gd^{(m)}_{\sig{\Sigma}}$, and, if $p\leq m$, $\Gd^{(p)}_{\sig{\Sigma}}\circ \Tr^{(p,m)}_{\sig{\Sigma}}=\Tr^{\gl(p,m)}_{\sig{\Sigma}}\circ\Gd^{(m)}_{\sig{\Sigma}}$, by Propositions~\ref{PCompatGdDs} and~\ref{PCompatGdTr}. Hence both functors have the same truncations and, by the uniqueness in Proposition~\ref{PUnivPropGlCatAlg}, they are equal, so that we can take $\alpha$ to be the identity.

\textsf{(The natural isomorphism $\beta$)} By Proposition~\ref{PUfOmega}, $\mathrm{in}^{[m]}$ is a natural isomorphism from $\Uf^{(m)}_{\sig{\Sigma}}\circ \Tr^{\gl(m,\omega)}_{\sig{\Sigma}}$ to $\Tr^{(m,\omega)}_{\sig{\Sigma}}\circ\Uf^{(\omega)}_{\sig{\Sigma}}$.

This completes the proof.
\end{proof}



% \appendix
%
% \section{Many-sorted sets and many-sorted algebras}\label{S0C}
%
% In this section we collect the basic facts, without proofs, about many-sorted sets and many-sorted (total) algebras, with special emphasis on the free many-sorted algebra  that we will need afterwards.
%
% \subsection{Many-sorted sets}
%
% \begin{assumption}
% From now on $S$ stands for a set of sorts in $\boldsymbol{\mathcal{U}}$, fixed once and for all.
% \end{assumption}
%
% \begin{definition}
% An $S$-\emph{sorted set} is a mapping $A = (A_{s})_{s\in S}$ from $S$ to $\boldsymbol{\mathcal{U}}$. If $A$ and $X$ are $S$-sorted sets, then we will say that $X$ is an 
% $S$-\emph{sorted subset} or, to abbreviate, a \emph{subset} of $A$, denoted by $X\subseteq A$, if, for every $s\in S$, $X_{s}\subseteq A_{s}$. We will denote by $\mathrm{Sub}(A)$ the set of all $S$-sorted sets $X$ such that $X\subseteq A$.
% \end{definition}
%
% \begin{definition}
% Let $A$ and $B$ be $S$-sorted sets.  The \emph{cartesian product of} $A$ \emph{and} $B$, denoted by $A\times B$, is the $S$-sorted set $\left(A_{s}\times B_{s}\right)_{s\in S}.$
%
% Let $\Phi$ be an $S$-sorted set. We will say that $\Phi$ is an $S$-\emph{sorted relation from} $A$ \emph{to} $B$ if $\Phi\subseteq A\times B$. Thus, for every $s\in S$, $\Phi_{s}\subseteq A_{s}\times B_{s}$. We denote by $\Rel(A,B)$ the set of all $S$-sorted relations from $A$ to $B$. If $A=B$, then we write $\Rel(A)$ instead of $\Rel(A,A)$ and call its elements $S$-\emph{sorted relations on} $A$. 
%
% The \emph{diagonal of} $A$, denoted by $\Delta_{A}$, is the $S$-sorted relation on $A$ defined, for every $s\in S$, as $\Delta_{A_{s}}$, and the \emph{codiagonal of} $A$, denoted by $\nabla_{A}$, is the $S$-sorted relation on $A$ defined, for every $s\in S$, as $\nabla_{A_{s}} = A_{s}\times A_{s}$.
%
% Let $\Phi$ be an $S$-sorted relation from $A$ to $B$ and $\Psi$ an $S$-sorted relation from $B$ to $C$, then the \emph{composition of} $\Phi$ \emph{and} $\Psi$, denoted by $\Psi\comp\Phi$, is the $S$-sorted relation from $A$ to $C$ defined, for every $s\in S$, as 
% $$
% \Psi_{s}\comp\Phi_{s}
% =\left\lbrace
% (x,z)\in A_{s}\times C_{s}
% \mid \exists\,y\in B_{s}\,
% \left(
% (x,y)\in \Phi_{s}\And (y,z)\in \Psi_{s}
% \right) 
% \right\rbrace.
% $$ 
% This composition is associative and the diagonal relation $\Delta_{A}$ is a neutral element for it. 
%
% Let $\Phi$ an $S$-sorted relation from $A$ to $B$. Then the \emph{inverse} of $\Phi$, denoted by $\Phi^{-1}$, is the $S$-sorted relation from $B$ to $A$ defined, for every $s\in S$, as
% $$
% \Phi_{s}^{-1}
% =\left\lbrace(y,x)\in B_{s}\times A_{s}\mid (x,y)\in \Phi_{s}
% \right\rbrace.
% $$
%
% Let $\Phi$ be an $S$-sorted relation from $A$ to $B$, $X$ an $S$-sorted subset of $A$, and $Y$ an $S$-sorted
% subset of $B$, then the \emph{direct image} % 
% of $X$ under $\Phi$, denoted by $\Phi[X]$, is the $S$-sorted subset of $B$ defined, for every $s\in S$, as
% $$
% \Phi[X]_{s} =\left\lbrace\,b\in B_ {s}\mid \exists\, x\in X_ {s}((x,b)\in \Phi_
% {s})\,\right\rbrace,
% $$
% therefore $\Phi[X] = (\Phi_ {s}[X_ {s}])_ {s\in S}$, and the \emph{inverse image} of $Y$ under $\Phi$, denoted by $\Phi^{-1}[Y]$, is the $S$-sorted subset of $A$ defined, for every $s\in S$, as
% $$
% \Phi^{-1}[Y]_{s} =\left\lbrace\,a\in A_ {s}\mid \exists\, y\in Y_ {s}((a,y)\in \Phi_
% {s})\,\right\rbrace,
% $$
% therefore $\Phi^{-1}[Y] = (\Phi^{-1}_{s}[Y_ {s}])_{s\in S}$. For an $S$-sorted relation $\Phi$ on an $S$-sorted set $A$ we will call $\Phi[A]$ and $\Phi^{-1}[A]$ the \emph{image} and the \emph{domain} of $\Phi$ respectively and we will denote them by $\mathrm{Im}(\Phi)$ and $\Dom(\Phi)$ respectively. Moreover, we will call 
% $\Dom(\Phi)\cup\mathrm{Im}(\Phi)$ the \emph{field} of $\Phi$ and we will denote it by $\mathrm{Fld}(\Phi)$.
%
% An $S$-\emph{sorted function from} $A$ \emph{to} $B$ is a functional $S$-sorted relation $F$ from $A$ to $B$, i.e., an $S$-sorted relation $F$ from $A$ to $B$ such that, for every $s\in S$, $F_{s}$ is a function from $A_{s}$ to $B_{s}$. We denote by $\Fnc(A,B)$ the set of all $S$-sorted functions from $A$ to $B$.  The composition of $S$-sorted functions, which is a particular case of the composition of $S$-sorted relations, is an $S$-sorted function.
%
% An $S$-\emph{sorted mapping from} $A$ \emph{to} $B$ is a triple $f = (A,F,B)$ where
% $F$ is an $S$-sorted function from $A$ to $B$. We denote by $\Hom(A,B)$ or by $B_{A}$ the set of all
% $S$-sorted mappings from $A$ to $B$. We consider the expressions $f\in\Hom(A,B)$, $f\in B_{A}$, and $f\colon A\mor B$ as synonymous. Moreover, given $f\colon A\mor B$ and $g\colon B\mor C$, $g\comp f = (A,G\comp F,C)$,
% the \emph{composition of} $f$ \emph{and} $g$, is an $S$-sorted mapping from $A$ to $C$, and $\id^{A} = (A,\Delta_{A},A)$, is an  $S$-sorted endomapping of $A$, the identity $S$-sorted mapping at $A$. We denote by
% $\mathsf{Set}^{S}$ the category of $S$-sorted sets and $S$-sorted mappings.
%
% Let $f\colon A\mor B$ be an $S$-sorted mapping. Then the mapping 
% $$
% f[\cdot]\colon\mathrm{Sub}(A)\mor\mathrm{Sub}(B),
% $$ 
% of $f$-\emph{direct image formation}, sends $X \in \mathrm{Sub}(A)$ to $f[X] = (f_{s}[X_{s}])_{s\in S} \in \mathrm{Sub}(B)$, and the mapping 
% $$
% f^{-1}[\cdot]\colon\mathrm{Sub}(B)\mor\mathrm{Sub}(A),
% $$ 
% of $f$-\emph{inverse image formation}, sends $Y \in \mathrm{Sub}(B)$ to $f^{-1}[Y] = (f_{s}^{-1}[Y_{s}] )_{s\in S} \in \mathrm{Sub}(A)$. 
%
% The \emph{image} of $f$, denoted by $\mathrm{Im}(f)$, is $f[A]$, i.e., $(f_{s}[A_{s}])_{s\in S}$. Moreover, if $X\subseteq A$, then the \emph{restriction of} $f$ \emph{to} $X$, denoted by $f{\upharpoonright}_{X}$, $f{\mid}_{X}$ or $\mathrm{res}_{X}(f)$, is $f\circ \mathrm{in}^{X,A}$, where $\mathrm{in}^{X,A} = (\mathrm{in}^{X_{s},A_{s}})_{s\in S}$ is the canonical embedding of $X$ into $A$.
% \end{definition}
%
% \begin{remark}
% To give an $S$-sorted mapping from $A$ to $B$, as in the just stated definition, is equivalent to give an $S$-indexed family $f = (f_{s})_{s\in S}$, where, for every $s$ in $S$, $f_{s}$ is a mapping from $A_{s}$ to  $B_{s}$. Thus, an $S$-sorted mapping from $A$ to $B$ is, essentially, an element of $\prod_{s\in S}\mathrm{Hom}(A_{s}, B_{s})$. 
% %We will denote by $\mathrm{Hom}(A,B)$ the set of all $S$-sorted mappings from $A$ to $B$. $S$-sorted sets and $S$-sorted mappings form a category which we will denote henceforth by $\mathsf{Set}^{S}$. %From now on, $\mathsf{Set}^{S}$ stands for the category of $S$-sorted sets and $S$-sorted mappings.
% \end{remark}
%
% \begin{proposition}\label{PDeltaPhiFunc}
% Let $S$ and $T$ be sets and $\varphi \colon S \mor T$ a mapping. We let $\Delta_{\varphi}$ stand for the assignment from $\mathsf{Set}^{T}$ to $\mathsf{Set}^{S}$ defined as follows:
% \begin{enumerate}
% \item
% for every $T$-sorted set $A = (A_{t})_{t \in T}$, $\Delta_{\varphi}(A)$ is the $S$-sorted set $A_{\varphi} = (A_{\varphi(s)})_{s \in S}$, and
% \item
% for every $T$-sorted sets $A$ and $B$, and every $T$-sorted mapping $f \colon A \mor B$, $\Delta_{\varphi}(f)$ is the $S$-sorted mapping $f_{\varphi} = (f_{\varphi(s)})_{s \in S}$.
% \end{enumerate}
% Then, $\Delta_{\varphi}$ is a covariant functor from $\mathsf{Set}^{T}$ to $\mathsf{Set}^{S}$.
% \end{proposition}
%
% \begin{proposition}\label{PMSetFunc}
% We let $\mathrm{MSet}$ stand for the assignment from $\mathsf{Set}$ to $\mathsf{Cat}$ defined as follows:
% \begin{enumerate}
% \item
% for every set $S$, $\mathrm{MSet}(S)$ is the category $\mathsf{Set}^{S}$, and
% \item
% for every mapping $\varphi \colon S \mor T$, $\mathrm{MSet}(f)$ is $\Delta_{\varphi}$ the functor from $\mathsf{Set}^{T}$ to $\mathsf{Set}^{S}$ introduced in Proposition~\ref{PDeltaPhiFunc}.
% \end{enumerate}
% Then $\mathrm{MSet}$ is a contravariant functor from $\mathsf{Set}$ to $\mathsf{Cat}$.
% \end{proposition}
%
% \begin{definition}
% Let $I$ be a set in $\boldsymbol{\mathcal{U}}$ and $(A^{i})_{i\in I}$ an $I$-indexed family of $S$-sorted sets. Then the \emph{product} of $(A^{i})_{i\in I}$, denoted by $\prod_{i\in I}A^{i}$, is the $S$-sorted set defined, for every $s\in S$, as $\left(\prod\nolimits_{i\in I}A^{i}\right)_{s} = \prod\nolimits_{i\in I}A^{i}_{s}$, where
% $$
% \textstyle
% \prod_{i\in I}A^{i}_{s} = \left\{(a_{i})_{i\in I}\in\mathrm{Fnc}\left(I,\bigcup_{i\in I}A^{i}_{s}\right)\mid \forall\,i\in I\,\left(a_{i}\in A^{i}_{s}\right)\right\}.
% $$
% For every $i\in I$, the \emph{i-th canonical projection}, $\mathrm{pr}^{i} = (\mathrm{pr}^{i}_{s})_{s\in S}$, is the $S$-sorted mapping from  $\prod_{i\in I}A^{i}$ to $A^{i}$ that, for every $s\in S$, sends $(a_{i})_{i\in I}$ in $\prod_{i\in I}A^{i}_{s}$ to $a_{i}$ in $A^{i}_{s}$. The ordered pair $(\prod_{i\in I}A^{i},(\mathrm{pr}^{i})_{i\in I})$ has the following universal property: For every $S$-sorted set $B$ and every $I$-indexed family of $S$-sorted mappings $(f^{i})_{i\in I}$, where, for every $i\in I$, $f^{i}$ is an $S$-sorted mapping from $B$ to $A^{i}$, there exists a unique $S$-sorted mapping $\left<f^{i}\right>_{i\in I}$ from $B$ to $\prod_{i\in I}A^{i}$ such that, for every $i\in I$, $\mathrm{pr}^{i}\circ \left<f^{i}\right>_{i\in I} = f^{i}$.
%
% The \emph{coproduct} of $(A^{i})_{i\in I}$, denoted by $\coprod_{i\in I}A^{i}$, is the $S$-sorted set defined, for every $s\in S$, as $\left(\coprod\nolimits_{i\in I}A^{i}\right)_{s} = \coprod\nolimits_{i\in I}A^{i}_{s}$, where
% $$
% \textstyle
% \coprod_{i\in I}A^{i}_{s} = \bigcup_{i\in I}(A^{i}_{s}\times\{i\}).
% $$
% For every $i\in I$, the \emph{i-th canonical injection}, $\mathrm{in}^{i}$, is the $S$-sorted mapping from $A^{i}$ to $\coprod_{i\in I}A^{i}$ that, for every $s\in S$, sends $a$ in $A^{i}_{s}$ to $(a,i)$ in $\coprod_{i\in I}A^{i}_{s}$. The ordered pair $(\coprod_{i\in I}A^{i},(\mathrm{in}^{i})_{i\in I})$ has the following universal property: For every $S$-sorted  set $B$ and every $I$-indexed family of $S$-sorted mappings $(f^{i})_{i\in I}$, where, for every $i\in I$, $f^{i}$ is an $S$-sorted mapping from $A^{i}$ to $B$, there exists a unique $S$-sorted mapping $[f^{i}]_{i\in I}$ from $\coprod_{i\in I}A^{i}$ to $B$ such that, for every $i\in I$, $ [f^{i}]_{i\in I}\circ\mathrm{in}^{i} = f^{i}$.
%
% The remaining set-theoretic operations on $S$-sorted sets: $\amalg$ (binary coproduct), $\bigcup$ (union), $\cup$ (binary union), $\bigcap$ (intersection), $\cap$ (binary intersection), $-$ (difference), and $\complement_{A}$ (complement of an $S$-sorted set in a fixed $S$-sorted $A$), are defined in a similar way, i.e., componentwise.
% \end{definition}
%
% \begin{definition}
% %An $S$-sorted set $A$ is \emph{subfinal} if $\mathrm{card}(A_{s})\leq 1$, for every $s\in S$.
% We will denote by $1^{S}$ the (standard) final $S$-sorted set of $\mathsf{Set}^{S}$, which is $1^{S} = (1)_{s\in S}$, and by $\varnothing^{S}$ the initial $S$-sorted set, which is $\varnothing^{S} = (\varnothing)_{s\in S}$. We shall abbreviate $1^{S}$ to $1$ and $\varnothing^{S}$ to $\varnothing$ when this is unlikely to cause confusion.
% \end{definition}
%
% %\begin{remark}
% %For every $S$-sorted set $A$, the ordered pairs $(B,f)$, where $B$ is an $S$-sorted set and $f$ a monomorphisms from $B$ to $A$, are classified by letting $(B,f) \equiv (C,g)$ if and only if there exists an isomorphism $h$ from $B$ to $C$ such that $f = g\circ h$, and the corresponding equivalence classes are called \emph{subobjects} of $A$. Then $\mathrm{Sub}(A)$ is isomorphic to the set of all subobjects $A$.
% %\end{remark}
%
% \begin{definition}
% Let $\delta$ be the mapping from $S\times \boldsymbol{\mathcal{U}}$ to $\boldsymbol{\mathcal{U}}^{S}$ that sends $(t,X)$ in $S\times \boldsymbol{\mathcal{U}}$ to the $S$-sorted set $\delta^{t,X} = (\delta^{t,X}_{s})_{s\in S}$ defined, for every $s\in S$, as follows: $\delta^{t,X}_{s} = X$, if $s = t$; $\delta^{t,X}_{s} = \varnothing$, otherwise. We will call the value of $\delta$ at $(t,X)$ the \emph{delta of Kronecker associated to} $(t,X)$. If $X = \{x\}$, then, for simplicity of notation, we will write $\delta^{t,x}$ instead of $\delta^{t,\{x\}}$. Moreover, for a sort $t$ in $S$, $\delta^{t,1}$, the delta of Kronecker associated to $(t,1)$, will be denoted by $\delta^{t}$ and called \emph{delta of Kronecker}.
%
% %Given a sort $t\in S$ we call \emph{delta of Kronecker in} $t$, the $S$-sorted set $\delta^{t} = (\delta^{t}_{s})_{s\in S}$ defined, for every $s\in S$, as follows: $\delta^{t}_{s} = 1$, if $s = t$; $\delta^{t}_{s} = \varnothing$, otherwise.
% %%\begin{equation*}
% %%\delta^{t}_{s} =
% %%\begin{cases}
% %%1, &\text{if $s = t$;}\\
% %%\varnothing, & \text{otherwise.}
% %%\end{cases}
% %%\end{equation*}
% %Let $t$ be a sort in $S$ and $X$ a set, then we will denote by $\delta^{t,X}$ the $S$-sorted set defined, for every $s\in S$, as follows: $\delta^{t,X}_{s} = X$, if $s = t$; $\delta^{t,X}_{s} = \varnothing$, otherwise.  If $X = \{x\}$, then, for simplicity of notation, we write $\delta^{t,x}$ instead of $\delta^{t,\{x\}}$.
% %%\begin{equation*}
% %%\delta^{t,X}_{s} =
% %%\begin{cases}
% %%X, &\text{if $s = t$;}\\
% %%\varnothing, & \text{otherwise.}
% %%\end{cases}
% %%\end{equation*}
% %
% %Let us note that $\delta^{t}$ is $\delta^{t,1}$, i.e., the deltas of Kronecker, are particular cases of the $S$-sorted sets $\delta^{t,X}$ (however, see the remark immediately below). Therefore we will use indistinctly  $\delta^{t}$ or $\delta^{t,1}$.
% \end{definition}
%
% \begin{remark}
% For a sort $t\in S$ and a set $X$, the $S$-sorted set $\delta^{t,X}$ is isomorphic to the $S$-sorted set $\coprod_{x\in X}\delta^{t}$, i.e., to the coproduct of the family $(\delta^{t})_{x\in X}$ which is constantly $\delta^{t}$.
%
% For every sort $t\in S$ we have a functor $\delta^{t,\cdot}$ from $\mathsf{Set}$ to $\mathsf{Set}^{S}$. In fact, for every set $X$, $\delta^{t,\cdot}(X) = \delta^{t,X}$, and, for every mapping $f\colon X\mor Y$, $\delta^{t,\cdot}(f) = \delta^{t,f}$, where, for $s\in S$, $\delta^{t,f}_{s} = \id^{\varnothing}$, if $s\neq t$, and $\delta^{t,f}_{t} = f$. Moreover, for every $t\in S$, the object mapping of the functor $\delta^{t,\cdot}$ is injective and $\delta^{t,\cdot}$ is full and faithful. Hence, for every $t\in S$, $\delta^{t,\cdot}$ is a full embedding from $\mathsf{Set}$ to $\mathsf{Set}^{S}$.
% %Note that the diagonal functor $\Delta_{S}\colon \mathsf{Set}\mor \mathsf{Set}^{S}$ is $\coprod_{t_{S}}\delta^{t,\cdot}$.
%
% The final object $1^{S}$ does not generate ($\equiv$ separate) the category $\mathsf{Set}^{S}$. However, the set $\{\,\delta^{s}\mid s\in S\,\}$, of the deltas of Kronecker, is a generating ($\equiv$ separating) set for the category $\mathsf{Set}^{S}$. Therefore, every $S$-sorted set $A$ can be represented as a coproduct of copowers of deltas of Kronecker, i.e., $A$ is naturally isomorphic to $\coprod_{s\in S}\mathrm{card}(A_{s})\boldsymbol{\cdot}\delta^{s}$, where, for every $s\in S$, $\mathrm{card}(A_{s})\boldsymbol{\cdot}\delta^{s}$ is the copower of the family $(\delta^{s})_{\alpha\in \mathrm{card}(A_{s})}$, i.e., the coproduct of $(\delta^{s})_{\alpha\in \mathrm{card}(A_{s})}$.
% To this we add the following facts: (1) $\{\,\delta^{s}\mid s\in S\,\}$ is the set of atoms of the Boolean algebra $\mathbf{Sub}(1^{S})$, of subobjects of $1^{S}$; (2) the Boolean algebras $\mathbf{Sub}(1^{S})$ and $\mathbf{Sub}(S)$ are isomorphic; (3) for every $s\in S$, $\delta^{s}$ is a projective object; and (4) for every $s\in S$, every $S$-sorted mapping from $\delta^{s}$ to another $S$-sorted set is a monomorphism.
%
% In view of the above, it must  be concluded that the deltas of Kronecker are of crucial importance for many-sorted sets and associated fields.
% \end{remark}
%
% Before proceeding any further, let us point out that it is no longer unusual to find in the works devoted to investigate both many-sorted algebras and many-sorted algebraic systems the following. (1) That an $S$-sorted set $A$ is defined in such a way that $\mathrm{Hom}(1^{S},A)\neq \varnothing$, or, what is equivalent, requiring that, for every $s\in S$, $A_{s}\neq \varnothing$. This has as an immediate consequence that the corresponding category is not even finite cocomplete. Since cocompleteness (and completeness) are desirable properties for a category, we exclude such a convention in our work (the admission of $\varnothing^{S}$ is crucial in many applications). And (2) that an $S$-sorted set $A$ must be such that, for every $s$, $t\in S$, if $s\neq t$, then $A_{s}\cap A_{t} = \varnothing$. We also exclude such a requirement (the possibility of a common underlying set for the different sorts is very important in many applications). The above conventions are possibly based on the untrue widespread belief that many-sorted equational logic and many-sorted first-order logic with equality are \emph{inessential} variations of equational logic and first-order logic with equality, respectively. One can find a definitive refutation to the just mentioned belief in~\cite{gm85} and \cite{m76}, regarding many-sorted equational logic, and in~\cite{Hook85}, with respect to many-sorted first-order logic with equality.
%
% %On this matter it seems that still is in force \emph{The Law of Conservation of Ignorance}, attributed by Kline, in~\cite{kl80}, on p.~88, to Cantor: ``A false conclusion once arrived at and widely accepted is not easily dislodged and the less it is understood the more tenaciously it is held.''
%
%
% %\begin{definition}
% %Let $f\colon A\mor B$ be an $S$-sorted mapping. Then the mapping $f[\cdot]\colon\mathrm{Sub}(A)\mor\mathrm{Sub}(B)$, of $f$-\emph{direct image formation}, sends $X \in \mathrm{Sub}(A)$ to $f[X] = (f_{s}[X_{s}])_{s\in S} \in \mathrm{Sub}(B)$, and the mapping $f^{-1}[\cdot]\colon\mathrm{Sub}(B)\mor\mathrm{Sub}(A)$, of $f$-\emph{inverse image formation}, sends $Y \in \mathrm{Sub}(B)$ to $f^{-1}[Y] = (f_{s}^{-1}[Y_{s}] )_{s\in S} \in \mathrm{Sub}(A)$. The \emph{kernel} of $f$, denoted by $\mathrm{Ker}(f)$, is $(\mathrm{Ker}(f_{s}))_{s\in S}$ and the \emph{image} of $f$, denoted by $\mathrm{Im}(f)$, is $f[A]$. Moreover, if $X\subseteq A$, then the \emph{restriction of} $f$ \emph{to} $X$, denoted by $f\!\!\upharpoonright_{X}$, $f|_{X}$ or $\mathrm{res}_{X}(f)$, is $f\circ \mathrm{in}_{X,A}$, where $\mathrm{in}_{X,A} = (\mathrm{in}_{X_{s},A_{s}})_{s\in S}$ is the canonical embedding of $X$ into $A$.
% %\begin{enumerate}
% %\item The mapping $f[\cdot]$ of $f$-\emph{direct image formation} is the mapping defined as follows:
% %$$
% %f[\cdot] \nfunction
% %{\mathrm{Sub}(A)}
% %{\mathrm{Sub}(B)}
% %{X}
% %{ f[X] = (f_{s}[X_{s}])_{s\in S}  }
% %$$
%
% %\item The mapping $f^{-1}[\cdot]$ of $f$-\emph{inverse image formation} is the mapping defined as follows:
% %$$
% %f^{-1}[\cdot] \nfunction
% %{\mathrm{Sub}(B)}
% %{\mathrm{Sub}(A)}
% %{Y}
% %{ f^{-1}[Y] = (f_{s}^{-1}[Y_{s}] )_{s\in S} }
% %$$
% %\end{enumerate}
% %\end{definition}
%
% %Before stating the following proposition, we recall that, for an $S$-sorted set $A$, the $S$-sorted set $(\mathrm{Sub}(A_{s}))_{s\in S}$, usually denoted by $A^{\wp}$, is the power object of $A$ in the topos $\mathsf{Set}^{S}$. Therefore, one should take great care not to confuse $\mathrm{Sub}(A)$, which is a \emph{set}, i.e., an object of $\mathsf{Set}$---naturally isomorphic to the set $\prod_{s\in S}2^{A_{s}}$---, and $A^{\wp}$, which is an $S$-\emph{sorted set}, i.e., an object of $\mathsf{Set}^{S}$---naturally isomorphic to the $S$-sorted set $(2^{A_{s}})_{s\in S}$. On the other hand, it is clear that there exists a natural isomorphism between $\mathrm{Sub}(A)$ and $\prod A^{\wp}$.
%
% %We next prove that every $S$-sorted mapping $f$ from an $S$-sorted set $X$ to the $S$-sorted set $(\mathrm{Sub}(Y_{s}))_{s\in S}$ canonically associated to an $S$-sorted set $Y$ has a distinguished extension up to an $S$-sorted  mapping $f^{\mathfrak{p}}$ from $(\mathrm{Sub}(X_{s}))_{s\in S}$ to $(\mathrm{Sub}(Y_{s}))_{s\in S}$ and that the canonical embedding of the category $\mathsf{Set}^{S}_{\wp,\mathrm{ca}}$, with objects the $S$-sorted sets $(\mathrm{Sub}(X_{s}))_{s\in S}$, for $X\in\boldsymbol{\mathcal{U}}^{S}$, and morphisms from $(\mathrm{Sub}(X_{s}))_{s\in S}$ to $(\mathrm{Sub}(Y_{s}))_{s\in S}$ the completely additive mappings (see below), into the category $\mathsf{Set}^{S}$ has a left adjoint.
%
% %\begin{proposition}
% %Let $X$ and $Y$ be $S$-sorted sets and $f$ an $S$-sorted mapping from $X$ to $Y^{\wp}$. Then there exists a unique $S$-sorted mapping $f^{\mathfrak{p}}$ from $(\mathrm{Sub}(X_{s}))_{s\in S}$ to $Y^{\wp}$ such that $f^{\mathfrak{p}}$ is completely additive, i.e., for every $s\in S$ and every $\mathcal{L}\subseteq \mathrm{Sub}(X_{s})$, $f^{\mathfrak{p}}_{s}(\bigcup\mathcal{L}) = \bigcup_{L\in \mathcal{L}}f^{\mathfrak{p}}_{s}(L)$, and $f^{\mathfrak{p}}\circ \{\cdot\}_{X} = f$, where $\{\cdot\}_{X}$ is the $S$-sorted mapping from $X$ to $X^{\wp}$ that, for every $s\in S$, sends $x\in X_{s}$ to $\{x\}\in\mathrm{Sub}(X_{s})$.
% %\end{proposition}
% %
% %\begin{proof}
% %The $S$-sorted mapping $f^{\mathfrak{p}}$ from $X^{\wp}$ to $Y^{\wp}$ that, for every $s\in S$, assigns to $L\subseteq X_{s}$ the set $\bigcup_{x\in L}f_{s}(x)\subseteq Y_{s}$ is completely additive and $f^{\mathfrak{p}}\circ \{\cdot\}_{X} = f$. Moreover, $f^{\mathfrak{p}}$ is clearly the unique $S$-sorted mapping from $X^{\wp}$ to $Y^{\wp}$ satisfying the aforementioned  conditions.
% %\end{proof}
%
% %\begin{corollary}
% %Let $\mathsf{Set}^{S}_{\wp,\mathrm{ca}}$ be the category whose objects are the $S$-sorted sets $X^{\wp}$, where $X\in\boldsymbol{\mathcal{U}}^{S}$, and in which the set of morphisms from $X^{\wp}$ to $Y^{\wp}$ is the set of the completely additive $S$-sorted mappings from $X^{\wp}$ to $Y^{\wp}$. Then from $\mathsf{Set}^{S}_{\wp,\mathrm{ca}}$ to $\mathsf{Set}^{S}$ we have a canonical inclusion, denoted by $\mathrm{In}_{\mathsf{Set}^{S}_{\wp,\mathrm{ca}},\mathsf{Set}^{S}}$, and, for every $S$-sorted set $X$, the ordered pair $(X^{\wp},\{\cdot\}_{X})$ is a universal morphism from $X$ to $\mathrm{In}_{\mathsf{Set}^{S}_{\wp,\mathrm{ca}},\mathsf{Set}^{S}}$.
% %\end{corollary}
%
% %\begin{definition}
% %We will denote by $(\cdot)^{\wp}$ the functor from $\mathsf{Set}^{S}$ to $\mathsf{Set}^{S}_{\wp,\mathrm{ca}}$ that sends an $S$-sorted set $X$ to $X^{\wp}$ and an $S$-sorted  mapping $f$ from $X$ to $Y$ to the $S$-sorted mapping $f^{\wp} = (\{\cdot\}_{Y}\circ f)^{\mathfrak{p}}$ from $X^{\wp}$ to $Y^{\wp}$ (the $S$-sorted mappings $\{\cdot\}_{X}$, for $X$ in $\boldsymbol{\mathcal{U}}^{S}$, are the components of the unit of the adjunction: $\mathrm{Hom}_{\mathsf{Set}^{S}_{\wp,\mathrm{ca}}}(X^{\wp},Y^{\wp})\cong
% %\mathrm{Hom}_{\mathsf{Set}^{S}}(X,Y^{\wp})$). Thus, $(\cdot)^{\wp}$ is a left adjoint of $\mathrm{In}_{\mathsf{Set}^{S}_{\wp,\mathrm{ca}},\mathsf{Set}^{S}}$.
% %\end{definition}
%
% %\begin{remark}
% %For an $S$-sorted mapping $f$ from $X$ to $Y$ the $S$-sorted mapping $f^{\wp}$ from $X^{\wp}$ to $Y^{\wp}$ sends, for every $s\in S$, $L\subseteq X_{s}$ to $f_{s}[L]\subseteq Y_{s}$, i.e., $f^{\wp} = (f_{s}[\cdot])_{s\in S}$. One should be careful not to confuse $f[\cdot]$, which is a \emph{mapping} from the \emph{set} $\mathrm{Sub}(X)$ to the \emph{set} $\mathrm{Sub}(Y)$, and $f^{\wp}$, which is an $S$-\emph{sorted mapping} from the $S$-\emph{sorted set} $X^{\wp} = (\mathrm{Sub}(X_{s}))_{s\in S}$ to the $S$-\emph{sorted set} $Y^{\wp} = (\mathrm{Sub}(Y_{s}))_{s\in S}$. On the other hand, it is evident that $f[\cdot]$ and $\prod f^{\wp}$ are \emph{essentially} the same mapping.
% %\end{remark}
%
% %Aquí propiedad universal del subconjunto, corestrictions, astrictions y coastrictions
%
% \begin{figure}
% $$\xymatrix{
%   & A \ar[dl]_{h} \ar[d]^{f} \\
% Y \ar @{{ +}{-}{>}}[r]_{\mathrm{in}^{Y,B}} & B
%  }
% $$
% \caption{The corestriction of a many-sorted mapping.}
% \label{FPUS}
% \end{figure}
%
%
% \begin{proposition}\label{PUS}
% Let $B$ be an $S$-sorted set, $Y$ a subset of $B$, and $f$ an $S$-sorted mapping from $A$ to $B$. Then the following statements are equivalent:
% \begin{enumerate}
% \item $\mathrm{Im}(f)\subseteq Y$.
% \item There exists an $S$-sorted mapping $h$ from $A$ to $Y$ such that the diagram
% in Figure~\ref{FPUS} 
% commutes.
% \end{enumerate}
% If one of the above equivalent statements holds, then we will call $h$, which is univocally determined, the \emph{corestriction of} $f$ \emph{to} $Y$ and we denote it by $f|^{Y}$ or $\mathrm{cores}_{Y}(f)$.
% \end{proposition}
%
% %\begin{remark}
% %Let $B$ be an $S$-sorted set and $Y$ a subset of $B$. Then the ordered pair $(Y,\mathrm{in}_{Y,B})$ is an equalizer of the $S$-sorted mappings $\mathrm{ch}_{Y}$ and $\kappa_{1}$ from $B$ to $2$, i.e., $\mathrm{ch}_{Y}\comp \mathrm{in}_{Y,B} = \kappa_{1}\comp \mathrm{in}_{Y,B}$ and, for every $S$-sorted set $A$ and every $S$-sorted mapping $f$ from $A$ to $B$, if $\mathrm{ch}_{Y,B}\comp f = \kappa_{1}\comp f$, then there exists a uniqye $S$-sorted mapping $h$ from $A$ to $Y$ such that $\mathrm{in}_{Y,B}\comp h = g$
% %\end{remark}
%
% \begin{proposition}
% Let $\mathsf{PSet}^{S}$ be the category whose objects are the \emph{$S$-sorted set pairs}, i.e., the ordered pairs pairs $(A,X)$ where $A$ is an $S$-sorted set and $X\subseteq A$, and in which the set of morphisms from $(A,X)$ to $(B,Y)$ is the set of all $S$-sorted mappings $f$ from $A$ to $B$ such that $f[X]\subseteq Y$. Let $G$ be the functor from $\mathsf{Set}^{S}$ to $\mathsf{PSet}^{S}$ whose object mapping sends $A$ to 
% $(A,A)$ and whose morphism mapping sends $f\colon A\mor B$ to $f\colon (A,A)\mor (B,B)$. Then, for every $S$-sorted pair $(B,Y)$, there exists a universal mapping from $G$ to $(B,Y)$, which is precisely the ordered pair $(Y,\mathrm{in}^{Y,B})$ with $\mathrm{in}^{Y,B}\colon (Y,Y)\mor (B,Y)$ the morphism in $\mathsf{PSet}^{S}$ associated to $\mathrm{in}^{Y,B}\colon Y\mor B$.
% \end{proposition}
%
% \begin{definition}
% Let $f$ be a $\mathsf{PSet}^{S}$-morphism from $(A,X)$ to $(B,Y)$. Then we denote by $f|_{X}^{Y}$, $\mathrm{bires}_{X,Y}(f)$, or, if no confusion can arise, $\widehat{f}$ the $S$-sorted mapping $\mathrm{cores}_{Y}(\mathrm{res}_{X}(f))$ (which is identical to $\mathrm{res}_{X}(\mathrm{cores}_{Y}(f))$). We will call this $S$-sorted mapping the \emph{birestriction of} $f$ to $X$ and $Y$. 
% \end{definition}
%
% \begin{definition}
% Let $A$ be an $S$-sorted set. Then the \emph{cardinal} of $A$, denoted by $\mathrm{card}(A)$, is $\mathrm{card}(\coprod A)$, i.e., the cardinal of the set $\coprod A = \bigcup_{s\in S}(A_{s}\times \{s\})$. An $S$-sorted set $A$ is \emph{finite} if $\mathrm{card}(A)<\aleph_{0}$. We will say that an $S$-sorted set $X$ is a \emph{finite} subset of $A$ if $X$ is finite and $X\subseteq A$. We will denote by $\mathrm{Sub}_{\mathrm{f}}(A)$ or, sometimes, by $\mathrm{Fin}(A)$ the set of all $S$-sorted sets $X$ in $\mathrm{Sub}(A)$ which are finite and, on some occasions, we will write $X\subseteq_{\mathrm{f}}A$ to indicate that $X\in \mathrm{Sub}_{\mathrm{f}}(A)$.
% \end{definition}
%
% \begin{remark}
% For an object $A$ in the topos $\mathsf{Set}^{S}$ the following assertions are equivalent: (1) $A$ is finite, (2) $A$ is a finitary object of $\mathsf{Set}^{S}$, and (3) $A$ is a strongly finitary object of $\mathsf{Set}^{S}$ (for the notions of finitary and strongly finitary object of a category see~\cite{hs73}, Exercise 22E, on p.~155).
% %An $S$-sorted set $A$ is finite if and only if the covariant hom-functor $\mathrm{H}(A,\cdot)\colon \mathsf{Set}^{S}\mor \mathsf{Set}$ is finitary, i.e., if and only if for every $\boldsymbol{\mathcal{U}}$-small upward-directed preordered set $\mathbf{I}$ and every functor $D$ from $\mathbf{C}(\mathbf{\mathbf{I}})$, the category canonically associated to $\mathbf{I}$, to $\mathsf{Set}^{S}$, if $((f^{i})_{i\in \mathrm{Ob}(\mathbf{C}(\mathbf{\mathbf{I}}))},L)$ is an inductive limit of $D$, then $((\mathrm{H}(A,f^{i}))_{i\in \mathrm{Ob}(\mathbf{C}(\mathbf{\mathbf{I}}))},\mathrm{H}(A,L))$ is an epi-sink. Moreover, the $S$-sorted set $A$ is finite if and only if the functor $\mathrm{H}(A,\cdot)$ is strongly finitary, i.e., if and only if for every $\boldsymbol{\mathcal{U}}$-small upward-directed preordered set $\mathbf{I}$ and every functor $D$ from $\mathbf{C}(\mathbf{\mathbf{I}})$ to $\mathsf{Set}^{S}$, if $((f^{i})_{i\in \mathrm{Ob}(\mathbf{C}(\mathbf{\mathbf{I}}))},L)$ is an inductive limit of $D$, then $((\mathrm{H}(A,f^{i}))_{i\in \mathrm{Ob}(\mathbf{C}(\mathbf{\mathbf{I}}))},\mathrm{H}(A,L))$ is an inductive limit.
%
% In $\mathsf{Set}^{S}$ there is another notion of finiteness: An $S$-sorted set $A$ is called $S$-\emph{finite} or \emph{locally finite}, abbreviated as $S\text{-}\mathrm{f}$, if and only if, for every $s\in S$, $A_{s}$ is finite.  We will denote by $\mathrm{Sub}_{S\text{-}\mathrm{f}}(A)$ the set of all $S$-sorted sets $X$ in $\mathrm{Sub}(A)$ which are $S$-finite and, on some occasions, we will write 
% $X\subseteq_{S\text{-}\mathrm{f}}A$ to indicate that $X\in \mathrm{Sub}_{S\text{-}\mathrm{f}}(A)$. Although, unless $S$ is finite, this notion of finiteness is not categorial in nature, however, it plays a relevant role in the field of many-sorted algebra and in computer science. 
% %(see below, after the definition of many-sorted algebra and when we deal with the congruences of locally finite index, respectively).
% \end{remark}
%
% \begin{remark}\label{RDelta}
% Let $\mathrm{Fin}^{S}$ be the functor from $\mathsf{Set}^{S}$ to $\mathsf{Set}^{S}$ that sends an $S$-sorted set $A$ to the $S$-sorted set $(\mathrm{Fin}(A))_{s\in S}$, which is constantly $\mathrm{Fin}(A)$, and an $S$-sorted mapping $f\colon A\mor B$ to the $S$-sorted mapping $\mathrm{Fin}^{S}(f)$ from $(\mathrm{Fin}(A))_{s\in S}$ to $(\mathrm{Fin}(B))_{s\in S}$ that, for every $s\in S$, sends $X\in \mathrm{Fin}(A)$ to $f[X]$ in $\mathrm{Fin}(B)$, i.e., to the $S$-sorted mapping $(f[\cdot])_{s\in S}$, which is constantly $f[\cdot]$. Then the family $(\delta^{S,A})_{A\in \boldsymbol{\mathcal{U}}}$, where, for every $A\in \boldsymbol{\mathcal{U}}$,  $\delta^{S,A}$ is the $S$-sorted mapping from $A$ to $(\mathrm{Fin}(A))_{s\in S}$ that, for every $s\in S$,  sends $a\in A_{s}$ to $\delta^{S,A}_{s}(a) = \delta^{s,a}\in \mathrm{Fin}(A)$, is a natural transformation from $\Id^{\mathsf{Set}^{S}}$ to $\mathrm{Fin}^{S}$. This natural transformation will be used later on to define, for an $s\in S$ and an $S$-sorted set $X$, the $S$-sorted set of the variables of an $(X,s)$-term.
% \end{remark}
%
% \begin{definition}
% Let $A$ be an $S$-sorted set. Then the \emph{support of} $A$, denoted by $\mathrm{supp}_{S}(A)$, is the set $\{s\in S\mid A_{s}\neq \varnothing\}$.
% %(i.e., the subset of $S$ on which $A$ does something nontrivial).
% \end{definition}
%
% \begin{remark}
% An $S$-sorted set $A$ is finite if and only if $\mathrm{supp}_{S}(A)$ is finite and, for every $s\in \mathrm{supp}_{S}(A)$, $A_{s}$ is finite.
% \end{remark}
%
% %In the following proposition we gather together the most interesting properties of the mapping   $\mathrm{supp}_{S}\colon\boldsymbol{\mathcal{U}}^{S}\mor \mathrm{Sub}(S)$, the support mapping for $S$, which sends an $S$-sorted set $A$ to $\mathrm{supp}_{S}(A)$.
% %
% %\begin{proposition}\label{propssupport}
% %Let $A$ and $B$ be two $S$-sorted sets, $I$ a set in $\boldsymbol{\mathcal{U}}$, and $(A^{i})_{i\in I}$ an $I$-indexed family of $S$-sorted sets. Then the following properties hold:
% %\begin{enumerate}
% %\item $\mathrm{Hom}(A,B)\neq \varnothing$ if and only if
% %      $\mathrm{supp}_{S}(A)\subseteq\mathrm{supp}_{S}(B)$. Therefore, if $A\subseteq B$, then $\mathrm{supp}_{S}(A)\subseteq\mathrm{supp}_{S}(B)$. Moreover, if $f$ is an $S$-sorted mapping from $A$ to $B$ and $X\subseteq A$, then $\mathrm{supp}_{S}(X) = \mathrm{supp}_{S}(f[X])$.
% %
% %\item If from $A$ to $B$ there exists a surjective $S$-sorted mapping $f$, then we have 
% %      $\mathrm{supp}_{S}(A) = \mathrm{supp}_{S}(B)$. Moreover, if $Y\subseteq B$, then it happens that $\mathrm{supp}_{S}(Y) = \mathrm{supp}_{S}(f^{-1}[Y])$.
% %
% %\item $\mathrm{supp}_{S}(\varnothing^{S}) = \varnothing$; $\mathrm{supp}_{S}(1^{S}) = S$; $\mathrm{supp}_{S}(\bigcup_{i\in I} A^{i}) = \bigcup_{i\in I}\mathrm{supp}_{S}(A^{i})$; $\mathrm{supp}_{S}(\coprod_{i\in I} A^{i}) = \bigcup_{i\in I}\mathrm{supp}_{S}(A^{i})$; $\mathrm{supp}_{S}(\prod_{i\in I}A^{i}) = \bigcap\nolimits_{i\in I}\mathrm{supp}_{S}(A_{i})$; if $I\neq \varnothing$, $\mathrm{supp}_{S}(\bigcap_{i\in I} A^{i}) \subseteq \bigcap_{i\in I}\mathrm{supp}_{S}(A^{i})$; and $\mathrm{supp}_{S}(A)-\mathrm{supp}_{S}(B)\subseteq\mathrm{supp}_{S}(A-B)$.
% %\end{enumerate}
% %\end{proposition}
%
% %We next define the concept of relation on a many-sorted set, the operation of formation of the inverse of a relation, and the operation of composition of relations. Thereafter we define the notion of equivalence relation on an many-sorted set and state the universal property of the corresponding quotient many-sorted set.
% %
% %\begin{definition}\label{DRel} 
% %Let $A$ be an $S$-sorted set. An $S$-\emph{sorted relation on} an $A$ is a subset $\Phi = (\Phi_{s})_{s\in S}$ of $A\times A = (A_{s}\times A_{s})_{s\in S}$. 
% %
% %Let $\Phi$ an $S$-sorted relation on $A$. Then the \emph{inverse} of $\Phi$, denoted by $\Phi^{-1}$, is the $S$-sorted relation $\Phi^{-1}=(\Phi^{-1}_{s})_{s\in S}$ where, for every $s\in S$,
% %$$
% %\Phi_{s}^{-1}
% %=\{(y,x)\in A_{s}\times A_{s}\mid (x,y)\in \Phi_{s}\}.
% %$$
% %
% %Let $\Phi$ and $\Psi$ be $S$-sorted relations on $A$. Then the \emph{composite} of $\Phi$ and $\Psi$, denoted by $\Psi\circ\Phi$, is the $S$-sorted relation $\Psi\circ\Phi = (\Psi_{s}\circ\Phi_{s})_{s\in S}$ where, for every $s\in S$,
% %$$
% %\Psi_{s}\circ\Phi_{s}
% %=\{(x,z)\in A_{s}\times A_{s}\mid \exists\,y\in A_{s}\,( (x,y)\in \Phi_{s} \And (y,z)\in \Psi_{s} )\}.
% %$$
% %This composition is associative and the diagonal relation $\Delta_{A}$ is a neutral element for it.
% %%In particular $\Delta_{A}\circ\Delta_{A}=\Delta_{A}$.
% %\end{definition}
%
% %%%%%%%%%%%%%%%%%%%%%%
%
% %\begin{definition}
% %Let $S$ be a set of sorts and $A$ an $S$-sorted set.
% %\begin{enumerate}
% %\item An $S$-sorted relation $\Phi$ on $A$ is an
% %$S$-\emph{sorted equivalence on} $A$, or, to simplify the
% %expression, an $S$-\emph{equivalence on} $A$, if, for every $s\in
% %S$, $\Phi_{s}$ is an equivalence relation on $A_{s}$. We will
% %consider, for $s\in S$ and $(a,b)\in A_{s}\times A_{s}$, as
% %synonymous $(a,b)\in\Phi_{s}$, $a\equiv b \pmod{\Phi_{s}}$, and
% %$a\equiv_{\Phi_{s}} b$. We denote by $\mathrm{Eqv}(A)$ the set of
% %all $S$-equivalences o $A$, and by $\mathbf{Eqv}(A)$ the ordered
% %set $(\mathrm{Eqv}(A),\subseteq)$. As for ordinary sets,
% %$\mathbf{Eqv}(A)$ is also an algebraic lattice, and we denote by
% %$\mathrm{Eg}_{A}$ the canonically associated algebraic closure
% %operator.  For $A$, we have that
% %$\mathrm{Eg}_{A}(\Phi)=(\mathrm{Eg}_{A_{s}}(\Phi_{s}))_{s\in S}$.
% %
% %\item An $S$-sorted relation $\Phi$ on $A$ is an $S$-sorted equivalence on $A$ of \emph{finite index} if the $S$-sorted set $A/\Phi$ is finite.
% %
% %\item An $S$-sorted relation $\Phi$ on $A$ is an $S$-sorted equivalence on $A$ of $S$-\emph{finite index} if the $S$-sorted set $A/\Phi$ is $S$-finite.
% %
% %\item Let $\Phi, \Psi\in\mathrm{Eqv}(A)$ be such that $\Phi\subseteq\Psi$.
% %      Then the \emph{quotient of} $\Psi$ \emph{by} $\Phi$, denoted by $\Psi/\Phi$,
% %      is the $S$-equivalence $(\Psi_{s}/\Phi_{s})_{s\in
% %      S}$ on $A/\Phi$ defined, for every $s\in S$, as:
% %      $$
% %      \Psi_{s}/\Phi_{s}=\{([a]_{\Phi_{s}},
% %      [b]_{\Phi_{s}})\in (A_{s}/\Phi_{s})\mid (a,b)\in\Psi_{s}\}
% %      $$
% %
% %\item Let $X\subseteq A$ be and $\Phi\in\mathrm{Eqv}(A)$. Then the
% %      $\Phi$-\emph{saturation of} $X$, denoted by
% %      $[X]^{\Phi}$, is the
% %      $S$-sorted set defined, for every $s\in S$, as:
% %      $$
% %      \textstyle[X]^{\Phi}_{s} = \{a\in A_{s}\mid
% %      X_{s}\cap [a]_{\Phi_{s}}\neq\vacio\}=
% %      \bigcup_{x\in X_{s}}[x]_{\Phi_{s}}.
% %      $$
% %\item Let $X\subseteq A$ be and $\Phi\in\mathrm{Eqv}(A)$. Then we say that $X$ is
% %      $\Phi$-\emph{saturated} if, and only if, $X = [X]^{\Phi}$.
% %\end{enumerate}
% %\end{definition}
% %
% %\begin{remark}
% %Let $S$ be a set of sorts and $A$ an $S$-sorted set. Then $A$ is subfinal if, and only if, $\Delta_{A} = \nabla_{A}$.
% %\end{remark}
% %
% %\begin{remark}
% %Let $S$ be a set of sorts, $A$ an $S$-sorted set, and $\Phi\in\mathrm{Eqv}(A)$. Then, for an $S$-sorted subset $X$ of $A$, we have that the $\Phi$-saturation of $X$ is $(\mathrm{pr}^{\Phi})^{-1}[\mathrm{pr}^{\Phi}[X]]$ (which always contains to $X$). Therefore, $X$ is $\Phi$-saturated if, and only if, $X \supseteq [X]^{\Phi}$.
% %\end{remark}
%
% %%%%%%%%%%%%%%%%%%%%%%
%
% %\begin{definition}
% %Let $f$ be an $S$-sorted mapping from $A$ to $B$ and $\Psi$ an $S$-sorted equivalence on $B$. Then will call 
% %$\mathrm{pr}^{\Psi}\circ f$ the \emph{astriction of}
% %$f$ \emph{to} $B/\Psi$ and we denote it by $\mathrm{ast}_{B/\Psi}(f)$.
% %\end{definition}
%
% We next define the notion of equivalence relation on a many-sorted set and state the universal property of the corresponding quotient many-sorted set.
%
% \begin{definition}\label{DEqvRel}
% An $S$-\emph{sorted equivalence relation on} an $S$-sorted set $A$ is an $S$-sorted relation $\Phi$ on $A$ such that, for every $s\in S$, $\Phi_{s}$ is an equivalence relation on $A_{s}$. Sometimes, we will abbreviate the expression ``$S$-sorted equivalence relation on $A$'' to ``$S$-sorted equivalence on $A$'', or, if no confusion can arise, to ``equivalence on $A$''. We will denote by $\mathrm{Eqv}(A)$ the set of all $S$-sorted equivalences on $A$ (which is an algebraic closure system on $A\times A$), by $\mathbf{Eqv}(A)$ the algebraic lattice  $(\mathrm{Eqv}(A),\subseteq)$, by $\nabla_{A}$ the greatest element of $\mathbf{Eqv}(A)$, and by $\Delta_{A}$ the least element of $\mathbf{Eqv}(A)$. As for ordinary sets,
% $\mathbf{Eqv}(A)$ is also an algebraic lattice, and we denote by
% $\mathrm{Eg}_{A}$ the canonically associated algebraic closure
% operator.  For $A$, we have that
% $\mathrm{Eg}_{A}(\Phi)=(\mathrm{Eg}_{A_{s}}(\Phi_{s}))_{s\in S}$.
%
% For an $S$-sorted equivalence relation $\Phi$ on $A$, the $S$-\emph{sorted quotient set of} $A$ \emph{by} $\Phi$, denoted by  $A/\Phi$, is $(A_{s}/\Phi_{s})_{s\in S} = (\{[x]_{\Phi_{s}}\mid x\in A_{s}\})_{s\in S} (\subseteq (\mathrm{Sub}(A_{s}))_{s\in S})$, where, for every $s\in S$ and every $x\in A_{s}$, $[x]_{\Phi_{s}}$, the \emph{equivalence class of} $x$ \emph{with respect to} $\Phi_{s}$ (or, the $\Phi$-\emph{equivalence class of} $x$) is $\{y\in A_{s}\mid (x,y)\in \Phi_{s}\}$, and $\mathrm{pr}^{\Phi}\colon A\mor A/\Phi$, the \emph{canonical projection from} $A$ \emph{to} $A/\Phi$, is the $S$-sorted mapping $(\mathrm{pr}^{\Phi}_{s})_{s\in S}$, where, for every $s\in S$, $\mathrm{pr}^{\Phi}_{s}$ is the canonical projection from $A_{s}$ to $A_{s}/\Phi_{s}$ (which sends $x$ in $A_{s}$ to $\mathrm{pr}^{\Phi}_{s}(x) = [x]_{\Phi_{s}}$, the $\Phi_{s}$-equivalence class of $x$, in $A_{s}/\Phi_{s}$).  Moreover, if $\Psi$ is an $S$-sorted equivalence on an $S$-sorted set $B$ and $f$ an $S$-sorted mapping from $A$ to $B$, then the \emph{astriction of} $f$ \emph{to} $B/\Psi$, denoted by $\mathrm{ast}_{B/\Psi}(f)$, is $\mathrm{pr}^{\Psi}\circ f$, where $\mathrm{pr}^{\Psi}$ is the canonical projection of $B$ onto $B/\Psi$. Sometimes, we will abbreviate the expression ``$S$-sorted sorted quotient set of $A$ by $\Phi$'' to ``$S$-sorted quotient of $A$ by $\Phi$'', or, if no confusion can arise, to ``quotient of $A$ by $\Phi$''
%
% %The ordered pair $(A/\Phi,\mathrm{pr}^{\Phi})$ has the following universal property: $\mathrm{Ker}(\mathrm{pr}^{\Phi})$ is $\Phi$ and, for every $S$-sorted set $B$ and every $S$-sorted mapping $f$ from $A$ to $B$, if $\mathrm{Ker}(f)\supseteq \Phi$, then there exists a unique $S$-sorted mapping $h$ from $A/\Phi$ to $B$ such that $h\circ\mathrm{pr}^{\Phi} = f$. We will call $h$ the \emph{coastriction of}
% %$f$ \emph{to} $A/\Phi$ and we denote it by $\mathrm{coast}_{A/\Phi}(f)$. In particular, if $\Psi$ is an  $S$-sorted equivalence relation on $A$ such that $\Phi\subseteq \Psi$, then we will denote by $\mathrm{p}^{\Phi,\Psi}$ the unique $S$-sorted mapping from $A/\Phi$ to $A/\Psi$ such that $\mathrm{p}^{\Phi,\Psi}\circ \mathrm{pr}^{\Phi} = \mathrm{pr}^\Psi$.
%
% %Let $\Phi$ and $\Psi$ be $S$-sorted equivalence on $A$ such that $\Phi\subseteq\Psi$. Then the \emph{quotient of} $\Psi$ \emph{by} $\Phi$, denoted by $\Psi/\Phi$, is the $S$-sorted equivalence $(\Psi_{s}/\Phi_{s})_{s\in S}$ on $A/\Phi$ defined, for every $s\in S$, as follows:
% %$$
% %\Psi_{s}/\Phi_{s}=\{([a]_{\Phi_{s}},
% %[b]_{\Phi_{s}})\in (A_{s}/\Phi_{s})\mid (a,b)\in\Psi_{s}\}.
% %$$
% \end{definition}
%
% %\begin{remark}
% %Let $A$ be an $S$-sorted set and $\Phi\in\mathrm{Eqv}(A)$. Then, by Proposition~\ref{propssupport}, $\mathrm{supp}_{S}(A) = \mathrm{supp}_{S}(A/\Phi)$.
% %\end{remark}
% \begin{figure}
% $$\xymatrix{
%   A \ar @{{+>}}[r]^{\mathrm{pr}^{\Phi}}\ar[dr]_{f} & A/\Phi \ar[d]^{h}\\
%  	                                              & B
%  }
% $$
% \caption{The coastriction of a many-sorted mapping.}
% \label{FCOAS}
% \end{figure}
%
% \begin{definition}
% Let $A$ be an $S$-sorted set and $\Phi\in \mathrm{Eqv}(A)$. Then a \emph{transversal of} $A/\Phi$ \emph{in} $A$ is a subset $X$ of $A$ such that, for every $s\in S$ and every $a\in A_{s}$, $\mathrm{card}(X_{s}\cap [a]_{\Phi_{s}}) = 1$.
% \end{definition}
%
% \begin{remark}
% For an $S$-sorted equivalence relation $\Phi$ on $A$, the set of all transversals of $A/\Phi$ in $A$ is isomorphic to the set of all cross-sections of $\mathrm{pr}_{\Phi}$, where an $S$-sorted mapping $f$ from $A/\Phi$ to $A$ is a cross-section of $\mathrm{pr}^{\Phi}$ if $\mathrm{pr}^{\Phi}\circ f = \id^{A/\Phi}$. Moreover, if $\Psi$ is another equivalence relation on $A$, $\Psi$ is a refinement of $\Phi$, i.e., $\Psi\subseteq \Phi$, and $X^{\Phi}$ is a transversal of $A/\Phi$ in $A$, then, for every $s\in S$ and every $a\in A_{s}$, there exists a unique $x\in X^{\Phi}_{s}$ such that $[a]_{\Psi_{s}}\subseteq [x]_{\Phi_{s}}$.
% \end{remark}
%
% We next define the concept of kernel of an $S$-sorted mapping.
%
% \begin{definition}
% let $f\colon A\mor B$ be an $S$-sorted mapping. Then the \emph{kernel of} $f$, denoted by $\mathrm{Ker}(f)$, is the $S$-equivalence on $A$ defined as:
% $$
% \mathrm{Ker}(f)=(\mathrm{Ker}(f_{s}))_{s\in S}.
% $$
% \end{definition}
%
% \begin{proposition}
% Let $A$ be an $S$-sorted set, $\Phi$ an $S$-sorted equivalence on $A$, and $f\colon A\mor B$ an $S$-sorted mapping. Then the following statements are equivalent:
% \begin{enumerate}
% \item $\Phi\subseteq\mathrm{Ker}(f)$.
% \item There exists an $S$-sorted mapping $h$ from $A/\Phi$ to $B$ such that the  diagram in Figure~\ref{FCOAS} 
% commutes.
% \end{enumerate}
% If one of the above equivalent statements holds, then we will call $h$, which is univocally determined, the \emph{coastriction of} $f$ \emph{to} $A/\Phi$ and we denote it by $\mathrm{coast}_{A/\Phi}(f)$.
% \end{proposition}
%
% \begin{proposition}
% Let $\mathsf{ClfdSet}^{S}$ be the category whose objects are the \emph{classified $S$-sorted sets}, i.e, the ordered pairs $(A,\Phi)$ where $A$ is an $S$-sorted set and $\Phi$ an $S$-sorted equivalence relation on $A$, and in which the set of morphisms from $(A,\Phi)$ to $(B,\Psi)$ is the set of all $S$-sorted mappings $f$ from $A$ to $B$ such that, for every $s\in S$ and every $(x,y)\in A^{2}_{s}$, if $(x,y)\in \Phi_{s}$, then $(f_{s}(x),f_{s}(y))\in \Psi_{s}$. Let $G$ be the functor from $\mathsf{Set}^{S}$ to $\mathsf{ClfdSet}^{S}$ whose object mapping sends $A$ to $(A,\Delta_{A})$ and whose morphism mapping sends $f\colon A\mor B$ to $f\colon (A,\Delta_{A})\mor (B,\Delta_{B})$. Then, for every classified $S$-sorted set $(A,\Phi)$, there exists a universal mapping from $(A,\Phi)$ to $G$, which is precisely the ordered pair $(A/\Phi,\mathrm{pr}^{\Phi})$ with $\mathrm{pr}^{\Phi}\colon (A,\Phi)\mor (A/\Phi,\Delta_{A/\Phi})$.
% \end{proposition}
%
% \begin{definition}
% Let $f$ be a $\mathsf{ClfdSet}^{S}$-morphism from $(A,\Phi)$ to $(B,\Psi)$. Then we denote by $\mathrm{biast}_{A/\Phi,B/\Psi}(f)$, $f^{\Phi,\Psi}$, or, if no confusion can arise, $\overline{f}$ the $S$-sorted mapping $\mathrm{ast}_{B/\Psi}(\mathrm{coast}_{A/\Phi}(f))$ (which is identical to $\mathrm{coast}_{A/\Phi}(\mathrm{ast}_{B/\Psi}(f))$). We will call this $S$-sorted mapping the \emph{biastriction of} $f$ to $A/\Phi$ and $B/\Psi$. 
% \end{definition}
%
% We next state the quadrangular and triangular factorizations of an $S$-sorted mapping.
%
% \begin{proposition}
% Let $f\colon A\mor B$ be an $S$-sorted mapping. Then we have the following quadrangular and triangular factorizations of $f$:
% $$\xymatrix{
% A
% \ar[r]^-{\mathrm{pr}^{\Ker(f)}}
% \ar[d]_{f^{\mathrm{e}}} &
% A/\Ker(f)
% \ar[d]^{f^{\mathrm{m}}}
% \ar[ld]_{f^{\mathrm{b}}}\\
% \mathrm{Im}(f)
% \ar[r]_{\mathrm{in}^{\mathrm{Im}(f)}} &
% B
% }
% $$
% where the involved $S$-sorted mappings are defined coordinatewise, i.e., for every $s\in S$, $(\mathrm{pr}^{\mathrm{Ker}(f)})_{s}$ is $\mathrm{pr}^{\mathrm{Ker}(f_{s})}$, the canonical projection from $A_{s}$ to $A_{s}/\Ker(f_{s})$, $f^{\mathrm{b}}_{s}$ the canonical isomorphism from $A_{s}/\mathrm{Ker}(f_{s})$ to $\mathrm{Im}(f_{s})$, $(\mathrm{in}^{\mathrm{Im}(f)})_{s}$ the
% canonical embedding $\mathrm{in}^{\mathrm{Im}(f_{s})}$ of $\mathrm{Im}(f_{s})$ into $B_{s}$, $f^{\mathrm{e}}_{s}$ the corestriction of $f_{s}$ to $\mathrm{Im}(f_{s})$, and $f^{\mathrm{m}}_{s}$ the mapping which sends $[a]_{\Ker(f_{s})}$ in $A_{s}/\Ker(f_{s})$ to $f_{s}(a)$ in $B_{s}$.
% \end{proposition}
%
% %\begin{definition}
% %Let $A$ be an $S$-sorted set, $X\subseteq A$, and $\Phi\in\mathrm{Eqv}(A)$. Then the $\Phi$-\emph{saturation of} $X$, denoted by $[X]^{\Phi}$, is the $S$-sorted set defined, for every $s\in S$, as:
% %$$
% %\textstyle[X]^{\Phi}_{s} = \{a\in A_{s}\mid
% %X_{s}\cap [a]_{\Phi_{s}}\neq\vacio\}=
% %\bigcup_{x\in X_{s}}[x]_{\Phi_{s}}.
% %$$
% %We will say that $X$ is $\Phi$-\emph{saturated} if, and only if, $X = [X]^{\Phi}$.
% %\end{definition}
% %
% %\begin{remark}
% %Let $S$ be a set of sorts, $A$ an $S$-sorted set, and $\Phi\in\mathrm{Eqv}(A)$. Then, for an $S$-sorted subset $X$ of $A$, we have that the $\Phi$-saturation of $X$ is $(\mathrm{pr}^{\Phi})^{-1}[\mathrm{pr}^{\Phi}[X]]$ (which always contains to $X$). Therefore, $X$ is $\Phi$-saturated if, and only if, $X \supseteq [X]^{\Phi}$.
% %\end{remark}
%
% \begin{definition}
% Let $A$ be an $S$-sorted set, $X$ a subset of $A$, and $\Phi\in\mathrm{Eqv}(A)$. Then the $\Phi$-\emph{saturation of} $X$ (or, the \emph{saturation of} $X$ \emph{with respect to} $\Phi$), denoted by $[X]^{\Phi}$, is the $S$-sorted set defined, for every $s\in S$, as follows:
% $$
% \textstyle[X]^{\Phi}_{s} = \{a\in A_{s}\mid
% X_{s}\cap [a]_{\Phi_{s}}\neq\varnothing\}= \bigcup_{x\in X_{s}}[x]_{\Phi_{s}} = [X_{s}]^{\Phi_{s}}.
% $$
% Let $X$ be a subset of $A$ and $\Phi\in\mathrm{Eqv}(A)$. Then we will say that $X$ is $\Phi$-\emph{saturated} if and only if $X = [X]^{\Phi}$. We will denote by $\Phi\text{-}\mathrm{Sat}(A)$ the subset of $\mathrm{Sub}(A)$ defined as $\Phi\text{-}\mathrm{Sat}(A) = \{X\in \mathrm{Sub}(A)\mid X = [X]^{\Phi}\}$.
% \end{definition}
%
% \begin{remark}
% Let $A$ be an $S$-sorted set and $\Phi\in\mathrm{Eqv}(A)$. Then, for a subset $X$ of
% $A$, we have that the $\Phi$-saturation of $X$ is $(\mathrm{pr}^{\Phi})^{-1}[\mathrm{pr}^{\Phi}[X]]$. Therefore, $X$ is $\Phi$-saturated if and only if $X \supseteq [X]^{\Phi}$. Besides, $X$ is $\Phi$-saturated if and only if there exists a $\mathcal{Y}\subseteq A/\Phi$ such that $X = (\mathrm{pr}^{\Phi})^{-1}[\mathcal{Y}]$.
% \end{remark}
%
% \begin{proposition}\label{PSatIncl}
% Let $A$ be an $S$-sorted set and $\Phi$, $\Psi\in\mathrm{Eqv}(A)$. Then
% $$
% \Phi\subseteq \Psi\, \text{if and only if }\, \forall X\subseteq A\;([[X]^{\Psi}]^{\Phi}=[X]^{\Psi}).
% $$
% Moreover, for a sort $s\in S$, we have that
% $$
% \Phi_{s}\subseteq \Psi_{s}\, \text{if and only if }\, \forall X\subseteq A_{s}\;([[X]^{\Psi_{s}}]^{\Phi_{s}}=[X]^{\Psi_{s}}).
% $$
% \end{proposition}
%
% \begin{proof}
% We restrict ourselves to proving the first assertion. Let us assume that $\Phi\subseteq \Psi$ and let $X$ be a subset of $A$. In order to prove that $[[X]^{\Psi}]^{\Phi}=[X]^{\Psi}$ it suffices to verify that $[[X]^{\Psi}]^{\Phi}\subseteq [X]^{\Psi}$. Let $s$ be an element of $S$. Then, by definition, $a\in [[X]^{\Psi}]^{\Phi}_{s}$ if and only if there exists some $b\in [X]^{\Psi}_{s}$ such that $a\in [b]_{\Phi_{s}}$. Since $\Phi\subseteq \Psi$, we have that $a\in [b]_{\Psi_{s}}$, therefore $a\in [X]^{\Psi}_{s}$.
%
% To prove the converse, let us assume that $\Phi\not\subseteq \Psi$. Then there exists some sort $s\in S$ and elements $a$, $b$ in $A_{s}$ such that $(a,b)\in\Phi_{s}$ and $(a,b)\not\in \Psi_{s}$. Hence $b$ does not belong to $[\delta^{s,[a]_{\Psi_{s}}}]^{\Psi}_s$, whereas it does belong to $[[\delta^{s,[a]_{\Psi_{s}}}]^{\Psi}]^{\Phi}_{s}$. It follows that $[\delta^{s,[a]_{\Psi_{s}}}]^{\Psi}\neq [[\delta^{s,[a]_{\Psi_{s}}}]^{\Psi}]^{\Phi}$.
% \end{proof}
%
% \begin{corollary}\label{CSatIncl}
% Let $A$ be an $S$-sorted set and $\Phi$, $\Psi\in\mathrm{Eqv}(A)$. If $\Phi\subseteq \Psi$, then  $\Psi\text{-}\mathrm{Sat}(A)\subseteq\Phi\text{-}\mathrm{Sat}(A)$. Moreover, for $s\in S$ and $L\subseteq A_{s}$, if $\Phi_{s}\subseteq \Psi_{s}$ and $L = [L]^{\Psi_{s}}$, then  $L = [L]^{\Phi_{s}}$.
% \end{corollary}
%
% \begin{remark}
% If, for an $S$-sorted set $A$, we denote by $(\cdot)\text{-}\mathrm{Sat}(A)$ the mapping from $\mathrm{Eqv}(A)$ to $\mathrm{Sub}(\mathrm{Sub}(A))$ which sends $\Phi$ to $\Phi\text{-}\mathrm{Sat}(A)$, then the above corollary means that $(\cdot)\text{-}\mathrm{Sat}(A)$ is an antitone ($\equiv$ order-reversing) mapping from the ordered set $(\mathrm{Eqv}(A),\subseteq)$ to the ordered set $(\mathrm{Sub}(\mathrm{Sub}(A)),\subseteq)$.
% \end{remark}
%
% \begin{proposition}\label{PSatNabla}
% Let $A$ be an $S$-sorted set and $X\subseteq A$. Then $X\in \nabla_{A}\text{-}\mathrm{Sat}(A)$ if and only if, for every $s\in S$, if $s\in \mathrm{supp}_{S}(X)$, then $X_{s} = A_{s}$.
% \end{proposition}
%
% \begin{proof}
% Let us suppose that there exists a $t\in S$ such that $X_{t}\neq \varnothing$ and $X_{t}\neq A_{t}$. Then, since $[X]_{t}^{\nabla_{A}} = \bigcup_{x\in X_{t}}[x]_{\nabla_{A_{t}}}$ and $X_{t}\neq \varnothing$, we have that, for some $y\in X_{t}$, $[y]_{\nabla_{A_{t}}} = A_{t}$. But $X_{t}\subset A_{t}$. Hence $[X]_{t}^{\nabla_{A}}\neq X_{t}$. Therefore $X\not\in \nabla_{A}\text{-}\mathrm{Sat}(A)$.
%
% The converse implication is straightforward.
% \end{proof}
%
% \begin{remark}
% Let $A$ be an $S$-sorted set. Then, from the above proposition, it follows that $\varnothing^{S}$, $A\in \nabla_{A}\text{-}\mathrm{Sat}(A)$. Moreover, for every subset $T$ of $S$, we have that $\bigcup_{t\in T}\delta^{t,A_{t}}\in \nabla_{A}\text{-}\mathrm{Sat}(A)$.
% \end{remark}
%
% \begin{proposition}
% Let $A$ be an $S$-sorted set, $X\subseteq A$, and $\Phi$, $\Psi\in\mathrm{Eqv}(A)$. Then
% $[X]^{\Phi\cap\Psi}\subseteq[X]^{\Phi}\cap[X]^{\Psi}$.
% \end{proposition}
%
% \begin{proof}
% Let $s$ be a sort in $S$ and $b\in[X]^{\Phi\cap \Psi}_{s}$. Then, by definition, there exists an $a\in X_{s}$ such that $(a,b)\in (\Phi\cap\Psi)_{s} = \Phi_{s}\cap\Psi_{s}$. Hence, $(a,b)\in \Phi_{s}$ and $(a,b)\in \Psi_{s}$. Therefore $b\in [X]^{\Phi}_{s}$ and $b\in[X]^{\Psi}_{s}$. Consequently, $b\in ([X]^{\Phi}\cap [X]^{\Psi})_{s}$. Thus $[X]^{\Phi\cap\Psi}\subseteq[X]^{\Phi}\cap[X]^{\Psi}$.
% \end{proof}
%
% \begin{corollary}
% Let $A$ be an $S$-sorted set and $\Phi$, $\Psi\in\mathrm{Eqv}(A)$. Then we have that  $\Phi\text{-}\mathrm{Sat}(A)\cap\Psi\text{-}\mathrm{Sat}(A)\subseteq(\Phi\cap\Psi)\text{-}\mathrm{Sat}(A)$.
% \end{corollary}
%
% We next state that the set $\Phi\text{-}\mathrm{Sat}(A)$ is the set of all fixed points of a suitable operator on $A$, i.e., of an endomapping of $\mathrm{Sub}(A)$.
%
% \begin{proposition}\label{PSatOp}
% Let $A$ be an $S$-sorted set and $\Phi\in\mathrm{Eqv}(A)$. Then the mapping $[\cdot]^{\Phi}$ from $\mathrm{Sub}(A)$ to $\mathrm{Sub}(A)$ that sends $X$ in $\mathrm{Sub}(A)$ to $[\cdot]^{\Phi}(X) = [X]^{\Phi}$ in $\mathrm{Sub}(A)$
% %defined as follows:
% %$$
% %[\cdot]^{\Phi} \nfunction
% %{\mathrm{Sub}(A)}
% %{\mathrm{Sub}(A)}
% %{X}
% %{ [\cdot]^{\Phi}(X) = [X]^{\Phi}  }
% %$$
% is a completely additive closure operator on $A$. Moreover, for every nonempty set $I$ in $\boldsymbol{\mathcal{U}}$ and every $I$-indexed family $(X^{i})_{i\in I}$ in $\mathrm{Sub}(A)$, $[\bigcap_{i\in I}X^{i}]^{\Phi} \subseteq \bigcap_{i\in I}[X^{i}]^{\Phi}$ (and, obviously, $[A]^{\Phi} = A$), and, for every $X\subseteq A$, if $X = [X]^{\Phi}$, then  $\complement_{A}X = [\complement_{A}X]^{\Phi}$. Besides, $[\cdot]^{\Phi}$ is uniform, i.e., is such that, for every $X$, $Y\subseteq A$, if  $\mathrm{supp}_{S}(X) = \mathrm{supp}_{S}(Y)$, then $\mathrm{supp}_{S}([X]^{\Phi}) = \mathrm{supp}_{S}([Y]^{\Phi})$---hence, in particular, $[\cdot]^{\Phi}$ is a uniform algebraic closure operator on $A$. And $\Phi\text{-}\mathrm{Sat}(A) = \mathrm{Fix}([\cdot]^{\Phi})$, where $\mathrm{Fix}([\cdot]^{\Phi})$ is the set of all fixed point of the operator $[\cdot]^{\Phi}$.
% \end{proposition}
%
% %\begin{remark}
% %Given a subset $X$ of $A$, if $X = [X]^{\Phi} = (\mathrm{pr}^{\Phi})^{-1}[\mathcal{Y}]$, for some $\mathcal{Y}\subseteq A/\Phi$, then $\complement_{A}X = (\mathrm{pr}^{\Phi})^{-1}[\complement_{A/\Phi}\mathcal{Y}]$. Hence $\complement_{A}X = [\complement_{A}X]^{\Phi}$.
% %\end{remark}
%
% \begin{proposition}\label{PSatCABA} Let $A$ be an $S$-sorted set and $\Phi\in\mathrm{Eqv}(A)$. Then the ordered pair
% $\Phi\text{-}\mathbf{Sat}(A) = (\Phi\text{-}\mathrm{Sat}(A),\subseteq)$ is a complete atomic Boolean algebra.
% %(for brevity a CABA).
% \end{proposition}
%
% \begin{proof}
% The proof is straightforward and we leave it to the reader. We only point out that the atoms of $\Phi\text{-}\mathbf{Sat}(A)$ are precisely the deltas of Kronecker $\delta^{t, [x]_{\Phi_t}}$, for some $t\in S$ and some $x\in A_{t}$, and that, obviously, every $\Phi$-saturated subset $X$ of $A$ is the join ($\equiv$ union) of all atoms smaller than $X$.
% \end{proof}
%
% \begin{proposition}\label{POptLift}
% Let $A$ be an $S$-sorted set, $(B,\Psi)$ a classified $S$-sorted set and $f\colon A\mor B$ an $S$-sorted mapping; situation that we will indicate by:
% $$
%   f\colon A\mor (B,S).
% $$
% Then there exists an \emph{optimal lift of} $\Psi$ \emph{through} $f$, i.e., there exists an $S$-sorted equivalence on $A$, denoted by $\mathrm{L}^{f}(\Psi)$, the \emph{optimal lift of} $\Psi$ \emph{through} $f$, such that 
% \begin{enumerate}
% \item $f\colon (A,\mathrm{L}^{f}(\Psi))\mor (B,\Psi)$ is a morphism and,
% \item for every  classified $S$-sorted set $(D,\Upsilon)$ and every $S$-sorted mapping $g\colon D\mor A$, 
% if $f\circ g\colon (D,\Upsilon)\mor (B,\Psi)$ is a morphism, then
% $g\colon (D,\Upsilon)\mor(A,\mathrm{L}^{f}(\Psi))$ is a morphism.
% \end{enumerate}
% Moreover, we have that:
% \begin{enumerate}
% \item For every $S$-sorted equivalence $\Phi\in \mathrm{Eqv}(A)$:
% $$
%   \mathrm{L}^{\id^{A}}(\Phi) = \Phi.
% $$
% \item If $f\colon A\mor B$, $g\colon B\mor C$ are $S$-sorted mappings and $\Omega\in
% \mathrm{Eqv}(C)$, then:
% $$
%   \mathrm{L}^{g\circ f}(\Omega) = \mathrm{L}^{f}(\mathrm{L}^{g}(\Omega)).
% $$
% \end{enumerate}
% \end{proposition}
%
% %\begin{proof}
% %Es suficiente tomar como $\mathrm{L}^{f}(R)$ la relación sobre $A$
% %definida como:
% %$$
% %  \mathrm{L}^{f}(R) = \{\,(x,y)\in A\times A\mid (f(x),f(y))\in R\,\}.
% %$$
% %\end{proof}
%
% \begin{remark}\label{ROptLiftDiag}
% Let $A$ be an $S$-sorted set, $(B,\Psi)$ a classified $S$-sorted set and $f\colon A\mor B$ an $S$-sorted mapping. Then there exists a unique $S$-sorted mapping $f^{\mathrm{L}^{f}(\Psi),\Psi}$ from $A/\mathrm{L}^{f}(\Psi)$ to $B/\Psi$ such that $\mathrm{pr}^{\Psi}\circ f = f^{\mathrm{L}^{f}(\Psi),\Psi}\circ \mathrm{pr}^{\mathrm{L}^{f}(\Psi)}$, i.e., the following diagram commutes
% $$\xymatrix{
% A
% \ar[r]^-{f}
% \ar[d]_{\mathrm{pr}^{\mathrm{L}^{f}(\Psi)}} &
% B
% \ar[d]^{\mathrm{pr}^{\Psi}}
% \\
% A/\mathrm{L}^{f}(\Psi)
% \ar[r]_{f^{\mathrm{L}^{f}(\Psi),\Psi}} &
% B/\Psi
% }
% $$
% \end{remark}
%
% \begin{proposition}\label{PClfd}
% The functor from $\mathsf{ClfdSet}$ to $\mathsf{ClfdSet}^{S}$, whose object mapping assigns to a classified set $(X,\Theta)$ the classified $S$-sorted set $((X)_{s\in S},(\Theta)_{s\in S})$ and whose morphism mapping sends a morphism $f$ from $(X,\Theta)$ to $(Y,\Omega)$ to the morphism $(f)_{s\in S}$ from $((X)_{s\in S},(\Theta)_{s\in S})$ to $((Y)_{s\in S},(\Omega)_{s\in S})$, has a left adjoint.
% \end{proposition}
%
% \begin{proof}
% Let $(A,\Phi)$ be a classified $S$-sorted set. Then $(\coprod A, \Phi^{\coprod A})$, where $\Phi^{\coprod A}$ is the equivalence relation on $\coprod A$ defined as
% $$
% \textstyle
% \Phi^{\coprod A} = \{((x,s),(y,t))\in ( \coprod A)^{2}\mid s = t \And (x,y)\in \Phi_ {s}\},
% $$
% together with the canonical embedding of $(A,\Phi)$ into $((\coprod A)_{s\in S},(\Phi^{\coprod A})_{s\in S})$ is a universal morphism from $(A,\Phi)$ to the  functor from $\mathsf{ClfdSet}$ to $\mathsf{ClfdSet}^{S}$. The details are left to the reader.
% \end{proof}
%
% \subsection{Free monoids}
%
% We next define the concept of free monoid on a set and several notions associated with it that will be used afterwards to construct the free algebra on an $S$-sorted set and to define various substitution operators.
%
% \begin{definition}
% Let $A$ be a set. The \emph{free monoid on} $A$, denoted by $\mathbf{A}^{\star}$, is $(A^{\star},\curlywedge,\lambda)$, where $A^{\star}$, the set of all \emph{words on} $A$, is $\bigcup_{n\in\mathbb{N}}\mathrm{Hom}(n,A)$, the set of all mappings $\mathbf{a}\colon n\mor A$ from some $n\in \mathbb{N}$ to $A$. A word $\mathbf{a}\in A^{\star}$ is usually denoted as a sequence $(a_{i})_{i\in\bb{\mathbf{a}}}$, where, for $i\in\bb{\mathbf{a}}$, $a_{i}$ is the letter in $A$ satisfying $\mathbf{a}(i)=a_{i}$. Furthermore, $\curlywedge$, the \emph{concatenation} of words on $A$, is the binary operation on $A^{\star}$ which sends a pair of words $(\mathbf{a},\mathbf{b})$ on $A$ to the mapping $\mathbf{a}\curlywedge \mathbf{b}$ from $\bb{\mathbf{a}}+\bb{\mathbf{b}}$ to $A$, where $\bb{\mathbf{a}}$ and $\bb{\mathbf{b}}$ are the lengths ($\equiv$ domains) of the mappings $\mathbf{a}$ and $\mathbf{b}$, respectively, defined as follows:
% $$
% \mathbf{a}\bconcat \mathbf{b}
% \nfunction
% {\bb{\mathbf{a}}+\bb{\mathbf{b}}}{A}
% {i}{
% \begin{cases}
%   a_{i}, & \text{if $0\leq i < \bb{\mathbf{a}}$;}\\
%   b_{i-\bb{\mathbf{a}}}, & \text{if $\bb{\mathbf{a}}\leq i < \bb{\mathbf{a}}+\bb{\mathbf{b}}$,}
% \end{cases}
%    }
% $$
% and $\lambda$, the \emph{empty word on} $A$, is the unique mapping from $\varnothing$ to $A$. We will denote by $\eta^{A}$ the mapping from $A$ to $A^{\star}$ that sends $a\in A$ to $(a)\in A^{\star}$, i.e., to the mapping $(a)\colon 1\mor A$ that sends $0$ to $a$. The ordered pair $(\mathbf{A}^{\star},\eta^{A})$ is a universal morphism from $A$ to the forgetful functor from the category $\mathsf{Mon}$, of monoids, to $\mathsf{Set}$.
% \end{definition}
%
% \begin{remark}
% For a word $\mathbf{a}\in A^{\star}$, $\bb{\mathbf{a}}$, the length of $\mathbf{a}$, is the value at $\mathbf{a}$ of the unique homomorphism $\bb{\,\cdot\,}$ from $\mathbf{A}^{\star}$ to $(\mathbb{N},+,0)$, the additive monoid of the natural numbers, such that $\bb{\,\cdot\,}\circ \eta^{A} = \kappa_{1}$, where $\kappa_{1}$ is the mapping from $A$ to $\mathbb{N}$ constantly $1$. Note that, for every $n\in \mathbb{N}$, $\bb{\,\cdot\,}$ sends $\mathbf{a}\in \mathrm{Hom}(n,A)$ to $n$. Thus, for the family of mappings $(\kappa_{n})_{n\in \mathbb{N}}$, where, for every $n\in \mathbb{N}$, $\kappa_{n}$ is the mapping from $\mathrm{Hom}(n,A)$ to $\mathbb{N}$ constantly $n$, and by applying the universal property of the coproduct, we have $\bb{\,\cdot\,} = [\kappa_{n}]_{n\in \mathbb{N}}$.
% \end{remark}
%
%
% %\begin{definition}
% %Let $\mathbf{c}$ be a word in $C^{\star}$, $k$, $l\in\bb{\mathbf{c}}$ such that $k\leq l$, and $\mathbf{c}'\in C^{\star}$. We will say that $\mathbf{c}'$ is a subword of $\mathbf{c}$ \emph{beginning at position} $k$ and \emph{ending at position} $l$ if there are two words $\mathbf{a}, \mathbf{b}\in C^{\star}$ such that $\mathbf{c} = \mathbf{a}\curlywedge\mathbf{c}'\curlywedge\mathbf{b}$, $\bb{\mathbf{a}} = k$, and $\bb{\mathbf{b}} = (\bb{\mathbf{c}}-1)-l$. For $\mathbf{c}\in C^{\star}$ and $k$, $l\in\bb{\mathbf{c}}$ such that $k\leq l$, we will denote by $\mathbf{c}^{k,l}$ the word $\mathbf{c}'$  beginning at position $k$ and ending at position $l$. In particular, we will denote by $\mathbf{c}^{0,l}$ the word $\mathbf{c}'$  beginning at position $0$ and ending at position $l$ (which is an initial segment of $\mathbf{c}$) and by $\mathbf{c}^{k,\bb{\mathbf{c}}-1}$ the word $\mathbf{c}'$  beginning at position $k$ and ending at position $\bb{\mathbf{c}}-1$ (which is a final segment of $\mathbf{c}$).
% %\end{definition}
% %
% %\begin{remark}
% %Let $\mathbf{c}$ be a word in $C^{\star}$ and $k$, $l\in\bb{\mathbf{c}}$ such that $k\leq l$. Then $\bb{\mathbf{c}^{k,l}}=l-k+1$. In particular, $\bb{\mathbf{c}^{0,k}}=k+1$ and $\bb{\mathbf{c}^{l,\bb{\mathbf{c}}-1}}=\bb{\mathbf{c}}-l$.
% %\end{remark}
%
% \begin{definition}\label{DSubw}
% Let $\mathbf{a}$ and $\mathbf{a}'$ be words in $A^{\star}$ and $k$, $l\in\bb{\mathbf{a}}$ such that $k\leq l$. We will say that $\mathbf{a}'$ is a \emph{subword} of $\mathbf{a}$ \emph{beginning at position} $k$ and \emph{ending at position} $l$ if there are words $\mathbf{b}$ and $\mathbf{c}$ such that $\mathbf{a} = \mathbf{b}\bconcat \mathbf{a}'\bconcat \mathbf{c}$, $\bb{\mathbf{b}} = k$, and $\bb{\mathbf{c}} = (\bb{\mathbf{a}}-1)-l$. For $\mathbf{a}\in A^{\star}$ and $k$, $l\in\bb{\mathbf{a}}$ such that $k\leq l$, we will denote by $\mathbf{a}^{k,l}$ the word $\mathbf{a}'$  beginning at position $k$ and ending at position $l$. In particular, we will denote by $\mathbf{a}^{0,l}$ the word $\mathbf{a}'$  beginning at position $0$ and ending at position $l$, which is an \emph{initial segment} of $\mathbf{a}$, and by $\mathbf{a}^{k,\bb{\mathbf{a}}-1}$ the word $\mathbf{a}'$  beginning at position $k$ and ending at position $\bb{\mathbf{a}}-1$, which is a \emph{final segment} of $\mathbf{a}$. We will call the word $\mathbf{a}^{0,\bb{\mathbf{a}}-2}$ the \emph{maximal prefix} of $\mathbf{a}$. Moreover, we will call $\mathbf{a}^{0,0}$, which is, essentially, the first letter of $\mathbf{a}$, the \emph{head} of $\mathbf{a}$ and $\mathbf{a}^{\bb{\mathbf{a}}-1,\bb{\mathbf{a}}-1}$, which is, essentially, the last letter of $\mathbf{a}$, the \emph{tail} of $\mathbf{a}$. A word $\mathbf{a}$ may have several subwords equal to $\mathbf{a}'$. In that case, the equation $\mathbf{a} = \mathbf{b}\bconcat \mathbf{a}'\bconcat \mathbf{c}$ has several solutions $(\mathbf{b},\mathbf{c})$.
% %The first subword of $w$ equal to $w'$ is called \emph{the first occurrence of }$w'$ \emph{in} $w$, the second subword of $w$ equal to $w'$ is \emph{the second occurrence of} $w'$ \emph{in} $w$ and so on.
% If the pairs $(\mathbf{b}_{i},\mathbf{c}_{i})$ $(i\in n)$ are all solutions of $\mathbf{a} = \mathbf{b}\bconcat \mathbf{a}'\bconcat \mathbf{c}$ and if $\bb{\mathbf{b}_{0}}< \bb{\mathbf{b}_{1}}<\cdots<\bb{\mathbf{b}_{n-1}}$, then $(\mathbf{b}_{i},\mathbf{c}_{i})$ determines the $i$-th occurrence of $\mathbf{a}'$ in $\mathbf{a}$. The solution $(\mathbf{b},\mathbf{c})$ in which either $\mathbf{b}$ or $\mathbf{c}$ is $\lambda$ is not excluded.
% %If $u = \lambda$, then we will say that $w'$ is an \emph{initial segment} of $w$, and if $v = \lambda$, then we will say that $w'$ is a \emph{final segment} of $w$.
% Let $\mathbf{a}$ be a word in $A^{\star}$ and $a\in A$. We will say that \emph{$a$ occurs in} $\mathbf{a}$ if there are words $\mathbf{b}$, $\mathbf{c}$ in $A^{\star}$ such that $\mathbf{a} = \mathbf{b}\bconcat(a)\bconcat \mathbf{c}$. Note that $a$ occurs in $\mathbf{a}$ if and only if there exists an $i\in \bb{\mathbf{a}}$ such that $\mathbf{a}(i) = a$. We will denote by $\bb{\mathbf{a}}_{a}$ the natural number $\mathrm{card}(\{i\in \bb{\mathbf{a}}\mid \mathbf{a}(i) = a\}) = \mathrm{card}(\mathbf{a}^{-1}[\{a\}])$, i.e., the number of occurrences of $a$ in $\mathbf{a}$. Moreover, we let $(i_{j})_{j\in\bb{\mathbf{a}}_{a}}$ stand for the  enumeration in ascending order of the occurrences of $a$ in $\mathbf{a}$. Thus $(i_{j})_{j\in\bb{\mathbf{a}}_{a}}$ is the order embedding of $(\bb{\mathbf{a}}_{a},<)$ into $(\bb{\mathbf{a}},<)$ defined  recursively as follows:
% $$i_{0} = \mathrm{min}\{i\in \bb{\mathbf{a}}\mid \mathbf{a}(i) = a\},$$
% and, for $j\in [1,\bb{\mathbf{a}}_{a}-1]$, 
% $$
% i_{j} = \mathrm{min}\{i\in \bb{\mathbf{a}}-\{i_{0},\ldots,i_{j-1}\}\mid \mathbf{a}(i) = a\}.
% $$
% If, for every $j\in \bb{\mathbf{a}}_{a}$,  the pairs $(\mathbf{b}_{i_{j}},\mathbf{c}_{i_{j}})$  are all solutions of $\mathbf{a} = \mathbf{b}_{i_{j}}\bconcat (a)\bconcat \mathbf{c}_{i_{j}}$ and if $\bb{\mathbf{b}_{i_{0}}}< \bb{\mathbf{b}_{i_{1}}}<\cdots<\bb{\mathbf{b}_{i_{\bb{\mathbf{a}}_{a}-1}}}$, then $(\mathbf{b}_{i_{j}},\mathbf{c}_{i_{j}})$ determines the $j$-th occurrence of $(a)$ in $\mathbf{a}$ and we will say that \emph{$a$ occurs at the $i_{j}$-th place of $\mathbf{a}$}.
% %For a word $w$ in $A^{\star}$ we let $\downarrow\!\! w$ stand for the $A$-sorted set defined, for every $a\in A$ as $(\downarrow\!\! w)_{a} = w^{-1}[\{a\}]$. Therefore, if $(i_{\alpha})_{\alpha\in\bb{w}_{a}}$ is the order embedding of $(\bb{w}_{a},<)$ into $(\bb{w},<)$ defined above, then $(\downarrow\!\! w)_{a} = \{i_{\alpha}\mid \alpha\in\bb{w}_{a}\}$, i.e., the image of $(i_{\alpha})_{\alpha\in\bb{w}_{a}}$. Sometimes, for a fixed $A$-countably infinite $A$-sorted set $V^{A} = (V^{A}_{a})_{a\in A}$, where, for every $a\in A$, $V^{A}_{a} = \{v_{a,n}\mid n\in \mathbb{N}\}$, we will identify $\downarrow\!\! w$ with the $A$-sorted set $(\downarrow\!\! v_{a,\bb{w}_{a}})_{a\in A}$, defined, for every $a\in A$, as $\downarrow\!\! v_{a,\bb{w}_{a}} = \{v_{a,\alpha}\mid \alpha\in \bb{w}_{a}\}$.
% \end{definition}
%
% %\begin{remark}
% %Let $w$ be a word in $A^{\star}$ and $k$, $l\in\bb{w}$ such that $k\leq l$. Then $\bb{w^{k,l}}=l-k+1$. In particular, $\bb{w^{0,k}}=k+1$ and $\bb{w^{l,\bb{w}-1}}=\bb{w}-l$.
% %\end{remark}
%
% \begin{remark}
% Let $a$ be an element of $A$. Then, for $\mathbf{a}\in A^{\star}$, $\bb{\mathbf{a}}_{a}$, the number of occurrences of $a$ in $\mathbf{a}$, is the value at $\mathbf{a}$ of the unique homomorphism $\bb{\,\cdot\,}_{a}$ from $\mathbf{A}^{\star}$ to $(\mathbb{N},+,0)$ such that $\bb{\,\cdot\,}_{a}\circ\eta^{A} = \delta_{a}$, where $\delta_{a}$ is the mapping from $A$ to $\mathbb{N}$ that sends $a\in A$ to $1\in\mathbb{N}$ and $b\in A-\{a\}$ to $0\in \mathbb{N}$. Therefore, since, for every $\mathbf{a}\in A^{\star}$, the $A$-indexed family $(\bb{\mathbf{a}}_{a})_{a\in A}$ in $\mathbb{N}$ is such that $\bb{\mathbf{a}}_{a} = 0$ for all but a finite number of elements $a$ in $A$, i.e., is such that $\mathrm{card}(\{a\in A\mid \bb{\mathbf{a}}_{a}\neq 0\})<\aleph_{0}$, we have that $\sum_{a\in A}\bb{\mathbf{a}}_{a}\in \mathbb{N}$ and, obviously, for every $\mathbf{a}\in A^{\star}$, $\bb{\mathbf{a}} = \sum_{a\in A}\bb{\mathbf{a}}_{a}$, i.e., $\bb{\,\cdot\,} = \sum_{a\in A}\bb{\,\cdot\,}_{a}$.
% \end{remark}
%
% %\begin{remark}
% %Let $w$ be a word in $A^{\star}$ and $X$ and $A$-sorted set. Then $X_{\downarrow w} = \mathrm{Hom}(\downarrow\!\! w,X)$, the set of all $A$-sorted mappings from $\downarrow\!\! w$ to $X$, is isomorphic to $X_{w} = \prod_{i\in\bb{w}}X_{w_{i}}$.
% %
% %On the other hand, if $X\in \mathrm{Sub}_{\mathrm{f}}(V^{A})$ and $m = (m_{a})_{a\in A}$ is the $S$-cardinal of $X$, then $m\in \mathbb{N}^{(A)}$, the underlying set of $\mathbf{N}^{(A)}$, the free abelian monoid on $A$. Therefore, since $\mathbf{N}^{(A)}$ is a homomorphic image of $\mathbf{A}^{\star}$, there exists a word $w$ in $A^{\star}$ such that $\downarrow\!\! w$ and $X$ are isomorphic. These isomorphisms are very useful when dealing with derived operations of the many-sorted algebras, with many-sorted terms and with equations between many-sorted terms.
% %\end{remark}
%
% %\begin{remark}
% %Let $w$ be a word in $A^{\star}$ and $a\in A$. Then the occurrences of $a$ in $w$ are linearly ordered, i.e., $(\{i_{\alpha}\mid \alpha\in \bb{w}_{a}\},<)$ is a finite linearly ordered set, where, for every $\alpha\in \bb{w}_{a}$, $i_{\alpha}$ indicates the position in which $a$ occurs in $w$ and, for every $\alpha, \beta\in \bb{w}_{a}$, $i_{\alpha}<i_{\beta}$ if and only if $\alpha<\beta$. Therefore, we can represent  $w$ as:
% %$$
% %w = u_{0}\bconcat (w(i_{0}))\bconcat u_{1}\bconcat (w(i_{1}))\bconcat\cdots\bconcat u_{\bb{w}_{a}-1}\bconcat (w(i_{\bb{w}_{a}-1}))\bconcat u,
% %$$
% %for univocally determined words $u_{0}$, \ldots, $u_{\bb{w}_{a}-1}$, $u$ in $A^{\star}$.
% %\end{remark}
%
% %\begin{definition}
% %Let $w$ and $w'$ be words in $A^{\star}$. We will say that $w'$ is a \emph{subword} of a word $w$ if there are words $u$ and $v$ such that $w = u\bconcat w'\bconcat v$. A word $w$ may have several subwords equal to $w'$. In that case, the equation $w = u\bconcat w'\bconcat v$ has several solutions $(u,v)$. The first subword of $w$ equal to $w'$ is called \emph{the first occurrence of }$w'$ \emph{in} $w$, the second subword of $w$ equal to $w'$ is \emph{the second occurrence of} $w'$ \emph{in} $w$ and so on. If the pairs $(u_{i},v_{i})$ $(i\in n)$ are all solutions of $w = u\bconcat w'\bconcat v$ and if $\bb{u_{0}}< \bb{u_{1}}<\cdots<\bb{u_{n-1}}$, then $(u_{i},v_{i})$ determines the $i$-th occurrence of $w'$ in $w$. The solution $(u,v)$ in which either $u$ or $v$ is $\lambda$ is not excluded.
% %\end{definition}
%
% %\begin{definition}
% %Let $w$ and $w''$ be words in $A^{\star}$. If $w'$ is subword of $w$ and $u_{i}\bconcat w'\bconcat v_{i}$ is the $i$-th occurrence of $w'$ in $w$, then \emph{the result of substituting the word} $w''$ \emph{in} $w$ \emph{for the $i$-th occurrence o}f $w'$ \emph{in} $w$ is the word $u_{i}\bconcat w''\bconcat v_{i}$.
% %%, denoted by $\left(\!\begin{smallmatrix}w'\\ w''\end{smallmatrix}\!\right)_{i}(w)$,
% %\end{definition}
%
% %\begin{definition} Let $w$ be a word in $A^{\star}$. Then we will denote by $\left(\!\begin{smallmatrix}a\\ \cdot\end{smallmatrix}\!\right)_{a\in A}(w)$ the mapping from $\prod_{a\in A}(A^{\star})^{\bb{w}_{a}}$ to $A^{\star}$ that assigns to $((q^{a}_{\alpha})_{\alpha\in \bb{w}_{a}})_{a\in A}$ in $\prod_{a\in A}(A^{\star})^{\bb{w}_{a}}$ the word $\left(\!\begin{smallmatrix}a\$q^{a}_{\alpha})_{\alpha\in \bb{w}_{a}}\end{smallmatrix}\!\right)_{a\in A}(w)$ in $A^{\star}$ obtained by substituting in $w$, for every $a\in A$ and every $\alpha\in \bb{w}_a$, $q^{a}_{\alpha}$ for the $i_{\alpha}$-th occurrence of $a$ in $w$. We call  $\left(\!\begin{smallmatrix}a\$q^{a}_{\alpha})_{\alpha\in \bb{w}_{a}}\end{smallmatrix}\!\right)_{a\in A}(w)$ \emph{the substitution of $(q_{\alpha})_{\alpha\in\bb{w}_{a}}$ for $a$ in $w$ for every $a$ in $A$}, and $\left(\!\begin{smallmatrix}a\\ \cdot\end{smallmatrix}\!\right)_{a\in A}(w)$ the \emph{substitution operator for $w$}. Let $c$ be an element of $A$. Then we will denote by $\left(\!\begin{smallmatrix}c\\ \cdot\end{smallmatrix}\!\right)(w)$ the mapping from $(A^{\star})^{\bb{w}_{c}}$ to $A^{\star}$ that assigns to $(q^{c}_{\alpha})_{\alpha\in \bb{w}_{c}}\in (A^{\star})^{\bb{w}_{c}}$ the term $\left(\!\begin{smallmatrix}c\$q^{c}_{\alpha})_{\alpha\in \bb{w}_{c}}\end{smallmatrix}\!\right)(w)$ in $A^{\star}$ obtained by substituting in $w$, for every $\alpha\in \bb{w}_{c}$, $q^{c}_{\alpha}$ for the $i_{\alpha}$-th occurrence of $c$ in $w$. We call $\left(\!\begin{smallmatrix}c\$q^{c}_{\alpha})_{\alpha\in \bb{w}_{c}}\end{smallmatrix}\!\right)(w)$ \emph{the substitution of $(q^{c}_{\alpha})_{\alpha\in\bb{w}_{c}}$ for $c$ in $w$}.
% %\end{definition}
%
% %\begin{remark}
% %Therefore, for every set $A$, we have obtained the family of substitution operators:
% %$$
% %\textstyle
% %(\left(\!\begin{smallmatrix}a\\ \cdot\end{smallmatrix}\!\right)_{a\in A}(w))_{w\in A^{\star}}\in \prod_{w\in A^{\star}}\mathrm{Hom}(\prod_{a\in A}(A^{\star})^{\bb{w}_{a}},A^{\star}).
% %$$
% %\end{remark}
%
% %\begin{proof}
% %If $\bb{w}_{a} = 0$, then the mapping $\left(\!\begin{smallmatrix}a\\ \cdot\end{smallmatrix}\!\right)(w)$ assigns to the unique element of $(A^{\star})^{0}$ precisely $w\in A^{\star}$. We now turn to the case $\bb{w}_{a} \neq 0$. For the word $w$ in $A^{\star}$ and the letter $a\in A$ let us consider the  factorization:
% %$$
% %w = u_{0}\bconcat (w(i_{0}))\bconcat u_{1}\bconcat (w(i_{1}))\bconcat\cdots\bconcat u_{\bb{w}_{a}-1}\bconcat (w(i_{\bb{w}_{a}-1}))\bconcat u,
% %$$
% %where, for every $\alpha\in \bb{w}_{a}$, $w(i_{\alpha}) = a$ and $u_{\alpha}$, $u\in A^{\star}$.
% %
% %Let us consider, for $(q_{\alpha})_{\alpha\in \bb{w}_{a}}\in (A^{\star})^{\bb{w}_{a}}$, the following diagram:
% %$$\xymatrix@C=65pt@R=40pt{
% %{} & 1 \ar[d]^{\kappa_{(q_{\alpha})_{\alpha\in \bb{w}_{a}}}} & \\
% %{} & (A^{\star})^{\bb{w}_{a}}\ar[rd]^*[l]{\mathrm{pr}_{0} = \mathrm{pr}_{\bb{w}_{a},0}}\ar[ld]_{\kappa_{u_{0}}}\ar[d]|{\left<\kappa_{u_{0}},%\mathrm{pr}_{0}\right>} & {} \\
% %A^{\star} & A^{\star}\times A^{\star}\ar[l]^{\mathrm{pr}_{0}}\ar[r]_{\mathrm{pr}_{1}}\ar[d]^{\bconcat} & A^{\star} \\
% %{} & A^{\star} & {}
% %   }
% %$$
% %in which the value of the mapping $\kappa_{(q_{i})_{i\in \bb{w}_{a}}}$ from $1$ to $(A^{\star})^{\bb{w}_{a}}$ at $0$ is $(q_{\alpha})_{\alpha\in %\bb{w}_{a}}$, the value of the mapping $\kappa_{u_{0}}$ from $(A^{\star})^{\bb{w}_{a}}$ to $A^{\star}$ for every argument is $u_{0}$, and $\mathrm{pr}_{0} = \mathrm{pr}_{\bb{w}_{a},0}$ assigns to $(q_{\alpha})_{\alpha\in \bb{w}_{a}}$ precisely $q_{0}$. The value of the mapping $\bconcat\circ \left<\kappa_{u_{0}},\mathrm{pr}_{0}\right>\circ \kappa_{(q_{\alpha})_{\alpha\in \bb{w}_{a}}}$ from $1$ to $A^{\star}$ at $0$ is $u_{0}\bconcat q_{0}$.
% %
% %We next consider the diagram:
% %$$\xymatrix@C=65pt@R=40pt{
% %{} & 1 \ar[d]^{\kappa_{(q_{\alpha})_{\alpha\in \bb{w}_{a}}}} & \\
% %{} & (A^{\star})^{\bb{w}_{a}}\ar[rd]^{\kappa_{u_{1}}}\ar[ld]_*[r]{\bconcat\circ \left<\kappa_{u_{0}},\mathrm{pr}_{0}\right>}\ar[d]|(0.6){\left<%\bconcat\circ \left<\kappa_{u_{0}},\mathrm{pr}_{0}\right>,\kappa_{u_{1}}\right>} & {} \\
% %A^{\star} & A^{\star}\times A^{\star}\ar[l]^{\mathrm{pr}_{0}}\ar[r]_{\mathrm{pr}_{1}}\ar[d]^{\bconcat} & A^{\star} \\
% %{} & A^{\star} & {}
% %   }
% %$$
% %in which the value of the mapping $\kappa_{u_{1}}$ from $(A^{\star})^{\bb{w}_{a}}$ to $A^{\star}$ for every argument is $u_{1}$. The value of the mapping $\bconcat\circ \left<\bconcat\circ \left<\kappa_{u_{0}},\mathrm{pr}_{0}\right>,\kappa_{u_{1}}\right>\circ \kappa_{(q_{\alpha})_{\alpha\in \bb{w}_{a}}}$ from $1$ to $A^{\star}$ at $0$ is $u_{0}\bconcat q_{0}\bconcat u_{1}$.
% %
% %Now let us consider the diagram:
% %$$\xymatrix@C=110pt@R=40pt{
% %{} & 1 \ar[d]^{\kappa_{(q_{\alpha})_{\alpha\in \bb{w}_{a}}}} & \\
% %{} & (A^{\star})^{\bb{w}_{a}}\ar[rd]^*[l]{\mathrm{pr}_{1} = \mathrm{pr}_{\bb{w}_{a},1}}\ar[ld]_*[r]{\bconcat\circ \left<\bconcat\circ\left<\kappa_{u_{0}},\mathrm{pr}_{0}\right>,\kappa_{u_{1}}\right>}\ar[d]|(0.6){\left<\bconcat\circ \left<\bconcat\circ\left<\kappa_{u_{0}},\mathrm{pr}_{0}\right>,\kappa_{u_{1}}\right>,\mathrm{pr}_{1}\right>} & {} \\
% %A^{\star} & A^{\star}\times A^{\star}\ar[l]^{\mathrm{pr}_{0}}\ar[r]_{\mathrm{pr}_{1}}\ar[d]^{\bconcat} & A^{\star} \\
% %{} & A^{\star} & {}
% %   }
% %$$
% %in which $\mathrm{pr}_{1} = \mathrm{pr}_{\bb{w}_{a},1}$ assigns to $(q_{\alpha})_{\alpha\in \bb{w}_{a}}$ precisely $q_{1}$. The value of the mapping $\bconcat\circ \left<\bconcat\circ \left<\bconcat\circ\left<\kappa_{u_{0}},\mathrm{pr}_{0}\right>,\kappa_{u_{1}}\right>,\mathrm{pr}_{1}\right>\circ \kappa_{(q_{\alpha})_{\alpha\in \bb{w}_{a}}}$ from $1$ to $A^{\star}$ at $0$ is $u_{0}\bconcat q_{0}\bconcat u_{1}\bconcat q_{1}$.
% %
% %It is obvious that by iterating this procedure a finite number of times one finally obtains
% %$$
% %u_{0}\bconcat q_{0}\bconcat u_{1}\bconcat q_{1}\bconcat\cdots\bconcat u_{\bb{w}_{a}-1}\bconcat q_{\bb{w}_{a}-1}\bconcat u.
% %$$
% %And this constitutes, for $\bb{w}_{a} \neq 0$, the definition of the mapping $\left(\!\begin{smallmatrix}a\\ \cdot\end{smallmatrix}\!\right)(w)$ from $(A^{\star})^{\bb{w}_{a}}$ to $A^{\star}$ at $(q_{\alpha})_{\alpha\in \bb{w}_{a}}\in (A^{\star})^{\bb{w}_{a}}$.
% %\end{proof}
%
% %We will denote by $\left(\!\begin{smallmatrix}a\$q_{\alpha})_{\alpha\in \bb{w}_{a}}\end{smallmatrix}\!\right)(w)$ the value of $\left(\!\begin{smallmatrix}a\\ \cdot\end{smallmatrix}\!\right)(w)$ at $(q_{\alpha})_{\alpha\in \bb{w}_{a}}$. Let us observe that, in virtue of the universal property of the coproduct, using the family of substitution operator $(\left(\!\begin{smallmatrix}a\\ \cdot\end{smallmatrix}\!\right)(w))_{(w,a)\in A^{\star}\times A}$, we obtain the operator:
% %$$
% %\textstyle
% %[\left(\!\begin{smallmatrix}a\\ \cdot\end{smallmatrix}\!\right)(w)]_{(w,a)\in A^{\star}\times A}\colon
% %\coprod_{(w,a)\in A^{\star}\times A}(A^{\star})^{\bb{w}_{a}}\mor A^{\star}.
% %$$
%
% %Therefore leading to an operator:
% %$$
% %\left(\left(\!\begin{smallmatrix}a\\ \cdot\end{smallmatrix}\!\right)(K)\right)_{(K,a)\in \mathrm{Sub}(A^{\star})\times A}\colon\mathrm{Sub}(A^{\star})\times A\mor\mathrm{Sub}(A^{\star})^{\mathrm{Sub}(A^{\star})}
% %$$
% %that assigns to $(K,a)\in \mathrm{Sub}(A^{\star})\times A$ the application $\left(\!\begin{smallmatrix}a\\ \cdot\end{smallmatrix}\!\right)(K)$ from $\mathrm{Sub}(A^{\star})$ in $\mathrm{Sub}(A^{\star})$ which, to an $L\in \mathrm{Sub}(A^{\star})$, assigns the following subset of $A^{\star}$.
% %$$
% %\left(\!\begin{smallmatrix}a\\L\end{smallmatrix}\!\right)(K) = \{\left(\!\begin{smallmatrix}a\$q_{\alpha})_{\alpha\in \bb{w}_{a}}\end{smallmatrix}\!\right)(w)\mid w\in K \And (q_{\alpha})_{\alpha\in \bb{w}_{a}}\in L^{\bb{w}_{a}}\}.
% %$$
%
% %Following \cite{GS84}, on p. 73, Definition 4.3, we define the application
% %$$
% %(\left(\!\begin{smallmatrix}a\\L_{a}\end{smallmatrix}\!\right)_{a\in A})_{(L_{a})_{a\in A}\in \mathrm{Sub}(A^{\star})^{A}}\colon \mathrm{Sub}(A^{\star})^{A}\mor \mathrm{Sub}(A^{\star})^{A^{\star}}.
% %$$
% %Let $L = (L_{a})_{a\in A}\in \mathrm{Sub}(A^{\star})^{A}$. Then the operator $\left(\!\begin{smallmatrix}a\\L_{a}\end{smallmatrix}\!\right)_{a\in A}$ from $A^{\star}$ to $\mathrm{Sub}(A^{\star})$ assigns to $w\in A^{\star}$ the set of all words on $A$ obtained from $w$ by replacing every occurrence of each letter $a$ by a word from the corresponding language $L_{a}$.
%
% %Using the last mapping, we define
% %$$
% %(\left(\!\begin{smallmatrix}a\\L_{a}\end{smallmatrix}\!\right)^{\wp}_{a\in A})_{(L_{a})_{a\in A}\in \mathrm{Sub}(A^{\star})^{A}}\colon \mathrm{Sub}(A^{\star})^{A}\mor \mathrm{Sub}(A^{\star})^{\mathrm{Sub}(A^{\star})}.
% %$$
% %Let $L = (L_{a})_{a\in A}\in \mathrm{Sub}(A^{\star})^{A}$ and $K\in \mathrm{Sub}(A^{\star})$, then
% %$$
% %\textstyle
% %\left(\!\begin{smallmatrix}a\\ L_{a}\end{smallmatrix}\!\right)^{\wp}_{a\in A}(K) = \bigcup_{w\in K}\left(\!\begin{smallmatrix}a\\L_{a}\end{smallmatrix}\!\right)_{a\in A}(w),
% %$$
% %i.e., $\left(\!\begin{smallmatrix}a\\L_{a}\end{smallmatrix}\!\right)^{\wp}_{a\in A}$ is the canonical extension of $\left(\!\begin{smallmatrix}a\\L_{a}\end{smallmatrix}\!\right)_{a\in A}$ to $\mathrm{Sub}(A^{\star})$.
%
% %We also consider the operator:
% %$$
% %\mathrm{Sub}(A^{\star})\mor \mathrm{Hom}(\mathrm{Sub}(A^{\star})^{A},\mathrm{Sub}(A^{\star}))
% %$$
% %that assigns to a $K\in \mathrm{Sub}(A^{\star})$ the mapping $\left(\!\begin{smallmatrix}a\\ \cdot\end{smallmatrix}\!\right)_{a\in A}(K)$ from $\mathrm{Sub}(A^{\star})^{A}$ to $\mathrm{Sub}(A^{\star})$ that sends $L = (L_{a})_{a\in A}\in \mathrm{Sub}(A^{\star})^{A}$ to the subset
% %$$
% %\left(\!\begin{smallmatrix}a\\L_{a}\end{smallmatrix}\!\right)_{a\in A}(K) = \{\left(\!\begin{smallmatrix}a\$q_{\alpha})_{\alpha\in \bb{w}_{a}}\end{smallmatrix}\!\right)_{a\in A}(w)\mid w\in K \And \forall a\in A\, ((q_{\alpha})_{\alpha\in \bb{w}_{a}}\in L_{a}^{\bb{w}_{a}})\}.
% %$$
% %of $A^{\star}$. The words in $\left(\!\begin{smallmatrix}a\\L_{a}\end{smallmatrix}\!\right)_{a\in A}(K)$ are obtained
%
% \subsection{Artinian ordered sets}
%
% In this subsection we recall the notion of Artinian ordered set.
%
% \begin{definition}\label{DPosArt}
% A \emph{preordered} set is an ordered pair $\mathbf{A}=(A,\leq)$ in which $A$ is a set and $\leq$ a reflexive and transitive binary relation on $A$. An \emph{ordered} set is an ordered pair $\mathbf{A}=(A,\leq)$ in which $A$ is a set and $\leq$ a reflexive, antisymmetric and transitive binary relation on $A$. 
% A \emph{strictly ordered} set is an ordered pair $\mathbf{A}=(A,<)$ in which $A$ is a set and $<$ an irreflexive and transitive binary relation on $A$ (recall that, for every set $A$, there exists a  biunivocal correspondence between the set of the orderings on $A$ and the set of the strict orderings on $A$). We will say that an ordered set $\mathbf{A}=(A,\leq)$ is \emph{Artinian} if in $\mathbf{A}$ there is not any strictly decreasing $\omega_{0}$-chain or, what is equivalent, if every nonempty subset of $A$ has a minimal element. Finally, a subset $X$ of $A$ is a \emph{lower subset} of $\mathbf{A}$ if it is downward closed, i.e., if, for every $x\in X$ and every $a\in A$, if $a\leq x$, then $a\in X$.
% \end{definition}
%
% \begin{remark}
% An ordered set $\mathbf{A}$ is Artinian if and only if, for every $\omega_{0}$-indexed family $(a_{n})_{n\in\omega_{0}}$ in $A$, if, for every $n\in\omega_{0}$, $a_{n+1}\leq a_{n}$, then there exists an $m\in\omega_{0}$ such that, for every $p\in \omega_{0}$, $a_{m} = a_{m+p}$. In other words, an ordered set $\mathbf{A}$ is Artinian if and only if every decreasing $\omega_{0}$-chain in $\mathbf{A}$ becomes stationary. Moreover, an ordered set $\mathbf{A}$ is Artinian if and only if, for every $X\subseteq A$, if $\mathrm{Min}(\mathbf{A})$, the set of all minimals of $\mathbf{A}$, is included in $X$ and, for every $a\in A$, if from $\{b\in A\mid b<a\}\subseteq X$ it follows that $a\in X$, then $X=A$. This last condition enables us to carry out not only proofs by Artinian induction but also definitions by Artinian recursion  (see~\cite{agk67} for details). The latter will be fundamental to this work, specifically when we define the $S$-sorted Curry-Howard mappings.
% \end{remark}
%
% \begin{remark}
% Let us recall that in Set Theory (see, e.g.,~\cite{end77}) a binary relation $\Phi$ on a set $A$ is \emph{well-founded} if every nonempty subset $X$ of $A$ has a $\Phi$-minimal element, i.e., there exists an $m\in X$ such that, for every $x\in X$, $(x,m)\nin \Phi$. If the binary relation $\Phi$ on $A$ is well-founded, then it is irreflexive, i.e., for every $a\in A$, $(a,a)\nin \Phi$ (otherwise, $\Phi$ would not be well-founded). On the other hand, the binary relation $\Phi$ on $A$ is well-founded if and only if there is no $\omega_{0}$-family $(a_{n})_{n\in \mathbb{N}}$ in $A$ such that, for every $n\in \mathbb{N}$, $(a_{n+1},a_{n})\in \Phi$. Finally, if the binary relation $\Phi$ on $A$ is well-founded, then $\Phi^{+}$, its transitive closure---which is $\bigcup_{n\in \mathbb{N}-1}\Phi^{n}$---, is also well-founded and, consequently, $\Phi^{+}$ is irreflexive and transitive on $A$ (i.e., $\Phi^{+}$ is an strict order on $A$) and $(A,\Phi)$ is an Artinian ordered set.
% \end{remark}
%
% %\begin{remark}
% %The class of the Artinian partially ordered sets is strictly more wide than the class of the well-ordered sets. Moreover, for such a class of partially ordered sets one has, as for well-ordered sets, both a principle of proof by Artinian induction and a principle of definition by Artinian recursion (see~\cite{agk67} for details).
% %\end{remark}
%
% %Lo que se dice en la siguiente remark no es cierto en general
%
% %\begin{remark}\label{RPosArt}
% %If $\mathbf{A}$ is an Artinian ordered set and $\Phi$ an equivalence relation on $A$, then $\mathrm{pr}^{\Phi}$, the canonical projection of $A$ onto $A/{\Phi}$, has a co-optimal lift (for this notion see~\cite{em76}). In fact, the relation $\leq_{A/\Phi}$ on $A/{\Phi}$ defined, for every $[a]_{\Phi}$, $[b]_{\Phi}$ in $A/{\Phi}$, as
% %$$
% %([a]_{\Phi} {{\leq}_{A/\Phi}} [b]_{\Phi}) \text{ if and only if } \exists \,x, y\in A\, (x\leq y \And [x]_{\Phi} = [a]_{\Phi} \And [y]_{\Phi} = [b]_{\Phi}),
% %$$
% %is an Artinian order structure on $A/{\Phi}$. Thus $\mathbf{A}/{\Phi} = (A/\Phi,\leq_{A/\Phi})$ is an Artinian ordered set.
% %\end{remark}
%
%
% \subsection{Many-sorted algebras}
%
% Our next aim is to provide those notions from the field of many-sorted universal algebra that will be used afterwards.
% We will specially focus on the constructive description of the subalgebra generating many-sorted operator, the congruence generating many-sorted operator, the many-sorted translations, the (absolutely) free many-sorted algebra on a many-sorted set, and the substitution operators associated to a free many-sorted algebra.
%
% \begin{convention} In what follows, for a set of sorts $S$, an arbitrary word on $S^{\star}$ will be denoted by $\mathbf{s}$, i.e., a lower case bold type $s$. The letter $s$ will be used to represent an arbitrary letter in $S$.
% \end{convention}
%
% \begin{definition}\label{DSig}
% An $S$-\emph{sorted signature} is a mapping $\Sigma$ from $S^{\star}\times S$ to $\boldsymbol{\mathcal{U}}$ which sends a pair $(\mathbf{s},s)\in S^{\star}\times S$ to the set $\Sigma_{\mathbf{s},s}$ of the \emph{formal operations} of \emph{arity} $\mathbf{s}$, \emph{sort} (or \emph{coarity}) $s$, and \emph{rank} (or \emph{biarity}) $(\mathbf{s},s)$.
% %Sometimes we will write $\sigma\colon w\mor s$ to indicate that the formal operation $\sigma$ belongs to
% %$\Sigma_{w,s}$.
% %If $\Sigma$ and $\Lambda$ are $S$-sorted signatures, then a \emph{morphism from} $\Sigma$ \emph{to} $\Lambda$ is an $S^{\star}\times S$-indexed family $d = (d_{\mathbf{s},s})_{(\mathbf{s},s)\in S^{\star}\times S}$, where, for every $(\mathbf{s},s)\in S^{\star}\times S$, $d_{\mathbf{s},s}$ is a mapping from $\Sigma_{\mathbf{s},s}$ to $\Lambda_{\mathbf{s},s}$. $S$-sorted signatures and morphisms between $S$-sorted signatures form a category which we will denote henceforth by $\mathsf{Sig}(S)$.
%
% %From now on, $\mathsf{Sig}(S)$ stands for the category of $S$-sorted signatures and morphisms between $S$-sorted signatures.
% \end{definition}
%
% %\begin{remark}
% %For every set of sorts $S$, the category $\mathsf{Sig}(S)$ is $\mathsf{Set}^{S^{\star}\times S}$.
% %\end{remark}
%
% %%%%%%%%%%%%%%%%%%%%%%%%%%%%%%%%
% %%%%%%%%%%%%%%%%%%%%%%%%%%%%%%%%
%
%
% \begin{assumption}
% From now on $\Sigma$ stands for an $S$-sorted signature, fixed once and for all.
% \end{assumption}
%
% %\begin{remark}
% %To give an $S$-sorted signature as above is equivalent to give an ordered triple $(\Sigma,\mathrm{ar},\mathrm{car})$ where $\Sigma$ is a set (of operation symbols), $\mathrm{ar}$ a mapping from $\Sigma$ to $S^{\star}$ (the arity mapping), and $\mathrm{car}$ a mapping from $\Sigma$ to $S$ (the the coarity mapping). More accurately, the category $\mathsf{Set}^{S^{\star}\times S}$ is equivalent to the category $\mathsf{Sig}(S)$ whose objects are the just considered ordered triples $(\Sigma,\mathrm{ar},\mathrm{car})$ and in which the set of morphisms from $(\Sigma,\mathrm{ar},\mathrm{car})$ to $(\Sigma',\mathrm{ar}',\mathrm{car}')$ is the set of all mappings $d$ from $\Sigma$ to $\Sigma'$ such that $\mathrm{ar}'\circ d = \mathrm{ar}$ and $\mathrm{car}'\circ d = \mathrm{car}$. In all work with many-sorted signatures, we shall feel free to deal directly either with an $S$-sorted signature or with the correlated ordered triple, whichever is most convenient for the work at hand.
% %
% %%Therefore, although both notions, literally, are different, we will use them indistinctly.
% %\end{remark}
%
% We shall now give precise definitions of the concepts of many-sorted algebra and homomorphism between many-sorted algebras.
%
% \begin{definition}\label{DAlg}
% The $S^{\star}\times S$-sorted set of the \emph{finitary operations on} an $S$-sorted set $A$ is $(\mathrm{Hom}(A_{\mathbf{s}},A_{s}))_{(\mathbf{s},s)\in S^{\star}\times S}$, where, for every $\mathbf{s}\in S^{\star}$, $A_{\mathbf{s}} = \prod_{j\in \bb{\mathbf{s}}}A_{s_{j}}$, with $\bb{\mathbf{s}}$ denoting the length of the word $\mathbf{s}$ (if $\mathbf{s} = \lambda$, then $A_{\lambda}$ is a final set). Sometimes we let $\mathrm{O}_{\mathrm{H}_{S}}(A)$ stand for 
% $(\mathrm{Hom}(A_{\mathbf{s}},A_{s}))_{(\mathbf{s},s)\in S^{\star}\times S}$. A \emph{structure of} $\Sigma$-\emph{algebra on} an $S$-sorted set  $A$ is a family $(F_{\mathbf{s},s})_{(\mathbf{s},s)\in S^{\star}\times S}$, denoted by $F$, where, for $(\mathbf{s},s)\in S^{\star}\times S$, $F_{\mathbf{s},s}$ is a mapping from $\Sigma_{\mathbf{s},s}$ to $\mathrm{Hom}(A_{\mathbf{s}},A_{s})$ (if $(\mathbf{s},s) = (\lambda,s)$ and $\sigma\in \Sigma_{\lambda,s}$, then $F_{\lambda,s}(\sigma)$ picks out an element of $A_{s}$).  A \emph{many-sorted} $\Sigma$-\emph{algebra} (or, to abbreviate, $\Sigma$-\emph{algebra}) is a pair $(A,F)$, denoted by $\mathbf{A}$, where $A$ is an $S$-sorted set and $F$ a structure of $\Sigma$-algebra on $A$. For a pair $(\mathbf{s},s)\in S^{\star}\times S$ and a formal operation $\sigma\in \Sigma_{\mathbf{s},s}$, in order to simplify the notation, the operation $F_{\mathbf{s},s}(\sigma)$ from $A_{\mathbf{s}}$ to $A_{s}$ will be written as $F_{\sigma}$. In some cases, to avoid mistakes, we will denote by $F^{\mathbf{A}}$ the structure of $\Sigma$-algebra on $A$, and, for $(\mathbf{s},s)\in S^{\star}\times S$ and $\sigma\in \Sigma_{\mathbf{s},s}$, by $F^{\mathbf{A}}_{\sigma}$, or simply by $\sigma^{\mathbf{A}}$, the corresponding operation. Moreover, for $s\in S$ and $\sigma\in\Sigma_{\lambda,s}$, we will, usually, denote by $\sigma^{\mathbf{A}}$ the value of the mapping $F^{\mathbf{A}}_{\sigma}\colon A_{\lambda}\mor A_{s}$ at the unique element in $A_{\lambda}$.
%
% A $\Sigma$-\emph{homomorphism} (or, to abbreviate, \emph{homomorphism}) from $\mathbf{A}$ to $\mathbf{B}$, where $\mathbf{B} = (B,G)$, is a triple $(\mathbf{A},f,\mathbf{B})$, denoted by $f\colon \mathbf{A}\mor \mathbf{B}$, where $f$ is an $S$-sorted mapping from $A$ to $B$ such that, for every $(\mathbf{s},s)\in S^{\star}\times S$, every  $\sigma\in \Sigma_{\mathbf{s},s}$, and every $(a_{j})_{j\in \bb{\mathbf{s}}}\in A_{\mathbf{s}}$, we have 
% $
% f_{s}(\sigma^{\mathbf{A}}((a_{j})_{j\in \bb{\mathbf{s}}})) = \sigma^{\mathbf{B}}(f_{\mathbf{s}}((a_{j})_{j\in \bb{\mathbf{s}}})),
% $
% where $f_{\mathbf{s}}$ is the mapping $\prod_{j\in \bb{\mathbf{s}}}f_{s_{j}}$ from $A_{\mathbf{s}}$ to $B_{\mathbf{s}}$ that sends $(a_{j})_{j\in \bb{\mathbf{s}}}$ in $A_{\mathbf{s}}$ to $(f_{s_{j}}(a_{j}))_{j\in \bb{\mathbf{s}}}$ in $B_{\mathbf{s}}$. For a $\Sigma$-algebra $\mathbf{A}$ the triple $(\mathbf{A},\id^{A},\mathbf{A})$, where $\id^{A}$ is the identity at $A$, which is a homomorphism from $\mathbf{A}$ to $\mathbf{A}$, will be denoted by $\id^{\mathbf{A}}$ and will be called the \emph{identity} homomorphism at $\mathbf{A}$. We will denote by $\mathsf{Alg}(\Sigma)$ the category of $\Sigma$-algebras and homomorphisms and by $\mathrm{Alg}(\Sigma)$ the set of objects  of $\mathsf{Alg}(\Sigma)$.
%
%
%
% We will denote by $\mathbf{1}^{S}$ or, to abbreviate, by $\mathbf{1}$, the (standard) final $\Sigma$-algebra.
% \end{definition}
%
% We add some further results concerning direct products and homomorphisms, which will be of great importance later on.
%
% \begin{definition}\label{DAlgProd}
% Let $(\mathbf{A}^{i})_{i\in I}$ be a family of $\Sigma$-algebras, where we agree that, for every $i\in I$, $\mathbf{A}^{i}=(A^{i},F^{i})$. Then
% \begin{enumerate}
% \item The \emph{product}
%       of $(\mathbf{A}^{i})_{i\in I}$, $\prod_{i\in I}\mathbf{A}^{i}$, is
%       the  $\Sigma$-algebra which has as $S$-sorted underlying set $\prod_{i\in I}A^{i}$, 
%       and where, for every $(\mathbf{s},s)\in S^{\star}\times S$ and every $\sigma\in \Sigma_{\mathbf{s},s}$,
%       the structural operation $\sigma^{\prod_{i\in I}\mathbf{A}^{i}}$ is defined as follows:
%       $$
%       \sigma^{\prod_{i\in I}\mathbf{A}^{i}}
%       \nfunction
%       { (\prod_{i\in I}A^{i})_{\mathbf{s}}}
%       {\prod_{i\in I}A^{i}_{s}}
%       {(a_{j})_{j\in \bb{\mathbf{s}}}}
%       {(\sigma^{\mathbf{A}_{i}}((a_{j}(i))_{j\in \bb{\mathbf{s}}})_{i\in I}}.
%       $$
% \item For every $i\in I$, the $i$-th \emph{canonical projection}
%       is the homomorphism from $\prod_{i\in I}\mathbf{A}^{i}$ to $\mathbf{A}^{i}$
%       determined by the $S$-sorted mapping $\pr^{i}$ which, for every $s\in S$, is
%       defined as follows
%       $$
%       \pr^{i}_{s}
%       \nfunction
%       {\prod_{i\in I}A^{i}_{s}}
%       {A^{i}_{s}}
%       {(a_{i})_{i\in I}}
%       {a_{i}}
%       $$
% \end{enumerate}
% \end{definition}
%
% \begin{proposition}\label{PAlgProd}
% Let $(\mathbf{A}^{i})_{i\in I}$ be a family of $\Sigma$-algebras.
% Then the pair $(\prod_{i\in I}\mathbf{A}^{i},(\pr^{i})_{i\in I})$
% is a product in $\mathsf{Alg}(\Sigma)$.
% \end{proposition}
%
% We next introduce the support of a many-sorted algebra and the definition of finiteness of a many-sorted algebra.
%
% \begin{definition}
% Let $\mathbf{A}$ be a $\Sigma$-algebra. Then the \emph{support of} $\mathbf{A}$, denoted by $\mathrm{supp}_{S}(\mathbf{A})$, is $\mathrm{supp}_{S}(A)$, i.e., the support of the underlying $S$-sorted set $A$ of $\mathbf{A}$.
% \end{definition}
%
% \begin{remark}
% The set $\{\mathrm{supp}_{S}(\mathbf{A})\mid \mathbf{A}\in \mathrm{Alg}(\Sigma)\}$ is a closure system on $S$.
% \end{remark}
%
% \begin{definition}
% Let $\mathbf{A}$ be a $\Sigma$-algebra. We will say that $\mathbf{A}$ is \emph{finite} if $A$, the underlying $S$-sorted set of $\mathbf{A}$, is finite.
% \end{definition}
%
% \begin{remark}
% In $\mathsf{Alg}(\Sigma)$, as it was the case with $\mathsf{Set}^{S}$, there is another notion of finiteness: A $\Sigma$-algebra  $\mathbf{A}$ is called $S$-\emph{finite} or \emph{locally finite}, abbreviated as $S\text{-}\mathrm{f}$, if and only if the underlying $S$-sorted set of $\mathbf{A}$ is $S$-finite. As was noted above this notion of finiteness plays a relevant role in the field of many-sorted algebra, e.g., to define $S$-finite, also called locally finite, terms and to distinguish, on the one hand, between many-sorted varieties and finitary many-sorted  varieties, and, on the other hand, between many-sorted quasivarieties and finitary many-sorted quasivarieties (see, e.g., \cite{gm85} and \cite{{m76}}).
% \end{remark}
%
% %We next define the notion of subelement preorder on a many-sorted algebra.
% %
% %\begin{definition}\label{DOrd}
% %Let $\mathbf{A}=(A,F)$ be a $\Sigma$-algebra. Then $<_{\mathbf{A}}$ denotes the binary relation on $\coprod A$ consisting of the ordered pairs $((b,t),(a,s))\in(\coprod A)^{2}$ for which there exists a word $\mathbf{s}\in S^{\star}-\{\lambda\}$, an operation symbol
% %$\sigma\in\Sigma_{\mathbf{s},s}$, and a family of elements $(x_{j})_{j\in\bb{\mathbf{s}}}\in A_{\mathbf{s}}$ such that $\sigma^{\mathbf{A}}((x_{j})_{j\in\bb{\mathbf{s}}}) = a$ and, for some $j\in \bb{\mathbf{s}}$, $s_{j}=t$ and  $x_{j}=b$. We will denote by $\leq_{\mathbf{A}}$ the reflexive and transitive closure of $<_{\mathbf{A}}$, i.e., the preorder on $\coprod A$  generated by $<_{\mathbf{A}}$.  If $((b,t),(a,s))\in \leq_{\mathbf{A}}$ we will say that $b$ is a \emph{subelement} of $a$. The preorder $\leq_{\mathbf{A}}$ will be called the \emph{subelement preorder} on $\mathbf{A}$.
% %\end{definition}
%
% %\begin{definition}
% %Let $\mathbf{A}=(A,F)$ be a partial $\Sigma$-algebra. Then $<_{\mathbf{A}}$ denotes the binary relation on $\coprod A$ consisting of the ordered pairs $((b,t),(a,s))\in(\coprod A)^{2}$ for which there exists a $\mathbf{s}\in S^{\star}-\{\lambda\}$, a
% %$\sigma\in\Sigma_{\mathbf{s},s}$, and an $x\in \Dom(F_{\sigma})$ such that $F_{\sigma}(x) = a$ and, for some $j\in \bb{\mathbf{s}}$, $s_{j}=t$ and  $x_{j}=b$. We will denote by $\leq_{\mathbf{A}}$ the reflexive and transitive closure of $<_{\mathbf{A}}$, i.e., the preorder on $\coprod A$  generated by $<_{\mathbf{A}}$.
% %\end{definition}
%
% %\begin{remark}
% %The subelement preorder $\leq_{\mathbf{A}}$ on $\coprod A$ is defined by letting $((b,t),(a,s))\in \leq_{\mathbf{A}}$ mean that $s = t$ and $a = b$ or there exists an $n\in \mathbb{N}-\{0\}$, a $\mathbf{c}\in S^{\star}$, and a family $(p_{i})_{i\in \bb{\mathbf{c}}}\in A_{\mathbf{c}}$ such that $\bb{\mathbf{c}} = n+1$, $p_{0} = b$, $p_{n} = a$, $c_{0} = t$, $c_{n} = s$, and, for every $i\in n$, $((p_{i},c_{i}),(p_{i+1},c_{i+1}))\in <_{\mathbf{A}}$.
% %\end{remark}
%
%
%
%
% %We next define the subset many-sorted algebra associated to a many-sorted algebra and prove, in particular, that this is the object mapping of an endofunctor of $\mathsf{Alg}(\Sigma)$. We note that, in Section~3, the subset many-sorted algebra associated to a free many-sorted algebra will appear as the codomain of a substitution operator which will be used to prove recognizability results.
%
% %\begin{proposition}\label{subsetalg}
% %Let $(\cdot)^{\wp}$ be the mapping that sends (1) a $\Sigma$-algebra $\mathbf{A}$ to the $\Sigma$-algebra $\mathbf{A}^{\wp} = (A^{\wp},F^{\wp})$ where $A^{\wp}$ is $(\mathrm{Sub}(A_{s}))_{s\in S}$ and where, for every $(w,s)\in S^{\star}\times S$ and every $\sigma\in \Sigma_{w,s}$, $F_{\sigma}^{\wp}$ is the mapping from $A^{\wp}_{w} = \prod_{i\in \bb{w}}\mathrm{Sub}(A_{w_{i}})$ to $\mathrm{Sub}(A_{s})$ that sends $(L_{i})_{i\in \bb{w}}$ in $A^{\wp}_{w}$ to  $\{F_{\sigma}((x_{i})_{i\in \bb{w}})\mid (x_{i})_{i\in \bb{w}}\in \prod_{i\in \bb{w}}L_{i}\}$ in $\mathrm{Sub}(A_{s})$, and (2) a homomorphism $f$ from $\mathbf{A}$ to $\mathbf{B}$ to the $S$-sorted mapping $f^{\wp} = (f_{s}[\cdot])_{s\in S}$ from $A^{\wp}$ to $B^{\wp}$. Then $(\cdot)^{\wp}$ is an endofunctor of $\mathsf{Alg}(\Sigma)$. The $\Sigma$-algebra $\mathbf{A}^{\wp}$ is called the \emph{subset algebra} associated to $\mathbf{A}$ (this notion, but for single-sorted algebras, is due to Mezei and Wright, cf.~\cite{MW67}, Definition 2.2). Moreover, there exists a pointwise monomorphic natural transformation $\{\cdot\}^{\Sigma}$ from $\Id_{\mathsf{Alg}(\Sigma)}$ to $(\cdot)^{\wp}$, displayed as:
% %$$
% %\xymatrix@C=20ex@R=8ex{
% %\mathsf{Alg}(\Sigma)
% %\ar@/^20pt/[r]^{\Id_{\mathsf{Alg}(\Sigma)}}="f"
% %  \ar@/^-20pt/[r]_{(\cdot)^{\wp}}="g" &
% %\mathsf{Alg}(\Sigma)
% %\ar @{} "f";"g" |{\dir{=>}}^{\,\{\cdot\}^{\Sigma}}
% %}
% %$$
% %\end{proposition}
% %
% %\begin{proof}
% %To show that $(\cdot)^{\wp}$ is an endofunctor of $\mathsf{Alg}(\Sigma)$ it suffices to verify that, for every homomorphism $f\colon \mathbf{A}\mor \mathbf{B}$, $f^{\wp}$ is, actually, a homomorphism from $\mathbf{A}^{\wp}$ to $\mathbf{B}^{\wp}$. Let $f$ be a homomorphism from $\mathbf{A}$ to $\mathbf{B}$, $(w,s)\in S^{\star}\times S$, and $\sigma\in \Sigma_{w,s}$. Then the following diagram commutes
% %$$
% %\xymatrix@C=60pt{
% %A^{\wp}_{w} \ar[d]_-{F_{\sigma}^{\wp}}\ar[r]^-{f_{{w}}[\cdot]}&
% %B^{\wp}_{w} \ar[d]^-{G_{\sigma}^{\wp}}\\
% %\mathrm{Sub}(A_{s})\ar[r]_-{f_{s}[\cdot]}&\mathrm{Sub}(B_{s})
% %}
% %$$
% %that is, for every sequence $(L_{i})_{i\in \bb{w}}$ in $A^{\wp}_{w}$, we have 
% %$$
% %f_{s}[F_{\sigma}^{\wp}((L_{i})_{i\in \bb{w}})] = G_{\sigma}^{\wp}((f_{w_{i}}[L_{i}])_{i\in \bb{w}}).
% %$$
% %
% %Now let $\{\cdot\}^{\Sigma}$ be the mapping from $\mathrm{Alg}(\Sigma)$ to $\mathrm{Mor}(\mathsf{Alg}(\Sigma))$ that assigns to a $\Sigma$-algebra $\mathbf{A}$ the $S$-sorted mapping $\{\cdot\}^{\Sigma}_{\mathbf{A}}$ from $A$ to $A^{\wp}$ that, for every $s\in S$, sends $a$ in $A_{s}$ to $\{a\}$ in $A^{\wp}_{s}$. It is easily seen that, for every $\Sigma$-algebra $\mathbf{A}$, $\{\cdot\}^{\Sigma}_{\mathbf{A}}$ is an injective homomorphism from $\mathbf{A}$ to $\mathbf{A}^{\wp}$ and that $\{\cdot\}^{\Sigma} = (\{\cdot\}^{\Sigma}_{\mathbf{A}})_{\mathbf{A}\in \mathrm{Alg}(\Sigma)}$ is, in fact, a natural transformation from $\Id_{\mathsf{Alg}(\Sigma)}$ to $(\cdot)^{\wp}$
% %\end{proof}
%
%
% %To state the principle of proof by algebraic induction we need the concept of generation which in turn needs the one of a closed subset of an algebra.
%
% \subsection{Subalgebras}\label{SAlgSub}
% We shall now go on to define the notion of subalgebra of a $\Sigma$-algebra $\mathbf{A}$, the principle of proof by Algebraic Induction, and the subalgebra generating operator for $\mathbf{A}$. The constructive description of the subalgebra generating operator for $\mathbf{A}$ will allow us to introduce the Borel subsets of a many-sorted algebra and the height of an element with respect to a given subset. This latter notion will be compared with the subelement preorder on $\mathbf{A}$.
%
% \begin{definition}\label{DAlgSub}
% Let $\mathbf{A}$ be a $\Sigma$-algebra and $X\subseteq A$. Given $(\mathbf{s},s)\in S^{\star}\times S$ and $\sigma\in\Sigma_{\mathbf{s},s}$, we will say that $X$ is \emph{closed under the operation} $\sigma^{\mathbf{A}}\colon A_{\mathbf{s}}\mor A_{s}$ if, for every $(a_{j})_{j\in\bb{\mathbf{s}}}\in X_{\mathbf{s}}$, $\sigma^{\mathbf{A}}((a_{j})_{j\in\bb{\mathbf{s}}})\in X_{s}$. We will say that $X$ is a \emph{closed subset} of $\mathbf{A}$ if $X$ is closed under the operations of $\mathbf{A}$. We will denote by $\mathrm{Cl}(\mathbf{A})$ the set of all closed subsets of $\mathbf{A}$ (which is an algebraic closure system on $A$) and by  $\mathbf{Cl}(\mathbf{A})$ the algebraic lattice $(\mathrm{Cl}(\mathbf{A}),\subseteq)$. We will say that a $\Sigma$-algebra $\mathbf{B}$ is a \emph{subalgebra} of $\mathbf{A}$ if $B\subseteq A$ and the canonical embedding of $B$ into $A$ determines an embedding of $\mathbf{B}$ into $\mathbf{A}$. We will denote by $\mathrm{Sub}(\mathbf{A})$ the set of all subalgebras of $\mathbf{A}$. Since $\mathrm{Cl}(\mathbf{A})$ and $\mathrm{Sub}(\mathbf{A})$ are isomorphic, we shall feel free to deal either with a closed subset of $\mathbf{A}$ or with the correlated subalgebra of $\mathbf{A}$, whichever is most convenient for the work at hand.
% \end{definition}
%
%
%
% \begin{definition}
% Let $\mathbf{A}$ be a $\Sigma$-algebra. Then we denote by $\mathrm{Sg}_{\mathbf{A}}$ the many-sorted closure operator on $A$ defined as follows:
% $$\textstyle
%   \mathrm{Sg}_{\mathbf{A}}\nfunction
%   {\mathrm{Sub}(A)}
%   {\mathrm{Sub}(A)}
%   {X}
%   {\bigcap \{\,C\in\mathrm{Sub}(\mathbf{A})\mid X\subseteq C\,\}}
% $$
%
% We call $\mathrm{Sg}_{\mathbf{A}}$ the \emph{subalgebra generating many-sorted operator on} $A$ \emph{determined by} $\mathbf{A}$. For every $X\subseteq A$, we call $\mathrm{Sg}_{\mathbf{A}}(X)$ the \emph{subalgebra of} $\mathbf{A}$ \emph{generated by} $X$. Moreover, if $X\subseteq A$ is such that  $\mathrm{Sg}_{\mathbf{A}}(X) = A$, then we say that $X$ is an $S$-sorted set of \emph{generators} of $\mathbf{A}$, or that $X$ \emph{generates} $\mathbf{A}$. Besides, we say that $\mathbf{A}$ is \emph{finitely generated} if there exists an $S$-sorted  subset $X$ of $A$ such that $X$ \emph{generates} $\mathbf{A}$ and $\mathrm{card}(X)<\aleph_{0}$.
% \end{definition}
%
% \begin{remark}
% Let $\mathbf{A}$ be a $\Sigma$-algebra. Then the algebraic closure operator $\mathrm{Sg}_{\mathbf{A}}$ is uniform, i.e., for every  $X$, $Y\subseteq A$, if $\mathrm{supp}_{S}(X) = \mathrm{supp}_{S}(Y)$, then we have $\mathrm{supp}_{S}(\mathrm{Sg}_{\mathbf{A}}(X)) = \mathrm{supp}_{S}(\mathrm{Sg}_{\mathbf{A}}(Y))$.
% \end{remark}
%
%
% We next recall the Principle of Proof by Algebraic Induction. %This principle will be used in Section~3.
%
% \begin{proposition}[Principle of Proof by Algebraic Induction]\label{PPAI}
% Let $\mathbf{A}$ be a  $\Sigma$-alge\-bra generated by $X$. Then to prove that a subset $Y$ of $A$ is equal to $A$ it suffices to show: (1) $X\subseteq Y$ (algebraic induction basis) and (2) $Y$ is a subalgebra of $\mathbf{A}$ (algebraic induction step).
% \end{proposition}
%
% \begin{proposition}\label{PAlgSg}
% Let $\mathbf{A}$ be a $\Sigma$-algebra. Then the many-sorted closure operator $\mathrm{Sg}_{\mathbf{A}}$ on $A$ is algebraic, i.e., for every $S$-sorted subset $X$ of $A$, $\mathrm{Sg}_{\mathbf{A}}(X) =
% \bigcup_{K\subseteq_{\mathrm{f}} X}\mathrm{Sg}_{\mathbf{A}}(K)$.
% \end{proposition}
%
% For a $\Sigma$-algebra $\mathbf{A}$ we next provide another, more constructive, description of the algebraic many-sorted closure operator $\Sg_{\mathbf{A}}$, which will allow us to define, for a $\Sigma$-algebra $\mathbf{A}$ and a subset $X$ of $A$, the Borel sets of $X$ in $\mathbf{A}$.
%
% \begin{definition}
% Let $\Sigma$ be an $S$-sorted signature and $\mathbf{A}$ a $\Sigma$-algebra.
% \begin{enumerate}
% \item We denote by $\mathrm{E}_{\mathbf{A}}$ the many-sorted operator on
%       $A$ that assigns to an $S$-sorted subset $X$ of
%       $A$, $\mathrm{E}_{\mathbf{A}}(X) =
%       X\cup\bigl(\;\bigcup_{\sigma\in\Sigma_{\neq\lambda,s}}
%       \sigma^{\mathbf{A}}[X_{\mathrm{ar}(\sigma)}]\bigr)_{s\in S}$, where, for $s\in
%       S$, $\Sigma_{\neq\lambda,s}$ is the set of all many-sorted formal
%       operations $\sigma$ such that $\mathrm{ar}(\sigma)\in S^{\star}-\{\lambda\}$,
%       $\mathrm{car}(\sigma) = s$, and if $\mathrm{ar}(\sigma) = \mathbf{s}$, then
%       $X_{\mathrm{ar}(\sigma)} = \prod_{i\in \bb{\mathbf{s}}}X_{s_{i}}$.
%
% \item If $X\subseteq A$, then we define the family
%       $(\mathrm{E}^{n}_{\mathbf{A}}(X))_{n\in \mathbb{N}}$ in $\mathrm{Sub}(A)$,
%       recursively, as follows:
%       \begin{align*}
%       \mathrm{E}_{\mathbf{A}}^{0}(X) &=
%       \textstyle
%       X\cup (\bigcup_{\sigma\in \Sigma_{\lambda,s}}\sigma^{\mathbf{A}}[A_{\lambda}])_{s\in S}\text{,} \\
%       \mathrm{E}_{\mathbf{A}}^{n+1}(X) &=
%       \mathrm{E}_{\mathbf{A}}(\mathrm{E}_{\mathbf{A}}^{n}(X)) \text{, $n\geq 0$.}
%       \end{align*}
% \item We denote by $\mathrm{E}_{\mathbf{A}}^{\omega}$ the many-sorted operator on
%       $A$ that assigns to an $S$-sorted subset $X$ of
%       $A$, $\mathrm{E}_{\mathbf{A}}^{\omega}(X) = \bigcup_{n\in
%       \mathbb{N}}\mathrm{E}_{\mathbf{A}}^{n}(X)$.
%       %defined as:
% %     $$\mathrm{E}_{\mathbf{A}}^{\omega}
% %               \nfunction{\mathrm{Sub}(A)}{\mathrm{Sub}(A)}
% %               {X}
% %               {\union\nolimits_{n\in \mathbb{N}}\mathrm{E}_{\mathbf{A}}^{n}(X)}
% %       $$
% \end{enumerate}
% \end{definition}
%
% \begin{remark}
% Since, for $s\in S$ and $\sigma\in\Sigma_{\lambda,s}$, $\sigma^{\mathbf{A}}$ is the value of $F^{\mathbf{A}}_{\sigma}\colon A_{\lambda}\mor A_{s}$ at the unique element in $A_{\lambda}$, it follows  that
% $$
% \textstyle
% \bigl(\;\bigcup_{\sigma\in\Sigma_{\lambda,s}}\sigma^{\mathbf{A}}[A_{\lambda}]\bigr)_{s\in S} =
% (\{\sigma^{\mathbf{A}}\mid \sigma\in \Sigma_{\lambda,s}\})_{s\in S}
% $$
% \end{remark}
%
% \begin{proposition}\label{PSgE}
% Let $\mathbf{A}$ be a $\Sigma$-algebra and $X\subseteq A$, then $\Sg_{\mathbf{A}}(X) = \mathrm{E}_{\mathbf{A}}^{\omega}(X)$.
% \end{proposition}
%
% \begin{proof}
% Let us first prove that $\mathrm{Sg}_{\mathbf{A}}(X)\subseteq \mathrm{E}_{\mathbf{A}}^{\omega}(X)$. For this, because, by definition, $\mathrm{Sg}_{\mathbf{A}}(X)$ is the smallest subalgebra of $\mathbf{A}$ containing $X$, it suffices to verify  that $\mathrm{E}_{\mathbf{A}}^{\omega}(X)$ contains $X$ and that $\mathrm{E}_{\mathbf{A}}^{\omega}(X)$ is a subalgebra of $\mathbf{A}$. 
%
% Since, by definition, $\mathrm{E}_{\mathbf{A}}^{0}(X) = X\cup (\bigcup_{\sigma\in \Sigma_{\lambda,s}}\sigma^{\mathbf{A}}[A_{\lambda}])_{s\in S}$, we have that
% $X\subseteq \mathrm{E}_{\mathbf{A}}^{\omega}(X)$. Moreover, (1) if $s\in S$ and $\sigma\in\Sigma_{\lambda,s}$, then, since 
% $$
% \textstyle
% \mathrm{E}_{\mathbf{A}}^{0}(X) 
% = X\cup 
% (\bigcup_{\sigma\in \Sigma_{\lambda,s}}\sigma^{\mathbf{A}}[A_{\lambda}])_{s\in S},$$
% we have that $\sigma^{\mathbf{A}}\in \mathrm{E}_{\mathbf{A}}^{\omega}(X)_{s}$, and (2) if $\mathbf{s}\in S^{\star}-\{\lambda\}$, $s\in S$, $\sigma\in \Sigma_{\mathbf{s},s}$, and $(a_{j})_{j\in\bb{\mathbf{s}}}\in \mathrm{E}_{\mathbf{A}}^{\omega}(X)_{\mathbf{s}}$, then, for every $j\in \bb{\mathbf{s}}$, there exists an $n_{j}\in \mathbb{N}$ such that $a_{j}\in \mathrm{E}_{\mathbf{A}}^{n_{j}}(X)_{s_{i}}$. But the family
% $(\mathrm{E}^{n}_{\mathbf{A}}(X))_{n\in \mathbb{N}}$ is an ascending chain, thus there exists a $n\in \bb{\mathbf{s}}$ such that, for every $j\in \bb{\mathbf{s}}$, $\mathrm{E}_{\mathbf{A}}^{n_{j}}(X)\subseteq
% \mathrm{E}_{\mathbf{A}}^{n}(X)$. Therefore, for every $j\in \bb{\mathbf{s}}$, $a_{j}\in \mathrm{E}_{\mathbf{A}}^{n_{k}}(X)_{s_{j}}$, hence ${\sigma}^{\mathbf{A}}((a_{j})_{j\in\bb{\mathbf{s}}})\in
% \mathrm{E}_{\mathbf{A}}^{n+1}(X)_{s}$. Whence ${\sigma}^{\mathbf{A}}((a_{j})_{j\in\bb{\mathbf{s}}})\in \mathrm{E}_{\mathbf{A}}^{\omega}(X)_{s}$. Consequently, $\mathrm{Sg}_{\mathbf{A}}(X)\subseteq \mathrm{E}_{\mathbf{A}}^{\omega}(X)$.
%
% We next prove, by finite induction, that $\mathrm{E}_{\mathbf{A}}^{\omega}(X)\subseteq
% \mathrm{Sg}_{\mathbf{A}}(X)$. 
%
% \textsf{Base step of the induction.}
%
% Since, by definition, $\mathrm{E}_{\mathbf{A}}^{0}(X) =  X\cup (\bigcup_{\sigma\in \Sigma_{\lambda,s}}\sigma^{\mathbf{A}}[A_{\lambda}])_{s\in S}$, 
% we, obviously, have that $\mathrm{E}_{\mathbf{A}}^{0}(X)\subseteq\mathrm{Sg}_{\mathbf{A}}(X)$. 
%
% \textsf{Inductive step of the induction.}
%
% Let us suppose that, for $n\geq 0$, the requirement $\mathrm{E}_{\mathbf{A}}^{n}(X)\subseteq \mathrm{Sg}_{\mathbf{A}}(X)$ is satisfied. Then, since, by definition,   $\mathrm{E}_{\mathbf{A}}^{n+1}(X) = \mathrm{E}_{\mathbf{A}}(\mathrm{E}_{\mathbf{A}}^{n}(X))$, to prove that
% $\mathrm{E}_{\mathbf{A}}^{n+1}(X)\subseteq \mathrm{Sg}_{\mathbf{A}}(X)$, it suffices to verify that $\mathrm{E}_{\mathbf{A}}^{n}(X)\subseteq\mathrm{Sg}_{\mathbf{A}}(X)$ and that 
% $$
% \textstyle
% (\bigcup_{\sigma\in\Sigma_{\neq\lambda,s}}
% \sigma^{\mathbf{A}}[\mathrm{E}_{\mathbf{A}}^{n}(X)_{\mathrm{ar}(\sigma)}])_{s\in S}
% \subseteq\mathrm{Sg}_{\mathbf{A}}(X).$$
%
% That $\mathrm{E}_{\mathbf{A}}^{n}(X)\subseteq\mathrm{Sg}_{\mathbf{A}}(X)$ follows directly from the induction hypothesis. Let $(\mathbf{s}, s)$ be an element of $(S^{\star}-\{\lambda\})\times S$, $\sigma$ an operation symbol in $\Sigma_{\mathbf{s},s}$, and $(a_{j})_{j\in\bb{\mathbf{s}}}\in
% \mathrm{E}_{\mathbf{A}}^{n}(X)_{\mathbf{s}}$. Then,
% $(a_{j})_{j\in\bb{\mathbf{s}}}\in \mathrm{Sg}_{\mathbf{A}}(X)_{\mathbf{s}}$, thus ${\sigma}^{\mathbf{A}}((a_{j})_{j\in\bb{\mathbf{s}}})\in \mathrm{Sg}_{\mathbf{A}}(X)_{s}$. Therefore $\sigma^{\mathbf{A}}[\mathrm{E}_{\mathbf{A}}^{n}(X)_{\mathrm{ar}(\sigma)}]\subseteq\mathrm{Sg}_{\mathbf{A}}(X)_{s}$. Hence $\bigcup_{\sigma\in\Sigma_{\neq\lambda,s}}\sigma^{\mathbf{A}}	[\mathrm{E}_{\mathbf{A}}^{n}(X)_{\mathrm{ar}(\sigma)}]\subseteq \mathrm{Sg}_{\mathbf{A}}(X)_{s}$. Consequently, $\mathrm{E}_{\mathbf{A}}^{\omega}(X)\subseteq\mathrm{Sg}_{\mathbf{A}}(X)$.
% \end{proof}
%
%
% The last proposition allows us to define the height of an element of type $s\in S$ in $\mathrm{Sg}_{\mathbf{A}}(X)_{s}$ as the smallest $n\in \mathbb{N}$ for which the element belongs to $\mathrm{E}_{\mathbf{A}}^{n}(X)_{s}$.
%
% \begin{definition} Let $\mathbf{A}$ be a $\Sigma$-algebra, $X\subseteq A$, $s$ a sort in $S$, and $a$ an element in $\mathrm{Sg}_{\mathbf{A}}(X)_{s}$. Then we define the \emph{height of $a$ with respect to $\mathbf{A}$ and $X$}, denoted by $\mathrm{hgt}^{\mathbf{A}, X}_{s}(a)$, as:
% $$
% \mathrm{hgt}^{\mathbf{A}, X}_{s}(a) = \mathrm{min}\{n\in \mathbb{N}\mid a\in \mathrm{E}_{\mathbf{A}}^{n}(X)_{s}\}.
% $$
% %the smallest $n\in\mathbb{N}$ for which $a\in \mathrm{E}_{\mathbf{A}}^{n}(X)_{s}$.
% \end{definition}
%
% \begin{remark}
% Let $\mathbf{A}$ be a $\Sigma$-algebra and $X\subseteq A$. Then $\mathrm{hgt}^{\mathbf{A},X} = (\mathrm{hgt}^{\mathbf{A},X}_{s})_{s\in S}$ is an $S$-sorted mapping from $\mathrm{Sg}_{\mathbf{A}}(X)$ to $(\mathbb{N})_{s\in S}$.
% \end{remark}
%
% %\begin{remark}
% %Let $\Sigma$ be an $S$-sorted signature and $X$ an $S$-sorted set. On the $S$-sorted set $(\mathbb{N})_{s\in S}$ we consider the following $\Sigma$-algebra structure: If $\sigma\in \Sigma_{\lambda,s}$, for some $s\in S$, then $F_{\sigma}\colon \mathbb{N}^{0}\mor \mathbb{N}$ is identified to $0$; and if $\sigma\in \Sigma_{\mathbf{s},s}$, for some $(\mathbf{s},s)\in (S^{\star}-\{\lambda\})\times S$, then $F_{\sigma}\colon \mathbb{N}^{\bb{\mathbf{s}}}\mor \mathbb{N}$ sends $(n_{i})_{i\in\bb{\mathbf{s}}}$ to $1+\mathrm{max}\{n_{i}\mid i\in\bb{\mathbf{s}}\}$. Then there exists a unique homomorphism $\mathrm{hgt}^{\mathbf{T}_{\Sigma}(X),X}\colon \mathbf{T}_{\Sigma}(X)\mor ((\mathbb{N})_{s\in S},F)$ such that, for every $s\in S$ and every $x\in X_{s}$, $\mathrm{hgt}_{s}^{\mathbf{T}_{\Sigma}(X),X}(x) = 0$. For every $s\in S$ and every $P\in \mathrm{T}_{\Sigma}(X)_{s}$, we have 
% %$$
% %\mathrm{hgt}_{s}^{\mathbf{T}_{\Sigma}(X),X}(P) = \mathrm{min}\{n\in \mathbb{N}\mid P\in \E_{\mathbf{T}_{\Sigma}(X)}^{n}(X)_{s}\}.
% %$$
% %\end{remark}
%
% We next define, for a $\Sigma$-algebra $\mathbf{A}$ and a subset $X\subseteq A$, the family of the Borel sets of $X$ in $\mathbf{A}$.
%
% \begin{definition} Let $\mathbf{A}$ be a $\Sigma$-algebra and $X\subseteq A$. Then we define the family $(\mathrm{B}^{n}_{\mathbf{A}}(X))_{n\in \mathbb{N}}$ in $\mathrm{Sub}(A)$, recursively, as follows:
%       \begin{align*}
%       \mathrm{B}_{\mathbf{A}}^{0}(X) &= \mathrm{E}^{0}_{\mathbf{A}}(X)\text{,} \\
%        \mathrm{B}_{\mathbf{A}}^{n+1}(X) &=
%       \mathrm{E}_{\mathbf{A}}^{n+1}(X) -  \mathrm{E}_{\mathbf{A}}^{n}(X)\text{, $n\geq 0$.}
%       \end{align*}
% For a natural number $n\in\mathbb{N}$, the $S$-sorted set $\mathrm{B}^{n}_{\mathbf{A}}(X)$ will be called the \emph{$n$-th strict Borel set of $X$ in $\mathbf{A}$}, abbreviated to \emph{$n$-th Borel set of $X$ in $\mathbf{A}$}. 
% \end{definition}
%
% \begin{remark}
% Let $\mathbf{A}$ be a $\Sigma$-algebra and $X\subseteq A$. Then  
% $$
% \textstyle
% \mathrm{B}_{\mathbf{A}}^{0}(X) = X\cup (\bigcup_{\sigma\in \Sigma_{\lambda,s}}\sigma^{\mathbf{A}}[A_{\lambda}])_{s\in S},
% $$ 
% while the $1$-th Borel set of $X$ in $\mathbf{A}$ is
% $$
% \textstyle
% \mathrm{B}_{\mathbf{A}}^{1}(X) = \bigl(\;\bigcup_{\tau\in\Sigma_{\neq\lambda,s}}
%       \tau^{\mathbf{A}}[(X\cup (\bigcup_{\sigma\in \Sigma_{\lambda,s}}\sigma^{\mathbf{A}}[A_{\lambda}])_{s\in S})_{\mathrm{ar}(\sigma)}]\bigr)_{s\in S}.
% $$
% \end{remark}
%
% We now characterize, for every natural number $n$, the $n$-th Borel set of a set $X$ in a $\Sigma$-algebra $\mathbf{A}$ by means of the $S$-sorted mapping $\mathrm{hgt}^{\mathbf{A},X}$.
%
% \begin{proposition} Let $\mathbf{A}$ be a $\Sigma$-algebra, $X\subseteq A$, $n\in \mathbb{N}$ and $s\in S$. Then it happens that:
% $$
% \mathrm{B}^{n}_{\mathbf{A}}(X)_{s}
% =
% \{a\in\mathrm{Sg}_{\mathbf{A}}(X)_{s}\mid \mathrm{hgt}^{\mathbf{A},X}_{s}(a)=n\}.
% $$
% Therefore, for every $n\in \mathbb{N}$, $\mathrm{B}^{n}_{\mathbf{A}}(X) = ((\mathrm{hgt}^{\mathbf{A},X}_{s})^{-1}[\{n\}])_{s\in S}$.
% \end{proposition}
%
% We next define for a $\Sigma$-algebra $\mathbf{A}$ several binary relations on 
% $\coprod A$.
%
% \begin{definition}\label{DAlgOrd} 
% Let $\mathbf{A}$ be a $\Sigma$-algebra.
% We let $\prec_{\mathbf{A}}$ denote the binary relation on $\coprod A$ consisting of the ordered pairs $((b,t),(a,s))\in(\coprod A)^{2}$ for which there exists a word $\mathbf{s}\in S^{\star}-\{\lambda\}$, an operation symbol $\sigma\in \Sigma_{\mathbf{s},s}$ and a family of elements $(a_{j})_{j\in\bb{\mathbf{s}}}\in A_{\mathbf{s}}$ such that $a=\sigma^{\mathbf{A}}((a_{j})_{j\in\bb{\mathbf{s}}})$ and, for some $j\in\bb{\mathbf{s}}$, $s_{j}=t$ and $b=a_{j}$. If $((b,t),(a,s))\in \prec_{\mathbf{A}}$, also written as $(b,t)\prec_{\mathbf{A}}(a,s)$, then we will say that $(b,t)$ $\prec_{\mathbf{A}}$-\emph{precedes} $(a,s)$ or, when no confusion can arise, that $b$ \emph{precedes} $a$. We will call $\prec_{\mathbf{A}}$ the \emph{algebraic predecessor relation} on $\mathbf{A}$.
%
% We will denote by $\leq_{\mathbf{A}}$ the reflexive and transitive closure of $\prec_{\mathbf{A}}$, i.e., the preorder on $\coprod A$ generated by $\prec_{\mathbf{A}}$, and by $<_{\mathbf{A}}$ the transitive closure of $\prec_{\mathbf{A}}$. If $((b,t),(a,s))\in \leq_{\mathbf{A}}$, also written as $(b,t)\leq_{\mathbf{A}}(a,s)$, then we will say that $b$ is a \emph{subelement} of $a$. The preorder $\leq_{\mathbf{A}}$ will be called the \emph{algebraic predecessor preorder} on $\mathbf{A}$.
% \end{definition}
%
%
% \begin{remark}\label{RAlgOrd}
% The preorder $\leq_{\mathbf{A}}$ on $\coprod A$ is $\bigcup_{n\in\mathbb{N}}\prec_{\mathbf{A}}^{n}$, where $\prec_{\mathbf{A}}^{0}$ is the diagonal of $\coprod A$ and, for $n\in\mathbb{N}$, 
% $$\prec_{\mathbf{A}}^{n+1} = \prec_{\mathbf{A}}^{n}\circ \prec_{\mathbf{A}}.$$ 
% Thus, for every $((b,t),(a,s))\in (\coprod A)^{2}$, $(b,t)\leq_{\mathbf{A}}(a,s)$ if and only if $s = t$ and $b = a$ or there exists  a natural number $m\in\mathbb{N}-\{0\}$, a word $\mathbf{c}\in S^{\star}$ of length $\bb{\mathbf{c}}=m+1$, and a family of elements $(r_{k})_{k\in\bb{\mathbf{c}}}$ in $\mathrm{A}_{\mathbf{c}}$, such that $c_{0}=t$,  $r_{0}=b$, $c_{m}=s$, $r_{m}=a$ and, for every 
% $k\in m$, $(r_{k},c_{k})\prec_{\mathrm{A}}(r_{k+1},c_{k+1})$. Moreover, $<_{\mathbf{A}}$, the transitive closure of $\prec_{\mathbf{A}}$, is $\bigcup_{n\in\mathbb{N}-
% \{0\}}\prec_{\mathbf{A}}^{n}$.
% \end{remark}
%
% We conclude this subsection by showing how $\prec_{\mathbf{A}}$ relates to the height of an element of the underlying $S$-sorted set of $\mathbf{A}$.
%
% \begin{proposition} 
% Let $\mathbf{A}$ be a $\Sigma$-algebra, $s,t$ sorts in $S$ and $a,b$ elements in $A_{s}$ and $A_{t}$, respectively, such that $(b,t)$ $\prec_{\mathbf{A}}$-precedes $(a,s)$, thus there exists a word 
% $\mathbf{s}\in S^{\star}-\{\lambda\}$, an operation symbol $\sigma\in \Sigma_{\mathbf{s},s}$ and a family of elements $(a_{j})_{j\in\bb{\mathbf{s}}}\in A_{\mathbf{s}}$ such that $a=\sigma^{\mathbf{A}}((a_{j})_{j\in\bb{\mathbf{s}}})$ and, for some $j\in\bb{\mathbf{s}}$, $s_{j}=t$ and $b=a_{j}$. Then, for every $S$-sorted subset $X\subseteq A$ it holds that
% $$
% \mathrm{hgt}^{\mathbf{A},X}_{s}(a)\geq
% \mathrm{max}\{
% \mathrm{hgt}^{\mathbf{A},X}_{s_{j}}(a_{j})\mid j\in\bb{\mathbf{s}}
% \}+1.
% $$
% In particular, 
% $
% \mathrm{hgt}^{\mathbf{A},X}_{s}(a)\geq
% \mathrm{hgt}^{\mathbf{A},X}_{t}(b)+1
% $.
% \end{proposition}
%
% \subsection{Congruences}\label{SCong}
% Our next goal is to define the concepts of congruence on a $\Sigma$-algebra and of quotient of a $\Sigma$-algebra by a congruence on it. Moreover, we recall the notion of kernel of a homomorphism between $\Sigma$-algebras, the universal property of the quotient of a $\Sigma$-algebra by a congruence on it, and, for a $\Sigma$-algebra, the congruence generating operator for it. 
%
% \begin{definition}\label{DCong}
% Let $\mathbf{A}$ be a $\Sigma$-algebra and $\Phi$ an $S$-sorted equivalence on $A$. We will say that $\Phi$ is an $S$-\emph{sorted congruence on} $\mathbf{A}$ if, for every $(\mathbf{s},s)\in (S^{\star}-\{\lambda\})\times S$, every $\sigma\in \Sigma_{\mathbf{s},s}$, and every $(a_{j})_{j\in\bb{\mathbf{s}}},(b_{j})_{j\in\bb{\mathbf{s}}}\in A_{\mathbf{s}}$, if, for every $j\in \bb{\mathbf{s}}$, $(a_{j}, b_{j})\in\Phi_{s_{j}}$, then $(\sigma^{\mathbf{A}}((a_{j})_{j\in\bb{\mathbf{s}}}), \sigma^{\mathbf{A}}((b_{j})_{j\in\bb{\mathbf{s}}}))\in \Phi_{s}$.
% %we have that
% %$$
% %\frac
% %{\forall i\in \lvert w\rvert\mathrm{, }\,\, (a_{i}, b_{i})\in\Phi_{w_{i}} }
% %{(F_{\sigma}(a), F_{\sigma}(b))\in \Phi_{s}}\cdot
% %$$
% Sometimes, we will abbreviate the expression ``$S$-sorted congruence on $\mathbf{A}$'' to ``congruence on $\mathbf{A}$''. We will denote by $\mathrm{Cgr}(\mathbf{A})$ the set of all $S$-sorted congruences on $\mathbf{A}$ (which is an algebraic closure system on $A\times A$), by $\mathbf{Cgr}(\mathbf{A})$ the algebraic lattice  $(\mathrm{Cgr}(\mathbf{A}),\subseteq)$, by $\nabla_{\mathbf{A}}$ the greatest element of $\mathbf{Cgr}(\mathbf{A})$, and by $\Delta_{\mathbf{A}}$ the least element of $\mathbf{Cgr}(\mathbf{A})$.
%
% For a congruence $\Phi$ on $\mathbf{A}$, the \emph{quotient $\Sigma$-algebra of} $\mathbf{A}$ \emph{by} $\Phi$, denoted by $\mathbf{A}/\Phi$, is the $\Sigma$-algebra $(A/\Phi, F^{\mathbf{A}/\Phi})$, where, for every  $(\mathbf{s},s)\in S^{\star}\times S$ and every $\sigma\in \Sigma_{\mathbf{s},s}$, the operation $F^{\mathbf{A}/\Phi}_{\sigma}$ from $(A/\Phi)_{\mathbf{s}}$ to $A_{s}/\Phi_{s}$, also denoted, to simplify, by $\sigma^{\mathbf{A}/\Phi}$, sends $([a_{j}]_{\Phi_{s_{j}}})_{j\in\bb{\mathbf{s}}}$ in $(A/\Phi)_{\mathbf{s}}$ to $[\sigma^{\mathbf{A}}((a_{j})_{j\in \bb{\mathbf{s}}})]_{\Phi_{s}}$ in $A_{s}/\Phi_{s}$,  %
% %$$
% %F_{\sigma}
% %\nfunction
% %{(A/\Phi)_{w}} {A_{s}/\Phi_{s}}
% %{([a_{i}]_{\Phi_{w_{i}}})_{i\in\lvert w\rvert}}
% %{[F_{\sigma}((a_{i})_{i\in \lvert w\rvert})]_{\Phi_{s}}}
% %$$
% and the \emph{canonical projection} from $\mathbf{A}$ to $\mathbf{A}/\Phi$, denoted by $\mathrm{pr}^{\Phi}\colon \mathbf{A}\mor \mathbf{A}/\Phi$, is the full surjective homomorphism determined by the projection from $A$ to $A/\Phi$.
% The ordered pair $(\mathbf{A}/\Phi,\mathrm{pr}^{\Phi})$ has the following universal property: $\mathrm{Ker}(\mathrm{pr}^{\Phi})$ is $\Phi$ and, for every $\Sigma$-algebra $\mathbf{B}$ and every homomorphism $f$ from $\mathbf{A}$ to $\mathbf{B}$, if $\Phi\subseteq \mathrm{Ker}(f)$, then there exists a unique homomorphism $h$ from $\mathbf{A}/\Phi$ to $\mathbf{B}$ such that $h\circ\mathrm{pr}^{\Phi} = f$. In particular, if $\Psi$ is a congruence on $A$ such that $\Phi\subseteq \Psi$, then we will denote by $\mathrm{p}^{\Phi,\Psi}$ the unique homomorphism from $\mathbf{A}/\Phi$ to $\mathbf{A}/\Psi$ such that $\mathrm{p}^{\Phi,\Psi}\circ \mathrm{pr}^{\Phi} = \mathrm{pr}^{\Psi}$. %Moreover, $\Phi$ is a closed congruence on $\mathbf{A}$ if and only if $\mathrm{pr}^{\Phi}$ is a closed homomorphism.
% \end{definition}
%
% %\begin{remark}
% %Let $\mathbf{A}$ be a $\Sigma$-algebra and $\Phi$ an $S$-sorted equivalence on $A$. Then $\Phi$ is a congruence on $\mathbf{A}$ if and only if $\Phi$ is a subalgebra of $\mathbf{A}\times \mathbf{A}$.
% %\end{remark}
%
%
% \begin{remark}
% Let $\mathsf{ClfdAlg}(\Sigma)$ be the category whose objects are the \emph{classified $\Sigma$-algebras}, i.e, the ordered pairs $(\mathbf{A},\Phi)$ where $\mathbf{A}$ is a $\Sigma$-algebra and $\Phi$ a congruence on $\mathbf{A}$, and in which the set of morphisms from $(\mathbf{A},\Phi)$ to $(\mathbf{B},\Psi)$ is the set of all homomorphisms $f$ from $\mathbf{A}$ to $\mathbf{B}$ such that, for every $s\in S$ and every $(x,y)\in A^{2}_{s}$, if $(x,y)\in \Phi_{s}$, then $(f_{s}(x),f_{s}(y))\in \Psi_{s}$. Let $G$ be the functor from $\mathsf{Alg}(\Sigma)$ to $\mathsf{ClfdAlg}(\Sigma)$ whose object mapping sends $\mathbf{A}$ to $(\mathbf{A},\Delta_{\mathbf{A}})$ and whose morphism mapping sends $f\colon \mathbf{A}\mor \mathbf{B}$ to $f\colon (\mathbf{A},\Delta_{\mathbf{A}})\mor (\mathbf{B},\Delta_{\mathbf{B}})$. Then, for every classified $\Sigma$-algebra $(\mathbf{A},\Phi)$, there exists a universal mapping from $(\mathbf{A},\Phi)$ to $G$, which is precisely the ordered pair $(\mathbf{A}/\Phi,\mathrm{pr}^{\Phi})$ with $\mathrm{pr}^{\Phi}\colon (\mathbf{A},\Phi)\mor (\mathbf{A}/\Phi,\Delta_{\mathbf{A}/\Phi})$.
% \end{remark}
%
% %\begin{proposition}
% %Let $f$ be a homomorphism from $\mathbf{A}$ to $\mathbf{B}$. Then $\mathrm{Ker}(f)$ is a congruence on $\mathbf{A}$.
% %\end{proposition}
%
% %\begin{proposition}
% %Let $\mathbf{A}$ be a $\Sigma$-algebra and $\Phi\in \mathrm{Cgr}(\mathbf{A})$. Then the pair $(\mathrm{pr}^{\Phi},\mathbf{A}/\Phi)$ is such that:
% %\begin{enumerate}
% %\item $\mathrm{Ker}(\mathrm{pr}^{\Phi}) = \Phi$.
% %\item \emph{(Universal property)} For every homomorphism $f\colon\mathbf{A}\mor \mathbf{B}$, if $\Phi\subseteq \mathrm{Ker}(f)$, then there exists a unique homomorphism $\mathrm{p}^{\Phi,\mathrm{Ker}(f)}$ from $\mathbf{A}/\Phi$ to $\mathbf{B}$ such that $f = \mathrm{p}^{\Phi,\mathrm{Ker}(f)}\circ \mathrm{pr}^{\Phi}$.
% %\end{enumerate}
% %\end{proposition}
%
% %Given a $\Sigma$-algebra $\mathbf{A}$ and two congruences $\Phi$ and $\Psi$ on $\mathbf{A}$, if $\Phi\subseteq \Psi$, then, in the sequel, unless otherwise stated, $\mathrm{p}^{\Phi,\Psi}$ stands for the unique homomorphism from $\mathbf{A}/\Phi$ to $\mathbf{A}/\Psi$ such that $\mathrm{p}^{\Phi,\Psi}\circ \mathrm{pr}^{\Phi} = \mathrm{pr}^{\Psi}$.
%
% %\begin{remark}
% %Let $\mathbf{A}$ be a $\Sigma$-algebra. Then, for the congruence $\nabla_{\mathbf{A}}$, we have that $\mathbf{A}/\nabla_{\mathbf{A}}$ is isomorphic to a subalgebra of $\mathbf{1}$. In the single-sorted case, if $\boldsymbol{\varnothing}$ is an algebra, then $\boldsymbol{\varnothing}/\nabla_{\boldsymbol{\varnothing}}$ is $\boldsymbol{\varnothing}$, which is a subalgebra of $\mathbf{1}$, and $\mathrm{Sub}(\mathbf{1})$ is $\{\boldsymbol{\varnothing},\mathbf{1}\}$. In the many-sorted case, for a set of sorts $S$ such that $\mathrm{card}(S)\geq 2$ and an $S$-sorted signature $\Sigma$ such that, for every $s\in S$, $\Sigma_{\lambda,s} = \varnothing$, we have that $\boldsymbol{\varnothing}^{S}$ is a $\Sigma$-algebra and that $\boldsymbol{\varnothing}^{S}/\nabla_{\boldsymbol{\varnothing}^{S}}$ is $\boldsymbol{\varnothing}^{S}$. However, in contrast with what happens in the single-sorted case, in the many-sorted case, for a set of sorts $S$ and an $S$-sorted signature $\Sigma$ satisfying the above conditions, there may be $\Sigma$-algebras $\mathbf{A}$ such that $\varnothing\subset\mathrm{supp}_{S}(\mathbf{A})\subset S$. Hence, for such a type of $\Sigma$-algebras, the quotient $\Sigma$-algebra $\mathbf{A}/\nabla_{\mathbf{A}}$ will be isomorphic to a subalgebra in $\mathrm{Sub}(\mathbf{1}) - \{\boldsymbol{\varnothing}^{S},\mathbf{1}\}$.
% %\end{remark}
%
% We next recall the congruence generating operator for a $\Sigma$-algebra and provide a constructive description for it.
%
% \begin{definition}\label{DCongOpInt} Let $\mathbf{A}$ be a $\Sigma$-algebra. Then we denote by $\mathrm{Cg}_{\mathbf{A}}$ the many-sorted closure operator on $A\times A$ defined as follows:
%
% $$\textstyle
%   \mathrm{Cg}_{\mathbf{A}}\nfunction
%   {\mathrm{Sub}(A\times A)}
%   {\mathrm{Sub}(A\times A)}
%   {\Phi}
%   {\bigcap \{\,\Psi\in\mathrm{Con}(\mathbf{A})\mid \Phi\subseteq \Psi\,\}}
% $$
%
% We call $\mathrm{Cg}_{\mathbf{A}}$ the \emph{congruence generating many-sorted operator on} $A\times A$ \emph{determined by} $\mathbf{A}$. For every relation $\Phi\subseteq A\times A$, we call $\mathrm{Cg}_{\mathbf{A}}(\Phi)$ the \emph{congruence on} $\mathbf{A}$ \emph{generated by} $\Phi$. 
% \end{definition}
%
%
% For a $\Sigma$-algebra $\mathbf{A}$ we next provide another, more constructive, description of the algebraic many-sorted closure operator $\mathrm{Cg}_{\mathbf{A}}$.
%
% \begin{definition}\label{DCongOpC} 
% Let $\Sigma$ be an $S$-sorted signature and $\mathbf{A}$ a $\Sigma$-algebra.
% \begin{itemize}
% \item[(1)] We denote by $\mathrm{C}_{\mathbf{A}}$ the many-sorted operator on $A\times A$ that assigns to an $S$-sorted relation $\Phi\subseteq A\times A$, the $S$-sorted relation 
% $$
% \textstyle
% \mathrm{C}_{\mathbf{A}}(\Phi)
% =(\Phi\circ\Phi)
% \cup
% \left(
% \bigcup_{\sigma\in\Sigma_{\neq\lambda,s}}
% \sigma^{\mathbf{A}}\times \sigma^{\mathbf{A}}
% \left[
% \Phi_{\mathrm{ar}(\sigma)}
% \right]
% \right)_{s\in S},
% $$
% where, for $s\in S$, $\Sigma_{\neq\lambda,s}$ is the set of all many-sorted formal operations $\sigma$ such that $\mathrm{ar}(\sigma)\in S^{\star}-\{\lambda\}$, $\mathrm{car}(\sigma)=s$, and if $\mathrm{ar}(\sigma)=\mathbf{s}$, then $\Phi_{\mathrm{ar}(\sigma)}=\prod_{j\in\bb{\mathbf{s}}}\Phi_{s_{j}}$. Let us note that 
% $$
% \textstyle
% \Phi_{\mathrm{ar}(\sigma)}\subseteq \prod_{j\in\bb{\mathbf{s}}}A_{s_{j}}^{2}\cong (\prod_{j\in\bb{\mathbf{s}}}A_{s_{j}})^{2}.
% $$
% \item[(2)] If $\Phi\subseteq A\times A$, then we define the family $(\mathrm{C}^{n}_{\mathbf{A}}(\Phi))_{n\in\mathbb{N}}$ in $\mathrm{Sub}(A\times A)$, recursively, as follows:
% \begin{align*}
% \mathrm{C}^{0}_{\mathbf{A}}(\Phi)&=\Phi\cup\Phi^{-1}\cup\Delta_{\mathrm{A}},\\
% \mathrm{C}^{n+1}_{\mathbf{A}}(\Phi)&=\mathrm{C}_{\mathbf{A}}(\mathrm{C}_{\mathbf{A}}^{n}(\Phi)),\, n\geq 0.
% \end{align*}
% \item[(3)] We denote by $\mathrm{C}^{\omega}_{\mathbf{A}}$ the many-sorted operator on $A\times A$ that assigns to an $S$-sorted relation $\Phi\subseteq A\times A$, $\mathrm{C}^{\omega}_{\mathbf{A}}(\Phi)=\bigcup_{n\in\mathbb{N}}\mathrm{C}^{n}_{\mathbf{A}}(\Phi)$.
% \end{itemize}
% \end{definition}
%
% \begin{proposition}\label{PCongOpC} 
% Let $\mathbf{A}$ be a $\Sigma$-algebra and $\Phi\subseteq A\times A$ a relation on $A$. Then, for every $n\in\mathbb{N}$, we have that 
% \begin{itemize}
% \item[(1)] $\mathrm{C}^{n}_{\mathbf{A}}(\Phi)$ is a reflexive relation;
% \item[(2)]  $\mathrm{C}^{n}_{\mathbf{A}}(\Phi)\subseteq\mathrm{C}^{n+1}_{\mathbf{A}}(\Phi)$;
% \item[(3)] $\mathrm{C}^{n}_{\mathbf{A}}(\Phi)$ is a symmetric relation.
% \end{itemize}
% \end{proposition}
%
% \begin{proof}
% We prove the first item by induction on $n\in\mathbb{N}$.
%
% \textsf{Base step of the induction for the reflexivity property.}
%
% Since $\Delta_{\mathrm{A}}\subseteq \mathrm{C}^{0}_{\mathbf{A}}(\Phi)$, we have that $\mathrm{C}^{0}_{\mathbf{A}}(\Phi)$ is a reflexive relation.
%
% \textsf{Inductive step of the induction  for the reflexivity property.}
%
% Let us assume that the statement holds for $n\in\mathbb{N}$, i.e., that $\mathrm{C}^{n}_{\mathbf{A}}(\Phi)$ is a reflexive relation.
%
% We now prove that the statement holds for $n+1$. By the inductive hypothesis, $\mathrm{C}^{n}_{\mathbf{A}}(\Phi)$ is a reflexive relation, i.e., $\Delta_{\mathrm{A}}\subseteq \mathrm{C}^{n}_{\mathbf{A}}(\Phi)$. Hence 
% $$
% \Delta_{\mathrm{A}}\circ \mathrm{C}^{n}_{\mathbf{A}}(\Phi)\subseteq \mathrm{C}^{n}_{\mathbf{A}}(\Phi)\circ \mathrm{C}^{n}_{\mathbf{A}}(\Phi) \quad\text{and}\quad \Delta_{\mathrm{A}}\circ \Delta_{\mathrm{A}}\subseteq \Delta_{\mathrm{A}}\circ \mathrm{C}^{n}_{\mathbf{A}}(\Phi).
% $$ 
% But $\Delta_{\mathrm{A}}\subseteq \Delta_{\mathrm{A}}\circ \Delta_{\mathrm{A}}$. Consequently $\Delta_{\mathrm{A}}\subseteq \mathrm{C}^{n}_{\mathbf{A}}(\Phi)\circ \mathrm{C}^{n}_{\mathbf{A}}(\Phi)$. Moreover, we have that $\mathrm{C}^{n}_{\mathbf{A}}(\Phi)\circ \mathrm{C}^{n}_{\mathbf{A}}(\Phi)\subseteq \mathrm{C}^{n+1}_{\mathbf{A}}(\Phi)$. Thus $\Delta_{\mathrm{A}}\subseteq\mathrm{C}^{n+1}_{\mathbf{A}}(\Phi)$
%
% Therefore, $\mathrm{C}^{n+1}_{\mathbf{A}}(\Phi)$ is a reflexive relation. This finishes the proof of the first item.
%
% We now prove the second item. By (1), for every $n\in\mathbb{N}$, the relation $\mathrm{C}^{n}_{\mathbf{A}}(\Phi)$ is reflexive. Hence $\Delta_{\mathrm{A}}\subseteq \mathrm{C}^{n}_{\mathbf{A}}(\Phi)$. Since $\Delta_{\mathrm{A}}$ is a neutral element for the composition of relations, we have that 
% $$
% \mathrm{C}^{n}_{\mathbf{A}}(\Phi)
% =
% \mathrm{C}^{n}_{\mathbf{A}}(\Phi)\circ\Delta_{\mathrm{A}}
% \subseteq \mathrm{C}^{n}_{\mathbf{A}}(\Phi)\circ\mathrm{C}_{\mathbf{A}}^{n}(\Phi)
% \subseteq \mathrm{C}^{n+1}_{\mathbf{A}}(\Phi).
% $$
% Hence, $\mathrm{C}^{n}_{\mathbf{A}}(\Phi)\subseteq\mathrm{C}^{n+1}_{\mathbf{A}}(\Phi)$. This finishes the proof of the second item.
%
% We now prove the third statement by induction on $\mathrm{n}\in\mathbb{N}$.
%
% \textsf{Base step of the induction for the symmetry property.}
%
% Let us recall that 
% $\mathrm{C}^{0}_{\mathbf{A}}(\Phi)=\Phi\cup\Phi^{-1}\cup\Delta_{\mathrm{A}}$. Hence, the inverse of  $\mathrm{C}^{0}_{\mathbf{A}}(\Phi)$, which is given by 
% $
% \left(
% \mathrm{C}^{0}_{\mathbf{A}}(\Phi)
% \right)^{-1}
% =\Phi^{-1}\cup\Phi\cup\Delta_{\mathrm{A}}
% $, is included in $\mathrm{C}^{0}_{\mathbf{A}}(\Phi)$. Hence, $\mathrm{C}^{0}_{\mathbf{A}}(\Phi)$ is a symmetric relation.
%
% \textsf{Inductive step of the induction for the symmetry property.}
%
% Let us assume that the statement holds for $n\in\mathbb{N}$, i.e., that $\mathrm{C}^{n}_{\mathbf{A}}(\Phi)$ is a symmetric relation.
%
% We now prove the statement for $n+1$. Let $s$ be a sort in $s$ and let  $(x,y)$ be a pair in $\mathrm{C}^{n+1}_{\mathbf{A}}(\Phi)$. Let us recall that $\mathrm{C}^{n+1}_{\mathbf{A}}(\Phi)$ is given by
% $$
% \textstyle
% \mathrm{C}^{n+1}_{\mathbf{A}}(\Phi)=
% \mathrm{C}_{\mathbf{A}}(\mathrm{C}^{n}_{\mathbf{A}}(\Phi))
% =
% (\mathrm{C}^{n}_{\mathbf{A}}(\Phi)\circ\mathrm{C}^{n}_{\mathbf{A}}(\Phi))
% \cup
% \left(
% \bigcup_{\sigma\in\Sigma_{\neq\lambda,s}}
% \sigma^{\mathbf{A}}\times \sigma^{\mathbf{A}}
% \left[
% \mathrm{C}^{n}_{\mathbf{A}}(\Phi)
% \right]
% \right)_{s\in S},
% $$
% hence the pair $(x,y)$ is either (1) a pair  in $\mathrm{C}^{n}_{\mathbf{A}}(\Phi)_{s}\circ\mathrm{C}^{n}_{\mathbf{A}}(\Phi)_{s}$ or (2) a pair in $\bigcup_{\sigma\in\Sigma_{\neq\lambda,s}}
% \sigma^{\mathbf{A}}\times \sigma^{\mathbf{A}}
% \left[
% \mathrm{C}^{n}_{\mathbf{A}}(\Phi)
% \right]
% $.
%
% If (1), then exists an element $z\in A_{s}$ such that the pairs $(x,z)$ and $(z,y)$ are in $\mathrm{C}^{n}_{\mathbf{A}}(\Phi)_{s}$. By induction, $\mathrm{C}^{n}_{\mathbf{A}}(\Phi)_{s}$ is a symmetric relation. Hence, the pairs $(z,x)$ and $(y,z)$ belong to $\mathrm{C}^{n}_{\mathbf{A}}(\Phi)_{s}$. Thus, the pair $(y,x)$ also belongs to $\mathrm{C}^{n+1}_{\mathbf{A}}(\Phi)_{s}$.
%
% If (2), then we can find a non-empty word $\mathbf{s}\in S^{\star}-\{\lambda\}$, an operation symbol $\sigma\in \Sigma_{\mathbf{s},s}$ and a family of pairs $((x_{j},y_{j}))_{j\in\bb{\mathbf{s}}}$ in $\mathrm{C}^{n}_{\mathbf{A}}(\Phi)_{\mathbf{s}}$ such that $x=\sigma^{\mathbf{A}}((x_{j})_{j\in\bb{\mathbf{s}}})$ and $y=\sigma^{\mathbf{A}}((y_{j})_{j\in\bb{\mathbf{s}}})$. By induction we have that, for every $j\in\bb{\mathbf{s}}$, $\mathrm{C}^{n}_{\mathbf{A}}(\Phi)_{s_{j}}$ is a symmetric relation. Hence, for every $j\in\bb{\mathbf{s}}$, the pair $(y_{j},x_{j})$ belongs to $\mathrm{C}^{n}_{\mathbf{A}}(\Phi)_{s_{j}}$. Thus, the pair $(\sigma^{\mathbf{A}}((y_{j})_{j\in\bb{\mathbf{s}}}), \sigma^{\mathbf{A}}((x_{j})_{j\in\bb{\mathbf{s}}}))$, which is equal to $(y,x)$, also belongs to $\mathrm{C}^{n+1}_{\mathbf{A}}(\Phi)_{s}$.
%
% In any case, the pair $(y,x)$ belongs to $\mathrm{C}^{n+1}_{\mathbf{A}}(\Phi)_{s}$.
% Hence, $\mathrm{C}^{n+1}_{\mathbf{A}}(\Phi)$ is a symmetric relation. This finishes the proof of the third item.
% \end{proof}
%
% \begin{proposition}\label{PCongIntC} 
% Let $\mathbf{A}$ be a $\Sigma$-algebra and $\Phi\subseteq A\times A$. Then
% $\mathrm{Cg}_{\mathbf{A}}(\Phi)=\mathrm{C}_{\mathbf{A}}^{\omega}(\Phi)$.
% \end{proposition}
%
% \begin{proof}
% We first prove the inclusion from left to right. To do so, we prove that $\mathrm{C}_{\mathbf{A}}^{\omega}(\Phi)$ is a $\Sigma$-congruence on $\mathbf{A}$.
%
% By Proposition~\ref{PCongOpC}, $\Delta_{A}\subseteq\mathrm{C}^{0}_{\mathbf{A}}(\Phi)$, it follows that $\Delta_{A}\subseteq\bigcup_{n\in\mathbb{N}}\mathrm{C}^{n}_{\mathbf{A}}(\Phi)=\mathrm{C}^{\omega}_{\mathbf{A}}(\Phi)$. Hence, $\mathrm{C}^{\omega}_{\mathbf{A}}(\Phi)$ is a reflexive relation. 
%
% Regarding symmetry, let $s$ be a sort in $S$ and let $(x,y)$ be a pair in $\mathrm{C}^{\omega}_{\mathbf{A}}(\Phi)_{s}$, then there exists an $n\in\mathbb{N}$ for which $(x,y)$ is a pair in $\mathrm{C}^{n}_{\mathbf{A}}(\Phi)_{s}$. By Proposition~\ref{PCongOpC}, $\mathrm{C}^{n}_{\mathbf{A}}(\Phi)_{s}$ is a symmetric relation. Hence $(y,x)\in \mathrm{C}^{n}_{\mathbf{A}}(\Phi)_{s}\subseteq \mathrm{C}^{\omega}_{\mathbf{A}}(\Phi)_{s}$. Therefore $\mathrm{C}^{\omega}_{\mathbf{A}}(\Phi)$ is a symmetric relation. 
%
% Regarding transitivity, let $s$ be a sort in $S$ and let $(x,y)$ and $(y,z)$ be a pair of terms in $\mathrm{C}^{\omega}_{\mathbf{A}}(\Phi)_{s}$, then there exists an $m$ and an $n\in\mathbb{N}$ for which $(x,y)$ is a pair in $\mathrm{C}^{m}_{\mathbf{A}}(\Phi)_{s}$ and $(y,z)$ is a pair in $\mathrm{C}^{n}_{\mathbf{A}}(\Phi)_{s}$. Let $r=\mathrm{max}(m,n)$ then, in virtue of Proposition~\ref{PCongOpC}, the pairs $(x,y)$ and $(y,z)$ are in $\mathrm{C}^{r}_{\mathbf{A}}(\Phi)_{s}$. It follows that $(x,z)$ is a pair in $\mathrm{C}^{r}_{\mathbf{A}}(\Phi)_{s}\circ\mathrm{C}^{r}_{\mathbf{A}}(\Phi)_{s}\subseteq \mathrm{C}^{r+1}_{\mathbf{A}}(\Phi)_{s}\subseteq \mathrm{C}^{\omega}_{\mathbf{A}}(\Phi)_{s}$. Therefore $\mathrm{C}^{\omega}_{\mathbf{A}}(\Phi)$ is a transitive relation. 
%
% Finally, we check that $\mathrm{C}^{m}_{\mathbf{A}}(\Phi)$ is compatible with the operations in $\Sigma$. Let $\mathbf{s}$ be a non-empty word in $S^{\star}-\{\lambda\}$ and $s$ a sort in $S$. Let $\sigma$ be an operation symbol in $\Sigma_{\mathbf{s},s}$. Let $(x_{j})_{j\in\bb{\mathbf{s}}}$ and $(y_{j})_{j\in\mathbf{s}}$ be a pair of families in $\mathbf{A}_{\mathbf{s}}$ such that, for every $j\in\bb{\mathbf{s}}$, the pair $(x_{j},y_{j})$ is in $\mathrm{C}^{\omega}_{\mathbf{A}}(\Phi)_{s_{j}}$. Then there exists, for every $j\in\bb{\mathbf{s}}$, a natural number $n_{j}\in\mathbb{N}$ such that $(x_{j},y_{j})$ is a pair in $\mathrm{C}^{n_{j}}_{\mathbf{A}}(\Phi)_{s_{j}}$. Let $r=\mathrm{max}\{n_{j}\mid j\in\bb{\mathbf{s}}\}$ then, in virtue of Proposition~\ref{PCongOpC}, for every $j\in\bb{\mathbf{s}}$, the pair $(x_{j},y_{j})$ is in $\mathrm{C}^{r}_{\mathbf{A}}(\Phi)_{s_{j}}$. Therefore  
% $
% (
% \sigma^{\mathbf{A}}
% (
% (
% x_{j}
% )_{j\in\bb{\mathbf{s}}}
% ),
% \sigma^{\mathbf{A}}
% (
% (
% y_{j}
% )_{j\in\bb{\mathbf{s}}}
% )
% )
% $ is a pair in $\sigma^{\mathbf{A}}\times \sigma^{\mathbf{A}}[\mathrm{C}^{r}_{\mathbf{A}}(\Phi)_{\mathbf{s}}]$,
% which is a subset of $\mathrm{C}^{r+1}_{\mathbf{A}}(\Phi)_{s}\subseteq \mathrm{C}^{\omega}_{\mathbf{A}}(\Phi)_{s}$. 
%
% All in all, we can affirm that $\mathrm{C}^{\omega}_{\mathbf{A}}(\Phi)$ is a congruence on $\mathbf{A}$.
%
% Moreover $\mathrm{C}^{\omega}_{\mathbf{A}}(\Phi)$ contains $\Phi$. This is so because $\Phi\subseteq \mathrm{C}^{0}_{\mathbf{A}}(\Phi)$. From these facts, we can affirm that $\mathrm{Cg}_{\mathbf{A}}(\Phi)\subseteq\mathrm{C}^{\omega}_{\mathbf{A}}(\Phi)$.
%
% To prove the other inclusion, we proceed by induction on $n\in \mathbb{N}$.
%
% \textsf{Base step of the induction.} 
%
% Since $\mathrm{Cg}_{\mathbf{A}}(\Phi)$ is a congruence on $\mathbf{A}$ containing $\Phi$, we can affirm that $\Delta_{A}$, $\Phi$ and $\Phi^{-1}$ are subsets of $\mathrm{Cg}_{\mathbf{A}}(\Phi)$. Hence $\mathrm{C}^{0}_{\mathbf{A}}(\Phi)\subseteq \mathrm{Cg}_{\mathbf{A}}(\Phi)$.
%
% \textsf{Inductive step of the induction.}
%
% Let us assume that the statement holds for  $n\in\mathbb{N}$, i.e., we have that
% $\mathrm{C}^{n}_{\mathbf{A}}(\Phi)\subseteq \mathrm{Cg}_{\mathbf{A}}(\Phi)$. 
%
% We now prove the statement for $n+1$. 
% By the inductive hypothesis, the relation $\mathrm{C}^{n}_{\Phi}$ is included in $\mathrm{Cg}_{\mathbf{A}}(\Phi)$. Since $\mathrm{Cg}_{\mathbf{A}}(\Phi)$ is a transitive relation, we can affirm that $(\mathrm{C}^{n}_{\Phi}\circ \mathrm{C}^{n}_{\Phi})$ is included in $\mathrm{Cg}_{\mathbf{A}}(\Phi)$. 
%
% On the other hand, for every non-empty word $\mathbf{s}\in S-\{\lambda\}$, every sort $s\in S$, every operation symbol $\sigma$ in $\Sigma_{\mathbf{s},s}$, if $(x_{j},y_{j})_{j\in\bb{\mathbf{s}}}$ is a family of pairs in $\mathrm{C}^{n}_{\mathbf{A}}(\Phi)_{\mathbf{s}}$, since $\mathrm{Cg}_{\mathbf{A}}(\Phi)$ is a $\Sigma$-congruence and, by the 
% inductive hypothesis, for every $j\in\bb{\mathbf{s}}$, $\mathrm{C}^{n}_{\mathbf{A}}(\Phi)_{s_{j}}\subseteq \mathrm{Cg}_{\mathbf{A}}(\Phi)_{s_{j}}$, we have that the pair $(\sigma^{\mathbf{A}}((x_{j})_{j\in\bb{\mathbf{s}}}, \sigma^{\mathbf{A}}((y_{j})_{j\in\bb{\mathbf{s}}})$ belongs to $\mathrm{Cg}_{\mathbf{A}}(\Phi)_{s}$.
%
% All in all, we can affirm that $\mathrm{C}^{\omega}_{\mathbf{A}}(\Phi)\subseteq \mathrm{Cg}_{\mathbf{A}}(\Phi)$.
% Hence, $\mathrm{Cg}_{\mathbf{A}}(\Phi)=\mathrm{C}^{\omega}_{\mathbf{A}}(\Phi)$.
%
% This finishes the proof.
% \end{proof}
%
% \begin{proposition}
% Let $f$ be a surjective homomorphism from the $\Sigma$-algebra $\mathbf{A}$ to the $\Sigma$-algebra 
% $\mathbf{B}$ and $\Phi\subseteq A\times A$. Then 
% $$
% \mathrm{Cg}_{\mathbf{B}}(f^{2}[\Phi]) = f^{2}[\mathrm{Ker}(f)\vee\mathrm{Cg}_{\mathbf{A}}(\Phi)].
% $$
% \end{proposition}
%
% \subsection{Elementary translations and translations}\label{STrans}
%
% We next define for a $\Sigma$-algebra the concepts of elementary translation and of translation with respect to it, and provide, by using the just mentioned concepts, two characterizations of the congruences on a $\Sigma$-algebra. To this we add that the concept of translation will play an essential role in the work at hand. 
% %To this we add that the concept of translation will allow us to define the concept of congruence cogenerated by an $S$-sorted subset of the underlying $S$-sorted set of a $\Sigma$-algebra, which will be all-important in our congruence based proofs of the recognizability theorems for free many-sorted algebras, in Section 3.
%
% \begin{definition}\label{DETrans}
%
% Let $\mathbf{A}$ be a $\Sigma$-algebra and $t\in S$. Then we will denote by $\mathrm{Etl}_{t}(\mathbf{A})$ the subset $(\mathrm{Etl}_{t}(\mathbf{A})_{s})_{s\in S}$ of $(\mathrm{Hom}(A_{t},A_{s}))_{s\in S}$ defined, for every $s\in S$, as follows: for every mapping $T\in \mathrm{Hom}(A_{t},A_{s})$, $T\in \mathrm{Etl}_{t}(\mathbf{A})_{s}$ if and only if the following condition holds
% \begin{enumerate}
% \item There is a word $\mathbf{s}\in  S^{\star}-\{\lambda\}$, an index $k\in \lvert \mathbf{s} \rvert$, an operation symbol $\sigma\in \Sigma_{\mathbf{s},s}$, a family $(a_{j})_{j\in i}\in\prod_{j\in k}A_{s_{j}}$, and a family $(a_{l})_{l\in \lvert \mathbf{s} \rvert-(k+1)} \in\prod_{l\in \lvert \mathbf{s} \rvert-(k+1)}A_{s_{l}}$ (recall that $k+1 = \{0, 1,\ldots,k\}$ and that $\lvert \mathbf{s} \rvert-(k+1) = \{k+1,\ldots,\lvert \mathbf{s} \rvert-1\}$) such that $s_{k} = t$ and, for every $x\in A_{t}$, 
% $$T(x) =
% \sigma^{\mathbf{A}}(a_{0},\ldots,a_{k-1},x,a_{k+1},\ldots,a_{\lvert \mathbf{s} \rvert-1}).$$ 
% \end{enumerate}
% We will sometimes use the following presentation of the elementary translations, which consists in adding an underlined space to denote where the variable will be placed
% $$T =
% \sigma^{\mathbf{A}}(a_{0},\ldots,a_{k-1},\underline{\quad},a_{k+1},\ldots,a_{\lvert \mathbf{s} \rvert-1}).$$ 
% In this case we will say that $T$ is an \emph{elementary translation of type} $\sigma$.
% We will call the elements of $\mathrm{Etl}_{t}(\mathbf{A})_{s}$ the $t$-\emph{elementary translations of sort} $s$ \emph{for} $\mathbf{A}$.
% \end{definition}
%
% \begin{definition}\label{DTrans}
%
% Let $\mathbf{A}$ be a $\Sigma$-algebra and $t\in S$. Then we will denote by $\mathrm{Tl}_{t}(\mathbf{A})$ the subset $(\mathrm{Tl}_{t}(\mathbf{A})_{s})_{s\in S}$ of $(\mathrm{Hom}(A_{t},A_{s}))_{s\in S}$ defined, for every $s\in S$, as follows: for every mapping $T\in \mathrm{Hom}(A_{t},A_{s})$, $T\in \mathrm{Tl}_{t}(\mathbf{A})_{s}$ if and only if there is an $n\in \mathbb{N}-1$, a word $\mathbf{w}\in S^{n+1}$, and a family $(T_{j})_{j\in n}$ such that $w_{0} = t$, $w_{n} = s$, for every $j\in n$, $T_{j}\in \mathrm{Etl}_{w_{j}}(\mathbf{A})_{w_{j+1}}$, and $T = T_{n-1}\circ\cdots\circ T_{0}$. For translations, as for words on an alphabet, we have the notion of subtranslation of a translation, which is the counterpart of that of subword of a word. In particular, for a translation as above we will let $T'$ stand for the composition $T_{n-2}\circ \cdots \circ T_{0}$ and we will call it the \emph{maximal prefix of $T$}, and we will represent $T$ as $T_{n-1}\circ T'$ or under the form: 
% $$
% T=\sigma^{\mathbf{A}}\left(
% a_{0},\cdots, a_{k-1}, T',
% a_{k+1}, \cdots, a_{\mathbf{s}-1}
% \right),
% $$
% where $T_{n-1}=\sigma^{\mathbf{A}}(a_{0},\ldots,a_{k-1},\underline{\quad},a_{k+1},\ldots,a_{\lvert \mathbf{s} \rvert-1})$. 
% %This second representation means that, for every $a\in A_{t}$, $T(a)=T_{n-1}(T'(a))$. 
% The underlined space notation can be extended to translations as well. We will say that $T$ is a \emph{translation of type $\sigma$} if the elementary translation $T_{n-1}$ is of type $\sigma$. We will call $n$ the \emph{height} of $T$ and we will denote this fact by $\bb{T}=n$. In this regard, elementary translations have height $1$ and, if $T$ is a translation of height $n$, i.e., $\bb{T}=n$, then its maximal prefix has height $n-1$, i.e., $\bb{T'}=n-1$. We will call the elements of $\mathrm{Tl}_{t}(\mathbf{A})_{s}$ the $t$-\emph{translations of sort} $s$ \emph{for} $\mathbf{A}$. In addition, for every $t\in S$, the mapping  $\id^{A_{t}}$ will be viewed as an element of $\mathrm{Tl}_{t}(\mathbf{A})_{t}$. The identity translation has no associated type and we will consider that it has height $0$, i.e., $\bb{\id^{A_{t}}}=0$. Moreover, since we consider that, for every $t\in S$, the identity mapping $\id^{A_{t}}$ is a traslation, we agree that an elementary translation $T\colon A_{t}\mor A_{s}$, of sort $s$ for $\mathbf{A}$, has as maximal prefix the identity at the domain.
% \end{definition}
%
%
% \begin{remark}
% The $S\times S$-sorted set $(\mathrm{Tl}_{t}(\mathbf{A})_{s})_{(t,s)\in S\times S}$ determines a category $\mathsf{Tl}(\mathbf{A})$ whose object set is $S$ and in which, for every $(t,s)\in S\times S$, $\mathrm{Hom}_{\mathsf{Tl}(\mathbf{A})}(t,s)$, the hom-set from $t$ to $s$, is $\mathrm{Tl}_{t}(\mathbf{A})_{s}$. Therefore, for every $t\in S$, $\mathrm{End}_{\mathsf{Tl}(\mathbf{A})}(t)$ is equipped with a structure of monoid.
% \end{remark}
%
% %\begin{remark}
% %Let $X$ be an $S$-sorted set, $\mathbf{A}$ a $\Sigma$-algebra, and $M\in \mathrm{Sub}(A)$. Then $\mathbf{O}_{X}(\mathbf{A})$, the $\Sigma$-algebra of the $X$-\emph{ary} \emph{operations on} $\mathbf{A}$, is $\mathbf{A}^{\mathrm{Hom}(X,A)}$, i.e., the direct $\mathrm{Hom}(X,A)$-power of $\mathbf{A}$, and $\mathsf{Alg}_{X}(\mathbf{A},M)$, the $\Sigma$-algebra of the $X$-\emph{ary algebraic operations with constants in} $M$ \emph{on} $\mathbf{A}$, is the subalgebra of $\mathbf{O}_{X}(\mathbf{A})$ generated by the subfamily $\mathcal{P}^{A}_{X}\cup \mathcal{K}^{A,M}_{X}$ of $\mathrm{O}_{X}(A)$, where $\mathcal{P}^{A}_{X}$ is $(\{\,\mathrm{pr}^{A}_{X,s,x} \mid x\in X_{s}\,\})_{s\in S}$, with $\mathrm{pr}^{A}_{X,s,x}$ the mapping from $\mathrm{Hom}(X,A)$ to $A_{s}$ which sends $a\in \mathrm{Hom}(X,A)$ to $a_{s}(x)$, for every $s\in S$ and $x\in X_{s}$,
% %%$
% %%\mathcal{P}^{A}_{X} = (\mathcal{P}^{A}_{X,s})_{s\in S} =
% %%(\{\,\mathrm{pr}^{A}_{X,s,x} \mid x\in X_{s}\,\})_{s\in S},
% %%$
% %%and, for every $s\in S$ and $x\in X_{s}$, $\mathrm{pr}^{A}_{X,s,x}$ is the mapping from $\mathrm{Hom}(X,A)$ to $A_{s}$ which sends $a\in \mathrm{Hom}(X,A)$ to $a_{s}(x)$,
% %and $\mathcal{K}^{A,M}_{X}$ is $(\{\,\kappa^{a}_{X,s} \mid a\in M_{s}\,\})_{s\in S}$, with $\kappa^{a}_{X,s}$ the mapping from $\mathrm{Hom}(X,A)$ to $A_{s}$ that is constantly $a$, for every $s\in S$ and $a\in M_{s}$.
% %%$
% %%\mathcal{K}^{A,M}_{X} = (\mathcal{K}^{A,M}_{X,s})_{s\in S}
% %%= (\{\,\kappa^{a}_{X,s} \mid a\in M_{s}\,\})_{s\in S},
% %%$
% %%and, for every $s\in S$ and $a\in M_{s}$, $\kappa^{a}_{X,s}$ is the mapping from $\mathrm{Hom}(X,A)$ to $A_{s}$ constantly $a$.
% %If, for $M = A$, we write $\mathsf{Alg}_{X}(\mathbf{A})$ for $\mathsf{Alg}_{X}(\mathbf{A},A)$, and we call it the $\Sigma$-algebra of the $X$-\emph{ary algebraic operations on} $\mathbf{A}$, then, for $X = \delta^{t}$, we have that $\mathrm{Tl}_{t}(\mathbf{A})$ is, up to isomorphism, a part of $\mathrm{Alg}_{\delta^{t}}(\mathbf{A})$.
% %Note that, for $M = \varnothing^{S}$, $\mathsf{Alg}_{X}(\mathbf{A},\varnothing^{S})$ is $\mathbf{Ter}_{X}(\mathbf{A})$, the $\Sigma$-algebra of the $X$-\emph{ary term operations on} $\mathbf{A}$, and that it is a subalgebra of $\mathsf{Alg}_{X}(\mathbf{A})$ since $\mathbf{Ter}_{X}(\mathbf{A})$ is the subalgebra of $\mathbf{O}_{X}(\mathbf{A})$ generated by the subfamily $\mathcal{P}^{A}_{X}$ of $\mathrm{O}_{X}(A)$.
% %\end{remark}
%
% %Given a $\Sigma$-algebra $\mathbf{A}$ and a translation $T\in \mathrm{Tl}_{t}(\mathbf{A})_{s}$ we next define, as\-so\-ci\-at\-ed to the mappings $T[\cdot]\colon \mathrm{Sub}(A_{t})\mor \mathrm{Sub}(A_{s})$ and $T^{-1}[\cdot]\colon \mathrm{Sub}(A_{s})\mor \mathrm{Sub}(A_{t})$, operators: $T[\cdot]$ and $T^{-1}[\cdot]$ from $\mathrm{Sub}(A)$ to $\mathrm{Sub}(A)$, $T[\cdot]$ from $\mathrm{Sub}(A_{t})$ to $\mathrm{Sub}(A)$, and $T^{-1}[\cdot]$ from $\mathrm{Sub}(A_{s})$ to $\mathrm{Sub}(A)$. This will be used below in the characterization of the recognizable languages.
%
% %\begin{definition}\label{DAntiTrans}
% %Let $\mathbf{A}$ be a $\Sigma$-algebra, $L\subseteq A$, $s$, $t\in S$, $M\subseteq A_{t}$, $N\subseteq A_{s}$, and $T\in \mathrm{Tl}_{t}(\mathbf{A})_{s}$. Then
% %\begin{enumerate}
% %\item $T[L]$ is the subset of $A$ defined as follows: $T[L]_{s} = T[L_{t}]$ and $T[L]_{u} = \varnothing$, if $u\neq s$. Therefore, $T[L] = \delta^{s,T[L_{t}]}$;
% %\item $T^{-1}[L]$ is the subset of $A$ defined as follows: $T^{-1}[L]_{t} = T^{-1}[L_{s}]$ and $T^{-1}[L]_{u} = \varnothing$, if $u\neq t$. Therefore, $T^{-1}[L] = \delta^{t,T^{-1}[L_{s}]}$;
% %\item $T[M]$ is $\delta^{s,T[M]} (= T[\delta^{t,M}])$; and
% %\item $T^{-1}[N]$ is $\delta^{t,T^{-1}[N]} (= T^{-1}[\delta^{s,N}])$.
% %\end{enumerate}
% %\end{definition}
%
% %\begin{remark}
% %Let $\mathbf{A}$ be a $\Sigma$-algebra, $K$, $L\subseteq A$, $s$, $t\in S$, $M\subseteq A_{t}$, $N\subseteq A_{s}$, and $T\in \mathrm{Tl}_{t}(\mathbf{A})_{s}$. Then $T[M]\subseteq \delta^{s,N}$ if and only if $\delta^{t,M}\subseteq T^{-1}[N]$. Moreover, $T[K]\subseteq \delta^{s,L_{s}}$ if and only if $\delta^{t,K_{t}}\subseteq T^{-1}[L]$. Had the operators $T[\cdot]$ and $T^{-1}[\cdot]$ been defined (for $u\in S-\{s,t\}$) differently, then the suitably modified counterparts of the just stated connections would not be met.
% %\end{remark}
%
% The following are the above-mentioned characterizations of the congruences on a $\Sigma$-algebra. We note that in~\cite{m76}, on p.~199, it was announced without proof a proposition similar to that set out below.
%
%
% \begin{proposition}\label{PTransCong}
% Let $\mathbf{A}$ be a $\Sigma$-algebra and $\Phi$ an $S$-sorted equivalence on $A$. Then the following conditions are equivalent:
% \begin{enumerate}
% \item $\Phi$ is a congruence on $\mathbf{A}$.
% \item $\Phi$ is closed under the elementary translations on
% $\mathbf{A}$, i.e., for every every $t$, $s\in S$, every $x$, $y\in A_{t}$, and every $T\in \mathrm{Etl}_{t}(\mathbf{A})_{s}$, if $(x,y)\in \Phi_{t}$, then $(T(x),T(y))\in \Phi_{s}$.
% \item $\Phi$ is closed under the translations on $\mathbf{A}$, i.e., for every every $t$, $s\in S$, every $x$, $y\in A_{t}$, and every $T\in \mathrm{Tl}_{t}(\mathbf{A})_{s}$, if $(x,y)\in \Phi_{t}$, then $(T(x),T(y))\in \Phi_{s}$.
% \end{enumerate}
% \end{proposition}
%
% \begin{proof}
% Let us first prove that (1) and (2) are equivalent.
%
% Let us suppose that $\Phi$ is a congruence on $\mathbf{A}$. We want to show that $\Phi$ is closed under the elementary translations on $\mathbf{A}$. Let $t$ and $s$ be elements of $S$ and $T$ a $t$-elementary translation of sort $s$ for $\mathbf{A}$. Then $T\colon A_{t}\mor A_{s}$ and there is a word $\mathbf{s}\in S^{\star}-\{\lambda\}$, an index $k\in \lvert \mathbf{s} \rvert$, an operation symbol $\sigma\in
% \Sigma_{\mathbf{s},s}$, a family $(a_{j})_{j\in k}\in\prod_{j\in k}A_{s_{j}}$, and a family $(a_{l})_{l\in \lvert \mathbf{s} \rvert-(k+1)}
% \in\prod_{kl\in \lvert \mathbf{s} \rvert-(k+1)}A_{s_{l}}$ such that $s_{k} = t$ and, for every $z\in A_{t}$, $T(z) =
% \sigma^{\mathbf{A}}(a_{0},\ldots,a_{k-1},z,a_{k+1},\ldots,a_{\lvert \mathbf{s} \rvert-1})$. Let $x$ and $y$ be elements of $A_{t}$ such that $(x,y)\in \Phi_{t}$. Since, for every $j\in k$, $(a_{j},a_{j})\in \Phi_{s_{j}}$, for every $l\in \lvert \mathbf{s} \rvert-(k+1)$, $(a_{l},a_{l})\in \Phi_{s_{l}}$, and, in addition, $(x,y)\in \Phi_{t} = \Phi_{s_{k}}$, then $(T(x),T(y))\in \Phi_{s}$.
%
% Reciprocally, let us suppose that, for every $t$, $s\in S$, every $x$, $y\in A_{t}$, and every $T\in \mathrm{Etl}_{t}(\mathbf{A})_{s}$, if $(x,y)\in \Phi_{t}$, then $(T(x),T(y))\in \Phi_{s}$. We want to show that  $\Phi$ is a congruence on $\mathbf{A}$. Let $(\mathbf{s},s)\in (S^{\star}-\{\lambda\})\times S$, $\sigma\in\Sigma_{\mathbf{s},s}$,
% and $(a_{i})_{i\in\lvert \mathbf{s} \rvert}$ and $ (b_{i})_{i\in\lvert \mathbf{s} \rvert}$ be families in $ A_{\mathbf{s}}$ such that, for every $i\in \lvert \mathbf{s} \rvert$ we have that $(a_{i}, b_{i})\in \Phi_{s_{i}}$. We now define, for every $i\in\lvert \mathbf{s} \rvert$, $T_{i}$, the $s_{i}$-elementary translation of sort $s$ for $\mathbf{A}$, as the mapping from $A_{s_{i}}$ to $A_{s}$ which sends $x\in A_{s_{i}}$ to $\sigma^{\mathbf{A}}(b_{0},\ldots,b_{i-1},x,a_{i+1},\ldots,a_{\lvert \mathbf{s} \rvert-1}) \in A_{s}$. Then $\sigma^{\mathbf{A}}(a_{0},\ldots,a_{\lvert \mathbf{s} \rvert-1}) = T_{0}(a_{0})$ and $(T_{0}(a_{0}),T_{0}(b_{0}))\in \Phi_{s_{0}}$. But $T_{0}(b_{0}) = T_{1}(a_{1})$ and $(T_{1}(a_{1}),T_{1}(b_{1}))\in \Phi_{s_{1}}$. By proceeding in the same way we, finally, come to $T_{\lvert \mathbf{s} \rvert-2}(b_{\lvert \mathbf{s} \rvert-2}) = T_{\lvert \mathbf{s} \rvert-1}(a_{\lvert \mathbf{s} \rvert-1})$, $(T_{\lvert \mathbf{s} \rvert-1}(a_{\lvert \mathbf{s} \rvert-1}), T_{\lvert \mathbf{s} \rvert-1}(b_{\lvert \mathbf{s} \rvert-1}))\in \Phi_{s_{\lvert \mathbf{s} \rvert-1}}$, and $T_{\lvert \mathbf{s} \rvert-1}(b_{\lvert \mathbf{s} \rvert-1})\! =\! \sigma^{\mathbf{A}}(b_{0},\ldots,b_{\lvert \mathbf{s} \rvert-1})$. Therefore $(\sigma^{\mathbf{A}}((a_{j})_{j\in\bb{\mathbf{s}}}), \sigma^{\mathbf{A}}((b_{j})_{j\in\bb{\mathbf{s}}}))\in \Phi_{s}$.
%
% We shall now proceed to verify that (2) and (3) are equivalent.
%
% Since every elementary translations on $\mathbf{A}$ is a translation on $\mathbf{A}$, it is obvious that if $\Phi$ is closed under the translations on $\mathbf{A}$, then $\Phi$ is closed under the elementary translations on $\mathbf{A}$.
%
% Reciprocally, let us suppose that $\Phi$ is closed under the elementary translations on $\mathbf{A}$. We want to show that $\Phi$ is closed under the translations on $\mathbf{A}$. Let $t$ and $s$ be elements of $S$, $x$, $y$ elements of $A_{t}$, $T\in \mathrm{Tl}_{t}(\mathbf{A})_{s}$, and let us suppose that $(x,y)\in \Phi_{t}$. Then there is an $n\in \mathbb{N}-1$, a word $\mathbf{w}\in S^{n+1}$, and a family $(T_{j})_{j\in n}$ such that $w_{0} = t$, $w_{n} = s$, for every $j\in n$, 
% $T_{j}\in \mathrm{Etl}_{w_{j}}(\mathbf{A})_{w_{j+1}}$,
% and $T = T_{n-1}\circ\cdots\circ T_{0}$. Then, from $(x,y)\in \Phi_{t} = \Phi_{w_{0}}$, we infer that $(T_{0}(x),T_{0}(y))\in \Phi_{w_{1}}$. By proceeding in the same way we, finally, come to $(T_{n-1}(\ldots(T_{0}(x))\ldots),T_{n-1}(\ldots(T_{0}(y))\ldots))\in \Phi_{s} = \Phi_{w_{n}}$, i.e., to $(T(x),T(y))\in \Phi_{s}$.
% \end{proof}
%
% %\begin{remark}
% %The notion of congruence on a $\Sigma$-algebra $\mathbf{A}$ can also be characterized by means of (1) all $\downarrow\! w$-ary term operations on $\mathbf{A}$, where $\downarrow\! w = (w^{-1}[\{s\}])_{s\in S}$ and $w\in S^{\star}-\{\lambda\}$, and (2) all $\delta^{t}$-ary algebraic operations on $\mathbf{A}$, where $t\in S$.
% %\end{remark}
%
% We next investigate the relationships between the translations and the homomorphisms between $\Sigma$-algebras. The first result is a direct consequence of Proposition~\ref{PTransCong}.
%
% \begin{corollary}\label{CTransCong} Let $\mathbf{A}$ and $\mathbf{B}$ be two $\Sigma$-algebras and $f\colon \mathbf{A}\mor \mathbf{B}$ a $\Sigma$-homomorphism. For $s,t\in S$, let $T$ be a translation on $A$ in $\mathrm{Tl}_{t}(\mathbf{A})_{s}$. Let $x,y\in A_{t}$ be such that $f_{t}(x)=f_{t}(y)$, then $f_{s}(T(x))=f_{s}(T(y))$.
% \end{corollary}
% \begin{proof}
% It follows directly from Proposition~\ref{PTransCong} and from the fact that 	$\mathrm{Ker}(f)$ is a $\Sigma$-congruence on $\mathbf{A}$.
% \end{proof}
%
%
% Now we elucidate the relationships that exist between translations and homomorphisms.
%
%
% \begin{proposition}\label{PTransHom}
% Let $\mathbf{A}$ and $\mathbf{B}$ be two $\Sigma$-algebras and $f\colon \mathbf{A}\mor \mathbf{B}$ a $\Sigma$-homomorphism. Then, for every $t$, $s\in S$ and every $T\in \mathrm{Tl}_{t}(\mathbf{A})_{s}$, there exists a $T^{f}\in\mathrm{Tl}_{t}(\mathbf{B})_{s}$ such that $f_{s}\circ T = T^{f}\circ f_{t}$. Moreover, if $f$ is an epimorphism, then, for every $t$, $s\in S$ and every $U\in\mathrm{Tl}_{t}(\mathbf{B})_{s}$, there exists a $T\in \mathrm{Tl}_{t}(\mathbf{A})_{s}$ such that $T^{f} = U$.
% \end{proposition}
%
% \begin{proof}
% If $T$ is a $t$-elementary translation of sort $s$ for $\mathbf{A}$, for some $t$ and $s$ in $S$, then there is a word $\mathbf{s}\in S^{\star}-\{\lambda\}$, an index $k\in \lvert \mathbf{s} \rvert$, an operation symbol $\sigma\in \Sigma_{\mathbf{s},s}$, a family $(a_{j})_{j\in k}\in\prod_{j\in k}A_{s_{j}}$, and a family $(a_{l})_{l\in \lvert \mathbf{s} \rvert-(k+1)}\in\prod_{l\in \lvert \mathbf{s} \rvert-(k+1)}A_{s_{l}}$ such that $s_{k} = t$ and, for every $x\in A_{t}$, $T(x) = \sigma^{\mathbf{A}}(a_{0},\ldots,a_{k-1},x,a_{k+1},\ldots,a_{\lvert \mathbf{s} \rvert-1})$. Then it suffices to take as $T^{f}$ precisely the mapping from $B_{t}$ to $B_{s}$ defined, for every $y\in B_{t}$, as follows:
% $$
% T^{f}(y) = \sigma^{\mathbf{B}}(f_{s_{0}}(a_{0}),\ldots,f_{s_{k-1}}(a_{k-1}),y,f_{s_{k+1}}(a_{k+1}),\ldots,f_{s_{\lvert \mathbf{s} \rvert-1}}(a_{\lvert \mathbf{s} \rvert-1})).
% $$
%
%
% If $T\in\mathrm{Tl}_{t}(\mathbf{A})_{s}$, for some $t$ and $s$ in $S$, and $T$ is not a $t$-elementary translation of sort $s$ for $\mathbf{A}$, then there is an $n\in \mathbb{N}-1$, a word $\mathbf{w}\in S^{n+1}$, and a family $(T_{j})_{j\in n}$ such that $w_{0} = t$, $w_{n} = s$, for every $j\in n$, 
% $T_{j}\in \mathrm{Etl}_{w_{j}}(\mathbf{A})_{w_{j+1}}$
% and $T = T_{n-1}\circ\cdots\circ T_{0}$. Then it suffices to take as $T^{f}$ precisely the mapping from $B_{t}$ to $B_{s}$ defined as
% $T^{f} = T_{n-1}^{f}\circ\cdots\circ T_{0}^{f}$. Let us note that if $T$ is $\id^{{A}_{t}}$, for some $t\in S$, then it suffices to take as $T^{f}$ precisely $\id^{{B}_{t}}$.
%
% We next prove that if $f$ is an epimorphism, then, for every $t$, $s\in S$ and every $U\in\mathrm{Tl}_{t}(\mathbf{B})_{s}$, there exists a $T\in \mathrm{Tl}_{t}(\mathbf{A})_{s}$ such that $T^{f} = U$. If $U$ is a $t$-elementary translation of sort $s$ for $\mathbf{B}$, for some $t$ and $s$ in $S$, then and there is a word $\mathbf{s}\in S^{\star}-\{\lambda\}$, an index $k\in \lvert \mathbf{s} \rvert$, an operation symbol $\sigma\in\Sigma_{\mathbf{s},s}$, a family $(b_{j})_{j\in k}\in\prod_{j\in k}B_{s_{j}}$, and a family $(b_{l})_{l\in \lvert \mathbf{s} \rvert-(k+1)}\in\prod_{l\in \lvert \mathbf{s} \rvert-(k+1)}B_{s_{l}}$ such that $s_{k} = t$ and, for every $y\in B_{t}$, $U(y) = \sigma^{\mathbf{B}}(b_{0},\ldots,b_{k-1},y,b_{k+1},\ldots,b_{\lvert \mathbf{s} \rvert-1})$. Then, since $f$ is an epimorphism, there exists a family $(a_{j})_{j\in k}\in\prod_{j\in k}A_{s_{j}}$ and a family $(a_{l})_{l\in \lvert \mathbf{s} \rvert-(l+1)}\in\prod_{l\in \lvert \mathbf{s} \rvert-(k+1)}A_{s_{l}}$ such that, for every $j\in k$, $f_{s_{j}}(a_{j}) = b_{j}$, and, for every $l\in \lvert \mathbf{s} \rvert-(k+1)$, $f_{s_{l}}(a_{l}) = b_{l}$. Then, after fixing $(a_{j})_{j\in k}$ and $(a_{l})_{l\in \lvert \mathbf{s} \rvert-(k+1)}$, it suffices to take as $T$ precisely the mapping from $A_{t}$ to $A_{s}$ defined, for every $x\in A_{t}$, as follows:
% $$
% T(x) = \sigma^{\mathbf{A}}(a_{0},\ldots,a_{k-1},x,a_{k+1},\ldots,a_{\lvert \mathbf{s} \rvert-1}).
% $$
%
% If $U\in\mathrm{Tl}_{t}(\mathbf{B})_{s}$, for some $t$ and $s$ in $S$, and $U$ is not a $t$-elementary translation of sort $s$ for $\mathbf{B}$, then there is an $n\in \mathbb{N}-1$, a word $\mathbf{w}\in S^{n+1}$, and a family $(U_{j})_{j\in n}$ such that $w_{0} = t$, $w_{n} = s$, for every $j\in n$
% $U_{j}\in \mathrm{Etl}_{w_{j}}(\mathbf{A})_{w_{j+1}}$, 
% and $U = U_{n-1}\circ\cdots\circ U_{0}$. Then, after choosing, for every $i\in n$, a $T_{i}$ such that $T_{i}^{f} = U_{i}$, it suffices to take as $T$ precisely the mapping from $A_{t}$ to $A_{s}$ defined as $T = T_{n-1}\circ\cdots\circ T_{0}$.
% \end{proof}
% %Thus we have two mappings: $\mathrm{Etl}$ which assigns to every $\Sigma$-algebra $\mathbf{A}$ the $S\times S$-sorted set $(\mathrm{Etl}_{t}(\mathbf{A})_{s})_{(t,s)\in S\times S}$, and $\mathrm{Tll}$ which assigns to every $\Sigma$-algebra $\mathbf{A}$ the $S\times S$-sorted set $(\mathrm{Tl}_{t}(\mathbf{A})_{s})_{(t,s)\in S\times S}$. Moreover, every homomorphism $f\colon \mathbf{A}\mor B$ determines two $S\times S$-sorted mappings: one from $(\mathrm{Etl}_{t}(\mathbf{A})_{s})_{(t,s)\in S\times S}$ to $(\mathrm{Etl}_{t}(\mathbf{B})_{s})_{(t,s)\in S\times S}$ and another from $(\mathrm{Tl}_{t}(\mathbf{A})_{s})_{(t,s)\in S\times S}$ to $(\mathrm{Tl}_{t}(\mathbf{B})_{s})_{(t,s)\in S\times S}$. And these $S\times S$-sorted mappings are compatible.
%
% \begin{corollary}\label{PTransHomComp}
% Let $\mathbf{A}$ and $\mathbf{B}$ be two $\Sigma$-algebras, $f\colon \mathbf{A}\mor \mathbf{B}$ a $\Sigma$-ho\-mo\-mor\-phism, $t$, $s\in S$, $T\in \mathrm{Tl}_{t}(\mathbf{A})_{s}$, $\mathbf{w}\in S^{\star}$ and $\sigma\in \Sigma_{\mathbf{w},t}$. Then $f_{s}\circ T\circ \sigma^{\mathbf{A}} = T^{f}\circ \sigma^{\mathbf{B}}\circ f_{\mathbf{w}}$.
% \end{corollary}
%
% \subsection{
% \texorpdfstring
% {The free $\Sigma$-algebra on an $S$-sorted set}
% {The free algebra on a many-sorted set}
% }\label{SAlgFree}
%
% We next state that the forgetful functor $\mathrm{G}_{\Sigma}$ from $\mathsf{Alg}(\Sigma)$ to
% $\mathsf{Set}^{S}$ has a left adjoint $\mathbf{T}_{\Sigma}$ which assigns to an $S$-sorted set $X$ the free $\Sigma$-algebra $\mathbf{T}_{\Sigma}(X)$ on $X$. We later state the universal property of the free many-sorted algebra and provide a characterization by means of the notion of Dedekind Peano-algebras. Furthermore, for free many-sorted algebras, we establish that the subelement preorder on them is an Artinian order and, for a term in a free many-sorted algebra, we define the sorted set of the subterms of it.
%
% Let us note that in what follows, to construct the algebra of $\Sigma$-rows in $X$, and the free $\Sigma$-algebra on $X$, since neither the $S$-sorted signature $\Sigma$ nor the $S$-sorted set $X$ are subject to any constraint, coproducts must necessarily be used.
%
% \begin{definition}
% Let $X$ be an $S$-sorted set. The \emph{algebra of} $\Sigma$-\emph{rows in} $X$, denoted by $\mathbf{W}_{\Sigma}(X)$, is defined as follows:
% \begin{enumerate}
% \item The underlying $S$-sorted set of $\mathbf{W}_{\Sigma}(X)$, written as $\mathrm{W}_{\Sigma}(X)$, is precisely the $S$-sorted set $((\coprod\Sigma \amalg \coprod X)^{\star})_{s\in S}$, i.e., the mapping from $S$ to $\boldsymbol{\mathcal{U}}$ constantly $(\coprod\Sigma \amalg \coprod X)^{\star}$, where $(\coprod\Sigma \amalg \coprod X)^{\star}$ is the set of all words on the set $\coprod\Sigma \amalg \coprod X$, i.e., on the set
%       $$
%       \textstyle
%       [(\bigcup_{(\mathbf{s},s)\in S^{\star}\times S}(\Sigma_{\mathbf{s},s}\times\{(\mathbf{s},s)\}))\times\{0\}]\cup
%       [(\bigcup_{s\in S}(X_{s}\times\{s\}))\times\{1\}].
%       $$
%
% \item For every $(\mathbf{s},s)\in S^{\star}\times S$, and every $\sigma\in\Sigma_{\mathbf{s},s}$, the structural operation $F_{\sigma}$ associated to $\sigma$ is the mapping from $\mathrm{W}_{\Sigma}(X)_{\mathbf{s}}$ to ${\mathrm{W}_{\Sigma}(X)}_{s}$ which sends $(P_{j})_{j\in\bb{\mathbf{s}}} \in \mathrm{W}_{\Sigma}(X)_{\mathbf{s}}$ to $(\sigma)\curlywedge\concat_{j\in\bb{\mathbf{s}}}P_{j} \in {\mathrm{W}_{\Sigma}(X)}_{s}$, where, for every $(\mathbf{s},s)\in S^{\star}\times S$, and every $\sigma\in\Sigma_{\mathbf{s},s}$,
%     %defined as follows:
% %      $$
% %      F_{\sigma}
% %      \nfunction
% %      {\mathrm{W}_{\Sigma}(X)_{w}}
% %      {{\mathrm{W}_{\Sigma}(X)}_{s}}
% %      {(P_{i})_{i\in\lvert w \rvert}}
% %      {(\sigma)\curlywedge\concat_{i\in\lvert w \rvert}P_{i}}
% %      $$
%       $(\sigma)$ stands for $(((\sigma,(\mathbf{s},s)),0))$, which is the value at $\sigma$ of the canonical mapping from $\Sigma_{\mathbf{s},s}$ to $(\coprod\Sigma \amalg \coprod X)^{\star}$.
% %      mappings indicated in the following figure
% %      $$\xymatrix@C=2.50pc@R=1pc{
% %      \Sigma_{w,s} \ar[r]^-{\inc_{\Sigma_{w,s}}} & \coprod \Sigma
% %      \ar[r]^-{\inc_{\coprod\Sigma}} &
% %      \coprod\Sigma\amalg\coprod X
% %      \ar[r]^-{\eta_{\coprod\Sigma\amalg\coprod X}} &
% %      (\coprod\Sigma \amalg \coprod X)^{\star} %\\
% %      \sigma \ar@{|->}[r]  & (\sigma,(w,s)) \ar@{|->}[r] &
% %      ((\sigma,(w,s)),0) \ar@{|->}[r]  &
% %      (((\sigma,(w,s)),0))\equiv(\sigma)
% %      }
% %      $$
% \end{enumerate}
% \end{definition}
%
% \begin{definition}\label{DAlgFree}
% The \emph{free} $\Sigma$-\emph{algebra on} an $S$-sorted set $X$, denoted by $\mathbf{T}_{\Sigma}(X)$, is the $\Sigma$-algebra determined by $\mathrm{Sg}_{\mathbf{W}_{\Sigma}(X)}((\{(x)\mid x\in X_{s}\})_{s\in S})$, the subalgebra of $\mathbf{W}_{\Sigma}(X)$ generated by $(\{(x)\mid x\in X_{s}\})_{s\in S}$, where, for every $s\in S$ and every $x\in X_{s}$, $(x)$ stands for $(((x,s),1))$, which is the value at $x$ of the canonical mapping from $X_{s}$ to $(\coprod\Sigma \amalg \coprod X)^{\star}$.
% %mappings indicated in the following figure
% %$$
% %\xymatrix@C=3pc@R=1pc{
% %X_{s} %\ar `u[r] `[rrr]^{\eta_{X}} [rrr]
% %      %\ar@(u,u)[rrr]^{\eta_{X}}
% %      \ar[r]^-{\inc_{X_{s}}} &
% %\coprod X \ar[r]^-{\inc_{\coprod X}} &
% %\coprod\Sigma\amalg\coprod X
% %\ar[r]^-{\eta_{\coprod\Sigma\amalg\coprod X}} &
% %(\coprod\Sigma \amalg \coprod X)^{\star} %\\
% %%x \ar@{|->}[r]  & (x,s) \ar@{|->}[r]  & ((x,s),1) \ar@{|->}[r]  &
% %%(((x,s),1))\equiv(x)
% %}
% %$$
% We will denote by $\mathrm{T}_{\Sigma}(X)$ the underlying $S$\nobreakdash-sorted set of $\mathbf{T}_{\Sigma}(X)$ and, for $s\in S$, we will call the elements of $\mathrm{T}_{\Sigma}(X)_{s}$ \emph{terms of type} $s$ \emph{with variables in} $X$ or  $(X,s)$-\emph{terms}.
% \end{definition}
%
% \begin{remark}
% Since $(\{(x)\mid x\in X_{s}\})_{s\in S}$ is a generating subset of $\mathbf{T}_{\Sigma}(X)$, to prove that a subset $\mathcal{T}$ of $\mathrm{T}_{\Sigma}(X)$ is equal to $\mathrm{T}_{\Sigma}(X)$ it suffices, by Proposition~\ref{PPAI}, to show: (1) $(\{(x)\mid x\in X_{s}\})_{s\in S}\subseteq \mathcal{T}$ (algebraic induction basis) and (2) $\mathcal{T}$ is a subalgebra of $\mathbf{T}_{\Sigma}(X)$ (algebraic induction step).
% \end{remark}
%
% In the many-sorted case we have, as in the single-sorted case, the following characterization of the elements of $\mathrm{T}_{\Sigma}(X)_{s}$, for $s\in S$.
%
% \begin{proposition}\label{PTermChar}
% Let $X$ be an $S$-sorted set. Then, for every $s\in S$ and every $P\in
% \mathrm{T}_{\Sigma}(X)_{s}$, we have that $P$ is a term of type $s$ with variables in $X$ if and only if $P = (x)$, for a unique $x\in X_{s}$, or $P = (\sigma)$, for a unique $\sigma\in\Sigma_{\lambda,s}$, or $P = (\sigma)\curlywedge\concat(P_{j})_{j\bb{\mathbf{s}}}$, for a unique $\mathbf{s}\in S^{\star}-\{\lambda\}$, a unique $\sigma\in\Sigma_{\mathbf{s},s}$, and a unique family $(P_{j})_{j\in\bb{\mathbf{s}}}\in\mathrm{T}_{\Sigma}(X)_{\mathbf{s}}$.
%
% Moreover, the three possibilities are mutually exclusive. From now on, for simplicity of notation, we will write $x$, $\sigma^{\mathbf{T}_{\Sigma}(X)}$, and $\sigma^{\mathbf{T}_{\Sigma}(X)}(P_{0},\ldots,P_{\bb{\mathbf{s}}-1})$ or $\sigma^{\mathbf{T}_{\Sigma}(X)}((P_{j})_{j\in \bb{\mathbf{s}}})$ instead of $(x)$, $(\sigma)$, and $(\sigma)\curlywedge\concat(P_{j})_{j\in\bb{\mathbf{s}}}$, respectively. Thus, in particular, the structural operation $F_{\sigma}$ (or more accurately $F_{\sigma}^{\mathbf{T}_{\Sigma}(X)}$) associated with $\sigma$ is identified with $\sigma^{\mathbf{T}_{\Sigma}(X)}$.
% \end{proposition}
%
% Since the terms of the last class of the proposition just stated, as well as the first letter of them, will be constantly used in this work, in the following definition, we give them special names.
%
% \begin{definition}
% If a term $P\in \mathrm{T}_{\Sigma}(X)_{s}$ of type $(X,s)$ is of the form $$\sigma^{\mathbf{T}_{\Sigma}(X)}((P_{j})_{j\in \bb{\mathbf{s}}}),$$ for a unique $\mathbf{s}\in S^{\star}-\{\lambda\}$, a unique $\sigma\in\Sigma_{\mathbf{s},s}$, and a unique family $(P_{j})_{j\in\bb{\mathbf{s}}}\in\mathrm{T}_{\Sigma}(X)_{\mathbf{s}}$, then it will be called a \emph{complex term} and the first letter of it (which is the operation symbol $\sigma$) will be called the \emph{head} of $P$ and will be denoted by $\mathrm{hd}(P)$.
% \end{definition}
%
%
% We next distinguish, from among the non-empty finite families of terms of the same type with variables in a given $S$-sorted set, those that are  head-constant. And we do such a thing because in this work some specific types of  head-constant families of terms will play an important role. 
%
% \begin{definition} \label{DHeadCt}
% Let $X$ be an $S$-sorted set, $s\in S$, $n\in \mathbb{N}-1$, and $(P_{i})_{i\in n}\in \mathrm{T}_{\Sigma}(X)_{s}^{n}$. If, for every $i\in n$, $P_{i}$ is a complex term and, for every $i$, $j\in n$, $\mathrm{hd}(P_{i}) = \mathrm{hd}(P_{j})$, then we will say that the family $(P_{i})_{i\in n}$ is \emph{head-constant}.  
% Therefore, $(P_{i})_{i\in n}$ is complex and head-constant if, and only if, there exists a unique word $\mathbf{s}\in S^{\star}-\{\lambda\}$, a unique operation
% symbol $\sigma\in \Sigma_{\mathbf{s},s}$, and a unique family of path terms $((P_{i,j})_{j\in\bb{\mathbf{s}}})_{i\in n}\in \mathrm{T}_{\Sigma}(X)_{\mathbf{s}}^{n}$ such that, for every $i\in n$,
% $$
% P_{i}=\sigma^{\mathbf{T}_{\Sigma}(X)}((P_{i,j})_{j\in\bb{\mathbf{s}}}).
% $$ 
% \end{definition}
%
%
%
% %Moreover, to shorten notation, for every $(w,s)\in S^{\star}\times S$, and every $\sigma\in\Sigma_{w,s}$, we will let $\sigma$ stand for $F_{\sigma}$, the structural operation of $\mathbf{T}_{\Sigma}(X)$ associated to $\sigma$.
%
% From Proposition~\ref{PTermChar} follows, immediately, the universal property of the free $\Sigma$-algebra on an $S$-sorted set $X$, as stated in the subsequent proposition.
%
% \begin{proposition}\label{PPropUniv}
% For every $S$-sorted set $X$, the pair $(\eta^{X},\mathbf{T}_{\Sigma}(X))$, where $\eta^{X}$, the
% \emph{insertion of (the $S$-sorted set of generators)} $X$ \emph{into} $\mathrm{T}_{\Sigma}(X)$, is the co-restric\-tion to
% $\mathrm{T}_{\Sigma}(X)$ of the canonical embedding of $X$ into $\mathrm{W}_{\Sigma}(X)$, has the following universal property: for every $\Sigma$-algebra $\mathbf{A}$ and every $S$-sorted mapping $f\colon X\mor A$, there exists a unique homomorphism $f^{\sharp}\colon\mathbf{T}_{\Sigma}(X)\mor\mathbf{A}$ such that $f^{\sharp}\circ \eta^{X} = f$.
% \end{proposition}
%
% \begin{proof}
% For every $s\in S$ and every $(X,s)$-term $P$, the $s$-th coordinate $f^{\sharp}_{s}$ of $f^{\sharp}$ is defined recursively as follows: $f^{\sharp}_{s}(x) = f_{s}(x)$, if $P = x$; $f^{\sharp}_{s}(\sigma^{\mathbf{T}_{\Sigma}(X)}) = \sigma^{\mathbf{A}}$, if $P = \sigma^{\mathbf{T}_{\Sigma}(X)}$; and, finally,  $f^{\sharp}_{s}(\sigma^{\mathbf{T}_{\Sigma}(X)}((P_{j})_{j\in\bb{\mathbf{s}}})) = \sigma^{\mathbf{A}}((f^{\sharp}_{s_{j}}(P_{j}))_{j\in\bb{\mathbf{s}}})$, if $P = \sigma^{\mathbf{T}_{\Sigma}(X)}((P_{j})_{j\in\bb{\mathbf{s}}})$.
% %$$
% %P\,\lmapsto\,f^{\sharp}_{s}(P) =
% %\begin{cases}
% %  f_{s}(x),&\text{if }P = x; \\
% %  \sigma^{\mathbf{A}},&\text{if }P = \sigma;\\
% %  F_{\sigma}^{\mathbf{A}}(f^{\sharp}_{w(0)}(P_{0}),\ldots,
% %  f^{\sharp}_{w(\lvert w \rvert-1)}(P_{\lvert w \rvert-1})),
% %  &\text{if }P = \sigma(P_{0},\ldots,P_{\lvert w \rvert-1}).
% %\end{cases}
% %$$
% \end{proof}
%
% The just stated proposition allows us to carry out definitions by algebraic recursion on a free many-sorted algebra as indeed we will be doing throughout  this paper.
%
% %The just stated proposition provides a sound basis for the principle of definition by algebraic recursion.
%
% \begin{corollary} \label{CTermAdj}
% The functor $\mathbf{T}_{\Sigma}$, which sends an $S$-sorted set $X$ to $\mathbf{T}_{\Sigma}(X)$ and an $S$-sorted mapping $f$ from $X$ to $Y$ to $f^{@} (= (\eta^{Y}\circ f)^{\sharp})$, the unique homomorphism from $\mathbf{T}_{\Sigma}(X)$ to $\mathbf{T}_{\Sigma}(Y)$ such that $f^{@}\circ\eta^{X} = \eta^{Y}\circ f$, is left adjoint for the forgetful functor $\mathrm{G}_{\Sigma}$ from $\mathsf{Alg}(\Sigma)$ to $\mathsf{Set}^{S}$.
% %$$
% %\xymatrix@=5pc{ \mathsf{Alg}(\Sigma)
% %  \ar@<1.5ex>[r]^{\mathrm{G}_{\Sigma}}
% %  \ar@{}[r]|{\uadj}
% % &
% %\mathsf{Set}^{S}
% %  \ar@<1.5ex>[l]^{\mathbf{T}_{\Sigma}}
% %}
% %$$
% \end{corollary}
%
%
% We introduce two small results that will be of interest later on.
%
% \begin{proposition}\label{PPropUnivId}
% Let  $X$ be an $S$-sorted set. Then $(\eta^{X})^{\sharp} = \id^{\mathbf{T}_{\Sigma}(X)}$ . 
% \end{proposition}
%
% \begin{figure}
%     \centering
%     \begin{tikzpicture}
% [ACliment/.style={-{To [angle'=45, length=5.75pt, width=4pt, round]}},
% RHACliment/.style={right hook-{To [angle'=45, length=5.75pt, width=4pt, round]}, font=\scriptsize}, 
% scale=1
% ]
%
% \node[] (X) at (0,0) {$X$};
% \node[] (T) at (2,0) {$\mathbf{T}_{\Sigma}(X)$};
% \node[] (A) at (2,-2) {$\mathbf{A}$};
% \node[] (B) at (2,-4) {$\mathbf{B}$};
%
% \draw[ACliment]  (X) to node [above ] {$\eta^{X}$} (T);
% \draw[ACliment]  (X) to node [below left] {$f$} (A);
% \draw[ACliment]  (T) to node [right] {$f^{\sharp}$} (A);
% \draw[ACliment]  (A) to node [right ] {$g$} (B);
%
%
% \draw[ACliment, rounded corners]
% (T.east) -- 
% ($(T)+(1,0)$) -- node [right] {$(g\circ f)^{\sharp}$} 
% ($(B)+(1,0)$) -- ($(B.east)+(0,0)$)
% ;
%
% \end{tikzpicture}
%     \caption{The mappings of Proposition~\ref{PPropUnivComp}.}
%     \label{FPropUnivComp}
% \end{figure}
%
% \begin{proposition}\label{PPropUnivComp}
% Let  $X$ be an $S$-sorted set, $\mathbf{A}$ and $\mathbf{B}$ $\Sigma$-algebras, $f$ a mapping from $X$ to $A$ and $g$ a $\Sigma$-homomorphism from $\mathbf{A}$ to $\mathbf{B}$. Then $(g \circ f)^{\sharp} = g \circ f^{\sharp}$. See the diagram in Figure~\ref{FPropUnivComp}. 
% \end{proposition}
% \begin{proof}
% By associativity and the universal property of $f^{\sharp}$, we have that
% $$
% (g \circ f^{\sharp}) \circ \eta^{X}
% =
% g \circ (f^{\sharp} \circ \eta^{X})
% =
% g \circ f.
% $$
%
% Thus, the result follows by the uniqueness of $(g \circ f)^{\sharp}$.
% \end{proof}
%
% \begin{remark}
% Every many-sorted \emph{algebra} in a variety is isomorphic to a quotient of a free algebra in the variety. 
% \end{remark}
%
% %\begin{remark}
% %The set $\{\mathbf{T}_{\Sigma}(\delta^{s})\mid s\in S\}$ is a set of separators ($\equiv$ generators) for $\mathsf{Alg}(\Sigma)$. Actually, if $\boldsymbol{\mathcal{V}}$ is a variety of $\Sigma$-algebras, then $\{\mathbf{T}_{\boldsymbol{\mathcal{V}}}(\delta^{s})\mid s\in S\}$, where, for every $s\in S$, $\mathbf{T}_{\boldsymbol{\mathcal{V}}}(\delta^{s})$ is the free $\boldsymbol{\mathcal{V}}$-algebra on $\delta^{s}$, is a set of separators for $\boldsymbol{\mathcal{V}}$ (this is so because $G_{\boldsymbol{\mathcal{V}}}$, the forgetful from $\boldsymbol{\mathcal{V}}$ to $\mathsf{Set}^{S}$, is faithful, has a left adjoint, and $\{\delta^{s}\mid s\in S\}$ is a set of separators  for $\mathsf{Set}^{S}$).
% %\end{remark}
%
% %\begin{conventions}
% %To simplify notation, terms of the form $(x)$, for some $x\in X$, will be abbreviated to $x$ and will be called \emph{variables}; terms of the form $(\sigma)$, for some $\sigma\in\Sigma_{0}$, will be abbreviated to $\sigma$ and will be called \emph{constants}; and terms of the form $(\sigma)\curlywedge\concat_{i\in n}P_{i}$, for some $n>0$, $\sigma\in\Sigma_{n}$, and $(P_{i})_{i\in n}\in \mathrm{T}_{\Sigma}(X)^{n}$ will be abbreviated to $\sigma((P_{i})_{i\in n})$. Let us note that the former proposition justifies that the different terms are sufficiently distinct from each other so as to prevent any confusion that could arise. Moreover, for abbreviation, we let, for every $n\in \mathbb{N}$ and every $\sigma\in\Sigma_{n}$, $\sigma$ stand for $F_{\sigma}^{\mathbf{T}_{\Sigma}(X)}$, the structural operation of $\mathbf{T}_{\Sigma}(X)$ associated to $\sigma$.
% %\end{conventions}
%
% It is possible to give an internal characterization of the free algebras by means of the Dedekind-Peano algebras.
%
%
% \begin{definition}
% Let $\mathbf{A}$ be a $\Sigma$-algebra.  We will say that $\mathbf{A}$ is a \emph{Dedekind-Peano $\Sigma$-algebra}, abbreviated to DP $\Sigma$-\emph{algebra} when this is unlikely to cause confusion, if the following axioms hold
% \begin{enumerate}
% \item[DP1.] For every $(\mathbf{s},s)\in S^{\star}\times S$ and every $\sigma\in\Sigma_{\mathbf{s},s}$, $\sigma^{\mathbf{A}}\colon A_{\mathbf{s}}\mor A_{s}$ is injective.
%
% \item[DP2.] For every $s\in S $ and every $\sigma, \tau \in\Sigma_{\cdot,s}$, if $\sigma\neq\tau$, then $\mathrm{Im}(\sigma^{\mathbf{A}})\cap \mathrm{Im}(\tau^{\mathbf{A}}) = \varnothing$.
%
% \item[DP3.] $\mathrm{Sg}_{\mathbf{A}}(A-(\bigcup_{\sigma\in\Sigma_{\cdot,s}}\mathrm{Im}(\sigma^{\mathbf{A}})_{s\in S}) = A$.
% \end{enumerate}
% We call the $S$-sorted set $A-(\bigcup_{\sigma\in\Sigma_{\cdot,s}}\mathrm{Im}(\sigma^{\mathbf{A}})_{s\in S})$
% the \emph{basis of Dedekind-Peano of} $\mathbf{A}$, and we denote it by $\mathrm{B}(\mathbf{A})$.
% \end{definition}
%
% We remark that the axioms DP1 and DP2 are equivalent to the following axiom
% \begin{enumerate}
% \item[DP4.] For every $s\in S$, $\sigma, \tau \in\Sigma_{\cdot,s}$, $a\in A_{\mathrm{ar}(\sigma)}$, $b\in
% A_{\mathrm{ar}(\tau)}$, if $\sigma^{\mathbf{A}}(a)=\tau^{\mathbf{A}}(b)$ then $\sigma=\tau$ and $a=b$.
% \end{enumerate}
%
% Moreover, if $\mathbf{A}$ is a DP $\Sigma$-algebra, then $\mathrm{B}(\mathbf{A})$ is $\bigcap\{X\subseteq A\mid \mathrm{Sg}_{\mathbf{A}}(X)=A\}$
%
% \begin{proposition}
% Let $\mathbf{A}$ be a DP $\Sigma$-algebra. Then $\mathbf{A}$ is isomorphic
% to $\mathbf{T}_{\Sigma}(\mathrm{B}(\mathbf{A}))$.
% \end{proposition}
%
% \begin{proposition}
% Let $X$ be an $S$-sorted set. Then $\mathbf{T}_{\Sigma}(X)$ is a DP $\Sigma$-algebra.
% \end{proposition}
%
%
%
%
%
%
%
%
%
% \begin{proposition}
% Let $\mathbf{A}$ be a DP $\Sigma$-algebra. Then $\leq_{\mathbf{A}}$ is antisymmetric and does not have strictly descending $\omega_{0}$-chains, i.e., is an Artinian order.
% \end{proposition}
%
% \begin{corollary}
% Let $X$ be an $S$-sorted set. Then $\leq_{\mathbf{T}_{\Sigma}(X)}$ is antisymmetric and does not have strictly descending $\omega_{0}$-chains, i.e., is an Artinian order.
% \end{corollary}
%
% \begin{remark}\label{RTermOrd}
% Let us recall from Definition~\ref{DAlgOrd} that $\prec_{\mathbf{T}_{\Sigma}(X)}$ denotes the binary relation on $\coprod \mathrm{T}_{\Sigma}(X)$ consisting of the ordered pairs $((Q,t),(P,s))\in(\coprod \mathrm{T}_{\Sigma}(X))^{2}$ for which there exists a word $\mathbf{s}\in S^{\star}-\{\lambda\}$, an operation symbol $\sigma\in \Sigma_{\mathbf{s},s}$ and a family of terms $(P_{j})_{j\in\bb{\mathbf{s}}}\in \mathrm{T}_{\Sigma}(X)_{\mathbf{s}}$ such that $P=\sigma^{\mathbf{T}_{\Sigma}(X)}((P_{j})_{j\in\bb{\mathbf{s}}})$ and, for some $j\in\bb{\mathbf{s}}$, $s_{j}=t$ and $Q=P_{j}$. 
%
% We denote by $\leq_{\mathbf{T}_{\Sigma}(X)}$ the reflexive and transitive closure of $\prec_{\mathbf{T}_{\Sigma}(S)}$, i.e., the preorder on $\coprod \mathrm{T}_{\Sigma}(X)$  generated by $\prec_{\mathbf{T}_{\Sigma}(X)}$, and by $<_{\mathbf{T}_{\Sigma}(X)}$ the transitive closure of $\prec_{\mathbf{T}_{\Sigma}(X)}$.
%
% Note that the binary relation $\prec_{\mathbf{T}_{\Sigma}(X)}$ on $\coprod \mathrm{T}_{\Sigma}(X)$ is irreflexive.
%
% The preorder $\leq_{\mathbf{T}_{\Sigma}(X)}$ on $\coprod \mathrm{T}_{\Sigma}(X)$ is $\bigcup_{n\in\mathbb{N}}\prec_{\mathbf{T}_{\Sigma}(X)}^{n}$, where $\prec_{\mathbf{T}_{\Sigma}(X)}^{0}$ is the diagonal of $\coprod \mathrm{T}_{\Sigma}(X)$ and, for $n\in\mathbb{N}$, 
% $$\prec_{\mathbf{T}_{\Sigma}(X)}^{n+1} = \prec_{\mathbf{T}_{\Sigma}(X)}^{n}\circ \prec_{\mathbf{T}_{\Sigma}(X)}.$$ 
%
% Thus, for every $((Q,t),(P,s))\in (\mathbf{T}_{\Sigma}(X))^{2}$, $(Q,t)\leq_{\mathbf{T}_{\Sigma}(X)}(P,s)$ if and only if $s = t$ and $Q=P$ or there exists  a natural number $m\in\mathbb{N}-\{0\}$, a word $\mathbf{c}\in S^{\star}$ of length $\bb{\mathbf{c}}=m+1$, and a family of terms $(R_{k})_{k\in\bb{\mathbf{c}}}$ in $\mathrm{T}_{\Sigma}(X)_{\mathbf{c}}$, such that $c_{0}=t$,  $R_{0}=Q$, $c_{m}=s$, $R_{m}=P$ and, for every $k\in m$, $(R_{k},c_{k})\prec_{\mathrm{T}_{\Sigma}(X)}(R_{k+1},c_{k+1})$. Moreover, 
% $<_{\mathbf{T}_{\Sigma}(X)}$, the transitive closure of $\prec_{\mathbf{T}_{\Sigma}(X)}$, is 
% $\bigcup_{n\in\mathbb{N}-\{0\}}\prec_{\mathbf{T}_{\Sigma}(X)}^{n}$.
% \end{remark}
%
% %\begin{definition}
% %Let $\mathbf{A}$ be a DP $\Sigma$-algebra, $s\in S$, and $a\in A_{s}$. Then the $S$-sorted set of all \emph{subterms of} $a$, denoted by $\mathrm{Subt}_{\mathbf{A}}(a)$ or, simply, by $\mathrm{Subt}_{\Sigma}(a)$ when no confusion can arise, is defined as follows:
% %$$
% %\mathrm{Subt}_{\Sigma}(a) = (\{b\in A_{t} \mid (b,t)
% %\leq_{\mathbf{A}} (a,s)\})_{t \in S }.
% %$$
% %\end{definition}
%
% \begin{definition}\label{DSubt}
% Let $X$ be an $S$-sorted set, $s\in S$, and $P\in\mathrm{T}_{\Sigma}(X)_{s}$. Then the $S$-sorted set of all \emph{subterms of} $P$, denoted by $\mathrm{Subt}_{\mathbf{T}_{\Sigma}(X)}(P)$ or, simply, by $\mathrm{Subt}_{\Sigma}(P)$ when no confusion can arise, is defined as follows:
% $$
% \mathrm{Subt}_{\Sigma}(P) = (\{Q\in \mathrm{T}_{\Sigma}(X)_{t} \mid (Q,t)
% \leq_{\mathbf{T}_{\Sigma}(X)} (P,s)\})_{t \in S }.
% $$
% \end{definition}
%
%
% \begin{remark}\label{RSubt}
% $\mathrm{Subt}_{\Sigma}(P)\in \mathrm{Sub}_{\mathrm{f}}(\mathrm{T}_{\Sigma}(X))$. Moreover, $\mathrm{Subt}_{\Sigma}(P)$ can also be characterized as the smallest subset $\mathcal{L}$ of $\mathrm{T}_{\Sigma}(X)$ which satisfies the following conditions: 
% \begin{enumerate}
% \item $P\in \mathcal{L}_{s}$; and
% \item for every $(\mathbf{s},s)\in S^{\star}\times S$, every $\sigma\in \Sigma_{\mathbf{s},s}$, and every $(Q_{j})_{j\in\bb{\mathbf{s}}}\in \mathrm{T}_{\Sigma}(X)_{\mathbf{s}}$, if $\sigma^{\mathbf{T}_{\Sigma}(X)}((Q_{j})_{j\in\bb{\mathbf{s}}})\in \mathcal{L}_{s}$, then, for every $j\in\bb{\mathbf{s}}$, $Q_{j}\in \mathcal{L}_{s_{j}}$.
% \end{enumerate}
% Note that the second condition is exactly the converse of the defining condition of the concept of subalgebra of a $\Sigma$-algebra.
% \end{remark}
%
% \begin{definition} Let $X$ be an $S$-sorted set and let $\mathbf{T}_{\Sigma}(X)$ be the free many-sorted $\Sigma$-algebra on $X$. For a natural number $n\in\mathbb{N}$ we will use the shorthand $\mathrm{B}^{n}_{\Sigma}(X)$ for the $n$-th Borel set of $X$ in $\mathbf{T}_{\Sigma}(X)$, that is 
% $$
% \mathrm{B}^{n}_{\Sigma}(X)=\mathrm{B}^{n}_{\mathbf{T}_{\Sigma}(X)}(X).
% $$
% Let us note that the elements of $\mathrm{T}_{\Sigma}(X)-\mathrm{B}^{0}_{\Sigma}(X)$ are precisely the complex terms.
% \end{definition}
%
% Following this we associate to every term for $\Sigma$ of type $(X,s)$ its $S$-sorted set of variables. %This will be used, in Section 3, in the proof of diverse propositions.
%
% \begin{definition}
% Let $X$ be an $S$-sorted set.  Then $\mathbf{Fin}(X)$ is the $\Sigma$-algebra which has as underlying
% $S$-sorted set $\mathrm{Fin}^{S}(X) = (\mathrm{Sub}_{\mathrm{f}}(X))_{s\in S}$, i.e., the $S$-sorted set constantly $\mathrm{Sub}_{\mathrm{f}}(X)$, and, for every $(\mathbf{s},s)\in S^{\star}\times S$ and every $\sigma\in \Sigma_{\mathbf{s},s}$, as operation $\sigma^{\mathbf{Fin}(X)}\colon \mathrm{Sub}_{\mathrm{f}}(X)^{\bb{\mathbf{s}}}\mor
% \mathrm{Sub}_{\mathrm{f}}(X)$ that one defined, for every $(K^{j})_{j\in\bb{\mathbf{s}}}\in \mathrm{Sub}_{\mathrm{f}}(X)^{\bb{\mathbf{s}}}$, as $\sigma^{\mathbf{Fin}(X)}((K^{j})_{j\in\bb{\mathbf{s}}})= \bigcup_{j\in\bb{\mathbf{s}}}K^{j}$. Let
% $\delta^{S,X}=(\delta^{S,X}_{s})_{s\in S}$ be the $S$-sorted mapping from $X$ to $\mathrm{Fin}^{S}(X)$ defined in Remark~\ref{RDelta}. 
% Then we will denote by $\mathrm{Var}^{S,X}$ the unique homomorphism $\ext{(\delta^{S,X})}$ from $\mathbf{T}_{\Sigma}(X)$ to
% $\mathbf{Fin}(X)$ such that the diagram in Figure~\ref{FVar} commutes. 
% For a sort $s\in S$ and a term $P\in \mathrm{T}_{\Sigma}(X)_{s}$ we will call $\mathrm{Var}^{S,X}_{s}(P) (\in \mathrm{Sub}_{\mathrm{f}}(X))$ the $S$-sorted set of \emph{variables of} $P$ (or \emph{occurring in} $P$). Moreover, when this is unlikely to cause confusion, we will write $\mathrm{Var}^{X}(P)$ or, simply, $\mathrm{Var}(P)$ for $\mathrm{Var}^{S,X}_{s}(P)$.
% \end{definition}
%
% \begin{remark}
% Let $s$ be a sort in $S$ and $P\in \mathrm{T}_{\Sigma}(X)_{s}$. Then $\mathrm{Var}^{S,X}_{s}(x) = \delta^{s,x}$, if $P = x$, for a unique $x\in X_{s}$; $\mathrm{Var}^{S,X}_{s}(\sigma^{\mathbf{T}_{\Sigma}(X)}) = \varnothing^{S}$, if $P = \sigma^{\mathbf{T}_{\Sigma}(X)}$, for a unique $\sigma \in \Sigma_{\lambda,s}$; and, finally, $\mathrm{Var}^{S,X}_{s}(\sigma^{\mathbf{T}_{\Sigma}(X)}((P_{j})_{j\in\bb{\mathbf{s}}})) = \bigcup_{j\in\bb{\mathbf{s}}}\mathrm{Var}^{S,X}_{s_{j}}(P_{j})$, if $P = \sigma^{\mathbf{T}_{\Sigma}(X)}((P_{j})_{j\in\bb{s}})$, for a unique $\mathbf{s}\in S^{\star}-\{\lambda\}$, a unique $\sigma \in \Sigma_{\mathbf{s},s}$ and a unique family $(P_{j})_{j\in\bb{\mathbf{s}}}\in \mathrm{T}_{\Sigma}(X)_{\mathbf{s}}$.
% \end{remark}
%
% \begin{figure}
% $$
% \xymatrix{ X \ar[r]^-{\eta^{X}}
%   \ar[rd]_-{\delta^{S,X}}
%  &
% \mathrm{T}_{\Sigma}(X)
%   \ar[d]^{\ext{(\delta^{S,X})} = \mathrm{Var}^{S,X} = (\mathrm{Var}^{S,X}_{s})_{s\in S}} \\
%  &
% {\mathrm{Fin}^{S}}(X) = (\mathrm{Sub}_{\mathrm{f}}(X))_{s\in S}}
% $$
% \caption{Variables.}
% \label{FVar}
% \end{figure}
%
% %\subsection{Substitutions}
% %
% %In this subsection we define several substitution operators associated to a free many-sorted algebra which will be used  extensively afterwards. In particular we state in precise terms the notion of subterm substitution.
% %
% %\begin{definition}\label{DSubstOp}
% %Let $X$ be an $S$-sorted set, $s\in S$, and $P\in \mathrm{T}_{\Sigma}(X)_{s}$. Then we will denote by $\left(\!\begin{smallmatrix}(x,t)\\ \cdot\end{smallmatrix}\!\right)_{(x,t)\in \coprod X}(P)$ the mapping from $\prod_{(x,t)\in \coprod X}\mathrm{T}_{\Sigma}(X)_{t}^{\bb{P}_{x}}$ to $\mathrm{T}_{\Sigma}(X)_{s}$ that assigns to $\left((Q^{x,t}_{\alpha})_{\alpha\in \bb{P}_{x}}\right)_{(x,t)\in \coprod X}$ in $\prod_{(x,t)\in \coprod X}\mathrm{T}_{\Sigma}(X)_{t}^{\bb{P}_{x}}$ the term
% %\begin{equation}\label{ESubstOp}
% %\left(\!\begin{smallmatrix}(x,t)\$Q^{x,t}_{\alpha})_{\alpha\in \bb{P}_{x}}\end{smallmatrix}\!\right)_{(x,t)\in \coprod X}(P)
% %\end{equation}
% %in $\mathrm{T}_{\Sigma}(X)_{s}$ obtained by substituting in $P$, for every $t\in S$, every $x\in X_{t}$, and every $\alpha\in \bb{P}_x$, $Q^{x,t}_{\alpha}$ for the $i_{\alpha}$-th occurrence of $x$ in $P$. We will call \ref{ESubstOp} \emph{the substitution of $(Q^{x,t}_{\alpha})_{\alpha\in\bb{P}_{x}}$ for $x$ in $P$ for every $t\in S$ and every $x$ in $X_{t}$}.%, and $\left(\!\begin{smallmatrix}(x,t)\\ \cdot\end{smallmatrix}\!\right)_{(x,t)\in \coprod X}(P)$ the \emph{global substitution operator for $P$}.
% %
% %Let $u$ be a sort in $S$ and $z\in X_{u}$. Then we will denote by $\left(\!\begin{smallmatrix}z\\ \cdot\end{smallmatrix}\!\right)(P)$ the mapping from $\mathrm{T}_{\Sigma}(X)_{u}^{\bb{P}_{z}}$ to $\mathrm{T}_{\Sigma}(X)_{s}$ that assigns to $(Q_{\alpha})_{\alpha\in \bb{P}_{z}}$ in $\mathrm{T}_{\Sigma}(X)_{u}^{\bb{P}_{z}}$ the term $\left(\!\begin{smallmatrix}z\$Q_{\alpha})_{\alpha\in \bb{P}_{z}}\end{smallmatrix}\!\right)(P)$ in $\mathrm{T}_{\Sigma}(X)_{s}$ obtained by substituting in $P$, for every $\alpha\in \bb{P}_{z}$, $Q_{\alpha}$ for the ${\alpha}$-th occurrence of $z$ in $P$. We will call $\left(\!\begin{smallmatrix}z\$Q_{\alpha})_{\alpha\in \bb{P}_{z}}\end{smallmatrix}\!\right)(P)$ \emph{the substitution of $(Q_{\alpha})_{\alpha\in\bb{P}_{z}}$ for $z$ in $P$}.
% %
% %Let $u$ be a sort in $S$, $z\in X_{u}$ and $\alpha\in\bb{P}_{z}$. Then we will denote by $\left(\!\begin{smallmatrix}z,\alpha\\ \cdot\end{smallmatrix}\!\right)(P)$ the mapping from $\mathrm{T}_{\Sigma}(X)_{u}$ to $\mathrm{T}_{\Sigma}(X)_{s}$ that assigns to $Q$ in $\mathrm{T}_{\Sigma}(X)_{u}$ the term $\left(\!\begin{smallmatrix}z,\alpha\\Q\end{smallmatrix}\!\right)(P)$ in $\mathrm{T}_{\Sigma}(X)_{s}$ obtained by substituting in $P$, $Q$ for the ${\alpha}$-th occurrence of $z$ in $P$. We will call $\left(\!\begin{smallmatrix}z,\alpha\\Q\end{smallmatrix}\!\right)(P)$ \emph{the $\alpha$-th substitution of $Q$ for $z$ in $P$}.
% %\end{definition}
% %
% %\begin{remark}
% %We have that $\left(\!\begin{smallmatrix}z\\ \cdot\end{smallmatrix}\!\right)(P)$ is, essentially, the restriction of $\left(\!\begin{smallmatrix}(x,t)\\ \cdot\end{smallmatrix}\!\right)_{(x,t)\in \coprod X}(P)$ to $\mathrm{T}_{\Sigma}(X)_{u}^{\bb{P}_{z}}$. Similarly,  $\left(\!\begin{smallmatrix}z,\alpha\\ \cdot\end{smallmatrix}\!\right)(P)$ is, essentially, the restriction of $\left(\!\begin{smallmatrix}z\\ \cdot\end{smallmatrix}\!\right)(P)$ to $\mathrm{T}_{\Sigma}(X)_{u}$.
% %\end{remark}
% %
% %We now state that, for every $s\in S$ and every $P\in \mathrm{T}_{\Sigma}(X)$, $\left(\!\begin{smallmatrix}(x,t)\\ \cdot\end{smallmatrix}\!\right)_{(x,t)\in \coprod X}(P)$ is, actually, a mapping from $\prod_{(x,t)\in \coprod X}\mathrm{T}_{\Sigma}(X)_{t}^{\bb{P}_{x}}$ to $\mathrm{T}_{\Sigma}(X)_{s}$. For a proof of this result see~\cite{CVCL18}, p.~28, Proposition~3.5.
% %
% %\begin{proposition}
% %Let $X$ be an $S$-sorted set, $s\in S$, and $P$ a term in $\mathrm{T}_{\Sigma}(X)_{s}$. Then, for every $\left((Q^{x,t}_{\alpha})_{\alpha\in \bb{P}_{x}}\right)_{(x,t)\in \coprod X}$ in $\prod_{(x,t)\in \coprod X}\mathrm{T}_{\Sigma}(X)_{t}^{\bb{P}_{x}}$, we have that $\left(\!\begin{smallmatrix}(x,t)\$Q^{x,t}_{\alpha})_{\alpha\in \bb{P}_{x}}\end{smallmatrix}\!\right)_{(x,t)\in \coprod X}(P)$ is a term in $\mathrm{T}_{\Sigma}(X)_{s}$. Therefore $\left(\!\begin{smallmatrix}(x,t)\\ \cdot\end{smallmatrix}\!\right)_{(x,t)\in \coprod X}(P)$ is a mapping from $\prod_{(x,t)\in \coprod X}\mathrm{T}_{\Sigma}(X)_{t}^{\bb{P}_{x}}$ to $\mathrm{T}_{\Sigma}(X)_{s}$.
% %\end{proposition}
% %
% %\begin{corollary}
% %Let $X$ be an $S$-sorted set, $u\in S$, $z\in X_{u}$, and $P$ a term in $\mathrm{T}_{\Sigma}(X)_{s}$. Then, for every
% %$(Q_{\alpha})_{\alpha\in \bb{P}_{z}}$ in $\mathrm{T}_{\Sigma}(X)_{u}^{\bb{P}_{z}}$, we have that  $\left(\!\begin{smallmatrix}z\$Q_{\alpha})_{\alpha\in \bb{P}_{z}}\end{smallmatrix}\!\right)(P)$ is a term in $\mathrm{T}_{\Sigma}(X)_{s}$. Therefore $\left(\!\begin{smallmatrix}z\\ \cdot\end{smallmatrix}\!\right)(P)$ is a mapping from $\mathrm{T}_{\Sigma}(X)^{\bb{P}_{z}}_{u}$ to $\mathrm{T}_{\Sigma}(X)_{s}$. Moreover, for an $\alpha\in\bb{P}_{z}$, $\left(\!\begin{smallmatrix}z,\alpha\\ \cdot\end{smallmatrix}\!\right)(P)$ is a mapping from $\mathrm{T}_{\Sigma}(X)_{u}$ to $\mathrm{T}_{\Sigma}(X)_{s}$.
% %\end{corollary}
% %
% %Substitutions can also be performed on proper subterms.
% %
% %\begin{figure}
% %$$
% %\begin{aligned}
% %\xymatrix@C=40pt{
% %X
% %\ar[r]^-{\eta^{X}}
% %\ar[d]_-{\mathrm{in}^{X}}
% %&
% %\mathrm{T}_{\Sigma}(X)
% %\ar[d]^-{(\mathrm{in}^{X})^{@}}
% %\\
% %X\cup\delta^{M,t}
% %\ar[r]_-{\eta^{X\cup\delta^{M,t}}}
% %&
% %\mathrm{T}_{\Sigma}(X\cup\delta^{M,t})
% %}
% %%\end{aligned}
% %\quad
% %%\begin{aligned}
% %\xymatrix@C=50pt{
% %X\cup\delta^{M,t}
% %\ar[dr]_-{\mathrm{in}^{X\cup\delta^{M,t}}\,\,\,\,}
% %\ar[r]^-{\eta^{X\cup\delta^{M,t}}}
% %&
% %\mathrm{T}_{\Sigma}(X\cup\delta^{M,t})
% %\ar[d]^-{(\mathrm{in}^{X\cup\delta^{M,t}})^{\sharp}}
% %\\
% %&
% %\mathrm{T}_{\Sigma}(X)
% %}
% %\end{aligned}
% %$$
% %\caption{Subterm substitution.}
% %\label{FSubstOp}
% %\end{figure}
% %
% %\begin{definition}
% %Let $X$ be an $S$-sorted set, $s,t\in S$, $P\in\mathrm{T}_{\Sigma}(X)_{s}$, and 
% %$M\in\mathrm{Subt}_{\Sigma}(P)_{t}$. Then there exists a unique homomorphism $(\mathrm{in}^{X})^{@}$ from $\mathbf{T}_{\Sigma}(X)$ to $\mathbf{T}_{\Sigma}(X\cup\delta^{M,t})$ and a unique homomorphism $(\mathrm{in}^{X\cup\delta^{M,t}})^{\sharp}$ from $\mathbf{T}_{\Sigma}(X\cup\delta^{M,t})$ to $\mathbf{T}_{\Sigma}(X)$ such that the  diagrams in Figure~\ref{FSubstOp} commute.
% %
% %Therefore $(\mathrm{in}^{X\cup\delta^{M,t}})^{\sharp}\circ (\mathrm{in}^{X})^{@} = \id^{\mathbf{T}_{\Sigma}(X)}$. Hence the set $((\mathrm{in}^{X\cup\delta^{M,t}})^{\sharp})_{s}^{-1}[\{P\}]$ is not empty. Then we let $\xi_{M}(P)$ be the term in $\mathrm{T}_{\Sigma}(X\cup\delta^{M,t})_{s}$ such that
% %\begin{enumerate}
% %\item $(\mathrm{in}^{X\cup\delta^{M,t}})^{\sharp}_{s}(\xi_{M}(P)) = P$ \text{ and}
% %\item $M\not\in(\mathrm{Subt}_{\Sigma}(\xi_{M}(P))\cap (\mathrm{T}_{\Sigma}(X\cup\delta^{M,t})-(X\cup\delta^{M,t})))_{t}$.
% %\end{enumerate}
% %(Informally speaking, $\xi_{M}(P)$ is $(\mathrm{in}^{X})^{@}_{s}(P)$ but changing the status of all occurrences of $M$ in $P$ from subterm to variable.)
% %
% %After having picked out the term $\xi_{M}(P)$, we let $\bb{P}_{M}$ be the set $\bb{\xi_{M}(P)}_{M}$ and, for $\alpha\in\bb{P}_{M}$, we say that $M$ occurs at the $\alpha$-th place of $P$ if $M$ occurs at the $\alpha$-th place of $\xi_{M}(P)$.
% %
% %We will denote by $\left(\!\begin{smallmatrix}M\\ \cdot\end{smallmatrix}\!\right)(P)$ the mapping from $\mathrm{T}_{\Sigma}(X)_{t}^{\bb{P}_{M}}$ to $\mathrm{T}_{\Sigma}(X)_{s}$ that assigns to $(N_{\alpha})_{\alpha\in \bb{P}_{M}}$ in $\mathrm{T}_{\Sigma}(X)_{t}^{\bb{P}_{M}}$ the term $\left(\!\begin{smallmatrix}M\$N_{\alpha})_{\alpha\in \bb{P}_{M}}\end{smallmatrix}\!\right)(P)$ in $\mathrm{T}_{\Sigma}(X)_{s}$ obtained as:
% %$$
% %\left(\!\begin{smallmatrix}M\$N_{\alpha})_{\alpha\in \bb{P}_{M}}\end{smallmatrix}\!\right)(P)=
% %(\mathrm{in}^{X\cup\delta^{M,t}})^{\sharp}_{s}
% %\left(
% %\left(\!\begin{smallmatrix}M\$(\mathrm{in}^{X})^{@}_{t}(N{\alpha}))_{\alpha\in \bb{P}_{M}}\end{smallmatrix}\!\right)(\xi_{M}(P))
% %\right).
% %$$
% %That is, $\left(\!\begin{smallmatrix}M\$N_{\alpha})_{\alpha\in \bb{P}_{M}}\end{smallmatrix}\!\right)(P)$ is obtained
% %by applying $(\mathrm{in}^{X\cup\delta^{M,t}})^{\sharp}_{s}$ to an ordinary substitution $\left(\!\begin{smallmatrix}M\$\mathrm{in}^{X})^{@}_{t}(N_{\alpha}))_{\alpha\in \bb{P}_{M}}\end{smallmatrix}\!\right)(\xi_{M}(P))$. We will call $\left(\!\begin{smallmatrix}M\$N_{\alpha})_{\alpha\in \bb{P}_{M}}\end{smallmatrix}\!\right)(P)$ \emph{the substitution of $(N_{\alpha})_{\alpha\in \bb{P}_{M}}$ for $M$ in $P$}, and $\left(\!\begin{smallmatrix}M\\ \cdot\end{smallmatrix}\!\right)(P)$ the \emph{local substitution operator for $M$ in $P$}.
% %
% %For $\alpha\in \bb{P}_{M}$, we will denote by  $\left(\!\begin{smallmatrix}M,\alpha\\ \cdot\end{smallmatrix}\!\right)(P)$ the mapping from $\mathrm{T}_{\Sigma}(X)_{t}$ to $\mathrm{T}_{\Sigma}(X)_{s}$ that assigns to $N$ in $\mathrm{T}_{\Sigma}(X)_{t}$ the term $\left(\!\begin{smallmatrix}M,\alpha\\N\end{smallmatrix}\!\right)(P)$ in $\mathrm{T}_{\Sigma}(X)_{s}$ obtained as follows.
% %$$
% %\left(\!\begin{smallmatrix}M,\alpha\\N\end{smallmatrix}\!\right)(P)=
% %(\mathrm{in}^{X\cup\delta^{M,t}})^{\sharp}_{s}
% %\left(
% %\left(\!\begin{smallmatrix}M,\alpha\$\mathrm{in}^{X})^{@}_{t}(N)\end{smallmatrix}\!\right)(\xi_{M}(P))
% %\right)
% %$$
% %That is, $\left(\!\begin{smallmatrix}M,\alpha\\N\end{smallmatrix}\!\right)(P)$ is obtained by applying $(\mathrm{in}^{X\cup\delta^{M,t}})^{\sharp}_{s}$ to an ordinary $\alpha$-th substitution $\left(\!\begin{smallmatrix}M,\alpha\$\mathrm{in}^{X})^{@}_{t}(N)\end{smallmatrix}\!\right)(\xi_{M}(P))$. We will call $\left(\!\begin{smallmatrix}M,\alpha\\N\end{smallmatrix}\!\right)(P)$ \emph{the $\alpha$-th substitution of $N$ for $M$ in $P$}, and $\left(\!\begin{smallmatrix}M\\ \cdot\end{smallmatrix}\!\right)(P)$ the \emph{$\alpha$-th local substitution operator for $M$ in $P$}.
% %\end{definition}
% %
% %\begin{remark} Note that, in the particular case in which $M\in X_{t}$, i.e., $M$ is already a variable, $\xi_{M}(P) = P$ and all of the previous substitutions refer to the original intended meaning.
% %\end{remark}
% %
% %\begin{remark}
% %For every $((P,s),(M,t))\in \bigcup_{(P,s)\in \coprod\mathrm{T}_{\Sigma}(X)}(\{(P,s)\}\times \coprod\mathrm{Subt}_{\Sigma}(P))$, we have that $\left(\!\begin{smallmatrix}M\\ \cdot\end{smallmatrix}\!\right)(P)$ is a mapping from $\mathrm{T}_{\Sigma}(X)_{t}^{\bb{P}_{M}}$ to $\mathrm{T}_{\Sigma}(X)_{s}$. Therefore, the family of these local substitution operators is an element of
% %$$
% %\textstyle
% %\prod_{((P,s),(M,t))\in \bigcup_{(P,s)\in \coprod\mathrm{T}_{\Sigma}(X)}(\{(P,s)\}\times \coprod
% %\mathrm{Subt}_{\Sigma}(P))}\mathrm{Hom}\left(\mathrm{T}_{\Sigma}(X)_{t}^{\bb{P}_{M}},\mathrm{T}_{\Sigma}(X)_{s}\right).
% %$$
% %%Hacer lo mismo para $\left(\!\begin{smallmatrix}M,\alpha\\ \cdot\end{smallmatrix}\!\right)(P)$.
% %\end{remark}
% %
% %We give an end to this section by showing that substituting terms by terms in the kernel of a homomorphism does not alter the overall evaluation of the resulting terms.
% %
% %\begin{lemma}\label{LSubstOp}
% %Let $X$ be an $S$-sorted set, $s,t\in S$, $P\in\mathrm{T}_{\Sigma}(X)_{s}$, $M\in\mathrm{T}_{\Sigma}(X)_{t}$, $\mathbf{A}$ a $\Sigma$-algebra, and $g$ a homomorphism from $\mathbf{T}_{\Sigma}(X)$ to $\mathbf{A}$. Let $\beta=\bb{P}_{M}$. Then, for every pair of families $(N_{\alpha})_{\alpha\in \beta}$, $(N'_{\alpha})_{\alpha\in \beta}$ in $ \mathrm{T}_{\Sigma}(X)_{t}^{\beta}$ satisfying that, for every $\alpha\in\beta$ we have that $g_{t}(N_{\alpha})=g_{t}(N'_{\alpha})$, then
% %$$
% %g_{s}\left(\left(\!\begin{smallmatrix}M\$N_{\alpha})_{\alpha\in \beta}\end{smallmatrix}\!\right)(P)\right)
% %=
% %g_{s}\left(\left(\!\begin{smallmatrix}M\$N'_{\alpha})_{\alpha\in \beta}\end{smallmatrix}\!\right)(P)\right).
% %$$
% %\end{lemma}
% %
% %\begin{proof}
% %Before providing a complete proof of the lemma, we will study two important cases that will allow us to consider extra assumptions on the relation between $P$ and $M$.
% %
% %\textsf{Case~1.} Assume that $M\not\in\mathrm{Subt}_{\Sigma}(P)_{t}$, i.e., we are in the case in which $M$ is not a subterm of type $t$ of $P$, then the substitution above leaves $P$ invariant. In this case, the following equalities hold
% %$$
% %g_{s}\left(\left(\!\begin{smallmatrix}M\$N_{\alpha})_{\alpha\in \beta}\end{smallmatrix}\!\right)(P)\right)
% %=g_{s}(P)=
% %g_{s}\left(\left(\!\begin{smallmatrix}M\$N'_{\alpha})_{\alpha\in \beta}\end{smallmatrix}\!\right)(P)\right).
% %$$
% %
% %Therefore, we can work under the following assumption.
% %
% %\begin{assumption}\label{ASubstOpI} $M$ is a subterm of type $t$ in $\mathrm{Subt}_{\Sigma^{\mathrm{c},\mathcal{A}}}(P)_{t}$.
% %\end{assumption}
% %
% %\textsf{Case~2.} Assume that $t=s$ and $M=P$. In this case, the number of occurrences of $M$ as a subterm of type $s$ of $P$ is $1$. Hence $\bb{P}_{M}=1$. Thus for the given terms $N$ and $N'$ in $\mathrm{T}_{\Sigma}(X)_{s}$, since they satisfy that $g_{s}(N)=g_{s}(N')$, we have that
% %$$
% %g_{s}\left(\left(\!\begin{smallmatrix}M\\ N\end{smallmatrix}\!\right)(P)\right)
% %=g_{s}(N)=g(N')=
% %g_{s}\left(\left(\!\begin{smallmatrix}M\\N'\end{smallmatrix}\!\right)(P)\right).
% %$$
% %
% %Therefore, we can work under the following assumption.
% %
% %\begin{assumption}\label{ASubstOpII} If $t=s$, then $M\neq P$.
% %\end{assumption}
% %
% %We prove the general case by induction on the height of the term $P$ in $\mathrm{T}_{\Sigma}(X)_{s}$.
% %
% %\textsf{Base step of the inducion.}
% %
% %Let $P$ be a term in $\mathrm{B}^{0}_{\Sigma}(X)_{s}$, then 
% %$
% %\mathrm{Subt}_{\Sigma}(P)=\delta^{P,s},
% %$ that this, the unique subterm we can extract from  $P$ is $P$ itself. This case follows either from Case~1 or~2.
% %
% %\textsf{Inductive step of the induction.}
% %
% %Assume that the statement holds for terms $P$ of height up to $n$. That is, for every pair of  sorts $s,t\in S$ for every $P\in\mathrm{B}^{n}_{\Sigma}(X)_{s}$, every term $M\in\mathrm{T}_{\Sigma}(X)_{t}$, every  $\Sigma$-algebra $\mathbf{A}$, every  homomorphism  $g$ from $\mathbf{T}_{\Sigma}(X)$ to $\mathbf{A}$ and  every pair of families $(N_{\alpha})_{\alpha\in \beta}$, $(N'_{\alpha})_{\alpha\in \beta}$ in $ \mathrm{T}_{\Sigma}(X)_{t}^{\beta}$, where $\beta=\bb{P}_{M}$, satisfying that, for every $\alpha\in\beta$, $g_{t}(N_{\alpha})=g_{t}(N'_{\alpha})$, then
% %$$
% %g_{s}\left(\left(\!\begin{smallmatrix}M\$N_{\alpha})_{\alpha\in \beta}\end{smallmatrix}\!\right)(P)\right)
% %=
% %g_{s}\left(\left(\!\begin{smallmatrix}M\$N'_{\alpha})_{\alpha\in \beta}\end{smallmatrix}\!\right)(P)\right).
% %$$
% %
% %We now prove the statement for the case $n+1$. Let $s,t$ be sorts in $S$, $P$ a term in $\mathrm{B}^{n+1}_{\Sigma}(X)_{s}$, $M$ a term in $\mathrm{T}_{\Sigma}(X)_{t}$, $\mathbf{A}$ a $\Sigma$-algebra,  $g$ a $\Sigma$-homomorphism from $\mathbf{T}_{\Sigma}(X)$ to $\mathbf{A}$. Let $\beta=\bb{P}_{M}$ and let $(N_{\alpha})_{\alpha\in \beta}$, $(N'_{\alpha})_{\alpha\in \beta}$ a pair of families in $ \mathrm{T}_{\Sigma}(X)_{t}^{\beta}$ satisfying that, for every $\alpha\in\beta$ it is the case that $g_{t}(N_{\alpha})=g_{t}(N'_{\alpha})$.
% %
% %Since $P$ is a term in $\mathrm{B}^{n+1}_{\Sigma}(X)_{s}$, it follows that there exists a unique word $\mathbf{s}\in S^{\star}$, a unique operation symbol $\sigma\in\Sigma_{\mathbf{s},s}$ and unique family of terms $((P_{j})_{j\in\bb{\mathbf{s}}})$ in $\mathrm{B}^{n}_{\Sigma}(X)_{\mathbf{s}}$ for which $P=\sigma^{\mathbf{T}_{\Sigma}(X)}((P_{j})_{j\in\bb{\mathbf{s}}})$.
% %
% %Under Assumption~\ref{ASubstOpI}, $M$ is a subterm of type $t$ in $\mathrm{Subt}_{\Sigma}(P)_{t}$ and, under~ Assumption~\ref{ASubstOpII}, in case $t=s$, we have that $M\neq P$. From this configuration it follows that, for $M$ being a subterm of type $t$ in $P$ whilst not being equal to $P$ in case $t=s$, then $M$ has to be a subterm of type $t$ amongst the subterms $P_{j}$, for $j\in\bb{\mathbf{s}}$,
% %$$
% %\textstyle
% %M\in\bigcup_{j\in\bb{\mathbf{s}}}\mathrm{Subt}_{\Sigma}\left(P_{j}\right)_{t}.
% %$$
% %
% %Using, for every $j\in\bb{\mathbf{s}}$, the shorthand $\beta_{j}=\bb{P_{j}}_{M}$, it follows from the previous discussion that 
% %$$
% %\textstyle
% %\beta=\bb{P}_{M}=\sum_{j\in\bb{\mathbf{s}}}\bb{P_{j}}_{M}=\sum_{j\in\bb{\mathbf{s}}}\beta_{j}
% %$$
% %
% %Moreover, we have that 
% %\allowdisplaybreaks
% %\begin{align*}
% %g_{s}\left(
% %\left(
% %\!\begin{smallmatrix}
% %M\$N_{\alpha})_{\alpha\in \beta}
% %\end{smallmatrix}\!\right)
% %\left(
% %P
% %\right)\right)
% %&=
% %g_{s}\left
% %(\left
% %(\!\begin{smallmatrix}
% %M\$N_{\alpha})_{\alpha\in \beta}
% %\end{smallmatrix}\!\right)
% %\left(
% %\sigma^{\mathbf{T}_{\Sigma}(X)}\left(\left(
% %P_{j}
% %\right)_{j\in\bb{\mathbf{s}}}\right)
% %\right)\right)
% %\tag{1}
% %\\&=
% %g_{s}\left(
% %\sigma^{\mathbf{T}_{\Sigma}(X)}
% %\left(
% %\left(
% %\left(\!\begin{smallmatrix}
% %M\\
% %\left(N_{\alpha+\sum_{k\in j}\beta_{j}}\right)_{\alpha\in \beta_{j}}
% %\end{smallmatrix}\!\right)
% %\left(
% %P_{j}
% %\right)\right)_{j\in\bb{\mathbf{s}}}
% %\right)
% %\right)
% %\tag{2}
% %\\&=
% %\sigma^{\mathbf{A}}
% %\left(
% %\left(g_{s_{j}}
% %\left(
% %\left(\!\begin{smallmatrix}
% %M\\
% %\left(
% %N_{\alpha+\sum_{k\in j}\beta_{j}}
% %\right)_{\alpha\in \beta_{j}}
% %\end{smallmatrix}\!\right)\left(
% %P_{j}
% %\right)\right)
% %\right)_{j\in\bb{\mathbf{s}}}
% %\right)
% %\tag{3}
% %\\&=
% %\sigma^{\mathbf{A}}
% %\left(
% %\left(g_{s_{j}}
% %\left(
% %\left(\!\begin{smallmatrix}
% %M\\
% %\left(
% %N'_{\alpha+\sum_{k\in j}\beta_{j}}
% %\right)_{\alpha\in \beta_{j}}
% %\end{smallmatrix}\!\right)\left(
% %P_{j}
% %\right)\right)
% %\right)_{j\in\bb{\mathbf{s}}}
% %\right)
% %\tag{4}
% %\\&=
% %g_{s}\left(
% %\sigma^{\mathbf{T}_{\Sigma}(X)}
% %\left(
% %\left(
% %\left(\!\begin{smallmatrix}
% %M\\
% %\left(
% %N'_{\alpha+\sum_{k\in j}\beta_{j}}
% %\right)_{\alpha\in \beta_{j}}
% %\end{smallmatrix}\!\right)\left(
% %P_{j}
% %\right)\right)_{j\in\bb{\mathbf{s}}}
% %\right)
% %\right)
% %\tag{5}
% %\\&=
% %g_{s}\left(
% %\left(
% %\!\begin{smallmatrix}
% %M\$N'_{\alpha})_{\alpha\in \beta}
% %\end{smallmatrix}\!\right)\left(
% %\sigma^{\mathbf{T}_{\Sigma}(X)}
% %\left(\left(
% %P_{j}\right)_{j\in\bb{\mathbf{s}}}\right)
% %\right)\right)
% %\tag{6}
% %\\&=
% %g_{s}\left(\left(\!\begin{smallmatrix}M\$N'_{\alpha})_{\alpha\in \beta}\end{smallmatrix}\!\right)(P)\right).
% %\tag{7}
% %\end{align*}
% %
% %The first equality unravels the term $P$; the second equality applies the substitution operator; the third equality follows from the fact that $g$ is a $\Sigma$-homomorphism; the fourth equality follows by induction. In fact, since, for every $j\in\bb{\mathbf{s}}$, the term $P_{j}$ has height at most $n$, i.e., $P_{j}\in\mathrm{B}^{n}_{\Sigma}(X)_{s_{j}}$, we have that, for the pair of families of terms $(N_{\alpha+\sum_{k\in j}\beta_{k}})_{\alpha\in\beta_{j}}$ and $(N'_{\alpha+\sum_{k\in j}\beta_{k}})_{\alpha\in\beta_{j}}$ in $\mathrm{T}_{\Sigma}(X)_{t}^{\beta_{j}}$, given that, for every $\alpha\in\beta_{j}$, they satisfy that $g_{t}(N_{\alpha+\sum_{k\in j}\beta_{k}})=g_{t}(N'_{\alpha+\sum_{k\in j}\beta_{k}})$, it follows, by induction, that
% %$$
% %g_{s_{j}}
% %\left(
% %\left(\!\begin{smallmatrix}
% %M\$N_{\alpha+\sum_{k\in j}\beta_{j}})_{\alpha\in \beta_{j}}
% %\end{smallmatrix}\!\right)(
% %P_{j}
% %)\right)
% %=
% %g_{s_{j}}
% %\left(
% %\left(\!\begin{smallmatrix}
% %M\$N'_{\alpha+\sum_{k\in j}\beta_{j}})_{\alpha\in \beta_{j}}
% %\end{smallmatrix}\!\right)(
% %P_{j}
% %)\right);
% %$$
% %The fifth equality follows from the fact that $g$ is a $\Sigma$-homomorphism; the sixth equation recovers the substitution operator; finally, the last equality recovers the term $P$.
% %
% %This finishes the proof.
% %\end{proof}
% %
% %As a consequence of the last lemma we have the following corollary, that states that every congruence on $\mathbf{T}_{\Sigma}(X)$ is compatible with the substitution operator.
% %
% %\begin{corollary}\label{CSubstOp}
% %Let $X$ be an $S$-sorted set, $s,t\in S$, $P\in\mathrm{T}_{\Sigma}(X)_{s}$, $M\in\mathrm{T}_{\Sigma}(X)_{t}$, and  $\beta=\bb{P}_{M}$. If $\Phi$ is a $\Sigma$-congruence on $\mathbf{T}_{\Sigma}(X)$ then, for every pair of families of terms $(N_{\alpha})_{\alpha\in \beta}$, $(N'_{\alpha})_{\alpha\in \beta}$ in $ \mathrm{T}_{\Sigma}(X)_{t}^{\beta}$ satisfying that, for every $\alpha\in\beta$, $[N_{\alpha}]_{\Phi_{t}}=[N'_{\alpha}]_{\Phi_{t}}$, then
% %$$
% %\left[\left(\!\begin{smallmatrix}M\$N_{\alpha})_{\alpha\in \beta}\end{smallmatrix}\!\right)(P)\right]_{\Phi_{s}}
% %=
% %\left[\left(\!\begin{smallmatrix}M\$N'_{\alpha})_{\alpha\in \beta}\end{smallmatrix}\!\right)(P)\right]_{\Phi_{s}}.
% %$$
% %\end{corollary} 
%
%
% \subsection{Derived operations}
%
% Next we define, for an $S$-sorted signature $\Sigma$, a $\Sigma$-algebra $\mathbf{A}$ and a word $\mathbf{s}$ on $S$, the concept of many-sorted $\mathbf{s}$-ary operation on $\mathbf{A}$, that of many-sorted $\mathbf{s}$-ary derived operation on $\mathbf{A}$, and the procedure of realization of terms $P$ of type $(\vs{\mathbf{s}},s)$ as derived operations $P^{\mathbf{A}}$ on $\mathbf{A}$ (for the meaning of the phrase  ``term of type $(\vs{\mathbf{s}},s)$'' see below). But before doing that, we investigate the relationships between some functors defined over a category associated with the free monoid on the set of sorts $S$.
%
% \begin{definition}\label{DLR}
% Let $\mathbf{S}^{\star}$ be the free monoid on $S$ and $\mathsf{C}(\mathbf{S}^{\star})$ the category whose objects are the words $\mathbf{s}$ on $S$ and whose morphisms from $\mathbf{s}$ to $\mathbf{s}'$ are the isotone and injective mappings 
% $\varphi$ from $\bb{\mathbf{s}}$ to $\bb{\mathbf{s}'}$ such that $\mathbf{s}'\circ \varphi= \mathbf{s}$. Then we let $L$ stand for the contravariant functor from  
% $\mathsf{C}(\mathbf{S}^{\star})$ to $\mathsf{Set}^{\mathsf{Set}^{S}}$ whose object mapping sends $\mathbf{s}$ to the functor $L(\mathbf{s})$ from $\mathsf{Set}^{S}$ to 
% $\mathsf{Set}$ that sends an $S$-sorted set $A$ to $A_{\mathbf{s}} = \prod_{i\in \bb{\mathbf{s}}}A_{s_{i}}$ and an $S$-sorted mapping $f\colon A\mor B$ to the unique mapping $\prod_{i\in \bb{\mathbf{s}}}f_{s_{i}}$, denoted simply as $f_{\mathbf{s}}$, from $A_{\mathbf{s}}$ to $B_{\mathbf{s}}$ such that, for every $i\in \bb{\mathbf{s}}$, the following diagram
% \begin{center}
% \begin{tikzpicture}
% [ACliment/.style={-{To [angle'=45, length=5.75pt, width=4pt, round]}},
% RHACliment/.style={right hook-{To [angle'=45, length=5.75pt, width=4pt, round]}, font=\scriptsize}, 
% scale=1
% ]
%
% \node[] (AS) at (0,0) {$A_{\mathbf{s}}$};
% \node[] (As) at (0,-2) {$A_{s_{i}}$};
% \node[] (BS) at (4,0) {$B_{\mathbf{s}}$};
% \node[] (Bs) at (4,-2) {$B_{s_{i}}$};
%
% \draw[ACliment]  (AS) to node [above ] {$f_{\mathbf{s}}$} (BS);
% \draw[ACliment]  (As) to node [below ] {$f_{s_{i}}$} (Bs);
% \draw[ACliment]  (AS) to node [left ] {$\mathrm{pr}_{i}$} (As);
% \draw[ACliment]  (BS) to node [right ] {$\mathrm{pr}_{i}$} (Bs);
% \end{tikzpicture}
% \end{center}
% commutes; and whose morphism mapping sends $\varphi\colon \mathbf{s}\mor \mathbf{s}'$ to the natural transformation $L(\varphi)$ from $L(\mathbf{s}')$ to $L(\mathbf{s})$ that assigns to an $S$-sorted set $A$ the mapping $L(\varphi)_{A}$ from $A_{\mathbf{s}'}$ to $A_{\mathbf{s}}$ that sends $a\in A_{\mathbf{s}'}$ to $a\circ\varphi = a_{\varphi}\in A_{\mathbf{s}}$. 
%
% Moreover, we let $R$ stand for the contravariant functor from $\mathsf{C}(\mathbf{S}^{\star})$ to $\mathsf{Set}^{\mathsf{Set}^{S}}$ whose object mapping sends $\mathbf{s}$ to the functor $R(\mathbf{s})$ from $\mathsf{Set}^{S}$ to $\mathsf{Set}$ that sends an $S$-sorted set $A$ to $A_{\vs{\mathbf{s}}} = \mathrm{Hom}(\vs{\mathbf{s}},A)$, where 
% $\vs{\mathbf{s}}$ is the $S$-sorted set defined, for every $s\in S$, as 
% $(\vs{\mathbf{s}})_{s} = \mathbf{s}^{-1}[\{s\}]$, and an $S$-sorted mapping $f\colon A\mor B$ to the mapping 
% $\mathrm{Hom}(\vs{\mathbf{s}},f)$, also denoted $f_{\vs{\mathbf{s}}}$, from $\mathrm{Hom}(\vs{\mathbf{s}},A)$ to $\mathrm{Hom}(\vs{\mathbf{s}},B)$ that sends $a\in \mathrm{Hom}(\vs{\mathbf{s}},A)$ to $f\circ a\in \mathrm{Hom}(\vs{\mathbf{s}},B)$, i.e., $R(\mathbf{s})$ is the hom functor $\mathrm{Hom}(\vs{\mathbf{s}},\bigcdot)$; and whose morphism mapping sends $\varphi\colon \mathbf{s}\mor \mathbf{s}'$ to the natural transformation $R(\varphi)$ from $R(\mathbf{s}')$ to $R(\mathbf{s})$ that assigns to an $S$-sorted set $A$ the mapping $R(\varphi)_{A}$ from $A_{\vs{\mathbf{s}'}}$ to $A_{\vs{\mathbf{s}}}$ that sends $a\in A_{\vs{\mathbf{s}'}}$ to $a\circ\varphi = a_{\varphi}\in A_{\vs{\mathbf{s}}}$.
% \end{definition}
%
% \begin{remark}
% If $\varphi$ is a morphism from $\mathbf{s}$ to $\mathbf{s}'$, then $\mathbf{s}'$ can be represented as:
% $$ 
% \mathbf{s}' = \mathbf{u}_{0}\bconcat s_{0}\bconcat \mathbf{u}_{1}
% \bconcat\ldots \bconcat
% \mathbf{u}_{\bb{\mathbf{s}}-1}\bconcat s_{\bb{\mathbf{s}}-1}\bconcat \mathbf{u}_{\bb{\mathbf{s}}},
% $$
% for a univocally determined family $(\mathbf{u}_{i})_{i\in \bb{\mathbf{s}}+1}$ of words on $S$ (and where, for every $i\in \bb{\mathbf{s}}$, $s_{i} = s'_{\varphi(i)}$). And, reciprocally, if for two words $\mathbf{s}$ and 
% $\mathbf{s}'$, the word $\mathbf{s}'$ can be represented as a concatenation of words as just said and satisfying the letters $s_{i}$ appropriate conditions, then there exists a unique morphism from $\mathbf{s}$ to $\mathbf{s}'$ which induces the given representation of $\mathbf{s}'$.
% \end{remark}
%
% \begin{proposition}\label{isovs}
% Let $\mathbf{s}$ be an element of $S^{\star}$ and $A$ an $S$-sorted set. Then there exists an isomorphism $\theta_{\mathbf{s},A}$ from $A_{\vs{\mathbf{s}}}$
% to $A_{\mathbf{s}}$. Moreover, for every $S$-sorted mapping $f\colon
% A\mor B$, the following diagram
% \begin{center}
% \begin{tikzpicture}
% [ACliment/.style={-{To [angle'=45, length=5.75pt, width=4pt, round]}},
% RHACliment/.style={right hook-{To [angle'=45, length=5.75pt, width=4pt, round]}, font=\scriptsize}, 
% scale=1
% ]
%
% \node[] (RA) at (0,0) {$A_{\vs{\mathbf{s}}}$};
% \node[] (RB) at (0,-2) {$B_{\vs{\mathbf{s}}}$};
% \node[] (LA) at (4,0) {$A_{\mathbf{s}}$};
% \node[] (LB) at (4,-2) {$B_{\mathbf{s}}$};
%
% \draw[ACliment]  (RA) to node [above ] {$\theta_{\mathbf{s},A}$} (LA);
% \draw[ACliment]  (RB) to node [below ] {$\theta_{\mathbf{s},B}$} (LB);
% \draw[ACliment]  (RA) to node [left ] {$f_{\vs{\mathbf{s}}}$} (RB);
% \draw[ACliment]  (LA) to node [right ] {$f_{\mathbf{s}}$} (LB);
% \end{tikzpicture}
% \end{center}
% commutes. Therefore the family $\theta = (\theta_{\mathbf{s}})_{\mathbf{s}\in S^{\star}}$ is a natural isomorphism from the functor $R$ to the functor $L$. 
% \end{proposition}
%
% %\begin{proof}
% %The isomorphism $\theta_{\mathbf{w},A}$ and its inverse are given by: %
% %$$
% %\function
% %  {A_{\vs{\mathbf{w}}}} {A_{\mathbf{w}}}
% %  {a} {\nfunction
% %          {\bb{\mathbf{w}}} {\union_{i\in\bb{\mathbf{w}}}A_{w(i)}}
% %          {i} {a_{w(i)}(i)} }
% %\quad
% %\function
% %  {A_{\mathbf{w}}} {A_{\vs{\mathbf{w}}}}
% %  {x} { \Bigl( \,
% %          \function
% %            {(\vs{\mathbf{w}})_{s}} {A_{s}}
% %            {i} {x_{i}}
% %            \Bigr)_{s\in S}
% %     }
% %$$
% %\end{proof}
%
%
%
% \begin{definition}
% We let $\vs{\,\bigcdot\,}$ stand for the functor from $\mathsf{C}(\mathbf{S}^{\star})$ to $\mathsf{Set}^{S}$ whose object mapping sends $\mathbf{s}$ to $\vs{\mathbf{s}}$ and whose morphism mapping sends $\varphi\colon \mathbf{s}\mor \mathbf{s}'$ to the $S$-sorted mapping $\vs{\varphi}$ from $\vs{\mathbf{s}}$ to $\vs{\mathbf{s}'}$ defined, for every $s\in S$ and every $i\in (\vs{\mathbf{s}})_{s}$, as  $(\vs{\varphi})_{s}(i) = \varphi(i)$. 
%
% Let $V^{S} = (V^{S}_{s})_{s\in S}$ be an arbitrary, but fixed, $S$-sorted set of variables with a countably infinity of variables in each component, say    $V^{S}_{s} = \{\,v^{s}_{n}\mid n\in \mathbb{N}\,\}$, for every $s\in S$. Then we let $V^{S}_{\,\bigcdot\,}$ stand for the functor from $\mathsf{C}(\mathbf{S}^{\star})$ to $\mathsf{Set}^{S}$ whose object mapping sends $\mathbf{s}$ to the finite subset $V^{S}_{\mathbf{s}}$ of $V^{S}$, where $V^{S}_{\mathbf{s}}$ is defined, for every $s\in S$, as
% $
% V^{S}_{\mathbf{s},s} = \{\,v^{s}_{i}\in V^{S}_{s}\mid i\in \bb{\mathbf{s}}_{s}\,\}
% $ 
% (observe that $V^{S}_{\mathbf{s},s}$ is empty for those sorts $s$ that do
% not occur in the word $\mathbf{s}$); and whose morphism mapping sends $\varphi\colon \mathbf{s}\mor \mathbf{s}'$ to the $S$-sorted mapping $V^{S}_{\varphi}$ from 
% $V^{S}_{\mathbf{s}}$ to $V^{S}_{\mathbf{s}'}$ defined, for every $s\in S$ and every 
% $j\in \bb{\mathbf{s}}_{s}$, as $V^{S}_{\varphi,s}(v^{s}_{j}) = 
% v^{s}_{j'}$, where, by Definition~\ref{DSubw}, for  $(i_{j})_{j\in\bb{\mathbf{s}}_{s}}$, the enumeration in ascending order of the occurrences of $s$ in $\mathbf{s}$, and $(i_{k})_{k\in\bb{\mathbf{s}'}_{s}}$, the enumeration in ascending order of the occurrences of $s$ in $\mathbf{s}'$, $j'$ is the unique element in $\bb{\mathbf{s}'}_{s}$ such that  $\varphi(i_{j}) = i_{j'}$.
% \end{definition}
%
% \begin{proposition}\label{isovsV}
% Let $\mathbf{s}$ be an element of $S^{\star}$. Then there exists an isomorphism 
% $\vartheta_{\mathbf{s}}$ from $\vs{\mathbf{s}}$
% to $V^{S}_{\mathbf{s}}$. Moreover, for every morphism $\varphi\colon
% \mathbf{s}\mor \mathbf{s}'$, the following diagram
% \begin{center}
% \begin{tikzpicture}
% [ACliment/.style={-{To [angle'=45, length=5.75pt, width=4pt, round]}},
% RHACliment/.style={right hook-{To [angle'=45, length=5.75pt, width=4pt, round]}, font=\scriptsize}, 
% scale=1
% ]
%
% \node[] (s) at (0,0) {$\vs{\mathbf{s}}$};
% \node[] (s') at (0,-2) {$\vs{\mathbf{s}'}$};
% \node[] (Vs) at (4,0) {$V^{S}_{\mathbf{s}}$};
% \node[] (Vs') at (4,-2) {$V^{S}_{\mathbf{s}'}$};
%
% \draw[ACliment]  (s) to node [above ] {$\vartheta_{\mathbf{s}}$} (Vs);
% \draw[ACliment]  (s') to node [below ] {$\vartheta_{\mathbf{s}'}$} (Vs');
% \draw[ACliment]  (s) to node [left ] {$\vs{\varphi}$} (s');
% \draw[ACliment]  (Vs) to node [right ] {$V^{S}_{\varphi}$} (Vs');
% \end{tikzpicture}
% \end{center}
% commutes. Therefore the family $\vartheta = (\vartheta_{\mathbf{s}})_{\mathbf{s}\in S^{\star}}$ is a natural isomorphism from the functor $\vs{\,\bigcdot\,}$ to the functor $V^{S}_{\,\bigcdot\,}$. 
% \end{proposition}
%
%
% \begin{remark}
% In virtue of the previous propositions, for every word $\mathbf{s}\in S^{\star}$ and every $S$-sorted set $A$, the sets $A_{\vs{\mathbf{s}}}$, $A_{\mathbf{s}}$ and $A_{V^{S}_{\mathbf{s}}}$ are naturally  isomorphic.
% In what follows, if there is no ambiguity, we will not distinguish notationally between the elements of $A_{\vs{\mathbf{s}}}$ and those of $A_{\mathbf{s}}$ or those of $A_{V^{S}_{\mathbf{s}}}$. Also, for simplicity of notation, sometimes, instead of using $V^{S}_{\mathbf{s}}$ we will use 
% $\vs{\mathbf{s}}$.
% \end{remark}
%
% \begin{definition}\label{wdop}
% Let $\mathbf{s}$ be an element of $S^{\star}$, $\mathbf{A}$ a $\Sigma$-algebra, $s\in S$ and $P\in \mathrm{T}_{\Sigma}(\vs{\mathbf{s}})_{s}$ a term for $\Sigma$ of type $(\vs{\mathbf{s}},s)$. Then
% \begin{enumerate}
% \item the $\Sigma$-algebra of the \emph{many-sorted} $\vs{\mathbf{s}}$-\emph{ary} \emph{operations on} $\mathbf{A}$, $\mathbf{O}_{\vs{\mathbf{s}}}(\mathbf{A})$, is 
% $\mathbf{A}^{A_{\vs{\mathbf{s}}}}$, i.e., the direct $A_{\vs{\mathbf{s}}}$-power of 
% $\mathbf{A}$, where $A_{\vs{\mathbf{s}}}$ is $\mathrm{Hom}(\vs{\mathbf{s}},A)$, the (ordinary) set of the $S$-sorted mappings from $\vs{\mathbf{s}}$ to $A$ (which is naturally isomorphic to $A_{\mathbf{s}}$). From now on, to shorten terminology, we will speak of $\mathbf{s}$-\emph{ary} \emph{operations on} $\mathbf{A}$ instead of \emph{many-sorted} $\vs{\mathbf{s}}$-\emph{ary} \emph{operations on} $\mathbf{A}$, and, for brevity, we will write $\mathbf{O}_{\mathbf{s}}(\mathbf{A})$ instead of $\mathbf{O}_{\vs{\mathbf{s}}}(\mathbf{A})$;
%
% \item the $\Sigma$-algebra of the \emph{many-sorted} $\mathbf{s}$-\emph{ary} \emph{derived operations on} $\mathbf{A}$, $\mathbf{D}_{\vs{\mathbf{s}}}(\mathbf{A})$, is the subalgebra of $\mathbf{O}_{\vs{\mathbf{s}}}(\mathbf{A})$ generated by the subfamily
% $$
% \mathcal{P}^{A}_{\vs{\mathbf{s}}} = (\mathcal{P}^{A}_{\vs{\mathbf{s}},s})_{s\in S} =
% (\{\,\mathrm{pr}^{A}_{\vs{\mathbf{s}},s,i} \mid i\in (\vs{\mathbf{s}})_{s}\,\})_{s\in S}
% $$
% of $\mathrm{O}_{\vs{\mathbf{s}}}(A) = A^{A_{\vs{\mathbf{s}}}}$, where, for every $s\in S$ and $i\in (\vs{\mathbf{s}})_{s}$, $\mathrm{pr}^{A}_{\vs{\mathbf{s}},s,i}$ is the mapping from $A_{\vs{\mathbf{s}}}$ to $A_{s}$ which sends $a\in A_{\vs{\mathbf{s}}}$ to $a_{s}(i)$ (note that the set $\{\,\mathrm{pr}^{A}_{\vs{\mathbf{s}},s,i} \mid i\in (\vs{\mathbf{s}})_{s}\,\}$ has $\mathrm{card}((\vs{\mathbf{s}})_{s})$ projections if $A_{\vs{\mathbf{s}}} \neq \varnothing$ and 
% $\mathrm{card}(A_{s}) > 1$, or a unique mapping if $A_{\vs{\mathbf{s}}} = \varnothing$ or $\mathrm{card}(A_{s}) = 1$). From now on, to shorten terminology, we will speak of $\mathbf{s}$-\emph{ary} \emph{derived operations on} $\mathbf{A}$ instead of \emph{many-sorted} 
% $\vs{\mathbf{s}}$-\emph{ary} \emph{derived operations on} $\mathbf{A}$, and, for brevity, we will write $\mathbf{D}_{\mathbf{s}}(\mathbf{A})$ instead of $\mathbf{D}_{\vs{\mathbf{s}}}(\mathbf{A})$;
%
% \item we denote by $\mathrm{d}^{\mathbf{s},\mathbf{A}}$ the unique homomorphism from 
% $\mathbf{T}_{\Sigma}(\vs{\mathbf{s}})$ to $\mathbf{O}_{\mathbf{s}}(\mathbf{A})$ such that $\mathrm{pr}^{A}_{\vs{\mathbf{s}}} = \mathrm{d}^{\mathbf{s},\mathbf{A}}\circ \eta^{\vs{\mathbf{s}}}$, where $\mathrm{pr}^{A}_{\vs{\mathbf{s}}}$ is the $S$-sorted mapping $(\mathrm{pr}^{A}_{\vs{\mathbf{s}},s})_{s\in S}$ from $\vs{\mathbf{s}}$ to 
% $\mathrm{O}_{\mathbf{s}}(A)$, with $\mathrm{pr}^{A}_{\vs{\mathbf{s}},s} = (\mathrm{pr}^{A}_{\vs{\mathbf{s}},s,i})_{i\in (\vs{\mathbf{s}})_{s}}$, for every $s\in S$. Furthermore, $P^{\mathbf{A}}$ denotes the image of $P$ under $\mathrm{d}^{\mathbf{s},\mathbf{A}}_{s}$, and we call the mapping $P^{\mathbf{A}}$ from $A_{\vs{\mathbf{s}}}$ to $A_{s}$, the \emph{derived operation on} $\mathbf{A}$ \emph{determined by} $P$, or, less frequently, the \emph{term realization of} $P$ \emph{on} $\mathbf{A}$. Moreover, to simplify the notation, we will also denote by $\mathrm{d}^{\mathbf{s},\mathbf{A}}$ the corestriction of the homomorphism $\mathrm{d}^{\mathbf{s},\mathbf{A}}\colon\mathbf{T}_{\Sigma}(\vs{\mathbf{s}})\mor \mathbf{O}_{\mathbf{s}}(\mathbf{A})$ to the subalgebra $\mathbf{D}_{\mathbf{s}}(\mathbf{A})$ of $\mathbf{O}_{\mathbf{s}}(\mathbf{A})$.
% \end{enumerate}
% \end{definition}
%
% \begin{remark}
% What we have called \emph{derived operations on} $\mathbf{A}$, following the terminology in Cohn~\cite{pC81}, pp. 145--149, are also known, for those following that one in Gr\"{a}tzer~\cite{gG08}, pp. 37--45, and J\'{o}nsson~\cite{bJ72}, pp. 83--87, as \emph{polynomial operations of} $\mathbf{A}$.
% \end{remark}
%
% %\begin{remark}
% %For $\eta^{\vs{\mathbf{s}}}\colon \vs{\mathbf{s}}\mor \mathrm{T}_{\Sigma}(\vs{\mathbf{s}})$, $\mathrm{d}^{\vs{\mathbf{s}},\mathbf{A}}[\eta^{\vs{\mathbf{s}}}[\vs{\mathbf{s}}]]$, the (direct) image of $\eta^{\vs{\mathbf{s}}}[\vs{\mathbf{s}}]$ under $\mathrm{d}^{\vs{\mathbf{s}},\mathbf{A}}$, is also $\mathbf{D}_{\mathbf{s}}(\mathbf{A})$, i.e., the derived operations on an algebra are the same as the operations determined by the terms built from variables and formal operations denoting the primitive
% %operations of the algebra.
% %\end{remark}
%
% %Since, for every word $\mathbf{w}\in 
% %S^{\star}$ and every $S$-sorted set $A$, the sets $A_{\vs{\mathbf{w}}}$, 
% %$A_{\mathbf{w}}$ and $A_{V^{S}_{\mathbf{w}}}$ are isomorphic, we will use, indistinctly,
% %$\mathbf{O}_{\vs{\mathbf{w}}}(\mathbf{A})$, 
% %$\mathbf{O}_{\mathbf{w}}(\mathbf{A})$ or 
% %$\mathbf{O}_{V^{S}_{\mathbf{w}}}(\mathbf{A})$; and   
% %$\mathbf{D}_{\vs{\mathbf{w}}}(\mathbf{A})$, 
% %$\mathbf{D}_{\mathbf{w}}(\mathbf{A})$ or
% %$\mathbf{D}_{V^{S}_{\mathbf{w}}}(\mathbf{A})$.
%
% %\begin{remark}
% %For every finite subset $X$ of $V^{S}$ there exists an $n_{X}\in \mathbb{N}$ and a word $\mathbf{s}=(s_{i})_{i\in n_{X}}$ in $S^{\star}$ such that $X\cong  V^{S}_{\mathbf{s}}$. Thus, for the finitary derived operations on a $\Sigma$-algebra $\mathbf{A}$, i.e., for the $X$-\emph{ary} derived operators on 
% %$\mathbf{A}$, with $X$ a finite subset of of $V^{S}$, and without loss of generality, instead of finite subsets of $V^{S}$ we can consider initial sections of $V^{S}$, i.e., subsets of the form $V^{S}_{\mathbf{s}}$, with $\mathbf{s}$ in $S^{\star}$. What we have just said is also valid for arbitrary finite $S$-sorted sets.
% %\end{remark}
%
% \begin{proposition}[Reciprocity law]\label{Rlw}
% Let $\mathbf{s}$ be an element of $S^{\star}$, $s\in S$, 
% $P\in \mathrm{T}_{\Sigma}(\vs{\mathbf{s}})_{s}$, $\mathbf{A}$ a $\Sigma$-algebra and $a$ an $S$-sorted mapping from $\vs{\mathbf{s}}$ to $A$. Then 
% $$
% a_{s}^{\sharp}(P) = P^{\mathbf{A}}(a).
% $$  
% \end{proposition}
%
% %\begin{corollary}\label{CTermHom}
% %    Let $\mathbf{s}$ be an element of $S^{\star}$, $s \in S$, $P \in \mathrm{T}_{\Sigma}(\vs{\mathbf{s}})_{s}$, $\mathbf{A}$ and $\mathbf{B}$ two $\Sigma$-algebras, $f \colon \mathbf{A} \mor \mathbf{B}$ a $\Sigma$-homomorphism and $(a_{j})_{j \in \bb{\mathbf{s}}}$ an element in $A_{\mathbf{s}}$. Then
% %    $$
% %        f_{s} \left(
% %        P^{\mathbf{A}}\left(
% %            (a_{j})_{j \in \bb{\mathbf{s}}}
% %        \right)
% %        \right)
% %        =
% %        P^{\mathbf{B}}\left(
% %        \left(
% %        f_{s_{j}}
% %            (a_{j})_{j \in \bb{\mathbf{s}}}
% %        \right)
% %        \right).
% %    $$
% %\end{corollary}
% %
% % \begin{proof}
% %     The following chain of equalities holds
% %     \begin{align*}
% %         f_{s} \left(
% %         P^{\mathbf{A}}\left(
% %             (a_{j})_{j \in \bb{\mathbf{s}}}
% %         \right)
% %         \right)
% %         &=
% %         f_{s} \left(
% %         \left(
% %         \left(a_{j}\right)_{j \in \bb{\mathbf{s}}}
% %         \right)^{\sharp}_{s}
% %         (P)
% %         \right)
% %         \tag{1}
% %         \\&=
% %         \left(
% %         \left(f_{s_{j}}\left(a_{j}\right)\right)_{j \in \bb{\mathbf{s}}}
% %         \right)^{\sharp}_{s}
% %         (P)
% %         \tag{2}
% %         \\&=
% %         P^{\mathbf{B}}\left(
% %         \left(
% %         f_{s_{j}}
% %             (a_{j})_{j \in \bb{\mathbf{s}}}
% %         \right)
% %         \right).
% %         \tag{3}
% %     \end{align*}
% % \end{proof}
%
% %\newpage
% \subsection{Derivors}
%
% In this subsection, after defining the variety of Hall algebras, we introduce the notion of derivor between many-sorted signatures. This will allow us, using the homomorphisms between Hall algebras, to obtain $\mathsf{Sig}_{\mathfrak{d}}$, the category of many-sorted signatures and derivors. Next, after defining a suitable contravariant functor from $\mathsf{Sig}_{\mathfrak{d}}$ to $\mathsf{Cat}$, the category of $\boldsymbol{\mathcal{U}}$-locally small categories and functors, we obtain, by means of the Grothendieck construction, $\mathsf{Alg}_{\mathfrak{d}}$, the category with objects the ordered pairs $((S,\Sigma),(A,F))$, with $(S,\Sigma)$ a many-sorted signature and $(A,F)$ a $\Sigma$-algebra, and morphisms from $((S,\Sigma),(A,F))$ to  $((T,\Lambda),(B,G))$ the ordered pairs $((\varphi,d),f)$, with $(\varphi,d)$ a derivor from $(S,\Sigma)$ to $(T,\Lambda)$ and $f$ a homomorphism of $\Sigma$-algebras from $(A,F)$ to $(B_{\varphi},G^{(\varphi,d)})$, a canonically derived algebra of $(B,G)$. 
%
% Before defining the notion of Hall algebra, we recall that
% \begin{enumerate}
% \item a finitary specification is an ordered triple $(S,\Sigma,\mathcal{E})$, where $S$ is a set of sorts, $\Sigma$ an $S$-sorted signature, and $\mathcal{E}\subseteq\mathrm{Eq}_{V^{S}}(\Sigma)=
% (\mathrm{T}_{\Sigma}(X)_{s}^{2})_{(X,s)\in \mathrm{Sub}_{\mathrm{f}}(V^{S})\times S}$, i.e., a set of finitary  $\Sigma$-equations, where $V^{S}$, the $S$-sorted set of the variables, is a fixed $S$-countably infinite $S$-sorted set; and that
% \item a $\Sigma$-algebra $\mathbf{A}$ is a $(S,\Sigma,\mathcal{E})$-algebra if $\mathbf{A}\models^{\Sigma} \mathcal{E}$, i.e., if, for every $(X,s)\in \mathrm{Sub}_{\mathrm{f}}(V^{S})\times S$ and every $(P,Q)\in \mathcal{E}_{X,s}$, $\mathbf{A}$ is a model of $(P,Q)$, in symbols $\mathbf{A}\models^{\Sigma}_{X,s} (P,Q)$, which in turn means that, for every $a\in \mathrm{Hom}(X,A)$, $a^{\sharp}_{s}(P) = a^{\sharp}_{s}(Q)$.
% \end{enumerate}
%
% \begin{definition}
% Let $S$ be a set of sorts and let $V^{\mathrm{H}_{S}} = (V^{\mathrm{H}_{S}}_{\mathbf{s},s})_{(\mathbf{s},s)\in S^{\star}\times S}$ be a fixed $S^{\star}\times S$-sorted set of variables with a countably infinity of variables in each component, say $V^{\mathrm{H}_{S}}_{\mathbf{s},s} = \{\,v^{\mathbf{s},s}_{n}\mid n\in \mathbb{N}\,\}$ (which is isomorphic to 
% $\{(\mathbf{s},s)\}\times\mathbb{N}$), for every $(\mathbf{s},s)\in S^{\star}\times S$.  A \emph{Hall algebra for} $S$ is a $\mathrm{H}_{S} = (S^{\star}\times S,\Sigma^{\mathrm{H}_{S}},\mathcal{E}^{\mathrm{H}_{S}})$-algebra, where $\Sigma^{\mathrm{H}_{S}}$ is the $S^{\star}\times S$-sorted signature, i.e., the $(S^{\star}\times S)^{\star}\times (S^{\star}\times S)$-sorted set, defined as follows:
% \begin{enumerate}
% \item[$\mathrm{HS}_{1}$.] For every $\mathbf{s}\in S^{\star}$ and $i\in\bb{\mathbf{s}}$,
%       $$
%       \pi^{\mathbf{s}}_{i}\colon\lambda\mor(\mathbf{s},s_{i}),
%       $$
%       where $\bb{\mathbf{s}}$ is the \emph{length} of the word $\mathbf{s}$ and $\lambda$ the
%       \emph{empty word} in the underlying set of the free monoid on
%       $S^{\star}\times S$.
%
% \item[$\mathrm{HS}_{2}$.] For every $\mathbf{r}$, $\mathbf{s}\in S^{\star}$ and $s\in S$,
%       $$
%       \xi_{\mathbf{r},\mathbf{s},s}\colon ((\mathbf{s},s),(\mathbf{r},s_{0}),\ldots,(\mathbf{r},s_{\bb{\mathbf{s}}-1}))\mor(\mathbf{r},s);
%       $$
% \end{enumerate}
% while $\mathcal{E}^{\mathrm{H}_{S}}$ is the sub-$(S^{\star}\times
% S)^{\star}\times (S^{\star}\times S)$-sorted set of
% $\mathrm{Eq}(\Sigma^{\mathrm{H}_{S}})$, where
% $$
% \mathrm{Eq}\left(\Sigma^{\mathrm{H}_{S}}\right) =
% \left(\mathrm{T}_{\Sigma^{\mathrm{H}_{S}}}(\vs{\ol{\mathbf{s}}})_{(\mathbf{r},s)}^{2}\right)_{(\ol{\mathbf{s}},(\mathbf{r},s))\in
% (S^{\star}\times S)^{\star}\times (S^{\star}\times S)},
% $$
% defined as follows:
% \begin{enumerate}
% \item[$\mathrm{H}_{1}$.] \emph{Projection}.
%       For every $\mathbf{r}$, $\mathbf{s}\in S^{\star}$ and $i\in\bb{\mathbf{s}}$, the equation
%       $$
%       \xi_{\mathbf{r},\mathbf{s},s_{i}}(\pi^{\mathbf{s}}_{i},v^{\mathbf{r},s_{0}}_{0},\ldots,
%            v^{\mathbf{r},s_{\bb{\mathbf{s}}-1}}_{\bb{\mathbf{s}}-1})=v^{\mathbf{r},s_{i}}_{i}
%       $$
%       of type
%       $(((\mathbf{r},s_{0}),\ldots,(\mathbf{r},s_{\bb{\mathbf{s}}-1})),(\mathbf{r},s_{i})).$
% \item[$\mathrm{H}_{2}$.] \emph{Identity}.
%       For every $\mathbf{r}\in S^{\star}$ and $j\in \bb{\mathbf{r}}$, the equation
%       $$
%       \xi_{\mathbf{r},\mathbf{r},r_{j}}(v^{\mathbf{r},r_{j}}_{j},\pi^{\mathbf{r}}_{0},\ldots,\pi^{\mathbf{r}}_{\bb{\mathbf{r}}-1})=
%            v^{\mathbf{r},r_{j}}_{j}
%       $$
%       of type
%       $(((\mathbf{r},r_{j})),(\mathbf{r},r_{j})).$
% \item[$\mathrm{H}_{3}$.] \emph{Associativity}.
%       For every $\mathbf{q}$, $\mathbf{r}$, $\mathbf{s}\in S^{\star}$ and $s\in S$, the
%       equation
% \begin{multline*}
% \xi_{\mathbf{q},\mathbf{r},s}(
%       \xi_{\mathbf{r},\mathbf{s},s}(v^{\mathbf{s},s}_{0},v^{\mathbf{r},s_{0}}_{1},\ldots,
%           v^{\mathbf{r},s_{\bb{\mathbf{s}}-1}}_{\bb{\mathbf{s}}}),
%           v^{\mathbf{q},r_{0}}_{\bb{\mathbf{s}}+1},
%           \ldots,v^{\mathbf{q},r_{\bb{\mathbf{r}}-1}}_{\bb{\mathbf{s}}+\bb{\mathbf{r}}}) = \\
%        \xi_{\mathbf{q},\mathbf{s},s}(v^{\mathbf{s},s}_{0},
%           \xi_{\mathbf{q},\mathbf{r},s_{0}}(v^{\mathbf{r},s_{0}}_{1},v^{\mathbf{q},r_{0}}_{\bb{\mathbf{s}}+1},
%                    \ldots,v^{\mathbf{q},r_{\bb{\mathbf{r}}-1}}_{\bb{\mathbf{s}}+\bb{\mathbf{r}}}),
%                    \\
%           \ldots, 
%        \xi_{\mathbf{q},\mathbf{r},s_{\bb{\mathbf{s}}-1}}(v^{\mathbf{r},s_{\bb{\mathbf{s}}-1}}_{\bb{\mathbf{s}}},
%                    v^{\mathbf{q},r_{0}}_{\bb{\mathbf{s}}+1},
%        \ldots,v^{\mathbf{q},r_{\bb{\mathbf{r}}-1}}_{\bb{\mathbf{s}}+\bb{\mathbf{r}}}))
% \end{multline*}
%      of type
%      $(((\mathbf{s},s),(\mathbf{r},s_{0}),\ldots,(\mathbf{r},s_{\bb{\mathbf{s}}-1}),
%             (\mathbf{q},r_{0}),\ldots,(\mathbf{q},r_{\bb{\mathbf{r}}-1})),(\mathbf{q},s)).$
% \end{enumerate}
%
% We call the formal constants $\pi^{\mathbf{s}}_{i}$ \emph{projections}, and the formal operations $\xi_{\mathbf{r},\mathbf{s},s}$ \emph{substitution operators}. Furthermore, we will denote by $\mathsf{Alg}(\mathrm{H}_{S})$ the category of Hall algebras for $S$ and homomorphisms between Hall algebras.  Since $\mathsf{Alg}(\mathrm{H}_{S})$ is a variety, the forgetful functor $\mathrm{G}_{\mathrm{H}_{S}}$ from $\mathsf{Alg}(\mathrm{H}_{S})$ to $\mathsf{Set}^{S^{\star}\times S}$ has a left adjoint $\mathbf{T}_{\mathrm{H}_{S}}$, situation denoted by $\mathbf{T}_{\mathrm{H}_{S}}\dashv \mathrm{G}_{\mathrm{H}_{S}}$, or diagrammatically by
% $$
% \xymatrix@=5pc{ \mathsf{Alg}(\mathrm{H}_{S})
% \ar@<1.5ex>[r]^{\mathrm{G}_{\mathrm{H}_{S}}} \ar@{}[r]|{\uadj} &
% \mathsf{Set}^{S^{\star}\times S}
% \ar@<1.5ex>[l]^{\mathbf{T}_{\mathrm{H}_{S}}} }
% $$
% which assigns to an $S^{\star}\times S$-sorted set $\Sigma$ the corresponding free Hall algebra
% $\mathbf{T}_{\mathrm{H}_{S}}(\Sigma)$.
% \end{definition}
%
% \begin{remark}
% The category $\mathsf{Alg}(\mathrm{H}_{S})$, of Hall algebras for $S$, is an algebraic version of many-sorted clone theory.
% \end{remark}
%
% \begin{remark}
% From $\mathrm{H}_{3}$, for $\mathbf{s}=\lambda$, the empty word on $S$, we get the invariance of constant functions axiom in~\cite{gm85}: For every $\mathbf{q}$, $\mathbf{r}\in S^{\star}$ and $s\in S$, we have the equation
% $$
% \xi_{\mathbf{q},\mathbf{r},s}(\xi_{\mathbf{r},\lambda,s}(v^{\lambda,s}_{0}),v^{\mathbf{q},r_{0}}_{1},
% \ldots,v^{\mathbf{q},r_{\bb{\mathbf{r}}-1}}_{\bb{\mathbf{r}}})=
% \xi_{\mathbf{q},\lambda,s}(v^{\lambda,s}_{0})
% $$
% of type $(((\lambda,s),(\mathbf{q},r_{0}),\ldots,(\mathbf{q},r_{\bb{\mathbf{r}}-1})),(\mathbf{q},s))$.
% \end{remark}
%
% For every $S$-sorted set $A$, $\mathrm{O}_{\mathrm{H}_{S}}(A) = (\mathrm{Hom}(A_{\mathbf{s}},A_{s}))_{(\mathbf{s},s)\in S^{\star}\times S}$, the $S^{\star}\times S$-sorted set of operation for $A$, is naturally equipped with a structure of Hall algebra, as stated in the following proposition, if we realize the projections as the true projections and the substitution operators as the generalized composition of mappings.
%
%
%
%
%
% \begin{proposition}
% Let $A$ be an $S$-sorted set and $\mathbf{O}_{\mathrm{H}_{S}}(A)$ the $\Sigma^{\mathrm{H}_{S}}$-algebra with underlying many-sorted set $\mathrm{O}_{\mathrm{H}_{S}}(A)$ and algebraic structure defined as follows:
% \begin{enumerate}
% \item For every $\mathbf{s}\in S^{\star}$ and $i\in\bb{\mathbf{s}}$,
%       $(\pi^{\mathbf{s}}_{i})^{\mathbf{O}_{\mathrm{H}_{S}}(A)} =
%       \mathrm{pr}^{A}_{\mathbf{s},i}\colon A_{\mathbf{s}}\mor A_{s_{i}}$.
%
% \item For every $\mathbf{r},\mathbf{s}\in S^{\star}$ and $s\in S$,
%       $\xi_{\mathbf{r},\mathbf{s},s}^{\mathbf{O}_{\mathrm{H}_{S}}(A)}$ is defined, for every $f\in
%       A_{s}^{ A_{\mathbf{s}} }$ and $g\in A^{A_{\mathbf{r}}}_{\mathbf{s}}$, as
%       $\xi_{\mathbf{r},\mathbf{s},s}^{\mathbf{O}_{\mathrm{H}_{S}}(A)}(f,g_{0},\ldots,g_{\bb{\mathbf{s}}-1})
%       = f\circ\langle g_{i}\rangle_{i\in\bb{\mathbf{s}} }$, where $\langle
%       g_{i}\rangle_{i\in\bb{\mathbf{s}}}$ is the unique mapping from $A_{\mathbf{r}}$ to
%       $A_{\mathbf{s}}$ such that, for every $i\in \bb{\mathbf{s}}$, we have that
%       $$
%       \mathrm{pr}^{A}_{\mathbf{s},i}\circ\langle g_{i}\rangle_{i\in\bb{\mathbf{s}}} = g_{i}.
%       $$
% \end{enumerate}
% Then $\mathbf{O}_{\mathrm{H}_{S}}(A)$ is a Hall algebra, the \emph{Hall algebra for} $(S,A)$.
% \end{proposition}
%
% \begin{remark}
% The closed sets of the Hall algebra $\mathbf{O}_{\mathrm{H}_{S}}(A)$ for $(S,A)$ are precisely the clones of (many-sorted) operations on the $S$-sorted set $A$. On the other hand, every $\Sigma$-algebra $\mathbf{A}$ has associated a Hall algebra. In fact, it suffices to consider $\mathbf{O}_{\mathrm{H}_{S}}(A)$, denoted by $\mathbf{O}_{\mathrm{H}_{S}}(\mathbf{A})$. Moreover, the finitary derived operations on $\mathbf{A}$ is a subalgebra of the Hall algebra $\mathbf{O}_{\mathrm{H}_{S}}(\mathbf{A})$.
% \end{remark}
%
%
%
% For every $S$-sorted signature $\Sigma$, $\mathrm{Tm}_{\mathrm{H}_{S}}(\Sigma) = (\mathrm{T}_{\Sigma}(\vs{\mathbf{s}})_{s})_{(\mathbf{s},s)\in S^{\star}\times S}$ is also e\-quipped with a structure of Hall algebra that formalizes the concept of substitution as stated in the following proposition.
%
% \begin{proposition}
% Let $\Sigma$ be an $S$-sorted signature and $\mathbf{Tm}_{\mathrm{H}_{S}}(\Sigma)$ the $\Sigma^{\mathrm{H}_{S}}$-algebra
% with underlying many-sorted set $\mathrm{Tm}_{\mathrm{H}_{S}}(\Sigma)$ and algebraic structure defined as follows:
% \begin{enumerate}
% \item For every $\mathbf{s}\in S^{\star}$ and $i\in\bb{\mathbf{s}}$,
%       $(\pi^{\mathbf{s}}_{i})^{\mathbf{Tm}_{\mathrm{H}_{S}}(\Sigma)}$
%       is the image under $\eta^{\vs{\mathbf{s}}}_{s_{i}}$ of
%       the variable $v_{i}^{s_{i}}$, where $\eta^{\vs{\mathbf{s}}} =
%       (\eta^{\vs{\mathbf{s}}}_{s})_{s\in S}$ is the canonical embedding
%       of $\vs{\mathbf{s}}$ into $\mathrm{T}_{\Sigma}(\vs{\mathbf{s}})$. Sometimes,
%       to abbreviate, we will write $\pi^{\mathbf{s}}_{i}$ instead of
%       $(\pi^{\mathbf{s}}_{i})^{\mathbf{Tm}_{\mathrm{H}_{S}}(\Sigma)}$.
%
% \item For every $\mathbf{r},\mathbf{s}\in S^{\star}$ and $s\in S$,
%       $\xi_{\mathbf{r},\mathbf{s},s}^{\mathbf{Tm}_{\mathrm{H}_{S}}(\Sigma)}$
%       is the mapping
%       $$\xi_{\mathbf{r},\mathbf{s},s}^{\mathbf{Tm}_{\mathrm{H}_{S}}(\Sigma)}\nfunction
%       {\mathrm{T}_{\Sigma}(\vs{\mathbf{s}})_{s} \times
%       \mathrm{T}_{\Sigma}(\vs{\mathbf{r}})_{s_{0}} \times \cdots \times
%       \mathrm{T}_{\Sigma}(\vs{\mathbf{r}})_{s_{\bb{\mathbf{s}}-1}}}
%       {\mathrm{T}_{\Sigma}(\vs{\mathbf{r}})_{s}}
%       {(P,(Q_{i})_{i\in\bb{\mathbf{s}}})}
%       {\left(\left(\!\begin{smallmatrix}v^{s_{i}}_{i}\\Q_{i}\end{smallmatrix}\!\right)_{i\in \bb{\mathbf{s}}}\right)^{\sharp}_{s}(P)}
%      $$
%      where, for $\left(\!\begin{smallmatrix}v^{s_{i}}_{i}\\Q_{i}\end{smallmatrix}\!\right)_{i\in \bb{\mathbf{s}}}$, the $S$-sorted mapping from $\vs{\mathbf{s}}$
%      to $\mathrm{T}_{\Sigma}(\vs{\mathbf{r}})$ canonically associated to the family
%      $(Q_{i})_{i\in\bb{\mathbf{s}}}$, $\left(\left(\!\begin{smallmatrix}v^{s_{i}}_{i}\\Q_{i}\end{smallmatrix}\!\right)_{i\in \bb{\mathbf{s}}}\right)^{\sharp}$, also denoted by $\mathcal{Q}^{\sharp}$, is the
%      unique homomorphism from $\mathbf{T}_{\Sigma}(\vs{\mathbf{s}})$ into
%      $\mathbf{T}_{\Sigma}(\vs{\mathbf{r}})$ such that
%      $$
%      \left(\left(\!\begin{smallmatrix}v^{s_{i}}_{i}\\Q_{i}\end{smallmatrix}\!\right)_{i\in \bb{\mathbf{s}}}\right)^{\sharp}\circ \eta^{\vs{\mathbf{s}}} = \left(\!\begin{smallmatrix}v^{s_{i}}_{i}\\Q_{i}\end{smallmatrix}\!\right)_{i\in \bb{\mathbf{s}}}.
%      $$
%      Sometimes, to abbreviate, we will write $\xi_{\mathbf{r},\mathbf{s},s}$ instead of
%      $\xi_{\mathbf{r},\mathbf{s},s}^{\mathbf{Tm}_{\mathrm{H}_{S}}(\Sigma)}$.
% \end{enumerate}
%
% Then $\mathbf{Tm}_{\mathrm{H}_{S}}(\Sigma)$ is a Hall algebra, the \emph{Hall algebra for} $(S,\Sigma)$.
% \end{proposition}
%
% \begin{remark}
% For every $\Sigma$-algebra $\mathbf{A}$ there exists a homomorphism of Hall algebras $\mathrm{d}^{\mathbf{A}} = (\mathrm{d}^{\mathbf{A}}_{\mathbf{\mathbf{s}},s})_{(\mathbf{\mathbf{s}},s)\in S^{\star}\times S}$ from the Hall algebra $\mathbf{Tm}_{\mathrm{H}_{S}}(\Sigma)$ to the Hall algebra $\mathbf{O}_{\mathrm{H}_{S}}(\mathbf{A})$ and its image is $\mathbf{D}_{\mathrm{H}_{S}}(\mathbf{A})$, the Hall subalgebra of $\mathbf{O}_{\mathrm{H}_{S}}(\mathbf{A})$ of the finitary derived operations on $\mathbf{A}$, which, for every $(\mathbf{\mathbf{s}},s)\in S^{\star}\times S$, is such that 
% $\mathrm{D}_{\mathrm{H}_{S}}(\mathbf{A})_{\mathbf{\mathbf{s}},s} = 
% \mathrm{D}_{\mathbf{\mathbf{s}}}(\mathbf{A})_{s}$, where $\mathrm{D}_{\mathbf{\mathbf{s}}}(\mathbf{A})_{s}$ is the $s$-th coordinate of the underlying $S$-sorted of 
% $\mathbf{D}_{\mathbf{\mathbf{s}}}(\mathbf{A})$, the $\Sigma$-algebra of the $\mathbf{\mathbf{s}}$-\emph{ary} derived operations on $\mathbf{A}$, introduced in Definition~\ref{wdop}.
% \end{remark}
%
% %\begin{proposition}
% %Let $\Sigma$ be an $S$-sorted signature. Then $\mathbf{Tm}_{\mathrm{H}_{S}}(\Sigma)^{\wp}$ is a Hall algebra.
% %\end{proposition}
%
% %recall that the underlying $S^{\star}\times S$-sorted of $\mathbf{Tm}_{\mathrm{H}_{S}}(\Sigma)^{\wp}$ is $(\mathrm{Sub}(\mathrm{T}_{\Sigma}(\vs{w})_{s}))_{(w,s)\in S^{\star}\times S}$
%
% Our next goal is to prove that, for every $S^{\star}\times S$-sorted set $\Sigma$, $\mathbf{T}_{\mathrm{H}_{S}}(\Sigma)$, the free Hall algebra on $\Sigma$, is isomorphic to $\mathbf{Tm}_{\mathrm{H}_{S}}(\Sigma)$. We remark that the existence of this isomorphism is interesting because it enables us to get a more tractable description of the terms in $\mathbf{T}_{\mathrm{H}_{S}}(\Sigma)$.
% %We remark that the existence of this isomorphism is interesting because it enables us, on the one hand, to get a more tractable description of the terms in $\mathbf{T}_{\mathrm{H}_{S}}(\Sigma)$, and, on the other hand, as we will show afterwards, to state, for every $\Sigma$-algebra $\mathbf{A}$, taking into account the adjunction $\mathbf{T}_{\mathrm{H}_{S}}\dashv \mathrm{G}_{\mathrm{H}_{S}}$, the existence of a homomorphism of Hall algebras
% %$\Tr^{\mathbf{A}}$ from $\mathbf{Tm}_{\mathrm{H}_{S}}(\Sigma)$ to $\mathbf{O}_{\mathrm{H}_{S}}(\mathbf{A}) = \mathbf{O}_{\mathrm{H}_{S}}(A)$ such that $\mathrm{Th}_{\Sigma}(\mathbf{A})$, the finitary $\Sigma$-equational theory determined by $\mathbf{A}$, is
% %precisely $\mathrm{Ker}(\Tr^{\mathbf{A}})$, the kernel of the homomorphism $\Tr^{\mathbf{A}}$.
%
% To attain the goal just stated we begin by defining, for a Hall algebra $\mathbf{A}$, an $S$-sorted signature $\Sigma$, an
% $S^{\star}\times S$-mapping $f\colon\Sigma\mor A$, and a word $\mathbf{r}\in S^{\star}$, the concept of derived $\Sigma$-algebra of
% $\mathbf{A}$ for $(f,\mathbf{r})$, since it will be used afterwards in the proof of the isomorphism between
% $\mathbf{T}_{\mathrm{H}_{S}}(\Sigma)$ and $\mathbf{Tm}_{\mathrm{H}_{S}}(\Sigma)$.
%
%
%
% \begin{definition}
% Let $\mathbf{A}$ be a Hall algebra and $\Sigma$ an $S$-sorted signature. Then, for every $f\colon\Sigma\mor A$ and $\mathbf{r}\in
% S^{\star}$, $\mathbf{A}^{f,\mathbf{r}}$, the \emph{derived} $\Sigma$-\emph{algebra} \emph{of} $\mathbf{A}$ \emph{for} $(f,\mathbf{r})$,
% is the $\Sigma$-algebra with underlying $S$-sorted set $A^{f,\mathbf{r}} = (A_{\mathbf{r},s})_{s\in S}$ and algebraic structure $F^{f,\mathbf{r}}$, defined, for every $(\mathbf{s},s)\in S^{\star}\times S$, as
% $$
% F^{f,\mathbf{r}}_{\mathbf{s},s} \nfunction
% {\Sigma_{\mathbf{s},s}}{\mathrm{O}_{\mathbf{s}}(A^{f,\mathbf{r}})_{s}} {\sigma}{\nfunction
%     {\prod_{i\in \bb{\mathbf{s}}}A_{\mathbf{u},s_{i}}}
%     {A_{\mathbf{u},s}}
%     {(a_{0},\ldots,a_{\bb{\mathbf{s}}-1})}
%     {\xi_{\mathbf{r},\mathbf{s},s}^{\mathbf{A}}(f_{(\mathbf{s},s)}(\sigma),a_{0},\ldots,a_{\bb{\mathbf{s}}-1})}
% }
% $$
% where $\mathrm{O}_{\mathbf{s}}(A^{f,\mathbf{r}})_{s} = A_{\mathbf{r},s}^{\prod_{i\in \bb{\mathbf{s}}}A_{\mathbf{r},s_{i}}}$.
%
% Furthermore, we will denote by $p^{\mathbf{r}}$ the $S$-sorted mapping from $\vs{\mathbf{r}}$ to $A^{f,\mathbf{r}}$ defined, for every $s\in S$ and $i\in\bb{\mathbf{r}}$, as $p^{\mathbf{r}}_{s}(v^{s}_{i}) = (\pi^{\mathbf{r}}_{i})^{\mathbf{A}}$, and by $(p^{\mathbf{r}})^{\sharp}$ the unique homomorphism from
% $\mathbf{T}_{\Sigma}(\vs{\mathbf{r}})$ to $\mathbf{A}^{f,\mathbf{r}}$ such that $(p^{\mathbf{r}})^{\sharp}\circ \eta^{\vs{\mathbf{r}}} = p^{\mathbf{r}}$.
% \end{definition}
%
% \begin{remark}
% For a $\Sigma$-algebra $\mathbf{B}=(B,G)$, we have that $G\colon\Sigma\mor\mathrm{O}_{\mathrm{H}_{S}}(B)$ and
% $\mathbf{B}\iso\mathbf{O}_{\mathrm{H}_{S}}(B)^{G,\lambda}$, where $\lambda$ is the empty word on $S$.  Besides, for every $\mathbf{r}\in S^{\star}$, we have that $\mathbf{B}^{B_{\mathbf{r}}}$, the direct $B_{\mathbf{r}}$-power of $\mathbf{B}$, is isomorphic to $\mathbf{O}_{\mathrm{H}_{S}}(B)^{G,\mathbf{r}}$.
% \end{remark}
%
%
%
% \begin{lemma}\label{L:aux-app}
% Let $\Sigma$ be an $S$-sorted signature, $\mathbf{A}$ a Hall algebra, $f\colon\Sigma\mor A$ and $\mathbf{q}\in S^{\star}$.  Then, for
% every $(\mathbf{r},s)\in S^{\star}\times S$, $P\in \mathrm{T}_{\Sigma}(\vs{\mathbf{r}})_{s}$ and $a\in \prod_{i\in \bb{\mathbf{r}}}A_{\mathbf{q},r_{i}}$, we have that
% $$
% P^{\mathbf{A}^{f,\mathbf{q}}}(a_{0},\ldots,a_{\bb{\mathbf{r}}-1}) =
% \xi_{\mathbf{q},\mathbf{r},s}^{\mathbf{A}}((p^{\mathbf{r}})^{\sharp}_{s}(P),a_{0},\ldots,a_{\bb{\mathbf{r}}-1}).
% $$
% \end{lemma}
%
% \begin{proof}
% By algebraic induction on the complexity of $P$.  If $P$ is a variable $v_{i}^{s}$, with $i\in\bb{\mathbf{r}}$, then
% \begin{align*}
%     v_{i}^{s,\mathbf{A}^{f,\mathbf{q}}}(a_{0},\ldots,a_{\bb{\mathbf{r}}-1})\ &=
%     a^{\sharp}_{r_{i}}(v^{s}_{i})\\
%     &=
%     a_{i} \\
%     &=
%     \xi_{\mathbf{q},\mathbf{r},s}^{\mathbf{A}}((\pi^{\mathbf{r}}_{i})^{\mathbf{A}},a_{0},\ldots,a_{\bb{\mathbf{r}}-1} )
%     \tag{by $\mathrm{H}_{1}$} \\
%     &=
%     \xi_{\mathbf{q},\mathbf{r},s}^{\mathbf{A}}((p^{w})^{\sharp}_{s}(v^{s}_{i}),a_{0},\ldots,a_{\bb{\mathbf{r}}-1} ).
% \end{align*}
% Let us assume that $P = \sigma(Q_{0},\ldots,Q_{\bb{x}-1})$, with $\sigma\in \Sigma_{\mathbf{s},s}$ and that, for every $j\in\bb{x}$, $Q_{j}\in \mathrm{T}_{\Sigma}(\vs{\mathbf{r}})_{s_{j}}$ fulfils the induction hypothesis. Then we have that
% \begin{flushleft}
% $(\sigma(Q_{0},\ldots,Q_{\bb{\mathbf{s}}-1}))^{\mathbf{A}^{f,\mathbf{q}}}
%       (a_{0},\ldots,a_{\bb{\mathbf{r}}-1})$
% \allowdisplaybreaks
% \begin{align*}
%     &=
%     \sigma^{\mathbf{A}^{f,\mathbf{q}}}
%        (
%        Q_{0}^{\mathbf{A}^{f,\mathbf{q}}}(a_{0},\ldots,a_{\bb{\mathbf{r}}-1}),
%        \ldots,
%        Q_{\bb{\mathbf{s}}-1}^{\mathbf{A}^{f,\mathbf{q}}}(a_{0},\ldots,a_{\bb{\mathbf{r}}-1})
%        )\\
%     &=
%     \xi_{\mathbf{q},\mathbf{s},s}^{\mathbf{A}}
%        (
%        f(\sigma),
%        Q_{0}^{\mathbf{A}^{f,\mathbf{q}}}(a_{0},\ldots,a_{\bb{\mathbf{r}}-1}),
%        \ldots,
%        Q_{\bb{\mathbf{s}}-1}^{\mathbf{A}^{f,\mathbf{q}}}(a_{0},\ldots,a_{\bb{\mathbf{r}}-1})
%        )\\
%    &=
%     \xi_{\mathbf{q},\mathbf{s},s}^{\mathbf{A}}
%        (
%        f(\sigma),
%       \xi_{\mathbf{q},\mathbf{r},s_{0}}^{\mathbf{A}}
%             ((p^{\mathbf{r}})^{\sharp}_{s_{0}}(Q_{0}),a_{0},\ldots,a_{\bb{\mathbf{r}}-1}),
%        \ldots, 
% \\&\qquad\qquad\qquad\qquad\qquad\qquad
% \xi_{\mathbf{q},\mathbf{r},s_{\bb{\mathbf{s}}-1}}^{\mathbf{A}}
%             ((p^{\mathbf{r}})^{\sharp}_{s_{\bb{\mathbf{s}}-1}}(Q_{\bb{\mathbf{s}}-1}),a_{0},\ldots,a_{\bb{\mathbf{r}}-1})
%        ) \tag{by ind.}
% \\&=
%     \xi_{\mathbf{q},\mathbf{r},s}^{\mathbf{A}}
%        (
%          \xi_{\mathbf{r},\mathbf{s},s}^{\mathbf{A}}
%            (f(\sigma),
%             (p^{\mathbf{r}})^{\sharp}_{s_{0}}(Q_{0}),
%             \ldots,
%             (p^{\mathbf{r}})^{\sharp}_{s_{\bb{\mathbf{s}}-1}}(Q_{\bb{\mathbf{s}}-1})
%            ),
%          a_{0},
%          \ldots,
%          a_{\bb{\mathbf{r}}-1}
%        ) \tag{by $\mathrm{H}_{3}$}   \\
%     &=
%     \xi_{\mathbf{q},\mathbf{r},s}^{\mathbf{A}}
%        (
%        \sigma^{\mathbf{A}_{\mathbf{r}}}
%           (
%           (p^{\mathbf{r}})^{\sharp}_{s_{0}}(Q_{0}),
%           \ldots,
%           (p^{\mathbf{r}})^{\sharp}_{s_{\bb{\mathbf{s}}-1}}(Q_{\bb{\mathbf{s}}-1})
%           ),
%        a_{0},
%        \ldots,
%        a_{\bb{\mathbf{r}}-1}
%        ) \\
%     &=
%     \xi_{\mathbf{q},\mathbf{r},s}^{\mathbf{A}}
%        (
%        (p^{\mathbf{r}})^{\sharp}_{s}(\sigma,Q_{0},\ldots,Q_{\bb{\mathbf{s}}-1}),
%        a_{0},
%        \ldots,
%        a_{\bb{\mathbf{r}}-1}
%        ) \\
%     &=
%     \xi_{\mathbf{q},\mathbf{r},s}^{\mathbf{A}}
%        (
%        (p^{\mathbf{r}})^{\sharp}_{s}(P),
%        a_{0},
%        \ldots,
%        a_{\bb{\mathbf{r}}-1}
%        ).
% \end{align*}
% \end{flushleft}
% This completes the proof.
% \end{proof}
%
% Next we prove that, for every $S^{\star}\times S$-sorted set $\Sigma$, the Hall algebra for $(S,\Sigma)$ is isomorphic to the free Hall algebra on $\Sigma$.
%
%
% \begin{proposition}\label{iso:FrH-TerH-app}
% Let $\Sigma$ be an $S$-sorted signature, i.e., an $S^{\star}\times S$-sorted set.  Then the Hall algebra $\mathbf{Tm}_{\mathrm{H}_{S}}(\Sigma)$ is isomorphic to $\mathbf{T}_{\mathrm{H}_{S}}(\Sigma)$, the free Hall algebra on $\Sigma$.
% \end{proposition}
%
% \begin{proof}
% It is enough to prove that $\mathbf{Tm}_{\mathrm{H}_{S}}(\Sigma)$ has the universal property of $\mathbf{T}_{\mathrm{H}_{S}}(\Sigma)$. Therefore we have to specify an $S^{\star}\times S$-sorted mapping $\eta^{\Sigma}$ from $\Sigma$ to $\mathrm{Tm}_{\mathrm{H}_{S}}(\Sigma)$ such that, for every Hall algebra $\mathbf{A}$ and $S^{\star}\times S$-sorted mapping $f$ from $\Sigma$ to $A$, there is a unique homomorphism $\widehat{f}$ from $\mathbf{Tm}_{\mathrm{H}_{S}}(\Sigma)$ to $\mathbf{A}$ such that
% $\widehat{f}\circ \eta^{\Sigma} = f$.  Let $\eta^{\Sigma}$ be the $S^{\star}\times S$-sorted mapping defined, for every
% $(\mathbf{s},s)\in S^{\star}\times S$, as%
% $$
% \eta^{\Sigma}_{\mathbf{s},s}\nfunction
% {\Sigma_{\mathbf{s},s}}
% {\mathrm{T}_{\Sigma}(\vs{\mathbf{s}})_{s}}
% {\sigma}
% {\sigma\left(v^{s_{0}}_{0},\ldots,v^{s_{\bb{\mathbf{s}}-1}}_{\bb{\mathbf{s}}-1}\right)}
% $$
% Let $\mathbf{A}$ be a Hall algebra, $f\colon\Sigma\mor A$ an $S^{\star}\times S$-sorted mapping and $\widehat{f}$ the $S^{\star}\times S$-sorted mapping from $\mathrm{Tm}_{\mathrm{H}_{S}}(\Sigma)$ to $A$ defined, for every $(\mathbf{s},s)\in S^{\star}\times S$, as $\widehat{f}_{\mathbf{s},s} = (p^{\mathbf{s}})^{\sharp}_{s}$, where, we recall, $(p^{\mathbf{s}})^{\sharp}$ is the unique homomorphism from $\mathbf{T}_{\Sigma}(\vs{\mathbf{s}})$ to $\mathbf{A}^{f,\mathbf{s}}$ such that $(p^{\mathbf{s}})^{\sharp}\circ \eta^{\vs{\mathbf{s}}} = p^{\mathbf{s}}$. Then $\widehat{f}$ is a homomorphism of Hall algebras, because, on the one hand, for $\mathbf{s}\in S^{\star}$ and $i\in\bb{\mathbf{s}}$ we have that%
% \begin{align*}
%     \widehat{f}_{\mathbf{s},s_{i}}((\pi^{\mathbf{s}}_{i})^{\mathbf{Tm}_{\mathrm{H}_{S}}(\Sigma)})
%     &=
%     \widehat{f}_{\mathbf{s},s_{i}}(v^{s_{i}}_{i}) \\
%     &=
%     p^{\mathbf{s}}_{s_{i}}(v^{s_{i}}_{i}) \\
%     &=
%     (\pi^{\mathbf{s}}_{i})^{\mathbf{A}},
% \end{align*}
% and, on the other hand, for $P\in \mathrm{T}_{\Sigma}(\vs{\mathbf{s}})_{s}$ and
% $\mathcal{Q} = (Q_{i})_{i\in \bb{\mathbf{s}}}\in \mathrm{T}_{\Sigma}(\vs{\mathbf{r}})_{\mathbf{s}}$ we have that %
% \begin{flushleft}
% $\widehat{f}_{\mathbf{r},s}(
%       \xi_{\mathbf{r},\mathbf{s},s}^{\mathbf{Tm}_{\mathrm{H}_{S}}(\Sigma)}
%         (P,Q_{0},\ldots,Q_{\bb{\mathbf{s}}-1})
%                    )$
% \begin{align*}
% \quad    &=
%     (p^{\mathbf{r}})^{\sharp}_{s}(
%       \mathcal{Q}^{\sharp}_{s}(P)
%                      ) \\
%     &=
%     ((p^{\mathbf{r}})^{\sharp}\circ \mathcal{Q})^{\sharp}_{s} (P)
%     \tag{$(p^{\mathbf{r}})^{\sharp}\circ\mathcal{Q}^{\sharp} =
%     ((p^{\mathbf{r}})^{\sharp}\circ \mathcal{Q})^{\sharp}$}
%     \\
%     &=
%     P^{\mathbf{A}^{f,\mathbf{r}}}
%       ((p^{\mathbf{r}})^{\sharp}_{s_{0}}(Q_{0}),\ldots,
%         (p^{\mathbf{r}})^{\sharp}_{s_{\bb{\mathbf{s}}-1}}(Q_{\bb{\mathbf{s}}-1})
%       ) \\
%     &=
%     \xi_{\mathbf{r},\mathbf{s},s}^{\mathbf{A}}
%       ((p^{\mathbf{r}})^{\sharp}_{s}(P),
%        (p^{\mathbf{r}})^{\sharp}_{s_{0}}(Q_{0}),\ldots,
%        (p^{\mathbf{r}})^{\sharp}_{s_{\bb{\mathbf{s}}-1}}(Q_{\bb{\mathbf{s}}-1})
%       )
%     \tag{by Lemma~\ref{L:aux-app}}
%       \\
%     &=
%     \xi_{\mathbf{r},\mathbf{s},s}^{\mathbf{A}}
%       ( \widehat{f}_{\mathbf{s},s}(P),
%         \widehat{f}_{\mathbf{r},s_{0}}(Q_{0}),\ldots,
%         \widehat{f}_{\mathbf{r},s_{\bb{\mathbf{s}}-1}}(Q_{\bb{\mathbf{s}}-1})
%       ).
% \end{align*}
% \end{flushleft}
% Therefore the $S^{\star}\times S$-sorted mapping $\widehat{f}$ is a homomorphism. Furthermore, $\widehat{f}\circ \eta^{\Sigma} = f$, because, for every $\mathbf{s}\in S^{\star}$, $s\in S$, and $\sigma\in \Sigma_{\mathbf{s},s}$, we have that%
% \begin{align*}
% \widehat{f}_{\mathbf{s},s}(\eta^{\Sigma}_{\mathbf{s},s}(\sigma)) &=
%     (p^{\mathbf{s}})^{\sharp}_{s}(\sigma(v^{s_{0}}_{0},\ldots,v^{s_{\bb{\mathbf{s}}-1}}_{\bb{\mathbf{s}}-1}))
%     \\&=
%     \sigma^{\mathbf{A}_{\mathbf{s}}}
%       (p^{\mathbf{s}}_{{s}_{0}}(v^{s_{0}}_{0}),
%         \ldots,
%        p^{\mathbf{s}}_{{s}_{\bb{\mathbf{s}}-1}}(v^{s_{\bb{\mathbf{s}}-1}}_{\bb{\mathbf{s}}-1})
%       ) \\
%     &=
%     \xi_{\mathbf{s},\mathbf{s},s}^{\mathbf{A}}
%       ( f_{(\mathbf{s},s)}(\sigma),
%         (\pi^{\mathbf{s}}_{0})^{\mathbf{A}},
%         \ldots,
%         (\pi^{\mathbf{s}}_{\bb{\mathbf{s}}-1})^{\mathbf{A}}
%       ) \\
%     &=
%     f_{\mathbf{s},s}(\sigma) \tag{by $\mathrm{H}_{2}$}.
% \end{align*}
% It is obvious that $\widehat{f}$ is the unique homomorphism such that $\widehat{f}\circ \eta^{\Sigma} = f$.  Henceforth $\mathbf{Tm}_{\mathrm{H}_{S}}(\Sigma)$ is isomorphic to $\mathbf{T}_{\mathrm{H}_{S}}(\Sigma)$.
% \end{proof}
%
%
% This isomorphism, together with the adjunction $\mathbf{T}_{\mathrm{H}_{S}}\dashv\mathrm{G}_{\mathrm{H}_{S}}$,  has as an immediate consequence that, for every $S$-sorted set $A$ and every $S$-sorted signature $\Sigma$, the sets $\mathrm{Hom}({\Sigma},\mathrm{G}_{\mathrm{H}_{S}}(\mathbf{O}_{\mathrm{H}_{S}}(A)))$, in 
% $\mathsf{Set}^{S^{\star}\times S}$, and $\mathrm{Hom}(\mathbf{Tm}_{\mathrm{H}_{S}}(\Sigma),\mathbf{O}_{\mathrm{H}_{S}}(A))$, in $\mathsf{Alg}(\mathrm{H}_{S})$, are naturally isomorphic.
%
% Actually, for an $S$-sorted set $A$ and an $S$-sorted signature $\Sigma$,
% \begin{enumerate}
% \item the mapping that assigns to a structure of $\Sigma$-algebra $F$ on $A$ (i.e., an $S^{\star}\times S$-sorted mapping $F$ from $\Sigma$ to $\mathrm{O}_{\mathrm{H}_{S}}(A)$) the unique  homomorphism of Hall algebras $\widehat{F} = (\widehat{F}_{w,s})_{(w,s)\in S^{\star}\times S}$ from $\mathbf{Tm}_{\mathrm{H}_{S}}(\Sigma)$ to $\mathbf{O}_{\mathrm{H}_{S}}(A)$ such that $\widehat{F}\circ \eta^{\Sigma} = F$, together with
% \item the mapping that assigns to a homomorphism $h$ from $\mathbf{Tm}_{\mathrm{H}_{S}}(\Sigma)$ to $\mathbf{O}_{\mathrm{H}_{S}}(A)$, essentially, the algebraic structure $\mathrm{G}_{\mathrm{H}_{S}}(h)\circ \eta^{\Sigma}$ on $A$, where $\eta^{\Sigma}$ is the canonical embedding of $\Sigma$ into $\mathrm{T}_{\mathrm{H}_{S}}(\Sigma)$, 
% \end{enumerate}
% which are natural in both variables and mutually inverse bijections, provide the mentioned natural isomorphism.
%
% \begin{remark}
% For an $S$-sorted set $A$, a structure of $\Sigma$-algebra $F$ on $A$ and a word $\mathbf{s}\in S^{\star}$, the subfamily $(\widehat{F}_{\mathbf{s},s})_{s\in S}$ of $\widehat{F} = (\widehat{F}_{\mathbf{s},s})_{(\mathbf{s},s)\in S^{\star}\times S}$,
% denoted by $\mathrm{d}^{(A,F)} = (\mathrm{d}^{(A,F)}_{\mathbf{s},s})_{s\in S}$, is the unique homomorphism from $\mathbf{T}_{\Sigma}(\vs{\mathbf{s}})$ to $(A,F)^{A_{\mathbf{s}}}$, the direct $A_{\mathbf{s}}$-power of $(A,F)$, such that $\mathrm{d}^{(A,F)}_{\mathbf{s}}\circ \eta^{\vs{\mathbf{s}}} = \mathrm{p}^{A}_{\mathbf{s}}$, where $\mathrm{p}^{A}_{\mathbf{s}}$ is the $S$-sorted mapping from $\vs{\mathbf{s}}$ to $A^{A_{\mathbf{s}}}$ defined, for every $s\in S$ and every $v^{s}_{i}\in(\vs{\mathbf{s}})_{s}$, as $\mathrm{p}^{A}_{\mathbf{s},s}(v^{s}_{i}) = \mathrm{pr}^{A}_{\mathbf{s},i}$.
% \end{remark}
%
% For a set of sorts $S$ and an $S$-sorted signature $\Sigma$, the derived operation symbols of $\Sigma$, i.e., the elements of the components $\mathrm{T}_{\Sigma}(\vs{\mathbf{s}})_{s}$ of $\mathrm{Tm}_{\mathrm{H}_{S}}(\Sigma)$ can be considered as the operation symbols of a new signature. The mappings $\varphi$ from a set of sorts $S$ to another $T$, together with the $S^{\star}\times S$-sorted mappings $d$ that assign to operation symbols of an $S$-sorted signature $\Sigma$ derived operations symbols of a $T$-sorted signature $\Lambda$, i.e., terms of the components $\mathrm{T}_{\Lambda}(\vs{\varphi^{\star}(\mathbf{s})})_{\varphi(s)}$ of $\mathrm{Tm}_{\mathrm{H}_{T}}(\Lambda)_{\varphi^{\star}\times\varphi}$, form a new class of morphisms between many-sorted signatures which we will call \emph{derivors}.
%
% \begin{definition}\label{DDerivor-app}
% Let $\sig{\Sigma} = (S,\Sigma)$ and $\sig{\Lambda} = (T,\Lambda)$ be many-sorted signatures. A \emph{derivor} from $\sig{\Sigma}$ to $\sig{\Lambda}$ is an ordered pair $\mathbf{d} = (\varphi,d)$, where $\varphi\colon S\mor T$ and  $d\colon\Sigma\mor\mathrm{Tm}_{\mathrm{H}_{T}}(\Lambda)_{\varphi^{\star}\times\varphi}$. Thus, for every $(\mathbf{s},s)\in S^{\star}\times S$, %
% $$
% d_{\mathbf{s},s}\colon \Sigma_{\mathbf{s},s}\mor\mathrm{Tm}_{\mathrm{H}_{T}}(\Lambda)_{\varphi^{\star}(\mathbf{s}),\varphi(s)}=
%     \mathrm{T}_{\Lambda}(\vs{\varphi^{\star}(\mathbf{s})})_{\varphi(s)}.
% $$
% Hence, for every $\sigma\in \Sigma_{\mathbf{s},s}$, $d_{\mathbf{s},s}(\sigma)\in \mathrm{T}_{\Lambda}(\vs{\varphi^{\star}(\mathbf{s})})_{\varphi(s)}$, i.e., $d_{\mathbf{s},s}(\sigma)$ is a term of sort $\varphi(s)$ with variables in $\vs{\varphi^{\star}(\mathbf{s})}$.
% \end{definition}
%
%
%
% \begin{remark}\label{Rdop-app}
% Let $\varphi$ be a mapping from the set of sorts $S$ to the set of sorts $T$ and let $V^{T}$ be a fixed $T$-sorted set with a countably infinity of variables in each component, say $V^{T}_{t} = \{v^{t}_{n}\mid n\in \mathbb{N}\}$ (which is isomorphic to $\{t\}\times\mathbb{N}$). Then, for every $\mathbf{s}\in S^{\star}$, we let $\vs{\varphi^{\star}(\mathbf{s})} = ((\vs{\varphi^{\star}(\mathbf{s})})_{t})_{t\in T}$ stand for the $T$-sorted set  defined, for every $t\in T$, as:
% $$
% (\vs{\varphi^{\star}(\mathbf{s})})_{t} = \{v^{t}_{0},\ldots,v^{t}_{\bb{\varphi^{\star}(\mathbf{s})}_{t}-1}\} (\cong \varphi^{\star}(\mathbf{s})^{-1}[\{t\}] = \{i\in \bb{\varphi^{\star}(\mathbf{s})}\mid \varphi(w_{i}) = t\}).
% $$
% Thus, for every $t\in T$, $(\vs{\varphi^{\star}(\mathbf{s})})_{t}$ is an initial segment of $V^{T}_{t}$ and 
% $\vs{\varphi^{\star}(\mathbf{s})}$ is a finite subset of $V^{T}$. The latter is fulfilled because $\mathrm{supp}_{T}(\vs{\varphi^{\star}(\mathbf{s})}) = \mathrm{Im}(\varphi^{\star}(\mathbf{s}))$ and $\bb{\varphi^{\star}(\mathbf{s})} = \bb{\mathbf{s}}$ is finite. Moreover, $\mathrm{supp}_{T}(\vs{\varphi^{\star}(\mathbf{s})}) =\varphi[\mathrm{supp}_{S}(\vs{\mathbf{s}})]$. Recall  that $\varphi^{\star}(\mathbf{s}) = (\varphi(s_{i}))_{i\in \bb{\mathbf{s}}}$ and that, since $\varphi^{\star}(\lambda) = \lambda$, $\vs{\varphi^{\star}(\lambda)} = \varnothing^{T} = (\varnothing)_{t\in T}$.
%
% If $\mathbf{d} = (\varphi,d)$ is a derivor from $\sig{\Sigma}$ to $\sig{\Lambda}$, $\mathbf{B}$ a 
% $\sig{\Lambda}$-algebra, $(\mathbf{s},s)\in S^{\star}\times S$ and $\sigma\in \Sigma_{\mathbf{s},s}$, then the derived operation (or term operation or polynomial operation) of type $(\varphi^{\star}(\mathbf{s}),\varphi(s))$ determined by the term $d_{\mathbf{s},s}(\sigma)$ in $\mathbf{B}$, denoted by $d_{\mathbf{s},s}(\sigma)^{\mathbf{B}}$ or by $d(\sigma)^{\mathbf{B}}$ for short, is a mapping from $B_{\varphi^{\star}(\mathbf{s})}$ to $B_{\varphi(s)}$. In particular, for a $T$-sorted set $Y$, we have that the derived operation (or term operation or polynomial operation) of type $(\varphi^{\star}(\mathbf{s}),\varphi(s))$ determined by the term $d_{\mathbf{s},s}(\sigma)$ in $\mathbf{T}_{\Lambda}(Y)$, denoted by $d_{\mathbf{s},s}(\sigma)^{\mathbf{T}_{\Lambda}(Y)}$ or by  $d(\sigma)^{\mathbf{T}_{\Lambda}(Y)}$ for short, is a mapping from $\mathrm{T}_{\Lambda}(Y)_{\varphi^{\star}(\mathbf{s})}$ to $\mathrm{T}_{\Lambda}(Y)_{\varphi(s)}$. Notice that, for every $s\in S$ and every $\sigma\in \Sigma_{\lambda,s}$, $d_{\lambda,s}(\sigma)\in \mathrm{T}_{\Lambda}(\varnothing^{T})_{\varphi(s)}$, i.e., $d_{\lambda,s}(\sigma)$ is a closed term of sort $\varphi(s)$.  Therefore, for every $\sig{\Lambda}$-algebra $\mathbf{B}$, $d_{\lambda,s}(\sigma)^{\mathbf{B}}$ will be a mapping from $B_{\lambda}$ to $B_{\varphi(s)}$, i.e., essentially, a distinguished element of $B_{\varphi(s)}$. In particular, for a $T$-sorted set $Y$, $d_{\lambda,s}(\sigma)^{\mathbf{T}_{\Lambda}(Y)}$ will be a mapping from $\mathrm{T}_{\Lambda}(Y)_{\lambda}$ to $\mathrm{T}_{\Lambda}(Y)_{\varphi(s)}$, i.e., essentially, a distinguished element of $\mathrm{T}_{\Lambda}(Y)_{\varphi(s)}$. Let us note that $d(\sigma)^{\mathbf{B}}$ is the value of $\mathrm{d}^{\mathbf{B}}_{\varphi^{\star}(\mathbf{s}),s}$ at $d_{\mathbf{s},s}(\sigma)$ and that $d(\sigma)^{\mathbf{T}_{\Lambda}(Y)}$ is the value of $\mathrm{d}^{\mathbf{T}_{\Lambda}(Y)}_{\varphi^{\star}(\mathbf{s}),s}$ at $d_{\mathbf{s},s}(\sigma)$, where 
% $\mathrm{d}^{\mathbf{B}}_{\varphi^{\star}(\mathbf{s})}$ is the unique homomorphism from $\mathbf{T}_{\Lambda}(\vs{\varphi^{\star}(\mathbf{s})})$ to $\mathbf{B}^{B_{\varphi^{\star}(\mathbf{s})}}$ for which the left diagram below commutes, and $\mathrm{d}^{\mathbf{T}_{\Lambda}(Y)}_{\varphi^{\star}(\mathbf{s})}$ the unique homomorphism from $\mathbf{T}_{\Lambda}(\vs{\varphi^{\star}(\mathbf{s})})$ to $\mathbf{T}_{\Lambda}(Y)^{\mathrm{T}_{\Lambda}(Y)_{\varphi^{\star}(\mathbf{s})}}$ for which the right diagram below commutes.
% $$
% \begin{aligned}
% \xymatrix@C=40pt@R=33pt{
% \vs{\varphi^{\star}(\mathbf{s})}
%   \ar[r]^-{\eta^{\vs{\varphi^{\star}(\mathbf{s})}}}
%   \ar[rd]_-{\mathrm{p}^{B}_{\varphi^{\star}(\mathbf{s})}} &
% \mathrm{T}_{\Lambda}(\vs{\varphi^{\star}(\mathbf{s})}) \ar[d]^{\mathrm{d}^{\mathbf{B}}_{\varphi^{\star}(\mathbf{s})}}
%   \\
% & B^{B_{\varphi^{\star}(\mathbf{s})}}
% }
% %
% \end{aligned}
% \quad \text{} \qquad\;
% \begin{aligned}
% %
% \xymatrix@C=40pt@R=33pt{
% \vs{\varphi^{\star}(\mathbf{s})}
%   \ar[r]^-{\eta^{\vs{\varphi^{\star}(\mathbf{s})}}}
%   \ar[rd]_-{\mathrm{p}^{\mathrm{T}_{\Lambda}(Y)}_{\varphi^{\star}(\mathbf{s})}} &
% \mathrm{T}_{\Lambda}(\vs{\varphi^{\star}(\mathbf{s})}) \ar[d]^{\mathrm{d}^{\mathbf{T}_{\Lambda}(Y)}_{\varphi^{\star}(\mathbf{s})}}
%   \\
% & \mathrm{T}_{\Lambda}(Y)^{\mathrm{T}_{\Lambda}(Y)_{\varphi^{\star}(\mathbf{s})}}
% }
% %
% \end{aligned}
% $$
% Recall that, for every $t\in T$ and every $i\in (\vs{\varphi^{\star}(\mathbf{s})})_{t}$ (thus $(\varphi^{\star}(\mathbf{s}))(i) = \varphi(s_{i}) = t$), $\mathrm{p}^{B}_{\varphi^{\star}(\mathbf{s}),t}(i)$ is $\mathrm{pr}^{B}_{\varphi^{\star}(\mathbf{s}),i}$, i.e., the $\varphi^{\star}(\mathbf{s})$-ary, $i$-th canonical projection from $B_{\varphi^{\star}(\mathbf{s})}$ to $B_{t} = B_{(\varphi^{\star}(\mathbf{s}))(i)}$.
% \end{remark}
%
% \begin{remark}
% While derivors are not the most general type of morphism that might be considered between many-sorted signatures---for instance, one could consider polyderivors, see~\cite{CVST10b}---, they are an important class of such morphisms. One reason for its relevance is its formal properties (see below), another that there are many mathematical examples of them which are of interest (see~\cite{CVST10b}).
% \end{remark}
%
% Since by, Proposition~\ref{iso:FrH-TerH-app}, $\mathbf{Tm}_{\mathrm{H}_{T}}(\Lambda)$ is isomorphic to 
% $\mathbf{T}_{\mathrm{H}_{T}}(\Lambda)$, the derivors can be defined, alternative, but equivalently, as ordered pairs  $\mathbf{d} = (\varphi,d)$ with $d\colon\Sigma\mor \mathrm{T}_{\mathrm{H}_{T}}(\Lambda)$.
%
% On the other hand, every mapping $\varphi\colon S\mor T$ determines a functor from the category  $\mathsf{Alg}(\mathrm{H}_{T})$ to the category $\mathsf{Alg}(\mathrm{H}_{S})$, so $\mathrm{Tm}_{\mathrm{H}_{T}}(\Lambda)_{\varphi^{\star}\times\varphi}$ is in its turn e\-quipped with a structure of Hall algebra for $S$, which will allow us, in particular, to define the composition of derivors. We next show the existence of such a functor by defining a morphism of algebraic presentations from  $(\Sigma^{\mathrm{H}_{S}},\mathcal{E}^{\mathrm{H}_{S}})$ to $(\Sigma^{\mathrm{H}_{T}},\mathcal{E}^{\mathrm{H}_{T}})$.
%
%
%
% \begin{proposition}
% Let $\varphi\colon S\mor T$ be a mapping between sets of sorts and $h^{\varphi}$ the $S^{\star}\times S$-sorted mapping from $\Sigma^{\mathrm{H}_{S}}$ to $\Sigma^{\mathrm{H}_{T}}_{\varphi^{\star}\times \varphi}$ defined as follows: 
% \begin{enumerate}
% \item For every $\mathbf{s}\in S^{\star}$ and every $i\in\bb{\mathbf{s}}$, $h^{\varphi}(\pi^{\mathbf{s}}_{i}) = \pi^{\varphi^{\star}(\mathbf{s})}_{i}$.
%
% \item  For every $\mathbf{r},\,\mathbf{s}\in S^{\star}$ and every $s\in S$,
%   $h^{\varphi}(\xi_{\mathbf{r},\mathbf{s},s}) = \xi_{\varphi^{\star}(\mathbf{r}),\varphi^{\star}(\mathbf{s}),\varphi(s)}$.
% \end{enumerate}
% Then $(\varphi^{\star}\times \varphi, h^{\varphi}) \colon(S^{\star}\times S,\Sigma^{\mathrm{H}_{S}},\mathcal{E}^{\mathrm{H}_{S}})\mor
% (T^{\star}\times T,\Sigma^{\mathrm{H}_{T}},\mathcal{E}^{\mathrm{H}_{T}})$ is a morphism of algebraic presentations.
% \end{proposition}
%
% The morphisms of algebraic presentations determine functors in the opposite direction between the associated categories of algebras. Therefore, each mapping $\varphi\colon S\mor T$ between sets of sorts, determines a functor $(\varphi^{\star}\times \varphi, h^{\varphi})^{\ast}$ from $\mathsf{Alg}(\mathrm{H}_{T})$ to $\mathsf{Alg}(\mathrm{H}_{S})$, which transforms Hall algebras for $T$ into Hall algebras for $S$. The action of the functor on the free Hall algebra on a $T$-sorted signature $\Lambda$ is a Hall algebra for $S$, whose underlying $S^{\star}\times S$-sorted set is
% $\mathrm{Tm}_{\mathrm{H}_{T}}(\Lambda)_{\varphi^{\star}\times \varphi}$.
%
% If $\mathbf{d}\colon\sig{\Sigma}\mor\sig{\Lambda}$ is a derivor, then
% $d\colon \Sigma\mor \mathrm{Tm}_{\mathrm{H}_{T}}(\Lambda)_{\varphi^{\star}\times\varphi}$ determines a homomorphism of Hall algebras $d^{\sharp}\colon\mathbf{Tm}_{\mathrm{H}_{S}}(\Sigma)\mor
% \mathbf{Tm}_{\mathrm{H}_{T}}(\Lambda)_{\varphi^{\star}\times\varphi}$. Thus, for every $(\mathbf{s},s)\in S^{\star}\times S$, $d^{\sharp}_{\mathbf{s},s}$ sends terms in $\mathrm{T}_{\Sigma}(\vs{\mathbf{s}})_{s}$ to terms in $\mathrm{T}_{\Lambda}(\vs{\varphi^{\star}(\mathbf{s})})_{\varphi(s)}$.
%
% \begin{definition}\label{DDerCompId-app}
% Let $\mathbf{d}\colon\sig{\Sigma}\mor\sig{\Lambda}$ and $\mathbf{e}\colon\sig{\Lambda}\mor\sig{\Omega}$ be  derivors. Then $\mathbf{e}\circ \mathbf{d} = (\psi,e)\circ(\varphi,d)$, the \emph{composition} of $\mathbf{d}$ and  $\mathbf{e}$, is the derivor $(\psi\circ\varphi,e^{\sharp}_{\varphi^{\star}\times
% \varphi}\circ d)$, where $e^{\sharp}_{\varphi^{\star}\times\varphi}\circ d$ is obtained from
% $$
% \begin{aligned}
% \xymatrix@C=40pt@R=30pt{
% \Lambda
%   \ar[r]^-{\eta^{\Lambda}}
%   \ar[rd]_-{e} &
% \mathrm{Tm}_{\mathrm{H}_{T}}(\Lambda) \ar[d]^{e^{\sharp}}
%   \\
% & \mathrm{Tm}_{\mathrm{H}_{U}}(\Omega)_{\psi^{\star}\times \psi}
% }
% %
% \end{aligned}
% \qquad \text{as} \qquad\;\;
% \begin{aligned}
% %
% \xymatrix@C=40pt@R=30pt{
% \mathrm{Tm}_{\mathrm{H}_{T}}(\Lambda)_{\varphi^{\star}\times \varphi}
% \ar[d]^{e^{\sharp}_{\varphi^{\star}\times \varphi}} &
% \Sigma \ar[l]_-{d} \\
% {\mathrm{Tm}_{\mathrm{H}_{U}}(\Omega)_{\psi^{\star}\times\psi}}_{\varphi^{\star}\times\varphi}
% }
% \end{aligned}
% $$
% being $e^{\sharp}$ the canonical extension of $e$ to the free Hall algebra on $\Lambda$.
%
% For every many-sorted signature $\sig{\Sigma} = (S,\Sigma)$, the \emph{identity} at $(S,\Sigma)$
% is $(\id^{S},\eta^{\Sigma})$.
% \end{definition}
%
% The just stated definition allows us to form a category of many-sorted signatures whose morphisms are the derivors.
%
% \begin{proposition}\label{PSigdCat-app}
% The many-sorted signatures together with the derivors constitute a category, denoted by $\mathsf{Sig}_{\mathfrak{d}}$.
% \end{proposition}
%
%
%
% \begin{remark}
% Let $\mathsf{Sig}$ be the category whose objects are the many-sorted signatures and whose morphisms from $\sig{\Sigma}$ to $\sig{\Lambda}$ are the ordered pairs $(\varphi,d)$, where $\varphi$ is a mapping from $S$ to $T$ and $d$ an $S^{\star}\times S$-sorted mapping from $\Sigma$ to $\Lambda_{\varphi^{\star}\times \varphi}$ (thus a mapping in $\mathsf{Sig}(S)$). Note that $\mathsf{Sig}$ is the category obtained by means of the Grothendieck construction for an  appropriate contravariant functor from $\mathsf{Set}$ to $\mathsf{Cat}$.
% We next show that the category $\mathsf{Sig}_{\mathfrak{d}}$ can be obtained as the Kleisli category for a suitable monad. In fact, for every  set of sorts $S$, we have the adjoint situation $\mathbf{T}_{\mathrm{H}_{S}}\dashv\mathrm{G}_{\mathrm{H}_{S}}$
% %$$
% %\xymatrix@=5pc{ \mathsf{Alg}(\mathrm{H}_{S})
% %\ar@<1.5ex>[r]^{\mathrm{G}_{\mathrm{H}_{S}}} \ar@{}[r]|{\uadj} &
% %\mathbf{Set}^{S^{\star}\times S}
% %\ar@<1.5ex>[l]^{\mathbf{T}_{\mathrm{H}_{S}}} }
% %$$
% and, thus, a monad on $\mathsf{Sig}(S)$ which we will denote by $\mathbb{T}_{\mathrm{H}_{S}} = (\mathrm{T}_{\mathrm{H}_{S}},\eta^{\mathrm{H}_{S}},\mu^{\mathrm{H}_{S}})$. Then the ordered triple $\boldsymbol{\mathfrak{d}} = (\mathfrak{d},\eta,\mu)$ in which (1) $\mathfrak{d}$ is the endofunctor at $\mathsf{Sig}$ that sends $(S,\Sigma)$ to $(S,\mathrm{T}_{\mathrm{H}_{S}}(\Sigma))$ and a morphism $(\varphi,d)$ from $(S,\Sigma)$ to  $(T,\Lambda)$ to the morphism $(\varphi,d^{\sharp})$ from $(S,\mathrm{T}_{\mathrm{H}_{S}}(\Sigma))$ to  $(T,\mathrm{T}_{\mathrm{H}_{T}}(\Lambda))$, (2) $\eta$ the natural transformation from $\Id_{\mathsf{Sig}}$ to  $\mathfrak{d}$ that sends $(S,\Sigma)$ to $\eta_{(S,\Sigma)} = (\id,\eta^{\mathrm{H}_{S}}_{\Sigma})$, and (3) $\mu$ the natural transformation from $\mathfrak{d}\circ \mathfrak{d}$ to $\mathfrak{d}$ that sends $(S,\Sigma)$ to $\mu_{(S,\Sigma)} = (\id,\mu^{\mathrm{H}_{S}}_{\Sigma})$, is a monad on $\mathsf{Sig}$ and $\mathsf{Kl}(\boldsymbol{\mathfrak{d}})$, the Kleisli category for $\boldsymbol{\mathfrak{d}}$, is isomorphic to $\mathsf{Sig}_{\mathfrak{d}}$. This shows, in our opinion, the mathematical naturalness of the notion of derivor. Moreover, by defining a suitable notion of transformation between derivors one can equip the category $\mathsf{Sig}_{\mathfrak{d}}$ with a structure of $2$-category (this was, in fact, already done in~\cite{CVST10b} for polyderivors).
%
% \end{remark}
%
%
%
% \begin{remark}
% Since $\mathsf{Sig}$ has coproducts, $\mathsf{Kl}(\boldsymbol{\mathfrak{d}})\cong \mathsf{Sig}_{\mathfrak{d}}$ has coproducts. In fact, let $(\sig{\Sigma}^{i})_{i\in I}$ be an $I$-indexed family in $\mathrm{Ob}(\mathsf{Kl}(\boldsymbol{\mathfrak{d}}))$ and $(\coprod_{i\in I}\sig{\Sigma}^{i},(\mathrm{in}^{\sig{\Sigma}^{i}})_{i\in I})$ the coproduct of $(\sig{\Sigma}^{i})_{i\in I}$ in $\mathsf{Sig}$, where  the first coordinate of $\coprod_{i\in I}\sig{\Sigma}^{i}$ is $\coprod_{i\in I}S_{i}$ and the second coordinate of $\coprod_{i\in I}\sig{\Sigma}^{i}$, denoted by $\coprod_{i\in I}\mathrm{\Sigma}^{i}$, is defined, for every $(((s_{i_{j}},i_{j})_{j\in n}),(s,i))\in (\coprod_{i\in I}S_{i})^{\star}\times \coprod_{i\in I}S_{i}$, as follows: 
% \begin{equation*}
% \textstyle
% (\coprod_{i\in I}\mathrm{\Sigma}^{i})_{(((s_{i_{j}},i_{j})_{j\in n}),(s,i))} = 
% \begin{cases}
% \Sigma^{i}_{\mathbf{s},s}, \text{\,\,if $\exists\,\mathbf{s}\in S_{i}^{\star}$ such that $(s_{i_{j}},i_{j})_{j\in n} = \mathrm{in}_{i}^{\star}(\mathbf{s})$;}\\
% \varnothing, \text{\,\,otherwise}.
% \end{cases}
% \end{equation*}
% Let us recall that $\mathrm{in}_{i}^{\star}$ is the monoid homomorphism from the free monoid on $S_{i}$ to the free monoid on $\coprod_{i\in I}S_{i}$ that sends $\mathbf{s} = (s_{j})_{j\in n}\in S_{i}^{\star}$ to 
% $(s_{j},i)_{j\in n}\in (\coprod_{i\in I}S_{i})^{\star}$. Therefore, 
% $(s_{i_{j}},i_{j})_{j\in n} = \mathrm{in}_{i}^{\star}(\mathbf{s})$ if and only if, for every $j\in n$, $i_{j} = i$ and $s_{i_{j}} = w_{j}$.
% Then $(\boldsymbol{\mathfrak{d}}(\coprod_{i\in I}\sig{\Sigma}^{i}),(\eta_{\coprod_{i\in I}\sig{\Sigma}^{i}}\circ \mathrm{in}^{\sig{\Sigma}^{i}})_{i\in I})$ will be the required coproduct of $(\sig{\Sigma}^{i})_{i\in I}$ in $\mathsf{Kl}(\boldsymbol{\mathfrak{d}})$. This property of the category $\mathsf{Sig}_{\mathfrak{d}}$ will be essential in our work because to be able to compare Curry-Howard towers it will be necessary, as an intermediate step, to construct, by means of the formation of coproducts, from given many-sorted  signatures $(\sig{\Sigma}^{i})_{i\in I}$, a new many-sorted signature $\coprod_{i\in I}\sig{\Sigma}^{i}$ (finally, from $\coprod_{i\in I}\sig{\Sigma}^{i}$, we will construct the many-sorted signature $\biguplus_{i\in I}\sig{\Sigma}^{i}$ which will be the one of interest for us).
%
% %$\sig{\Sigma}$, $\sig{\Omega}_{\mathrm{c}}$, $\sig{\Omega}_{\mathrm{c},\mathcal{A}}$, $\sig{\Lambda}$, and $\sig{\Omega}_{\mathrm{c},\mathcal{B}}$ and $\mathsf{Sig}_{\mathfrak{d}}$-morphisms 
% %$\mathbf{d}$ from $\sig{\Sigma}$ to $\sig{\Lambda}$ and $\mathbf{e}$ from $\sig{\Omega}_{\mathrm{c},\mathcal{A}}$ to $\sig{\Omega}_{\mathrm{c},\mathcal{B}}$, as, e.g.,  
% %$$
% %\textstyle
% %\mathbf{d}\amalg\mathbf{id}_{\sig{\Omega}_{\mathrm{c}}}\colon \sig{\Sigma}\amalg\sig{\Omega}_{\mathrm{c}}\mor \sig{\Lambda}\amalg\sig{\Omega}_{\mathrm{c}} 
% %\quad\text{and}\quad
% %\mathbf{d}\amalg\mathbf{e}\colon \sig{\Sigma}\amalg\sig{\Omega}_{\mathrm{c},\mathcal{A}}\mor \sig{\Lambda}\amalg\sig{\Omega}_{\mathrm{c},\mathcal{B}}.
% %$$
% \end{remark}
%
% %\newpage
% \subsection{Adjunctions and the Grothendieck construction}\label{SAdjGrot}
%
% We recall the notions of adjunction and of the Grothendieck construction, which we use in the next subsection. We let ``$\ast$'' and ``$\circ$'' stand for the horizontal and vertical composition of natural transformations, respectively; we denote by $\Id_{\mathsf{X}}$ the identity functor at $\mathsf{X}$; and we denote by $\id_{G}$ the identity natural transformation at $G\colon \mathsf{A}\mor \mathsf{X}$.
%
% \begin{definition}\label{DAdj}
% An \emph{adjointness} is a $6$-tuple $(\mathsf{A},\mathsf{X},G,F,\eta,\varepsilon)$ where $\mathsf{A}$ and $\mathsf{X}$ are categories, $G$ a functor from $\mathsf{A}$ to $\mathsf{X}$, $F$ a functor from $\mathsf{X}$ to $\mathsf{A}$, $\eta$ a natural transformation from $\Id_{\mathsf{X}}$ to $G\circ F$ and $\varepsilon$ a natural transformation from $F\circ G$ to $\Id_{\mathsf{A}}$ subject to the so called triangular identities
% $$
% (G\ast \varepsilon)\circ (\eta\ast G) = \id_{G} \text{ and }
% (\varepsilon\ast F)\circ (F\ast\eta) = \id_{F}.
% $$
% If $(\mathsf{A},\mathsf{X},G,F,\eta,\varepsilon)$ is an adjointness, then we will say that $F$ is a \emph{left adjoint} of $G$, that $G$ is a \emph{right adjoint} of $F$, that $\eta$ is the \emph{unit} of the adjunction and that $\varepsilon$ is the \emph{counit} of the adjunction. We will also speak of an \emph{adjunction from} $\mathsf{X}$ \emph{to} $\mathsf{A}$, meaning by it a $4$-tuple $(G,F,\eta,\varepsilon)$ such that $(\mathsf{A},\mathsf{X},G,F,\eta,\varepsilon)$ is an adjointness, and we will write $F\dashv G$ to indicate that $F$ is a left adjoint of $G$.
% \end{definition}
%
% An adjunction determines an equivalence between the full subcategories of its fixed points (see~\cite{PT91}).
%
% \begin{proposition}\label{UOp}
% Let $(G,F,\eta,\varepsilon)$ be an adjunction from a category $\mathsf{X}$ to another $\mathsf{A}$. Then the categories $\mathsf{Fix}(\eta)$ and $\mathsf{Fix}(\varepsilon)$, where $\mathsf{Fix}(\eta)$ is the full subcategory of $\mathsf{X}$ determined by
% $$
% \{x\in \mathrm{Ob}(\mathsf{X})\mid \eta_{x}\text{ is an isomorphism}\}
% $$
% and $\mathsf{Fix}(\varepsilon)$ the full subcategory of $\mathsf{A}$ determined by
% $$
% \{a\in \mathrm{Ob}(\mathsf{A})\mid \varepsilon_{a}\text{ is an isomorphism}\},
% $$
% are equivalent, i.e., there exists a functor from $\mathsf{Fix}(\eta)$ to $\mathsf{Fix}(\varepsilon)$ which is an equivalence (one can say, following~\cite{Lk81}, that the equivalence between $\mathsf{Fix}(\eta)$ and $\mathsf{Fix}(\varepsilon)$ is the ``unity of opposites'').
% \end{proposition}
%
% \begin{definition}\label{DSplitIndexed}
% Let $\mathsf{I}$ be a category. A \emph{split} $\mathsf{I}$-\emph{indexed category} is a contravariant functor $F$ from $\mathsf{I}$ to $\mathsf{Cat}$. Given an object $i\in \mathrm{Ob}(\mathsf{I})$, we write $\mathsf{F}_{i}$ for the category $F(i)$, and given a morphism $\varphi\in \mathrm{Hom}_{\mathsf{I}}(i,j)$, we write $F_{\varphi}$ for the functor $F(\varphi)\colon \mathsf{F}_{j}\mor \mathsf{F}_{i}$. A \emph{split indexed category} is an ordered pair $(\mathsf{I},F)$ in which $\mathsf{I}$ is a category and $F$ a split $\mathsf{I}$-indexed category.
% \end{definition}
%
% We next define the Grothendieck construction, which assigns to a split indexed category a split fibration (see~\cite{Gro71} and, for split fibrations, \cite{BW12}, pp.~327--329). We follow the notation of~\cite{JY21}, Chapter~10, but add the domain of $F$ as a superscript to indicate the contravariance of $F$.
%
% \begin{definition}\label{DGConst}
% Let $(\mathsf{I},F)$ be a split indexed category. Then the \emph{Grothen\-dieck construction} at $(\mathsf{I},F)$ is $(\int^{\mathsf{I}}F,\pi_{F})$, where $\int^{\mathsf{I}}F$ is the category defined as follows:
% \begin{enumerate}
% \item $\mathrm{Ob}(\int^{\mathsf{I}}F) = \bigcup_{i\in\mathrm{Ob}(\mathsf{I})}(\{i\}\times\mathrm{Ob}(\mathsf{F}_{i}))$.
% \item For every $(i,x), (j,y)\in \mathrm{Ob}(\int^{\mathsf{I}}F)$, $\mathrm{Hom}_{\int^{\mathsf{I}}F}((i,x),(j,y))$ is the set of all ordered pairs $(\varphi,f)$, where $\varphi\in \mathrm{Hom}_{\mathsf{I}}(i,j)$ and $f\in \mathrm{Hom}_{\mathsf{F}_{i}}(x,F_{\varphi}(y))$.
% \item For every $(i,x)\in \mathrm{Ob}(\int^{\mathsf{I}}F)$, the identity morphism at $(i,x)$ is $(\id_{i},\id_{x})$.
% \item\label{DGrotCons} For every $(i,x), (j,y), (k,z)\in \mathrm{Ob}(\int^{\mathsf{I}}F)$, every $(\varphi,f)\colon (i,x)\mor (j,y)$, and every $(\psi,g)\colon (j,y)\mor (k,z)$, the composite morphism $(\psi,g)\circ (\varphi,f)$ from $(i,x)$ to $(k,z)$ is
% $$
% (\psi,g)\circ (\varphi,f) = (\psi\circ \varphi,F_{\varphi}(g)\circ f)\colon (i,x)\mor (k,z).
% $$
% Notice that $f\colon x\mor F_{\varphi}(y)$ and $g\colon y\mor F_{\psi}(z)$, hence, since $F_{\varphi}$ is a functor from $\mathsf{F}_{j}$ to $\mathsf{F}_{i}$, $F_{\varphi}(g)\colon F_{\varphi}(y)\mor F_{\varphi}(F_{\psi}(z))$. Therefore $F_{\varphi}(g)\circ f\colon x\mor F_{\varphi}(F_{\psi}(z)) = F_{\psi\circ\varphi}(z)$.
% \end{enumerate}
% And $\pi_{F}$ is the functor from $\int^{\mathsf{I}}F$ to $\mathsf{I}$ that sends an object $(i,x)$ of $\int^{\mathsf{I}}F$ to $i$ and a morphism $(\varphi,f)\colon (i,x)\mor (j,y)$ of $\int^{\mathsf{I}}F$ to $\varphi\colon i\mor j$. The functor $\pi_{F}$ is a split fibration, the \emph{canonical split fibration} from $\int^{\mathsf{I}}F$ to $\mathsf{I}$ (determined by $(\mathsf{I},F)$).
% \end{definition}
%
% \subsection{\texorpdfstring{The category $\mathsf{Alg}_{\mathfrak{d}}$}
% {The category Algd}}
%
%
% In this section, we associate to each derivor $\mathbf{d}\colon\sig{\Sigma}\mor\sig{\Lambda}$ a functor from the category of $\Lambda$-algebras to the category of $\Sigma$-algebras. This construction can be extended to a contravariant functor from $\mathsf{Sig}_{\mathfrak{d}}$ to $\mathsf{Cat}$. We will use this functor to obtain, through the Grothendieck construction, the category $\mathsf{Alg}_{\mathfrak{d}}$. Moreover, we prove the existence of two opposite functors between $\mathsf{Sig}_{\mathfrak{d}}$  and $\mathsf{Alg}_{\mathfrak{d}}$ that give rise to an adjunction. We conclude the section by proving that the category $\mathsf{Sig}_{\mathfrak{d}}$ is equivalent to a full subcategory of $\mathsf{Alg}_{\mathfrak{d}}$, denoted by $\mathsf{Tw}^{(-1)}_{\mathfrak{d}}$, which we will call towers at level $-1$ in this work.
%
% We start by associating a functor to every derivor.
%
% \begin{proposition}\label{PFunSig}
% Let $\mathbf{d}\colon\sig{\Sigma}\mor\sig{\Lambda}$ be a morphism in
% $\mathsf{Sig}_{\mathfrak{d}}$. Then $\mathrm{Alg}_{\mathfrak{d}}(\mathbf{d})$, also denoted by $\mathbf{d}^{\ast}_{\mathfrak{d}}$ for short, is the functor from  $\mathsf{Alg}(\sig{\Lambda})$ to $\mathsf{Alg}(\sig{\Sigma})$ that sends $(B,G)$ to $(B_{\varphi},G^{\mathbf{d}})$ and a homomorphism $f$ from $(B,G)$ to $(B',G')$ to the homomorphism $f_{\varphi}$ from $(B_{\varphi},G^{\mathbf{d}})$ to  $(B'_{\varphi},{G'}^{\mathbf{d}})$,
% %definido como %
% %$$\xymatrixcolsep={65pt} %
% %  \xymatrixrowsep={35pt} %
% %\functor
% %{\mathsf{Alg}(T,\Lambda)}{\mathrm{Alg}_{\mathfrak{d}}(\varphi,d)}{\mathsf{Alg}(S,\Sigma)}
% %{(B,G)}{(B_{\varphi},G^{(\varphi,d)})}
% %{f}{f_{\varphi}}
% %{(B',G')}{(B'_{\varphi},{G'}^{(\varphi,d)})}
% %$$ %
% where, for every $\sig{\Lambda}$-algebra $(B,G)$,
% $G^{\mathbf{d}} = G^{\sharp}_{\varphi^{\star}\times \varphi}\circ d$
% is obtained from
% % $$
% % \xymatrix{
% % & & \Sigma \ar[d]^{d} \\
% % \Lambda \ar[r]^{\eta_{\Lambda}}
% %         \ar[rd]_{G} &
% % \Ter_{\Hall_{T}}(\Lambda) \ar[d]^{\ext{G}}
% %             \ar@{}[r]|{\text{como}}
% %             &
% % \Ter_{\Hall_{T}}(\Lambda)_{\varphi^{\star}\times \varphi}
% % \ar[d]^{\ext{G}_{\varphi^{\star}\times \varphi}} \\
% % & \Op_{\Hall_{T}}(B)  &
% % \Op_{\Hall_{T}}(B)_{\varphi^{\star}\times \varphi}=\Op_{\Hall_{T}}(B_{\varphi})
% % }$$
% $$
% \begin{aligned}
% \xymatrix@C=40pt@R=33pt{
% \Lambda
%   \ar[r]^-{\eta^{\Lambda}}
%   \ar[rd]_-{G} &
% \mathrm{Tm}_{\mathrm{H}_{T}}(\Lambda) \ar[d]^{G^{\sharp}}
%   \\
% & \mathrm{O}_{\mathrm{H}_{T}}(B)
% }
% %
% \end{aligned}
% \qquad\quad \text{as} \quad
% \begin{aligned}
% %
% \xymatrix@C=40pt@R=33pt{
% \mathrm{Tm}_{\mathrm{H}_{T}}(\Lambda)_{\varphi^{\star}\times \varphi}
% \ar[d]^{G^{\sharp}_{\varphi^{\star}\times \varphi}} &
% \Sigma \ar[l]_-{d}\\
% {\mathrm{O}_{\mathrm{H}_{T}}(B)_{\varphi^{\star}\times\varphi}}=
% \mathrm{O}_{\mathrm{H}_{S}}(B_{\varphi})
% }
% \end{aligned}
% $$
% \end{proposition}
%
% \begin{proof}
% For every $\sig{\Lambda}$-algebra $(B,G)$, $G\colon\Lambda\mor\mathrm{O}_{\mathrm{H}_{T}}(B)$, and since  $\mathbf{O}_{\mathrm{H}_{T}}(B)$ is a Hall algebra, $G$  can be extended up to the free Hall algebra on $\Lambda$. Moreover, we have that $\mathrm{O}_{\mathrm{H}_{T}}(B)_{\varphi^{\star}\times \varphi} = \mathrm{O}_{\mathrm{H}_{S}}(B_{\varphi})$ since, for every $(\mathbf{s},s)\in S^{\star}\times S$ it holds that
% \allowdisplaybreaks
% \begin{align*}
%     (\mathrm{O}_{\mathrm{H}_{T}}(B)_{\varphi^{\star}\times\varphi})_{\mathbf{s},s} &=
%     \mathrm{O}_{\mathrm{H}_{T}}(B)_{\varphi^{\star}(\mathbf{s}),\varphi(s)} \\
%     &= B_{\varphi^{\star}(\mathbf{s})}\mor B_{\varphi(s)} \\
%     &= B_{(\varphi(s_{0}),\ldots,\varphi(s_{\bb{\mathbf{s}}-1}))}\mor B_{\varphi(s)} \\
%     &= \textstyle\prod(B_{\varphi(s_{i})}\mid i\in\bb{\mathbf{s}}) \mor B_{\varphi(s)} \\
%     &= \textstyle\prod((B_{\varphi})_{s_{i}}\mid i\in\bb{\mathbf{s}}) \mor (B_{\varphi})_{s} \\
%     &= (B_{\varphi})_{\mathbf{s}} \mor (B_{\varphi})_{s} \\
%     &= \mathrm{O}_{\mathrm{H}_{S}}(B_{\varphi})_{\mathbf{s},s},
% \end{align*}
% thus, so defined, $G^{(\varphi,d)}$ is  an algebraic structure on $B_{\varphi}$.
%
% Let $f\colon(B,G)\mor(B',G')$ be a homomorphism of $\sig{\Lambda}$-algebras, $(\mathbf{s},s)\in S^{\star}\times S$, and $\sigma\in\Sigma_{\mathbf{s},s}$. Then  $f_{\varphi}\colon(B_{\varphi},G^{\mathbf{d}}) \mor (B'_{\varphi},{G'}^{\mathbf{d}})$ is a homomorphism of $\sig{\Sigma}$-algebras, because $G^{\mathbf{d}}(\sigma)$ is a term operation and, consequently, the following diagram
% \begin{center}
% \begin{tikzpicture}
% [ACliment/.style={-{To [angle'=45, length=5.75pt, width=4pt, round]}},
% RHACliment/.style={right hook-{To [angle'=45, length=5.75pt, width=4pt, round]}, font=\scriptsize}, 
% scale=1
% ]
%
% \node[] (s) at (0,0) {${B_{\varphi}}_{\mathbf{s}}$};
% \node[] (s') at (0,-2) {${B'_{\varphi}}_{\mathbf{s}}$};
% \node[] (Vs) at (4,0) {$V^{S}_{\mathbf{s}}$};
% \node[] (Vs') at (4,-2) {$V^{S}_{\mathbf{s}'}$};
%
% \draw[ACliment]  (s) to node [above ] {$G^{(\varphi,d)}(\sigma)$} (Vs);
% \draw[ACliment]  (s') to node [below ] {$G'^{(\varphi,d)}(\sigma)$} (Vs');
% \draw[ACliment]  (s) to node [left ] {${f_{\varphi}}_{\mathbf{s}}$} (s');
% \draw[ACliment]  (Vs) to node [right ] {${f_{\varphi}}_{s}$} (Vs');
% \end{tikzpicture}
% \end{center}
% commutes. Moreover, $(g\circ f)_{\varphi}=g_{\varphi}\circ f_{\varphi}$, so that  $\mathrm{Alg}_{\mathfrak{d}}(\mathbf{d})$ is a functor.
% \end{proof}
%
%
% The next proposition establishes that for any derivation $\mathbf{d}=(\varphi,d)$ from $\sig{\Sigma}=(S,\Sigma)$ to $\sig{\Lambda}=(T,\Lambda)$ and any $\Lambda$-algebra $\mathbf{B}$,  every $\Lambda$-closed subset of $\mathbf{B}$ induces a $\Sigma$-closed subset of the derived algebra $\mathbf{d}^{\ast}_{\mathfrak{d}}(\mathbf{B})$ and every $\Lambda$-congruence of $\mathbf{B}$ induces a $\Sigma$-congruence of the derived algebra $\mathbf{d}^{\ast}_{\mathfrak{d}}(\mathbf{B})$.
%
%
% \begin{proposition}\label{PDerClCong}
% Let $\sig{\Sigma}=(S,\Sigma)$ and $\sig{\Lambda}=(T,\Lambda)$ be many-sorted signatures and $\mathbf{d}=(\varphi,d)$ a derivor from $\sig{\Sigma}$ to $\sig{\Lambda}$. Let $\mathbf{B}=(B,G)$ be a $\Lambda$-algebra and $\mathbf{d}^{\ast}_{\mathfrak{d}}(\mathbf{B})=(B_{\varphi}, G^{\mathbf{d}})$ the derived $\Sigma$-algebra.
% \begin{enumerate}
%     \item If $Y$ is a $\Lambda$-closed subset of $\mathbf{B}$, then $Y_{\varphi}$ is a $\Sigma$-closed subset of $\mathbf{d}^{\ast}_{\mathfrak{d}}(\mathbf{B})$;
%     \item If $\Phi$ is a $\Lambda$-congruence of $\mathbf{B}$, then $\Phi_{\varphi}$ is a $\Sigma$-congruence of $\mathbf{d}^{\ast}_{\mathfrak{d}}(\mathbf{B})$.
% \end{enumerate}
% \end{proposition}
%
%
% \begin{proof}
%     (1)
%     Let $Y \subseteq B$ be a $\Lambda$-closed subset of $\mathbf{B}$ and let $\mathbf{Y}$ be the corresponding subalgebra. Thus, the canonical inclusion mapping $\mathrm{in}^{Y, B}$ is a $\Lambda$-homomorphism from $\mathbf{Y}$ to $\mathbf{B}$. Therefore, it follows from Proposition~\ref{PFunSig}, that $\mathbf{d}_{\mathfrak{d}}^{\ast}(\mathrm{in}^{Y,B}) = \mathrm{in}^{Y, B}_{\varphi} = \mathrm{in}^{Y_{\varphi}, B_{\varphi}}$ is a $\Sigma$-homomorphism from $\mathbf{d}_{\mathfrak{d}}^{\ast}(\mathbf{Y})$ to $\mathbf{d}_{\mathfrak{d}}^{\ast}(\mathbf{B})$. Moreover, since $Y_{\varphi}$ is a subset of $B_{\varphi}$, it follows that $\mathbf{d}_{\mathfrak{d}}^{\ast}(\mathbf{Y})$ is a subalgebra of $\mathbf{d}_{\mathfrak{d}}^{\ast}(\mathbf{B})$. Therefore, the $S$-sorted set $Y_{\varphi}$ is a $\Sigma$-closed subset of $\mathbf{d}_{\mathfrak{d}}^{\ast}(\mathbf{B})$.
%
%     (2)
%     Let $\Phi \subseteq B \times B$ be a $\Lambda$-congruence of $\mathbf{B}$. Thus, the projection mapping $\mathrm{pr}^{\Phi}$ is a $\Lambda$-homomorphism from $\mathbf{B}$ to $\mathbf{B}/\Phi$. Therefore, it follows from Proposition~\ref{PFunSig}, that $\mathbf{d}_{\mathfrak{d}}^{\ast}(\mathrm{pr}^{\Phi}) = \mathrm{pr}^{\Phi}_{\varphi} = \mathrm{pr}^{\Phi_{\varphi}}$ is a $\Sigma$-homomorphism from $\mathbf{d}_{\mathfrak{d}}^{\ast}(\mathbf{B})$ to $\mathbf{d}_{\mathfrak{d}}^{\ast}(\mathbf{B}/\Phi)$. Hence, its kernel, $\mathrm{Ker}(\mathrm{pr}^{\Phi_{\varphi}}) = \Phi_{\varphi}$ is a $\Sigma$-congruence of $\mathbf{d}^{\ast}_{\mathfrak{d}}(\mathbf{B})$.
% \end{proof}
%
%
%
%
%
%
%
%
%
%
%
%
%
%
%
%
%
%
% The previous construction can be extended up to a contravariant functor from the category $\mathsf{Sig}_{\mathfrak{d}}$ to the category $\mathsf{Cat}$.
%
% \begin{proposition}\label{PAlgFun}
% From $\mathsf{Sig}_{\mathfrak{d}}$ to $\mathsf{Cat}$ there exists a contravariant functor,
% denoted by $\mathrm{Alg}_{\mathfrak{d}}$, that sends $\sig{\Sigma}=(S,\Sigma)$ to $\mathsf{Alg}(\sig{\Sigma})$ and a morphism  $\mathbf{d}=(\varphi,d)$ from $\sig{\Sigma}=(S,\Sigma)$ to $\sig{\Lambda}=(T,\Lambda)$ to the functor $\mathrm{Alg}_{\mathfrak{d}}(\mathbf{d})$ from  $\mathsf{Alg}(\sig{\Lambda})$ to $\mathsf{Alg}(\sig{\Sigma})$.
% \end{proposition}
%
% \begin{proof}
% Given $(\varphi,d)\colon(S,\Sigma) \mor (T,\Lambda)$ and $(\psi,e)\colon (T,\Lambda)\mor (U,\Omega)$, we show that %
% % $$
% % \xymatrix{
% % (S,\Sigma) \ar[r]^{(\varphi,d)} &
% % (T,\Lambda) \ar[r]^{(\psi,e)} &
% % (U,\Omega)
% % }
% % $$
% $\mathrm{Alg}_{\mathfrak{d}}(\varphi,d)\circ\mathrm{Alg}_{\mathfrak{d}}(\psi,e)= \mathrm{Alg}_{\mathfrak{d}}((\psi,e)\circ(\varphi,d))$.
%
% Let $(A,F)$ be a $(U,\Omega)$-algebra. Then ${A_{\psi}}_{\varphi}=A_{\psi\circ\varphi}$. Moreover, we have that %
% %
% \allowdisplaybreaks
% \begin{align*}
% {F^{(\psi,e)}}^{(\varphi,d)}
%   &=
% (F^{\sharp}_{\psi^{\star}\times\psi}\circ e)^{(\varphi,d)} \\
%   &=
% (F^{\sharp}_{\psi^{\star}\times\psi}\circ e)^{\sharp}_{\varphi^{\star}\times\varphi}\circ d \\
%   &=
% ({F^{\sharp}_{\psi^{\star}\times\psi}}_{\varphi^{\star}\times\varphi}\circ
%     e^{\sharp}_{\varphi^{\star}\times\varphi})\circ d \\
%   &=
% {F^{\sharp}_{\psi^{\star}\times\psi}}_{\varphi^{\star}\times\varphi}\circ
%     (e^{\sharp}_{\varphi^{\star}\times\varphi}\circ d )\\
%   &=
% F^{\sharp}_{(\psi\circ\varphi)^{\sharp}\times(\psi\circ\varphi)}
%     \circ (e^{\sharp}_{\varphi^{\star}\times\varphi}\circ d )\\
%   &=
% F^{((\psi\circ\varphi),e^{\sharp}_{\varphi^{\star}\times\varphi}\circ d)} \\
%   &=
% F^{(\psi,e)\circ(\varphi,d)}
% \end{align*}
% %
% thus $({A_{\psi}}_{\varphi},{F^{(\psi,e)}}^{(\varphi,d)})= (A_{\psi\circ\varphi},F^{(\psi,e)\circ (\varphi,d)})$.  Finally, if $f$ is a homomorphism of $(U,\Omega)$-algebras, then ${f_{\psi}}_{\varphi} = f_{\psi\circ\varphi}$.
% \end{proof}
%
%
% From the definition of the functor $\mathrm{Alg}_{\mathfrak{d}}$, for every derivor $\mathbf{d}\colon\sig{\Sigma}\mor\sig{\Lambda}$, it is obvious that the following diagram %
% \begin{center}
% \begin{tikzpicture}
% [ACliment/.style={-{To [angle'=45, length=5.75pt, width=4pt, round]}},
% RHACliment/.style={right hook-{To [angle'=45, length=5.75pt, width=4pt, round]}, font=\scriptsize}, 
% scale=1
% ]
%
% \node[] (s) at (0,0) {$\mathsf{Alg}(\sig{\Sigma})$};
% \node[] (s') at (0,-2) {$\mathsf{Alg}(\sig{\Lambda})$};
% \node[] (Vs) at (4,0) {$\mathsf{Set}^{S}$};
% \node[] (Vs') at (4,-2) {$\mathsf{Set}^{T}$};
%
% \draw[ACliment]  (s) to node [above ] {$\G_{\sig{\Sigma}}$} (Vs);
% \draw[ACliment]  (s') to node [below ] {$\G_{\sig{\Lambda}}$} (Vs');
% \draw[ACliment]  (s') to node [left ] {$\mathrm{Alg}_{\mathfrak{d}}(\mathbf{d})$} (s);
% \draw[ACliment]  (Vs') to node [right ] {$\Delta_{\varphi}$} (Vs);
% \end{tikzpicture}
% \end{center}
% commutes.
%
% \begin{definition}\label{DAlgCat}
% The category $\mathsf{Alg}_{\mathfrak{d}}$ is $\int^{\mathsf{Sig}_{\scriptscriptstyle\mathfrak{d}}}\mathrm{Alg}_{\mathfrak{d}}$, i.e., the category obtained by means of the Grothendieck construction applied to the contravariant functor $\mathrm{Alg}_{\mathfrak{d}}$. Thus the category $\mathsf{Alg}_{\mathfrak{d}}$ has as objects the ordered pairs $(\sig{\Sigma},\mathbf{A})$, with $\sig{\Sigma}$ a many-sorted signature and $\mathbf{A}$ a $\Sigma$-algebra, and as morphisms from $(\sig{\Sigma},\mathbf{A})$ to $(\sig{\Lambda},\mathbf{B})$ the ordered pairs $(\mathbf{d},h)$, with $\mathbf{d}$ a derivor from $\sig{\Sigma}$ to $\sig{\Lambda}$ and $h$ a homomorphism of $\Sigma$-algebras from $\mathbf{A}$ to $\mathbf{d}^{\ast}_{\mathfrak{d}}(\mathbf{B})$.
%
% Let $P^{(-1)}_{\mathfrak{d}}$ be equal to $\pi_{\mathrm{Alg}_{\mathfrak{d}}}$, the canonical split fibration from $\mathsf{Alg}_{\mathfrak{d}}$ to $\mathsf{Sig}_{\mathfrak{d}}$ introduced in Definition~\ref{DGConst}. We recall that $P^{(-1)}_{\mathfrak{d}}\colon \mathsf{Alg}_{\mathfrak{d}}\mor \mathsf{Sig}_{\mathfrak{d}}$ is the functor that sends and object $(\sig{\Sigma},\mathbf{A})$ in $\mathsf{Alg}_{\mathfrak{d}}$ to $\sig{\Sigma}$ in $\mathsf{Sig}_{\mathfrak{d}}$ and a morphism $(\mathbf{d},h)\colon (\sig{\Sigma},\mathbf{A})\mor (\sig{\Lambda},\mathbf{B})$ in $\mathsf{Alg}_{\mathfrak{d}}$ to $\mathbf{d}\colon \sig{\Sigma}\mor \sig{\Lambda}$ in $\mathsf{Sig}_{\mathfrak{d}}$.
%
% %Moreover, we let $\mathsf{FAlg}_{\mathfrak{d}}$ stand for the full subcategory of $\mathsf{Alg}_{\mathfrak{d}}$ whose objects are the ordered pairs $(\sig{\Sigma},\mathbf{T}_{\Sigma}(X))$, with $\sig{\Sigma} = (S,\Sigma)$ a many-sorted signature and $X$ an $S$-sorted set.
% \end{definition}
%
% \begin{remark}
% There exists a pseudo-functor $\mathrm{Alg}_{\boldsymbol{\mathfrak{d}}}$ from the 2-category $\mathsf{Sig}_{\boldsymbol{\mathfrak{d}}}$ to the 2-category $\mathsf{Cat}$, contravariant in the 1-cells, i.e., the derivors, and covariant in the 2-cells, i.e., the transformations between derivors (this was already done in~\cite{CVST10b} for polyderivors).
% \end{remark}
%
% We will prove the existence of a functor from $\mathsf{Sig}_{\mathfrak{d}}$ to $\mathsf{Alg}_{\mathfrak{d}}$ that is left adjoint to the canonical split fibration $P^{(-1)}_{\mathfrak{d}}$.  To this end, we first introduce the functor $J^{(-1)}_{\mathfrak{d}}$, from $\mathsf{Sig}_{\mathfrak{d}}$ to $\mathsf{Alg}_{\mathfrak{d}}$.
%
% \begin{proposition}\label{PJmu}
% Let $J^{(-1)}_{\mathfrak{d}}$ be the assignment from $\mathsf{Sig}_{\mathfrak{d}}$ to $\mathsf{Alg}_{\mathfrak{d}}$ defined as follows
% \begin{enumerate}
%     \item To every many-sorted signature $\sig{\Sigma}=(S,\Sigma)$ in $\mathsf{Sig}_{\mathfrak{d}}$, 
%     $$J^{(-1)}_{\mathfrak{d}}(\sig{\Sigma})=\left(\sig{\Sigma}, \mathbf{T}_{\Sigma}\left(\varnothing^{S}\right)\right),$$
%     where $\varnothing^{S}=(\varnothing)_{s\in S}$;
%     \item To every pair of many-sorted signatures $\sig{\Sigma}=(S,\Sigma)$ and $\sig{\Lambda}=(T,\Lambda)$, and every derivor $\mathbf{c}=(\varphi,c)$ from $\sig{\Sigma}$ to $\sig{\Lambda}$ in $\mathsf{Sig}_{\mathfrak{d}}$, 
%     $$J^{(-1)}_{\mathfrak{d}}(\mathbf{c})=\left(\mathbf{c},\left(!^{\mathbf{c}^{\ast}_{\mathfrak{d}}(\mathbf{T}_{\Lambda}(\varnothing^{T}))}\right)^{\sharp}\right),$$
%     where $!^{\mathbf{c}^{\ast}_{\mathfrak{d}}(\mathbf{T}_{\Lambda}(\varnothing^{T}))}=(!^{\mathrm{T}_{\Lambda}(\varnothing^{T})_{\varphi(s)}})_{s\in S}$ is the $S$-sorted mapping from $\varnothing^{S}$ to $\mathbf{c}^{\ast}_{\mathfrak{d}}(\mathbf{T}_{\Lambda}(\varnothing^{T}))$, whose $s$-th component, for $s\in S$,
%     is $!^{\mathrm{T}_{\Lambda}(\varnothing^{T})_{\varphi(s)}}$, the unique mapping from $\varnothing$ to $\mathrm{T}_{\Lambda}(\varnothing^{T})_{\varphi(s)}$, and $(!^{\mathbf{c}^{\ast}_{\mathfrak{d}}(\mathbf{T}_{\Lambda}(\varnothing^{T}))})^{\sharp}$ is the unique $\Sigma$-ho\-mo\-mor\-phism from $\mathbf{T}_{\Sigma}(\varnothing^{S})$ to $\mathbf{c}^{\ast}_{\mathfrak{d}}(\mathbf{T}_{\Lambda}(\varnothing^{T}))$ satisfying that 
%     $$(!^{\mathbf{c}^{\ast}_{\mathfrak{d}}(\mathbf{T}_{\Lambda}(\varnothing^{T}))})^{\sharp}\circ \eta^{\varnothing^{S}}=\,!^{\mathbf{c}^{\ast}_{\mathfrak{d}}(\mathbf{T}_{\Lambda}(\varnothing^{T}))}.$$
%
%     Note that $\eta^{\varnothing^{S}}=\,!^{\mathrm{T}_{\Sigma}(\varnothing^{S})}=\,(!^{\mathrm{T}_{\Sigma}(\varnothing^{S})_{s}})_{s\in S}$.
% \end{enumerate}
% Then, $J^{(-1)}_{\mathfrak{d}}$ is a functor from $ \mathsf{Sig}_{\mathfrak{d}}$ to $\mathsf{Alg}_{\mathfrak{d}}$.
% \end{proposition}
% \begin{proof}
%     That $J^{(-1)}_{\mathfrak{d}}$ sends objects and morphisms in $\mathsf{Sig}_{\mathfrak{d}}$ to objects and morphisms in  $\mathsf{Alg}_{\mathfrak{d}}$ is clear from the definition. 
%
%     We will prove that $J^{(-1)}_{\mathfrak{d}}$ preserves identities and compositions.
%
% \textsf{$J^{(-1)}_{\mathfrak{d}}$ preserves identities.}
%
% Let $\sig{\Sigma}=(S,\Sigma)$ be a many-sorted signature in $\mathsf{Sig}_{\mathfrak{d}}$. We need to prove that 
% \begin{equation*}
%     J^{(-1)}_{\mathfrak{d}}\left(\id^{\sig{\Sigma}}_{\mathsf{Sig}_{\mathfrak{d}}}\right)=\id^{J^{(-1)}_{\mathfrak{d}}(\sig{\Sigma})}_{\mathsf{Alg}_{\mathfrak{d}}}.
%     \tag{Id}\label{EPJmuId}
% \end{equation*}
%
% The following chain of equalities holds
% \allowdisplaybreaks
% \begin{align*}
%     J^{(-1)}_{\mathfrak{d}}\left(\id^{\sig{\Sigma}}\right)&=
%     \left(
%     \id^{\sig{\Sigma}},
%     \left(
%     !^{\left(
%     \id^{\sig{\Sigma}}
%     \right)^{\ast}_{\mathfrak{d}}\left(\mathbf{T}_{\Sigma}\left(\varnothing^{S}\right)\right)}
%     \right)^{\sharp}
%     \right)
%     \tag{1}
%     \\&=
%     \left(
%     \id^{\sig{\Sigma}},
%     \left(
%     !^{\mathbf{T}_{\Sigma}\left(\varnothing^{S}\right)}
%     \right)^{\sharp}
%     \right)
%     \tag{2}
%     \\&=
%     \left(
%     \id^{\sig{\Sigma}},
%     \id^{\mathbf{T}_{\Sigma}\left(\varnothing^{S}\right)}
%     \right)
%     \tag{3}
%     \\&=
%     \id^{\left(\sig{\Sigma},\mathbf{T}_{\Sigma}\left(\varnothing^{S}\right)\right)}
%     \tag{4}
%     \\&=
%     \id^{J^{(-1)}_{\mathfrak{d}}(\sig{\Sigma})}.
%     \tag{5}
% \end{align*}
%
% In the just stated chain of equalities, the first equality unravels the definition of $J^{(-1)}_{\mathfrak{d}}$ at $\id^{\sig{\Sigma}}$; the second equality follows from Proposition~\ref{PAlgFun}; the third equality follows from the fact that $(!^{\mathbf{T}_{\Sigma}(\varnothing^{S})})^{\sharp}\circ\eta^{\varnothing^{S}}=!^{\mathbf{T}_{\Sigma}(X)}=\eta^{\varnothing^{S}}$, so $(
%     !^{\mathbf{T}_{\Sigma}(\varnothing^{S})}
%     )^{\sharp}=\id^{\mathbf{T}_{\Sigma}(\varnothing^{S})}$; the fourth equality recovers the identity mapping in $\mathsf{Alg}_{\mathfrak{d}}$ at $(\sig{\Sigma},\mathbf{T}_{\Sigma}(\varnothing^{S}))$; finally the last equality recovers the definition of $J^{(-1)}_{\mathfrak{d}}(\sig{\Sigma})$.
%
% We conclude that Equation~\ref{EPJmuId} holds, that is, $J^{(-1)}_{\mathfrak{d}}$ preserves identities.
%
% \textsf{$J^{(-1)}_{\mathfrak{d}}$ preserves compositions.}
%
% Let $\sig{\Sigma}=(S,\Sigma)$, $\sig{\Lambda}=(T,\Lambda)$, and $\sig{\Omega}=(U,\Omega)$ be three many-sorted signatures in $\mathsf{Sig}_{\mathfrak{d}}$ and $\mathbf{c}\colon \sig{\Sigma}\mor \sig{\Lambda}$ and $\mathbf{d}\colon \sig{\Lambda}\mor \sig{\Omega}$ be two derivors, with $\mathbf{c}=(\varphi,c)$ and $\mathbf{d}=(\psi,d)$. We need to prove that 
% \begin{equation*}
%     J^{(-1)}_{\mathfrak{d}}\left(
%     \mathbf{d}\circ \mathbf{c}
%     \right)
%     =
%     J^{(-1)}_{\mathfrak{d}}\left(
%     \mathbf{d}\right)
%     \circ 
%     J^{(-1)}_{\mathfrak{d}}\left(
%     \mathbf{c}\right).
%     \tag{Comp}\label{EPJmuComp}
% \end{equation*}
%
% Note that following chain of equalities holds
% \allowdisplaybreaks
% \begin{align*}
%     J^{(-1)}_{\mathfrak{d}}\left(
%     \mathbf{d}\circ \mathbf{c}
%     \right)
%     &=
%     \left(
%     \mathbf{d}\circ \mathbf{c},
%     \left(
%     !^{
%     \left(\mathbf{d}\circ \mathbf{c}\right)^{\ast}_{\mathfrak{d}}
%     \left(
%     \mathbf{T}_{\Omega}\left(\varnothing^{U}\right)
%     \right)
%     }
%     \right)^{\sharp}
%     \right)
%     \tag{1}
%     \\&=
%     \left(
%     \mathbf{d}\circ \mathbf{c},
%     \left(
%     !^{\mathrm{T}_{\Omega}(\varnothing^{U})_{\psi\circ\varphi}}
%     \right)^{\sharp}
%     \right)
%     \tag{2}
%     \\&=
%     \left(
%     \mathbf{d}\circ \mathbf{c},
%     \left(
%     \left(
%     !^{\mathrm{T}_{\Omega}(\varnothing^{U})_{\psi}}
%     \right)^{\sharp}_{\varphi}
%     \circ \,
%     !^{\mathrm{T}_{\Lambda}(\varnothing^{T})_{\varphi}}
%     \right)^{\sharp}
%     \right)
%     \tag{3}
%     \\&=
%     \left(
%     \mathbf{d}\circ \mathbf{c},
%     \left(
%     !^{\mathrm{T}_{\Omega}(\varnothing^{U})_{\psi}}
%     \right)^{\sharp}_{\varphi}
%     \circ 
%     \left(
%     !^{\mathrm{T}_{\Lambda}(\varnothing^{T})_{\varphi}}
%     \right)^{\sharp}
%     \right)
%     \tag{4}
%     \\&=
%     \left(
%     \mathbf{d},\left(
%     !^{\mathrm{T}_{\Omega}(\varnothing^{U})_{\psi}}
%     \right)^{\sharp}
%     \right)
%     \circ 
%     \left(
%     \mathbf{c},
%     \left(
%     !^{\mathrm{T}_{\Lambda}(\varnothing^{T})_{\varphi}}
%     \right)^{\sharp}
%     \right)
%     \tag{5}
%     \\&=
%     \left(
%     \mathbf{d},
%     \left(
%     !^{\mathbf{d}^{\ast}_{\mathfrak{d}}\left(\mathbf{T}_{\Omega}(\varnothing^{U})\right)}
%     \right)^{\sharp}
%     \right)
%     \circ 
%     \left(
%     \mathbf{c},
%     \left(
%     !^{\mathbf{c}^{\ast}_{\mathfrak{d}}\left(\mathbf{T}_{\Lambda}(\varnothing^{T})\right)}
%     \right)^{\sharp}
%     \right)
%     \tag{6}
%     \\&=
%     J^{(-1)}_{\mathfrak{d}}\left(
%     \mathbf{d}\right)
%     \circ 
%     J^{(-1)}_{\mathfrak{d}}\left(
%     \mathbf{c}\right).
%     \tag{7}
% \end{align*}
%
% In the just stated chain of equalities, the first equality unravels the definition of $J^{(-1)}_{\mathfrak{d}}$ at $\mathbf{d}\circ \mathbf{c}$;
% the second equality unravels the image of the derivor $(\mathbf{d}\circ\mathbf{c})^{\ast}_{\mathfrak{d}}$ at $\mathbf{T}_{\Omega}(\varnothing^{U})$; the third equality follows from the fact that $!^{\mathbf{T}_{\Omega}(\varnothing^{U})_{\psi\circ\varphi}}$ is the unique mapping from $\varnothing^{S}$ to $\mathbf{T}_{\Omega}(\varnothing^{U})_{\psi\circ\varphi}$, thus
% $$
% !^{\mathrm{T}_{\Omega}(\varnothing^{U})_{\psi\circ\varphi}}
% =
% \left(
% !^{\mathrm{T}_{\Omega}(\varnothing^{U})_{\psi}}
% \right)^{\sharp}_{\varphi}
% \circ\, 
% !^{\mathrm{T}_{\Lambda}(\varnothing^{T})_{\varphi}};
% $$
% the fourth equality follows from Proposition~\ref{PPropUnivComp}; the fifth equality recovers the composition in $\mathsf{Alg}_{\mathfrak{d}}$; the sixth equality recovers the image of the derivor $\mathbf{d}^{\ast}_{\mathfrak{d}}$ at $\mathbf{T}_{\Omega}(\varnothing^{U})$ and the image of the derivor $\mathbf{c}^{\ast}_{\mathfrak{d}}$ at $\mathbf{T}_{\Lambda}(\varnothing^{T})$; finally, the last equality recovers the definition of $J^{(-1)}_{\mathfrak{d}}(\mathbf{d})$ and $J^{(-1)}_{\mathfrak{d}}(\mathbf{c})$.
%
% We conclude that Equation~\ref{EPJmuComp} holds, that is, $J^{(-1)}_{\mathfrak{d}}$ preserves compositions.
%
% This completes the proof.
% \end{proof}
%
% We are now in position to prove the aforementioned adjunction.
%
% \begin{proposition}
% The functor $J^{(-1)}_{\mathfrak{d}}$ is left adjoint to the functor $P^{(-1)}_{\mathfrak{d}}$. See the diagram in Figure~\ref{FAdjmu}.
% \end{proposition}
%
% \begin{proof}
% We will provide the remaining items to complete the adjunction. In this regard, see Definition~\ref{DAdj}.
%
%
% \textsf{Unit of the adjunction.}
%
% Note that
% $
% P^{(-1)}_{\mathfrak{d}} \circ J^{(-1)}_{\mathfrak{d}}
% =
% \Id^{\mathsf{Sig}_{\mathfrak{d}}}.
% $ Let $\sig{\Sigma}$ be a many-sorted signature in $\mathsf{Sig}_{\mathfrak{d}}$. We consider $\eta^{\sig{\Sigma}}$ to be $\id^{\sig{\Sigma}}$, the identity derivor at $\sig{\Sigma}$. Thus, $\eta^{(-1)} = (\id^{\sig{\Sigma}})_{\sig{\Sigma} \in \mathrm{Ob}(\mathsf{Sig}_{\mathfrak{d}})}$ is the identity natural transformation from $\Id^{\mathsf{Sig}_{\mathfrak{d}}}$ to $P^{(-1)}_{\mathfrak{d}} \circ J^{(-1)}_{\mathfrak{d}}$.
%
% \textsf{Counit of the adjunction.}
%
% Let $(\sig{\Sigma}, \mathbf{A})$ be an element in $\mathsf{Alg}_{\mathfrak{d}}$. Note that
% $$
% \left(J^{(-1)}_{\mathfrak{d}} \circ P^{(-1)}_{\mathfrak{d}} \right)(\sig{\Sigma}, \mathbf{A})
% =
% J^{(-1)}_{\mathfrak{d}} (\sig{\Sigma})
% =
% (\sig{\Sigma}, \mathbf{T}_{\Sigma}(\varnothing^{S})).
% $$
%
% Therefore, we denote by $\varepsilon^{(\sig{\Sigma}, \mathbf{A})}$ the ordered pair 
% $
% ( \id^{\sig{\Sigma}}, (!^{\mathbf{A}})^{\sharp} )
% $  
% where $\id^{\sig{\Sigma}}$ is the identity derivor at the many-sorted signature $\sig{\Sigma}$ and $!^{\mathbf{A}} = (!^{A_{s}})_{s \in S}$ is the $S$-sorted mapping from $\varnothing^{S}$ to $\mathbf{A}$, whose $s$-th component, for $s \in S$, is $!^{A_{s}}$, the unique mapping from $\varnothing$ to $A_{s}$, and $(!^{\mathbf{A}})^{\sharp}$ is the unique $\Sigma$-homomorphism from $\mathbf{T}_{\Sigma}(\varnothing^{S})$ to $\mathbf{A} = (\id^{\sig{\Sigma}})^{\ast}_{\mathfrak{d}}(\mathbf{A})$ satisfying that $(!^{\mathbf{A}})^{\sharp} \circ \eta^{\varnothing^{S}} =\, !^{\mathbf{A}}$. Note that $\varepsilon^{(\sig{\Sigma}, \mathbf{A})}$ is a morphism from $(\sig{\Sigma}, \mathbf{T}_{\Sigma}(\varnothing^{S}))$ to $(\sig{\Sigma}, \mathbf{A})$ in $\mathsf{Alg}_{\mathfrak{d}}$.
%
% Finally, we consider the family $\varepsilon^{(-1)} = (\varepsilon^{(\sig{\Sigma}, \mathbf{A})})_{(\sig{\Sigma}, \mathbf{A}) \in \mathrm{Ob}(\mathsf{Alg}_{\mathfrak{d}})}$.
%
% \begin{claim}\label{CEpsilonNT}
%     $\varepsilon^{(-1)}$ is a natural transformation from $J^{(-1)}_{\mathfrak{d}} \circ P^{(-1)}_{\mathfrak{d}}$ to $\Id^{\mathsf{Alg}_{\mathfrak{d}}}$.
% \end{claim}
% %We now show that $\varepsilon$ is a natural transformation from $J^{(-1)}_{\mathfrak{d}} \circ P^{(-1)}_{\mathfrak{d}}$ to $\Id^{\mathsf{Alg}_{\mathfrak{d}}}$.
%
% Let $(\mathbf{c}, h)$ be a morphism from $(\sig{\Sigma}, \mathbf{A})$ to $(\sig{\Lambda}, \mathbf{B})$ in $\mathsf{Alg}_{\mathfrak{d}}$. We need to prove that
% \begin{equation}
% \tag{NT}
% \label{EqEpsilonNT}
% (\mathbf{c}, h) \circ \varepsilon^{(\sig{\Sigma}, \mathbf{A})}
% =
% \varepsilon^{(\sig{\Lambda}, \mathbf{B})} \circ \left(\left(J^{(-1)}_{\mathfrak{d}} \circ P^{(-1)}_{\mathfrak{d}}\right) \left( \mathbf{c}, h \right)\right).
% \end{equation}
%
% Before proceeding further, let us note that
% $$
% \left(J^{(-1)}_{\mathfrak{d}} \circ P^{(-1)}_{\mathfrak{d}}\right) \left( \mathbf{c}, h \right)
% =
% J^{(-1)}_{\mathfrak{d}} (\mathbf{c})
% =
% \left(\mathbf{c},\left(!^{\mathbf{c}^{\ast}_{\mathfrak{d}}(\mathbf{T}_{\Lambda}(\varnothing^{T}))}\right)^{\sharp}\right).
% $$
% Moreover, the following chain of equalities holds
% \allowdisplaybreaks
% \begin{align*}
%     (\mathbf{c}, h) \circ \varepsilon^{(\sig{\Sigma}, \mathbf{A})}
%     &=
%     (\mathbf{c}, h) \circ \left( \id^{\sig{\Sigma}}, (!^{\mathbf{A}})^{\sharp} \right)
%     \tag{1}
%     \\
%     &=
%     \left( \mathbf{c} \circ \id^{\sig{\Sigma}}, h \circ (!^{\mathbf{A}})^{\sharp} \right)
%     \tag{2}
%     \\
%     &=
%     \left( \id^{\sig{\Lambda}} \circ \mathbf{c}, h \circ (!^{\mathbf{A}})^{\sharp} \right)
%     \tag{3}
%     \\
%     &=
%     \left( \id^{\sig{\Lambda}} \circ \mathbf{c}, (h \circ !^{\mathbf{A}})^{\sharp} \right)
%     \tag{4}
%     \\
%     &=
%     \left( \id^{\sig{\Lambda}} \circ \mathbf{c}, (!^{\mathbf{c}^{\ast}_{\mathfrak{d}}(\mathbf{B})})^{\sharp} \right)
%     \tag{5}
%     \\
%     &=
%     \left( \id^{\sig{\Lambda}} \circ \mathbf{c}, \left( \left(!^{\mathbf{B}}\right)^{\sharp}_{\varphi} \circ !^{\mathbf{c}_{\mathfrak{d}}^{\ast}(\mathbf{T}_{\Lambda}(\varnothing^{T}))} \right)^{\sharp} \right)
%     \tag{6}
%     \\
%     &=
%     \left( \id^{\sig{\Lambda}} \circ \mathbf{c},  \left(!^{\mathbf{B}}\right)^{\sharp}_{\varphi} \circ \left( !^{\mathbf{c}_{\mathfrak{d}}^{\ast}(\mathbf{T}_{\Lambda}(\varnothing^{T}))} \right)^{\sharp} \right)
%     \tag{7}
%     \\
%     &=
%     \left( \id^{\sig{\Lambda}}, (!^{\mathbf{B}})^{\sharp} \right) \circ \left( \mathbf{c}, (!^{\mathbf{c}_{\mathfrak{d}}^{\ast}(\mathbf{T}_{\Lambda}(\varnothing^{T}))})^{\sharp} \right)
%     \tag{8}
%     \\
%     &=
%     \varepsilon^{(\sig{\Lambda}, \mathbf{B})} \circ \left(\left(J^{(-1)}_{\mathfrak{d}} \circ P^{(-1)}_{\mathfrak{d}}\right) \left( \mathbf{c}, h \right)\right).
%     \tag{9}
% \end{align*}
%
% In the just stated chain of equalities, the first equality unravels the definition of the morphism $\varepsilon^{(\sig{\Sigma}, \mathbf{A})}$;
% %
% the second equality unravels the definition of the composition of morphisms in $\mathsf{Alg}_{\mathfrak{d}}$. In this regard, see Definition~\ref{DGConst};
% %
% the third equality follows from the fact that, according to Proposition~\ref{PSigdCat-app}, $\mathsf{Sig}_{\mathfrak{d}}$ is a category;
% %
% the fourth equality follows from Proposition~\ref{PPropUnivComp};
% %
% the fifth and sixth equalities follows from the fact that the mapping $!^{\mathbf{c}_{\mathfrak{d}}^{\ast}(\mathbf{B})}$ is the unique mapping from $\varnothing^{S}$ to $\mathbf{c}_{\mathfrak{d}}^{\ast}(\mathbf{B})$, thus
% $$
% h \circ\, !^{\mathbf{A}}
% =\,
% !^{\mathbf{c}_{\mathfrak{d}}^{\ast}(\mathbf{B})}
% =
% \left(!^{\mathbf{B}}\right)^{\sharp}_{\varphi} \circ\, !^{\mathbf{c}_{\mathfrak{d}}^{\ast}(\mathbf{T}_{\Lambda}(\varnothing^{T}))};
% $$
% %
% the seventh equality follows from Proposition~\ref{PPropUnivComp};
% %
% the eighth equality recovers the composition of morphisms in $\mathsf{Alg}_{\mathfrak{d}}$;
% %
% finally, the last equality recovers the definition of the morphisms $\varepsilon^{(\sig{\Lambda}, \mathbf{B})}$ and $(J^{(-1)}_{\mathfrak{d}} \circ P^{(-1)}_{\mathfrak{d}}) \left( \mathbf{c}, h \right)$.
%
% Thus, Claim~\ref{CEpsilonNT} follows.
%
% \textsf{Triangular identities.}
%
% We first prove the following triangular  identity 
% \begin{equation*}
%     \left(P^{(-1)}_{\mathfrak{d}} \ast \varepsilon^{(-1)}\right) \circ \left(\eta^{(-1)} \ast P^{(-1)}_{\mathfrak{d}}\right) = \id^{P^{(-1)}_{\mathfrak{d}}}.
%     \tag{TId1}\label{TIdPmu}
% \end{equation*}
%
% Let $(\sig{\Sigma},\mathbf{A})$ be an element in $\mathsf{Alg}_{\mathfrak{d}}$. On the one hand, we have that
% $$
% P^{(-1)}_{\mathfrak{d}} \left(\varepsilon^{(\sig{\Sigma},\mathbf{A})}\right)
% =
% P^{(-1)}_{\mathfrak{d}} \left(\id^{\sig{\Sigma}}, \left( !^{\mathbf{A}} \right)^{\sharp} \right)
% =
% \id^{\sig{\Sigma}}.
% $$
% On the other hand, we have that 
% $
% \eta^{P^{(-1)}_{\mathfrak{d}}(\sig{\Sigma},\mathbf{A})}
% =
% \eta^{\sig{\Sigma}}
% =
% \id^{\sig{\Sigma}}.
% $ 
%
% All in all, we conclude that
% $$
% P^{(-1)}_{\mathfrak{d}} \left(\varepsilon^{(\sig{\Sigma},\mathbf{A})}\right)
% \circ
% \eta^{P^{(-1)}_{\mathfrak{d}}(\sig{\Sigma},\mathbf{A})}
% =
% \id^{\sig{\Sigma}}
% =
% \id^{P^{(-1)}_{\mathfrak{d}}(\sig{\Sigma},\mathbf{A})}.
% $$
%
% This proves the triangular identity~\ref{TIdPmu}.
%
% We now prove the triangular identity
% \begin{equation*}
%     \left(\varepsilon^{(-1)} \ast J^{(-1)}_{\mathfrak{d}}\right) \circ \left(J^{(-1)}_{\mathfrak{d}} \ast \eta^{(-1)}\right) = \id^{J^{(-1)}_{\mathfrak{d}}}.
%     \tag{TId2}\label{TIdJmu}
% \end{equation*}
%
%
% Let $\sig{\Sigma}$ be a many-sorted signature in $\mathsf{Sig}_{\mathfrak{d}}$. On the one hand, we have that
% $$
% \varepsilon^{J^{(-1)}_{\mathfrak{d}}(\sig{\Sigma})}
% =
% \varepsilon^{(\sig{\Sigma}, \mathbf{T}_{\Sigma}(\varnothing^{S}))}
% =
% \left( \id^{\sig{\Sigma}}, \left(!^{\mathbf{T}_{\Sigma}(\varnothing^{S})}\right)^{\sharp} \right)
% =
% \id^{J^{(-1)}_{\mathfrak{d}}(\sig{\Sigma})}.
% $$
%
% On the other hand, we have that
% $
% J^{(-1)}_{\mathfrak{d}} (\eta^{\sig{\Sigma}})
% =
% J^{(-1)}_{\mathfrak{d}} (\id^{\sig{\Sigma}})
% =
% \id^{J^{(-1)}_{\mathfrak{d}} (\sig{\Sigma})}.
% $ 
%
% All in all, we conclude that
% $$
% \left(\varepsilon^{J^{(-1)}_{\mathfrak{d}}(\sig{\Sigma})}\right) \circ \left(J^{(-1)}_{\mathfrak{d}} (\eta^{\sig{\Sigma}})\right)
% =
% \id^{J^{(-1)}_{\mathfrak{d}} (\sig{\Sigma})}.
% $$
%
% This proves the triangular identity~\ref{TIdJmu}.
%
% We conclude that the functor $J^{(-1)}_{\mathfrak{d}}$ is left adjoint to the functor $P^{(-1)}_{\mathfrak{d}}$. 
% \end{proof}
%
% \begin{figure}
%     \centering
%     \begin{tikzpicture}
% [ACliment/.style={-{To [angle'=45, length=5.75pt, width=4pt, round]}},
% RHACliment/.style={right hook-{To [angle'=45, length=5.75pt, width=4pt, round]}, font=\scriptsize}, 
% scale=1
% ]
%
% \node[] (Smu) at (0,0) {\textcolor{white}{$AAAA$}};
% \node[] () at (0,0) {$\mathsf{Sig}_{\mathfrak{d}}$};
% \node[] (Amu) at (4,0) {\textcolor{white}{$AAAA$}};
% \node[] () at (4,0) {$\mathsf{Alg}_{\mathfrak{d}}$};
% \node[] () at (2,0) {\rotatebox{90}{$\vdash$}};
%
% \draw[ACliment]  (Smu.40) to node [above] {$J^{(-1)}_{\mathfrak{d}}$} (Amu.140);
% \draw[ACliment]  (Amu.220) to node [below] {$P^{(-1)}_{\mathfrak{d}}$} (Smu.320);
%
%
% \end{tikzpicture}
%     \caption{Adjunctions at level $-1$.}
%     \label{FAdjmu}
% \end{figure}
%
% \begin{remark}
% 	Following Proposition~\ref{UOp}, the categories $\mathsf{Fix}(\eta^{(-1)})$ and $\mathsf{Fix}(\varepsilon^{(-1)})$ are equivalent. Note that, since $\eta^{(-1)}$ is the identity natural transformation, it follows that $\mathsf{Fix}(\eta^{(-1)})$ is equal to $\mathsf{Sig}_{\mathfrak{d}}$. On the other hand, following the definition of $\varepsilon^{(-1)}$, for every object $(\sig{\Sigma}, \mathbf{A})$ of $\mathsf{Alg}_{\mathfrak{d}}$, with $\sig{\Sigma}=(S,\Sigma)$, the component $\varepsilon^{(\sig{\Sigma}, \mathbf{A})}$  is an isomorphism if and only if $\mathbf{A} \cong \mathbf{T}_{\Sigma}(\varnothing^{S})$. Thus, $\mathsf{Fix}(\varepsilon^{(-1)})$ is the full subcategory of $\mathsf{Alg}_{\mathfrak{d}}$ determined by
% 	$$
% 	\left\{ \left(\sig{\Sigma}, \mathbf{A}\right) \in \mathrm{Ob}(\mathsf{Alg}_{\mathfrak{d}}) \mid \mathbf{A} \cong \mathbf{T}_{\Sigma}(\varnothing^{S}) \right\}.
% 	$$
%     We will denote $\mathsf{Fix}(\varepsilon^{(-1)})$ by $\mathsf{Tw}^{(-1)}_{\mathfrak{d}}$ and we will call its elements \emph{towers at level $-1$}. As a consequence of all the above, we have that categories $\mathsf{Sig}_{\mathfrak{d}}$ and $\mathsf{Tw}^{(-1)}_{\mathfrak{d}}$ are equivalent.
% \end{remark}
%
% \section{Many-sorted partial algebras}\label{S0D}
%
% In this section we gather together the basic facts about many-sorted partial algebras that we will need afterwards. Specifically, we define the notion of many-sorted partial algebra, we introduce the homomorphisms between many-sorted partial algebras and several classes of them, as e.g., the closed and the full homomorphisms. Moreover, we define, in connection with the homomorphisms and their properties, several classes of subobjects and quotient objects of the many-sorted partial algebras. We also consider universal solutions, as e.g., free completions, and, finally, the concept of equation and the varieties of many-sorted partial algebras. 
%
% \subsection{Many-sorted partial algebras}
%
%
% \begin{definition} A \emph{structure of partial} $\Sigma$-\emph{algebra on} an $S$-sorted set $A$ is a family $(F_{\mathbf{s},s})_{(\mathbf{s},s)\in S^{\star}\times S}$, denoted by $F$, where, for $(\mathbf{s},s)\in S^{\star}\times S$, $F_{\mathbf{s},s}$ is a mapping from $\Sigma_{\mathbf{s},s}$ to $\mathrm{Hom}_{\mathrm{p}}(A_{\mathbf{s}},A_{s})$. If $\Dom(F_{\mathbf{s},s}(\sigma)) = A_{\mathbf{s}}$, then $F_{\mathbf{s},s}(\sigma)$ is a \emph{total operation} from $A_{\mathbf{s}}$ to $A_{s}$. If  $\Dom(F_{\mathbf{s},s}(\sigma)) = \varnothing$, then $F_{\mathbf{s},s}(\sigma)$ is called a \emph{discrete operation} or an \emph{empty (partial) operation} from $A_{\mathbf{s}}$ to $A_{s}$. If $(\mathbf{s},s) = (\lambda,s)$, $\sigma\in \Sigma_{\lambda,s}$, and $\Dom(F_{\lambda,s}(\sigma))\neq\varnothing$, then $F_{\lambda,s}(\sigma)$ picks out an element of $A_{s}$. Therefore a partial operation $F_{\lambda,s}(\sigma)$ from $A_{\lambda}$ to $A_{s}$, also called a \emph{partial constant} $F_{\lambda,s}(\sigma)$ from $A_{\lambda}$ to $A_{s}$, is either discrete or distinguishes exactly an element of $A_{s}$.  A \emph{many-sorted partial} $\Sigma$-\emph{algebra} (or, to abbreviate, \emph{partial $\Sigma$-\emph{algebra}}) is a pair $(A,F)$, abbreviated to $\mathbf{A}$, where $A$ is an $S$-sorted set and $F$ a structure of partial $\Sigma$-algebra on $A$. For a pair $(\mathbf{s},s)\in S^{\star}\times S$ and a formal operation $\sigma\in \Sigma_{\mathbf{s},s}$, in order to simplify the notation, the partial operation $F_{\mathbf{s},s}(\sigma)$ from $A_{\mathbf{s}}$ to $A_{s}$ will be written as $F_{\sigma}$. In some cases, to avoid mistakes, we will denote by $F^{\mathbf{A}}$ the structure of partial $\Sigma$-algebra on $A$, and for $(\mathbf{s},s)\in S^{\star}\times S$ and $\Sigma\in \Sigma_{\mathbf{s},s}$, by $F^{\mathbf{A}}_{\sigma}$, or simply by $\sigma^{\mathbf{A}}$, the corresponding partial operation.  
%
%
% A $\Sigma$-\emph{homomorphism} (or, to abbreviate, \emph{homomorphism}) from $\mathbf{A}$ to $\mathbf{B}$, where $\mathbf{B} = (B,G)$, is a triple $(\mathbf{A},f,\mathbf{B})$, abbreviated to $f\colon \mathbf{A}\mor \mathbf{B}$, where $f$ is an $S$-sorted mapping from $A$ to $B$ such that, for every $(\mathbf{s},s)\in S^{\star}\times S$, every $\sigma\in \Sigma_{\mathbf{s},s}$, and every $(a_{j})_{j\in \bb{\mathbf{s}}}\in A_{\mathbf{s}}$, if $(a_{j})_{j\in \bb{\mathbf{s}}}\in\Dom(\sigma^{\mathbf{A}})$, then $(f_{s_{j}}(a_{j}))_{j\in \bb{\mathbf{s}}}\in\Dom(\sigma^{\mathbf{B}})$ and
% $
% f_{s}(\sigma^{\mathbf{A}}((a_{j})_{j\in \bb{\mathbf{s}}})) = \sigma^{\mathbf{B}}((f_{s_{j}}(a_{j})_{j\in \bb{\mathbf{s}}})).
% $
% For a partial $\Sigma$-algebra $\mathbf{A}$ the triple $(\mathbf{A},\id^{A},\mathbf{A})$, where $\id^{A}$ is the identity at $A$, which is a homomorphism from $\mathbf{A}$ to $\mathbf{A}$, will be denoted by $\id^{\mathbf{A}}$ and will be called the \emph{identity} homomorphism at $\mathbf{A}$.
% We will denote by $\mathsf{PAlg}(\Sigma)$ the category of partial $\Sigma$-algebras and homomorphisms  and by $\mathrm{PAlg}(\Sigma)$ the set of objects  of $\mathsf{PAlg}(\Sigma)$.
% \end{definition}
%
%
%
% \begin{remark}\label{RDiscAdj}
% The forgetful functor $G_{\mathsf{PAlg}(\Sigma)}$ from the category $\mathsf{PAlg}(\Sigma)$ to the category $\mathsf{Set}^{S}$ has a left adjoint, denoted by $\mathbf{D}_{\Sigma}$, which sends an $S$-sorted set $A$ to $\mathbf{D}_{\Sigma}(A)$, the \emph{discrete} many-sorted partial $\Sigma$-algebra on $A$, defined as follows: Its underlying $S$-sorted set is $A$, and, for every $(\mathbf{s},s)\in S^{\star}\times S$ and every $\sigma\in \Sigma_{\mathbf{s},s}$, $\Dom(\sigma^{\mathbf{D}_{\Sigma}(A)}) = \varnothing$. To verify that 
% $\mathbf{D}_{\Sigma}\dashv G_{\mathsf{PAlg}(\Sigma)}$ it suffices to note that, for every many-sorted partial 
% $\Sigma$-algebra $\mathbf{B}$, it happens that every $S$-sorted mapping $f$ from $X$ to the underlying $S$-sorted set of $\mathbf{B}$ is a homomorphism from $\mathbf{D}_{\Sigma}(X)$ to $\mathbf{B}$.  
% \end{remark}
%
% \begin{remark}
% The category $\mathsf{Alg}(\Sigma)$ is a full subcategory of the category $\mathsf{PAlg}(\Sigma)$. We shall prove later that the canonical full embedding of $\mathsf{Alg}(\Sigma)$ into $\mathsf{PAlg}(\Sigma)$ has a left adjoint, the free completion of a partial $\Sigma$-algebra, which will play a key role in this paper. 
% \end{remark}
%
% Since in this work some of the fundamental results consists in the proof that certain homomorphisms between many-sorted partial algebras are isomorphisms, we next state a useful characterization of them. 
% (which show some similarity with the homeomorphisms, i.e., bijective and bicontinuous applications, between topological spaces). Further characterisations of the isomorphisms using the notions of closed homomorphism and complete homomorphism will be included below.
%
%
%
% \begin{proposition}\label{PPAlgIso}
% Let $f$ be a homomorphism from $\mathbf{A}$ to $\mathbf{B}$. Then the following statements are equivalent:
% \begin{enumerate}
% \item $f$ is an isomorphism from $\mathbf{A}$ to $\mathbf{B}$, i.e., $f$ is a section and a retraction.
% \item $f$ is bijective and $(\mathbf{B},f^{-1},\mathbf{A})$, abbreviated to $f^{-1}\colon \mathbf{B}\mor \mathbf{A}$, is also a homomorphism from $\mathbf{B}$ to $\mathbf{A}$.
% \end{enumerate}
% \end{proposition}
%
% \begin{remark}
% Let $\mathbf{A}$ be a partial $\Sigma$-algebra. Then $(\mathbf{D}_{\Sigma}(A),\id^{A},\mathbf{A})$ is a bijective homomorphism, but it is an isomorphism if, and only if, $\mathbf{A}$ is also discrete.
% \end{remark}
%
% To establish some of the results that follow, it is useful to have the concept of product of a family of many-sorted partial algebras.
%
% \begin{definition}
% Let $(\mathbf{A}^{i})_{i\in I}$ be a family of partial $\Sigma$-algebras. Then
% \begin{enumerate}
% \item The \emph{product}
%       of $(\mathbf{A}^{i})_{i\in I}$, $\prod_{i\in I}\mathbf{A}^{i}$, is
%       the partial $\Sigma$-algebra which has as $S$-sorted underlying set $\prod_{i\in I}A^{i}$, 
%       and where, for every $(\mathbf{s},s)\in S^{\star}\times S$ and every $\sigma\in \Sigma_{\mathbf{s},s}$,
%       the structural operation $\sigma^{\prod_{i\in I}\mathbf{A}^{i}}$ is defined as follows:
%       \begin{enumerate}
%       \item $\Dom(\sigma^{\prod_{i\in I}\mathbf{A}^{i}}) = \{(a_{j})_{j\in \bb{\mathbf{s}}}\in 
%       (\prod_{i\in I}A^{i})_{\mathbf{s}}\mid \forall\,i\in I\, 
%       ((a_{j}(i))_{j\in \bb{\mathbf{s}}}\in \Dom(\sigma^{\mathbf{A}^{i}}))\}$.
%       \item
%       $$
%       \sigma^{\prod_{i\in I}\mathbf{A}^{i}}
%       \nfunction
%       {\Dom(\sigma^{\prod_{i\in I}\mathbf{A}^{i}})}
%       {\prod_{i\in I}A^{i}_{s}}
%       {(a_{j})_{j\in \bb{\mathbf{s}}}}
%       {(\sigma^{\mathbf{A}^{i}}((a_{j}(i))_{j\in \bb{\mathbf{s}}})_{i\in I}}
%       $$
%       \end{enumerate}
% \item For every $i\in I$, the $i$-th \emph{canonical projection}
%       is the homomorphism from $\prod_{i\in I}\mathbf{A}^{i}$ to $\mathbf{A}^{i}$
%       determined by the $S$-sorted mapping $\pr^{i}$ which, for every $s\in S$, is
%       defined as follows
%       $$
%       \pr^{i}_{s}
%       \nfunction
%       {\prod_{i\in I}A^{i}_{s}}
%       {A^{i}_{s}}
%       {(a_{i})_{ i\in I}}
%       {a_{i}}
%       $$
% \end{enumerate}
% \end{definition}
%
% \begin{proposition}\label{PPAlgProd}
% Let $(\mathbf{A}^{i})_{i\in I}$ be a family of partial $\Sigma$-algebras.
% Then the pair $(\prod_{i\in I}\mathbf{A}^{i},(\pr^{i})_{i\in I})$
% is a product in $\mathsf{PAlg}(\Sigma)$.
% \end{proposition}
%
% \subsection{Subalgebras}
% In what follows we define the notions of subalgebra, relative subalgebra, and weak subalgebra of a partial $\Sigma$-algebra $\mathbf{A}$ and the subalgebra generating operator for $\mathbf{A}$. Furthermore, we characterize the subalgebras, the relative subalgebras and the weak subalgebras by means of the closed homomorphisms, the full homomorphisms and the homomorphisms, respectively.
%
% \begin{definition}\label{PPAlgSub}
% Let $\mathbf{A}$ be a partial $\Sigma$-algebra and $X\subseteq A$. Given $(\mathbf{s},s)\in S^{\star}\times S$ and $\sigma\in\Sigma_{\mathbf{s},s}$, we will say that $X$ is \emph{closed under the partial operation} $\sigma^{\mathbf{A}}\colon A_{\mathbf{s}}\dmor A_{s}$ if, for every $(a_{j})_{j\in\bb{\mathbf{s}}}\in X_{\mathbf{s}}$, if $(a_{j})_{j\in\bb{\mathbf{s}}}\in\Dom(\sigma^{\mathbf{A}})$, then $\sigma^{\mathbf{A}}((a_{j})_{j\in\bb{\mathbf{s}}})\in X_{s}$. We will say that $X$ is a \emph{closed subset} of $\mathbf{A}$ if $X$ is closed under the partial operations of $\mathbf{A}$. We will denote by $\mathrm{Cl}(\mathbf{A})$ the set of all closed subsets of $\mathbf{A}$ (which is an algebraic closure system on $A$) and by  $\mathbf{Cl}(\mathbf{A})$ the algebraic lattice $(\mathrm{Cl}(\mathbf{A}),\subseteq)$.
%
% We will say that a partial $\Sigma$-algebra $\mathbf{B} $ is a \emph{subalgebra} of $\mathbf{A}$ if $B$ is a closed subset of $\mathbf{A}$, for every $(\mathbf{s},s)\in S^{\star}\times S$ and every $\sigma\in\Sigma_{\mathbf{s},s}$, $\Dom(\sigma^{\mathbf{B}}) = \Dom(\sigma^{\mathbf{A}})\cap B_{\mathbf{s}}$ and, for every  $(b_{j})_{j\in\bb{\mathbf{s}}}\in \Dom(\sigma^{\mathbf{B}})$, $\sigma^{\mathbf{B}}((b_{j})_{j\in\bb{\mathbf{s}}}) = \sigma^{\mathbf{A}}((b_{j})_{j\in\bb{\mathbf{s}}})$. We will denote by $\mathrm{Sub}(\mathbf{A})$ the set of all subalgebras of $\mathbf{A}$ and by $\mathbf{Sub}(\mathbf{A})$ the ordered set $(\mathrm{Sub}(\mathbf{A}),\subseteq)$. Since $\mathbf{Cl}(\mathbf{A})$ and $\mathbf{Sub}(\mathbf{A})$ are isomorphic, we shall feel free to deal either with a closed subset of $\mathbf{A}$ or with the correlated subalgebra of $\mathbf{A}$, whichever is most convenient for the work at hand.
% \end{definition}
%
%
% \begin{remark}
% Let $\mathbf{A}$ and $\mathbf{B}$ be partial $\Sigma$-algebras such that $B\subseteq A$. Then $\mathbf{B} $ is a subalgebra of $\mathbf{A}$ if and only if, for every $(\mathbf{s},s)\in S^{\star}\times S$ and every $\sigma\in\Sigma_{\mathbf{s},s}$, $\Gamma_{\!\sigma^{\mathbf{B}}} = \Gamma_{\!\sigma^{\mathbf{A}}}\cap (B_{\mathbf{s}}\times A_{s})$.
% \end{remark}
%
%
% \begin{remark}
% Let $\mathbf{A}$ be a partial $\Sigma$-algebra. If, for every $s\in S$ and every $\sigma\in \Sigma_{\lambda,s}$, $\Dom(\sigma^{\mathbf{A}}) = \varnothing$, then $\varnothing^{S}\in \mathrm{Cl}(\mathbf{A})$. Let us note that, in such a case, $\boldsymbol{\varnothing}^{S}$, the subalgebra of $\mathbf{A}$ canonically associated to $\varnothing^{S}$, is such that, for every $(\mathbf{s},s)\in S^{\star}\times S$ and every $\sigma\in\Sigma_{\mathbf{s},s}$, the partial operation from $\varnothing^{S}_{\mathbf{s}}$ to $\varnothing^{S}_{s} = \varnothing$ associated to $\sigma$ is discrete.
% \end{remark}
%
% We next define the notion of closed homomorphism from a partial $\Sigma$-algebra to another which will allow us, among other things, to provide another characterization of the isomorphisms and of the subalgebras of a partial $\Sigma$-algebra.
%
% \begin{definition}
% Let $\mathbf{A}$ and $\mathbf{B}$ be partial $\Sigma$-algebras and $f$ a homomorphism from $\mathbf{A}$ to $\mathbf{B}$. We will say that $f$ is \emph{closed} if, for every $(\mathbf{s},s)\in S^{\star}\times S$, every $\sigma\in\Sigma_{\mathbf{s},s}$, and every $(a_{j})_{j\in\bb{\mathbf{s}}}\in A_{\mathbf{s}}$, if $(f_{s_{j}}(a_{j})_{j\in\bb{\mathbf{s}}})\in \Dom(\sigma^{\mathbf{B}})$, then $(a_{j})_{j\in\bb{\mathbf{s}}}\in \Dom(\sigma^{\mathbf{A}})$, i.e., $\Dom(\sigma^{\mathbf{A}}) \supseteq f_{\mathbf{s}}^{-1}[\Dom(\sigma^{\mathbf{B}})]$ (therefore $\Dom(\sigma^{\mathbf{A}}) = f_{\mathbf{s}}^{-1}[\Dom(\sigma^{\mathbf{B}})]$).
% \end{definition}
%
% \begin{remark}
% A homomorphism is closed if and only if the domains of the partial operations in the source are exactly the inverse images of the domains of the corresponding operations in the target.
% \end{remark}
%
% \begin{remark}
% Every homomorphism between $\Sigma$-algebras is a closed homomorphism.
% \end{remark}
%
% \begin{proposition}
% Let $f$ be a homomorphism from $\mathbf{A}$ to $\mathbf{B}$. Then the following statements are equivalent:
% \begin{enumerate}
% \item $f$ is an isomorphism from $\mathbf{A}$ to $\mathbf{B}$, i.e., $f$ is a section and a retraction.
% \item $f$ is bijective and closed, i.e., for every $(\mathbf{s},s)\in S^{\star}\times S$, every $\sigma\in \Sigma_{\mathbf{s},s}$, and every $(a_{j})_{j\in \bb{\mathbf{s}}}\in A_{\mathbf{s}}$, the following holds: $(a_{j})_{j\in \bb{\mathbf{s}}}\in \Dom(\sigma^{\mathbf{A}})$ if, and only if, $ (f_{s_{j}}(a_{j}))_{j\in \bb{\mathbf{s}}}\in \Dom(\sigma^{\mathbf{B}})$.
% \end{enumerate}
% \end{proposition}
%
% \begin{proposition}
% Let $\mathbf{A}$ and $\mathbf{B}$ be partial $\Sigma$-algebras, $f$ a homomorphism from $\mathbf{A}$ to $\mathbf{B}$, and $X$ and $Y$ closed subsets of $\mathbf{A}$ and $\mathbf{B}$, respectively. Then:
% \begin{enumerate}
% \item $f^{-1}[Y]$ is a closed subset of $\mathbf{A}$.
% \item If $f$ is closed, then $f[X]$ is a closed subset of $\mathbf{B}$.
% \end{enumerate}
% \end{proposition}
%
% \begin{remark}
% For a homomorphism $f$ from a partial $\Sigma$-algebra $\mathbf{A}$ to another $\mathbf{B}$ and a subset $X$ of $A$, we have that $f[\mathrm{Sg}_{\mathbf{A}}(X)] \subseteq \mathrm{Sg}_{\mathbf{B}}(f[X])$ (the converse inclusion does not holds in general; however, if $f$ is closed, then $f[\mathrm{Sg}_{\mathbf{A}}(X)] = \mathrm{Sg}_{\mathbf{B}}(f[X])$.
% \end{remark}
%
% We now characterize the subalgebras of a partial $\Sigma$-algebra by means of the closed homomorphisms.
%
% \begin{proposition}
% Let $\mathbf{A}$ and $\mathbf{B}$ be partial $\Sigma$-algebras such that $B\subseteq A$. Then $\mathbf{B}$ is a subalgebra of $\mathbf{A}$ if and only if $\mathrm{in}^{\mathbf{B},\mathbf{A}} = (\mathbf{B},\mathrm{in}^{B,A},\mathbf{A})$ is a closed homomorphism from $\mathbf{B}$ to $\mathbf{A}$.
% \end{proposition}
%
% \begin{definition}
% Let $\mathbf{A}$ be a partial $\Sigma$-algebra. Then we denote by $\mathrm{Sg}_{\mathbf{A}}$ the algebraic closure operator canonically associated to the algebraic closure system $\mathrm{Cl}(\mathbf{A})$ on $A$ and we call it the \emph{subalgebra  generating operator for} $\mathbf{A}$. Moreover, if $X\subseteq A$, then we call $\mathrm{Sg}_{\mathbf{A}}(X)$ the \emph{subalgebra of} $\mathbf{A}$ \emph{generated by} $X$, and if $X$ is such that $\mathrm{Sg}_{\mathbf{A}}(X) = A$, then we will say that $X$ is a \emph{generating subset of} $\mathbf{A}$. Besides, $\mathbf{Sg}_{\mathbf{A}}(X)$ denotes the partial  $\Sigma$-algebra determined by $\mathrm{Sg}_{\mathbf{A}}(X)$.
% \end{definition}
%
% \begin{remark}
% Let $\mathbf{A}$ be a partial $\Sigma$-algebra. Then the algebraic closure operator $\mathrm{Sg}_{\mathbf{A}}$ is uniform, i.e., for every  $X$, $Y\subseteq A$, if $\mathrm{supp}_{S}(X) = \mathrm{supp}_{S}(Y)$, then we have $\mathrm{supp}_{S}(\mathrm{Sg}_{\mathbf{A}}(X)) = \mathrm{supp}_{S}(\mathrm{Sg}_{\mathbf{A}}(Y))$.
% \end{remark}
%
% We next recall the proof by algebraic induction for partial $\Sigma$-algebras.
%
% \begin{proposition}[Principle of Proof by Algebraic Induction]\label{PPAII}
% Let $\mathbf{A}$ be a partial $\Sigma$-algebra generated by $X$. Then to prove that a subset $Y$ of $A$ is equal to $A$ it suffices to show: (1) $X\subseteq Y$ (algebraic induction basis) and (2) $Y$ is a closed subset subalgebra of $\mathbf{A}$ (algebraic induction step).
% \end{proposition}
%
% We next state the principle of extension of identities. This principle, which is fundamental to elucidate the equality of two coterminal homomorphisms, wil be used on several occasions in this work.
%
% \begin{proposition}
% Let $f,g\colon \mathbf{A}\mor \mathbf{B}$ be homomorphisms between partial $\Sigma$-algebras and let $X$ be a subset of $\mathbf{A}$. If $f\!\!\upharpoonright_{X} = g\!\!\upharpoonright_{X}$, then $f\!\!\upharpoonright_{\mathbf{Sg}_{\mathbf{A}}(X)} = g\!\!\upharpoonright_{\mathbf{Sg}_{\mathbf{A}}(X)}$. In particular, if $X$ is a generating subset of $\mathbf{A}$ and $f\!\!\upharpoonright_{X} = g\!\!\upharpoonright_{X}$, then $f = g$.
% \end{proposition}
%
% \begin{definition}\label{PPAlgSubRel}
% Let $\mathbf{A}$ be a partial $\Sigma$-algebra. We will say that a partial $\Sigma$-algebra $\mathbf{B} $ is a \emph{relative subalgebra} of $\mathbf{A}$ if $B\subseteq A$ and, for every $(\mathbf{s},s)\in S^{\star}\times S$ and every $\sigma\in\Sigma_{\mathbf{s},s}$, $\Dom(\sigma^{\mathbf{B}}) \subseteq \Dom(\sigma^{\mathbf{A}})$ and, for every  $(b_{j})_{j\in\bb{\mathbf{s}}}\in B_{\mathbf{s}}$, if $(b_{j})_{j\in\bb{\mathbf{s}}}\in\Dom(\sigma^{\mathbf{A}})$ and $\sigma^{\mathbf{A}}((b_{j})_{j\in\bb{\mathbf{s}}})\in B_{s}$, then $(b_{j})_{j\in\bb{\mathbf{s}}}\in\Dom(\sigma^{\mathbf{B}})$ and $\sigma^{\mathbf{B}}((b_{j})_{j\in\bb{\mathbf{s}}}) = \sigma^{\mathbf{A}}((b_{j})_{j\in\bb{\mathbf{s}}})$. We will denote by $\mathrm{Sub}_{\mathrm{r}}(\mathbf{A})$ the set of all relative subalgebras of $\mathbf{A}$.
% \end{definition}
%
%
% \begin{remark}
% Let $\mathbf{A}$ and $\mathbf{B}$ be partial $\Sigma$-algebras such that $B\subseteq A$. Then $\mathbf{B} $ is a relative subalgebra of $\mathbf{A}$ if and only if, for every $(\mathbf{s},s)\in S^{\star}\times S$ and every $\sigma\in\Sigma_{\mathbf{s},s}$, $\Gamma_{\!\sigma^{\mathbf{B}}} = \Gamma_{\!\sigma^{\mathbf{A}}}\cap (B_{\mathbf{s}}\times B_{s})$.
% \end{remark}
%
% \begin{remark}\label{PPAlgLift}
% Let $\mathbf{A} = (A,F)$ be a partial $\Sigma$-algebra and $X\subseteq A$. Then there exists a unique structure of partial $\Sigma$-algebra $\mathrm{L}^{\mathrm{in}^{X,A}}(F)$ on $X$, the \emph{optimal lift of} $F$ \emph{with respect to } $\mathrm{in}^{X,A}$, such that, for every partial $\Sigma$-algebra $\mathbf{C}$ and every $S$-sorted mapping $h$ from $C$ to $X$, if $\mathrm{in}^{X,A}\circ h$ is a homomorphism from $\mathbf{C}$ to $\mathbf{A}$, then $h$ is a homomorphism from $\mathbf{C}$ to $(X,\mathrm{L}^{\mathrm{in}^{X,A}}(F))$. In fact, it suffices to take, for every $(\mathbf{s},s)\in S^{\star}\times S$ and every $\sigma\in \Sigma_{\mathbf{s},s}$, as partial operation associated to $\sigma$ the partial mapping $\mathrm{L}^{\mathrm{in}^{X,A}}(F)_{\sigma}$ from $X_{\mathbf{s}}$ to $X_{s}$ defined as follows:
% $$
% \Dom(\mathrm{L}^{\mathrm{in}^{X,A}}(F)_{\sigma}) = \{(x_{j})_{j\in\bb{\mathbf{s}}}\in X_{\mathbf{s}}\mid (x_{j})_{j\in\bb{\mathbf{s}}}\in \Dom(\sigma^{\mathbf{A}}) \And \sigma^{\mathbf{A}}((x_{j})_{j\in\bb{\mathbf{s}}})\in X_{s}\}
% $$
% and, for every $(x_{j})_{j\in\bb{\mathbf{s}}}\in \Dom(\mathrm{L}^{\mathrm{in}^{X,A}}(F)_{\sigma})$, $\mathrm{L}^{\mathrm{in}^{X,A}}(F)_{\sigma}((x_{j})_{j\in\bb{\mathbf{s}}}) = \sigma^{\mathbf{A}}((x_{j})_{j\in\bb{\mathbf{s}}})$. This shows that every subset $X$ of $A$ is the underlying $S$-sorted set of exactly one relative subalgebra of $\mathbf{A}$, and that the structure of partial $\Sigma$-algebra on $X$ is uniquely determined by the structure of partial $\Sigma$-algebra on $A$ and the specification of the subset $X$ of $A$.
%
% Let us note that if the partial $\Sigma$-algebra $\mathbf{A}$ is \emph{total}, then, for every $(\mathbf{s},s)\in S^{\star}\times S$ and every $\sigma\in \Sigma_{\mathbf{s},s}$, we have that
% $$
% \Dom(\mathrm{L}^{\mathrm{in}^{X,A}}(F)_{\sigma}) = \{(x_{j})_{j\in\bb{\mathbf{s}}}\in X_{\mathbf{s}}\mid \sigma^{\mathbf{A}}((x_{j})_{j\in\bb{\mathbf{s}}})\in X_{s}\} = 
% X_{\mathbf{s}}\cap (\sigma^{\mathbf{A}})^{-1}[X_{s}].
% $$
% Therefore, since $X_{\mathbf{s}}\cap (\sigma^{\mathbf{A}})^{-1}[X_{s}]\subseteq X_{\mathbf{s}}$ (but not necessarily $X_{\mathbf{s}}\subseteq (\sigma^{\mathbf{A}})^{-1}[X_{s}]$), we have that $(X,\mathrm{L}^{\mathrm{in}^{X,A}}(F))$ will, in general, be a \emph{partial} $\Sigma$-algebra.
% \end{remark}
%
% We next define the notion of full homomorphism from a partial $\Sigma$-algebra to another which will allow us, among other things, to provide another characterization of the isomorphisms and of the relative subalgebras of a $\Sigma$-algebra.
%
% \begin{definition}
% Let $\mathbf{A}$ and $\mathbf{B}$ be partial $\Sigma$-algebras and $f$ a homomorphism from $\mathbf{A}$ to $\mathbf{B}$. We will say that $f$ is \emph{full} if, for every $(\mathbf{s},s)\in S^{\star}\times S$, every $\sigma\in\Sigma_{\mathbf{s},s}$, and every $(a_{j})_{j\in\bb{\mathbf{s}}}\in A_{\mathbf{s}}$, if $((f_{s_{j}}(a_{j}))_{j\in\bb{\mathbf{s}}})\in \Dom(\sigma^{\mathbf{B}})$ and $\sigma^{\mathbf{B}}((f_{s_{j}}(a_{j}))_{j\in\bb{\mathbf{s}}})\in f_{s}[A_{s}]$, then there exists an $(a'_{j})_{j\in\bb{\mathbf{s}}}\in \Dom(\sigma^{\mathbf{A}})$ such that $((f_{s_{j}}(a'_{j}))_{j\in\bb{\mathbf{s}}}) = ((f_{s_{j}}(a_{j}))_{j\in\bb{\mathbf{s}}})$.
% \end{definition}
%
% \begin{remark}\label{PPAlgFullIni}
% Those homomorphisms that are full and injective are characterizable as those that are initial and injective, where a homomorphism $f$ from a partial $\Sigma$-algebra $\mathbf{A}$ to another $\mathbf{B}$ is said to be \emph{initial} if, for every partial $\Sigma$-algebra $\mathbf{C}$ and every S-sorted mapping $h$ from $C$ to $A$, if $f\circ h$ is a homomorphism from $\mathbf{C}$ to $\mathbf{B}$, then $h$ is a homomorphism from $\mathbf{C}$ to $\mathbf{A}$. 
% \end{remark}
%
% \begin{proposition}
% Let $f$ be a homomorphism from $\mathbf{A}$ to $\mathbf{B}$. Then the following statements are equivalent:
% \begin{enumerate}
% \item $f$ is an isomorphism from $\mathbf{A}$ to $\mathbf{B}$, i.e., $f$ is a section and a retraction.
% \item $f$ is bijective and full.
% \end{enumerate}
% \end{proposition}
%
% \begin{proposition}
% Let $\mathbf{A}$ and $\mathbf{B}$ be partial $\Sigma$-algebras such that $B\subseteq A$. Then $\mathbf{B}$ is a relative subalgebra of $\mathbf{A}$ if and only if $\mathrm{in}^{\mathbf{B},\mathbf{A}} = (\mathbf{B},\mathrm{in}^{B,A},\mathbf{A})$ is a full homomorphism from $\mathbf{B}$ to $\mathbf{A}$.
% \end{proposition}
%
% %Hablar aquí de levantamiento optimal y después, cuando se tengan las congruencias, de levantamiento co-optimal.
% %
% %\begin{proposition}\label{PPAlgLift}
% %Let $\mathbf{A} = (A,F^{\mathbf{A}})$ be a partial $\Sigma$-algebra and $X\subseteq A$. Then there exists a structure of partial $\Sigma$-algebra $\mathrm{L}^{\mathrm{in}^{X,A}}(F^{\mathbf{A}})$ on $X$, the \emph{optimal lift of} $F^{\mathbf{A}}$ \emph{with respect to }$\mathrm{in}^{X,A}$, such that, for every partial $\Sigma$-algebra $\mathbf{C}$ and every mapping $g$ from $C$ to $X$, if $\mathrm{in}^{X,A}\circ g$ is a homomorphism from $\mathbf{C}$ to $\mathbf{A}$, then $g$ is a homomorphism from $\mathbf{C}$ to $(X,\mathrm{L}^{\mathrm{in}^{X,A}}(F^{\mathbf{A}}))$.
% %\end{proposition}
% %
% %\begin{proof}
% %It suffices to take, for every $(w,s)\in S^{\star}\times S$ and every $\sigma\in \Sigma_{w,s}$, as partial operation associated to $\sigma$ the partial mapping $\mathrm{L}^{\mathrm{in}^{X,A}}(F^{\mathbf{A}})_{\sigma}$ from $X_{w}$ to $\X_{s}$ defined as follows:
% %$$
% %\Dom(\mathrm{L}^{\mathrm{in}^{X,A}}(F^{\mathbf{A}})_{\sigma}) = \{x\in X_{w}\mid b\in \Dom(F^{\mathbf{A}}_{\sigma}) \And F^{\mathbf{A}}_{\sigma}(x)\in X_{s}\}
% %$$
% %and, for every $x\in \Dom(\mathrm{L}^{\mathrm{in}^{X,A}}(F^{\mathbf{A}})_{\sigma})$, $\mathrm{L}^{\mathrm{in}^{X,A}}(F^{\mathbf{A}})_{\sigma}(b) = F^{\mathbf{A}}_{\sigma}(x)$.
% %\end{proof}
%
% \begin{definition}\label{PPAlgSubWeak}
% Let $\mathbf{A}$ be a partial $\Sigma$-algebra. We will say that a partial $\Sigma$-algebra $\mathbf{B} $ is a \emph{weak subalgebra} of $\mathbf{A}$ if $B\subseteq A$ and, for every $(\mathbf{s},s)\in S^{\star}\times S$, every $\sigma\in\Sigma_{\mathbf{s},s}$, and every  $(b_{j})_{j\in\bb{\mathbf{s}}}\in B_{\mathbf{s}}$, if $(b_{j})_{j\in\bb{\mathbf{s}}}\in\Dom(\sigma^{\mathbf{B}})$, then $(b_{j})_{j\in\bb{\mathbf{s}}}\in\Dom(\sigma^{\mathbf{A}})$ and $\sigma^{\mathbf{B}}((b_{j})_{j\in\bb{\mathbf{s}}}) = \sigma^{\mathbf{A}}((b_{j})_{j\in\bb{\mathbf{s}}})$. We will denote by $\mathrm{Sub}_{\mathrm{w}}(\mathbf{A})$ the set of all weak subalgebras of $\mathbf{A}$.
% \end{definition}
%
%
% \begin{remark}
% Let $\mathbf{A}$ and $\mathbf{B}$ be partial $\Sigma$-algebras such that $B\subseteq A$. Then $\mathbf{B}$ is a weak subalgebra of $\mathbf{A}$ if and only if, for every $(\mathbf{s},s)\in S^{\star}\times S$ and every $\sigma\in\Sigma_{\mathbf{s},s}$, $\Gamma_{\!\sigma^{\mathbf{B}}} \subseteq \Gamma_{\!\sigma^{\mathbf{A}}}$.
% \end{remark}
%
% \begin{proposition}
% Let $\mathbf{A}$ and $\mathbf{B}$ be partial $\Sigma$-algebras such that $B\subseteq A$. Then $\mathbf{B}$ is a weak subalgebra of $\mathbf{A}$ if and only if $\mathrm{in}^{\mathbf{B},\mathbf{A}} = (\mathbf{B},\mathrm{in}^{B,A},\mathbf{A})$ is a homomorphism from $\mathbf{B}$ to $\mathbf{A}$.
% \end{proposition}
%
% \begin{remark}
% Since every closed homomorphism is a full homomorphism and every full homomorphism is a homomorphism, it follows that every subalgebra is a relative subalgebra and every relative subalgebra is a weak subalgebra.  Moreover, every relative subalgebra on a closed subset is a subalgebra.
% \end{remark}
%
% \begin{proposition}
% Let $\mathbf{A}$ and $\mathbf{B}$ be partial $\Sigma$-algebras and $f\colon A\mor B$ a mapping from $A$ to $B$. Then:
% \begin{enumerate}
% \item If $f\colon \mathbf{A}\mor \mathbf{B}$ is a homomorphism, then $\Gamma\!_{f}$ is a closed subset of $\mathbf{A}\times \mathbf{B}$.
% \item If $\mathbf{B}$ is a $\Sigma$-algebra and $\Gamma\!_{f}$ is a closed subset of $\mathbf{A}\times \mathbf{B}$, then $f$ is a homomorphism from $\mathbf{A}$ to $\mathbf{B}$.
% \end{enumerate}
% \end{proposition}
%
% \begin{remark}
% Notice the (partial) analogy of the just stated proposition and the closed graph theorem in point-set topology.
% \end{remark}
%
% \begin{lemma}
% Let $f\colon \mathbf{A}\mor \mathbf{B}$ be a homomorphism, $X\subseteq A$, $g\colon X\mor B$ such that $\Gamma\!_{g}\subseteq \Gamma\!_{f}$, and $\mathbf{X}$ the relative subalgebra of $\mathbf{A}$ on $X$. Then:
% \begin{enumerate}
% \item $g\colon \mathbf{X}\mor \mathbf{B}$ is a homomorphism.
% \item $\Gamma\!_{g}$ is a closed subset of $\mathbf{A}\times \mathbf{B}$ if and only if $X$ is a closed subset of $\mathbf{A}$.
% \end{enumerate}
% \end{lemma}
%
% \begin{corollary}
% Let $f\colon \mathbf{A}\mor \mathbf{B}$ be a homomorphism and $X\subseteq A$. Then the following statements are equivalent:
% \begin{enumerate}
% \item $\mathrm{Sg}_{\mathbf{A}}(X) = A$.
% \item $\mathrm{Sg}_{\mathbf{A}\times \mathbf{B}}(\Gamma\!_{f\upharpoonright_{X}}) = \Gamma\!_{f}$.
% \end{enumerate}
% \end{corollary}
%
% We state the following lemma because it is used, along with other results, in the proof of the generalized recursion theorem, which we will also call Schmidt's construction.
%
% \begin{lemma}
% Let $\mathbf{A}$, $\mathbf{B}$, $\mathbf{C}$, and $\mathbf{D}$ be partial $\Sigma$-algebras, $\mathbf{C}$  a weak subalgebra of $\mathbf{A}$ such that $\mathrm{Sg}_{\mathbf{A}}(C) = A$, $\mathbf{D}$ a relative subalgebra of $\mathbf{B}$, and $f$ a homomorphism from $\mathbf{C}$ to $\mathbf{D}$ which allows a homomorphic extension $\overline{f}$ from $\mathbf{A}$ to $\mathbf{B}$. Then:
% \begin{enumerate}
% \item $\mathrm{Sg}_{\mathbf{A}\times\mathbf{D}}(\Gamma\!_{f})$ is the underlying function of a
%       homomorphism $f^{\flat}$ from
%       $$\mathbf{Dom}(\mathrm{Sg}_{\mathbf{A}\times\mathbf{D}}(\Gamma\!_{f}))$$ to $\mathbf{D}$
%       (where $\mathbf{Dom}(\mathrm{Sg}_{\mathbf{A}\times\mathbf{D}}(\Gamma\!_{f}))$ is the
%       relative subalgebra of $\mathbf{A}$ whose underlying many-sorted set is
%       $\Dom(\mathrm{Sg}_{\mathbf{A}\times\mathbf{D}}(\Gamma\!_{f}))$) and
%       $\mathbf{Dom}(\mathrm{Sg}_{\mathbf{A}\times\mathbf{D}}(\Gamma\!_{f}))$ is generated by $C$.
% \item If $\mathbf{E}$ is a $C$-generated relative subalgebra of $\mathbf{A}$ and
%       $g\colon \mathbf{E}\mor \mathbf{D}$ is a homomorphic extension of $f$, then
%       $\Gamma\!_{g}\subseteq \Gamma\!_{f^{\flat}}$.
% \item If $\mathbf{A}$ is a $\Sigma$-algebra, then
%       $f^{\flat}$ is a closed homomorphism from
%       $\mathbf{Dom}(\mathrm{Sg}_{\mathbf{A}\times\mathbf{D}}(\Gamma\!_{f}))$ to $\mathbf{D}$;
%       and if $f^{\flat}$ is closed and $g\colon \mathbf{E}\mor \mathbf{D}$ is any closed
%       homomorphic extension of $f$ such that $\mathbf{E}$ is a $C$-generated relative
%       subalgebra of $\mathbf{A}$, then $g = f^{\flat}$.
% \end{enumerate}
% \end{lemma}
%
% \begin{definition}
% Let $\mathbf{A}$ and $\mathbf{B}$ be partial $\Sigma$-algebras and $f$ a homomorphism from $\mathbf{A}$ to $\mathbf{B}$. We will say that $f$ is \emph{dense} if $f[A]$ generates $\mathbf{B}$, i.e., if $\mathrm{Sg}_{\mathbf{B}}(f[A]) = B$.
% \end{definition}
%
% \begin{lemma}
% Let $\mathbf{A}$ and $\mathbf{B}$ be partial $\Sigma$-algebras and $f$ a homomorphism from $\mathbf{A}$ to $\mathbf{B}$. Then $f$ is an epimorphism if and only if $f$ is dense.
% \end{lemma}
%
% \begin{remark}
% If a homomorphism between partial $\Sigma$-algebras is dense and closed, then it is surjective (because the image of a closed homomorphism is a closed subset of the target).
% \end{remark}
%
% \begin{proposition}
% Let $\mathbf{A}$ and $\mathbf{B}$ be partial $\Sigma$-algebras and $f$ a homomorphism from $\mathbf{A}$ to $\mathbf{B}$. Then there exists a unique epimorphism $f^{\mathrm{e}}$ from $\mathbf{A}$ to $\mathbf{Sg}_{\mathbf{B}}(f[A])$ such that $f = \mathrm{in}^{\mathbf{Sg}_{\mathbf{B}}[f[A]]}\circ f^{\mathrm{e}}$, where $\mathrm{in}^{\mathbf{Sg}_{\mathbf{B}}(f[A])}$ is the closed embedding of $\mathbf{Sg}_{\mathbf{B}}(f[A])$ into $\mathbf{B}$. Moreover, (1) $f^{\mathrm{e}}$ is an isomorphism if and only if $f$ is closed and injective; and (2) for every partial $\Sigma$-algebra $\mathbf{C}$, every epimorphism $g$ from $\mathbf{A}$ to $\mathbf{C}$ and every closed injective homomorphism $h$ from $\mathbf{C}$ to $\mathbf{B}$, if $f = h\circ g$, then there exists a unique isomorphism $k$ from $\mathbf{Sg}_{\mathbf{B}}(f[A])$ to $\mathbf{C}$ such that $g = k\circ f^{\mathrm{e}}$ and $h \circ k = \mathrm{in}^{\mathbf{Sg}_{\mathbf{B}}(f[A])}$.
% \end{proposition}
%
% Therefore $(\mathrm{Epimorphisms}, \mathrm{Closed\, and\, injective\, homomorphisms})$ is a factorization system in $\mathsf{PAlg}(\Sigma)$.
%
%
% \subsection{Congruences}
%
% Congruences provide an internal description of the homomorphic images from a many-sorted partial $\Sigma$-algebra and closed congruences provide an internal description of the closed homomorphic images from a many-sorted partial $\Sigma$-algebra. Moreover, closed congruences are fundamental for the description of the varieties of many-sorted partial $\Sigma$-algebras defined by existentially conditional existence equations.
%
% \begin{definition}\label{DPAlgCong}
% Let $\mathbf{A}$ be a partial $\Sigma$-algebra and $\Phi$ an $S$-sorted equivalence on $A$. We will say that $\Phi$ is an $S$-\emph{sorted congruence on} (or, to abbreviate, a \emph{congruence on}) $\mathbf{A}$ if, for every $(\mathbf{s},s)\in (S^{\star}-\{\lambda\})\times S$, every $\sigma\in \Sigma_{\mathbf{s},s}$,
% and every $(a_{j})_{j\in\bb{\mathbf{s}}},(b_{j})_{j\in\bb{\mathbf{s}}}\in A_{\mathbf{s}}$, if, for every $j\in \bb{\mathbf{s}}$, $(a_{j}, b_{j})\in\Phi_{s_{j}}$, $(a_{j})_{j\in \bb{\mathbf{s}}}\in\Dom(\sigma^{\mathbf{A}})$, and $(b_{j})_{j\in \bb{\mathbf{s}}}\in\Dom(\sigma^{\mathbf{A}})$, then $(\sigma^{\mathbf{A}}((a_{j})_{j\in\bb{\mathbf{s}}}), \sigma^{\mathbf{A}}((b_{j})_{j\in\bb{\mathbf{s}}}))\in \Phi_{s}$.
% %we have that
% %$$
% %\frac
% %{\forall i\in \lvert w\rvert\mathrm{, }\,\, (a_{i}, b_{i})\in\Phi_{w_{i}} }
% %{(F_{\sigma}(a), F_{\sigma}(b))\in \Phi_{s}}\cdot
% %$$
% Sometimes, we will abbreviate the expression ``$S$-sorted congruence on $\mathbf{A}$'' to ``congruence on $\mathbf{A}$''. We will denote by $\mathrm{Cgr}(\mathbf{A})$ the set of all $S$-sorted congruences on $\mathbf{A}$, which is an algebraic closure system on $A\times A$, by $\mathrm{Cg}_{\mathbf{A}}$ the corresponding algebraic closure operator, by $\mathbf{Cgr}(\mathbf{A})$ the algebraic lattice $(\mathrm{Cgr}(\mathbf{A}),\subseteq)$, by $\nabla_{\mathbf{A}}$ the greatest element of $\mathbf{Cgr}(\mathbf{A})$, and by $\Delta_{\mathbf{A}}$ the least element of $\mathbf{Cgr}(\mathbf{A})$.
%
% For a congruence $\Phi$ on $\mathbf{A}$, the \emph{quotient partial $\Sigma$-algebra of} $\mathbf{A}$ \emph{by} $\Phi$, denoted by $\mathbf{A}/\Phi$, is the partial $\Sigma$-algebra $(A/\Phi, F^{\mathbf{A}/\Phi})$, where, for every  $(\mathbf{s},s)\in S^{\star}\times S$ and every $\sigma\in \Sigma_{\mathbf{s},s}$, the domain of the partial operation $F_{\sigma}^{\mathbf{A}/\Phi}$, also denoted, to simplify, by $\sigma^{\mathbf{A}/{\Phi}}$, from $(A/\Phi)_{\mathbf{s}}$ to $A_{s}/\Phi_{s}$ is the set
% $$
% \{([a_{j}]_{\Phi_{s_{j}}})_{j\in\bb{\mathbf{s}}}\in (A/\Phi)_{\mathbf{s}}\mid \exists\, (a'_{j})_{j\in\bb{\mathbf{s}}}\in A_{\mathbf{s}}\,\text{such that}\,\forall\, j\in\bb{\mathbf{s}}\,((a_{j},a'_{j})\in\Phi_{s_{j}})\},
% $$
% and if $([a_{j}]_{\Phi_{s_{j}}})_{j\in\bb{\mathbf{s}}}$ in $\Dom(\sigma^{\mathbf{A}/{\Phi}})$, then
% $$
% \sigma^{\mathbf{A}/{\Phi}}(([a_{j}]_{\Phi_{s_{j}}})_{j\in\bb{\mathbf{s}}}) = [\sigma^{\mathbf{A}}((a_{j})_{j\in \bb{\mathbf{s}}})]_{\Phi_{s}},
% $$  %
% %$$
% %F_{\sigma}
% %\nfunction
% %{(A/\Phi)_{w}} {A_{s}/\Phi_{s}}
% %{([a_{i}]_{\Phi_{w_{i}}})_{i\in\lvert w\rvert}}
% %{[F_{\sigma}((a_{i})_{i\in \lvert w\rvert})]_{\Phi_{s}}}
% %$$
% and the \emph{canonical projection} from $\mathbf{A}$ to $\mathbf{A}/\Phi$, denoted by $\mathrm{pr}^{\Phi}\colon \mathbf{A}\mor \mathbf{A}/\Phi$, is the full and surjective homomorphism determined by the projection from $A$ to $A/\Phi$.
% The ordered pair $(\mathbf{A}/\Phi,\mathrm{pr}^{\Phi})$ has the following universal property: $\mathrm{Ker}(\mathrm{pr}^{\Phi})$ is $\Phi$ and, for every $\Sigma$-algebra $\mathbf{B}$ and every homomorphism $f$ from $\mathbf{A}$ to $\mathbf{B}$, if $ \Phi\subseteq\mathrm{Ker}(f)$, then there exists a unique homomorphism $h$ from $\mathbf{A}/\Phi$ to $\mathbf{B}$ such that $h\circ\mathrm{pr}^{\Phi} = f$. In particular, if $\Psi$ is a congruence on $A$ such that $\Phi\subseteq \Psi$, then we will denote by $\mathrm{p}^{\Phi,\Psi}$ the unique homomorphism from $\mathbf{A}/\Phi$ to $\mathbf{A}/\Psi$ such that $\mathrm{p}^{\Phi,\Psi}\circ \mathrm{pr}^{\Phi} = \mathrm{pr}^\Psi$.
%
% We will say that $\Phi$ is an $S$-\emph{sorted closed congruence on} (or, to abbreviate, a \emph{closed congruence on}) $\mathbf{A}$ if $\Phi$ is a congruence on $\mathbf{A}$ and, for every $(\mathbf{s},s)\in (S^{\star}-\{\lambda\})\times S$, every $\sigma\in \Sigma_{\mathbf{s},s}$,
% and every $(a_{j})_{j\in\bb{\mathbf{s}}},(b_{j})_{j\in\bb{\mathbf{s}}}\in A_{\mathbf{s}}$, if, for every $j\in \bb{\mathbf{s}}$, $(a_{j}, b_{j})\in\Phi_{s_{j}}$, and $(a_{j})_{j\in \bb{\mathbf{s}}}\in\Dom(\sigma^{\mathbf{A}})$, then $(b_{j})_{j\in \bb{\mathbf{s}}}\in\Dom(\sigma^{\mathbf{A}})$. We will denote by $\mathrm{Cgr}_{\mathrm{c}}(\mathbf{A})$ the set of all $S$-sorted closed congruences on $\mathbf{A}$. The congruence $\Delta_{\mathbf{A}}$ is a closed congruence on $\mathbf{A}$, while $\nabla_{\mathbf{A}}$ is closed if and only if, for every $(\mathbf{s},s)\in (S^{\star}\times S$ and every $\sigma\in \Sigma_{\mathbf{s},s}$, $\sigma^{\mathbf{A}}$ is either total or discrete. If $\Phi$ and $\Psi$ are closed congruences on $\mathbf{A}$, then $\mathrm{Cg}_{\mathbf{A}}(\Phi\cup\Psi)$, their supremum in $\mathbf{Cgr}(\mathbf{A})$, is a closed congruence on $\mathbf{A}$ and it is the least equivalence relation containing them, i.e., 
% $
% \mathrm{Cg}_{\mathbf{A}}(\Phi\cup\Psi) = \bigcup_{n\in \mathbb{N}}(\Psi\circ\Phi)^{n}
% $ 
% (let us point out that, in general, the supremum of two non-closed congruences is not necessarily the least equivalence relation containing them). The set $\mathrm{Cgr}_{\mathrm{c}}(\mathbf{A})$ is a principal ideal of $\mathbf{Cgr}(\mathbf{A})$ precisely the one determined by $\bigcup_{\Phi\in \mathrm{Cgr}_{\mathrm{c}}(\mathbf{A})}\Phi$, which is the largest closed congruence on $\mathbf{A}$.
% \end{definition}
%
% \begin{remark}
% Let $\mathbf{A}$ be a partial $\Sigma$-algebra and $\Phi$ an $S$-sorted equivalence on $A$. Then $\Phi$ is a congruence on $\mathbf{A}$ if and only if $\Phi$ is a subalgebra of $\mathbf{A}\times \mathbf{A}$.
% \end{remark}
%
% \begin{remark}\label{RPAlgHomFinal}
% Let $\mathbf{A}$ be a partial $\Sigma$-algebra and $\Phi$ an $S$-sorted congruence on $A$, then the full and surjective homomorphism $\mathrm{pr}^{\Phi}$ from $\mathbf{A}$ to $\mathbf{A}/\Phi$ is final and surjective, where a homomorphism $f$ from a partial $\Sigma$-algebra $\mathbf{A}$ to another $\mathbf{B}$ is said to be \emph{final} if, for every partial $\Sigma$-algebra $\mathbf{C}$ and every S-sorted mapping $h$ from $B$ to $C$, if $h\circ f$ is a homomorphism from $\mathbf{A}$ to $\mathbf{C}$, then $h$ is a homomorphism from $\mathbf{B}$ to $\mathbf{C}$. For a homomorphism $f$ from a partial $\Sigma$-algebra $\mathbf{A}$ to another $\mathbf{B}$ the property of being final is very interesting because it allows one to state that an $S$-sorted mapping $h$ from $B$ to $C$, the underlying set of another partial $\Sigma$-algebra $\mathbf{C}$, is a homomorphism if its composition with the subjacent application of the homomorphism $f$ is a homomorphism from $\mathbf{A}$ to $\mathbf{C}$.
% Let us point out that a homomorphism $f$ from a partial $\Sigma$-algebra $\mathbf{A} $ to another $\mathbf{B} $ is final if and only if, for every $(\mathbf{s},s)\in S^{\star}\times S$ and every $\sigma\in \Sigma_{\mathbf{s},s}$, $\Gamma_{\!\sigma^{\mathbf{B}}} = (f_{\mathbf{s}}\times f_{s})[\Gamma_{\!\sigma^{\mathbf{A}}}]$.
% \end{remark}
%
%
% \begin{remark}
% Let $\mathbf{A}$ be a partial $\Sigma$-algebra. If $\Phi$ is an $S$-sorted closed congruence on $A$, then $\mathrm{pr}^{\Phi}\colon \mathbf{A}\mor \mathbf{A}/\Phi$, is a closed and surjective homomorphism.
% \end{remark}
%
% \begin{remark}
% Let $\mathsf{ClfdPAlg}(\Sigma)$ be the category whose objects are the \emph{classified partial $\Sigma$-algebras}, i.e., the ordered pairs $(\mathbf{A},\Phi)$ where $\mathbf{A}$ is a partial $\Sigma$-algebra and $\Phi$ a congruence on $\mathbf{A}$, and in which the set of morphisms from $(\mathbf{A},\Phi)$ to $(\mathbf{B},\Psi)$ is the set of all homomorphisms $f$ from $\mathbf{A}$ to $\mathbf{B}$ such that, for every $s\in S$ and every $(x,y)\in A^{2}_{s}$, if $(x,y)\in \Phi_{s}$, then $(f_{s}(x),f_{s}(y))\in \Psi_{s}$. Let $G$ be the functor from $\mathsf{PAlg}(\Sigma)$ to $\mathsf{ClfdPAlg}(\Sigma)$ whose object mapping sends $\mathbf{A}$ to $(\mathbf{A},\Delta_{\mathbf{A}})$ and whose morphism mapping sends $f\colon \mathbf{A}\mor \mathbf{B}$ to $f\colon (\mathbf{A},\Delta_{\mathbf{A}})\mor (\mathbf{B},\Delta_{\mathbf{B}})$. Then, for every classified partial $\Sigma$-algebra $(\mathbf{A},\Phi)$, there exists a universal mapping from $(\mathbf{A},\Phi)$ to $G$, which is precisely the ordered pair $(\mathbf{A}/\Phi,\mathrm{pr}^{\Phi})$ with $\mathrm{pr}^{\Phi}\colon (\mathbf{A},\Phi)\mor (\mathbf{A}/\Phi,\Delta_{\mathbf{A}/\Phi})$.
% \end{remark}
%
% \begin{remark}\label{RPAlgLiftCo}
% Let $\mathbf{A} = (A,F)$ be a partial $\Sigma$-algebra and $\Phi$ an $S$-sorted equivalence on $A$. Then there exists a unique structure of partial $\Sigma$-algebra $\mathrm{L}^{\mathrm{pr}^{\Phi}}(F)$ on $A/\Phi$, the \emph{co-optimal lift of} $F$ \emph{with respect to} $\mathrm{pr}^{\Phi}$, such that, for every partial $\Sigma$-algebra $\mathbf{C}$ and every mapping $h$ from $A/\Phi$ to $C$, if $h\circ\mathrm{pr}^{\Phi}$ is a homomorphism from $\mathbf{A}$ to, $\mathbf{C}$ then $h$ is a homomorphism from $\mathbf{C}$ to $(A/\Phi,\mathrm{L}^{\mathrm{pr}^{\Phi}}(F))$ if and only if $\Phi$ is an $S$-sorted congruence on $\mathbf{A}$. This shows that, for every $S$-sorted congruence $\Phi$ on $\mathbf{A}$ $A$, $A/\Phi$ is the underlying $S$-sorted set of exactly one $\Sigma$-algebra, and that the structure of partial $\Sigma$-algebra on $A/\Phi$ is uniquely determined by the structure of partial $\Sigma$-algebra on $A$ and the specification of the $S$-sorted congruence $\Phi$ on $A$. What we have just said is, in fact, a particular case of the following: for a partial $\Sigma$-algebra $\mathbf{A} = (A,F)$ and an $S$-sorted mapping $f$ from $A$ to $B$, there exists a co-optimal lift of $F$ with respect to $f$ if and only if $\mathrm{Ker}(f)$ is a congruence on $\mathbf{A}$.
% \end{remark}
%
%
% We next state the weak homomorphis theorem.
%
% \begin{proposition}
% Let $\mathbf{A}$, $\mathbf{B}$, and $\mathbf{C}$ be partial $\Sigma$-algebras, $f$ a full and surjective homomorphism from $\mathbf{A}$ to $\mathbf{B}$, and $g$ a homomorphism from $\mathbf{A}$ to $\mathbf{C}$.
% If $\mathrm{Ker}(g)\supseteq \mathrm{Ker}(f)$, then there exists a unique homomorphism $h$ from $\mathbf{B}$ to $\mathbf{C}$ such that $g = h\circ f$ (obviously, if there exists a homomorphism $h$ from $\mathbf{B}$ to $\mathbf{C}$ such that $g = h\circ f$, then  $\mathrm{Ker}(g)\supseteq \mathrm{Ker}(f)$). Moreover, (1) $h$ is a full homomorphism if and only if $g$ is a full homomorphism; (2) $h$ is an isomorphism if and only if $g$ is a surjective and full homomorphism and $\mathrm{Ker}(g) = \mathrm{Ker}(f)$; (3) if $g$ is closed, then $h$ is closed; (4) if $f$ is closed, then $h$ is closed if and only if $g$ is closed.
% \end{proposition}
%
% From the just stated proposition we obtain the weak isomorphism theorem.
%
% \begin{corollary}
% Let $\mathbf{A}$, $\mathbf{B}$ be partial $\Sigma$-algebras, and $f$ a full and surjective homomorphism from $\mathbf{A}$ to $\mathbf{B}$. Then $\mathbf{B}$ and $\mathbf{A}/\mathrm{Ker}(f)$ are isomorphic.
% \end{corollary}
%
% \begin{proposition}
% Let $\mathbf{A}$ and $\mathbf{B}$ be partial $\Sigma$-algebras and $f$ a homomorphism from $\mathbf{A}$ to $\mathbf{B}$. Then there exists a unique injective homomorphism $f^{\mathrm{m}}$ from $\mathbf{A}/\mathrm{Ker}(f)$ to $\mathbf{B}$ such that $f = f^{\mathrm{m}}\circ \mathrm{pr}^{\mathrm{Ker}(f)}$. Moreover, (1) $f^{\mathrm{m}}$ is an isomorphism if and only if $f$ is full and surjective; and (2) for every partial $\Sigma$-algebra $\mathbf{C}$, every full and surjective homomorphism $g$ from $\mathbf{A}$ to $\mathbf{C}$ and every injective homomorphism $h$ from $\mathbf{C}$ to $\mathbf{B}$, if $f = h\circ g$, then there exists a unique isomorphism $k$ from $\mathbf{C}$ to $\mathbf{A}/\mathrm{Ker}(f)$ such that $k\circ g =  \mathrm{pr}^{\mathrm{Ker}(f)}$ and $f^{\mathrm{m}}\circ k = h$.
% \end{proposition}
%
% Therefore $(\mathrm{Full\, and\, surjective\, homomorphisms}, \mathrm{Monomorphisms})$ is a factorization system in $\mathsf{PAlg}(\Sigma)$.
%
% \subsection{Partial Dedekind-Peano Algebras}
%
% We next define for many-sorted partial algebras the counterpart of the notion of Dedekind-Peano algebra. Moreover, we state the principle of definition by algebraic recursion for free many-sorted $\Sigma$-algebras, i.e., many-sorted DP-algebras, with respect to many-sorted partial $\Sigma$-algebras, which, for many-sorted partial algebras, is more important than the principle of definition by algebraic recursion for partial  Dedekind-Peano algebra with respect to many-sorted algebras.
%
% \begin{definition}\label{DPAlgDP}
% Let $\mathbf{A}$ be a partial $\Sigma$-algebra. We will say that $\mathbf{A}$ is a \emph{partial Dedekind-Peano $\Sigma$-algebra}, abbreviated to $\mathrm{PDP}$-\emph{algebra} when this is unlikely to cause confusion, if the following axioms hold
% \begin{enumerate}
% \item[PDP1.] For every $(\mathbf{s},s)\in S^{\star}\times S$ and every $\sigma\in\Sigma_{\mathbf{s},s}$, $\sigma^{\mathbf{A}}\colon A_{\mathbf{s}}\dmor A_{s}$ is injective, i.e., for every $(a_{j})_{j\in\bb{\mathbf{s}}}$, $(b_{j})_{j\in\bb{\mathbf{s}}}\in \Dom(\sigma^{\mathbf{A}})$, if $\sigma^{\mathbf{A}}((a_{j})_{j\in\bb{\mathbf{s}}}) = \sigma^{\mathbf{A}}((b_{j})_{j\in\bb{\mathbf{s}}})$, then $(a_{j})_{j\in\bb{\mathbf{s}}}=(b_{j})_{j\in\bb{\mathbf{s}}}$.
%
% \item[PDP2.] For every $s\in S $ and every $\sigma, \tau \in\Sigma_{\cdot,s}$, if $\sigma\neq\tau$, then $\mathrm{Im}(\sigma^{\mathbf{A}})\cap \mathrm{Im}(\tau^{\mathbf{A}}) = \varnothing$, i.e., $\sigma^{\mathbf{A}}[\Dom(\sigma^{\mathbf{A}})]\cap \tau^{\mathbf{A}}[\Dom(\tau^{\mathbf{A}})] = \varnothing$.
%
% \item[PDP3.] $\mathrm{Sg}_{\mathbf{A}}(A-(\bigcup_{\sigma\in\Sigma_{\cdot,s}}\mathrm{Im}(\sigma^{\mathbf{A}}))_{s\in S}) = A$.
% \end{enumerate}
% We call the $S$-sorted set $A-(\bigcup_{\sigma\in\Sigma_{\cdot,s}}\mathrm{Im}(\sigma^{\mathbf{A}}))_{s\in S}$
% the \emph{basis of Dedekind-Peano of} $\mathbf{A}$, and we denote it by $\mathrm{B}(\mathbf{A})$.
% \end{definition}
%
% \begin{remark}
% Every partial Dedekind-Peano algebra can be obtained, up to isomorphism, as a weak subalgebra of a Dedekind-Peano algebra. But unlike for many-sorted algebras in a variety, for many-sorted partial algebras what we have is that every many-sorted partial algebra in an $\mathrm{E}$-variety is a \emph{weak} homomorphic image of a free many-sorted partial algebra, yet not necessarily a \emph{full} or \emph{closed} homomorphic image. And in $\mathrm{ECE}$-varieties or $\QE$-varieties the more general concept of universal solution becomes of more importance than the one of freeness. The notions of $\mathrm{E}$-, $\mathrm{ECE}$-, and $\QE$-variety will be defined later on.
% \end{remark}
%
% \begin{remark}
% A discrete partial $\Sigma$-algebra is always a PDP-algebra on its underlying $S$-sorted set. Moreover, every weak subalgebra of a PDP-algebra is a PDP-algebra.
% \end{remark}
%
% We next state the principle of the definition by algebraic recursion for many-sorted PDP-algebras with respect to many-sorted algebras.
%
% \begin{proposition}\label{PPAlgDP}
% Let $\mathbf{A}$ be a PDP-algebra, $\mathrm{B}(\mathbf{A})$ its basis, $\mathbf{B}$ a $\Sigma$-algebra, and $f$ a mapping from $\mathrm{B}(\mathbf{A})$ to $B$. Then there exists a unique homomorphism $f^{\sharp}$ from $\mathbf{A}$ to $\mathbf{B}$ such that $f^{\sharp}\circ \mathrm{in}^{\mathrm{B}(\mathbf{A}),\mathbf{A}} = f$.
% \end{proposition}
%
% \begin{proof}
% It suffices to take as $f^{\sharp}$ precisely $(\mathbf{A},\Gamma\!_{f^{\sharp}},\mathbf{B})$, where $\Gamma\!_{f^{\sharp}}$, the underlying $S$-sorted function of $f^{\sharp}$ is $\mathrm{Sg}_{\mathbf{A}\times \mathbf{B}}(\Gamma\!_{f})$.
% \end{proof}
%
% We next state the principle of definition by algebraic recursion for free many-sorted $\Sigma$-algebras, i.e., many-sorted DP-algebras, with respect to many-sorted partial $\Sigma$-algebras.
%
% \begin{proposition}[Proposition~30, pp.~29--32,~\cite{CVCL23}]\label{PPAlgPartial}
% Let $X$ be an $S$-sorted set, $\mathbf{A}$ a partial $\Sigma$-algebra and $f$ an $S$-sorted mapping from $X$ to $A$. Then there exists a unique homomorphism 
% $$
% f^{\partial}\colon \boldsymbol{\partial}(f)\mor \mathbf{A}
% $$ 
% such that
% \begin{enumerate}
% \item $\eta^{X}[X]\subseteq \partial(f)$, where $\eta^{X}$ is the canonical embedding of $X$ into $\mathrm{T}_{\Sigma}(X)$.
% \item $\boldsymbol{\partial}(f)$ is an $X(\cong\eta^{X}[X])$-generated relative subalgebra of $\mathbf{T}_{\Sigma}(X)$.
% \item $f^{\partial}\circ \mathrm{in}^{X,\boldsymbol{\partial}(f)} = f$, where $\mathrm{in}^{X,\boldsymbol{\partial}(f)}$ is the canonical embedding of $X$ into $\boldsymbol{\partial}(f)$.
% \item $f^{\partial}$ is a closed homomorphism.
% \item $f^{\partial}$ is the largest homomorphic extension of $f$ to an $X$-generated relative subalgebra of $\mathbf{T}_{\Sigma}(X)$ with codomain $\mathbf{A}$.
% \end{enumerate}
% \end{proposition}
% \begin{proof}
% Let $\mathbf{A}^{\infty}$ be the $\Sigma$-algebra defined as follows: The underlying $S$-sorted set of $\mathbf{A}^{\infty}$ is $(A_{s}\cup\{A_{s}\})_{s\in S}$ and, for every $(\mathbf{s},s)\in S^{\star}\times S$ and every $\sigma\in \Sigma_{\mathbf{s},s}$, the operation $\sigma^{\mathbf{A}^{\infty}}\colon A^{\infty}_{\mathbf{s}}\mor
% A^{\infty}_{s}$ associated to $\sigma$ is defined as
% $$
% \sigma^{\mathbf{A}^{\infty}}
% \nfunction
% {A^{\infty}_{\mathbf{s}}}{A^{\infty}_{s}}
% {(a_{j})_{j\in \bb{\mathbf{s}}}}{
% \begin{cases}
% \sigma^{\mathbf{A}}((a_{j})_{j\in \bb{\mathbf{s}}}), & \text{if } (a_{j})_{j\in \bb{\mathbf{s}}}\in \Dom(\sigma^{\mathbf{A}});\\
%   A_{s}, & \text{otherwise.}
% \end{cases}
%    }
% $$
%
% We will call $\mathbf{A}^{\infty}$ the \emph{one-point per sort completion} of the partial $\Sigma$-algebra $\mathbf{A}$. Notice that $\mathbf{A}$ is a relative subalgebra of $\mathbf{A}^{\infty}$, that $\mathbf{A}^{\infty}$ need not be generated by $A$ (since some of the new elements are not necessarily accessible from $A$) and that the canonical embedding $\mathrm{in}^{A,A^{\infty}}$ of $A$ into $A^{\infty}$ is the underlying mapping of a full and injective homomorphism, denoted by $\mathrm{in}^{\mathbf{A},\mathbf{A}^{\infty}}$, from $\mathbf{A}$ to $\mathbf{A}^{\infty}$.
% Then, by the universal property of the free algebra, there exists a unique homomorphism 
% $$
% (\mathrm{in}^{\mathbf{A},\mathbf{A}^{\infty}}\circ f)^{\sharp}\colon\mathbf{T}_{\Sigma}(X)\mor\mathbf{A}^{\infty}
% $$ 
% such that $(\mathrm{in}^{\mathbf{A},\mathbf{A}^{\infty}}\circ f)^{\sharp}\circ\eta^{X} = \mathrm{in}^{\mathbf{A},\mathbf{A}^{\infty}}\circ f$. Let $\partial(f)$ be $((\mathrm{in}^{A,A^{\infty}}\circ f)^{\sharp})^{-1}[A]$ and let $f^{\partial}$ be $(\mathrm{in}^{A,A^{\infty}}\circ f)^{\sharp}\bigr|_{((\mathrm{in}^{A,A^{\infty}}\circ f)^{\sharp})^{-1}[A]}^{A}$, i.e., the birestriction of $(\mathrm{in}^{A,A^{\infty}}\circ f)^{\sharp}$ to $((\mathrm{in}^{A,A^{\infty}}\circ f)^{\sharp})^{-1}[A]$ and $A$. Then $\boldsymbol{\partial}(f)$, the relative subalgebra of $\mathbf{T}_{\Sigma}(X)$ on $\partial(f)$, together with the homomorphism determined by  $f^{\partial}$, which, with the customary abuse of notation, we denote by the same symbol, satisfy the desired conditions.
%
% To better understand the mappings at play, we provide the commutative diagram in Figure~\ref{FPAlgPartial}.
% \end{proof}
%
%
%
% \begin{figure}
% $$\xymatrix@C=60pt@R=40pt{
%     X
%     \ar@/_1pc/[ddr]_{f}
%     \ar@/^1pc/[drr]^{\eta^{X}}
%     \ar[dr]|-{\mathrm{in}^{X,\boldsymbol{\partial}(f)}} \\
%      & \boldsymbol{\partial}(f)
%      \ar[d]_{f^{\partial}}
%      \ar[r]^-{\mathrm{in}^{\boldsymbol{\partial}(f),\mathbf{T}_{\Sigma}(X)}}  &
%      \mathbf{T}_{\Sigma}(X)
%      \ar[d]^{(\mathrm{in}^{\mathbf{A},\mathbf{A}^{\infty}}\circ f)^{\sharp}} \\
%      &
%      \mathbf{A}
%      \ar[r]_{\mathrm{in}^{\mathbf{A},\mathbf{A}^{\infty}}} &
%      \mathbf{A}^{\infty}
%    }
% $$
% \caption{Recursion Theorem for DP-algebras with respect to partial algebras.}
% \label{FPAlgPartial}
% \end{figure}
%
% \begin{remark}
% Let us note that, by the Axiom of Regularity, we have that, for every $s\in S$, $A_{s}\cap\{A_{s}\} = \varnothing$. Moreover, for every $(\mathbf{s},s)\in S^{\star}\times S$ and every $\sigma\in \Sigma_{\mathbf{s},s}$, $\Gamma_{\sigma^{\boldsymbol{\partial}(f)}}$, the underlying partial function of $\sigma^{\boldsymbol{\partial}(f)}$, the partial operation from $\partial(f)_{\mathbf{s}}$ to $\partial(f)_{s}$ associated to $\sigma$, is
% $$
% \Gamma_{\!\sigma^{\mathbf{T}_{\Sigma}(X)}}\cap (\partial(f)_{\mathbf{s}}\times \partial(f)_{s}).
% $$
% On the other hand, $\Gamma\!_{f^{\partial}} = \mathrm{Sg}_{\mathbf{T}_{\Sigma}(X)\times \mathbf{A}}((\eta^{X}\times \id^{A})[\Gamma\!_{f}])$, i.e., the underlying function of $f^{\partial}$, is the subalgebra of the partial $\Sigma$-algebra $\mathbf{T}_{\Sigma}(X)\times \mathbf{A}$ generated by the image under the $S$-sorted mapping $\eta^{X}\times \id^{A}\colon X\times A\mor\mathrm{T}_{\Sigma}(X)\times \mathrm{A}$, of the underlying function of $f$. Finally, we have that
% \begin{enumerate}
% \item For every $s\in S$ and every $x\in X_{s}$, $f^{\partial}_{s}(x) = f_{s}(x)$.
% \item For every $s\in S$ and every $\sigma\in\Sigma_{\lambda,s}$, if $\Dom(\sigma^{\mathbf{A}})\neq\varnothing$, then $\sigma\in \partial(f)_{s}$ and $f^{\partial}_{s}(\sigma^{\boldsymbol{\partial}(f)}) = \sigma^{\mathbf{A}}$.
% \item For every $(\mathbf{s},s)\in (S^{\star}-\{\lambda\})\times S$, every $\sigma\in \Sigma_{\mathbf{s},s}$, and every $(P_{j})_{j\in \bb{\mathbf{s}}}\in \mathrm{T}_{\Sigma}(X)_{\mathbf{s}}$, if, for every $j\in\bb{\mathbf{s}}$, $P_{j}\in \partial(f)_{s_{j}}$ and the fami\-ly $(f^{\partial}_{s_{j}}(P_{j}))_{j\in \bb{\mathbf{s}}}\in \Dom(\sigma^{\mathbf{A}})$, then $\sigma^{\boldsymbol{\partial}(f)}((P_{j})_{j\in \bb{\mathbf{s}}})\in \partial(f)_{s}$ and it holds that $$f^{\partial}_{s}(\sigma^{\boldsymbol{\partial}(f)}((P_{j})_{j\in \bb{\mathbf{s}}})) = \sigma^{\mathbf{A}}((f^{\partial}_{s_{j}}(P_{j})))_{j\in \bb{\mathbf{s}}}).$$
% \end{enumerate}
% \end{remark}
%
% \begin{remark}
% The construction that assigns to a partial $\Sigma$-algebra $\mathbf{A}$ the $\Sigma$-algebra $\mathbf{A}^{\infty}$ is not the object mapping of a functor left adjoint to the inclusion functor $\mathrm{In}_{\mathsf{Alg}(\Sigma)}$ from $\mathsf{Alg}(\Sigma)$ to $\mathsf{PAlg}(\Sigma)$. However, the just mentioned construction is the object mapping of a functor $(\cdot)^{\infty}$ from the category $\mathsf{PAlg}(\Sigma)_{\mathrm{c}}$ of many-sorted partial $\Sigma$-algebras and closed homomorphisms to the category $\mathsf{Alg}(\Sigma)$.
% \end{remark}
%
% \subsection{
% \texorpdfstring
% {The free completion of a many-sorted partial $\Sigma$-algebra}
% {The free completion of a many-sorted partial algebra}
% }
%
% We next prove, as a consequence of a well-known theorem about adjoint functors, that the inclusion functor $\mathrm{In}_{\mathsf{Alg}(\Sigma)}$ from $\mathsf{Alg}(\Sigma)$ to $\mathsf{PAlg}(\Sigma)$ has a left adjoint, the (absolutely) free completion functor. We point out that the free completion of a many-sorted partial $\Sigma$-algebra is one of the most useful tools of the theory of partial algebras. At the end of this subsection, and related to the free completion functor, we state the generalized recursion theorem of Schmidt (see~\cite{sch70}), which we will also call the Schmidt construction. This construction is fundamental in the field of many-sorted partial algebras.
%
%
% \begin{figure}
% $$\xymatrix{
%    \mathsf{A}
%    \ar[r]^{G}
%    \ar[dr]_{U}
%    & \mathsf{B} \ar[d]^{V} \\
%    & \mathsf{C}
%     }
% $$
% \caption{Conditions for the existence of a left adjoint.}
% \label{FPAlgCondAdj}
% \end{figure}
%
% \begin{proposition}
% If the diagram of categories and functors in Figure~\ref{FPAlgCondAdj}
% commutes, and the following conditions are satisfied:
% \begin{enumerate}
% \item[(i)] $\mathsf{A}$ is complete, well-powered and co-well-powered,
% \item[(ii)] $G$ preserves limits,
% \item[(iii)] $U$ has a left-adjoint,
% \item[(iv)] $V$ is faithful,
% \end{enumerate}
% then $G$ has a left adjoint.
% \end{proposition}
%
% \begin{figure}
% $$\xymatrix{
%    \mathsf{Alg}(\Sigma)
%    \ar[r]^{\mathrm{In}_{\mathsf{Alg}(\Sigma)}}
%    \ar[dr]_{G_{\mathsf{Alg}(\Sigma)}}
%    & \mathsf{PAlg}(\Sigma) \ar[d]^{G_{\mathsf{PAlg}(\Sigma)}} \\
%    & \mathsf{Set}^{S}
%     }
% $$
% \caption{Conditions for the existence of the free completion.}
% \label{FFreeComp}
% \end{figure}
%
% \begin{corollary}
% The diagram in Figure~\ref{FFreeComp}
% commutes and
% \begin{enumerate}
% \item[(i)] $\mathsf{Alg}(\Sigma)$ is complete, well-powered and co-well-powered,
% \item[(ii)] $\mathrm{In}_{\mathsf{Alg}(\Sigma)}$ preserves limits,
% \item[(iii)] $G_{\mathsf{Alg}(\Sigma)}$ has a left-adjoint ($\mathbf{T}_{\Sigma}$),
% \item[(iv)] $G_{\mathsf{PAlg}(\Sigma)}$ is faithful.
% \end{enumerate}
% Therefore $\mathrm{In}_{\mathsf{Alg}(\Sigma)}$ has a left adjoint $\mathbf{F}_{\Sigma}$, the \emph{free completion} functor.
% \end{corollary}
%
% We next provide an explicit construction of the free completion of a partial $\Sigma$-algebra.
%
% % La demostración hace uso del principle of the definition by algebraic recursion for PDP-algebras with respect to many-sorted algebras.
%
% \begin{proposition}\label{PFreeComp}
% Let $\mathbf{A}$ be a partial $\Sigma$-algebra. Then there exists a $\Sigma$-algebra $\mathbf{F}_{\Sigma}(\mathbf{A})$, the free completion of $\mathbf{A}$, and a homomorphism $\eta^{\mathbf{A}}$ from $\mathbf{A}$ to $\mathbf{F}_{\Sigma}(\mathbf{A})$ such that, for every $\Sigma$-algebra $\mathbf{B}$ and every homomorphism $f$ from $\mathbf{A}$ to $\mathbf{B}$, there exists a unique homomorphism $f^{\mathrm{fc}}$ from $\mathbf{F}_{\Sigma}(\mathbf{A})$ to $\mathbf{B}$ such that $f = f^{\mathrm{fc}}\circ \eta^{\mathbf{A}}$.
% \end{proposition}
%
% \begin{proof}
% Let $\mathbf{T}_{\Sigma}(A)$ be the free $\Sigma$-algebra on $A$ and, for every $(\mathbf{s},s)\in S^{\star}\times S$ and every $\sigma\in \Sigma_{\mathbf{s},s}$, let $\overline{F}_{\sigma}$ be the mapping from $\mathrm{T}_{\Sigma}(A)_{\mathbf{s}}$ to $\mathrm{T}_{\Sigma}(A)_{s}$ defined as:
% $$
% \overline{F}_{\sigma}
% \nfunction
% {\mathrm{T}_{\Sigma}(A)_{\mathbf{s}}}{\mathrm{T}_{\Sigma}(A)_{s}}
% {(P_{j})_{j\in \bb{\mathbf{s}}}}{
% \begin{cases}
%  F^{\mathbf{A}}_{\sigma}((P_{j})_{j\in \bb{\mathbf{s}}}), & \text{if } (P_{j})_{j\in \bb{\mathbf{s}}}\in \Dom(F^{\mathbf{A}}_{\sigma});\\
%  F^{\mathbf{T}_{\Sigma}(A)}_{\sigma}((P_{j})_{j\in \bb{\mathbf{s}}}), & \text{otherwise.}
% \end{cases}
%    }
% $$
% %Thus $\overline{F}_{\sigma}$ has been obtained  by activating $F^{\mathbf{A}}_{\sigma}$, the partial  operation of $\mathbf{A}$ associated to $\sigma$, when this is possible, and when this is not the case, retrieving a syntactic term.
% Let $\overline{\mathbf{T}}_{\Sigma}(\mathbf{A}) = (\mathrm{T}_{\Sigma}(A),\overline{F})$ be the resulting $\Sigma$-algebra. Then 
% $\mathbf{F}_{\Sigma}(\mathbf{A}) = (\mathrm{F}_{\Sigma}(\mathbf{A}),\overline{F}) = \mathbf{Sg}_{\overline{\mathbf{T}}_{\Sigma}(\mathbf{A})}(A)$, the subalgebra of $\overline{\mathbf{T}}_{\Sigma}(\mathbf{A})$ generated by $A$, together with $\eta^{\mathbf{A}}$, the canonical embedding of $\mathbf{A}$ into $\mathbf{F}_{\Sigma}(\mathbf{A})$ induced by the corestriction of the mapping $\eta^{A}\colon A\mor \mathrm{T}_{\Sigma}(A)$ to the underlying $S$-sorted set of $\mathbf{F}_{\Sigma}(\mathbf{A})$, is a universal morphism from $\mathbf{A}$ to $\mathrm{In}_{\mathbf{Alg}(\Sigma)}$. In fact, from $\mathbf{F}_{\Sigma}(\mathbf{A})$ we obtain a partial Dedekind-Peano $\Sigma$-algebra $\mathbf{F}^{\ast}_{\Sigma}(\mathbf{A}) = (\mathrm{F}_{\Sigma}(\mathbf{A}),F^{\ast})$ by defining, for every $(\mathbf{s},s)\in S^{\star}\times S$ and every $\sigma\in \Sigma_{\mathbf{s},s}$, $F_{\sigma}^{\ast}$ to be the partial mapping from $\mathrm{F}_{\Sigma}(\mathbf{A})_{\mathbf{s}}$ to $\mathrm{F}_{\Sigma}(\mathbf{A})_{s}$ whose domain of definition is $\mathrm{F}_{\Sigma}(\mathbf{A})_{\mathbf{s}}-\Dom(F_{\sigma}^{\mathbf{A}})$ and is such that, 
% for $(P_{j})_{j\in \bb{\mathbf{s}}}\in \Dom(F_{\sigma}^{\ast})$, 
% $F_{\sigma}^{\ast}((P_{j})_{j\in \bb{\mathbf{s}}}) = \overline{F}_{\sigma}((P_{j})_{j\in \bb{\mathbf{s}}})$. Then, given a $\Sigma$-algebra $\mathbf{B}$ and a homomorphism $f$ from $\mathbf{A}$ to $\mathbf{B}$, by Proposition~\ref{PPAlgDP}, 
% there exists a unique homomorphism $f^{\sharp}$ from $\mathbf{F}^{\ast}_{\Sigma}(\mathbf{A})$ to $\mathbf{B}$ such that $f^{\sharp}\circ\eta^{\mathbf{A}} = f$. We recall that 
% $\Gamma\!_{f^{\sharp}} = \mathrm{Sg}_{\mathbf{F}^{\ast}_{\Sigma}(\mathbf{A})\times \mathbf{B}}(\Gamma_{f})$. Since $f^{\sharp}\circ\eta^{\mathbf{A}} = f$ is a homomorphism from $\mathbf{A}$ to $\mathbf{B}$, $f^{\sharp}$ is a homomorphism from $\mathbf{F}_{\Sigma}(\mathbf{A})$ to $\mathbf{B}$. Then it suffices to take as  $f^{\mathsf{fc}}$ the homomorphism $f^{\sharp}$ but considered as a homomorphism from $\mathbf{F}_{\Sigma}(\mathbf{A})$ to $\mathbf{B}$, i.e., $\Gamma\!_{f^{\mathsf{fc}}} = \Gamma\!_{f^{\sharp}}$ and $\mathrm{d}_{0}(f^{\mathsf{fc}}) = \mathbf{F}_{\Sigma}(\mathbf{A})$.
% \end{proof}
%
% \begin{corollary} \label{CFreeAdj}
% The functor $\mathbf{F}_{\Sigma}$ from $\mathsf{PAlg}(\Sigma)$ to $\mathsf{Alg}(\Sigma)$, which sends a partial $\Sigma$-algebra $\mathbf{A}$ to $\mathbf{F}_{\Sigma}(\mathbf{A})$ and a homomorphism $f$ from $\mathbf{A}$ to $\mathbf{B}$ to $f^{@}$, the unique homomorphism $(\eta^{\mathbf{A}}\circ f)^{\mathrm{fc}}$ from $\mathbf{F}_{\Sigma}(\mathbf{A})$ to $\mathbf{F}_{\Sigma}(\mathbf{B})$ such that $(\eta^{\mathbf{B}}\circ f)^{\mathrm{fc}}\circ\eta^{\mathbf{A}} = \eta^{\mathbf{B}}\circ f$, is a left adjoint of the functor $\mathrm{In}_{\mathsf{Alg}(\Sigma)}$ from $\mathsf{Alg}(\Sigma)$ to $\mathsf{PAlg}(\Sigma)$.
% %$$
% %\xymatrix@=5pc{ \mathsf{Alg}(\Sigma)
% %  \ar@<1.5ex>[r]^{\mathrm{G}_{\Sigma}}
% %  \ar@{}[r]|{\uadj}
% % &
% %\mathsf{Set}^{S}
% %  \ar@<1.5ex>[l]^{\mathbf{T}_{\Sigma}}
% %}
% %$$
% We will call the $\Sigma$-homomorphism  $f^{@}$ the \emph{free completion of} $f$.
% \end{corollary}
%
% \begin{remark}
% For every $\Sigma$-algebra $\mathbf{A}$, $(\mathbf{A},\id^{\mathbf{A}})$ is a free completion of 
% $\mathbf{A}$, since it is a universal morphism from $\mathbf{A}$ to 
% $\mathrm{In}_{\mathsf{Alg}(\Sigma)}$. 
% \end{remark}
%
% \begin{remark}
% Since the functor $\mathrm{In}_{\mathsf{Alg}(\Sigma)}$ from $\mathsf{Alg}(\Sigma)$ to $\mathsf{PAlg}(\Sigma)$ is, clearly, full, faithful and injective on objects, we can assert that the functor $\mathbf{F}_{\Sigma}$ from $\mathsf{PAlg}(\Sigma)$ to $\mathsf{Alg}(\Sigma)$ is, not only a left adjoint, but also a left inverse of $\mathrm{In}_{\mathsf{Alg}(\Sigma)}$, i.e., the counit of the adjunction is the identity. We recall that $\eta$, the unit of the adjunction, determines, for every partial $\Sigma$-algebra $\mathbf{A}$ and every $\Sigma$-algebra $\mathbf{B}$ a natural isomorphism $\varphi_{\mathbf{A},\mathbf{B}}$ from 
% $\mathrm{Hom}(\mathbf{F}_{\Sigma}(\mathbf{A}),\mathbf{B})$ to 
% $\mathrm{Hom}(\mathbf{A},\mathrm{In}_{\mathsf{Alg}(\Sigma)}(\mathbf{B}))$ and that $\varepsilon^{\mathbf{A}}$, the value of the counit 
% $\varepsilon$ of the adjunction at a $\Sigma$-algebra $\mathbf{A}$, is 
% $\varphi^{-1}_{\mathrm{In}_{\mathsf{Alg}(\Sigma)}(\mathbf{A}),\mathbf{A}}(\id^{\mathbf{A}})$, which in this case is, precisely, $\id^{\mathbf{A}}$.
% \end{remark}
%
% \begin{remark}
% For every partial $\Sigma$-algebra $\mathbf{A}$, the injective homomorphism $\eta^{\mathbf{A}}$ from $\mathbf{A}$ to $\mathbf{F}_{\Sigma}(\mathbf{A})$, i.e., the value of the unit of the adjuction at $\mathbf{A}$, is, in addition, an epimorphism, hence a bimorphism. Therefore, the isomorphism-closed, full subcategory $\mathsf{Alg}(\Sigma)$ of $\mathsf{PAlg}(\Sigma)$ is monoreflective and epireflective.
% \end{remark}
%
% \begin{figure}
% \begin{center}
% \begin{tikzpicture}
% [ACliment/.style={-{To [angle'=45, length=5.75pt, width=4pt, round]}
% }, scale=0.7]
% %,font=\scriptsize}]
% \node[] (A) 	at (0,-1) 	[] {$\mathbf{A}$};
% \node[] (B) 	at (0,-4) 	[] {$\mathbf{B}$};
% \node[] (TA) 	at (10,2) 	[] {$\mathbf{T}_{\Sigma}(\mathbf{A})$};
% \node[] (TB) 	at (10,-7) 	[] {$\mathbf{T}_{\Sigma}(\mathbf{B})$};
% \node[] (TpA) 	at (7,.5) 	[] 
% {$\overline{\mathbf{T}}_{\Sigma}(\mathbf{A})$};
% \node[] (TpB) 	at (7,-5.5) 	[] 
% {$\overline{\mathbf{T}}_{\Sigma}(\mathbf{B})$};
% \node[] (FA) 	at (4,-1) 	[] {$\mathbf{F}_{\Sigma}(\mathbf{A})$};
% \node[] (FB) 	at (4,-4) 	[] {$\mathbf{F}_{\Sigma}(\mathbf{B})$};
% \draw[ACliment]  (A) 	to node [left]	{$f$} (B);
% \draw[ACliment]  (TA) 	to node [right]	{$f^{@} = (\eta^{\mathbf{B}}\circ f)^{\sharp}$} (TB);
% \draw[ACliment, bend left]  (A) 	to node [above]	{$\eta^{\mathbf{A}}$} (TA);
% \draw[ACliment, bend right]  (B) 	to node [below]	{$\eta^{\mathbf{B}}$} (TB);
% \draw[ACliment, bend left=20]  (A) 	to node [above]	{$\mathrm{in}^{\mathbf{A}}$} (TpA);
% \draw[ACliment, bend right=20]  (B) 	to node [below]	{$\mathrm{in}^{\mathbf{B}}$} (TpB);
% \draw[ACliment]  (TA) 	to node [above left]	{$(\mathrm{in}^{\mathbf{A}})^{\sharp}$} (TpA);
% \draw[ACliment]  (TB) 	to node [below left]	{$(\mathrm{in}^{\mathbf{B}})^{\sharp}$} (TpB);
% \draw[ACliment]  (A) 	to node [below]	{$\eta^{\mathbf{A}}$} (FA);
% \draw[ACliment]  (B) 	to node [above]	{$\eta^{\mathbf{B}}$} (FB);
% \draw[ACliment]  (FA) 	to node [right]	{$f^{@} = (\eta^{\mathbf{B}}\circ f)^{\mathrm{fc}}$} (FB);
% \draw[ACliment]  (FA) 	to node [above left]	{$\mathrm{in}^{\mathbf{F}_{\Sigma}(\mathbf{A})}$} (TpA);
% \draw[ACliment]  (FB) 	to node [below left]	{$\mathrm{in}^{\mathbf{F}_{\Sigma}(\mathbf{B})}$} (TpB);
% \draw[ACliment, bend left]  (TA) 	to node [below right]	{$(\eta^{\mathbf{A}})^{\sharp}$} (FA);
% \draw[ACliment, bend right]  (TB) 	to node [above right]	{$(\eta^{\mathbf{B}})^{\sharp}$} (FB);
% \end{tikzpicture}
% \end{center}
% \caption{The free completion functor.}
% \label{FFreeCompFun}
% \end{figure}
%
% \begin{remark}\label{RFreeDisc}
% Let $X$ be an $S$-sorted set. Then, for $\mathbf{D}_{\Sigma}(X)$, the discrete many-sorted partial $\Sigma$-algebra associated to an $S$-sorted set $X$, the three many-sorted $\Sigma$-algebras $\mathbf{F}_{\Sigma}(\mathbf{D}_{\Sigma}(X))$, $\overline{\mathbf{T}}_{\Sigma}(\mathbf{D}_{\Sigma}(X))$, and $\mathbf{T}_{\Sigma}(\mathbf{D}_{\Sigma}(X))$ are  naturally isomorphic to $\mathbf{T}_{\Sigma}(X)$, the free many-sorted $\Sigma$-algebra on $X$. This follows from the fact that 
% $\mathbf{D}_{\Sigma}\dashv G_{\mathsf{PAlg}(\Sigma)}$ (see Remark~\ref{RDiscAdj}), 
% $\mathbf{F}_{\Sigma}\dashv \mathrm{In}_{\mathsf{Alg}(\Sigma)}$ and 
% $G_{\mathsf{PAlg}(\Sigma)}\circ \mathrm{In}_{\mathsf{Alg}(\Sigma)} = G_{\mathsf{Alg}(\Sigma)}$.
% \end{remark}
%
% To better understand the working of the functor $\mathbf{F}_{\Sigma}$, we provide the commutative diagram in Figure~\ref{FFreeCompFun}.
%
% Let us note that for $f\colon \mathbf{D}_{\Sigma}(X)\mor \mathbf{B}$ the diagram in Figure~\ref{FFreeCompFun} becomes the commutative diagram in Figure~\ref{FFreeCompFunDisc}
%
% \begin{figure}
% \begin{center}
% \begin{tikzpicture}
% [ACliment/.style={-{To [angle'=45, length=5.75pt, width=4pt, round]}
% }, scale=0.7]
% %,font=\scriptsize}]
% \node[] (A) 	at (0,-1) 	[] {$\mathbf{D}_{\Sigma}(X)$};
% \node[] (B) 	at (0,-4) 	[] {$\mathbf{B}$};
% \node[] (TA) 	at (4,-1) 	[] {$\mathbf{T}_{\Sigma}(X)$};
% \node[] (TB) 	at (10,-7) 	[] {$\mathbf{T}_{\Sigma}(\mathbf{B})$};
% \node[] (TpB) 	at (7,-5.5) 	[] 
% {$\overline{\mathbf{T}}_{\Sigma}(\mathbf{B})$};
% \node[] (FB) 	at (4,-4) 	[] {$\mathbf{F}_{\Sigma}(\mathbf{B})$};
% \draw[ACliment]  (A) 	to node [left]	{$f$} (B);
% \draw[ACliment, bend left=45]  (TA) 	to node [above right]{$f^{@}  = (\eta^{\mathbf{B}}\circ f)^{\sharp}$} (TB);
% \draw[ACliment, bend right]  (B) 	to node [below]	{$\eta^{\mathbf{B}}$} (TB);
% \draw[ACliment, bend right=20]  (B) 	to node [below]	{$\mathrm{in}^{\mathbf{B}}$} (TpB);
% \draw[ACliment]  (TB) 	to node [below left]	{$(\mathrm{in}^{\mathbf{B}})^{\sharp}$} (TpB);
% \draw[ACliment]  (A) 	to node [below]	{$\eta^{\mathbf{A}}$} (TA);
% \draw[ACliment]  (B) 	to node [above]	{$\eta^{\mathbf{B}}$} (FB);
% \draw[ACliment]  (TA) 	to node [right]	{$f^{@}  = (\eta^{\mathbf{B}}\circ f)^{\mathrm{fc}}$} (FB);
% \draw[ACliment]  (FB) 	to node [below left]	{$\mathrm{in}^{\mathbf{F}_{\Sigma}(\mathbf{B})}$} (TpB);
% \draw[ACliment, bend right]  (TB) 	to node [above right]	{$(\eta^{\mathbf{B}})^{\sharp}$} (FB);
% \end{tikzpicture}
% \end{center}
% \caption{The free completion functor on a discrete $\Sigma$-algebra.}
% \label{FFreeCompFunDisc}
% \end{figure}
%
% %Let us recall that, for this particular case $f^{\mathrm{fc}}=(\overline{\eta}^{\mathbf{B}}\circ f)^{\sharp}$.
%
% \begin{remark}
% For every $S$-sorted set $X$, $\mathbf{T}_{\Sigma}(X)$ together with $\eta^{X}$ is the free completion of $\mathbf{D}_{\Sigma}(X)$.
% \end{remark}
%
% In the following proposition we state that the free completion of a partial $\Sigma$-algebra can also be characterized internally.
%
% \begin{proposition}\label{PFreeCompInt}
% Let $\mathbf{A}$ be a partial $\Sigma$-algebra and $\mathbf{B}$ a $\Sigma$-algebra such that $\mathbf{A}$ is a weak subalgebra of $\mathbf{B}$. Then $\mathbf{B}$ together with $\mathrm{in}^{\mathbf{A},\mathbf{B}}$, the canonical embedding of $\mathbf{A}$ into $\mathbf{B}$, is a free completion of $\mathbf{A}$ if, and only if, the following conditions hold:
% \begin{enumerate}
% \item[FC1.] For every $(\mathbf{s},s)\in S^{\star}\times S$, every $\sigma\in \Sigma_{\mathbf{s},s}$ and every $(b_{j})_{j\in \bb{\mathbf{s}}}\in B_{\mathbf{s}}$, if $\sigma^{\mathbf{B}}((b_{j})_{j\in \bb{\mathbf{s}}})\in A_{s}$, then $\sigma^{\mathbf{B}}((b_{j})_{j\in \bb{\mathbf{s}}}) = \sigma^{\mathbf{A}}((b_{j})_{j\in \bb{\mathbf{s}}})$ (thus, in particular, $(b_{j})_{j\in \bb{\mathbf{s}}}\in \Dom(\sigma^{\mathbf{A}})$).
% \item[FC2.] For every $s\in S$, every $\sigma, \tau \in\Sigma_{\cdot,s}$, every $(b_{j})_{j\in \bb{\mathbf{s}}}\in B_{\mathbf{s}}$, where $\mathbf{s} = \mathrm{ar}(\sigma)$, and every $(c_{k})_{k\in \bb{\mathbf{t}}}\in B_{\mathbf{s}}$, where $\mathbf{t} = \mathrm{ar}(\tau)$, if 
% $\sigma^{\mathbf{B}}((b_{j})_{j\in \bb{\mathbf{s}}}) = \tau^{\mathbf{B}}((c_{k})_{k\in \bb{\mathbf{t}}})\nin A_{s}$, then $\sigma = \tau$ and $(b_{j})_{j\in \bb{\mathbf{s}}} = (c_{k})_{k\in \bb{\mathbf{t}}}$ (i.e., ``outside of $\mathbf{A}$'' the second axiom of the notion of partial Dedekind-Peano is satisfied).
% \item[FC3.] $\mathrm{Sg}_{\mathbf{B}}(A) = B$.
% \end{enumerate}
% \end{proposition}
%
% \begin{remark}\label{PPAlgSubNorm}
% The condition $\mathrm{FC1}$ in Proposition~\ref{PFreeCompInt} entails that the weak subalgebra $\mathbf{A}$ of $\mathbf{B}$ is even a relative subalgebra of $\mathbf{B}$. Relative subalgebras satisfying the condition $\mathrm{FC1}$ will be called \emph{normal} according to Schmidt~\cite{sch70}.
% \end{remark}
%
% \begin{remark}
% There exists a natural transformation from the restriction of the functor $\mathbf{F}_{\Sigma}$ to $\mathsf{PAlg}(\Sigma)_{\mathrm{c}}$, the category of many-sorted partial $\Sigma$-algebras and closed homomorphisms, to the functor $(\cdot)^{\infty}$ from $\mathsf{PAlg}(\Sigma)_{\mathrm{c}}$ to $\mathsf{Alg}(\Sigma)$.
% \end{remark}
%
% \begin{figure}
% $$\xymatrix@C=60pt@R=40pt{
%     \mathbf{A}
%     \ar@/_1pc/[ddr]_{f}
%     \ar@/^1pc/[drr]^{\eta^{\mathbf{A}}}
%     \ar[dr]|-{\mathrm{in}^{\mathbf{A},\mathbf{Sch}(f)}} \\
%      & \mathbf{Sch}(f)
%      \ar[d]_{f^{\mathrm{Sch}}}
%      \ar[r]^-{\mathrm{in}^{\mathbf{Sch}(f)}}  &
%      \mathbf{F}_{\Sigma}(\mathbf{A})
%      \ar[d]^{f^{@}} \\
%      &
%      \mathbf{B}
%      \ar[r]_{\eta^{\mathbf{B}}} &
%      \mathbf{F}_{\Sigma}(\mathbf{B})
%    }
% $$
% \caption{The Schmidt algebra relative to a partial algebra.}
% \label{FSch}
% \end{figure}
%
% \begin{figure}
% \begin{center}
% \begin{tikzpicture}
% [ACliment/.style={-{To [angle'=45, length=5.75pt, width=4pt, round]}
% }]
% %,font=\scriptsize}]
% \node[] (A) 	at 	(-3,1.5) 	[] 	{$\mathbf{D}_{\Sigma}(X)$};
% \node[] (B) 	at 	(0,-2) 	[] 	{$\mathbf{B}$};
% \node[]	(S)		at 	(0,0)	[]	{$\mathbf{Sch}(f)$};
% \node[] (FA)	at	(4,0)	[]	{$\mathbf{T}_{\Sigma}(X)$};
% \node[] (FB)	at	(4,-2)	[]	{$\mathbf{F}_{\Sigma}(\mathbf{B})$};
% \draw[ACliment, bend right]  (A) 	to node [below left]	{$f$} (B);
% \draw[ACliment, bend left]  (A) 	to node [above right]	{$\eta^{\mathbf{D}_{\Sigma}(X)}$} (FA);
% \draw[ACliment]  (A) 	to node [above right]	{$\mathrm{in}^{\mathbf{D}_{\Sigma}(X),\mathbf{Sch}(f)}$} (S);
% \draw[ACliment]  (B) 	to node [below]	{$\eta^{\mathbf{B}}$} (FB);
% \draw[ACliment]  (S) 	to node [above]	{$\mathrm{in}^{\mathbf{Sch}(f)}$} (FA);
% \draw[ACliment]  (S) 	to node [right]	{$f^{\mathrm{Sch}}$} (B);
% \draw[ACliment]  (FA) 	to node [right]	{$f^{@}$} (FB);
% \end{tikzpicture}
% \end{center}
% \caption{The Schmidt algebra for the discrete case.}
% \label{FSchDisc}
% \end{figure}
%
% In connection with the functor $\mathbf{F}_{\Sigma}$ from $\mathsf{PAlg}(\Sigma)$ to $\mathsf{Alg}(\Sigma)$ we have the following generalized recursion theorem, which we will also call Schmidt's construction. The just mentioned theorem, as we will see later, is the basis of a model theoretic approach to the theory of many-sorted partial algebras. Moreover, relying on it, as we will see below, we associate to terms in a many-sorted algebra \emph{partial} term operations on the underlying many-sorted set of a many-sorted partial algebra, and we state the general homomorphism theorem of Schmidt. 
%
% %for every $S$-sorted set $X$, every $s\in S$ and every partial $\Sigma$-algebra $\mathbf{A}$, to every term $P$ in $\mathrm{T}_{\Sigma}(X)_{s}$ one can associate a partial term operation $P^{\mathbf{A}}$ on $\mathbf{A}$, i.e., a \emph{partial} mapping from $A_{X} = \mathrm{Hom}(X,A)$ to $A_{s}$, and, in addition, from it we get the general homomorphism theorem of Schmidt.
%
% \begin{proposition}[The Schmidt construction]\label{PSch}
% Let $f$ be a homomorphism from the partial $\Sigma$-algebra $\mathbf{A}$ to the partial $\Sigma$-algebra $\mathbf{B}$. Then there exists an $A$-generated relative subalgebra $\mathbf{Sch}(f)$ of $\mathbf{F}_{\Sigma}(\mathbf{A})$, the free completion of $\mathbf{A}$, and a closed homomorphism $f^{\mathrm{Sch}}\colon \mathbf{Sch}(f)\mor \mathbf{B}$ such that the  diagram in Figure~\ref{FSch} commutes, where $\mathrm{in}^{\mathbf{A},\mathbf{Sch}(f)}$ is the canonical inclusion of $\mathbf{A}$ into $\mathbf{Sch}(f)$. Moreover, $f^{\mathrm{Sch}}$ is the largest homomorphic extension of $f$ to an $A$-generated relative subalgebra of $\mathbf{F}_{\Sigma}(\mathbf{A})$, and it is the only closed one of this kind. In honour of J. Schmidt, who introduced these concepts in~\cite{sch70}, we will call $f^{\mathrm{Sch}}\colon \mathbf{Sch}(f)\mor \mathbf{B}$ the \emph{Schmidt closed} $\mathbf{A}$-\emph{initial extension} of $f$ and $\mathrm{Ker}(f^{\mathrm{Sch}})$, denoted by $\mathrm{SKer}(f)$, the \emph{Schmidt kernel} of $f$.
% \end{proposition}
%
% \begin{proof}
% It suffices to take as $f^{\mathrm{Sch}}$ precisely $(\mathbf{Sch}(f),\Gamma\!_{f^{\mathrm{Sch}}},\mathbf{B})$, where $\Gamma\!_{f^{\mathrm{Sch}}}$, the underlying $S$-sorted function of $f^{\mathrm{Sch}}$, is $\mathrm{Sg}_{\mathbf{F}_{\Sigma}(\mathbf{A})\times \mathbf{B}}(\Gamma\!_{f})$ and $\mathbf{Sch}(f) = (f^{@})^{-1}[\mathbf{B}]$, where, with the customary abuse of notation, we have identified $B$ with $\eta^{\mathbf{B}}[B]$.
% \end{proof}
%
% %\begin{proposition}
% %Let $f$ be a homomorphism from the partial $\Sigma$-algebra $\mathbf{A}$ to the partial $\Sigma$-algebra $\mathbf{B}$. Then $f^{\mathrm{Sch}}\colon \mathbf{Sch}(f)\mor \mathbf{B}$ is surjective if and only if $f$ is an epimorphism.
% %\end{proposition}
% %
% %\begin{proof}
% %The reason why this is so is that 
% %$$
% %\mathrm{Sg}_{\mathbf{B}}(f[A]) = \mathrm{Sg}_{\mathbf{B}}(f^{\mathrm{Sch}}[A]) = f^{\mathrm{Sch}}(\mathrm{Sg}_{\mathbf{Sch}(f)}(A)) = f^{\mathrm{Sch}}[\mathrm{Sch}(f)].
% %$$
% %\end{proof}
%
% \begin{remark}\label{RPAlgInit}
% For a $\Sigma$-algebra $\mathbf{A}$, the $A$-generated relative subalgebras of 
% $\mathbf{F}_{\Sigma}(\mathbf{A})$ are sometimes also called $\mathbf{A}$-\emph{initial segments} of $\mathbf{F}_{\Sigma}(\mathbf{A})$, and the closed congruences on $\mathbf{A}$-initial segments of $\mathbf{F}_{\Sigma}(\mathbf{A})$ are sometimes also called $\mathbf{A}$-\emph{initial congruences} of 
% $\mathbf{F}_{\Sigma}(\mathbf{A})$.
% \end{remark}
%
% \begin{remark}
% The canonical inclusion $\mathrm{in}^{\mathbf{A},\mathbf{Sch}(f)}$ of $\mathbf{A}$ into $\mathbf{Sch}(f)$ is an $\mathsf{Alg}(\Sigma)$-extendable epimorphism, where, we recall, a homomorphism $f\colon \mathbf{C}\mor \mathbf{D}$ in the category $\mathsf{PAlg}(\Sigma)$ is called $\mathsf{Alg}(\Sigma)$-\emph{extendable} if, for every $\Sigma$-algebra $\mathbf{E}$ in the category $\mathsf{Alg}(\Sigma)$ and every homomorphism $g\colon \mathbf{C}\mor \mathbf{E}$, there exists a homomorphism $h\colon \mathbf{D}\mor \mathbf{E}$ such that $h\circ f = g$.
% %
% %Ana Pasztor, in~$\mathrm{MR}0860269$, her review of \cite{kw86}, says, with respect to $\mathbf{Sch}(f)$ (our notation) the following: ``intuitively is the free \emph{partial} [we emphasize] completion of $\mathbf{A}$ to a maximal degree at which there is still a homomorphism $f^{\mathrm{Sch}}\colon \mathbf{Sch}(f)\mor \mathbf{B}$ (\emph{our notation}) extending $f$''.
% \end{remark}
%
% From Proposition~\ref{PSch} we obtain, for a homomorphism whose source is a discrete many-sorted partial algebra, the following corollary.
%
% \begin{corollary}\label{CPAlgPartial}
% Let $X$ be an $S$-sorted set, $\mathbf{B}$ a partial $\Sigma$-algebra, and $f$ an $S$-sorted mapping from $X$ to $B$ or, what is equivalent, a homomorphism from $\mathbf{D}_{\Sigma}(X)$ to $\mathbf{B}$. Then there exists an $X$-generated relative subalgebra $\mathbf{Sch}(f)$ of $\mathbf{T}_{\Sigma}(X)$, the free completion of $\mathbf{D}_{\Sigma}(X)$, and a closed homomorphism $$f^{\mathrm{Sch}}\colon \mathbf{Sch}(f)\mor \mathbf{B}$$ such that the diagram in Figure~\ref{FSchDisc} commutes, where $\mathrm{in}^{\mathbf{D}_{\Sigma}(X),\mathbf{Sch}(f)}$ is the canonical inclusion of $X$ into $\mathbf{Sch}(f)$. Moreover, $f^{\mathrm{Sch}}\colon \mathbf{Sch}(f)\mor \mathbf{B}$ is the largest homomorphic extension of $f$ to an $X$-generated relative subalgebra of $\mathbf{F}_{\Sigma}(X)$, and it is the only closed one of this kind.
% \end{corollary}
%
% Let us note that for $f\colon\mathbf{D}_{\Sigma}(X)\mor\mathbf{B}$ the diagram in Figure~\ref{FSch} becomes the diagram represented in Figure~\ref{FSchDisc}. Moreover, in this particular case, we have that $f^{@}=(\eta^{\mathbf{B}}\circ f)^{\sharp}$ and 
% $\mathbf{Sch}(f)=((\eta^{\mathbf{B}}\circ f)^{\sharp})^{-1}[\eta^{\mathbf{B}}[\mathbf{B}]]$.
%
% \begin{corollary}\label{CPAlgInitGen}
% Let $\mathbf{B}$ be a partial $\Sigma$-algebra, $X$ a generating subset of $\mathbf{B}$, and 
% $\mathrm{in}^{\mathbf{D}_{\Sigma}(X),\mathbf{B}}$ the canonical embedding of $\mathbf{D}_{\Sigma}(X)$ into $\mathbf{B}$. Then $\mathbf{B}$ is a closed homomorphic image of an $X$-generated relative subalgebra---namely of $\mathbf{Sch}(\mathrm{in}^{\mathbf{D}_{\Sigma}(X),\mathbf{B}})$---of $\mathbf{T}_{\Sigma}(X)$. (Therefore, for every $s\in S$, $\mathrm{card}(B_{s})\leq \mathrm{card}(\mathrm{T}_{\Sigma}(X)_{s})$.)
% \end{corollary}
%
% \begin{remark}\label{RSch}
% We would like to highlight that in the proof of Proposition~\ref{PSch} as well as in Corollary~\ref{PPAlgPartial}, we have made use of the free completion both in the domain and in the codomain. One of the main reasons for the development of Climent and Cosme~\cite{CVCL23}---which aimed to show the functoriality of a generalization of the Schmidt construction---was to understand the fact that these free completions can be replaced by others with similar behaviour. This led us, after defining a suitable category of completions and morphisms between them, to state and prove a generalization of the Schmidt construction. The reader can go to Proposition~\ref{PSch} to become aware of the generalisation of the Schmidt construction and to check the similarities between its proof and that of Proposition~\ref{PSch}.
% \end{remark}
%
% Also, on the basis of Proposition~\ref{PSch}, we associate to terms in a many-sorted algebra partial term operations, i.e., (generalized) partial operations, on the underlying many-sorted set of a many-sorted partial algebra. We point out that, as we shall see later on, partial term operations account for the semantics of the existence equations, the existentially conditioned existence equations, and the quasi-existence equations.
%
%
% \begin{definition}
% Let $X$ be an $S$-sorted set, $\mathbf{A}$ a partial $\Sigma$-algebra, $s\in S$ and $P\in \mathrm{T}_{\Sigma}(X)_{s}$. Then we denote by $P^{\mathbf{A}}$ the partial mapping from $A_{X} = \mathrm{Hom}(X,A)$ to $A_{s}$ defined as follows: 
% $$
% \Dom(P^{\mathbf{A}}) = \{a\in A_{X}\mid P\in \mathrm{Sch}(a)_{s}\}
% $$ 
% and, for every valuation $a\in \Dom(P^{\mathbf{A}})$ of $X$ in $A$, $P^{\mathbf{A}}(a) = a^{\mathrm{Sch}}_{s}(P)$. We will call $P^{\mathbf{A}}$ the \emph{partial term operation} on $\mathbf{A}$ induced by the term $P$. Moreover, we will say that $P$ is \emph{evaluable} in $\mathbf{A}$ with respect to a valuation $a\in A_{X}$ if and only if $P\in \mathrm{Sch}(a)_{s}$ (if and only if $a\in \Dom(P^{\mathbf{A}})$).
% \end{definition}
%
%
% \begin{remark}
% Let $X$ be an $S$-sorted set of variables, $\mathbf{A}$ a partial $\Sigma$-algebra, $s\in S$ and $P\in \mathrm{T}_{\Sigma}(X)_{s}$.  Then we have that: If $P = x$, then $P^{\mathbf{A}}$ is the mapping $\mathrm{pr}^{A}_{X,s,x}$ from $A_{X}$ to $A_{s}$ that sends $a\in A_{X}$ to $\mathrm{pr}^{A}_{X,s,x}(a) = a_{s}(x)$; if $P = \sigma^{\mathbf{T}_{\Sigma}(X)}$ for some constant operation symbol $\sigma\in \Sigma_{\lambda,s}$ and $\Dom(\sigma^{\mathbf{A}})\neq\varnothing$, then $P^{\mathbf{A}}$ is the mapping from $A_{X}$ to $A_{s}$ that sends $a\in A_{X}$ to $\sigma^{\mathbf{A}}\in A_{s}$, while if $\Dom(\sigma^{\mathbf{A}}) = \varnothing$, then $\Dom(P^{\mathbf{A}}) = \varnothing$ and $P^{\mathbf{A}}$ is the unique mapping from $\varnothing$ to $A_{s}$; finally, if $P = \sigma^{\mathbf{T}_{\Sigma}(X)}((P_{i})_{i\in \bb{\mathbf{s}}})$, for $(\mathbf{s},s)\in S^{\star}\times S$ and $\sigma\in \Sigma_{\mathbf{s},s}$, then $P^{\mathbf{A}}$ is the partial mapping from $A_{X}$ to $A_{s}$, defined as follows:
% $$
% \textstyle
% \Dom(P^{\mathbf{A}}) = \{a\in A_{X}\mid a\in \bigcap_{i\in \bb{\mathbf{s}}}\Dom(P_{i}^{\mathbf{A}}) \And (P^{\mathbf{A}}_{i}(a))_{i\in \bb{\mathbf{s}}}\in \Dom(\sigma^{\mathbf{A}})\},
% $$ 
% and, for every $a\in \Dom(P^{\mathbf{A}})$, $P^{\mathbf{A}}(a) = \sigma^{\mathbf{A}}((P^{\mathbf{A}}_{i}(a))_{i\in \bb{\mathbf{s}}})$. Thus $P^{\mathbf{A}}$ is the partial mapping
% $\sigma^{\mathbf{A}}\circ \left<P^{\mathbf{A}}_{i}\right>_{i\in\bb{\mathbf{s}}}$ from $A_{X}$ to $A_{s}$, where the partial mapping $\left<P^{\mathbf{A}}_{i}\right>_{i\in\bb{\mathbf{s}}}$ from $A_{X}$ to $A_{\mathbf{s}}$ has the following universal property: It is the unique partial mapping from $A_{X}$ to $A_{\mathbf{s}}$ such that, for every $i\in \bb{\mathbf{s}}$, $\mathrm{pr}_{\mathbf{s},i}\circ  \left<P^{\mathbf{A}}_{i}\right>_{i\in\bb{\mathbf{s}}} \leq P^{\mathbf{A}}_{i}$, where $\mathrm{pr}_{\mathbf{s},i}$ is the $i$-th canonical projection from $A_{\mathbf{s}}$ to $A_{s_{i}}$, and, for every partial mapping $h$ from $A_{X}$ to $A_{s}$, if, for every $i\in \bb{\mathbf{s}}$, $\mathrm{pr}_{\mathbf{s},i}\circ h \leq P^{\mathbf{A}}_{i}$, then $h\leq \left<P^{\mathbf{A}}_{i}\right>_{i\in\bb{\mathbf{s}}}$ (we recall that for two partial mappings $u, v\colon X\dmor Y$, $u\leq v$ if, and only if, 
% $\Dom(u)\subseteq \Dom(v)$ and, for every $x\in \Dom(u)$, $u(x) = v(x)$).
%
% Let us also note that the second case is a particular instance of the last case. 
% \end{remark}
%
% Another very important tool provided by the generalized recursion theorem is the general homomorphism theorem of Schmidt. As a preparation we present the following characterization of epimorphisms.
%
% \begin{lemma}
% Let $f$ be a homomorphism from the partial $\Sigma$-algebra $\mathbf{A}$ to the partial $\Sigma$-algebra $\mathbf{B}$. Then $f$ is an epimorphism if and only if $f^{\mathrm{Sch}}\colon \mathbf{Sch}(f)\mor \mathbf{B}$ is surjective.
% \end{lemma}
%
% \begin{proof}
% The reason why this is so is that 
% $$
% \mathrm{Sg}_{\mathbf{B}}(f[A]) = \mathrm{Sg}_{\mathbf{B}}(f^{\mathrm{Sch}}[A]) = f^{\mathrm{Sch}}[\mathrm{Sg}_{\mathbf{Sch}(f)}(A)] = f^{\mathrm{Sch}}[\mathrm{Sch}(f)].
% $$
% \end{proof}
%
% \begin{proposition}
% Let $f$ be a epimorphism from the partial $\Sigma$-algebra $\mathbf{A}$ to the partial $\Sigma$-algebra $\mathbf{B}$ and $g$ a homomorphism from $\mathbf{A}$ to the partial $\Sigma$-algebra $\mathbf{C}$. Then the following statements are equivalent:
% \begin{enumerate}
% \item[(i)] There exists a unique homomorphism $h$ from $\mathbf{B}$ to $\mathbf{C}$ such that $g = h\circ f$.
% \item[(ii)] $\mathrm{SKer}(f)\subseteq \mathrm{SKer}(g)$ (as binary relations on the underlying $S$-sorted of
% $\mathbf{F}_{\Sigma}(\mathbf{A})$).
% \end{enumerate}
% \end{proposition}
%
% \begin{corollary}
% Let $f$ be an epimorphism from the partial $\Sigma$-algebra $\mathbf{A}$ to the partial $\Sigma$-algebra $\mathbf{B}$. Then $\mathbf{B}\cong \mathbf{Sch}(f)/\mathrm{SKer}(f)$.
% \end{corollary}
%
% %Our next aim is to prove that there exists an endofunctor, denoted by $\mathrm{BS}$ in honour of Burmeister and Schmidt, of the twisted morphism category $\mathbf{Mor}_{\mathrm{tw}}(\mathsf{PAlg}(\Sigma))$ of $\mathsf{PAlg}(\Sigma)$, thus showing the naturalness of the previous construction. For this purpose, we first state the definition of $\mathbf{Mor}_{\mathrm{tw}}(\mathsf{PAlg}(\Sigma))$.
% %
% %\begin{definition}\label{DPAlgM} 
% %Let $\mathbf{Mor}_{\mathrm{tw}}(\mathsf{PAlg}(\Sigma))$ be the category with objects all $\Sigma$-homo\-morphisms between many-sorted partial $\Sigma$-algebras, and with morphisms from $f\colon \mathbf{A}\mor\mathbf{B}$ to $f'\colon\mathbf{A}'\mor\mathbf{B}'$ the ordered pairs $(g,h)$ where $g$ is a $\Sigma$-homo\-morphism from $\mathbf{A}$ to $\mathbf{A}'$ and $h$ a $\Sigma$-homo\-morphisms from $\mathbf{B}'$ to $\mathbf{B}$ such that $f=h\circ f'\circ g$.
% %%the following diagram commutes
% %%\begin{center}
% %%\begin{tikzpicture}[scale=0.8,
% %%ACliment/.style={-{To [angle'=45, length=5.75pt, width=4pt, round]}}]
% %%\tikzstyle{every state}=[very thick, draw=blue!50,fill=blue!20, inner sep=2pt,minimum size=20pt]
% %%\node[] (a) at (0,0) [] {$\mathbf{A}$};
% %%\node[] (ap) at (4,0) [] {$\mathbf{A}'$};
% %%\node[] (b) at (0,-2) [] {$\mathbf{B}$};
% %%\node[] (bp) at (4,-2) [] {$\mathbf{B}'$};
% %%
% %%\draw[ACliment]  (a) to node [left] {$f$} (b);
% %%\draw[ACliment]  (ap) to node [right] {$f'$} (bp);
% %%\draw[ACliment]  (a) to node [above] {$g$} (ap);
% %%\draw[ACliment]  (bp) to node [below] {$h$} (b);
% %%\end{tikzpicture}
% %%\end{center}
% %%
% %%In this category, for a $\Sigma$-homomorphism between many-sorted partial $\Sigma$-algebras $f\colon \mathbf{A}\mor\mathbf{B}$, the identity at $f$, written $\id^{f}$, is given by the pair $(\id^{\mathbf{A}},\id^{\mathbf{B}})$. 
% %%Let us note that  the diagram in Figure~\ref{FIdInPAlgM} commutes.
% %%\begin{center}
% %%\begin{tikzpicture}[scale=0.8,
% %%ACliment/.style={-{To [angle'=45, length=5.75pt, width=4pt, round]}}]
% %%\tikzstyle{every state}=[very thick, draw=blue!50,fill=blue!20, inner sep=2pt,minimum size=20pt]
% %%\node[] (a) at (0,0) [] {$\mathbf{A}$};
% %%\node[] (ap) at (4,0) [] {$\mathbf{A}$};
% %%\node[] (b) at (0,-2) [] {$\mathbf{B}$};
% %%\node[] (bp) at (4,-2) [] {$\mathbf{B}$};
% %%
% %%\draw[ACliment]  (a) to node [left] {$f$} (b);
% %%\draw[ACliment]  (ap) to node [right] {$f$} (bp);
% %%\draw[ACliment]  (a) to node [above] {$\id^{\mathbf{A}}$} (ap);
% %%\draw[ACliment]  (bp) to node [below] {$\id^{\mathbf{B}}$} (b);
% %%\end{tikzpicture}
% %%\end{center}
% %
% %%Furthermore, given three $\Sigma$-homomorphisms $f\colon \mathbf{A}\mor\mathbf{B}$, $f'\colon\mathbf{A}'\mor\mathbf{B}'$ and $f''\colon\mathbf{A}''\mor\mathbf{B}''$, and morphisms $(g,h)$ and $(g',h')$ in $\mathsf{PAlg}(\Sigma)^{\rightarrow}$ from, respectively, $f$ to $f'$ and $f'$ to $f''$, the composition of $(g,h)$ with $(g',h')$ is given by $(g'\circ g, h\circ h')$. 
% %%All in all the diagram in Figure~\ref{FCompInPAlgM} commutes.
% %%\begin{center}
% %%\begin{tikzpicture}[scale=0.8,
% %%ACliment/.style={-{To [angle'=45, length=5.75pt, width=4pt, round]}}]
% %%\tikzstyle{every state}=[very thick, draw=blue!50,fill=blue!20, inner sep=2pt,minimum size=20pt]
% %%\node[] (a) at (0,0) [] {$\mathbf{A}$};
% %%\node[] (ap) at (4,0) [] {$\mathbf{A}'$};
% %%\node[] (ap2) at (8,0) [] {$\mathbf{A}''$};
% %%\node[] (b) at (0,-2) [] {$\mathbf{B}$};
% %%\node[] (bp) at (4,-2) [] {$\mathbf{B}'$};
% %%\node[] (bp2) at (8,-2) [] {$\mathbf{B}''$};
% %%
% %%\draw[ACliment]  (a) to node [left] {$f$} (b);
% %%\draw[ACliment]  (ap) to node [right] {$f'$} (bp);
% %%\draw[ACliment]  (a) to node [above] {$g$} (ap);
% %%\draw[ACliment]  (bp) to node [below] {$h$} (b);
% %%\draw[ACliment, bend left]  (a) to node [above] {$g'\circ g$} (ap2);
% %%\draw[ACliment]  (ap2) to node [right] {$f''$} (bp2);
% %%\draw[ACliment]  (ap) to node [above] {$g'$} (ap2);
% %%\draw[ACliment]  (bp2) to node [below] {$h'$} (bp);
% %%\draw[ACliment, bend left]  (bp2) to node [below] {$h\circ h'$} (b);
% %%\end{tikzpicture}
% %%\end{center}
% %\end{definition}
% %
% %\begin{figure}
% %\begin{center}
% %\begin{tikzpicture}[scale=0.8,
% %ACliment/.style={-{To [angle'=45, length=5.75pt, width=4pt, round]}}]
% %\tikzstyle{every state}=[very thick, draw=blue!50,fill=blue!20, inner sep=2pt,minimum size=20pt]
% %\node[] (Sf) at (0,0) [] {$\mathbf{Sch}(f)$};
% %\node[] (FA) at (4,0) [] {$\mathbf{F}_{\Sigma}(\mathbf{A})$};
% %\node[] (B) at (0,-2) [] {$\mathbf{B}$};
% %\node[] (FB) at (4,-2) [] {$\mathbf{F}_{\Sigma}(\mathbf{B})$};
% %\node[] (A) at (-4,2)  [] {$\mathbf{A}$};
% %
% %\draw[ACliment]  (Sf) to node [left] {$f^{\mathrm{Sch}}$} (B);
% %\draw[ACliment]  (FA) to node [left] {$f^{@}$} (FB);
% %\draw[ACliment]  (Sf) to node [above] {$\mathrm{in}^{\mathbf{Sch}(f)}$} (FA);
% %\draw[ACliment]  (B) to node [below] {$\eta^{\mathbf{B}}$} (FB);
% %\draw[ACliment, bend left=15]  (A) to node [above right] {$\eta^{\mathbf{A}}$} (FA);
% %\draw[ACliment]  (A) to node [midway, fill=white] {$\mathrm{in}^{\mathbf{A}, \mathbf{Sch}(f)}$} (Sf);
% %\draw[ACliment, bend right=15]  (A) to node [below left] {$f$} (B);
% %
% %\node[] (Sfp) at (2,4) [] {$\mathbf{Sch}(f')$};
% %\node[] (FAp) at (6,4) [] {$\mathbf{F}_{\Sigma}(\mathbf{A'})$};
% %\node[] (Bp) at (2,2) [] {$\mathbf{B}'$};
% %\node[] (FBp) at (6,2) [] {$\mathbf{F}_{\Sigma}(\mathbf{B}')$};
% %\node[] (Ap) at (-2,6)  [] {$\mathbf{A}'$};
% %
% %\draw[ACliment]  (Sfp) to node [left] {$f'^{\mathrm{Sch}}$} (Bp);
% %\draw[ACliment]  (FAp) to node [right] {$f'^{@}$} (FBp);
% %\draw[ACliment]  (Sfp) to node [above] {$\mathrm{in}^{\mathbf{Sch}(f')}$} (FAp);
% %\draw[ACliment]  (Bp) to node [below] {$\eta^{\mathbf{B}'}$} (FBp);
% %\draw[ACliment, bend left=15]  (Ap) to node [above right]  {$\eta^{\mathbf{A}'}$} (FAp);
% %
% %\draw[ACliment]  (Ap) to node [midway, fill=white] {$\mathrm{in}^{\mathbf{A}', \mathbf{Sch}(f')}$} (Sfp);
% %\draw[ACliment, bend right=15]  (Ap) to node [below left] {$f'$} (Bp);
% %
% %\draw[ACliment] (A) to node [above left] {$g$} (Ap);
% %\draw[ACliment] (Bp) to node [below right] {$h$} (B);
% %\draw[ACliment] (FA) to node [above left, xshift=.35cm, yshift=.35cm] {$g^{@}$} (FAp);
% %\draw[ACliment] (FBp) to node [below right] {$h^{@}$} (FB);
% %\end{tikzpicture}
% %\end{center}
% %\caption{Pasting Schmidt diagrams.}
% %\label{FPastSch}
% %\end{figure}
% %
% %\begin{remark}\label{SchForTwoHoms}
% %Let $f\colon \mathbf{A}\mor\mathbf{B}$ and $f'\colon\mathbf{A}'\mor\mathbf{B}'$ be two objects of the category $\mathbf{Mor}_{\mathrm{tw}}(\mathsf{PAlg}(\Sigma))$ and let $(g,h)$ be a morphism from $f$ to $f'$. Then, for the homomorphisms $g^{@}$ from\ $\mathbf{F}_{\Sigma}(\mathbf{A})$ to $\mathbf{F}_{\Sigma}(\mathbf{A}')$ and $h^{@}$ from $\mathbf{F}_{\Sigma}(\mathbf{B}')$ to $\mathbf{F}_{\Sigma}(\mathbf{B})$, we have, by Corollary~\ref{CFreeAdj}, that $\eta^{\mathbf{A}'}\circ g=g^{@}\circ\eta^{\mathbf{A}}$ and $\eta^{\mathbf{B}}\circ h=h^{@}\circ\eta^{\mathbf{B}'}$. Thus, taking into account the Proposition~\ref{PSch}, the diagram in Figure~\ref{FPastSch} commutes.
% %
% %We suggest the reader to retain in his memory the previous commutativities since they will be used, in the proof of the following proposition, to justify certain equalities.
% %\end{remark}
% %
% %We next prove that the Schmidt construction stated in Proposition~\ref{PSch} is functorial.
% %
% %\begin{figure}
% %\begin{center}
% %\begin{tikzpicture}[scale=0.9,
% %ACliment/.style={-{To [angle'=45, length=5.75pt, width=4pt, round]}},
% %ACliment2/.style={|-{To [angle'=45, length=5.75pt, width=4pt, round]}}
% %]
% %\tikzstyle{every state}=[very thick, draw=blue!50,fill=blue!20, inner sep=2pt,minimum size=20pt]
% %\node[] (a) at (0,0) [] {$\mathbf{A}$};
% %\node[] (ap) at (4,0) [] {$\mathbf{A}'$};
% %\node[] (b) at (0,-2) [] {$\mathbf{B}$};
% %\node[] (bp) at (4,-2) [] {$\mathbf{B}'$};
% %
% %\draw[ACliment]  (a) to node [left] {$f$} (b);
% %\draw[ACliment]  (ap) to node [right] {$f'$} (bp);
% %\draw[ACliment]  (a) to node [above] {$g$} (ap);
% %\draw[ACliment]  (bp) to node [below] {$h$} (b);
% %
% %\node[] () at (5.75,-.6) [] {$\mathrm{BS}$};
% %\draw[ACliment2] (5,-1) to node [] {} (6.5,-1);
% %
% %\node[] (Sa) at (8,0) [] {$\mathbf{Sch}(f)$};
% %\node[] (Sap) at (12,0) [] {$\mathbf{Sch}(f')$};
% %\node[] (Sb) at (8,-2) [] {$\mathbf{B}$};
% %\node[] (Sbp) at (12,-2) [] {$\mathbf{B}'$};
% %
% %\draw[ACliment]  (Sa) to node [left] {$f^{\mathrm{Sch}}$} (Sb);
% %\draw[ACliment]  (Sap) to node [right] {$f'^{\mathrm{Sch}}$} (Sbp);
% %\draw[ACliment]  (Sa) to node [above] {$g^{@}\circ\mathrm{in}^{\mathbf{Sch}(f)}$} (Sap);
% %\draw[ACliment]  (Sbp) to node [below] {$h$} (Sb);
% %\end{tikzpicture}
% %\end{center}
% %\caption{The Burmeister-Schmidt functor.}
% %\label{FBS}
% %\end{figure}
% %
% %\begin{proposition}\label{PFuncSchmidt}  
% %Let $\mathrm{BS}$ be the mapping that sends
% %\begin{enumerate}
% %\item an object $f\colon\mathbf{A}\mor\mathbf{B}$ of $\mathbf{Mor}_{\mathrm{tw}}(\mathsf{PAlg}(\Sigma))$ to the object $f^{\mathrm{Sch}}\colon\mathbf{Sch}(f)\mor\mathbf{B}$ of the same category, and
% %\item a morphism $(g,h)$ of $\mathbf{Mor}_{\mathrm{tw}}(\mathsf{PAlg}(\Sigma))$ from the object $f$ to the object $f'$ to the morphism $(g^{@}\circ\mathrm{in}^{\mathbf{Sch}(f)},h)$ of the same category from the object $f^{\mathrm{Sch}}$ to the object $f'^{\mathrm{Sch}}$.
% %\end{enumerate}
% %We depict the action of $\mathrm{BS}$ on objects and morphisms in the diagram in Figure~\ref{FBS}.
% %
% %Then $\mathrm{BS}$, the \emph{Burmeister-Schmidt} functor, is an endofunctor of $\mathbf{Mor}_{\mathrm{tw}}(\mathsf{PAlg}(\Sigma))$.
% %\end{proposition}
% %
% %\begin{proof} 
% %The proof will consist of the following series of claims.
% %
% %\label{ClaimA}
% %\begin{claim}\label{CA0} The $\Sigma$-homomorphism $g^{@}\circ\mathrm{in}^{\mathbf{Sch}(f)}$ from $\mathbf{Sch}(f)$ to $\mathbf{F}_{\Sigma}(\mathbf{A}')$ co-restricts to $\mathbf{Sch}(f')$. Thus we need to prove that
% %$g^{@}[\mathrm{in}^{\mathbf{Sch}(f)}[\mathrm{Sch}(f)]]\subseteq \mathrm{Sch}(f')$. But 
% %$\mathrm{Sch}(f')=(f'^{@})^{-1}[\eta^{\mathbf{B}'}[B']]$. Hence, it suffices to prove that
% %$$
% %f'^{@}
% %\left[
% %g^{@}\left[
% %\mathrm{in}^{\mathbf{Sch}(f)}\left[
% %\mathrm{Sch}(f)
% %\right]\right]\right]\subseteq \eta^{\mathbf{B}'}[B'].
% %$$
% %\end{claim}
% %
% %Let us recall that $\mathbf{Sch}(f)$ is an $A$-generated relative subalgebra of $\mathbf{F}_{\Sigma}(\mathbf{A})$. Since $f'^{@}$, $g^{@}$, and $\mathrm{in}^{\mathbf{Sch}(f)}$ are $\Sigma$-homomorphisms, it suffices to prove that, for every sort $s\in S$ and every element $a\in A_{s}$, the following statement holds
% %$$
% %f'^{@}_{s}\left(
% %g^{@}_{s}\left(
% %\mathrm{in}^{\mathbf{Sch}(f)}_{s}\left(
% %\mathrm{in}^{\mathbf{A},\mathbf{Sch}(f)}_{s}\left(
% %a
% %\right)\right)\right)\right)\in \eta^{\mathbf{B}'}_{s}\left[
% %B_{s}\right].
% %$$
% %
% %But we have that the following chain of equalities holds
% %\allowdisplaybreaks
% %\begin{align*}
% %f'^{@}_{s}\left(
% %g^{@}_{s}\left(
% %\mathrm{in}^{\mathbf{Sch}(f)}_{s}\left(
% %\mathrm{in}^{\mathbf{A},\mathbf{Sch}(f)}_{s}\left(
% %a
% %\right)\right)\right)\right)&=
% %f'^{@}_{s}\left(
% %g^{@}_{s}\left(
% %\eta^{\mathbf{A}}_{s}\left(
% %a
% %\right)\right)\right)
% %\tag{1}
% %\\&=
% %f'^{@}_{s}\left(
% %\eta^{\mathbf{A}'}_{s}\left(
% %g_{s}\left(
% %a
% %\right)\right)\right)
% %\tag{2}
% %\\&=
% %\eta^{\mathbf{B}'}_{s}\left(
% %f'_{s}\left(
% %g_{s}\left(
% %a
% %\right)\right)\right).
% %\tag{3}
% %\end{align*}
% %
% %In the just stated chain of equalities, the first equality follows from the fact that $\mathrm{in}^{\mathbf{Sch}(f)}\circ\mathrm{in}^{\mathbf{A},\mathbf{Sch}(f)}=\eta^{\mathbf{A}}$; the second equality follows from the the fact that $\eta^{\mathbf{A}'}\circ g=g^{@}\circ\eta^{\mathbf{A}}$; and the third equality follows from the fact that $\eta^{\mathbf{B}'}\circ f'=f'^{@}\circ\eta^{\mathbf{A}}$.
% %
% %Let us note that, for every sort $s\in S$ and every element $a\in A_{s}$, the element $f'_{s}(g_{s}(a))$ belongs to $B'_{s}$. Therefore, we can affirm that, for every sort $s\in S$ and every element $a\in A_{s}$, the element
% %$f'^{@}_{s}
% %(g^{@}_{s}
% %(\mathrm{in}^{\mathbf{Sch}(f)}_{s}
% %(\mathrm{in}^{\mathbf{A},\mathbf{Sch}(f)}_{s}
% %(a))))$ is in $ \eta^{\mathbf{B}'}_{s}[B_{s}]$
% %
% %For simplicity of notation we continue to write $g^{@}\circ\mathrm{in}^{\mathbf{Sch}(f)}$, which is a $\Sigma$-homomorphism from $\mathbf{Sch}(f)$ to $\mathbf{F}_{\Sigma}(\mathbf{A}')$, for its co-restriction  to $\mathbf{Sch}(f')$.
% %
% %Claim~\ref{CA0} follows.
% %
% %\label{ClaimB}
% %\begin{claim}\label{CB0} For every morphism $(g,h)$ from $f$ to $f'$ in $\mathbf{Mor}_{\mathrm{tw}}(\mathsf{PAlg}(\Sigma))$, we have that $(g^{@}\circ\mathrm{in}^{\mathbf{Sch}(f)},h)$ is a morphism from $f^{\mathrm{Sch}}$ to $f'^{\mathrm{Sch}}$  in $\mathbf{Mor}_{\mathrm{tw}}(\mathsf{PAlg}(\Sigma))$. Thus we need to prove that
% %$$
% %h\circ f'^{\mathrm{Sch}}\circ g^{@}\circ \mathrm{in}^{\mathbf{Sch}(f)}=f^{\mathrm{Sch}}.
% %$$
% %\end{claim}
% %
% %Let us recall that $\mathbf{Sch}(f)$ is an $A$-generated relative subalgebra of $\mathbf{F}_{\Sigma}(\mathbf{A})$. Since $h$, $ f'^{\mathrm{Sch}}$, $g^{@}$, and $\mathrm{in}^{\mathbf{Sch}(f)}$ are $\Sigma$-homomorphisms, it suffices to prove that, for every sort $s\in S$, and every element $a\in A_{s}$, the following statement holds
% %$$
% %h_{s}\left(
% %f'^{\mathrm{Sch}}_{s}\left(
% %g^{@}_{s}\left(
% %\mathrm{in}^{\mathbf{Sch}(f)}_{s}\left(
% %\mathrm{in}^{\mathbf{A},\mathbf{Sch}(f)}_{s}\left(
% %a
% %\right)\right)\right)\right)\right)
% %=f^{\mathrm{Sch}}_{s}\left(
% %\mathrm{in}^{\mathbf{A},\mathbf{Sch}(f)}_{s}\left(
% %a
% %\right)\right).
% %$$
% %
% %The following chain of equalities holds
% %\begin{flushleft}
% %$h_{s}\left(
% %f'^{\mathrm{Sch}}_{s}\left(
% %g^{@}_{s}\left(
% %\mathrm{in}^{\mathbf{Sch}(f)}_{s}\left(
% %\mathrm{in}^{\mathbf{A},\mathbf{Sch}(f)}_{s}\left(
% %a
% %\right)\right)\right)\right)\right)$
% %\allowdisplaybreaks
% %\begin{align*}
% %\qquad
% %&=
% %h_{s}\left(
% %f'^{\mathrm{Sch}}_{s}\left(
% %g^{@}_{s}\left(
% %\eta^{\mathbf{A}}_{s}\left(
% %a
% %\right)\right)\right)\right)
% %\tag{1}
% %\\&=
% %h_{s}\left(
% %f'^{\mathrm{Sch}}_{s}\left(
% %\eta^{\mathbf{A}'}_{s}\left(
% %g_{s}\left(
% %a
% %\right)\right)\right)\right)
% %\tag{2}
% %\\&=
% %h_{s}\left(
% %f'^{\mathrm{Sch}}_{s}\left(
% %\mathrm{in}^{\mathbf{Sch}(f')}_{s}\left(
% %\mathrm{in}^{\mathbf{A}',\mathbf{Sch}(f')}_{s}\left(
% %g_{s}\left(
% %a
% %\right)\right)\right)\right)\right)
% %\tag{3}
% %\\&=
% %h_{s}\left(
% %f'^{\mathrm{Sch}}_{s}\left(
% %\mathrm{in}^{\mathbf{A}',\mathbf{Sch}(f')}_{s}\left(
% %g_{s}\left(
% %a
% %\right)\right)\right)\right)
% %\tag{4}
% %\\&=
% %h_{s}\left(
% %f'_{s}\left(
% %g_{s}\left(
% %a
% %\right)\right)\right)
% %\tag{5}
% %\\&=
% %f_{s}(a)
% %\tag{6}
% %\\&=
% %f^{\mathrm{Sch}}_{s}\left(
% %\mathrm{in}^{\mathbf{A},\mathbf{Sch}(f)}_{s}\left(
% %a
% %\right)\right).
% %\tag{7}
% %\end{align*}
% %\end{flushleft}
% %
% %In the just stated chain of equalities, the first equality follows from the fact that 
% %$\mathrm{in}^{\mathbf{Sch}(f)}\circ \mathrm{in}^{\mathbf{A},\mathbf{Sch}(f)}=\eta^{\mathbf{A}}$; the second equality follows from the fact that $\eta^{\mathbf{A}'}\circ g=g^{@}=\eta^{\mathbf{A}}$; the third equality follows from the fact that $\mathrm{in}^{\mathbf{Sch}(f')}\circ \mathrm{in}^{\mathbf{A}',\mathbf{Sch}(f')}=\eta^{\mathbf{A}'}$; the fourth equality follows from the fact that the value of $\mathrm{in}_{s}^{\mathbf{Sch}(f)}$ at the element $\mathrm{in}^{\mathbf{A}',\mathbf{Sch}(f')}_{s}(g_{s}(a))$ is this element itself; the fifth equality follows from the fact that $f'^{\mathrm{Sch}}\circ \mathrm{in}^{\mathbf{A}',\mathbf{Sch}(f')}=f'$; the sixth equation follows from the fact that $(g,h)$ is a morphism in $\mathbf{Mor}_{\mathrm{tw}}(\mathsf{PAlg}(\Sigma))$ from $f$ to $f'$, i.e., $f=h\circ f'\circ g$; and the seventh equality follows from the fact that $f^{\mathrm{Sch}}\circ \mathrm{in}^{\mathbf{A},\mathbf{Sch}(f)}=f$.
% %
% %Claim~\ref{CB0} follows.
% %
% %Therefore $\mathrm{BS}$ maps objects to objects and morphisms to morphisms.
% %
% %\label{ClaimC}
% %\begin{claim}\label{CC0} Let $f\colon\mathbf{A}\mor\mathbf{B}$ be a $\Sigma$-homomorphism between many-sorted partial $\Sigma$-algebras. We prove that $\mathrm{BS}$ sends the identity morphism of $f$ to the identity morphism of $\mathrm{BS}(f)$, i.e., that
% %$$
% %\mathrm{BS}(\id^{f})=\id^{\mathrm{BS}(f)}.
% %$$
% %\end{claim}
% %
% %
% %
% %We recall that $\id^{f}$, the identity at $f$ in $\mathbf{Mor}_{\mathrm{tw}}(\mathsf{PAlg}(\Sigma))$, is $(\id^{\mathbf{A}},\id^{\mathbf{B}})$. Moreover, $\mathrm{BS}(\id^{f}) = (\id^{\mathbf{A}}^{@}\circ\mathrm{in}^{\mathbf{Sch}(f)},\id^{\mathbf{B}})$. For this pair of $\Sigma$-homomorphisms the diagram on the left and immediately below commutes. On the other hand, $\mathrm{BS}(f)=f^{\mathrm{Sch}}$. Hence, the identity at $f^{\mathrm{Sch}}$ is $(\id^{\mathbf{Sch}(f)},\id^{\mathbf{B}})$. For this pair of $\Sigma$-homomorphisms the diagrams of Figure~\ref{FCC0} commute.
% %
% %\begin{figure}
% %\begin{center}
% %\begin{tikzpicture}[scale=0.8,
% %ACliment/.style={-{To [angle'=45, length=5.75pt, width=4pt, round]}}]
% %\tikzstyle{every state}=[very thick, draw=blue!50,fill=blue!20, inner sep=2pt,minimum size=20pt]
% %\node[] (a) at (0,0) [] {$\mathbf{Sch}(f)$};
% %\node[] (ap) at (4,0) [] {$\mathbf{Sch}(f)$};
% %\node[] (b) at (0,-2) [] {$\mathbf{B}$};
% %\node[] (bp) at (4,-2) [] {$\mathbf{B}$};
% %
% %\draw[ACliment]  (a) to node [left] {$f^{\mathrm{Sch}}$} (b);
% %\draw[ACliment]  (ap) to node [right] {$f^{\mathrm{Sch}}$} (bp);
% %\draw[ACliment]  (a) to node [above] {$\id^{\mathbf{A}}^{@}\circ\mathrm{in}^{\mathrm{Sch}(f)}$} (ap);
% %\draw[ACliment]  (bp) to node [below] {$\id^{\mathbf{B}}$} (b);
% %
% %
% %
% %\node[] (Sa) at (8,0) [] {$\mathbf{Sch}(f)$};
% %\node[] (Sap) at (12,0) [] {$\mathbf{Sch}(f)$};
% %\node[] (Sb) at (8,-2) [] {$\mathbf{B}$};
% %\node[] (Sbp) at (12,-2) [] {$\mathbf{B}$};
% %
% %\draw[ACliment]  (Sa) to node [left] {$f^{\mathrm{Sch}}$} (Sb);
% %\draw[ACliment]  (Sap) to node [right] {$f^{\mathrm{Sch}}$} (Sbp);
% %\draw[ACliment]  (Sa) to node [above] {$\id^{\mathbf{Sch}(f)}$} (Sap);
% %\draw[ACliment]  (Sbp) to node [below] {$\id^{\mathbf{B}}$} (Sb);
% %\end{tikzpicture}
% %\end{center}
% %\caption{The Burmeister-Schmidt functor on identities.}
% %\label{FCC0}
% %\end{figure}
% %
% %Therefore, in order to check that $\mathrm{BS}(\id^{f})=\id^{\mathrm{BS}(f)}$, it suffices to verify that
% %\[\id^{\mathbf{A}}^{@}\circ\mathrm{in}^{\mathbf{Sch}(f)}=\id^{\mathbf{Sch}(f)}.\]
% %
% %Let us recall that $\mathbf{Sch}(f)$ is an $A$-generated relative subalgebra of $\mathbf{F}_{\Sigma}(\mathbf{A})$. Since $\id^{\mathbf{A}}^{@}$ and $\mathrm{in}^{\mathbf{Sch}(f)}$ are $\Sigma$-homomorphisms, it suffices to prove that, for every sort $s\in S$ and every element $a\in A_{s}$, the following statement holds
% %$$
% %\id^{\mathbf{A}_{s}}^{@}\left(
% %\mathrm{in}^{\mathbf{Sch}(f)}_{s}\left(
% %\mathrm{in}^{\mathbf{A},\mathbf{Sch}(f)}_{s}\left(
% %a
% %\right)\right)\right)
% %=
% %\id^{\mathbf{Sch}(f)_{s}}\left(
% %\mathrm{in}^{\mathbf{A},\mathbf{Sch}(f)}_{s}\left(
% %a
% %\right)\right).
% %$$
% %
% %But we have that the following chain of equalities holds
% %\begin{align*}
% %\id^{\mathbf{A}_{s}}^{@}\left(
% %\mathrm{in}^{\mathbf{Sch}(f)}_{s}\left(
% %\mathrm{in}^{\mathbf{A},\mathbf{Sch}(f)}_{s}\left(
% %a
% %\right)\right)\right)
% %&=
% %\id^{\mathbf{A}_{s}}^{@}\left(
% %\eta^{\mathbf{A}}_{s}\left(
% %a
% %\right)\right)
% %\tag{1}
% %\\&=
% %\eta^{\mathbf{A}}_{s}\left(
% %\id^{\mathbf{A}_{s}}\left(
% %a
% %\right)\right)
% %\tag{2}
% %\\&=
% %\eta^{\mathbf{A}}_{s}\left(
% %a
% %\right)
% %\tag{3}
% %\\&=
% %\mathrm{in}^{\mathbf{Sch}(f)}_{s}\left(
% %\mathrm{in}^{\mathbf{A},\mathbf{Sch}(f)}_{s}\left(
% %a
% %\right)\right)
% %\tag{4}
% %\\&=
% %\id^{\mathbf{Sch}(f)_{s}}\left(
% %\mathrm{in}^{\mathbf{A},\mathbf{Sch}(f)}_{s}\left(
% %a
% %\right)\right).
% %\tag{5}
% %\end{align*}
% %
% %In the just stated chain of equalities, the first equality follows from the fact that 
% %$\mathrm{in}^{\mathbf{Sch}(f)}\circ \mathrm{in}^{\mathbf{A},\mathbf{Sch}(f)}=\eta^{\mathbf{A}}$; the second equality follows from the fact that $\eta^{\mathbf{A}}\circ \id^{\mathbf{A}}=\id^{\mathbf{A}}^{@}\circ \eta^{\mathbf{A}}$; the third equation applies the identity mapping at $a$; and  the fourth equality follows from the fact that 
% %$\mathrm{in}^{\mathbf{Sch}(f)}\circ \mathrm{in}^{\mathbf{A},\mathbf{Sch}(f)}=\eta^{\mathbf{A}}$; the fifth equation follows from the fact that $(\mathrm{in}^{\mathbf{A},\mathbf{Sch}(f)}_{s}(a)$ is an element of $\mathrm{Sch}_{s}(f)$.
% %
% %Claim~\ref{CC0} follows.
% %
% %\label{ClaimD}
% %\begin{claim}\label{CD0} Let $f\colon\mathbf{A}\mor\mathbf{B}$, $f'\colon\mathbf{A}'\mor\mathbf{B}'$ and $f''\colon\mathbf{A}''\mor\mathbf{B}''$ be $\Sigma$-homomorphisms between many-sorted partial $\Sigma$-algebras. Let $(g,h)$ and $(g',h')$ be morphisms in
% %$\mathbf{Mor}_{\mathrm{tw}}(\mathsf{PAlg}(\Sigma))$ from $f$ to $f'$ and from $f'$ to $f''$, respectively.
% %We prove that $\mathrm{BS}$ preserves the composition, i.e., that
% %$$
% %\mathrm{BS}\left(
% %(g',h')\circ (g,h)
% %\right)=\mathrm{BS}\left(
% %(g',h')\right)\circ\mathrm{BS}\left(
% %(g,h)\right).
% %$$
% %\end{claim}
% %
% %We recall that $(g',h')\circ (g,h)$, the composition in $\mathbf{Mor}_{\mathrm{tw}}(\mathsf{PAlg}(\Sigma))$, is $(g'\circ g, h\circ h')$. Moreover, $\mathrm{BS}(g'\circ g, h\circ h')$ is $((g'\circ g)^{@}\circ \mathrm{in}^{\mathbf{Sch}(f)}, h\circ h')$. For this pair of $\Sigma$-homomorphisms the outer rectangle of the diagram immediately below commutes. On the other hand, $\mathrm{BS}((g',h'))\circ\mathrm{BS}((g,h)) = (g'^{@}\circ\mathrm{in}^{\mathbf{Sch}(f')}\circ g^{@}\circ \mathrm{in}^{\mathbf{Sch}(f)}, h\circ h')$. For this pair of $\Sigma$-homomorphisms the inner rectangle of the diagram immediately below commutes. Let us recall that, since $(g,h)$ and $(g',h')$ are  morphisms in $\mathbf{Mor}_{\mathrm{tw}}(\mathsf{PAlg}(\Sigma))$ from $f$ to $f'$ and from $f'$ to $f''$, respectively, the inner small rectagles of the diagram immediately below of Figure~\ref{FCD0} also commute.
% %
% %\begin{figure}
% %\begin{center}
% %\begin{tikzpicture}[scale=0.7,
% %ACliment/.style={-{To [angle'=45, length=5.75pt, width=4pt, round]}}]
% %\tikzstyle{every state}=[very thick, draw=blue!50,fill=blue!20, inner sep=2pt,minimum size=20pt]
% %\node[] (a) at (0,0) [] {$\mathbf{Sch}(f)$};
% %\node[] (ap) at (6,0) [] {$\mathbf{Sch}(f')$};
% %\node[] (ap2) at (12,0) [] {$\mathbf{Sch}(f'')$};
% %\node[] (b) at (0,-3) [] {$\mathbf{B}$};
% %\node[] (bp) at (6,-3) [] {$\mathbf{B}'$};
% %\node[] (bp2) at (12,-3) [] {$\mathbf{B}''$};
% %\node [] () at (6,2) [] {$(g'\circ g)^{@}\circ\mathrm{in}^{\mathbf{Sch}(f)}$};
% %\node [] () at (6,-5) [] {$h\circ h'$};
% %
% %\draw[ACliment]  (a) to node [left] {$f^{\mathrm{Sch}}$} (b);
% %\draw[ACliment]  (ap) to node [right] {$f'^{\mathrm{Sch}}$} (bp);
% %\draw[ACliment]  (a) to node [above] {$g^{@}\circ \mathrm{in}^{\mathbf{Sch}(f)}$} (ap);
% %\draw[ACliment]  (bp) to node [below] {$h$} (b);
% %
% %\draw[ACliment]  (ap2) to node [right] {$f''^{\mathrm{Sch}}$} (bp2);
% %\draw[ACliment]  (ap) to node [above] {$g'^{@}\circ \mathrm{in}^{\mathbf{Sch}(f')}$} (ap2);
% %\draw[ACliment]  (bp2) to node [below] {$h'$} (bp);
% %
% %\draw[ACliment, rounded corners] (a.north)
% %--
% % ($(a.south)+(0,1)$)
% %|-
% %($(ap2.north)+(0,1)$)
% %-- (ap2.north);
% %
% %\draw[ACliment, rounded corners] (bp2.south)
% %--
% % ($(bp2.north)+(0,-1)$)
% %|-
% %($(b.south)+(0,-1)$)
% %-- (b.south);
% %\end{tikzpicture}
% %\end{center}
% %\caption{The Burmeister-Schmidt functor on compositions.}
% %\label{FCD0}
% %\end{figure}
% %
% %Therefore, in order to check that $\mathrm{BS}((g',h')\circ (g,h))=\mathrm{BS}((g',h'))\circ\mathrm{BS}((g,h))$, it suffices to verify that
% %$$
% %g'^{@}\circ\mathrm{in}^{\mathbf{Sch}(f')}\circ g^{@}\circ \mathrm{in}^{\mathbf{Sch}(f)}=\left(
% %g'\circ g\right)^{@}\circ \mathrm{in}^{\mathbf{Sch}(f)}.
% %$$
% %Let us recall that $\mathbf{Sch}(f)$ is an $A$-generated relative subalgebra of $\mathbf{F}_{\Sigma}(\mathbf{A})$. Since $g'^{@}$, $\mathrm{in}^{\mathbf{Sch}(f')}$, $g^{@}$, and $\mathrm{in}^{\mathbf{Sch}(f)}$ are $\Sigma$-homomorphisms, it suffices to check that, for every sort $s\in S$, and every element $a\in A_{s}$, the following equality holds
% %\begin{multline*}
% %g'^{@}_{s}\left(
% %\mathrm{in}^{\mathbf{Sch}(f')}_{s}\left(
% %g^{@}_{s}\left(
% %\mathrm{in}^{\mathbf{Sch}(f)}_{s}\left(
% %\mathrm{in}^{\mathbf{A},\mathbf{Sch}(f)}_{s}\left(
% %a
% %\right)\right)\right)\right)\right)
% %\\=
% %\left(g'\circ g\right)^{@}_{s}
% %\left(
% %\mathrm{in}^{\mathbf{Sch}(f)}_{s}\left(
% %\mathrm{in}^{\mathbf{A},\mathbf{Sch}(f)}_{s}\left(
% %a
% %\right)\right)\right).
% %\end{multline*}
% %
% %But we have that the following chain of equalities holds
% %\begin{flushleft}
% %$g'^{@}_{s}\left(
% %\mathrm{in}^{\mathbf{Sch}(f')}_{s}\left(
% %g^{@}_{s}\left(
% %\mathrm{in}^{\mathbf{Sch}(f)}_{s}\left(
% %\mathrm{in}^{\mathbf{A},\mathbf{Sch}(f)}_{s}\left(
% %a
% %\right)\right)\right)\right)\right)$
% %\allowdisplaybreaks
% %\begin{align*}
% %\qquad
% %&=
% %g'^{@}_{s}\left(
% %\mathrm{in}^{\mathbf{Sch}(f')}_{s}\left(
% %g^{@}_{s}\left(
% %\eta^{\mathbf{A}}_{s}\left(
% %a
% %\right)\right)\right)\right)
% %\tag{1}
% %\\&=
% %g'^{@}_{s}\left(
% %\mathrm{in}^{\mathbf{Sch}(f')}_{s}\left(
% %\eta^{\mathbf{A}'}_{s}\left(
% %g_{s}\left(
% %a
% %\right)\right)\right)\right)
% %\tag{2}
% %\\&=
% %g'^{@}_{s}\left(
% %\mathrm{in}^{\mathbf{Sch}(f')}_{s}\left(
% %\mathrm{in}^{\mathbf{Sch}(f')}_{s}\left(
% %\mathrm{in}^{\mathbf{A}',\mathbf{Sch}(f')}_{s}\left(
% %g_{s}\left(
% %a
% %\right)\right)\right)\right)\right)
% %\tag{3}
% %\\&=
% %g'^{@}_{s}\left(
% %\mathrm{in}^{\mathbf{Sch}(f')}_{s}\left(
% %\mathrm{in}^{\mathbf{A}',\mathbf{Sch}(f')}_{s}\left(
% %g_{s}\left(
% %a
% %\right)\right)\right)\right)
% %\tag{4}
% %\\&=
% %g'^{@}_{s}\left(
% %\eta^{\mathbf{A}'}_{s}\left(
% %g_{s}\left(
% %a
% %\right)\right)\right)
% %\tag{5}
% %\\&=
% %\eta^{\mathbf{A}''}\left(
% %g'_{s}\left(
% %g_{s}\left(
% %a
% %\right)\right)\right)
% %\tag{6}
% %\\&=
% %\left(g'\circ g\right)^{@}_{s}\left(
% %\eta^{\mathbf{A}}\left(
% %a
% %\right)\right)
% %\tag{7}
% %\\&=
% %\left(g'\circ g\right)^{@}_{s}\left(
% %\mathrm{in}^{\mathbf{Sch}(f)}_{s}\left(
% %\mathrm{in}^{\mathbf{A},\mathbf{Sch}(f)}_{s}\left(
% %a
% %\right)\right)\right).
% %\tag{8}
% %\end{align*}
% %\end{flushleft}
% %
% %In the just stated chain of equalities, the first equality follows from the fact that 
% %$\mathrm{in}^{\mathbf{Sch}(f)}\circ \mathrm{in}^{\mathbf{A},\mathbf{Sch}(f)}=\eta^{\mathbf{A}}$; the second equality follows from the fact that $\eta^{\mathbf{A}'}\circ g=g^{@}\circ\eta^{\mathbf{A}}$; the third equality follows from the fact that 
% %$\mathrm{in}^{\mathbf{Sch}(f)'}\circ \mathrm{in}^{\mathbf{A}',\mathbf{Sch}(f)'}=\eta^{\mathbf{A}'}$; the fourth equality follows from the fact that $\mathrm{in}^{\mathbf{Sch}(f)}$ is idempotent at every sort; the fifth equality follows from the fact that 
% %$\mathrm{in}^{\mathbf{Sch}(f)'}\circ \mathrm{in}^{\mathbf{A}',\mathbf{Sch}(f)'}=\eta^{\mathbf{A}'}$; the sixth equality follows from the fact that $\eta^{\mathbf{A}''}\circ g'=g'^{@}\circ\eta^{\mathbf{A}'}$; the seventh equality follows from the fact that $\eta^{\mathbf{A}''}\circ (g'\circ g)=(g'\circ g)^{@}\circ\eta^{\mathbf{A}}$; and the eighth equality follows from the fact that 
% %$\mathrm{in}^{\mathbf{Sch}(f)}\circ \mathrm{in}^{\mathbf{A},\mathbf{Sch}(f)}=\eta^{\mathbf{A}}$.
% %
% %Claim~\ref{CD0} follows.
% %
% %This completes the proof.
% %\end{proof}
%
%
% \subsection{Varieties of  many-sorted partial algebras}
%
% The concept of existence equation is necessary for the syntactic description of certain classes of many-sorted  partial algebras, e.g., the categories, the $n$-categories, etc. 
% %Just as a signature fixes the class of algebras over that signature, so an
% %equation serves to fix the class of algebras in which the terms of the equation have the same value.
%
% \begin{definition}
% Let $\Sigma$ be an $S$-sorted signature.
% \begin{enumerate}
% \item If $(X,s)\in\boldsymbol{\mathcal{U}}^{S}\times S$, then a
%       $\Sigma$-\emph{term of type} $(X,s)$ is an
%       element of $\mathrm{T}_{\Sigma}(X)_{s}$. If a $\Sigma$-term
%       of type $(X,s)$ is such that $X$ is $S$-finite, i.e., for every $s\in S$, $X_{s}$ is finite, then we will
%       call it $S$-\emph{finitary}, and if it is such that $X$ is
%       finite, i.e., $\mathrm{supp}_{S}(X)$ is finite and, for every $s\in \mathrm{supp}_{S}(X)$, $X_{s}$ is 
%       finite, then we will call it \emph{finitary}.
%
% \item We denote by $\mathrm{Tm}(\Sigma)$ the $\boldsymbol{\mathcal{U}}^{S}\times S$-sorted
%       set $(\mathrm{T}_{\Sigma}(X)_{s})_{(X,s)\in\boldsymbol{\mathcal{U}}^{S}\bprod S}$ 
%       of all $\Sigma$-terms, by $\mathrm{Tm}_{S\text{-}\mathrm{f}}(\Sigma)$
%       the $\boldsymbol{\mathcal{U}}^{S}_{S\text{-}\mathrm{f}}\times
%       S$-sorted set 
%       $(\mathrm{T}_{\Sigma}(X)_{s})_{(X,s)\in\boldsymbol{\mathcal{U}}^{S}_{S\text{-}\mathrm{f}}\bprod S}$ 
%       of all $S$-finitary terms, where $\boldsymbol{\mathcal{U}}^{S}_{S\text{-}\mathrm{f}}$
%       is the set of all $S$-sorted sets which are $S$-finite, and by
%       $\mathrm{Tm}_{\mathrm{f}}(\Sigma)$ the $\boldsymbol{\mathcal{U}}^{S}_{\mathrm{f}}\times
%       S$-sorted set 
%       $(\mathrm{T}_{\Sigma}(X)_{s})_{(X,s)\in\boldsymbol{\mathcal{U}}^{S}_{\mathrm{f}}\bprod S}$ of all finitary terms, where $\boldsymbol{\mathcal{U}}^{S}_{\mathrm{f}}$ is
%       the set of all $S$-sorted sets which are finite.
% \end{enumerate}
% \end{definition}
%
% For $(X,s)\in\boldsymbol{\mathcal{U}}^{S}\times S$, we next define the notions of existence equation of type $(X,s)$, of $S$-finitary existence equation of type $(X,s)$ and of finitary existence equation of type $(X,s)$.
%
% \begin{definition}
% Let $\Sigma$ be an $S$-sorted signature.
% \begin{enumerate}
% \item If $(X,s)\in\boldsymbol{\mathcal{U}}^{S}\times S$, then an \emph{existence equation}
%       (briefly $\mathrm{E}$-\emph{equation} or $\mathrm{Eeqt}$) of type $(X,s)$ is a pair $(P,Q)$, also denoted by
%       $P\eeq Q$, where $P$ and $Q$ are $\Sigma$-terms of type $(X,s)$.
%       If an $\mathrm{E}$-equation of type $(X,s)$ is such that $X$ is $S$-finite, then we will call it
%       $S$\emph{-finitary}, and if it is such that $X$ is finite, then we will call it \emph{finitary}.
%
% \item We denote by $\mathrm{Eeqt}(\Sigma)$ the $\boldsymbol{\mathcal{U}}^{S}\times S$-sorted set 
%       $\mathrm{Tm}(\Sigma)^{2}$ of
%       all $\mathrm{E}$-equations, by $\mathrm{Eeqt}_{S\text{-}\mathrm{f}}(\Sigma)$ the
%       $\boldsymbol{\mathcal{U}}^{S}_{S\text{-}\mathrm{f}}\times S$-sorted set 
%       $\mathrm{Tm}_{S\text{-}\mathrm{f}}(\Sigma)^{2}$ of all
%       $S$-finitary $\mathrm{E}$-equations, by $\mathrm{Eeqt}_{\mathrm{f}}(\Sigma)$
%       the $\boldsymbol{\mathcal{U}}^{S}_{\mathrm{f}}\times
%       S$-sorted set $\mathrm{Tm}_{\mathrm{f}}(\Sigma)^{2}$ of all finitary $\mathrm{E}$-equations, and by $\mathrm{Eeqt}(\Sigma)_{X}$
%       the $S$-sorted set $(\mathrm{T}_{\Sigma}(X)_{s}^{2})_{s\in S}$ of all $\mathrm{E}$-equations
%       with variables in an $S$-sorted set $X$.
% \item We will say that a family of $\mathrm{E}$-equations $\mathcal{E}\subseteq \mathrm{Eeqt}(\Sigma)$ is
%       $S$-\emph{finitary} (resp. \emph{finitary}) if
%       $\mathcal{E}\subseteq\mathrm{Eeqt}_{S\text{-} \mathrm{f}}(\Sigma)$
%       (resp. $\mathcal{E}\subseteq\mathrm{Eeqt}_{\mathrm{f}}(\Sigma)$).
% \end{enumerate}
% \end{definition}
%
% For $(X,s)\in\boldsymbol{\mathcal{U}}^{S}\times S$, we next define the ternary relation of satisfiability $\cdot \models^{\Sigma}_{X,s}\!\!\cdot\, [\cdot]$ between a partial $\Sigma$-algebra, an $\mathrm{E}$-equation of type $(X,s)$, and a valuation of $X$ in the underlying $S$-sorted set of the partial $\Sigma$-algebra, and the derived binary relation of validity $\cdot \models^{\Sigma}_{X,s}\!\!\cdot$ between a partial $\Sigma$-algebra and a $\mathrm{E}$-equation of type $(X,s)$. 
%
% \begin{definition}
% Let $\mathbf{A}$ be a partial $\Sigma$-algebra, $(X,s)\in\boldsymbol{\mathcal{U}}^{S}\times S$, $(P,Q)\in
% \mathrm{Eeqt}(\Sigma)_{X,s}$, $\mathcal{E}\subseteq \mathrm{Eeqt}(\Sigma)$, and $a\in A_{X} = \mathrm{Hom}(X,A)$.
% \begin{enumerate}
% \item We will say that $\mathbf{A}$ \emph{satisfies the $\mathrm{E}$-equation} $(P,Q)$
%       \emph{with respect to the valuation} $a$, denoted by 
%       $$
%       \mathbf{A}\models^{\Sigma}_{X,s}(P,Q) [a],
%       $$
%       if and only if $P\in \mathrm{Sch}(a)_{s}$, $Q\in \mathrm{Sch}(a)_{s}$,
%       and $a^{\mathrm{Sch}}_{s}(P) = a^{\mathrm{Sch}}_{s}(Q)$.
% \item We will say that $(P,Q)$ is \emph{valid} in $\mathbf{A}$, denoted by
%       $$
%       \mathbf{A}\models^{\Sigma}_{X,s}(P,Q),
%       $$ 
%       or by $\mathbf{A}\models(P,Q)$,
%       when this is unlikely to cause confusion, if, for every $a\in A_{X}$,
%       $\mathbf{A}\models^{\Sigma}_{X,s}(P,Q) [a]$.
% \item We will say that $\mathcal{E}$ is \emph{valid} in $\mathbf{A}$, or that $\mathbf{A}$ \emph{satisfies} $\mathcal{E}$, denoted by
%       $$
%       \mathbf{A}\models^{\Sigma}\mathcal{E},
%       $$ 
%       or by $\mathbf{A}\models\mathcal{E}$,
%       when this is unlikely to cause confusion, if, for every 
%       $(X,s)\in\boldsymbol{\mathcal{U}}^{S}\times S$ and every $(P,Q)\in\mathcal{E}_{X,s}$, 
%       $\mathbf{A}\models^{\Sigma}_{X,s}(P,Q)$.
% \item We will denote by $\mathrm{Mod}_{\Sigma}(\mathcal{E})$, or by $\mathrm{PAlg}(\Sigma,\mathcal{E})$, the set $$
% \{\mathbf{A}\in \mathrm{PAlg}(\Sigma)\mid \mathbf{A}\models^{\Sigma}\mathcal{E}\}
% $$
% of all partial $\Sigma$-algebras that satisfy $\mathcal{E}$.
% \end{enumerate}
% \end{definition}
%
% \begin{remark}
% Let $\mathbf{A}$ be a partial $\Sigma$-algebra, $(P,Q)\in\mathrm{Eeqt}(\Sigma)_{X,s}$, and $a\in A_{X}$.
% If $P\in \mathrm{Sch}(a)_{s}$ and $Q\in \mathrm{Sch}(a)_{s}$, then, since $P^{\mathbf{A}}(a) = a^{\mathrm{Sch}}_{s}(P)$ and $Q^{\mathbf{A}}(a) = a^{\mathrm{Sch}}_{s}(P)$, we have that
% $a^{\mathrm{Sch}}_{s}(P) = a^{\mathrm{Sch}}_{s}(Q)$ if and only if $P^{\mathbf{A}}(a) = Q^{\mathbf{A}}(a)$.
% \end{remark}
%
% \begin{definition}
% Let $\mathbf{A}$ be a partial $\Sigma$-algebra and $\mathcal{E}\subseteq \mathrm{Eeqt}(\Sigma)$.
% We will say that $\mathcal{E}$ is \emph{valid} in $\mathbf{A}$, denoted by
% $\mathbf{A}\models^{\Sigma} \mathcal{E}$, if, for every $(X,s)\in
% \boldsymbol{\mathcal{U}}^{S}\times S$ and every $(P,Q)\in \mathcal{E}_{X,s}$ we have that
% $\mathbf{A}\models^{\Sigma}_{X,s} (P,Q)$.
% \end{definition}
%
% In some cases it is convenient to consider instead of the
% $\boldsymbol{\mathcal{U}}^{S}\times S$-sorted set $\mathrm{Eeqt}(\Sigma)$ the set
% $\coprod\mathrm{Eeqt}(\Sigma) = \bigcup_{(X,s)\in\boldsymbol{\mathcal{U}}^{S}\times S}(\mathrm{T}_{\Sigma}(X)_{s}^{2}\times\{(X,s)\})$. Since
% $$\mathbf{\Sub}(\mathrm{Eeqt}(\Sigma))\iso\mathbf{\Sub}(\coprod\mathrm{Eeqt}(\Sigma)),$$ it follows that
% every family $\mathcal{E}\subseteq\mathrm{Eeqt}(\Sigma)$ has univocally associated a subset of $\coprod\mathrm{Eeqt}(\Sigma)$ and reciprocally. In the same way, sometimes, it is convenient to consider instead of the $S$-sorted set $\mathrm{Eeqt}(\Sigma)_{X}$ the set $\coprod\mathrm{Eeqt}(\Sigma)_{X}$.
%
% \begin{remark}
% $\mathrm{Sub}(\mathrm{Eeqt}(\Sigma))$ and $\prod_{X\in
% \boldsymbol{\mathcal{U}}^{S}}\mathrm{Sub}(\coprod\mathrm{Eeqt}(\Sigma)_{X})$ are naturally isomorphic. In fact, it suffices to consider the mapping that assigns to $\mathcal{E}\subseteq \mathrm{Eeqt}(\Sigma)$ precisely 
% $(\coprod_{s\in S}\mathcal{E}_{X,s})_{X\in \boldsymbol{\mathcal{U}}^{S}}$. 
% \end{remark}
%
% For $(X,(\mathbf{s},s))$ in $\boldsymbol{\mathcal{U}}^{S}\times (S^{\star}\times S)$, we next define the notions of existentially conditioned existence equation of type $(X,(\mathbf{s},s))$, of $S$-finitary existentially conditioned existence equation of type $(X,(\mathbf{s},s))$ and of finitary existentially conditioned existence equation of type $(X,(\mathbf{s},s))$.
%
% \begin{definition}
% Let $\Sigma$ be an $S$-sorted signature.
% \begin{enumerate}
% \item If $(X,(\mathbf{s},s))\in\boldsymbol{\mathcal{U}}^{S}\times (S^{\star}\times S)$, then an \emph{existentially conditioned existence equation} (briefly $\mathrm{ECE}$-\emph{equation} or $\mathrm{ECEeqt}$) of type $(X,(\mathbf{s},s))$ is a pair
% $$
% \left(\left(
% P_{j},P_{j}
% \right)_{j\in \bb{\mathbf{s}}},(P,Q)\right),
% $$
% also denoted by
% $$
% \textstyle
% \bigwedge_{j\in\bb{\mathbf{s}}}(P_{j},P_{j})\to (P,Q)\, \text{ or by }\,
% \bigwedge_{j\in\bb{\mathbf{s}}}(P_{j}\eeq P_{j})\to (P\eeq Q),
% $$
% where, for every $j\in \bb{\mathbf{s}}$, $P_{j}\in\mathrm{T}_{\Sigma}(X)_{s_{j}}$ and $(P,Q)\in \mathrm{T}_{\Sigma}(X)^{2}_{s}$.
% The pairs $(P_{j},P_{j})$, for $j\in \bb{\mathbf{s}}$, are called the \emph{premisses} and $(P,Q)$ the \emph{conclusion} of the $\mathrm{ECE}$-equation. If an $\mathrm{ECE}$-equation of type $(X,(\mathbf{s}, s))$ is such that $X$ is $S$-finite, then we will call it $S$-\emph{finitary}, and if it is such that $X$ is finite, then we will call it \emph{finitary}.
% \item We denote by $\mathrm{ECEeqt}(\Sigma)_{X}$ the $S^{\star}\times S$-sorted set of all 
%       $\mathrm{ECE}$-equations with
%       variables in $X$, by $\mathrm{ECEeqt}(\Sigma)$
%       the $\boldsymbol{\mathcal{U}}^{S}\times (S^{\star}\times S)$-sorted set of
%       all $\mathrm{ECE}$-equations, by $\mathrm{ECEeqt}_{S\text{-}\mathrm{f}}(\Sigma)$ the
%       $\boldsymbol{\mathcal{U}}^{S}_{S\text{-}\mathrm{f}}\times (S^{\star}\times S)$-sorted set of all
%       $S$-finitary $\mathrm{ECE}$-equations, and by $\mathrm{ECEeqt}_{\mathrm{f}}(\Sigma)$
%       the $\boldsymbol{\mathcal{U}}^{S}_{\mathrm{f}}\times (S^{\star}\times S)$-sorted set of all finitary
%       $\mathrm{ECE}$-equations.
% \item We will say that a family of $\mathrm{ECE}$-equations $\mathcal{E}\subseteq \mathrm{ECEeqt}(\Sigma)$ is
%       $S$-\emph{finitary} (resp. \emph{finitary}) if
%       $\mathcal{E}\subseteq\mathrm{ECEeqt}_{S\text{-} \mathrm{f}}(\Sigma)$
%       (resp. $\mathcal{E}\subseteq\mathrm{ECEeqt}_{\mathrm{f}}(\Sigma)$).
% \end{enumerate}
% \end{definition}
%
% In some cases it is convenient to consider instead of the
% $\boldsymbol{\mathcal{U}}^{S}\times (S^{\star}\times S)$-sorted set $\mathrm{ECEeqt}(\Sigma)$ the set $\coprod\mathrm{ECEeqt}(\Sigma)$, which is
% $$
% \textstyle
% \bigcup_{(X,(\mathbf{s},s))\in\boldsymbol{\mathcal{U}}^{S}\times (S^{\star}\times S)}
% \left(\left(\left(\prod_{j\in\bb{\mathbf{s}}}\Delta_{\mathrm{T}_{\Sigma}(X)_{s_{j}}}\right)\times\mathrm{T}_{\Sigma}(X)_{s}^{2}\right)\times\{(X,(\bb{\mathbf{s}},s))\}\right).
% $$
% Since
% $\mathbf{\Sub}(\mathrm{ECEeqt}(\Sigma))\iso\mathbf{\Sub}(\coprod\mathrm{ECEeqt}(\Sigma))$, it follows that
% every family $\mathcal{E}\subseteq\mathrm{ECEeqt}(\Sigma)$ has univocally associated a subset of $\coprod\mathrm{ECEeqt}(\Sigma)$ and reciprocally. In the same way, sometimes, it is convenient to consider instead of the $(S^{\star}\times S)$-sorted set $\mathrm{ECEeqt}(\Sigma)_{X}$ the set $\coprod\mathrm{ECEeqt}(\Sigma)_{X}$ which is
% $$
% \textstyle
% \bigcup_{(\mathbf{s},s)\in S^{\star}\times S}
% \left(\left(\left(
% \prod_{j\in\bb{\mathbf{s}}}\Delta_{\mathrm{T}_{\Sigma}(X)_{s_{j}}}
% \right)\times \mathrm{T}_{\Sigma}(X)^{2}_{s}
% \right)\times \{(\mathbf{s},s)\}\right),
% $$
% which, in its turn, is isomorphic to
% $$
% \textstyle
% \left(
% \coprod_{s\in S}\Delta_{\mathrm{T}_{\Sigma}(X)_{s}}
% \right)^{\star}\times \coprod_{s\in S}\mathrm{T}_{\Sigma}(X)_{s}^{2}.
% $$
%
% For $(X,(\mathbf{s},s))\in\boldsymbol{\mathcal{U}}^{S}\times (S^{\star}\times S)$, we next define the ternary relation of satisfiability $\cdot \models^{\Sigma}_{X,\mathbf{s},s}\!\!\cdot\, [\cdot]$ between a partial $\Sigma$-algebra, an $\mathrm{ECE}$-\emph{equation} of type $(X,(\mathbf{s},s))$, and a valuation of $X$ in the underlying $S$-sorted set of the partial $\Sigma$-algebra, and the derived binary relation of validity $\cdot \models^{\Sigma}_{X,\mathbf{s},s}\!\!\cdot$ between a partial $\Sigma$-algebra anf a $\mathrm{ECE}$-equation of type $(X,(\mathbf{s},s))$. 
%
% \begin{definition}
% Let $\mathbf{A}$ be a partial $\Sigma$-algebra, $(X,(\mathbf{s},s))\in\boldsymbol{\mathcal{U}}^{S}\times (S^{\star}\times S)$, $\bigwedge_{j\in\bb{\mathbf{s}}}(P_{j},P_{j})\to (P,Q)\in
% \mathrm{ECEeqt}(\Sigma)_{(X,(\mathbf{s},s))}$, $\mathcal{E}\subseteq \mathrm{ECEeqt}(\Sigma)$, and $a\in A_{X} = \mathrm{Hom}(X,A)$.
% \begin{enumerate}
% \item We will say that $\mathbf{A}$ \emph{satisfies the $\mathrm{ECE}$-equation} 
%       $\bigwedge_{j\in\bb{\mathbf{s}}}(P_{j},P_{j})\to (P,Q)$
%       \emph{with respect to the valuation} $a$ , denoted by 
%       $$
%       \mathbf{A}\models^{\Sigma}_{X,\mathbf{s},s}\bigwedge_{j\in\bb{\mathbf{s}}}(P_{j},P_{j})\to (P,Q) [a],
%       $$
%       if and only if, if, for every $j\in \bb{\mathbf{s}}$, 
%       $P_{j}\in \mathrm{Sch}(a)_{s_{j}}$ (i.e., $P^{\mathbf{A}}_{j}$ is defined at $a$), then 
%       $P\in \mathrm{Sch}(a)_{s}$, $Q\in \mathrm{Sch}(a)_{s}$,
%       and $a^{\mathrm{Sch}}_{s}(P) = a^{\mathrm{Sch}}_{s}(Q)$. In other words, if, for every 
%       $j\in \bb{\mathbf{s}}$, $\mathbf{A}\models^{\Sigma}_{X,s_{j}} (P_{j},P_{j})[a]$, then  
%       $\mathbf{A}\models^{\Sigma}_{X,s} (P,Q)[a]$.
% \item We will say that $\bigwedge_{j\in\bb{\mathbf{s}}}(P_{j},P_{j})\to (P,Q)$ is \emph{valid} in 
%       $\mathbf{A}$, denoted by
%       $$
%       \mathbf{A}\models^{\Sigma}_{X,\mathbf{s},s}\bigwedge_{j\in\bb{\mathbf{s}}}(P_{j},P_{j})\to (P,Q),
%       $$ 
%       or by $\mathbf{A}\models\bigwedge_{j\in\bb{\mathbf{s}}}(P_{j},P_{j})\to (P,Q)$,
%       when this is unlikely to cause confusion, if, for every $a\in A_{X}$,
%       $\mathbf{A}\models^{\Sigma}_{X,\mathbf{s},s}\bigwedge_{j\in\bb{\mathbf{s}}}(P_{j},P_{j})\to (P,Q) [a]$.
% \item We will say that $\mathcal{E}$ is \emph{valid} in $\mathbf{A}$, or that $\mathbf{A}$ \emph{satisfies} $\mathcal{E}$, denoted by
%       $$
%       \mathbf{A}\models^{\Sigma}\mathcal{E},
%       $$
%       or by $\mathbf{A}\models\mathcal{E}$,
%       when this is unlikely to cause confusion, if, for every 
%       $(X,(\mathbf{s},s))\in\boldsymbol{\mathcal{U}}^{S}\times (S^{\star}\times S)$ and every 
%       $\bigwedge_{j\in\bb{\mathbf{s}}}(P_{j},P_{j})\to (P,Q)\in
%       \mathcal{E}_{(X,(\mathbf{s},s))}$, 
%       $$
%       \mathbf{A}\models^{\Sigma}_{X,\mathbf{s},s}\bigwedge_{j\in\bb{\mathbf{s}}}(P_{j},P_{j})\to (P,Q).
%       $$
% \item We will denote by $\mathrm{Mod}_{\Sigma}(\mathcal{E})$, or by $\mathrm{PAlg}(\Sigma,\mathcal{E})$, the set 
% $$
% \{\mathbf{A}\in \mathrm{PAlg}(\Sigma)\mid \mathbf{A}\models^{\Sigma}\mathcal{E}\}
% $$
% of all partial $\Sigma$-algebras that satisfy $\mathcal{E}$.
% \end{enumerate}
% \end{definition}
%
% \begin{definition}
% Let $\mathbf{A}$ be a partial $\Sigma$-algebra and $\mathcal{E}\subseteq \mathrm{ECEeqt}(\Sigma)$.
% We will say that $\mathcal{E}$ is \emph{valid} in $\mathbf{A}$, denoted by
% $\mathbf{A}\models^{\Sigma} \mathcal{E}$, if, for every $(X,(\mathbf{s},s))\in
% \boldsymbol{\mathcal{U}}^{S}\times (S^{\star}\times S)$ and every $\bigwedge_{j\in\bb{\mathbf{s}}}(P_{j},P_{j})\to (P,Q)\in \mathcal{E}_{(X,(\mathbf{s},s))}$ we have that
% $\mathbf{A}\models^{\Sigma}_{X,\mathbf{s},s} \bigwedge_{j\in\bb{\mathbf{s}}}(P_{j},P_{j})\to (P,Q)$.
% \end{definition}
%
% For $(X,(\mathbf{s},s))$ in $\boldsymbol{\mathcal{U}}^{S}\times (S^{\star}\times S)$, we next define the notions of  quasi-existence equation of type $(X,(\mathbf{s},s))$, of $S$-finitary  quasi-existence equation of type $(X,(\mathbf{s},s))$ and of finitary  quasi-existence equation of type $(X,(\mathbf{s},s))$.
%
% \begin{definition}\label{DQEEq-app}
% Let $\Sigma$ be an $S$-sorted signature.
% \begin{enumerate}
% \item If $(X,(\mathbf{s},s))\in\boldsymbol{\mathcal{U}}^{S}\times (S^{\star}\times S)$, then a 
% \emph{quasi-existence equation} (briefly $\QE$-\emph{equation} or $\mathrm{QEeqt}$) of type $(X,(\mathbf{s},s))$ is a pair
% $$
% \left(\left(
% P_{j},Q_{j}
% \right)_{j\in \bb{\mathbf{s}}},(P,Q)\right),
% $$
% also denoted by
% $$
% \textstyle
% \bigwedge_{j\in\bb{\mathbf{s}}}(P_{j},Q_{j})\to (P,Q)\, \text{ or by }\,
% \bigwedge_{j\in\bb{\mathbf{s}}}(P_{j}\eeq Q_{j})\to (P\eeq Q),
% $$
% where, for every $j\in \bb{\mathbf{s}}$, $(P_{j},Q_{j})\in\mathrm{T}_{\Sigma}(X)^{2}_{s_{j}}$ and $(P,Q)\in \mathrm{T}_{\Sigma}(X)^{2}_{s}$.
% The pairs $(P_{j},Q_{j})$, for $j\in \bb{\mathbf{s}}$, are called the \emph{premisses} and $(P,Q)$ the \emph{conclusion} of the $\QE$-equation. If a $\QE$-equation of type $(X,(\mathbf{s}, s))$ is such that $X$ is $S$-finite, then we will call it $S$-\emph{finitary}, and if it is such that $X$ is finite, then we will call it \emph{finitary}.
% \item We denote by $\mathrm{QEeqt}(\Sigma)_{X}$ the $S^{\star}\times S$-sorted set of all 
%       $\QE$-equations with
%       variables in $X$, by $\mathrm{QEeqt}(\Sigma)$
%       the $\boldsymbol{\mathcal{U}}^{S}\times (S^{\star}\times S)$-sorted set of
%       all $\QE$-equations, by $\mathrm{QEeqt}_{S\text{-}\mathrm{f}}(\Sigma)$ the
%       $\boldsymbol{\mathcal{U}}^{S}_{S\text{-}\mathrm{f}}\times (S^{\star}\times S)$-sorted set of all
%       $S$-finitary $\QE$-equations, and by $\mathrm{QEeqt}_{\mathrm{f}}(\Sigma)$
%       the $\boldsymbol{\mathcal{U}}^{S}_{\mathrm{f}}\times (S^{\star}\times S)$-sorted set of all finitary
%       $\QE$-equations.
% \item We will say that a family of $\QE$-equations $\mathcal{E}\subseteq \mathrm{QEeqt}(\Sigma)$ is
%       $S$-\emph{finitary} (resp. \emph{finitary}) if
%       $\mathcal{E}\subseteq\mathrm{QEeqt}_{S\text{-} \mathrm{f}}(\Sigma)$
%       (resp. $\mathcal{E}\subseteq\mathrm{QEeqt}_{\mathrm{f}}(\Sigma)$).
% \end{enumerate}
% \end{definition}
%
% In some cases it is convenient to consider instead of the
% $\boldsymbol{\mathcal{U}}^{S}\times (S^{\star}\times S)$-sorted set $\mathrm{QEeqt}(\Sigma)$ the set $\coprod\mathrm{QEeqt}(\Sigma)$, which is
% $$
% \textstyle
% \bigcup_{(X,(\mathbf{s},s))\in\boldsymbol{\mathcal{U}}^{S}\times (S^{\star}\times S)}
% \left(\left(\left(\prod_{j\in\bb{\mathbf{s}}}\mathrm{T}_{\Sigma}(X)_{s_{j}}\right)\times\mathrm{T}_{\Sigma}(X)_{s}^{2}\right)\times\{(X,(\bb{\mathbf{s}},s))\}\right).
% $$
% Since
% $\mathbf{\Sub}(\mathrm{QEeqt}(\Sigma))\iso\mathbf{\Sub}(\coprod\mathrm{QEeqt}(\Sigma))$, it follows that
% every family $\mathcal{E}\subseteq\mathrm{QEeqt}(\Sigma)$ has univocally associated a subset of $\coprod\mathrm{QEeqt}(\Sigma)$ and reciprocally. In the same way, sometimes, it is convenient to consider instead of the $(S^{\star}\times S)$-sorted set $\mathrm{QEeqt}(\Sigma)_{X}$ the set $\coprod\mathrm{QEeqt}(\Sigma)_{X}$ which is
% $$
% \textstyle
% \bigcup_{(\mathbf{s},s)\in S^{\star}\times S}
% \left(\left(\left(
% \prod_{j\in\bb{\mathbf{s}}}\mathrm{T}_{\Sigma}(X)_{s_{j}}
% \right)\times \mathrm{T}_{\Sigma}(X)^{2}_{s}
% \right)\times \{(\mathbf{s},s)\}\right),
% $$
% which, in its turn, is isomorphic to
% $$
% \textstyle
% \left(
% \coprod_{s\in S}\mathrm{T}_{\Sigma}(X)^{2}_{s}
% \right)^{\star}\times \coprod_{s\in S}\mathrm{T}_{\Sigma}(X)_{s}^{2}.
% $$
%
% Since the semantics of the $\QE$-equations follows the same pattern that the semantics of the $\mathrm{ECE}$-equations, we leave the details to the reader. Let us just say that a partial $\Sigma$-algebra $\mathbf{A}$ satisfies a $\QE$-equation 
%       $\bigwedge_{j\in\bb{\mathbf{s}}}(P_{j},Q_{j})\to (P,Q)$
%       with respect to a valuation $a$, denoted by 
%       $$
%       \mathbf{A}\models^{\Sigma}_{X,\mathbf{s},s}\bigwedge_{j\in\bb{\mathbf{s}}}(P_{j},Q_{j})\to (P,Q) [a],
%       $$
%       if and only if, if, for every $j\in \bb{\mathbf{s}}$, 
%       $P_{j}\in \mathrm{Sch}(a)_{s_{j}}$, $Q_{j}\in \mathrm{Sch}(a)_{s_{j}}$,
%       and $a^{\mathrm{Sch}}_{s_{j}}(P_{j}) = a^{\mathrm{Sch}}_{s_{j}}(Q_{j})$, then 
%       $P\in \mathrm{Sch}(a)_{s}$, $Q\in \mathrm{Sch}(a)_{s}$,
%       and $a^{\mathrm{Sch}}_{s}(P) = a^{\mathrm{Sch}}_{s}(Q)$. In other words, if, for every 
%       $j\in \bb{\mathbf{s}}$, $\mathbf{A}\models^{\Sigma}_{X,s_{j}} (P_{j},Q_{j})[a]$, then  
%       $\mathbf{A}\models^{\Sigma}_{X,s} (P,Q)[a]$.
%
% %From now on, if there is no ambiguity, we will not distinguish notationally between both presentations.
%
% We next define, respecting the terminology used by Burmeister in~\cite{pb86}, the notions of finitary $\mathrm{E}$-, finitary $\mathrm{ECE}$-, and finitary $\QE$-varieties---which will be the type of varieties to which we will restrict ourselves in this work---and, for this purpose, we, on the one hand, chose a standard $S$-countably infinite $S$-sorted set of variables $V^{S} = (\{v^{s}_{n}\mid n\in\mathbb{N}\})_{s\in S}$ (one could take as $V^{S}$ the $S$-sorted set $(\{s\}\times\mathbb{N})_{s\in S}$) and, on the other hand, assume that all involved $S$-sorted sets of variables $X$ are finite subsets of $V^{S}$, i.e., $\mathrm{supp}_{S}(X)$ is finite and, for every $s\in \mathrm{supp}_{S}(X)$, $X_{s}$ is finite (therefore, $X$ will be isomorphic to $\downarrow\!\! w$, for some $w\in S^{\star}$).
%
% \begin{definition}
% We will say that a set $\boldsymbol{\mathcal{K}}$ of partial $\Sigma$-algebras is a \emph{finitary} $\mathrm{E}$-, \emph{finitary} $\mathrm{ECE}$-, or \emph{finitary} $\QE$-\emph{variety}, respectively, if there exists a set $\mathcal{E}$ of finitary $\mathrm{E}$-, $\mathrm{ECE}$-, or $\QE$-equations, respectively, such that $\boldsymbol{\mathcal{K}} = \mathrm{Mod}_{\Sigma}(\mathcal{E})$.
% \end{definition}
%
% \begin{remark}
% The above finitary varieties are also referred to as \emph{finitary} $\mathrm{E}$-, \emph{finitary} $\mathrm{ECE}$-, or \emph{finitary} $\QE$-\emph{equational class}, respectively. Moreover, for the characterization of the aforementioned varieties by closure properties with respect to algebraic constructions we refer the reader to~\cite{pb86}. 
% \end{remark}
%
% \begin{remark}
% Categories (in a fixed Grothendieck universe) constitute a finitary $\mathrm{ECE}$-variety and not only a finitary $\QE$-variety (see~\cite{pb02}).
% \end{remark}
%
% %\begin{definition}\hfill
% %\begin{enumerate}
% %\item Sea $\mathcal{E}\subseteq \mathrm{Eqt}(\Sigma)$.  La \emph{clase ecuacional
% %      determinada} por $\mathcal{E}$, $\mathrm{Mod}_{\Sigma}(\mathcal{E})$, consta de
% %      todas las $\Sigma$-álgebras $\mathbf{A}$ que satisfacen todas las
% %      ecuaciones de $\mathcal{E}$, i.e.,%
% %      $$
% %      \mathrm{Mod}_{\Sigma}(\mathcal{E})= \{ \mathbf{A} \in \mathsf{Alg}(\Sigma) \mid
% %      \mathbf{A} \models \mathcal{E} \}
% %     $$
% %
% %\item Sea $\mathcal{K}\subseteq\mathsf{Alg}(\Sigma)$.  La \emph{teoría ecuacional
% %      determinada} por $\mathcal{K}$, $\mathrm{Th}_{\Sigma}(\mathcal{K})$,
% %      consta de todas las $\Sigma$-ecuaciones válidas en todas las
% %      $\Sigma$-álgebras de $\mathcal{K}$, i.e.,%
% %      $$
% %      \mathrm{Th}_{\Sigma}(\mathcal{K})= \bigl(
% %      \{
% %      E\in \mathrm{Eqt}(\Sigma)_{X,s} \mid \forall\mathbf{A}\in\mathcal{K},\,
% %      \mathbf{A}\models E
% %      \}
% %      \bigr) _{(X,s)\in \boldsymbol{\mathcal{U}}^{S}\bprod S}
% %      $$
% %\end{enumerate}
% %\end{definition}
% %
% %\begin{proposition}
% %Sea $\Sigma$ una signatura algebraica, $\mathcal{E}$, $\mathcal{E'}$ dos
% %conjuntos de $\Sigma$-ecuaciones y $\mathcal{K}$, $\mathcal{K'}$
% %dos conjuntos de $\Sigma$-álgebras. Entonces se cumple que
% %\begin{enumerate}
% %\item Si $\mathcal{E}\subseteq\mathcal{E}'$,
% %      $\mathrm{Mod}_{\Sigma}(\mathcal{E'})\subseteq\mathrm{Mod}_{\Sigma}(\mathcal{E})$.
% %
% %\item Si $\mathcal{K}\subseteq \mathcal{K}'$,
% %      $\mathrm{Th}_{\Sigma}(\mathcal{K}')\subseteq\mathrm{Th}_{\Sigma}(\mathcal{K})$.
% %
% %\item $\mathcal{E}\subseteq\mathrm{Th}_{\Sigma}(\Mod_{\Sigma}(\mathcal{E}))$ y
% %      $\mathcal{K}\subseteq\mathrm{Mod}_{\Sigma}(\mathrm{Th}_{\Sigma}(\mathcal{K}))$.
% %\end{enumerate}
% %Por consiguiente, para cada signatura algebraica $\Sigma$, las
% %funciones $\mathrm{Th}_{\Sigma}$ y $\mathrm{Mod}_{\Sigma}$ forman una
% %\emph{conexión de Galois contravariante}.
% %\end{proposition}
% %
% %Para las categorías asociadas a los retículos de clases de
% %$\Sigma$-álgebras y familias de $\Sigma$-ecuaciones, se tiene la
% %adjunción
% %$$
% %\xymatrix@=5pc{ \mathbf{\Sub}(\mathsf{Alg}(\Sigma))^{\opp}
% %\ar@<1.5ex>[r]^{\mathrm{Th}_{\Sigma}} \ar@{}[r]|{\uadj} &
% %\mathbf{\Sub}(\mathrm{Eqt}(\Sigma)) \ar@<1.5ex>[l]^{\Mod_{\Sigma}} }
% %$$
% %en donde, para cada clase $\mathcal{K}$ de $\Sigma$-álgebras y
% %cada familia $\mathcal{E}$
% %de $\Sigma$-ecuaciones, se cumple que %
% %$$
% %\mathcal{K}\subseteq\mathrm{Mod}_{\Sigma}(\mathcal{E}) \text{ si y sólo si }
% %\mathcal{E}\subseteq \mathrm{Th}_{\Sigma}(\mathcal{K})
% %$$
%
% %\begin{definition}\hfill
% %\begin{enumerate}
% %\item El operador clausura sobre $\Eq(\Sigma)$ asociado a la
% %      conexión de Galois, $\Th_{\Sigma}\comp \Mod_{\Sigma}$,
% %      se denota como $\Cn_{\Sigma}$.  Los cerrados de
% %      $\Cn_{\Sigma}$ se denominan \emph{teorías ecuacionales}. Si
% %      $\ec{E}$ es una familia de $\Sigma$-ecuaciones y $E$ una
% %      $\Sigma$-ecuación, entonces $E$ es una
% %      \emph{consecuencia semántica} de \ec{E}, $\ec{E}\cons E$, si
% %      $\Mod_{\Sigma}(\ec{E})\incl\Mod_{\Sigma}(E)$, i.e., si
% %      $E\in\Cn_{\Sigma}(\ec{E})$.
% %
% %\item El operador clausura sobre $\mathsf{Alg}(\Sigma)$ asociado a la
% %      conexión de Galois, $\Mod_{\Sigma}\comp \Th_{\Sigma}$,
% %      se denota como $\Ec_{\Sigma}$.  Los cerrados de
% %      $\Ec_{\Sigma}$ se denominan \emph{clases ecuacionales}.  Si
% %      $\mathcal{K}$ es una clase de $\Sigma$-álgebras y $\mathbf{A}$ una
% %      $\Sigma$-álgebra, entonces $\mathbf{A}$ está en la clase ecuacional
% %      determinada por $\mathcal{K}$, $\mathcal{K}\models \mathbf{A}$, si
% %      $\Th_{\Sigma}(\mathcal{K})\incl\Th_{\Sigma}(\mathbf{A})$,
% %      i.e., si $\mathbf{A}\in\Ec_{\Sigma}(\mathcal{K})$.
% %\end{enumerate}
% %\end{definition}
% %
% %Las restricciones del operador $\Mod_{\Sigma}$ a
% %$\Eq(\Sigma)_{X}$, $\Eq_{S\text{-}\mathrm{f}}(\Sigma)$ y
% %$\Eq_{\mathrm{f}}(\Sigma)$ se denotan, respectivamente, como
% %$\Mod_{\Sigma,X}$, $\Mod_{\Sigma,S\text{-}\mathrm{f}}$ y
% %$\Mod_{\Sigma,\mathrm{f}}$. Asimismo, las correstricciones del
% %operador $\Th_{\Sigma}$ se denotan mediante $\Th_{\Sigma,X}$,
% %$\Th_{\Sigma,S\text{-}\mathrm{f}}$ y $\Th_{\Sigma,\mathrm{f}}$.  Todos
% %ellos forman, dos a dos, conexiones de Galois contravariantes que
% %determinan operadores clausura, y que se denotan con los
% %subíndices correspondientes.
% %
% %Una clase ecuacional $\mathcal{K}$ es $S$-\emph{finitaria} si
% %$\mathcal{K}$ es un cerrado de $\Ec_{\Sigma,S\text{-}\mathrm{f}}$,
% %\emph{finitaria} si es un cerrado de $\Ec_{\Sigma,\mathrm{f}}$ y
% %\emph{sobre} $X$ si $\mathcal{K}$ es un cerrado de
% %$\Ec_{\Sigma,X}$.
% %
% %\begin{proposition}
% %Sea $\mathcal{K}$ una clase de $\Sigma$-álgebras y \ec{E} una
% %familia de
% %$\Sigma$-ecua\-cio\-nes.  Entonces se cumple que:%
% %\begin{enumerate}
% %\item Para cada $(X,s)\in \boldsymbol{\mathcal{U}}^{S}\bprod S$,
% %$\Th_{\Sigma,X}(\mathcal{K})_{s}={\equiv^{\mathcal{K}}_{X,s}}$
% %
% %\item  $
% %\Th_{\Sigma}(\mathcal{K})=(\Th_{\Sigma,X}(\mathcal{K})_{s})_{
% %(X,s)\in \boldsymbol{\mathcal{U}}^{S}\bprod S} $
% %
% %\item $
% %\Mod_{\Sigma}(\ec{E})=\inter
% %(\Mod_{\Sigma,X}(\ec{E}_{X}))_{X\in\boldsymbol{\mathcal{U}}^{s}} $
% %\end{enumerate}
% %\DQEEqd
% %\end{proposition}
% %
% %En el álgebra homogénea, cada clase ecuacional lo es sobre un
% %conjunto infinito numerable, arbitrario pero fijo.  Para el
% %álgebra heterogénea, en cambio, existen clases ecuacionales que no
% %lo son sobre ningún $S$-conjunto de variables fijo.  Esto es
% %debido a que si dos $S$-conjuntos $X$ e $Y$ son tales que $X\incl
% %Y$, entonces se cumple que $\Eq(\Sigma,X)\incl\Eq(\Sigma,Y)$, pero
% %\emph{no toda $(\Sigma,Y)$-ecuación, con variables en $X$, válida
% %en un álgebra sigue siendo válida cuando se considera como una
% %$(\Sigma,X)$-ecuación} (la implicación inversa se cumple siempre),
% %una condición que sí se cumple para las ecuaciones homogéneas.
% %
% %\begin{proposition}
% %Sea $\mathbf{A}$ una $\Sigma$-álgebra, y $X$ e $Y$ dos
% %$S$-conjuntos tales que $X\subseteq Y$.  Si
% %$\mathbf{A}\models_{X,s}(P,Q)$, entonces
% %$\mathbf{A}\models_{Y,s}(P,Q)$. \DQEEqd
% %\end{proposition}
% %
% %La recíproca de la proposición anterior no es cierta en general.
% %Si tenemos que
% %$\mathrm{supp}_{S}(X)\subseteq\mathrm{supp}_{S}(A)\subset\mathrm{supp}_{S}(Y)$,
% %entonces $\mathbf{A}\models_{Y,s}(P,Q)$ se cumple vacuamente, sea
% %cual sea la estructura algebraica de $\mathbf{A}$, porque no
% %existe ninguna valoración de $Y$ en $A$.
% %
% %Sea $S = \{s,t\}$ un conjunto de sorts tal que $s\neq t$, $X$ el $S$-sorted set definido como: $X_{s} = \{x_{s,0},x_{s,1}\}$, $X_{t} = \varnothing$, e $Y$ el $S$-conjunto definido como: $Y_{s} = \{x_{s,0},x_{s,1}\}$, $Y_{t} = \{y_{t,0}\}$. Entonces, dada una $\Sigma$-álgebra $\mathbf{A}$, se cumple que: (1) la ecuación $x_{s,0} = x_{s,1}$, considerada respecto de las variables $X$, es válida en $\mathbf{A}$ si, y sólo si, $\mathrm{card}(A_{s})\leq 1$ y (2) la ecuación $x_{s,0} = x_{s,1}$, considerada respecto de las variables $Y$, es válida en $\mathbf{A}$ si, y sólo si, $\mathrm{card}(A_{s})\leq 1$ o $A_{t} = \varnothing$.
% %
% %\begin{proposition}
% %Sea $X$ un $S$-conjunto y $\mathbf{A}$ una $\Sigma$-álgebra.
% %Entonces las siguientes proposiciones son equivalentes:
% %\begin{enumerate}
% % \item $\mathbf{A}\models_{X,s}(P,Q)$,
% %
% %\item $\mathrm{supp}_{S}(X)-\mathrm{supp}_{S}(A)\neq \vacio$ o se cumple
% %      que $\mathbf{A}\models_{\mathrm{var}(P,Q),s}(P,Q)$, siendo
% %      $\mathrm{var}(P,Q)=\mathrm{var}(P)\cup\mathrm{var}(Q)$.
% %\end{enumerate}
% %\end{proposition}
% %
% %\begin{corollary}\label{SatisfSoportes}
% %Sean $X$ e $Y$ dos $S$-conjuntos tales que $\mathrm{supp}_{S}(X) =
% %\mathrm{supp}_{S}(Y)$. Entonces para cada $s\in S$, cada
% %$(P,Q)\in\mathrm{Eq}(\Sigma)_{X\cap Y,s}$, y cada $\Sigma$-álgebra
% %$\mathbf{A}$, se cumple que $\mathbf{A}\models_{X,s}(P,Q)$
% %exactamente si $\mathbf{A}\models_{Y,s}(P,Q)$.  \DQEEqd
% %\end{corollary}
%
%
% %Notación.
% %
% %Sean $\mathbf{A}$ y $\mathbf{X}$ dos categorías, $F$ un functor de $\mathbf{A}$ en $\mathbf{X}$ y $x$ un objeto de $\mathbf{X}$. Entonces denotamos por $(x\!\downarrow\! F)$ la categoría cuyos objetos son los pares ordenados $(f,a)$ en los que $a$ es un objeto de $\mathbf{A}$ y $f$ un morfismo de $\mathbf{X}$ de $x$ en $F(a)$; y en la que los morfismos de $(f,a)$ en $(g,b)$ son los triplos ordenados $((f,a),h,(g,b))$ en los que $h$ es un morfismo de $\mathbf{A}$ de $a$ en $b$ tal que $F(h)\circ f = g$. Dada la anterior situación, denotamos por $\mathrm{1st}(x\!\downarrow\! F)$ el conjunto de las primeras coordenadas del conjunto de los objetos de $(x\!\downarrow\! F)$.
% %
% %Cuando consideremos el functor de olvido $\mathrm{G}_{\boldsymbol{\mathcal{E}}^{\mathrm{c},\mathcal{A}}}$ de $\mathsf{PAlg}(\boldsymbol{\mathcal{E}}^{\mathrm{c},\mathcal{A}})$ en $\mathsf{PAlg}(\Sigma^{\mathrm{c},\mathcal{A}})$ y la many-sorted partial $\Sigma^{\mathrm{c},\mathcal{A}}$-algebra $\mathbf{T}^{\mathrm{c},\mathcal{A}}_{\Sigma}(X\amalg\coprod_{\beta}Q)^{2}$, escribiremos
% %
% %$\bigcap_{f\in \mathrm{1st}(\mathbf{T}^{\mathrm{c},\mathcal{A}}_{\Sigma}(X\amalg\coprod_{\beta}Q)^{2}\downarrow \mathrm{G}_{\boldsymbol{\mathcal{E}}^{\mathrm{c},\mathcal{A}}})}\mathbf{Sch}(f)$
%
% %%%%%%%%%%%%%%%%%%%%
% %%%%%%%%%%%%%%%%%%%%
%
% We next state the fundamental theorem on the existence of universal solutions for sets of partial algebras closed under subalgebras and products. But before doing that, for partial $\Sigma$-algebras $\mathbf{A}$ and $\mathbf{B}$, we make the convention that $f\colon \mathbf{B}\epi\mathbf{A}$ means that $f$ is a surjective homomorphism from $\mathbf{B}$ to $\mathbf{A}$ and that $f\colon \mathbf{A}\mon\mathbf{B}$ means that $f$ is a a closed and injective homomorphism from $\mathbf{A}$ to $\mathbf{B}$.
%
% \begin{definition}
% Let $\mathcal{K}$ be a set of partial $\Sigma$-algebras. Then
% \begin{enumerate}
% \item $\mathrm{I}(\mathcal{K}) = \{\mathbf{A}\in \mathrm{Alg}(\Sigma)\mid \exists\,\mathbf{B}\in \mathcal{K}\,(\mathbf{A}\cong\mathbf{B})\}$, i.e., $\mathrm{I}(\mathcal{K})$ is the set of all isomorphic copies of partial $\Sigma$-algebras in $\mathcal{K}$.
%
% \item $\mathrm{H}(\mathcal{K}) = \{\mathbf{A}\in \mathrm{Alg}(\Sigma)\mid \exists\,\mathbf{B}\in \mathcal{K}\,\exists\,f\colon \mathbf{B}\epi\mathbf{A}\}$, i.e., $\mathrm{H}(\mathcal{K})$ is the set of all homomorphic images of partial $\Sigma$-algebras in $\mathcal{K}$.
%
% \item $\mathrm{S}(\mathcal{K}) = \{\mathbf{A}\in \mathrm{Alg}(\Sigma)\mid \exists\,\mathbf{B}\in \mathcal{K}\,\exists\,f\colon \mathbf{A}\mon\mathbf{B}\}$, i.e., $\mathrm{S}(\mathcal{K})$ is the set of all isomorphic copies of subalgebras of partial $\Sigma$-algebras in $\mathcal{K}$.
%
% \item $\mathrm{P}(\mathcal{K}) = \{\mathbf{A}\in \mathrm{Alg}(\Sigma)\mid \exists\, I\in \boldsymbol{\mathcal{U}}\, \exists\, (\mathbf{A}^{i})_{i\in I}\in \mathcal{K}^{I}\,(\mathbf{A}\cong \prod_{i\in I}\mathbf{A}^{i})\}$, i.e., $\mathrm{P}(\mathcal{K})$ is the set of all partial $\Sigma$-algebras which are  isomorphic to products of families, indexed by some $\boldsymbol{\mathcal{U}}$-small set, of partial $\Sigma$-algebras in $\mathcal{K}$.
% \item We will say that $\mathcal{K}$ is (1) \emph{abstract} if $\mathrm{I}(\mathcal{K})\subseteq \mathcal{K}$, (2) \emph{closed under homomorphic images} if $\mathrm{H}(\mathcal{K})\subseteq \mathcal{K}$, (3) \emph{closed under subalgebras} if $\mathrm{S}(\mathcal{K})\subseteq \mathcal{K}$, and (4) \emph{closed under products} if $\mathrm{P}(\mathcal{K})\subseteq \mathcal{K}$. 
% %and (5) \emph{quasiprimitive} if $\mathcal{K} = \mathrm{S}(\mathrm{P}(\mathcal{K})) = \mathrm{I}(\mathrm{S}(\mathrm{P}(\mathcal{K})))$.
% \end{enumerate}
% \end{definition} 
%
% \begin{lemma}
% Let $\mathbf{A}$ be a partial $\Sigma$-algebra. Then:
% \begin{enumerate}
% \item The set of all $A$-generated relative subalgebras of $\mathbf{F}_{\Sigma}(\mathbf{A})$ is a closure system on $\mathrm{F}_{\Sigma}(\mathbf{A})$ with smallest element $\mathbf{A}$ and greatest element $\mathbf{F}_{\Sigma}(\mathbf{A})$.
% \item The set of all closed congruences on $A$-generated relative subalgebras of $\mathbf{F}_{\Sigma}(\mathbf{A})$ is a closure system on $\mathrm{F}_{\Sigma}(\mathbf{A})^{2}$ with smallest element $\Delta_{\mathbf{A}}$ and greatest element $\nabla_{{\mathbf{F}_{\Sigma}(\mathbf{A})}}$.
% \item For every nonempty family $(\Phi^{i})_{i\in I}$ of closed congruences on $A$-generated relative subalgebras of $\mathbf{F}_{\Sigma}(\mathbf{A})$, we have that
% $$
% \textstyle
% \mathrm{Fld}(\bigcap_{i\in I}\Phi^{i}) = \bigcap_{i\in I}\mathrm{Fld}(\Phi^{i}).
% $$
% \item The set of all closed congruences on $A$-generated relative subalgebras of $\mathbf{F}_{\Sigma}(\mathbf{A})$, ordered by inclusion, is upwards directed.
% \end{enumerate}
% \end{lemma}
%
% \begin{corollary}
% Let $\mathbf{A}$ be a partial $\Sigma$-algebra, $(\mathbf{B}^{i})_{i\in I}$ an $I$-indexed family of partial $\Sigma$-algebras, $(f^{i})_{i\in I}\in \prod_{i\in I}\mathrm{Hom}(\mathbf{A},\mathbf{B}^{i})$ and 
% $\left< f^{i}\right>_{i\in I}$ the unique homomorphism from $\mathbf{A}$ to 
% $\prod_{i\in I}\mathbf{B}^{i}$ such that, for every $i\in I$, 
% $\mathrm{pr}^{i}\circ \left< f^{i}\right>_{i\in I} = f^{i}$. Then
% \begin{enumerate}
% \item $\mathbf{Sch}(\left< f^{i}\right>_{i\in I}) = \bigcap_{i\in I}\mathbf{Sch}(f^{i})$.
% \item $\mathrm{SKer}(\left< f^{i}\right>_{i\in I}) = \bigcap_{i\in I}\mathrm{SKer}(f^{i})$.
% \end{enumerate}
% \end{corollary}
%
% \begin{lemma}\label{LCardinal}
% For every $S$-cardinal $(\mathfrak{n}_{s})_{s\in S}$, there exists a cardinal $\mathfrak{m}$ such that, for every partial $\Sigma$-algebra $\mathbf{B}$ and every $S$-sorted set $Y$, if $Y\subseteq B$, $\mathrm{Sg}_{\mathbf{B}}(Y) = B$ and $(\mathrm{card}(Y_{s}))_{s\in S}\leq (\mathfrak{n}_{s})_{s\in S}$, then $(\mathrm{card}(B_{s}))_{s\in S}\leq (\mathfrak{m})_{s\in S}$, i.e., for every $s\in S$, $\mathrm{card}(B_{s})\leq \mathfrak{m}$.
% \end{lemma}
%
% \begin{proof}
% It sufices to take $\mathfrak{m} = \mathrm{max}(\{\mathrm{card}(\Sigma),\aleph_{0}\}\cup\{\mathfrak{n}_{s}\mid s\in S\})$ (by virtue of Corollary~\ref{CPAlgInitGen} and the recursive definition of the free $\Sigma$-algebra on $Y$).
% \end{proof}
%
%
%
% %%%%%%%%%%%%%%%%%%%
% %%%%%%%%%%%%%%%%%%%
%
%
% \begin{theorem}\label{TFreeAdj}
% Let $\mathcal{K}$ be an abstract set of partial $\Sigma$-algebras closed under subalgebras and products. Then the inclusion functor $\mathrm{In}_{\boldsymbol{\mathcal{K}}}$ from $\boldsymbol{\mathcal{K}}$, the full subcategory of $\mathsf{PAlg}(\Sigma)$ determined by $\mathcal{K}$, to $\mathsf{PAlg}(\Sigma)$ has a left adjoint $\mathbf{T}_{\boldsymbol{\mathcal{K}}}$.
% \end{theorem}
%
% \begin{proof}
%
% The proof will consists in showing that for every partial $\Sigma$-algebra $\mathbf{A}$ there exists a universal morphism from $\mathbf{A}$ to $\mathrm{In}_{\boldsymbol{\mathcal{K}}}$.
%
% \textsf{First proof (via Schmidt)}.\hfill
%
% Let $\mathbf{A}$ be a partial $\Sigma$-algebra and $\equiv^{\mathcal{K}}_{\mathbf{A}}$ be
% $\bigcap_{f\in\bigcup_{\mathbf{K}\in\mathcal{K}}\mathrm{Hom}(\mathbf{A},\mathbf{K})}\mathrm{SKer}(f)$ (recall that $\mathrm{SKer}(f) = \mathrm{Ker}(f^{\mathrm{Sch}})$ and that it is a closed congruence on the $A$-generated relative subalgebras $\mathbf{Sch}(f)$ of $\mathbf{F}_{\Sigma}(\mathbf{A})$). For every
% $f\in\bigcup_{\mathbf{K}\in\mathcal{K}}\mathrm{Hom}(\mathbf{A},\mathbf{K})$, $\mathrm{Fld}(\mathrm{SKer}(f))$, the field of $\mathrm{SKer}(f)$, is $\mathrm{Sch}(f)$ (recall that, for an $S$-sorted equivalence relation $\Phi$ on an $S$-sorted set $Z$, since $\Phi$ is symmetric, $\mathrm{Fld}(\Phi) = \Dom(\Phi) = \mathrm{Im}(\Phi)$, and, since $\Phi$ is reflexive, $\mathrm{Fld}(\Phi) = Z$). Moreover
% $$
% \textstyle
% \mathrm{Fld}
% \left(
% \bigcap_{f\in\bigcup_{\mathbf{K}\in\mathcal{K}}\mathrm{Hom}(\mathbf{A},\mathbf{K})}\mathrm{SKer}(f)
% \right) =
% \bigcap_{f\in\bigcup_{\mathbf{K}\in\mathcal{K}}\mathrm{Hom}(\mathbf{A},\mathbf{K})}\mathrm{Fld}(\mathrm{SKer}(f)).
% $$
% Therefore
% $$
% \textstyle
% \mathrm{Fld}\left(
% \bigcap_{f\in\bigcup_{\mathbf{K}\in\mathcal{K}}\mathrm{Hom}(\mathbf{A},\mathbf{K})}\mathrm{SKer}(f)
% \right) =
% \bigcap_{f\in\bigcup_{\mathbf{K}\in\mathcal{K}}\mathrm{Hom}(\mathbf{A},\mathbf{K})}\mathrm{Sch}(f).
% $$
% From now on, $\mathbf{T}_{\boldsymbol{\mathcal{K}}}(\mathbf{A})$ stands for $\bigcap_{f\in\bigcup_{\mathbf{K}\in\mathcal{K}}\mathrm{Hom}(\mathbf{A},\mathbf{K})}\mathbf{Sch}(f)/{\equiv^{\mathcal{K}}_{\mathbf{A}}}$. Then $\mathbf{T}_{\boldsymbol{\mathcal{K}}}(\mathbf{A})$ together with the restriction to $\mathbf{A}$ of the  projection from $\bigcap_{f\in\bigcup_{\mathbf{K}\in\mathcal{K}}\mathrm{Hom}(\mathbf{A},\mathbf{K})}\mathbf{Sch}(f)$ to $\mathbf{T}_{\boldsymbol{\mathcal{K}}}(\mathbf{A})$ is a universal morphism from $\mathbf{A}$ to $\mathrm{In}_{\boldsymbol{\mathcal{K}}}$. On the set $\bigcup_{\mathbf{K}\in\mathcal{K}}\mathrm{Hom}(\mathbf{A},\mathbf{K})$ let $\equiv$ be the equivalence relation defined, for every $f$, $g\in \bigcup_{\mathbf{K}\in\mathcal{K}}\mathrm{Hom}(\mathbf{A},\mathbf{K})$ as:
% $$
% f\equiv g\,\,\, \text{if and only if}\,\,\, \mathrm{SKer}(f) = \mathrm{SKer}(g),
% $$
% and let $\mathcal{H}$ be a subset of $\bigcup_{\mathbf{K}\in\mathcal{K}}\mathrm{Hom}(\mathbf{A},\mathbf{K})$ such that, for every $f\in \bigcup_{\mathbf{K}\in\mathcal{K}}\mathrm{Hom}(\mathbf{A},\mathbf{K})$, $\mathrm{card}(\mathcal{H}\cap [f]_{\equiv}) = 1$. Then $\mathbf{T}_{\boldsymbol{\mathcal{K}}}(\mathbf{A})$ is isomorphic to a subalgebra of $\prod_{f\in \mathcal{H}}\mathbf{Sch}(f)/ \mathrm{SKer}(f)$.
%
% \textsf{Second proof (classical)}.\hfill
%
% Let $\mathbf{A}$ be a partial $\Sigma$-algebra with $(\mathrm{card}(A_{s}))_{s\in S} = (\mathfrak{n}_{s})_{s\in S}$. Then, by Lemma~\ref{LCardinal}, there exists a cardinal $\mathfrak{m}_{\mathbf{A}}$ such that $(\mathrm{card}(B_{s}))_{s\in S}\leq (\mathfrak{m}_{\mathbf{A}})_{s\in S}$, for all partial $\Sigma$-algebras $\mathbf{B}$ generated by $S$-sorted sets $X$ such that $(\mathrm{card}(X_{s}))_{s\in S}\leq (\mathfrak{n}_{s})_{s\in S}$.
%
% Let $\mathcal{K}_{\mathfrak{m}_{\mathbf{A}}}$ be the (small) set $\{\mathbf{K}\in \mathcal{K}\mid (\mathrm{card}(K_{s}))_{s\in S}\leq (\mathfrak{m}_{\mathbf{A}})_{s\in S}\}$. For $\mathbf{K}\in \mathcal{K}_{\mathfrak{m}_{\mathbf{A}}}$, let $(\mathbf{K}^{\mathrm{Hom}(\mathbf{A},\mathbf{K})},(\mathrm{pr}^{f})_{f\in \mathrm{Hom}(\mathbf{A},\mathbf{K})})$ be the product of $(\mathbf{K})_{f\in \mathrm{Hom}(\mathbf{A},\mathbf{K})}$. By the assumption on $\mathcal{K}$, $\mathbf{K}^{\mathrm{Hom}(\mathbf{A},\mathbf{K})}\in \mathcal{K}$. Then we denote by $\mathrm{p}^{\mathbf{K}}$ the unique homomorphism $\left<f\right>_{f\in\mathrm{Hom}(\mathbf{A},\mathbf{K})}$ from $\mathbf{A}$ to $\mathbf{K}^{\mathrm{Hom}(\mathbf{A},\mathbf{K})}$ such that, for every $f\in\mathrm{Hom}(\mathbf{A},\mathbf{K})$, $\mathrm{pr}^{f}\circ \left<f\right>_{f\in\mathrm{Hom}(\mathbf{A},\mathbf{K})} = f$.
%
% Next, let $(\prod_{\mathbf{K}\in \mathcal{K}_{\mathfrak{m}_{\mathbf{A}}}}\mathbf{K}^{\mathrm{Hom}(\mathbf{A},\mathbf{K})},(\mathrm{pr}^{\mathbf{K}})_{\mathbf{K}\in \mathcal{K}_{\mathfrak{m}_{\mathbf{A}}}})$ be the product of $(\mathbf{K}^{\mathrm{Hom}(\mathbf{A},\mathbf{K})})_{\mathbf{K}\in \mathcal{K}_{\mathfrak{m}_{\mathbf{A}}}}$.  By the assumption on $\mathcal{K}$, $\prod_{\mathbf{K}\in \mathcal{K}_{\mathfrak{m}_{\mathbf{A}}}}\mathbf{K}^{\mathrm{Hom}(\mathbf{A},\mathbf{K})}\in \mathcal{K}$. Then we denote by $\mathrm{q}$ the unique homomorphism $\left<\mathrm{p}^{\mathbf{K}}\right>_{\mathbf{K}\in \mathcal{K}_{\mathfrak{m}_{\mathbf{A}}}}$ from $\mathbf{A}$ to
% $\prod_{\mathbf{K}\in \mathcal{K}_{\mathfrak{m}_{\mathbf{A}}}}\mathbf{K}^{\mathrm{Hom}(\mathbf{A},\mathbf{K})}$ such that, for every $\mathbf{K}\in\mathcal{K}_{\mathfrak{m}_{\mathbf{A}}}$, $\mathrm{pr}^{\mathbf{K}}\circ \left<\mathrm{p}^{\mathbf{K}}\right>_{\mathbf{K}\in \mathcal{K}_{\mathfrak{m}_{\mathbf{A}}}} = \mathrm{p}^{\mathbf{K}}$. Let $\mathbf{T}_{\boldsymbol{\mathcal{K}}}(\mathbf{A})$ be the partial $\Sigma$-algebra induced by the subalgebra of $\prod_{\mathbf{K}\in \mathcal{K}_{\mathfrak{m}_{\mathbf{A}}}}\mathbf{K}^{\mathrm{Hom}(\mathbf{A},\mathbf{K})}$ generated by $\mathrm{q}[A]$, and $\theta^{\mathbf{A}}$ the unique epimorphism from $\mathbf{A}$ to $\mathbf{T}_{\boldsymbol{\mathcal{K}}}(\mathbf{A})$ such that $\mathrm{q} = \mathrm{in}^{\mathbf{T}_{\boldsymbol{\mathcal{K}}}(\mathbf{A})}\circ \theta^{\mathbf{A}}$. By the assumption on $\mathcal{K}$, $\mathbf{T}_{\boldsymbol{\mathcal{K}}}(\mathbf{A})\in \mathcal{K}$.
%
% Let $\mathbf{B}$ be a partial $\Sigma$-algebra belonging to $\mathcal{K}$ and $g$ a homomorphism from $\mathbf{A}$ to $\mathbf{B}$. Let $\mathbf{Sg}_{\mathbf{B}}(g[A])$ be the partial $\Sigma$-algebra induced by the subalgebra $\mathrm{Sg}_{\mathbf{B}}(g[A])$ of $\mathbf{B}$ generated by $g[A]$, and $g^{\mathrm{e}}$ the unique epimorphism from $\mathbf{A}$ to $\mathbf{Sg}_{\mathbf{B}}(g[A])$ such that
% $g = \mathrm{in}^{\mathbf{Sg}_{\mathbf{B}}(g[A])}\circ g^{\mathrm{e}}$.  By the assumption on $\mathcal{K}$, $\mathbf{Sg}_{\mathbf{B}}(g[A])\in \mathcal{K}$. Moreover, $(\mathrm{card}(g_{s}[A_{s}]))_{s\in S}\leq (\mathrm{card}(A_{s}))_{s\in S}$. Thus $(\mathrm{card}(\mathrm{Sg}_{\mathbf{B}}(g[A])_{s})_{s\in S}\leq (\mathfrak{m}_{\mathbf{A}})_{s\in S}$. Therefore $\mathbf{Sg}_{\mathbf{B}}(g[A])\in \mathcal{K}_{\mathfrak{m}_{\mathbf{A}}}$. Then
% $$
% \mathrm{in}^{\mathbf{Sg}_{\mathbf{B}}(g[A])}\circ\mathrm{pr}^{g^{\mathrm{e}}}\circ\mathrm{pr}^{\mathbf{Sg}_{\mathbf{B}}(g[A])}\circ\mathrm{in}^{\mathbf{T}_{\boldsymbol{\mathcal{K}}}(\mathbf{A})},
% $$
% denoted by $g^{\mathsf{p}}$, is the unique homomorphism from $\mathbf{T}_{\boldsymbol{\mathcal{K}}}(\mathbf{A})$ to $\mathbf{B}$ such that  $$g = g^{\mathsf{p}}\circ \theta^{\mathbf{A}}.$$
%
% This fact is depicted in Figure~\ref{FFreeK}.
% \end{proof}
%
% \begin{figure}
% $$\xymatrix@C=40pt@R=40pt{
% {} & \mathbf{B} & {} \\
% \mathbf{Sg}_{\mathbf{B}}(g[A])
% \ar[ru]^-{\mathrm{in}^{\mathbf{Sg}_{\mathbf{B}}(g[A])}} &
% \mathbf{A}
% \ar[u]^-{g}
% \ar[l]_-{g^{\mathrm{e}}}
% \ar[r]^-{\theta^{\mathbf{A}}}
% \ar[rd]^-{\mathrm{q}}
% \ar[d]^-{\mathrm{p}^{\mathbf{Sg}_{\mathbf{B}}(g[A])}} &
% \mathbf{T}_{\boldsymbol{\mathcal{K}}}(\mathbf{A})
% \ar[ul]_-{g^{\mathsf{p}}}
% \ar[d]^-{\mathrm{in}^{\mathbf{T}_{\boldsymbol{\mathcal{K}}}(\mathbf{A})}} \\
% {} &
% \mathbf{Sg}_{\mathbf{B}}(g[A])^{\mathrm{Hom}(\mathbf{A},\mathbf{Sg}_{\mathbf{B}}(g[A]))}
% \ar[lu]^-{\mathrm{pr}^{g^{\mathrm{e}}}}
% &
% \prod_{\mathbf{K}\in \mathcal{K}_{\mathfrak{m}_{\mathbf{A}}}}\mathbf{K}^{\mathrm{Hom}(\mathbf{A},\mathbf{K})}
% \ar[l]^-{\mathrm{pr}^{\mathbf{Sg}_{\mathbf{B}}(g[A])}}
%    }
% $$
% \caption{The free $\boldsymbol{\mathcal{K}}$-algebra.}
% \label{FFreeK}
% \end{figure}
%
% \begin{remark}
% For $\theta^{\mathbf{A}}$ we have $\mathrm{Ker}(\theta^{\mathbf{A}}) = \bigcap_{\mathbf{K\in \mathcal{K}}}\left(
% \bigcap_{f\in\mathrm{Hom}(\mathbf{A},\mathbf{K})}\mathrm{Ker}(f)\right)$.
% \end{remark}
%
% \begin{remark}
% The second proof in the just stated theorem is interesting because it, essentially, is the model of the proof of the General Adjoint Functor of P. Freyd (see, e.g., \cite{em76} for the proof of Freyd's theorem). 
%
% It seems to us appropriate to point out that in order to prove, via Freyd's theorem, that the inclusion functor $\mathrm{In}_{\boldsymbol{\mathcal{K}}}$ from $\boldsymbol{\mathcal{K}}$ to $\mathsf{PAlg}(\Sigma)$ has a left adjoint, it is sufficient to verify the following conditions: 
% \begin{enumerate}
% \item $\boldsymbol{\mathcal{K}}$ is $\boldsymbol{\mathcal{U}}$-locally small (i.e., has $\boldsymbol{\mathcal{U}}$-small hom sets),
% \item $\boldsymbol{\mathcal{K}}$ has $\boldsymbol{\mathcal{U}}$-small limits (i.e., has $\boldsymbol{\mathcal{U}}$-small products and equalizers),
% \item $\mathrm{In}_{\boldsymbol{\mathcal{K}}}$ preserves $\boldsymbol{\mathcal{U}}$-small products and equalizers and
% \item $\mathrm{In}_{\boldsymbol{\mathcal{K}}}$ sastisfies the solution set condition at $\mathbf{K}$, for every $\mathbf{K}$ in $\mathcal{K}$.
% \end{enumerate}
% The last condition, which is the most convoluted of them all, is settled in the above mentioned second proof, the remaining conditions are obviously satisfied.
% \end{remark}
%
% \begin{corollary}
% Let $\mathcal{K}$ be a set of partial $\Sigma$-algebras closed under the operators $\mathrm{I}$, $\mathrm{S}$, and $\mathrm{P}$. Then the functor $\mathbf{T}_{\Sigma,\boldsymbol{\mathcal{K}}} = \mathbf{T}_{\boldsymbol{\mathcal{K}}}\circ \mathbf{D}_{\Sigma}$ from $\mathsf{Set}^{S}$ to $\boldsymbol{\mathcal{K}}$ is a left adjoint of the functor $\mathrm{G}_{\Sigma}\circ \mathrm{In}_{\boldsymbol{\mathcal{K}}}$ from $\boldsymbol{\mathcal{K}}$ to $\mathsf{Set}^{S}$.
% \end{corollary}
%
% \begin{corollary}
% If $\mathcal{K}$ is a finitary $\mathrm{E}$-variety, a finitary $\mathrm{ECE}$-variety, or a finitary $\QE$-variety, then the inclusion functor $\mathrm{In}_{\boldsymbol{\mathcal{K}}}$ from $\boldsymbol{\mathcal{K}}$ to $\mathsf{PAlg}(\Sigma)$ has a left adjoint, and the functor $\mathrm{G}_{\Sigma}\circ \mathrm{In}_{\boldsymbol{\mathcal{K}}}$ from $\boldsymbol{\mathcal{K}}$ to $\mathsf{Set}^{S}$ has a left adjoint.
% \end{corollary}
%
% Let us conclude this section by recalling the following remark by Burmeister in~\cite{pb86}, on p.~68, on the higher complexity of the generation of algebraic structures in the case of partial algebras as compared to that of  total algebras: 
% \begin{quote}
% Especially for partial algebras, a combination of these two principles: stepwise generation of the elements and the structure on one side and stepwise enlarging of the congruence relations, form an important tool in connection with the generation of algebraic structures. In the theory of total algebras one usually already starts with a total ``free'' algebra and one only has to generate an appropriate congruence, while for partial algebras one additionally has to construct the elements and the admissible structure.
% \end{quote}
% This remark of Burmeister, as the reader will have the opportunity to verify, will be widely corroborated throughout this paper.
%
% %Notación
% %
% %$X{\scriptstyle\int}\int_{\beta}Q$
% %
% %$X\amalg\coprod_{\beta}Q$
% %
% %$\eta^{X{\scriptscriptstyle\int}\int_{\beta}Q}$
% %
% %${\scriptstyle\int}(X,\int_{\beta}Q)$
% %
% %$\eta^{{\scriptscriptstyle\int}(X,\int_{\beta}Q)}$
% %
% %$\eta^{X\amalg\coprod_{\beta}Q}$
%
% %Notación
%
% %Sea $\mathrm{in}^{\Sigma,\Sigma^{\mathrm{c},\mathcal{A}}}$ la inclusión canónica de $\Sigma$ en
% %$\Sigma^{\mathrm{c},\mathcal{A}}$. Para el morfismo $\mathbf{in}^{\Sigma,\Sigma^{\mathrm{c},\mathcal{A}}} = (\id^{S},\mathrm{in}^{\Sigma,\Sigma^{\mathrm{c},\mathcal{A}}})$ de $(S,\Sigma)$ en
% %$(S,\Sigma^{\mathrm{c},\mathcal{A}})$ y una partial $\Sigma^{\mathrm{c},\mathcal{A}}$-algebra $\mathbf{A}$ convenimos en denotar por $\mathbf{A}^{-\mathrm{c},\mathcal{A}}$ la partial $\Sigma$-algebra $(\mathbf{in}^{\Sigma,\Sigma^{\mathrm{c},\mathcal{A}}})^{\ast}(\mathbf{A})$.
%
% %%%%%%%%%%%%%%%%%%
% %%%%%%%%%%%%%%%%%%
%
% %\newpage
% %\subsection{A brief recall on free solutions to many-sorted partial algebras}\label{SABriefRecall}
%
% %Let us recall the construction of the free completion. For a given pair of many-sorted partial $\Sigma$-algebras $\mathbf{A}$ and $\mathbf{B}$ and for a $\Sigma$-homomorphism $f\colon\mathbf{A}\mor\mathbf{B}$, the free $\Sigma$-completion of $f$, denoted by $f^{\mathrm{fc}}$, is the unique $\Sigma$-algebra homomorphism $f^{\mathrm{fc}}\colon\mathbf{F}_{\Sigma}(\mathbf{A})\mor\mathbf{F}_{\Sigma}(\mathbf{A})$ making the following diagram commute.
% %
% %\begin{center}
% %\begin{tikzpicture}
% %[ACliment/.style={-{To [angle'=45, length=5.75pt, width=4pt, round]}
% %}, scale=0.8]
% %%,font=\scriptsize}]
% %\node[] (A) 	at (0,-1) 	[] {$\mathbf{A}$};
% %\node[] (B) 	at (0,-4) 	[] {$\mathbf{B}$};
% %\node[] (TA) 	at (10,2) 	[] {$\mathbf{T}_{\Sigma}(\mathbf{A})$};
% %\node[] (TB) 	at (10,-7) 	[] {$\mathbf{T}_{\Sigma}(\mathbf{B})$};
% %\node[] (TpA) 	at (7,.5) 	[] 
% %{$\overline{\mathbf{T}}_{\Sigma}(\mathbf{A})$};
% %\node[] (TpB) 	at (7,-5.5) 	[] 
% %{$\overline{\mathbf{T}}_{\Sigma}(\mathbf{B})$};
% %\node[] (FA) 	at (4,-1) 	[] {$\mathbf{F}_{\Sigma}(\mathbf{A})$};
% %\node[] (FB) 	at (4,-4) 	[] {$\mathbf{F}_{\Sigma}(\mathbf{B})$};
% %\draw[ACliment]  (A) 	to node [left]	{$f$} (B);
% %\draw[ACliment]  (TA) 	to node [right]	{$f^{@}$} (TB);
% %\draw[ACliment, bend left]  (A) 	to node [above]	{$\eta^{\mathbf{A}}$} (TA);
% %\draw[ACliment, bend right]  (B) 	to node [below]	{$\eta^{\mathbf{B}}$} (TB);
% %\draw[ACliment, bend left=20]  (A) 	to node [above]	{$\mathrm{in}^{\mathbf{A}}$} (TpA);
% %\draw[ACliment, bend right=20]  (B) 	to node [below]	{$\mathrm{in}^{\mathbf{B}}$} (TpB);
% %\draw[ACliment]  (TA) 	to node [above left]	{$(\mathrm{in}^{\mathbf{A}})^{\sharp}$} (TpA);
% %\draw[ACliment]  (TB) 	to node [below left]	{$(\mathrm{in}^{\mathbf{B}})^{\sharp}$} (TpB);
% %\draw[ACliment]  (A) 	to node [below]	{$\overline{\eta}^{\mathbf{A}}$} (FA);
% %\draw[ACliment]  (B) 	to node [above]	{$\overline{\eta}^{\mathbf{B}}$} (FB);
% %\draw[ACliment]  (FA) 	to node [right]	{$f^{\mathrm{fc}}$} (FB);
% %\draw[ACliment]  (FA) 	to node [above left]	{$\mathrm{in}^{\mathbf{F}_{\Sigma}(\mathbf{A})}$} (TpA);
% %\draw[ACliment]  (FB) 	to node [below left]	{$\mathrm{in}^{\mathbf{F}_{\Sigma}(\mathbf{B})}$} (TpB);
% %\draw[ACliment, bend left]  (TA) 	to node [below right]	{$(\overline{\eta}^{\mathbf{A}})^{\sharp}$} (FA);
% %\draw[ACliment, bend right]  (TB) 	to node [above right]	{$(\overline{\eta}^{\mathbf{B}})^{\sharp}$} (FB);
% %\end{tikzpicture}
% %\end{center}
% %
% %In case $X$ is an $S$-sorted set and we consider $\mathbf{D}_{\Sigma}(X)$, the discrete structure of many-sorted partial $\Sigma$-algebra on $X$, the three many-sorted $\Sigma$-algebras $\mathbf{F}_{\Sigma}(\mathbf{D}_{\Sigma}(X))$, $\overline{\mathbf{T}}_{\Sigma}(\mathbf{D}_{\Sigma}(X))$ and $\mathbf{T}_{\Sigma}(\mathbf{D}_{\Sigma}(X))$ are equal to $\mathbf{T}_{\Sigma}(X)$, that is, the free many-sorted $\Sigma$-algebra on $X$. Moreover, given any many-sorted partial $\Sigma$-algebra $\mathbf{B}$, every mapping $f$ from the set $X$ to the universe of $\mathbf{B}$ becomes a $\Sigma$-homomorphism $f\colon\mathbf{D}_{\Sigma}(X)\mor\mathbf{B}$. 
% %
% %The simplifications in this particular case translate to the above diagram retrieving the following one.
% %
% %\begin{center}
% %\begin{tikzpicture}
% %[ACliment/.style={-{To [angle'=45, length=5.75pt, width=4pt, round]}
% %}, scale=0.8]
% %%,font=\scriptsize}]
% %\node[] (A) 	at (0,-1) 	[] {$\mathbf{D}_{\Sigma}(X)$};
% %\node[] (B) 	at (0,-4) 	[] {$\mathbf{B}$};
% %\node[] (TA) 	at (4,-1) 	[] {$\mathbf{T}_{\Sigma}(X)$};
% %\node[] (TB) 	at (10,-7) 	[] {$\mathbf{T}_{\Sigma}(\mathbf{B})$};
% %\node[] (TpB) 	at (7,-5.5) 	[] 
% %{$\overline{\mathbf{T}}_{\Sigma}(\mathbf{B})$};
% %\node[] (FB) 	at (4,-4) 	[] {$\mathbf{F}_{\Sigma}(\mathbf{B})$};
% %\draw[ACliment]  (A) 	to node [left]	{$f$} (B);
% %\draw[ACliment, bend left=45]  (TA) 	to node [above right]	{$f^{@}$} (TB);
% %\draw[ACliment, bend right]  (B) 	to node [below]	{$\eta^{\mathbf{B}}$} (TB);
% %\draw[ACliment, bend right=20]  (B) 	to node [below]	{$\mathrm{in}^{\mathbf{B}}$} (TpB);
% %\draw[ACliment]  (TB) 	to node [below left]	{$(\mathrm{in}^{\mathbf{B}})^{\sharp}$} (TpB);
% %\draw[ACliment]  (A) 	to node [below]	{$\eta^{\mathbf{A}}$} (TA);
% %\draw[ACliment]  (B) 	to node [above]	{$\overline{\eta}^{\mathbf{B}}$} (FB);
% %\draw[ACliment]  (TA) 	to node [right]	{$f^{\mathrm{fc}}$} (FB);
% %\draw[ACliment]  (FB) 	to node [below left]	{$\mathrm{in}^{\mathbf{F}_{\Sigma}(\mathbf{B})}$} (TpB);
% %\draw[ACliment, bend right]  (TB) 	to node [above right]	{$(\overline{\eta}^{\mathbf{B}})^{\sharp}$} (FB);
% %\end{tikzpicture}
% %\end{center}
% %
% %Let us recall that, for this particular case $f^{\mathrm{fc}}=(\overline{\eta}^{\mathbf{B}}\circ f)^{\sharp}$.
%
% %Let $f$ be a $\Sigma$-homomorphism from a many-sorted partial $\Sigma$-algebra $\mathbf{A}$ to a many-sorted partial $\Sigma$-algebra $\mathbf{B}$. Then we can consider the $A$-generated relative subalgebra $\mathbf{Sch}(f)$ of the free $\Sigma$-completion $\mathbf{F}_{\Sigma}(\mathbf{A})$ and the closed $\Sigma$-homomorphism $f^{\mathrm{Sch}}\colon\mathbf{Sch}(f)\mor\mathbf{B}$. Let us remember that $f^{\mathrm{Sch}}$ is the unique closed $\Sigma$-homomorphism making the following diagram commute.
% %
% %\begin{center}
% %\begin{tikzpicture}
% %[ACliment/.style={-{To [angle'=45, length=5.75pt, width=4pt, round]}
% %}]
% %%,font=\scriptsize}]
% %\node[] (A) 	at 	(-3,1.5) 	[] 	{$\mathbf{A}$};
% %\node[] (B) 	at 	(0,-2) 	[] 	{$\mathbf{B}$};
% %\node[]	(S)		at 	(0,0)	[]	{$\mathbf{Sch}(f)$};
% %\node[] (FA)	at	(4,0)	[]	{$\mathbf{F}_{\Sigma}(\mathbf{A})$};
% %\node[] (FB)	at	(4,-2)	[]	{$\mathbf{F}_{\Sigma}(\mathbf{B})$};
% %\draw[ACliment, bend right]  (A) 	to node [below left]	{$f$} (B);
% %\draw[ACliment, bend left]  (A) 	to node [above right]	{$\overline{\eta}^{\mathbf{A}}$} (FA);
% %\draw[ACliment]  (A) 	to node [above right]	{$\mathrm{in}^{\mathbf{A},\mathbf{Sch}(f)}$} (S);
% %\draw[ACliment]  (B) 	to node [below]	{$\overline{\eta}^{\mathbf{B}}$} (FB);
% %\draw[ACliment]  (S) 	to node [above]	{$\mathrm{in}^{\mathbf{Sch}(f)}$} (FA);
% %\draw[ACliment]  (S) 	to node [right]	{$f^{\mathrm{Sch}}$} (B);
% %\draw[ACliment]  (FA) 	to node [right]	{$f^{\mathrm{fc}}$} (FB);
% %\end{tikzpicture}
% %\end{center}
% %
% %Let us recall that 
% %$$\mathbf{Sch}(f)=(f^{\mathrm{fc}})^{-1}[\overline{\eta}^{\mathbf{B}}[\mathbf{B}]].$$
%
% %In this paper, and for technical reasons, in addition to the usual (two-sorted) notion of category we will need to use a single-sorted definition of a category, which we recall below.
%
% %Dado un magma parcial $\mathbf{C} = (C,\circ)$ en el que $C$ es un conjunto y $\circ\colon C\times C\dmor C$ una aplicación parcial, un $u\in C$ es una \textbf{unidad} si, para cada $f\in C$, si $(u,f)\in \Dom(\circ)$, entonces $u\circ f = f$ y, para cada $g\in C$, si $(g,u)\in \Dom(\circ)$, entonces $g\circ u = g$.
%
% %Una categoría (definición de Ehresmann) es un par ordenado $\mathbf{C} = (C,\circ)$ en el que $C$ es un conjunto (de morfismos) y $\circ\colon C\times C\dmor C$ una aplicación parcial (composición de morfismos), tal que:
% %\begin{enumerate}
% %\item Para cada $f\in C$ existe una única unidad $\src(f)$ tal que $(f,\src(f))\in \Dom(\circ)$.
% %\item Para cada $f\in C$ existe una única unidad $\tgt(f)$ tal que $(\tgt(f),f)\in \Dom(\circ)$.
% %\item Para cada $f,g\in C$, si $(g,f)\in \Dom(\circ)$, entonces $\src(g) = \tgt(f)$, $\src(g\circ f) = \src(f)$, y $\tgt(g\circ f) = \tgt(g)$.
% %\item Para cualesquiera $f,g,h\in C$, si $(g,f)\in \Dom(\circ)$, $(h,g\circ f)\in \Dom(\circ)$, $(h,g)\in \Dom(\circ)$, y $(h\circ g, f)\in \Dom(\circ)$, entonces $h\circ (g\circ f) = (h\circ g)\circ f$.
% %\item Para cada $f,g\in C$, si $\src(g) = \tgt(f)$, entonces $(g,f)\in \Dom(\circ)$.
% %\end{enumerate}
%
% %Si $\mathbf{C}$ es una categoría, entonces las endoaplicaciones $\src$ y $\tgt$ de $C$ son tales que:
% %\begin{enumerate}
% %\item $\mathrm{Im}(\src) = \mathrm{Im}(\tgt) = C_{\mathrm{ob}}$.
% %\item Para cada $f\in C$, $\src(\src(f)) = \src(f)$, $\src(\tgt(f)) = \tgt(f)$, $\tgt(\src(f)) = \src(f)$, y  $\tgt(\tgt(f)) = \tgt(f)$, i.e., la bi-restricción de $\src$ a $C_{\mathrm{ob}}$ y $C_{\mathrm{ob}}$ es la identidad y la bi-restricción de $\tgt$ a $C_{\mathrm{ob}}$ y $C_{\mathrm{ob}}$ es la identidad.
% %\end{enumerate}
% %
% %Para cada $f\in C$ se cumple que $f\in C_{\mathrm{ob}}$ si y sólo si $f = \src(f)$ si y sólo si $f = \tgt(f)$.
%
% %\begin{definition}\label{DCatS}
% %A \emph{category} is an ordered quadruple $\mathbf{C} = (C,\circ,\src,\tgt)$ where $C$ is a set of \emph{morphisms}, $\circ\colon C\times C\dmor C$ a partial mapping of \emph{composition} of morphisms, and $\src$, $\tgt$ mappings from $C$ to $C$ which assign to every morphisms its \emph{source} and its \emph{target}, respectively, such that:
% %\begin{enumerate}
% %\item For every $f\in C$, $\src(\src(f)) = \src(f)$, $\src(\tgt(f)) = \tgt(f)$, $\tgt(\src(f)) = \src(f)$, and  $\tgt(\tgt(f)) = \tgt(f)$.
% %\item For every $f,g\in C$, $\src(g) = \tgt(f)$ if and only if $(g,f)\in \Dom(\circ)$.
% %\item For every $f,g\in C$, if $(g,f)\in \Dom(\circ)$, then $\src(g\circ f) = \src(f)$ and $\tgt(g\circ f) = \tgt(g)$.
% %\item For every $f\in C$, $f\circ \src(f) = f$ and $\tgt(f)\circ f = f$.
% %\item For every $f,g,h\in C$, if $(g,f)\in \Dom(\circ)$ and $(h,g\circ f)\in \Dom(\circ)$ or $(h,g)\in \Dom(\circ)$ and $(h\circ g, f)\in \Dom(\circ)$, then $h\circ (g\circ f) = (h\circ g)\circ f$.
% %\end{enumerate}
% %Then $f$ is an \emph{identity} if and only if $f = \src(f)$ or, equivalently, if and only if $f = \tgt(f)$.
% %\end{definition}
%
% %Las endoaplicaciones $\src$ y $\tgt$ de $C$ son idempotentes, $\mathrm{Im}(\src) = \mathrm{Im}(\tgt) = C_{\mathrm{ob}}$, y, para cada $f\in C$, se cumple que $f\in C_{\mathrm{ob}}$ si y sólo si $f = \src(f)$ si y sólo si $f = \tgt(f)$.
% %%For every $f\in C$, one should think of $\src(f)$ (resp. $\tgt(f)$) as if it were the identity morphism associated to the source of $f$ (resp. the identity morphism associated to the target of $f$).
%
% %For the above notion of category we have the following concept of functor.
% %
% %\begin{definition}
% %Let $\mathbf{C} = (C,\circ,\src,\tgt)$ and $\mathbf{C}' = (C',\circ',\src',\tgt')$ be two (single-sorted) categories. A functor from $\mathbf{C}$ to $\mathbf{C}'$ is a mapping $F$ from $C$ to $C'$ such that: 
% %\begin{enumerate}
% %\item For every $f\in C$, $F(\src(f)) = \src'(F(f))$.
% %\item For every $f\in C$, $F(\tgt(f)) = \tgt'(F(f))$.
% %\item For every $(g,f)\in \Dom(\circ)$, $F(g\circ f) = F(g)\circ'F(f)$.
% %\end{enumerate}
% %\end{definition}
% %
% %Añadir la definición de transformación natural.
% %
% %Lo que sigue es de Mac Lane~\cite{sM98} (y también se tiene una definición de $n$-category partiendo de un conjunto de $n$-cells y considerando a las $0$-cells, \ldots, $n-1$-cells como degenerate $n$-cells).
%
% %Similarly a $2$-category can be considered to be a single set $C$ considered as the set of $2$-cells. Then the previous $1$-cells (the arrows) and the $0$-cells (the objects) are just regarded as special ``degenerate'' $2$-cells. On the set $C$ of $2$-cells there are then two category structures, the ``horizontal'' structure $(\circ_{h},\mathrm{d}_{h},\mathrm{cd}_{h})$ and the ``vertical'' structure $(\circ_{v},\mathrm{d}_{v},\mathrm{cd}_{v})$. Each satisfies the axioms above for a category structure and in addition (1) every identity for the $0$-structure is an identity for the $1$-structure; (2) the two category structures commute with each other.
%
%
% %Una categoría es un par ordenado $\mathbf{C} = (C,\circ)$ en el que $C$ es un conjunto (de morfismos) y $\circ\colon C\times C\dmor C$ una aplicación parcial (composición de morfismos) tal que:
% %\begin{enumerate}
% %\item Para cualesquiera $f,g,h\in C$, the following conditions are equivalent:
% %\begin{enumerate}
% %\item $(h,g)\in \Dom(\circ)$ y $(g,f)\in \Dom(\circ)$,
% %\item $h\circ(g\circ f)\in \Dom(\circ)$, and
% %\item $(h\circ g)\circ f\in \Dom(\circ)$.
% %\end{enumerate}
% %\item Para cualesquiera $f,g,h\in C$, si $(h,g)\in \Dom(\circ)$ y $(g,f)\in \Dom(\circ)$, entonces $h\circ (g\circ f) = (h\circ g)\circ f$.
% %\item Para cada $f\in C$, existe una unidad $l$ en $\mathbf{C}$ y una unidad $r$ en $\mathbf{C}$ tales que $(l,f)\in \Dom(\circ)$ y $(f,r)\in \Dom(\circ)$.
% %\end{enumerate}
%
%
% %\subsection{Preliminaries on bicategories}
% %
% %This subsection is devoted to a description of general bicategial concepts needed for the study of many-sorted transductions. The material of this subsection is standard; it is recalled here for the convenience of the reader.
% %
% %Bicategories, lax functors ($\equiv$ morphisms) between bicategories.
% %
% %Lawvere, in~\cite{law96}, wrote: ``A bicategory is (apart from technical considerations concerning the precise sense in which composition is associative) a category enriched in Categories; that is, between any two objects there is (not just a class, but) a \emph{category} of morphisms (whose objects are called $1$-morphisms and whose morphisms are called $2$-morphisms of the bicategory) and composition is a \emph{functor} on the cartesian product category (for each triple of objects, or ``$0$-morphisms'', of the bicategory) \ldots Just as category theory can explicitly encapsulate much more mathematics than pure set theory, while yet remaining universal, so bicategories contain qualitatively more information than pure categories.''
% %
% %As usual, throughout this work, we will omit any of the indices appearing on the various symbols we introduce when this is unlikely to cause confusion.
% %
% %\begin{definition}
% %A \emph{bicategory} $\mathbf{C}$ consists of the following data:
% %\begin{enumerate}
% %\item A set $\mathrm{Ob}(\mathbf{C})$ called set of \emph{objects}, or $0$-\emph{cells} of $\mathbf{C}$.
% %\item For every $(A,B)\in \mathrm{Ob}(\mathbf{C})^{2}$, a category $\mathbf{C}(A,B)$. An object of $\mathbf{C}(A,B)$ is called a $1$-\emph{cell} or a \emph{morphism} of $\mathbf{C}$ and written $f\colon A\mor B$, while a morphism $\alpha\colon f\cel g$ of $\mathbf{C}(A,B)$ is called a $2$-\emph{cell} of $\mathbf{C}$. The composition sign $\circ$ of morphisms in $\mathbf{C}(A,B)$ will usually be omitted. For every object $f$ of $\mathbf{C}(A,B)$, the identity at $f$ is denoted by $\id^{f}\colon f\cel f$.
% %\item For every $(A,B,C)\in \mathrm{Ob}(\mathbf{C})^{3}$, a \emph{composition functor}:
% %$$
% %c_{A,B,C}\colon \mathbf{C}(B,C)\times \mathbf{C}(A,B)\mor \mathbf{C}(A,C).
% %$$
% % We write $g\circ f$ and $\beta\ast\alpha$ instead of $c_{A,B,C}(g,f)$ and $c_{A,B,C}(\beta,\alpha)$ for $(g,f)$ and $(\beta,\alpha)$ objects and morphisms of $\mathbf{C}(B,C)\times \mathbf{C}(A,B)$.
% %\item For every $A\in \mathrm{Ob}(\mathbf{C})$ a functor $\id^{A}$ from $\mathbf{1}$ to $\mathbf{C}(A,A)$, i.e., essentially, an object $\id^{A}$ of $\mathbf{C}(A,A)$ called \emph{identity morphism} of $A$. The identity morphism of $\id^{A}$ in $\mathbf{C}(A,A)$ is denoted by $\id^{\id^{A}}$ and is called \emph{identity $2$-cell} of $A$.
% %\item For every $(A,B,C,D)\in \mathrm{Ob}(\mathbf{C})^{4}$, a natural isomorphism $a_{A,B,C,D}$, called \emph{associativity isomorphism}, from the functor $c_{A,B,D}\circ (c_{B,C,D}\times \id^{\mathbf{C}(A,B)})$ to the functor $c_{A,C,D}\circ (\id^{\mathbf{C}(C,D)}\times c_{A,B,C})$. If $(h,g,f)$ is an object of $\mathbf{C}(C,D)\times \mathbf{C}(B,C)\times \mathbf{C}(A,B)$, then the isomorphism
% %    $$
% %    a_{A,B,C,D}(h,g,f)\colon (h\circ g)\circ f\cel h\circ(g\circ f)
% %    $$
% %    in $\mathbf{C}(A,D)$ is called the \emph{component} of $a_{A,B,C,D}$ at $(h,g,f)$.
% %\item For every $(A,B)\in \mathrm{Ob}(\mathbf{C})^{2}$, a natural isomorphism $r_{A,B}$, called \emph{right} identity, from the functor $c_{A,A,B}\circ (\id^{\mathbf{C}(A,B)}\times \id^{A})$ to the canonical isomorphism $\mathbf{C}(A,B)\times \mathbf{1}$ from $\mathbf{C}(A,B)\times \mathbf{1}$ to $\mathbf{C}(A,B)$. If $f$ is an object of $\mathbf{C}(A,B)$, then
% %    $$
% %    r_{A,B}(f)\colon f\circ \id^{A}\cel f.
% %    $$
% %\item For every $(A,B)\in \mathrm{Ob}(\mathbf{C})^{2}$, a natural isomorphism $l_{A,B}$, called \emph{left} identity, from the functor $c_{A,B,B}\circ (\id^{B}\times\id^{\mathbf{C}(A,B)})$ to the
% %    canonical isomorphism from $\mathbf{1}\times\mathbf{C}(A,B)$ to $\mathbf{C}(A,B)$. If $f$ is an object of $\mathbf{C}(A,B)$, then
% %    $$
% %    l_{A,B}(f)\colon \id^{B}\circ f\cel f.
% %    $$
% %\end{enumerate}
% %
% %The families of natural isomorphisms $a_{A,B,C,D}$, $r_{A,B}$, and $l_{A,B}$ are furthermore required to satisfy the following coherence axioms:
% %\begin{enumerate}
% %
% %\item[($\mathrm{AC}$)] Associativity coherence. If $(k,h,g,f)$ is an object of $\mathbf{C}(D,E)\times \mathbf{C}(C,D)\times \mathbf{C}(B,C)\times \mathbf{C}(A,B)$, then the following diagram commutes:
% %$$
% %\xymatrix@R=30pt@C=-17pt{
% %{} & ((k\circ h)\circ g)\circ f\ar@{=>}[dl]_-{a}\ar@{=>}[rr]^-{a\ast \id^{f}} & {} &
% %(k\circ(h\circ g))\circ f\ar@{=>}[rd]^-{a} & {}
% %\\
% %(k\circ h)\circ (g\circ f)\ar@{=>}[rrd]_-{a} & {} & {} & {} & k\circ((h\circ g)\circ f)\ar@{=>}[dll]^-*[l]{\id^{k}\ast a}
% %\\
% %{} & {} & k\circ(h\circ (g\circ f)) & {} & {}
% %}
% %$$
% %
% %\item[($\mathrm{IC}$)] Identity coherence. If $(g,f)$ is an object of $\mathbf{C}(B,C)\times \mathbf{C}(A,B)$, then the following diagram commutes:
% %$$
% %\xymatrix@C=0pt{
% %(g\circ \id^{B})\circ f\ar@{=>}[rr]^-{a}\ar@{=>}[rd]_-{r\ast \id^{f}} & {} &
% %g\circ (\id^{B}\circ f)\ar@{=>}[dl]^-{\id^{g}\ast l}
% %\\
% %{} & g\circ f & {}
% %}
% %$$
% %\end{enumerate}
% %\end{definition}
% %
% %Under the hypothesis that for every pair $(\gamma,\varphi)$ of coterminal mappings (= with the same codomain) we have chosen a pullback, we will denote by $\mathbf{Sp}(\mathsf{Set})$ the bicategory of \emph{spans} of $\mathsf{Set}$. The objects of $\mathbf{Sp}(\mathsf{Set})$ are the objects of $\mathsf{Set}$, i.e., the $\boldsymbol{\mathcal{U}}$-small sets. If $S$, and $T$ are two sets (in $\boldsymbol{\mathcal{U}}$), then the category $\mathbf{Sp}(\mathsf{Set})(S,T)$ has as objects, i.e., $1$-cells of $\mathbf{Sp}(\mathsf{Set})$, the ordered triples $(C,\beta,\gamma)$ \ldots A morphism in $\mathbf{Sp}(\mathsf{Set})(S,T)$ from $(C,\beta,\gamma)$ to $(C',\beta',\gamma')$, i.e., a $2$-cell of $\mathbf{Sp}(\mathsf{Set})$, is a mapping $\mu$ from $C$ to $C'$ such that \ldots Composition in $\mathbf{Sp}(\mathsf{Set})(S,T)$ is \ldots  The functors \ldots from \ldots to \ldots are defined by pullback. The identity morphism at $S$ is \ldots Finally $a$, $l$, and $r$ are given by the usual (and unique) isomorphisms of associativity and identity of pullbacks.
% %
% %The ``right'' notion of \emph{equality} between $1$-cells is that of $2$-isomorphism, or what is the same, equivalence (according to a general principle of category theory which claims that the only ``good'' notions are those invariant under equivalence).
% %
% %\begin{definition}
% %Let $\mathbf{C}$ and $\mathbf{C}'$ be two bicategories. A \emph{morphism} $(F,\varphi,\psi)$ from $\mathbf{C}$ to $\mathbf{C}'$ consists of the following data:
% %\begin{enumerate}
% %\item A mapping $F$ from $\mathrm{Ob}(\mathbf{C})$ to $\mathrm{Ob}(\mathbf{C}')$.
% %\item For every $(A,B)\in \mathrm{Ob}(\mathbf{C})^{2}$ a functor $F_{A,B}$ from $\mathbf{C}(A,B)$ to $\mathbf{C}'(F(A),F(B))$.
% %\item For every $(A,B,C)\in \mathrm{Ob}(\mathbf{C})^{3}$ a natural transformation
% %$$
% %\varphi_{A,B,C}\colon c'_{F(A),F(B),F(C)}\circ (F_{B,C}\times F_{A,B})\cel F_{A,C}\circ c_{A,B,C}.
% %$$
% %Thus, for every pair of objects $(g,f)$ of $\mathbf{C}(B,C)\times\mathbf{C}(A,B)$, we have that
% %$$
% %\varphi_{A,B,C}(g,f)\colon F(g)\circ F(f)\cel F(g\circ f).
% %$$
% %\item For every $A\in \mathrm{Ob}(\mathbf{C})$ a natural transformation
% %$$
% %\psi_{A}\colon\Id'_{F(A)}\cel F_{A,A}\circ \id^{A},
% %$$
% %i.e., essentially, a $2$-cell
% %$$
% %\psi_{A}\colon\Id'_{F(A)}\cel F(\id^{A}).
% %$$
% %\end{enumerate}
% %
% %These data are required to satisfy the following coherence axioms:
% %\begin{enumerate}
% %\item[($\mathrm{M.}0$)] If $(h,g,f)$ is an object of $\mathbf{C}(C,D)\times \mathbf{C}(B,C)\times \mathbf{C}(A,B)$, then the following diagram commutes:
% %$$
% %\xymatrix@C=50pt{
% %(F(h)\circ F(g))\circ F(f)\ar@{=>}[r]^-{\varphi\ast \id^{F(f)}}\ar@{=>}[d]_-{a'} & F(h\circ g)\circ F(f)\ar@{=>}[r]^-{\varphi} &
% %F((h\circ g)\circ f)\ar@{=>}[d]^-{F(a)}
% %\\
% %F(h)\circ (F(g)\circ F(f))\ar@{=>}[r]_-{\id^{F(h)}\ast \varphi} & F(h)\circ F( g\circ f)\ar@{=>}[r]_-{\varphi} &
% %F(h\circ (g\circ f))
% %}
% %$$
% %\item[($\mathrm{M.}1$)] If $f$ is an object of $\mathbf{C}(A,B)$, then the following diagrams commute:
% %$$
% %\xymatrix@C=50pt{
% %F(f)\circ \Id'_{F(A)}
% %  \ar@{=>}[r]^{\id^{F(f)}\ast\psi}
% %  \ar@{=>}[d]_{r'} &
% %F(f)\circ F(\id^{A})
% %  \ar@{=>}[r]^{\varphi} & F(f\circ \id^{A})
% %  \ar@{=>}[d]^{F(r)} \\
% %F(f)
% %  \ar@{=}[rr]_{} & {} &
% %F(f)
% %}
% %$$
% %
% %$$
% %\xymatrix@C=50pt{
% %\Id'_{F(B)}\circ F(f)
% %  \ar@{=>}[r]^{\psi\ast\id^{F(f)}}
% %  \ar@{=>}[d]_{l'} &
% %F(\id^{B})\circ F(f)
% %  \ar@{=>}[r]^{\varphi} & F(\id^{B}\circ f)
% %  \ar@{=>}[d]^{F(l)} \\
% %F(f)
% %  \ar@{=}[rr]_{} & {} &
% %F(f)
% %}
% %$$
% %\end{enumerate}
% %\end{definition}
% %
% %Hay que hablar de las fibred bicategories.
%
% %\subsection{Many-sorted transductions}
% %
% %As we will see, dropping the restriction of working with single-sorted algebras not only does not create problems (except the notational ones), but in fact brings to light new mathematical structures (paráfrasis de una frase de N. Gambino and A. Joyal in~\cite{GJ17}).
% %
% %The idea is always the same: To shed new light on the mathematics of the past.
% %
% %Dice P. Roquette, en la p.~1 de \cite{pr18}:
% %``Sometimes a result appears to be better understood if it
% %is generalized and freed from unnecessary assumptions, or if it is embedded into a
% %general theory which opens analogies to other fields of mathematics. Also, in order
% %to make further progress possible it is often convenient and sometimes necessary to
% %develop a framework, conceptual and notational, in which the known results become
% %trivial and almost self-evident, at least from a formalistic point of view.''
% %
% %In the following proposition, that is basic for a great deal of what
% %follows, for a mapping $\varphi$ from $S$ to $T$, we prove the
% %existence of an adjunction $\coprod_{\varphi}\ladj\Delta_{\varphi}$
% %from the category of $S$-sorted sets to the category of $T$-sorted
% %sets, as well as the existence of a contravariant functor
% %$\mathrm{MSet}$ and of a pseudo-functor $\mathrm{MSet}^{\ttcoprod}$
% %(related, respectively, to the right and left components of the
% %adjunction) from $\mathsf{Set}$ to $\mathbf{Cat}$.
% %
% %\begin{proposition}\label{Deltarightadjointcoproduct}
% %Let $\varphi\colon S\mor T$ be a mapping. Then there are functors
% %$\Delta_{\varphi}$ from $\mathsf{Set}^{T}$ to $\mathsf{Set}^{S}$
% %and $\coprod_{\varphi}$ from $\mathsf{Set}^{S}$ to $\mathsf{Set}^{T}$
% %such that $\coprod_{\varphi}\ladj\Delta_{\varphi}$. We write
% %$\theta^{\varphi}$, $\eta^{\varphi}$, and $\varepsilon^{\varphi}$,
% %respectively, for the natural isomorphism, the unit, and the counit
% %of the adjunction.
% %Besides, there exists a contravariant functor $\mathrm{MSet}$ from
% %$\mathsf{Set}$ to $\mathbf{Cat}$ which sends a set $S$ to the category
% %$\mathrm{MSet}(S) = \mathsf{Set}^{S}$, and a mapping $\varphi$ from
% %$S$ to $T$ to the functor $\Delta_{\varphi}$ from $\mathsf{Set}^{T}$
% %to $\mathsf{Set}^{S}$; and a pseudo-functor
% %$\mathrm{MSet}^{\ttcoprod}$ from $\mathsf{Set}$ to the
% %$2$\nobreakdash-category $\mathbf{Cat}$ given by the following data
% %\begin{enumerate}
% %\item The object mapping of $\mathrm{MSet}^{\ttcoprod}$ is that which
% %      sends a set $S$ to the category $\mathrm{MSet}^{\ttcoprod}(S) =
% %      \mathsf{Set}^{S}$.
% %
% %\item The morphism mapping of $\mathrm{MSet}^{\ttcoprod}$ is
% %      that which sends a mapping $\varphi$ from $S$ to $T$ to the
% %      functor $\mathrm{MSet}^{\ttcoprod}(\varphi) =
% %      \textstyle{\coprod}_{\varphi}$ from $\mathsf{Set}^{S}$ to
% %      $\mathsf{Set}^{T}$.
% %
% %\item For every $\varphi\colon S\mor T$ and $\psi\colon T\mor
% %      U$, the natural isomorphism $\gamma^{\varphi,\psi}$ from
% %      $\coprod_{\psi}\comp\coprod_{\varphi}$ to
% %      $\coprod_{\psi\comp\varphi}$ is that which is defined, for
% %      every $S$-sorted set $A$, as the $U$-sorted mapping from
% %      $\coprod_{\psi}(\coprod_{\varphi}A)$ to $\coprod_{\psi\comp\varphi}A$
% %      that in the $u$-th coordinate, with $u\in U$, is
% %      $((a,s),\varphi(s))\mapsto (a,s)$, if there
% %      exists an $s\in S$ such that $u=\psi(\varphi(s))$, and is the
% %      identity at $\vacio$, otherwise.
% %
% %\item For every set $S$, the natural isomorphism $\nu^{S}$ from
% %      $\Id_{\mathsf{Set}^{S}}$ to $\coprod_{\id_{S}}$ is that which is
% %      defined, for every $S$-sorted set $A$ and $s\in S$, as the
% %      canonical isomorphism from $A_{s}$ to $A_{s}\bprod\{s\}$.
% %      %(that sends $a$ in $A_{s}$ to the pair $(a,s)$ in $A_{s}\bprod\{s\}$).
% %\end{enumerate}
% %\end{proposition}
% %
% %\begin{proof}
% %Let $\Delta_{\varphi}$ be the functor from $\mathsf{Set}^{T}$ to
% %$\mathsf{Set}^{S}$ defined as follows: its object mapping sends each
% %$T$-sorted set $A$ to the $S$-sorted set $A_{\varphi} =
% %(A_{\varphi(s)})_{s\in S}$, i.e., the composite mapping $A\comp \varphi$;
% %its arrow mapping sends each $T$-sorted mapping $f\colon A\mor B$ to the
% %$S$-sorted mapping $f_{\varphi} =
% %(f_{\varphi(s)})_{s\in S}\colon A_{\varphi}\mor
% %B_{\varphi}$. Let $\coprod_{\varphi}$ be the functor from
% %$\mathsf{Set}^{S}$ to $\mathsf{Set}^{T}$ defined as follows:
% %its object mapping sends each $S$-sorted set $A$ to the
% %$T$-sorted set $\coprod_{\varphi}A =
% %(\coprod_{s\in\varphi^{-1}[t]}A_{s})_{t\in T}$; its arrow mapping
% %sends each $S$-sorted mapping
% %$f\colon A\mor B$ to the $T$-sorted mapping $\coprod_{\varphi}f =
% %(\coprod_{ s\in\varphi^{-1}[t]}f_{s})_{t\in T}\colon
% %\coprod_{\varphi}A\mor \coprod_{\varphi}B$. Then
% %the functor $\coprod_{\varphi}$ is a left adjoint for
% %$\Delta_{\varphi}$. Indeed, let $A$ be an $S$-sorted set, then
% %the pair $(\eta^{A}^{\varphi},\coprod_{\varphi}A)$, where
% %$\eta^{A}^{\varphi}$ is the $S$-sorted mapping from $A$ to
% %$\Delta_{\varphi}(\coprod_{\varphi}A) = (\tcoprod_{x\in
% %\varphi^{-1}[\varphi(s)]} A_{x} )_{s\in S}$ whose $s$-th coordinate,
% %for $s\in S$, is the canonical embedding from $A_{s}$ to
% %$\coprod_{x\in \varphi^{-1}[\varphi(s)]}A_{x}$, is a universal
% %morphism from $A$ to $\Delta_{\varphi}$.  This is so because, for a
% %$T$\nobreakdash-sorted set $B$ and an $S$-sorted mapping $f\colon
% %A\mor B_{\varphi}$, the $T$-sorted mapping $f^{\S}=(f^{\S}_{t})_{t\in
% %T}$ from $\coprod_{\varphi}A$ to $B$, where, for $t\in T$,
% %$f^{\S}_{t}$ is the unique mapping $[f_{s} ]_{s\in\varphi^{-1}[t]}$
% %from $\tcoprod_{s\in \varphi^{-1}[t] }A_{s}$ to $B_{t} =
% %B_{\varphi(s)}$ such that, for every $s\in \varphi^{-1}[t]$, the
% %following diagram commutes
% %$$
% %\xymatrix{ A_{s}
% %\ar[r]^-{\inc_{s}}
% %\ar[rd]_{f_{s}} &
% %\tcoprod_{s\in \varphi^{-1}[t] }A_{s}
% %\ar[d]^{[f_{s}]_{s\in\varphi^{-1}[t]} = f^{\S}_{t}} \\
% %     &
% %B_{t} = B_{\varphi(s)} }
% %$$
% %is such that $f = \Delta_{\varphi}(f^{\S})\comp \eta^{A}^{\varphi}$ and
% %unique with such a property.
% %
% %Otherwise stated, $\coprod_{\varphi}A$ is, simply,
% %$\mathrm{Lan}_{\varphi}A$, i.e., the left Kan extension of $A$ along
% %$\varphi$, recalling that every set is the set of objects of a
% %discrete category and every mapping between sets the object mapping of
% %a functor between discrete categories.
% %
% %Let us recall that, for each $S$-sorted set $X$ and each $T$-sorted set $A$,
% %$$
% %\textstyle\theta^{\varphi}_{X,A}\colon
% %\mathrm{Hom}_{\mathsf{Set}^{T}}(\coprod_{\varphi}X,A)\mor
% %\mathrm{Hom}_{\mathsf{Set}^{S}}(X,A_{\varphi}),
% %$$
% %the value of the natural isomorphism $\theta^{\varphi}$
% %at $(X,A)$, can be expressed in terms of $\eta_{X}^{\varphi}$,
% %for each $T$-sorted mapping $f\colon \coprod_{\varphi}X\mor A$,
% %as
% %$$
% %\theta^{\varphi}_{X,A}(f) =
% %f_{\varphi}\comp \eta_{X}^{\varphi}\colon X\mor A_{\varphi}.
% %$$
% %
% %To prove that $\mathrm{MSet}^{\ttcoprod}$ is a pseudo-functor, it is
% %enough to verify the coherence axioms. But given the situation
% %$$
% %\xymatrix{
% %S \ar[r]^{\varphi} &
% %T \ar[r]^{\psi} &
% %U \ar[r]^{\xi} &
% %X,
% %}
% %$$
% %the following diagrams commute
% %$$
% %\xymatrix@C=60pt{
% %\coprod_{\xi}\comp\coprod_{\psi}\comp\coprod_{\varphi}
% %  \ar[r]^{\id_{\coprod_{\xi}}\hcomp\gamma^{\varphi,\psi}}
% %  \ar[d]_{\gamma^{\psi,\xi}\hcomp\id_{\coprod_{\varphi}}} &
% %\coprod_{\xi}\comp\coprod_{\psi\comp\varphi}
% %  \ar[d]^{\gamma^{\psi\comp\varphi,\xi}} \\
% %\coprod_{\xi\comp\psi}\comp\coprod_{\varphi}
% %  \ar[r]_-{\gamma^{\varphi,\xi\comp\psi}} &
% %\coprod_{\xi\comp\psi\comp\varphi}
% %}
% %$$
% %
% %$$
% %\xymatrix{
% %\coprod_{\varphi}\comp\Id_{\mathsf{Set}^{S}}
% %  \ar[r]^{\id_{\coprod_{\varphi}}\hcomp\nu^{S}}
% %  \ar[d]_{\id_{\coprod_{\varphi}}} &
% %\coprod_{\varphi}\comp\coprod_{\id_{S}}
% %  \ar[d]^{\gamma^{\id_{S},\varphi}} \\
% %\coprod_{\varphi}
% %  \ar[r]_{\id_{\coprod_{\varphi}}} &
% %\coprod_{\varphi\comp\id_{S}}
% %}
% %\quad\quad
% %\xymatrix{
% %\Id_{\mathsf{Set}^{T}}\comp\coprod_{\varphi}
% %  \ar[r]^{\nu^{T}\hcomp\id_{\coprod_{\varphi}}}
% %  \ar[d]_{\id_{\coprod_{\varphi}}} &
% %\coprod_{\id_{T}}\comp\coprod_{\varphi}
% %  \ar[d]^{\gamma^{\varphi,\id_{T}}} \\
% %\coprod_{\varphi}
% %  \ar[r]_{\id_{\coprod_{\varphi}}} &
% %\coprod_{\id_{T}\comp\varphi}
% %}
% %$$
% %\end{proof}
% %
% %Henceforth, when dealing with a pseudo-functor we will restrict
% %ourselves to define explicitly only its object and morphism
% %mappings, if about the remaining data and conditions
% %there is not any doubt.
% %
% %Symbology.
% %
% %The empty word of any monoid will be denoted by $\lambda$.
% %
% %$S$, $T$, $U$, $C$, $D$ conjuntos de sorts.
% %
% %$X$: $S$-sorted set, $Y$: $T$-sorted set, $Z$: $U$-sorted set.
% %
% %$\Sigma$: $S$-sorted signature, $\Theta$: $T$-sorted signature, $\Omega$: $U$-sorted signature.
% %
% %$Q$: $C$-sorted set, $Q^{\boldsymbol{\cdot}}\subset Q$, $R$: $D$-sorted set, $R^{\boldsymbol{\cdot}}\subseteq R$.
% %
% %$\sig{\Sigma} = (S,\Sigma)$, $\sig{\Theta} = (T,\Theta)$, $\sig{\Xi} = (U,\Xi)$.
% %
% %$P, P', \ldots M, M', \ldots\in \mathrm{T}_{\Sigma}(X)_{s}$, $N, N', \ldots O, O', \ldots\in \mathrm{T}_{\Theta}(Y)_{t}$, \ldots
% %
% %
% %Spans: $(C,\beta,\gamma)\colon\xymatrix@1@C=19pt{S\ar[r]|(.44){\object@{|-}} & T}$, $(D,\varphi,\psi)\colon\xymatrix@1@C=19pt{T\ar[r]|(.44){\object@{|-}} & U}$.
% %$$
% %\xymatrix@C=20pt@R=10pt{
% %{} & C \ar[dl]_{\beta} \ar[dr]^-{\gamma} & {} \\
% %S & {} & T
% %}
% %$$
% %
% %$$
% %\xymatrix@C=20pt@R=10pt{
% %{} & C \ar[dl]_{\beta} \ar[dr]^-{\gamma} & {} & D \ar[dl]_{\varphi} \ar[dr]^-{\psi} & {} \\
% %S & {} & T & {} & U
% %}
% %$$
% %
% %Morfismo de un span en otro:
% %$$
% %\xymatrix@C=20pt@R=10pt{
% %{} & C \ar[dl]_{\beta} \ar[dr]^{\gamma}\ar@{=>}[dd]^{\mu} & {} \\
% %S & {} & T \\
% %{} & C' \ar[ul]^{\beta'} \ar[ur]_{\gamma'} & {}
% %}
% %$$
% %también denotado por $\mu\colon \xymatrix@1@C=19pt{(C,\beta,\gamma) \ar@{=>}[r] & (C',\beta',\gamma') }$.
% %
% %%$$
% %%\xymatrix@C=20pt@R=10pt{
% %%\ar[r]^a_b
% %%}
% %%$$
% %
% %Los objetos son: $(\sig{\Sigma},X) = (S,\Sigma,X)$, $(\sig{\Theta},Y) = (T,\Theta,Y)$, $(\sig{\Xi},Z) = (U,\Xi,Z)$. A tales pares ordenados los podríamos llamar \emph{(absolutely) free many-sorted presentations [without relations $\equiv$ relationless]} o \emph{lawless (many-sorted)  specifications}.
% %
% %Un morfismo de $(\sig{\Sigma},X)$ en $(\sig{\Theta},Y)$, también denominado un many-sorted term transducer de $(\sig{\Sigma},X)$ en $(\sig{\Theta},Y)$, es un triplo ordenado
% %$$
% %((\sig{\Sigma},X),((C,\beta,\gamma),Q,Q^{\boldsymbol{\cdot}},\mathcal{A}),(\sig{\Theta},Y)),
% %$$
% %denotado por $((C,\beta,\gamma),Q,Q^{\boldsymbol{\cdot}},\mathcal{A})\colon (\sig{\Sigma},X)\mor (\sig{\Theta},Y)$, o por $\boldsymbol{\mathcal{A}}\colon (\sig{\Sigma},X)\mor (\sig{\Theta},Y)$, en el que $(C,\beta,\gamma)$ es un span de $S$ en $T$, $Q$ un $C$-sorted set, el de los estados, $Q^{\boldsymbol{\cdot}}$ un subconjunto de $Q$, el de los estados finales, y $\mathcal{A}$, el $C$-sorted set de las producciones, un subconjunto finito del $C$-sorted set \ldots tal que \ldots; un morfismo de $(\sig{\Theta},Y)$ en $(\sig{\Omega},Z)$ es un triplo ordenado
% %$$
% %((\sig{\Theta},Y),((D,\varphi,\psi),R,R^{\boldsymbol{\cdot}},\mathcal{B}),(\sig{\Omega},Z)),
% %$$
% %denotado por $((D,\varphi,\psi),R,R^{\boldsymbol{\cdot}},\mathcal{B})\colon(\sig{\Theta},Y)\mor(\sig{\Omega},Z)$ o por
% %$\boldsymbol{\mathcal{B}}\colon (\sig{\Theta},Y)\mor(\sig{\Omega},Z)$.
% %
% %Si $\boldsymbol{\mathcal{A}}$ y $\boldsymbol{\mathcal{A}}'$ son morfismos de $(\sig{\Sigma},X)$ en $(\sig{\Theta},Y)$, entonces una $2$-célula del primero en el segundo es un par ordenado formado por un morfismo $\mu$ del span $(C,\beta,\gamma)$ en el span  $(C',\beta',\gamma')$ junto con un morfismo $g$ de $Q$ en $Q'_{\mu}$ o, lo que es equivalente, dado que $\coprod_{\mu}\dashv\Delta_{\mu}$, un morfismo $g^{\sharp}$ de $\coprod_{\mu}Q$ en $Q'$: $\dfrac{g\colon Q\mor Q'_{\mu}}{g^{\sharp}\colon \coprod_{\mu}Q\mor Q'}$. A partir de $g^{\sharp}\colon \coprod_{\mu}Q\mor Q'$, obtenemos un morfismo $\coprod_{\beta'}g^{\sharp}\colon \coprod_{\beta'}\coprod_{\mu}Q\mor \coprod_{\beta'}Q'$, i.e., un morfismo $\coprod_{\beta'}g^{\sharp}\colon \coprod_{\beta}Q\mor \coprod_{\beta'}Q'$, y un morfismo $\coprod_{\gamma'}g^{\sharp}\colon \coprod_{\gamma'}\coprod_{\mu}Q\mor \coprod_{\gamma'}Q'$, i.e., un morfismo $\coprod_{\gamma'}g^{\sharp}\colon \coprod_{\gamma}Q\mor \coprod_{\gamma'}Q'$. Y, de ellos, a su vez, un morfismo $\id^{X}\amalg\coprod_{\beta'}g^{\sharp}$, de $X\amalg\coprod_{\beta}Q$ en $X\amalg\coprod_{\beta'}Q'$ y un morfismo $\id^{Y}\amalg\coprod_{\gamma'}g^{\sharp}$, de $Y\amalg\coprod_{\gamma}Q$ en $Y\amalg\coprod_{\gamma'}Q'$. Por lo tanto tenemos el homomorfismo $(\id^{X}\amalg\coprod_{\beta'}g^{\sharp})^{@}$, de $\mathbf{T}_{\Sigma}(X\amalg\coprod_{\beta}Q)$ en $\mathbf{T}_{\Sigma}(X\amalg\coprod_{\beta'}Q')$ y el homomorfismo  $(\id^{Y}\amalg\coprod_{\gamma'}g^{\sharp})^{@}$, de $\mathbf{T}_{\Theta}(Y\amalg\coprod_{\gamma}Q)$ en $\mathbf{T}_{\Theta}(Y\amalg\coprod_{\gamma'}Q')$. De donde obtenemos the $C$-sorted mapping $((\id^{X}\amalg\coprod_{\beta'}g^{\sharp})^{@}_{\beta})\times g\times ((\id^{Y}\amalg\coprod_{\gamma'}g^{\sharp})^{@}_{\gamma})$ de $\mathrm{T}_{\Sigma}(X\amalg\coprod_{\beta}Q)_{\beta}\times Q\times \mathrm{T}_{\Theta}(Y\amalg\coprod_{\gamma}Q)_{\gamma}$ en $\mathrm{T}_{\Sigma}(X\amalg\coprod_{\beta'}Q')_{\beta}\times Q'_{\mu}\times \mathrm{T}_{\Theta}(Y\amalg\coprod_{\gamma'}Q')_{\gamma}$ (hay que tener en cuenta que $\beta = \beta'\circ\mu$, que $\gamma = \gamma'\circ\mu$, y que, por tener $\Delta_{\mu}$ un left adjoint, $\Delta_{\mu}$ preserves left limits, therefore, in particular, $\Delta_{\mu}$ preserves products) que nos permitiría transformar producciones de $\mathcal{A}$ en producciones de $\mathcal{A}'$. Falta añadir que hay compatibilidades entre $Q^{\boldsymbol{\cdot}}$ y $Q'^{\boldsymbol{\cdot}}$, así como entre $\mathcal{A}$ y $\mathcal{A}'$ (para esto último hay que tomar en cosideración que $\mathcal{A}'_{\mu}$ es una parte de $\mathrm{T}_{\Sigma}(X\amalg\coprod_{\beta'}Q')_{\beta}\times Q'_{\mu}\times \mathrm{T}_{\Theta}(Y\amalg\coprod_{\gamma'}Q')_{\gamma}$). Lo ``natural'' sería exigir:
% %\begin{enumerate}
% %\item $g[Q^{\boldsymbol{\cdot}}]\subseteq Q'^{{\boldsymbol{\cdot}}}_{\mu}$ \text{ y}
% %\item $\forall\,c\in C$, $\forall\,\mathfrak{p}\in \mathcal{A}_{c}$  $(((\id^{X}\amalg\coprod_{\beta'}g^{\sharp})^{@}_{\beta})\times g\times ((\id^{Y}\amalg\coprod_{\gamma'}g^{\sharp})^{@}_{\gamma})_{c}(\mathfrak{p})\in \mathcal{A}'_{\mu(c)})$.
% %\end{enumerate}
% %A estas $2$-células las podríamos llamar \emph{banales}. Estudiar entonces la relación que se establecería entre  $(C,\mathbf{C}(\boldsymbol{\mathcal{A}},c))_{c\in C}$ y $(C',\mathbf{C}(\boldsymbol{\mathcal{A}}',c'))_{c'\in C'}$ (para preservar o modificar la anterior definición). Enunciar una condición que sea suficiente para obtener una transformación.
% %
% %Denotamos a $((\id^{X}\amalg\coprod_{\beta'}g^{\sharp})^{@}_{\beta}$ por $g^{\beta', \beta}$ y a $(\id^{Y}\amalg\coprod_{\gamma'}g^{\sharp})^{@}_{\gamma}$ por $g^{\gamma', \gamma}$.
% %
% %
% %
% %
% %$$
% %\xymatrix@C=15ex@R=8ex{
% %(\sig{\Sigma},X)
% %\ar@/^15pt/[r]^{\boldsymbol{\mathcal{A}}}="f"
% %\ar@/^-15pt/[r]_{\boldsymbol{\mathcal{A}}'}="g" &
% %(\sig{\Theta},Y)
% %\ar @{} "f";"g" |{\dir{=>}}^{\,(\mu,g)}
% %}
% %$$
% %
% %$$
% %\xymatrix@C=20ex@R=8ex{
% %(\sig{\Sigma},X)
% %\ar@/^35pt/[r]^{\boldsymbol{\mathcal{A}}}="f"
% %  \ar[r]|{\boldsymbol{\mathcal{A}}'}="g"
% %  \ar@/^-37pt/[r]_{\boldsymbol{\mathcal{A}}''}="h" &
% %(\sig{\Theta},Y)
% %\ar @{} "f";"g" |{\dir{=>}}^{\,(\mu,g)}
% %\ar @{} "g";"h" |{\dir{=>}}^{\,(\mu',g')}
% %}
% %$$
% %
% %$$
% %\xymatrix@C=15ex@R=8ex{
% %(\sig{\Sigma},X)
% %\ar@/^15pt/[r]^{\boldsymbol{\mathcal{A}}}="f"
% %  \ar@/^-15pt/[r]_{\boldsymbol{\mathcal{A}}'}="g" &
% %(\sig{\Theta},Y)
% %\ar@/^15pt/[r]^{\boldsymbol{\mathcal{B}}}="f'"
% %  \ar@/^-15pt/[r]_{\boldsymbol{\mathcal{B}}'}="g'" &
% %(\sig{\Omega},Z)
% %\ar @{} "f";"g" |{\dir{=>}}^{\,(\mu,g)}
% %\ar @{} "f'";"g'" |{\dir{=>}}^{\,(\nu,h)}
% %}
% %$$
% %
% %Seguro que se obtiene una categoría y, posiblemente, una bicategoría.
% %
% %Creo que sería conveniente considerar a los many-sorted term transducers como la parte dinámica (dada, entre otros datos, por las producciones).
% %
% %A partir de $\boldsymbol{\mathcal{A}}\colon (\sig{\Sigma},X)\mor (\sig{\Theta},Y)$ se obtiene la bottom-up  many-sorted term transformation $\mathrm{BU}(\boldsymbol{\mathcal{A}})$ \ldots
% %
% %%\xymatrix@1@C=15pt{\ar@{=>}[r]^a_b&}
% %
% %Si $\Phi$ es una relación binaria en $A$ entonces $\Phi^{\ast}$, la clausura reflexiva y transitiva de $\Phi$, se define como: Para cada $(a,b)\in A^{2}$, $(a,b)\in \Phi^{\ast}$ si, y sólo si, $a = b$ o existe un $n\in \mathbb{N}-1$ y una $n+1$-familia $(x_{i})_{i\in n+1}\in A^{n+1}$ tal que $x_{0} = a$, $x_{n} = b$ y, para cada $i\in n$, $(x_{i},x_{i+1})\in \Phi$.
% %
% %Un span $(C,\beta,\gamma)$ puede ser tal que $C = S$ y que $\beta = \id^{S}$. Además, para un many-sorted term transducer de $(\sig{\Sigma},X)$ en $(\sig{\Theta},Y)$ cuyo span subyacente cumpla lo anterior, $Q$ y $Q^{\boldsymbol{\cdot}}$ pueden ser tales que \ldots (para obtener los tree homomorphisms).
% %$$
% %\xymatrix@C=20pt@R=10pt{
% %{} & S \ar[dl]_{\id^{S}} \ar[dr]^-{\gamma} & {} \\
% %S & {} & T
% %}
% %$$
% %
% %Desde la categoría (o bicategoría) $\mathbf{MTTrd}$, de los many-sorted term transducers (estamos nombrando a la categoría por sus morfismos, no por sus objetos), hasta la categoría (o bicategoría) $\mathbf{Rel}$ de las relaciones (estamos nombrando a la categoría por sus morfismos, no por sus objetos) se tiene un morfismo  que a $(\sig{\Sigma},X)$ le asigna $\mathrm{T}_{\Sigma}(X)$, a $\boldsymbol{\mathcal{A}}\colon (\sig{\Sigma},X)\mor (\sig{\Theta},Y)$ le asigna $\mathrm{FR}(\boldsymbol{\mathcal{A}})\subseteq \mathrm{T}_{\Sigma}(X)_{\beta}\times \mathrm{T}_{\Theta}(Y)_{\gamma}$ (realmente el span formado por $\mathrm{FR}(\boldsymbol{\mathcal{A}})$ y las proyecciones canónicas), y a una $2$-célula $(\mu,g)\colon\xymatrix@1@C=15pt{\boldsymbol{\mathcal{A}}\ar@{=>}[r]^{}& \boldsymbol{\mathcal{A}}'}$
% %%$(\mu,g)\colon \boldsymbol{\mathcal{A}} \cel \boldsymbol{\mathcal{A}}'$
% %la inclusión de $\mathrm{FR}(\boldsymbol{\mathcal{A}})$ en $\mathrm{FR}(\boldsymbol{\mathcal{A}}')$
% %$$
% %\xymatrix@C=15ex@R=8ex{
% %\mathrm{T}_{\Sigma}(X)
% %\ar@/^15pt/[r]^{\mathrm{FR}(\boldsymbol{\mathcal{A}})}="f"
% %  \ar@/^-15pt/[r]_{\mathrm{FR}(\boldsymbol{\mathcal{A}}')}="g" &
% %\mathrm{T}_{\Theta}(Y)
% %\ar @{} "f";"g" |{\dir{=>}}^{\,\mathrm{inc}}
% %}
% %$$
% %Hay que reconsiderar $\mathbf{Rel}$.
% %
% %Si $X$ es un $S$-sorted set e $Y$ un $T$-sorted set, entonces es natural considerar que una relación de $X$ en $Y$ es un subconjunto de $(X_{s}\times Y_{t})_{(s,t)\in S\times T}$. Si nos fijamos, aquí estamos considerando el span $(S\times T,\mathrm{pr}_{S},\mathrm{pr}_{T})\colon\xymatrix@1@C=19pt{S\ar[r]|(.44){\object@{|-}} & T}$,
% %%i.e.,
% %%$$
% %%\xymatrix@C=20pt@R=10pt{
% %%{} & S\times T \ar[dl]_{\mathrm{pr}_{S}} \ar[dr]^-{\mathrm{pr}_{T}} & {} \\
% %%S & {} & T,
% %%}
% %%$$
% %los $S\times T$-sorted sets $X_{\mathrm{pr}_{S}}$, $Y_{\mathrm{pr}_{T}}$, el producto cartesiano $X_{\mathrm{pr}_{S}}\times Y_{\mathrm{pr}_{T}} = (X_{s}\times Y_{t})_{(s,t)\in S\times T}$, y una parte de tal producto cartesiano. De modo que una relación de $X$ en $Y$ es un elemento de $\mathrm{Sub}(X_{\mathrm{pr}_{S}}\times Y_{\mathrm{pr}_{T}})$. Por otra parte, si $\Phi$ y $\Psi$ son relaciones de $X$ en $Y$, entonces de $\Phi$ en $\Psi$ hay una única $2$-célula, la inclusión canónica de $\Phi$ en $\Psi$, si $\Phi\subseteq\Psi$.
% %
% %En nuestra situación estamos considerando, para un par de conjuntos de tipos $(S,T)$, spans de tipo general $(C,\beta,\gamma)\colon\xymatrix@1@C=19pt{S\ar[r]|(.44){\object@{|-}} & T}$, con $\beta\colon C\mor S$ y $\gamma\colon C\mor T$. Entonces, dado un $S$-sorted set $X$ y un $T$-sorted set $Y$, consideramos los $C$-sorted sets $X_{\beta}$ y $Y_{\gamma}$, luego el producto cartesiano $X_{\beta}\times Y_{\gamma} = (X_{\beta(c)}\times Y_{\gamma(c)})_{c\in C}$, y, finalmente, subconjuntos de tal producto cartesiano. Tales subconjuntos serán las $1$-células de $X$ en $Y$. De modo que las $1$-células de $X$ en $Y$ son los elementos de $\bigcup_{(C,\beta,\gamma)\in\mathrm{Ob}(\mathbf{Sp}(\mathsf{Set})(S,T))}\mathrm{Sub}(X_{\beta}\times Y_{\gamma})$ (que contiene, en particular, a $\mathrm{Sub}(X\times Y) = \mathrm{Sub}(X_{\mathrm{pr}_{S}}\times Y_{\mathrm{pr}_{T}})$). Por otra parte, si $\mu$ es una $2$-célula del span $(C,\beta,\gamma)\colon\xymatrix@1@C=19pt{S\ar[r]|(.44){\object@{|-}} & T}$ en el span $(C',\beta',\gamma')\colon\xymatrix@1@C=19pt{S\ar[r]|(.44){\object@{|-}} & T}$, i.e., $\mu\colon \xymatrix@1@C=19pt{(C,\beta,\gamma) \ar@{=>}[r] & (C',\beta',\gamma') }$, $\Phi$ una $1$-célula, relativa a $(C,\beta,\gamma)$, de $X$ en $Y$, y $\Phi'$ una $1$-célula, relativa a $(C',\beta',\gamma')$, de $X$ en $Y$, entonces de  $\Phi$ en $\Phi'$ hay una única $2$-célula, la inclusión canónica de $\Phi$ en $\Phi'_{\mu}$, si $\Phi\subseteq\Phi'_{\mu}$ (así que de $\Phi$ en $\Phi'$ tendremos tantas $2$-células cuantas $2$-células $\mu$ tengamos del span  $(C,\beta,\gamma)\colon\xymatrix@1@C=19pt{S\ar[r]|(.44){\object@{|-}} & T}$ en el span $(C',\beta',\gamma')\colon\xymatrix@1@C=19pt{S\ar[r]|(.44){\object@{|-}} & T}$, siempre que se cumpla que $\Phi\subseteq\Phi'_{\mu}$).
% %
% %A la bicategoría correspondiente la podríamos llamar $\mathbf{MRel}_{\mathrm{sp}}$ (``$\mathbf{M}$'' por many-sorted, ``$\mathrm{sp}$'' porque su definición depende de los spans y las $2$-células entre spans).
% %
% %El conjunto de los objetos ($0$-cells) de $\mathbf{MRel}_{\mathrm{sp}}$ es $\bigcup_{S\in \boldsymbol{\mathcal{U}}}(\{S\}\times \mathrm{Hom}(S,\boldsymbol{\mathcal{U}}))$.
% %
% %Para cada par $((S,X),(T,Y))\in \mathrm{Ob}(\mathbf{MRel}_{\mathrm{sp}})^{2}$, el conjunto de los objetos ($1$-cells) de la categoría $\mathbf{MRel}_{\mathrm{sp}}((S,X),(T,Y))$ es:
% %$$
% %\textstyle
% %\bigcup_{(C,\beta,\gamma)\in\mathrm{Ob}(\mathbf{Sp}(\mathsf{Set})(S,T))}\mathrm{Sub}(X_{\beta}\times Y_{\gamma}),
% %$$
% %i.e., el conjunto formado por los subconjuntos $\Phi$ de $X_{\beta}\times Y_{\gamma}$, cuando $(C,\beta,\gamma)$ varía en $\mathrm{Ob}(\mathbf{Sp}(\mathsf{Set})(S,T))$.
% %
% %Sea $(\Phi,\Phi')$ un par de objetos de $\mathbf{MRel}_{\mathrm{sp}}((S,X),(T,Y))$, con $\Phi\subseteq X_{\beta}\times Y_{\gamma}$ y $\Phi'\subseteq X_{\beta'}\times Y_{\gamma'}$, siendo $(C,\beta,\gamma), (C',\beta',\gamma')\in\mathrm{Ob}(\mathbf{Sp}(\mathsf{Set})(S,T))$, y $\mu$ un morfismo de $(C,\beta,\gamma)$ en $(C',\beta',\gamma')$. Entonces, siendo $\mathrm{H}_{\mu}(\Phi,\Phi')$ el conjunto definido como:
% %$$
% %\mathrm{H}_{\mu}(\Phi,\Phi') =
% %    \begin{cases}
% %    \{\mathrm{in}_{\Phi,\Phi'_{\mu}}\}, &\text{si $\Phi\subseteq\Phi'_{\mu}$;} \\
% %    \varnothing,    &\text{si $\Phi\nsubseteq\Phi'_{\mu}$,}
% %    \end{cases}
% %$$
% %we define $\mathrm{Hom}_{\mathbf{MRel}_{\mathrm{sp}}((S,X),(T,Y))}(\Phi,\Phi')$, el conjunto de los morfismos de $\Phi$ en $\Phi'$, as:
% %$$
% %\textstyle
% %\bigcup_{\mu\in \mathrm{Hom}_{\mathbf{Sp}(\mathsf{Set})(S,T)}((C,\beta,\gamma), (C',\beta',\gamma'))}\mathrm{H}_{\mu}(\Phi,\Phi').
% %$$
% %
% %
% %Hay una proyección desde $\mathbf{MRel}_{\mathrm{sp}}$ hasta $\mathbf{Sp}(\mathsf{Set})$ (¿una fibración?).
% %
% %
% %Composición de las $1$-células. Preliminar:
% %$$
% %\xymatrix@C=20pt@R=10pt{
% %\textstyle
% %{} & {} & C\times_{T}D \ar[dl]_{\mathrm{p}_{C}} \ar[dr]^-{\mathrm{p}_{D}} & {} & {} \\
% %{} & C \ar[dl]_{\beta} \ar[dr]^-{\gamma} & {} & D \ar[dl]_{\varphi} \ar[dr]^-{\psi} & {} \\
% %S & {} & T & {} & U
% %}
% %$$
% %Si $\Phi$ es una $1$-célula, relativa a $(C,\beta,\gamma)$, de $X$ en $Y$, luego $\Phi\subseteq X_{\beta}\times Y_{\gamma}$, y $\Psi$ una $1$-célula, relativa a $(D,\varphi,\psi)$, de $Y$ en $Z$, luego $\Psi\subseteq Y_{\varphi}\times Z_{\psi}$, entonces $\Psi\circ \Phi$ es la $1$-célula, relativa a $(C\times_{T}D,\beta\circ \mathrm{p}_{C},\psi\circ \mathrm{p}_{D})$, de $X$ en $Z$ definida, para cada $(c,d)\in C\times_{T}D$, como $(\Psi\circ \Phi)_{c,d} = \Psi_{\psi(d)}\circ \Phi_{\beta(c)}$ (por ser $(c,d)\in C\times_{T}D$, se cumple que $\gamma(c) = \varphi(d)$, luego que $Y_{\gamma(c)} = Y_{\varphi(d)}$). De modo que $\Psi\circ \Phi\subseteq X_{\beta\circ \mathrm{p}_{C}}\times Z_{\psi\circ \mathrm{p}_{D}}$.
% %
% %Ahora falta definir la composición horizontal y vertical de las $2$-células, las unidades y verificar todos los axiomas.
% %
% %Composición vertical de las $2$-células. Preliminar:
% %$$
% %\xymatrix@C=35pt@R=25pt{
% %{} & C \ar[dl]_{\beta} \ar[dr]^{\gamma}\ar@{=>}[d]^{\mu} & {} \\
% %S & C' \ar[l]^{\beta'} \ar[r]_{\gamma'}\ar@{=>}[d]^{\mu'} & T \\
% %{} & C'' \ar[ul]^{\beta''} \ar[ur]_{\gamma''} & {}
% %}
% %$$
% %Si $\mu\colon \xymatrix@1@C=19pt{(C,\beta,\gamma) \ar@{=>}[r] & (C',\beta',\gamma')}$, $\mu'\colon \xymatrix@1@C=19pt{(C',\beta',\gamma') \ar@{=>}[r] & (C'',\beta'',\gamma'')}$, $\Phi$ es una $1$-célula, relativa a $(C,\beta,\gamma)$, de $X$ en $Y$, $\Phi'$ una $1$-célula, relativa a $(C',\beta',\gamma')$, de $X$ en $Y$, $\Phi''$ una $1$-célula, relativa a $(C'',\beta'',\gamma'')$, de $X$ en $Y$, y, además, $\Phi\subseteq\Phi'_{\mu}$ y $\Phi'\subseteq\Phi''_{\mu'}$, entonces de $\Phi$ en $\Phi''$ hay una única $2$-célula, la inclusión canónica de $\Phi$ en $\Phi''_{\mu'\circ \mu}$, que es la composición de la inclusión canónica de $\Phi$ en $\Phi'_{\mu}$ y de la inclusión canónica de $\Phi'_{\mu}$ en $\Phi''_{\mu'\circ\mu}$.
% %
% %Composición horizontal de las $2$-células. Preliminar:
% %$$
% %\xymatrix@C=20pt@R=10pt{
% %\textstyle
% %{} & {} & C\times_{T}D \ar[dl]_{\mathrm{p}_{C}} \ar[dr]^-{\mathrm{p}_{D}} & {} & {} \\
% %{} & C \ar[dl]_{\beta} \ar[dr]^-{\gamma}\ar@{=>}[dd]^{\mu} & {} & D \ar[dl]_{\varphi} \ar[dr]^-{\psi}\ar@{=>}[dd]^{\nu} & {} \\
% %S & {} & T & {} & U \\
% %{} & C' \ar[ul]^{\beta'} \ar[ur]_-{\gamma'} & {} & D' \ar[ul]^{\varphi'} \ar[ur]_-{\psi'} & {} \\
% %{} & {} & C'\times_{T}D' \ar[ul]^{\mathrm{p}_{C'}} \ar[ur]_-{\mathrm{p}_{D'}} & {} & {}
% %}
% %$$
% %De $C\times_{T}D$ en $C'$ tenemos el morfismo $\mu\circ \mathrm{p}_{C}$ y de $C\times_{T}D$ en $D'$ el morfismo $\nu\circ \mathrm{p}_{D}$. Ahora bien, puesto que $\gamma'\circ \mu\circ \mathrm{p}_{C} = \varphi'\circ \nu\circ \mathrm{p}_{D}$, entonces, en virtud de la propiedad universal del producto fibrado, existe un único morfismo $\nu\ast \mu$ de $C\times_{T}D$ en $C'\times_{T}D'$ tal que $\mathrm{p}_{C'} \circ (\nu\ast \mu) = \mu\circ \mathrm{p}_{C}$ y $\mathrm{p}_{D'} \circ (\nu\ast \mu) = \nu\circ \mathrm{p}_{D}$. Por lo tanto $(\beta'\circ\mathrm{p}_{C'})\circ(\nu\ast \mu) = \beta\circ \mathrm{p}_{C}$ y $(\psi'\circ\mathrm{p}_{D'})\circ(\nu\ast \mu) = \psi\circ \mathrm{p}_{D}$, i.e., el siguiente diagrama conmuta:
% %$$
% %\xymatrix@C=30pt@R=20pt{
% %{} & C\times_{T}D \ar[dl]_-{\beta\circ \mathrm{p}_{C}} \ar[dr]^-{\psi\circ\mathrm{p}_{D}}\ar@{=>}[dd]^{\nu\ast\mu} & {} \\
% %S & {} & U \\
% %{} & C'\times_{T}D' \ar[ul]^-{\beta'\circ \mathrm{p}_{C'}} \ar[ur]_-{\psi'\circ\mathrm{p}_{D'}} & {}
% %}
% %$$
% %
% %Identidades (Identities) para los objetos. Sea $(\sig{\Sigma},X)$ un objeto. Entonces es span subyacente de la identity at $(\sig{\Sigma},X)$, $\id^{(\sig{\Sigma},X)}\colon (\sig{\Sigma},X)\mor (\sig{\Sigma},X)$, is $(S,\id^{S},\id^{S})$, $Q = 1^{S}$, $Q^{\boldsymbol{\cdot}} = Q$
% %
% %Identidades horizontales. Preliminar:
% %
% %Dado $(C,\beta,\gamma)\colon\xymatrix@1@C=19pt{S\ar[r]|(.44){\object@{|-}} & T}$, se cumple que $(T,\id^{T},\id^{T})\colon\xymatrix@1@C=19pt{T\ar[r]|(.44){\object@{|-}} & T}$ es neutro por la izquierda y que $(S,\id^{S},\id^{S})\colon\xymatrix@1@C=19pt{S\ar[r]|(.44){\object@{|-}} & S}$ es neutro por la derecha para la composición horizontal de $2$-células.
% %
% %Identidades verticales. Preliminar:
% %
% %Si $\mu\colon \xymatrix@1@C=19pt{(C,\beta,\gamma) \ar@{=>}[r] & (C',\beta',\gamma') }$, entonces $\id^{C}\colon \xymatrix@1@C=19pt{(C,\beta,\gamma) \ar@{=>}[r] & (C,\beta,\gamma) }$ es neutro por la derecha y $\id^{C'}\colon \xymatrix@1@C=19pt{(C',\beta',\gamma') \ar@{=>}[r] & (C',\beta',\gamma') }$ es neutro por la izquierda para la composición vertical de $2$-células.
% %
% %Si $\boldsymbol{\mathcal{A}}$ y $\boldsymbol{\mathcal{A}}'$ son morfismos de $(\sig{\Sigma},X)$ en $(\sig{\Theta},Y)$ y $(\mu,g)$ es una $2$-célula del primero en el segundo (siendo $\mu$ un morfismo del span $(C,\beta,\gamma)$ en el span  $(C',\beta',\gamma')$ y $g$ un morfismo de $Q$ en $Q'_{\mu}$), entonces $(\id^{C},\id^{Q})$ es neutro por la derecha y $(\id^{C'},\id^{Q'})$ neutro por la izquierda para la composición vertical de $2$-células.
% %
% %En la bicategoría $\mathbf{MTTrd}$ (si es que lo es) podemos considerar, para los objetos las nociones de equivalencia y de adjunción. An adjunction between objects $(\sig{\Sigma},X)$ and $(\sig{\Theta},Y)$ consists of $1$-cells $\boldsymbol{\mathcal{G}}$ from $(\sig{\Sigma},X)$ to $(\sig{\Theta},Y)$, $\boldsymbol{\mathcal{F}}$ from $(\sig{\Theta},Y)$ to $(\sig{\Sigma},X)$, $2$-cells from the identity to $\boldsymbol{\mathcal{G}}\circ \boldsymbol{\mathcal{F}}$ and from $\boldsymbol{\mathcal{F}}\circ\boldsymbol{\mathcal{G}}$ to the identity satisfying the triangle equations. An equivalence between objects $(\sig{\Sigma},X)$ and $(\sig{\Theta},Y)$ consists of $1$-cells $\boldsymbol{\mathcal{G}}$ from $(\sig{\Sigma},X)$ to $(\sig{\Theta},Y)$, $\boldsymbol{\mathcal{F}}$ from $(\sig{\Theta},Y)$ to $(\sig{\Sigma},X)$, $2$-cells from the identity to $\boldsymbol{\mathcal{G}}\circ \boldsymbol{\mathcal{F}}$ and from $\boldsymbol{\mathcal{F}}\circ\boldsymbol{\mathcal{G}}$ to the identity which are invertible. Naturalmente, si $\boldsymbol{\mathcal{A}}$ es un morfismo (transducción) de $(\sig{\Sigma},X)$ en $(\sig{\Theta},Y)$, podremos plantear si tiene un adjunto por la izquierda o si es una equivalencia.
% %
% %Puesto que los morfismos entre bicategorías preservan adjunciones y equivalencias, en la bicategoría de relaciones tendríamos adjunciones y equivalencias inducidas.
% %
% %Classifying category of the bicategory $\mathbf{MTTrd}$ (Bénabou, Section 7.2, page 56). The objects are the objects of $\mathbf{MTTrd}$ and the arrows are isomorphism classes of $1$-cells of $\mathbf{MTTrd}$. Picard groupoid.
% %
% %Se tendrá que estudiar la relación entre los dos procedimientos: FR y RF (si ambos son lax-functors, ¿hay una lax natural transformation entre ambos?).
% %
% %
% %Si $\mathbf{A}$ y $\mathbf{B}$ son $\Sigma$-álgebras, entonces un meromorfismo de $\mathbf{A}$ en $\mathbf{B}$ es una subálgebra de $\mathbf{A}\times\mathbf{B}$. Si $\mathbf{T}_{\Theta}(Y)$ se pudiera convertir en una $\Sigma$-álgebra y $\mathbf{T}_{\Sigma}(X)$ en un álgebra, para una signatura relativa a $C$, entonces, a partir de $\mathrm{FR}(\boldsymbol{\mathcal{A}})$ se podría considerar la subálgebra generada por tal relación (el meromorfismo correspondiente) y estudiar si tal meromorfismo se puede obtener a partir de algún sistema de producción.
%
%
%
% %%%%%%%%%%%%%%%%%%%%%%%%%%%%%%%%%%%%%%%%%%%%%%%
\bibliographystyle{amsplain}
\bibliography{Biblio}

\end{document}
